% !TeX program = xelatex
% Chinese edition; build from the repository root.
\PassOptionsToPackage{colorlinks=true,linkcolor=blue,urlcolor=blue}{hyperref}



\documentclass[11pt]{article}

%\usepackage[top=1in, bottom=1in, left=1cm, right=1cm]{geometry}
%\usepackage{lipsum}
%\usepackage{background}
%\backgroundsetup{
%  position=current page.east,
%  angle=-90,
%  nodeanchor=east,
%  vshift=-5mm,
%  opacity=1,
%  scale=3,
%  contents=Confidential
%}

\usepackage{comment}
\usepackage{amsmath,amssymb,amsthm,bm}   % TeX4ht relies on these
\usepackage{mathtools}            % optional but harmless
\usepackage{rotating}

\usepackage[margin=1cm, footskip=5mm, paperwidth=8.5in, paperheight=11in]{geometry}
\usepackage[
  hypertexnames=false
]{hyperref}
\usepackage[UTF8,scheme=plain,fontset=fandol]{ctex}


\usepackage{mathtools}
\usepackage{amsfonts}
\usepackage{graphicx}
\usepackage{amsthm}

\usepackage{afterpage}
\usepackage{adjustbox}
\usepackage{placeins}

\usepackage{supertabular} 
\usepackage{titlesec}
\usepackage{array}
\usepackage{tabularx}
\usepackage{color}
\usepackage{algorithm}
\usepackage{algpseudocode}
\usepackage{subcaption}
\usepackage{hyperref}
\usepackage{pgfplots}
\usepackage{relsize}
\let\labelindent\relax
\usepackage{enumitem}
\usepackage{fancyhdr}
\usepackage{alltt}
\usepackage{soul}
\usepackage{fancyvrb}
\usepackage{xcolor}
%\usepackage[english,greek]{babel}
%\usepackage{ucs} 
%\usepackage[utf8x]{inputenc}
%\usepackage[usenames,dvipsnames]{xcolor}
%\usepackage{tikz}
%\usepackage{tkz-tab}
%\usepackage{caption}
%\usepackage{latexsym}
%\usepackage{amssymb}
\usepackage[margin=0.1cm]{subcaption}
\usepackage{multicol}
\usepackage{refcount}

% XeLaTeX uses native hyperref URL line breaking. % Break URL
\usepackage{multirow}
\usepackage[most]{tcolorbox}
%\usepackage[breakable,skins]{tcolorbox}

\tcbset{breakable}
\usepackage{esvect}
\usepackage{hyperref}
\usepackage{mathdots}
\usepackage{pifont} 
\usepackage{booktabs}
\usepackage{arydshln}
\usepackage{xcolor}
\def\UrlBreaks{\do\/\do-\do\_}
\usepackage{mathtools}  
\usepackage{array}
\usepackage{hyperref}
\usepackage{hyperref}
\usepackage{mathtools, nccmath}
%\usepackage[style=numeric,sorting=none]{biblatex}
%\usepackage[lite]{mtpro2}
\usepackage{titling}
\renewcommand\maketitlehooka{\null\mbox{}\vfill}
\renewcommand\maketitlehookd{\vfill\null}

\DeclarePairedDelimiter{\nint}\lfloor\rceil

\DeclarePairedDelimiter{\abs}\lvert\rvert
\usepackage{tocloft}
\usepackage{pdflscape}
\usepackage[all=normal, floats, bibnotes, wordspacing, charwidths, indent]{savetrees}
\makeatletter%
\captionsetup{belowskip=0pt}

\usepackage{enumitem}
\setlist{nolistsep}
%\setlist[itemize]{leftmargin=0.1}
\setlist[itemize]{leftmargin=*}
\setlist[enumerate]{leftmargin=*}

% align figures to the top margin
\makeatletter
\setlength{\@fptop}{0pt}

\newcommand\bunderline{}% check for being undefined
\DeclareRobustCommand\bunderline[1]{\mathord{\mathpalette\b@underline{#1}}}
\newcommand{\b@underline}[2]{%
  \sbox\z@{$\m@th#1\underline{#2}$}%
  \raisebox{-\dp\z@}{\scalebox{.5}[.25]{$\m@th#1[$}}%
  \copy\z@
  \raisebox{-\dp\z@}{\scalebox{.5}[.25]{$\m@th#1]$}}%
}


\makeatother

%\renewcommand{\figureautorefname}{foo}
%\renewcommand{\tableautorefname}{bar}

%\usepackage[all=normal, floats, bibnotes, wordspacing, charwidths, indent, lists]

%\captionsetup{belowskip=0pt}

\usepackage{graphicx}
\def\rotatecharone#1{\rotatebox[origin=c]{180}{#1}}

\def\rotatechartwo#1{\reflectbox{#1}}
\usepackage{hyperref}
\skip\footins 0.5\baselineskip %8.1pt plus 4pt minus 2pt
\floatsep 0.5\baselineskip
\textfloatsep 0.5\baselineskip
\intextsep 0.5\baselineskip 
\dbltextfloatsep 0.5\baselineskip  
\dblfloatsep 0.5\baselineskip 

%\raggedbottom

\belowcaptionskip 0pt

% Let TeX choose natural paragraph line counts for Chinese text.

\setlength{\abovecaptionskip}{0pt} % 

\newenvironment{myproof}[1][\proofname]{%
  \begin{proof}[#1]$ $\par\nobreak\ignorespaces
}{%
  \end{proof}
}
\makeatletter
\renewcommand{\sectionautorefname}{\S\@gobble}
\renewcommand{\subsectionautorefname}{\S\@gobble}
\renewcommand{\subsubsectionautorefname}{\S\@gobble}
\renewcommand{\appendixautorefname}{\S\@gobble}
\renewcommand{\appendixautorefname}{\S\@gobble}

\newcommand\mymathcolor[2]{{\color{#1}#2}}

\makeatother
\newcommand{\ignore}[1]{}

% Add additional comma in the math mode
%\DeclareMathSymbol{,}{\mathpunct}{letters}{"2C}
%\newcommand{\comma}{\mathpunct{\mkern\thinmuskip}}
%\mathcode`,=\string"8000
%{\catcode`,=\active \gdef,{\comma}}


\addtolength{\cftsecnumwidth}{15pt}
\addtolength{\cftsubsecnumwidth}{10pt}
\addtolength{\cftsubsubsecnumwidth}{12pt}

\makeatletter%
\usepackage{stackengine,xcolor}
% The unused pdfTeX-only outline macros are omitted in this XeLaTeX entry.
\usepackage{lmodern}
\usepackage{lipsum}


\makeatletter
\DeclareRobustCommand{\ccong}{\mathrel{\mathpalette\@verequiv\sim}}
\newcommand{\@verequiv}[2]{%
  \lower.5\p@\vbox{
    \lineskiplimit\maxdimen
    \lineskip-.5\p@
    \ialign{%
      $\m@th#1\hfil##\hfil$\crcr
      #2\crcr
      \equiv\crcr
    }%
  }%
}
\renewcommand*\env@matrix[1][c]{\hskip -\arraycolsep
  \let\@ifnextchar\new@ifnextchar
  \array{*\c@MaxMatrixCols #1}}

\newcommand{\hathat}{\tilde} 

%\newcommand{\hathat}[1]{% 
%\begingroup%
%  \let\macc@kerna\z@%
%  \let\macc@kernb\z@%
%  \let\macc@nucleus\@empty%
%  \hat{\raisebox{.2ex}{\vphantom{\ensuremath{#1}}}\smash{\hat{#1}}}%
%\endgroup%
%}

\newcommand{\revise}[1]{{{\textcolor{black}{#1}}}}

\makeatother

% \title{**\vspace*{\fill}**The Beginner's Textbook for Fully Homomorphic Encryption}

% \date{January 1, 2025}

% \author{Ronny Ko**\vspace*{\fill}**}


\hypersetup{breaklinks=true,unicode=true,pdftitle={全同态加密入门教材（中文译本）},pdfauthor={Ronny Ko},pdflang={zh-CN}}
\renewcommand{\contentsname}{目录}
\renewcommand{\refname}{参考文献}
\renewcommand{\figurename}{图}
\renewcommand{\tablename}{表}
\renewcommand{\proofname}{证明}
\renewcommand{\figureautorefname}{图}
\renewcommand{\tableautorefname}{表}
\renewcommand{\partautorefname}{部分}
\setlength{\emergencystretch}{2em}
\begin{document}
\title{\Huge{\textbf{全同态加密入门教材}}\\[0.6em]\Large{中文译本}\\[0.4em]\large{The Beginner\textquotesingle s Textbook for Fully Homomorphic Encryption}}
\author{
        \textbf{Ronny Ko} \thanks{\large \textit{致谢：}\\
        \textbf{Robin Geelen} (KU Leuven)\\  
        \textbf{Tianjian Yang}（北京大学）\\ \textbf{Yongwoo Lee}（仁荷大学）\\ \textbf{Nolan Carouge} (Grenoble INP Ensimag) \\
        \textbf{Navid Abapour}（萨里大学）
        }
        \\{LG Electronics Inc.}}%\\\texttt{\small{hajoon.ko@lge.com}}}
\date{\small 中文初译对应英文源码版本：\texttt{2a901038078f}}



\begin{titlingpage}
\maketitle
\end{titlingpage}


\clearpage

\section*{前言}
全同态加密（Fully Homomorphic Encryption，FHE）是一类允许直接对加密数据进行计算的密码方案，操作方式如同处理明文数据。对加密数据完成全部计算后，解密即可得到结果；正确解密要求所选参数和被求值电路满足该方案的噪声界及失败概率界。精确方案在指定的有限明文环中计算；CKKS 返回需要跟踪误差的近似结果，并非精确的实数计算结果。

FHE 支持加密数值的加法、乘法等基本运算。基于这些基本运算，可以构造更复杂的计算，包括减法、除法、逻辑门（例如 AND、OR、XOR、NAND、MUX），以及 ReLU、sigmoid 和三角函数（例如 sin、cos）等更复杂的数学函数。运算的含义及实现方式取决于明文定义域和允许的输入范围。例如，有限环中的除法要求分母可逆；实数倒数和激活函数通常需要有界区间近似或合适的离散电路。

FHE 可以支持隐私保护机器学习，使服务器通过机器学习模型处理客户端的加密数据。在安全假设和正确设计的协议下，不持有解密密钥的求值方无法从密文中恢复加密输入或输出。协议元数据、主动公开的结果和访问模式需要另外分析。
FHE 也可帮助保护智能合约数据的机密性。FHE 本身不证明求值方正确执行了指定计算；完整性和可验证性需要额外的协议机制。
FHE 的其他应用还包括安全外包数据分析、加密数据库查询、隐私保护搜索以及面向数字签名的高效多方计算等。

本书旨在帮助读者从数学层面理解 FHE 的工作原理，共分为以下四个部分：

$ $

\begin{itemize}
\item \textbf{\autoref{part:basic-math}：\nameref{part:basic-math}}介绍 FHE 所需的数学背景，包括群、域、阶、多项式环、分圆多项式、向量与矩阵、中国剩余定理、泰勒级数、多项式插值和快速傅里叶变换。

\item \textbf{\autoref{part:pqc}：\nameref{part:pqc}}介绍常见的格密码体制，包括 LWE、RLWE、GLWE、GLev 和 GGSW。

\item \textbf{\textbf{\autoref{part:generic-fhe}：\nameref{part:generic-fhe}}}介绍多种现有方案共同使用的 FHE 通用技术，例如同态加法、同态乘法、模数切换和密钥切换。


\item \textbf{\textbf{\autoref{part:fhe-schemes}：\nameref{part:fhe-schemes}}}介绍四种广泛使用的 FHE 方案：TFHE、BFV、CKKS 和 BGV，以及它们的 RNS 变体。
\end{itemize}

$ $

%These parts are designed in an incremental manner, and therefore understanding each part requires the understanding of its prior part(s). 

%$ $

本书英文原版提供\href{https://arxiv.org/abs/2503.05136}{\textbf{arXiv PDF}}和\href{https://fhetextbook.github.io}{\textbf{自动生成的动态网站}}（由\href{https://www.kodymirus.cz/overleaf-html-sample/main.html}{make4ht}生成）。作者还提供用于教学的 Python FHE 演示库，涵盖 TFHE、BFV、BGV 和 CKKS，可在\href{https://github.com/fhetextbook/fhetextbook.github.io}{\texttt{https://github.com/fhetextbook/fhetextbook.github.io}}获取。
英文原稿中的问题可在上游\href{https://github.com/fhetextbook/fhetextbook.github.io/issues}{\textbf{问题跟踪页面}}反馈。
\paragraph*{中文译本说明} 本译本保留原书的章节结构、公式、交叉引用和参考文献。图中已嵌入的英文文字及代码标识保留原样。本分支在完整初译后进行了逐项学术审校，已确认的勘误同步应用于中英文稿。审校范围、逐项依据、反例和剩余限制记录于 \texttt{audit/ACADEMIC\_REVIEW.md}；这不等同于完整的形式化验证或独立安全认证。



\subsubsection*{致谢}
特别感谢以下研究者：Robin Geelen（KU Leuven，\href{mailto:robin.geelen@esat.kuleuven.be}{robin.geelen@esat.kuleuven.be}）提供细致而投入的反馈；Yongwoo Lee（仁荷大学，\href{mailto:yongwoo@inha.ac.kr}{yongwoo@inha.ac.kr}）提供总体建议；Tianjian Yang（北京大学，\href{mailto:2300012738@stu.pku.edu.cn}{2300012738@stu.pku.edu.cn}）纠正全文中的大量笔误和其他错误；Nolan Carouge（Grenoble INP Ensimag，\href{mailto:nolan.carouge.pro@gmail.com}{nolan.carouge.pro@gmail.com}）校对第一部分；Navid Abapour（萨里大学，\href{mailto:n.abapour@surrey.ac.uk}{n.abapour@surrey.ac.uk}）校对第三和第四部分。

\thispagestyle{empty}

\newpage

\tableofcontents



%\newcommand{\subparagraph}{} % /usepackage{titlesec}
\titleformat*{\section}{\LARGE\bfseries}
\titleformat*{\subsection}{\Large\bfseries}
\titleformat*{\subsubsection}{\bfseries}
\titleformat*{\paragraph}{\normalsize\bfseries}
\titleformat*{\subparagraph}{\large\bfseries}

% page header foot 
%\usepackage{fancyhdr}
%\pagestyle{fancy}
%\lhead{Security and Privacy in Cyber-Physical Systems: Foundations and Applications}
%\rfoot{Copyright \textcopyright 2016 by Wiley}
% \thispagestyle{fancy}, after \maketitle
\newcommand{\para}[1]{\vspace{0.05in}\noindent{\bf{#1}}}


\newcommand{\white}[1]{{\textcolor{white}{#1}}} % phantom upgrade

\newcommand{\gap}[1]{\text{ } \text{#1} \text{ }}

\newcommand{\tboxlabel}[1]{$\bm{\langle}$总结~{#1}$\bm{\rangle}$}

\newcommand{\tboxdef}[1]{$\bm{\langle}$定义~{#1}$\bm{\rangle}$}

\newcommand{\tboxtheorem}[1]{$\bm{\langle}$定理~{#1}$\bm{\rangle}$}

\newtheorem{proposition}{命题}



\newtcolorbox[blend into=tables]{mytable}[2][]{float=htb, halign=center,  title={#2}, every float=\centering, #1}


% $\hat Y = \frac{1}{1 + e^{-Z}}$.

% $Z = {w_1 \cdot X_1 + w_2 \cdot X_2 + \dots + w_n \cdot X_n + b}$



\clearpage

%\section{Background}


\part{数学基础}
\label{part:basic-math}

\renewcommand{\thesection}{A-\arabic{section}}
\setcounter{section}{0}

本部分介绍数论的基本数学概念：群、域、阶、单位根、分圆多项式、多项式环和分解。这些概念是后量子密码的重要基础。

\clearpage

\section{模算术}
\label{sec:modulo}

\noindent \textbf{- 参考资料}

\href{https://www.youtube.com/watch?v=fz1vxq5ts5I}{YouTube -- 扩展欧几里得算法教程}。


\subsection{概述}
\label{subsec:modulo-overview} 

\begin{tcolorbox}[title={\textbf{\tboxdef{\ref*{subsec:modulo-overview}} 整数取模}}]

\begin{itemize}

\item \textbf{取模（modulo）}是计算一个数除以另一个数所得余数的运算。\textbf{modulo} 常缩写为 \textbf{mod}。

$ $

\item \textbf{$\bm{a}$ mod $\bm{q}$（即 $\bm{a} \bm{\text{ modulo } q}$）}是 $a$ 除以 $q$ 所得的余数，其值总是 $\{0, 1, 2, 3, \cdots, q-1\}$ 中的一个元素。例如，$7 \bmod 5 = 2$，因为 7 除以 5 的余数为 2。

$ $

\item \textbf{模数：}对于 $\bm{a}$ mod $\bm{q}$，我们将除数 $q$ 称为模数，而 \emph{取模}指的是相应的运算。

$ $

\item \textbf{模同余（$\bm{\equiv}$）：}若 $a$ 与 $b$ 除以 $q$ 后的余数相同，则称它们模 $q$ 同余（即 $a \equiv b \textbf{ mod } q$）。例如，$5 \equiv 12 \bmod 7$，因为 $5 \bmod 7 = 5$ 且 $12 \bmod 7 = 5$。在数学中，记号 $a \equiv b \bmod q$ 与 $a = b \pmod q$ 的含义相同，都表示 $a$ 除以 $q$ 的余数与 $b$ 除以 $q$ 的余数相同。注意，这种记号不同于 $a = b \bmod q$，后者表示 $a$ 等于 $b$ 除以 $q$ 的余数。

$ $

\item \textbf{同余与相等：}

$a \equiv b \bmod q \Longleftrightarrow a = b + k\cdot q$ \text{ }（其中 $k$ 为某个整数）

$ $

这意味着，$a$ 与 $b$ 模 $q$ 同余，当且仅当 $a$ 与 $b$ 之差是 $q$ 的某个整数倍。例如，$5 \equiv 12 \bmod 7 \Longleftrightarrow 5 = 12 + (-1)\cdot 7$

\end{itemize}

\end{tcolorbox}


\subsection{模运算}
\label{subsec:modulo-arithmetic}

模运算支持加法、减法和乘法。这些运算具有如下性质：

\begin{tcolorbox}[title={\textbf{\tboxtheorem{\ref*{subsec:modulo-arithmetic}.1} 模运算的性质}}]
对于任意整数 $x$，有：

\begin{enumerate}
\item \textbf{加法：}$a \equiv b \bmod q \Longleftrightarrow a + x\equiv b + x\bmod q$

\item \textbf{减法：}$a \equiv b \bmod q \Longleftrightarrow a - x\equiv b - x\bmod q$

\item \textbf{乘法：}$a \equiv b \bmod q \Longleftrightarrow a \cdot x \equiv b \cdot x \bmod q$。
上述等价关系成立的前提是 $\gcd(x,q)=1$。如果不满足这一假设，则只能保证蕴含关系 $a \equiv b \bmod q \Rightarrow a \cdot x \equiv b \cdot x \bmod q$ 成立。

\end{enumerate}

\end{tcolorbox}

\begin{proof}

$ $

对于任意整数 $x$，

\begin{enumerate}


\item \textbf{加法：}$a \equiv b \bmod q \Longleftrightarrow a = b + k q $（其中 $k$ 为某个整数）\textcolor{red}{ \text{ } \# $a$ 与 $b$ 之差是 $q$ 的某个整数倍}

$\Longleftrightarrow a + x = b + k\cdot q + x$

$\Longleftrightarrow a + x = b + x + k\cdot q$ \textcolor{red}{ $\rhd$ $a+x$ 与 $b+x$ 之差是 $q$ 的某个整数倍}

$\Longleftrightarrow a + x \equiv b + x \bmod q$

$ $

\item \textbf{减法：}$a \equiv b \bmod q \Longleftrightarrow a = b + k q $（其中 $k$ 为某个整数）

$\Longleftrightarrow a - x = b + k\cdot q - x$

$\Longleftrightarrow a - x = b - x + k\cdot q$ \textcolor{red}{ $\rhd$ $a-x$ 与 $b-x$ 之差是 $q$ 的某个整数倍}

$\Longleftrightarrow a - x \equiv b - x \bmod q$

$ $

\item \textbf{乘法：}$a \equiv b \bmod q \Longleftrightarrow a = b + k q $（其中 $k$ 为某个整数）

$\Longrightarrow a \cdot x = b \cdot x + k\cdot q \cdot x$

$\Longrightarrow a \cdot x = b \cdot x + k_x\cdot q$（其中 $k_x = k \cdot x$）\textcolor{red}{ $\rhd$ $a\cdot x$ 与 $b\cdot x$ 之差是 $q$ 的某个整数倍}

$\Longrightarrow a \cdot x \equiv b \cdot x \bmod q$

反过来，如果 $x$ 与 $q$ 互素（即 $\gcd(x,q)=1$），则 $x$ 在模 $q$ 下存在乘法逆元 $x^{-1}$。由 $a \cdot x \equiv b \cdot x \bmod q$ 可得

$\Longrightarrow a \cdot x \cdot x^{-1} \equiv b \cdot x \cdot x^{-1} \bmod q$

$\Longrightarrow a \equiv b \bmod q$

\end{enumerate}
\end{proof}

根据定理~\ref*{subsec:modulo-arithmetic}.1 中的模运算，还可以推导出如下性质：


\begin{tcolorbox}[title={\textbf{\tboxtheorem{\ref*{subsec:modulo-arithmetic}.2} 模运算的代数性质}}]

\begin{enumerate}
\item \textbf{结合律：}$(a \cdot b) \cdot c \equiv a \cdot (b \cdot c) \bmod q$

\item \textbf{交换律：}$(a \cdot b) \equiv (b \cdot a) \bmod q$

\item \textbf{分配律：}$(a \cdot (b + c)) \equiv ((a \cdot b) + (a \cdot c))  \bmod q$

\item \textbf{可替换性：}在模运算中，同余的数可以相互替换。

例如，假设 $(a \equiv b \bmod q)$ 且 $(c \equiv d \bmod q)$。那么，在模运算中，$a$ 与 $b$ 可以相互替换，$c$ 与 $d$ 也可以相互替换，如下所示：

$(a + c) \equiv (b + d) \equiv (a + d) \equiv (b + c) \bmod q$

$(a - c) \equiv (b - d) \equiv (a - d) \equiv (b - c) \bmod q$

$(a \cdot c) \equiv (b \cdot d) \equiv (a \cdot d) \equiv (b \cdot c) \bmod q$

\end{enumerate}

\end{tcolorbox}

定理~\ref*{subsec:modulo-arithmetic}.2 的证明与定理~\ref*{subsec:modulo-arithmetic}.1 类似，留给读者作为练习。


\subsection{逆元}
\label{subsec:modulo-inverse}

\begin{tcolorbox}[title={\textbf{\tboxdef{\ref*{subsec:modulo-inverse}} 模运算中的逆元}}]

在模 $q$ 下（即所有数除以 $q$ 后得到的余数组成的世界中），对于每个 $a \in \{0, 1, 2, \cdots, q-1\}$：

\begin{itemize}

\item $a$ 的\textbf{加法逆元}记为 $a_+^{-1}$，满足 $a + a_+^{-1} \equiv 0 \bmod q$。例如，在模 11 下，$3_+^{-1} = 8$，因为 $3 + 8 \equiv 0 \bmod 11$。

\item $a$ 的\textbf{乘法逆元}记为 $a_*^{-1}$，满足 $a \cdot a_*^{-1} \equiv 1 \bmod{q}$。这样的逆元存在，当且仅当 $\gcd(a,q)=1$。例如，在模 11 下，$3_*^{-1} = 4$，因为 $3 \cdot 4 \equiv 1 \bmod{11}$。

\end{itemize}

\end{tcolorbox}

\subsection{模除法}
\label{subsec:modulo-division}

在模运算中，\textit{模除法}与通常的数值除法不同。严格来说，并不存在一个独立的“模除法”运算，因为取模运算本身就只返回除法的余数。在实际使用中，“模除法”指的是在模逆元存在时，即满足 $\gcd(a,q)=1$ 时，乘以该逆元。在模 $q$ 下，用 $a$ 除 $b$ 的\textit{模除法}等价于计算\textit{模}乘法 $b \cdot a^{-1} \bmod q$。\textit{模除法}的结果与数值除法不同：\textit{模除法}总是得到整数（模 $q$ 的一个剩余类代表），因为它在模 $q$ 下将两个整数相乘；而数值除法的结果可以是实数。整数模 $q$ 的逆元可以通过扩展欧几里得算法计算（\href{https://www.youtube.com/watch?v=fz1vxq5ts5I}{YouTube 教程}）。

\subsection{中心化剩余表示}
\label{subsec:modulo-centered}

本节此前一直将余数视为正整数。例如，假设模 $q$ 的余数取自 $\{0, 1, \cdots, q-1\}$。这种表示称为标准（即无符号）剩余表示。与之对应，还可以采用有符号（即中心化）的剩余表示 $\left\{-\dfrac{q}{2}, -\dfrac{q}{2} + 1, \cdots, 0, \cdots, \dfrac{q}{2} - 2, \dfrac{q}{2} - 1\right\}$\footnote{这里假设 $q$ 为偶数。若 $q$ 为奇数，则剩余代表为 $\left\{ -\dfrac{q-1}{2}, -\dfrac{q-3}{2}, \cdots, 0, \cdots, \dfrac{q-3}{2}, \dfrac{q-1}{2}\right\}$。}，使剩余代表分布在 $0$ 两侧，而剩余代表的总数仍为 $q$。在这两种表示中，取模运算都将给定的值变换到相应的剩余代表范围内：(1) 如果给定值大于该范围的上界，就减去模数 $q$；(2) 如果该值小于范围的下界，就加上模数 $q$。这两种表示（标准表示与中心化表示）的区别仅在于上下界：标准剩余表示的下界和上界分别为 $0$ 与 $q-1$，而中心化剩余表示的下界和上界分别为 $-\dfrac{q}{2}$ 与 $\dfrac{q}{2} - 1$。标准剩余表示约定 $\mathbb{Z}_q = \{0, 1, \cdots, q-1\}$，中心化剩余表示则约定 $\mathbb{Z}_q = \left\{-\dfrac{q}{2}, -\dfrac{q}{2} + 1, \cdots, 0, \cdots, \dfrac{q}{2} - 2, \dfrac{q}{2} - 1\right\}$。

在两种表示中，加法、减法、乘法以及除数可逆时的除法都具有相同的性质。可以利用 \autoref{subsec:modulo-arithmetic} 中的推理证明这一点：在任一表示中，两个同余的剩余代表之差都是 $q$ 的整数倍。

逆元也具有相同的性质：$a$ 模 $q$ 的逆元为满足 $a \cdot a^{-1} \equiv 1 \bmod q$ 的 $a^{-1}$。

在某些情况下，使用有符号剩余表示很有帮助。先看标准（即无符号）剩余表示的例子：假设有表达式 $a + b \bmod q$，并且已知在某个应用中，$a + b$ 一定处于 $[0, q-1]$ 范围内（即 $0 \leq a + b \leq q-1$）。此时，$(a + b \bmod q)$ = $a + b$，因此可以省去取模运算，简化表达式。再看中心化（即有符号）剩余表示的另一个例子：假设有表达式 $a - b \bmod q$，并且已知在某个应用中，$a - b$ 一定处于 $\left[-\dfrac{q}{2}, \dfrac{q}{2} - 1\right]$ 范围内，则有 $(a - b \bmod q) = a - b$。但是，如果表达式 $a - b \bmod q$ 使用标准剩余表示，就不能省去取模运算，因为当 $a - b$ 为负数时，它小于标准剩余表示的下界（即 $0$），因此仍需进行模约减（即加上一个或多个 $q$）。

在 \autoref{subsec:rns-fastbconvex} 中，我们利用中心化剩余表示的这一有利性质设计 \textsf{FastBConvEx} 运算：在该算法中，由于已知 $-\dfrac{b_\alpha}{2} \leq \mu + u < \dfrac{b_\alpha}{2}$，因此可以将 $(\mu + u \bmod b_\alpha)$ 简化为 $\mu + u$。


\clearpage

\section{群}
\label{sec:group}
%\textbf{- First Read:} 
%\href{https://e.math.cornell.edu/people/belk/numbertheory/CyclotomicPolynomials.pdf}{Fields and Cyclotomic Polynomials}

\subsection{定义}
\label{subsec:group-def}


\begin{tcolorbox}[title={\textbf{\tboxdef{\ref*{subsec:group-def}} 群}}]
\noindent \textbf{\underline{集合中的元素}}
\begin{itemize}
\item \textbf{集合（$\mathbb{S}$）：}由元素组成的无序整体：$\mathbb{S} = \{a, b, c, \ldots\}$
\item \textbf{集合上的运算 $\bm{(+, \cdot)}$：}考虑定义在 $\mathbb{S}$ 上的两种二元运算：加法 $(+)$ 和乘法 $(\cdot)$。
\item \textbf{加法单位元（$0_{(+)}$，通常记为 $0$）：}如果元素 $i \in \mathbb{S}$ 对所有 $a \in \mathbb{S}$ 都满足 $i + a = a = a + i$，则称它为加法单位元。
\item \textbf{乘法单位元（$1_{(\cdot)}$，通常记为 $1$）：}如果元素 $i \in \mathbb{S}$ 对所有 $a \in \mathbb{S}$ 都满足 $i \cdot a = a = a \cdot i$，则称它为乘法单位元。
\item \textbf{加法逆元（$a^{-1}_{(+)}$）：}对于每个 $a \in \mathbb{S}$，它的加法逆元 $a^{-1}_{(+)}$（通常记为 $-a$）定义为满足 $a + a^{-1}_{(+)} = 0_{(+)} = a^{-1}_{(+)} + a$（即加法单位元）的元素。
\item \textbf{乘法逆元（$a^{-1}_{(\cdot)}$）：}对于每个关于 $(\cdot)$ 可逆的 $a \in \mathbb{S}$，它的乘法逆元 $a^{-1}_{(\cdot)}$（通常记为 $a^{-1}$）定义为满足 $a \cdot a^{-1}_{(\cdot)} = 1_{(\cdot)} = a^{-1}_{(\cdot)} \cdot a$（即乘法单位元）的元素。
\end{itemize}

$ $

\noindent \textbf{\underline{元素运算的性质}}
\begin{itemize}
\item \textbf{封闭性：}如果对于任意 $a, b \in \mathbb{S}$ 都有 $a + b \in \mathbb{S}$，则称集合 $\mathbb{S}$ 对运算 $(+)$ 封闭。同样，如果对于任意 $a, b \in \mathbb{S}$ 都有 $a \cdot b \in \mathbb{S}$，则称集合 $\mathbb{S}$ 对运算 $(\cdot)$ 封闭。
\item \textbf{结合律：}对于任意 $a,b,c \in \mathbb{S}$，有 $(a + b) + c = a + (b + c)$。
\item \textbf{交换律：}对于任意 $a,b \in \mathbb{S}$，有 $a + b = b + a$。
\item \textbf{分配律：}如果同时定义了 $(+)$ 和 $(\cdot)$（例如在环中），则有 $a \cdot (b + c) = (a \cdot b) + (a \cdot c)$，以及 $(a + b) \cdot c = a\cdot c + b\cdot c$。
\end{itemize}


$ $

\noindent \textbf{\underline{群及相关代数结构}}
\begin{itemize}
\item \textbf{半群：}半群是对某一种运算（$+$ 或 $\cdot$）封闭且满足结合律的元素集合。
\item \textbf{幺半群：}幺半群是具有单位元 $e$ 的半群；单位元与任意其他元素进行该运算时，都会保持该元素不变。

（例如，$0$ 是运算 $+$ 的单位元，$1$ 是运算 $\cdot$ 的单位元。）
\item \textbf{群：}群是每个元素都存在关于该运算的逆元的幺半群。
\item \textbf{阿贝尔群：}阿贝尔群是运算还满足交换律的群。
\end{itemize}
\end{tcolorbox}

\subsection{示例}
\label{subsec:group-ex}

$\mathbb{Z}$（即全体整数的集合）在加法（$+$）下构成阿贝尔群，因为：
\begin{itemize}
\item \textbf{封闭性：}对于任意整数 $a, b \in \mathbb{Z}$，$a + b = c$ 仍为整数（即 $a+b \in \mathbb{Z}$）。
\item \textbf{结合律：}对于任意整数 $a, b, c \in \mathbb{Z}$，有 $(a + b) + c = a + (b + c)$。
\item \textbf{单位元：}加法单位元为 0，因为对于任意 $a \in \mathbb{Z}$，有 $a + 0 = a$。
\item \textbf{逆元：}对于每个 $a \in \mathbb{Z}$，它的加法逆元为 $-a$，因为 $a + (-a) = 0$。
\item \textbf{交换律：}对于任意整数 $a, b \in \mathbb{Z}$，有 $a + b = b + a$。
\end{itemize}

$ $

\noindent $\mathbb{Z}$ 在乘法（$\cdot$）下构成幺半群，因为：
\begin{itemize}
\item \textbf{封闭性：}对于任意整数 $a, b \in \mathbb{Z}$，$a \cdot b = c$ 仍为整数（即 $a\cdot b \in \mathbb{Z}$）。
\item \textbf{结合律：}对于任意整数 $a, b, c \in \mathbb{Z}$，有 $(a \cdot b) \cdot c = a \cdot (b \cdot c)$。
\item \textbf{单位元：}乘法单位元为 1，因为对于任意 $a \in \mathbb{Z}$，有 $a \cdot 1 = a$。
\item \textbf{不具有逆元：}整数 $0$ 没有乘法逆元；对于 $a\ne0$，其倒数 $1/a$ 仅在 $a=\pm1$ 时为整数。因此，$(\mathbb{Z},\cdot)$ 是幺半群，但不是群。
\end{itemize}
$ $

\noindent $\mathbb{R}^\times$（即全体非零实数的集合）在乘法（$\cdot$）下构成阿贝尔群，因为：
\begin{itemize}
\item \textbf{封闭性：}对于任意实数 $a, b \in \mathbb{R}^\times$，$a \cdot b = c$ 仍为实数，且仍属于 $\mathbb{R}^\times$。
\item \textbf{结合律：}对于任意实数 $a, b, c \in \mathbb{R}^\times$，有 $(a \cdot b) \cdot c = a \cdot (b \cdot c)$。
\item \textbf{单位元：}乘法单位元为 1，因为对于任意实数 $a \in \mathbb{R}^\times$，有 $a \cdot 1 = a$。
\item \textbf{逆元：}对于每个实数 $a \in \mathbb{R}^\times$，它的乘法逆元为 $\dfrac{1}{a}$，后者是非零实数（$\in \mathbb{R}^{\times}$）。
\end{itemize}


\clearpage

\section{域}
\label{sec:field}
\textbf{- 参考资料：}
\href{https://e.math.cornell.edu/people/belk/numbertheory/CyclotomicPolynomials.pdf}{域与分圆多项式}~\cite{cyclotomic-polynomial}

\subsection{定义}
\label{subsec:field-def}

\begin{tcolorbox}[title={\textbf{\tboxdef{\ref*{subsec:field-def}} 域的定义}}]
\begin{itemize}
\item \textbf{环：}集合 $R$ 在加法 $(+)$ 下构成阿贝尔群，同时定义有封闭且满足结合律的乘法 $(\cdot)$，并且乘法对加法满足左右分配律：对于所有 $a,b,c\in R$，都有 $a\cdot(b+c)=a\cdot b+a\cdot c$ 和 $(a+b)\cdot c=a\cdot c+b\cdot c$。（乘法不一定满足交换律，例如矩阵乘法；除非特别说明是“含幺环”，否则不要求 $(\cdot)$ 的单位元存在。）
\item \textbf{域：}集合 $F$ 在 $(+)$ 下构成阿贝尔群，其非零元素 $F^\times=F\setminus\{0\}$ 在 $(\cdot)$ 下构成阿贝尔群，且乘法对加法满足分配律。
\item \textbf{伽罗瓦域（$\mathrm{GF}(p^n)$）：}具有 $p^n$ 个元素的域，其中 $p$ 为素数，$n$ 为正整数。
\item \textbf{$\mathbb{Z}_p$（$\mathbb{Z}/p\mathbb{Z}$）：}对于素数 $p$，集合 $\{0,1,\ldots,p-1\}$ 配备模 $p$ 的加法和乘法后构成有限域。更一般地，对于任意整数 $m\ge 2$，$\mathbb{Z}_m$ 是交换环；它构成域，当且仅当 $m$ 为素数。\end{itemize}
\end{tcolorbox}

$ $

\subsection{示例}
\label{subsec:field-ex}

$\mathbb{Z}$（全体整数的集合）是环，但不是域，因为并非它的所有元素都有乘法逆元（如 \autoref{subsec:group-ex} 所示）。

$ $

\noindent $\mathbb{R}$（全体实数的集合）是域。如 \autoref{subsec:group-ex} 所示，它在 $(+)$ 下构成阿贝尔群；其非零元素在 $(\cdot)$ 下构成阿贝尔群，且乘法对加法满足分配律。

$ $

\noindent $\mathbb{Z}_7 = \{0, 1, 2, 3, 4, 5, 6\}$ 是有限域，因为：
\begin{itemize}
\item \textbf{封闭性：}对于任意 $a,b \in \mathbb{Z}_7$，存在 $c_1,c_2 \in \mathbb{Z}_7$ 使得 $a+b \equiv c_1 \pmod{7}$ 且 $a\cdot b \equiv c_2 \pmod{7}$。
\item \textbf{结合律：}对于任意 $a,b,c \in \mathbb{Z}_7$，有 $(a+b)+c=a+(b+c)$ 和 $(a\cdot b)\cdot c=a\cdot(b\cdot c)$。
\item \textbf{交换律：}对于任意 $a,b \in \mathbb{Z}_7$，有 $a+b=b+a$ 和 $a\cdot b=b\cdot a$。
\item \textbf{分配律：}对于任意 $a,b,c \in \mathbb{Z}_7$，有 $(a+b)\cdot c=a\cdot c+b\cdot c$ 和 $a\cdot(b+c)=a\cdot b+a\cdot c$。
\item \textbf{单位元：}加法单位元为 $0$，乘法单位元为 $1$。
\item \textbf{逆元：}对于任意 $a \in \mathbb{Z}_7$，存在 $a' \in \mathbb{Z}_7$ 使得 $a+a' \equiv 0 \pmod{7}$（例如，$3$ 的加法逆元为 $4$，因为 $3+4\equiv0 \pmod{7}$）。对于任意 $a \in \mathbb{Z}_7^\times=\{1,\dots,6\}$，存在 $b \in \mathbb{Z}_7^{\times}$ 使得 $ab \equiv 1 \pmod{7}$（例如，$3\cdot5=15\equiv1 \pmod{7}$）。
\end{itemize}

\subsection{定理}
\label{subsec:field-theorem}

\begin{tcolorbox}[title={\textbf{\tboxtheorem{\ref*{subsec:field-theorem}} 域的相关定理}}]
\begin{enumerate}
\item \textbf{有限域的大小：}每个有限域（也称为伽罗瓦域）都具有 $p^n$ 个元素，其中 $p$ 为素数，$n$ 为正整数。反之，对于每个这样的素数幂 $p^n$，都存在阶为 $p^n$ 的有限域，并且在同构意义下唯一。
\item \textbf{同构域：}任意两个元素个数相同的有限域 $\mathbb{F}_1$ 和 $\mathbb{F}_2$ 都同构，即存在一个双射 $f:\mathbb{F}_1\to\mathbb{F}_2$，使得对于所有 $a,b\in\mathbb{F}_1$，都有 $f(a+b)=f(a)+f(b)$ 和 $f(ab)=f(a)f(b)$。
\end{enumerate}
\end{tcolorbox}


\clearpage

\section{阶}
\label{sec:order}
\textbf{- 参考资料：}
\href{https://e.math.cornell.edu/people/belk/numbertheory/CyclotomicPolynomials.pdf}{域与分圆多项式}~\cite{cyclotomic-polynomial}

\subsection{定义}
\label{subsec:order-def}

\begin{tcolorbox}[title={\textbf{\tboxdef{\ref*{subsec:order-def}} 阶的定义}}]
$\bm{\textsf{ord}_{\mathbb{F}}(a)}$：对于 $a \in \mathbb{F}^{\times}$（一个有限域，见 \autoref{subsec:field-def}），$a$ 的阶是使 $a^k = 1$ 成立的最小正整数 $k$。
\end{tcolorbox}

\subsection{定理}
\label{subsec:order-theorem}



\begin{tcolorbox}[title={\textbf{\tboxtheorem{\ref*{subsec:order-theorem}.1} 阶的性质（I）}}]
对于 $a \in \mathbb{F}^{\times}$ 和 $n \geq 1$，$a^n = 1$ 当且仅当 \textbf{\textsf{ord}}$_{\mathbb{F}}(a) \text{ } \mid \text{ } n$

（即 $\textsf{ord}_{\mathbb{F}}(a)$ 整除 $n$）。
\end{tcolorbox}

\begin{myproof}
    \begin{enumerate}
    \item \textit{正向证明：}若 $\textsf{ord}_{\mathbb{F}}(a) \text{ } | \text{ } n$，则对于 $\textsf{ord}_{\mathbb{F}}(a) = k$，其中 $k$ 为 $a$ 的阶，存在某个整数 $l$ 使得 $n = lk$。
    
    因此，$a^n = a^{lk} = (a^k)^l = 1^l = 1$。
    \item \textit{反向证明：}若 $a^n = 1$ 且 $\textsf{ord}_{\mathbb{F}}(a)=k$，将其写为 $n=qk+r$，其中 $0 \le r < k$。于是 $1=a^n=a^{qk+r}=(a^k)^q a^r=a^r$。由 $k$ 的最小性可知，必有 $r=0$，从而 $k \mid n$。
    \end{enumerate}
\end{myproof}


\begin{tcolorbox}[title={\textbf{\tboxtheorem{\ref*{subsec:order-theorem}.2} 阶的性质（II）}}]
若 $\textsf{ord}_{\mathbb{F}}(a) = k$，则对于任意 $n \geq 1$，都有 $\textsf{ord}_{\mathbb{F}}(a^n) = \dfrac{k} {\gcd(k, n)}$。
\end{tcolorbox}
\begin{myproof}
    \begin{enumerate}
    \item $a^k, a^{2k}, a^{3k}, \ldots = 1$. 
    \item 已知 $\textsf{ord}_{\mathbb{F}}(a^n) = m$，有 $(a^n)^m, (a^n)^{2m}, (a^n)^{3m}, \ldots = 1$。
    \item 根据阶的定义，$x=k$ 是满足 $a^x$ = 1 的最小值。因此，给定 $\textsf{ord}_{\mathbb{F}}(a^n) = m$，$m$ 就是使 $(a^n)^m = 1$ 成立的最小整数。注意，$(a^n)^m$ 应当是 $a^k$ 的倍数，也就是说，$mn$ 应当是 $k$ 的倍数。使 $mn$ 成为 $k$ 的倍数的最小整数 $m$ 为 $m = \dfrac{k}{\gcd(k, n)}$。
    \end{enumerate}
\end{myproof}

\begin{tcolorbox}[title={\textbf{\tboxtheorem{\ref*{subsec:order-theorem}.3} 阶的性质（III）}}]
假设 $k$ 整除 $n$。则 $\textsf{ord}_{\mathbb{F}}(a) = kn$ 当且仅当 $\textsf{ord}_{\mathbb{F}}(a^k) = n$。
\end{tcolorbox}
\begin{myproof}
\begin{enumerate}
    \item \textit{正向证明：}已知 $\textsf{ord}_{\mathbb{F}}(a) = kn$，由定理~\ref*{subsec:order-theorem}.2 可得 $\textsf{ord}_{\mathbb{F}}(a^k) = \dfrac{nk}{\gcd(k, nk)} = \dfrac{nk}{k} = n$。
    \item \textit{反向证明：}已知 $\textsf{ord}_{\mathbb{F}}(a^k)=n$，并令 $\textsf{ord}_{\mathbb{F}}(a)=m$。由定理~\ref*{subsec:order-theorem}.2 得到 $\textsf{ord}_{\mathbb{F}}(a^k)=\dfrac{m}{\gcd(m,k)}=n$，因此 $m=n\cdot\gcd(m,k)$（即 $m$ 是 $n$ 的某个倍数）。但由于 $k$ 整除 $n$，所以 $k$ 也整除 $m$。这意味着 $\gcd(m,k) = k$。因此，$\textsf{ord}_{\mathbb{F}}(a)=m=n\cdot\gcd(m,k)=nk$。
\end{enumerate}
\end{myproof}

\begin{tcolorbox}[title={\textbf{\tboxtheorem{\ref*{subsec:order-theorem}.4} 费马小定理}}]
给定素数 $|\mathbb{F}| = p$ 和 $a \in \mathbb{F}$，有 $a^p = a$。
\end{tcolorbox}
\begin{myproof}
    \begin{enumerate}
    \item 若 $a=0$，则 $a^p=a=0$。
    \item 若 $a\ne 0$，则 $a \in \mathbb{F}^{\times}$，即属于该域的乘法群，其大小为 $|\mathbb{F}^{\times}|=p-1$。
    根据拉格朗日定理（在有限群 $G$ 中，任意元素的阶都整除 $|G|$），$a$ 的阶整除 $p-1$，因此 $a^{p-1}=1$，从而 $a^p=a$。
    \end{enumerate}
\end{myproof}



\clearpage

\section{多项式环}
\label{sec:polynomial-ring}
\textbf{- 参考资料 1：}
\href{https://en.wikipedia.org/wiki/Polynomial_ring}{多项式环（维基百科）}~\cite{polynomial-ring}

\noindent \textbf{- 参考资料 2：}
\href{https://math.libretexts.org/Bookshelves/Combinatorics_and_Discrete_Mathematics/Applied_Discrete_Structures_(Doerr_and_Levasseur)/16%3A_An_Introduction_to_Rings_and_Fields/16.03%3A_Polynomial_Rings}{多项式环（LibreTexts）}~\cite{polynomial-rings}

\subsection{概述}
\label{subsec:poly-ring-overview}

\textbf{多项式环}是由多项式组成的集合，其中关于运算符 $(+, \cdot)$ 的多项式运算（例如 $f_1 + f_2$、$f_1 \cdot (f_2 - f_3)$、$f_1 + f_2 + f_4$）具有封闭性，并满足结合律、交换律和分配律。


多项式环 ${\mathbb{Z}_p[x] / (x^n + 1)}$ 是满足以下性质的所有多项式 $f_i$ 的集合：

\begin{tcolorbox}[title={\textbf{\tboxlabel{\ref*{subsec:poly-ring-overview}} 环}}]

对于满足 $f = c_0 + c_1x^1 + \cdots + c_{n-1}x^{n-1}$ 的多项式 $f \in \mathbb{Z}_p[x] / (x^n + 1)$：

\begin{itemize}
\item \textbf{系数环：}每个系数均满足 $c_j \in \mathbb{Z}_p$。

$ $

\item \textbf{次数界：}任意 $f' \in \mathbb{Z}_p[x]$ 都可写成：

$ $

$f'=(x^n+1)f_q+f_r,\qquad \deg f_r < n,$

$ $

因此，在商环 $\mathbb{Z}_p[x]/(x^n+1)$ 中有 $f' \equiv f_r \pmod{x^n+1}$。$f_q$ 和 $f_r$ 分别称为 $f'$ 除以 $x^n + 1$ 所得到的商多项式和余式。

$ $

\item \textbf{多项式同余：}两个同余的多项式属于同一个等价类，因此在多项式环中的多项式运算（$+, \cdot$）中可以相互替换。例如，如果：

$ $

$f' \equiv f_{r1} \in \mathbb{Z}_p[x] / (x^n + 1)$

$f{''} \equiv f_{r2} \in \mathbb{Z}_p[x] / (x^n + 1)$

$f_{r1} + f_{r2} \equiv f_{r3} \in \mathbb{Z}_p[x] / (x^n + 1)$

$ $

则多项式运算 $f' + f''$ 的结果与下式属于同一个等价类：

$f' + f'' \equiv f_{r1} + f_{r2} \equiv f_{r3} \in \mathbb{Z}_p[x] / (x^n + 1)$


\end{itemize}

为简化记号，将多项式环 $\mathbb{Z}_p[x] / (x^n + 1)$ 记为 $\mathcal{R}_{\langle n, p \rangle}$。
\end{tcolorbox}

回顾整数模环 $\mathbb{Z}_p$：任意 $b$ 都可以写成 $b = mp + r$，其中 $0\le r<p$，因此 $b \equiv r \pmod{p}$（商 $m$ 被消去）。类似地，在多项式环 $\mathcal{R}_{\langle n, p \rangle}$ 中，高次多项式 $f_{big}$ 可以除以多项式模数 $x^n + 1$，得到：

$f_{big} = (x^n + 1)\cdot(f_q) + f_r \equiv f_r \in \mathcal{R}_{\langle n, p \rangle}$

\noindent 其中，$f_q$ 是商多项式，$f_r$ 是余式。此时，$f_{big}$ 与 $f_r$ 同余，即属于同一个等价类。因此，可以消去 $f_q$，并在多项式环的多项式运算 $(+, \cdot)$ 中，用 $f_r$（即 $f_{big}$ 的简化形式）替代原多项式。
多项式环中的多项式简化（即约减），通过在 $f_{big}$ 中使用关系 $x^n \equiv -1$ 进行代换来完成，因为在该多项式环中有 $x^n + 1 \equiv 0$（这与模 $p$ 的整数环中用 $0$ 替换 $p$ 来约减一个数的情形相同）。在 $f_{big}$ 中不断使用关系 $x^n \equiv -1$ 进行代换，即可将高次多项式 $f_{big}$ 逐步简化为次数小于 $n$ 的多项式。

作为多项式模数，通常选择分圆多项式 $x^n + 1$ 作为除式，其中 $n$ 为 $2^p$，$p$ 为整数，因为这种选择有利于提高计算效率。

\subsubsection{示例}
\label{subsubsec:poly-ring-ex}
给定 $f \in \mathbb{Z}_7[x] / (x^2 + 1)$，假设 $f = x^4 + 3x^3 + 11x^2 + 6x + 10$。则有：

$ $

$f = (x^2)\cdot(x^2) + 3x\cdot(x^2) + 11x^2 + 6x + 10$ 

$ \equiv (-1)(-1) + 3x(-1) +  (11 \bmod 7)(-1) + 6x + (10 \bmod 7)$

$ = 3x \in \mathbb{Z}_7[x] / (x^2 + 1)$

$ $

\noindent 因此，在多项式环 $\mathbb{Z}_7[x] / (x^2 + 1)$ 中，$f(x) = x^4 + 3x^3 + 11x^2 + 6x + 10$ 与 $3x$ 等价（$\equiv$）。


\subsubsection{讨论}
\label{subsubsec:polynomial-ring-discuss}

\begin{table}[h]
\centering
\footnotesize
%\noindent\adjustbox{max width=\columnwidth}{
\begin{tabular}{|c||c|c|} % left align
\hline
&\textbf{整数模环} & \textbf{多项式环} \\
\hline
\hline
\textbf{{集合元素}} & 数 & 多项式 \\
\hline
\textbf{{环的记号}} & $\mathbb{Z}_p = \{0, 1, \ldots, p - 1\}$ & $\mathbb{Z}_p[x] / (x^n + 1)$ \\
\textbf{{及定义}} & $0$ 到 $p$ 之间的所有整数的集合 & 满足下式的所有多项式 $f$ 的集合：\\
&& $f= c_0 + c_1x^1 + c_2x^2 \cdots + c_{n-1}x^{n-1}$ \\
&& 其中每个系数满足 $c_i \in \mathbb{Z}_p$ \\
&& 且 $f$ 的次数小于 $n$ \\
\hline
\textbf{{支持的}}&$(+, \cdot)$& $(+, \cdot)$ \\
\textbf{{运算}}&（加法、乘法）&（加法、乘法）\\
\hline
\textbf{{（$+$）运算}} & 通常的数值加法 & $f_a = a_0 + a_1x^1 + a_2x^2 \cdots + a_{d_a-1}x^{d_a-1}$ \\
 & & $f_b = b_0 + b_1x^1 + b_2x^2 \cdots + b_{d_b-1}x^{d_b-1}$ \\
&& 则 $f_a + f_b$ 的计算方式为：\\
 & & $f_c = \sum\limits_{i=0}^{\textsf{max}(d_a,d_b)}(a_i+b_i)x^i$ \\
\hline
\textbf{{（$\cdot$）运算}} & 通常的数值乘法 & $f_a = a_0 + a_1x^1 + a_2x^2 \cdots + a_{d_a-1}x^{d_a-1}$ \\
 & & $f_b = b_0 + b_1x^1 + b_2x^2 \cdots + b_{d_b-1}x^{d_b-1}$ \\
&& 则 $f_a \cdot f_b$ 的计算方式为：\\
 & & $f_c = \sum\limits_{i=0}^{d_a+d_b}\sum\limits_{j=0}^{i}a_jb_{i-j}x^i$ \\
\hline
& 对于 $a, b \in \mathbb{Z}_p$ & 对于 $f_a, f_b \in \mathbb{Z}_p[x]/(x^n + 1)$，\\
\textbf{{交换律}} & $ a + b = b + a$ & $f_a + f_b = f_b + f_a$\\
\textbf{{结合律}} & $(a + b) + c = a + (b + c)$ & $(f_a + f_b) + f_c = f_a + (f_b + f_c)$\\
\textbf{{分配律}} & $a \cdot (b + c) = a\cdot b + a\cdot c$ & $f_a \cdot (f_b + f_c) = f_a\cdot f_b + f_a\cdot f_c$\\
\textbf{{封闭性}}& $a + b \equiv c \in \mathbb{Z}_p$，$a \cdot b \equiv d \in \mathbb{Z}_p$ & $f_a + f_b \equiv f_c \in \mathcal{R}_{\langle n, p \rangle}$，\\
&&$f_a \cdot f_b \equiv f_d \in \mathcal{R}_{\langle n, p \rangle}$\\
\hline
\textbf{{同余（$\equiv$）}} & 两个数模 $p$ 同余，即 $a \equiv b$，若： & 两个多项式在 $\mathbb{Z}_p[x] / (x^n + 1)$ 中同余，即 $f_a \equiv f_b$，若：\\
 & $(a \bmod p) = (b \bmod p)$   & $f'_a = f_a \bmod (x^n + 1) = \sum\limits_{i=0}^{d_a} a_ix^i$ \\
 & & $f'_b = f_b \bmod (x^n + 1) = \sum\limits_{i=0}^{d_b} b_ix^i$, \\
&& $d_a = d_b$ 且 $a_i \equiv b_i$，均在模 $p$ 下\\
&&对所有 $0 \leq i \leq d_a$ 成立\\
\hline
\end{tabular}%}
\centering
\caption{整数模环与多项式环的比较。
}
\label{tab:ring-comparison}
\end{table}

\para{同余：}两个同余的数属于同一个\textit{同余类}；两个同余的多项式也是如此。如果两个数学表达式的计算结果属于同一个同余类，那么它们的计算会在相应的模数范围内循环约减。这一性质对于密码方案很有用，因为加密和解密运算需要将数值限制在有限位数的二进制表示范围内。同余还可用于简化计算。例如，可以利用同余规则，将一个大数或复杂多项式\textit{规范化}为更简单的数或多项式，从而降低计算开销。

\para{多项式求值：}属于同一个同余类的两个数不一定相等。例如，在模 5 下有 $5 \equiv 10$，但这两个数并不相同。类似地，两个同余的多项式也不一定相同。在多项式环所支持的多项式运算（即 $(+, \cdot)$）中，例如 $f_1 + f_2$ 或 $f_1 \cdot (f_2 - f_3)$，两个同余的多项式可以相互替换；但是，将某个变量值 $x = a$ 代入它们求值时，结果不一定相同。例如，在 \autoref{subsubsec:poly-ring-ex} 的前一个例子中，两个多项式 $x^4 + 3x^3 + 11x^2 + 6x + 10$ 和 $3x$ 在多项式环 $\mathbb{Z}_7[x] / (x^2 + 1)$ 中同余。然而，在 $x = 0$ 处求值时，原多项式的结果为 10，而约减后（即简化后）的多项式结果为 0。这种求值差异源于同余的定义：若两个多项式 $M_1$ 和 $M_2$ 除以 $x^n + 1$ 后的余式相同（即对于某个多项式 $Q$，有 $M_1 = Q \cdot (x^n + 1) + M_2$），便称它们在 $x^n + 1$ 上同余（即 $M_1 \equiv M_2$）。因此，若在满足 $x^n = -1$ 的 $x$ 处求值，$x^n +1$ 项就为 0，此时 $M_1$ 和 $M_2$ 的求值结果都等于同一个多项式 $M_2$ 的求值结果。方程 $x^n = -1$ 的解称为 $n$ 次单位根，将在 \autoref{sec:roots} 和 \autoref{sec:cyclotomic-polynomial-integer-ring} 中介绍。归纳如下：

\begin{tcolorbox}[title={\textbf{\tboxlabel{\ref*{subsubsec:polynomial-ring-discuss}} 多项式环中的多项式求值}}]

假设多项式 $M_1$ 和 $M_2$ 在多项式环 $x^n + 1$ 中同余。这等价于存在某个多项式 $Q$，使得 $M_1 = Q \cdot (x^n + 1) + M_2$。如果 $X=x_i$ 是方程 $x^n + 1$ 的解（即 $X$ 是 $n$ 次单位根），则可以保证 $M_1(X)$ 与 $M_2(X)$ 的求值结果相同，即 $M_1(x_i) = M_2(x_i)$。

\end{tcolorbox}

\para{整数模环与多项式环：}这两种环具有许多共同特征，汇总于 \autoref{tab:ring-comparison}。


\subsection{系数旋转}
\label{subsec:coeff-rotation}

系数旋转是在多项式环中，将一个多项式的全部系数向左或向右移动的过程。若要将全部系数向左旋转 $h$ 个位置，可以将该多项式乘以 $x^{-h}$。

例如，假设有如下多项式：

$ $

$f(x) = c_0 + c_1x^1 + c_2x^2 + \cdots + c_hx^h + \cdots + c_{n-1}x^{n-1} \in \mathcal{R}_{\langle n, p \rangle}$

$ $

要将 $f$ 的所有系数向左移动 $h$ 个位置（即将 $f$ 的第 $h$ 个系数移至常数项），需要计算 $f \cdot x^{-h}$，如下所示：

$ $

\begin{tcolorbox}[title={\textbf{\tboxlabel{\ref*{subsec:coeff-rotation}} 多项式旋转}}]

给定如下 $(n-1)$ 次多项式：

$ $

$f(x) = c_0 + c_1x^1 + c_2x^2 + \cdots + c_hx^h + \cdots + c_{n-1}x^{n-1} \in \mathcal{R}_{\langle n, p \rangle}$

$ $

将 $f(x)$ 乘以 $x^{-h}$，即可把 $f(x)$ 的系数向左旋转 $h$ 个位置，具体如下：

$ $

$f(x) \cdot x^{-h} = c_0\cdot x^{-h} + c_1x^1 \cdot x^{-h} + c_2x^2 \cdot x^{-h} + \cdots + c_hx^h \cdot x^{-h} + \cdots + c_{n-1}x^{n-1} \cdot x^{-h}$
$\equiv c_h + c_{h+1}x + c_{h+2}x^2 + \cdots + c_{n-1}x^{n-1-h} - c_0x^{n - h} - \cdots - c_{h-1}x^{n-1} \in \mathcal{R}_{\langle n, p \rangle}$
\end{tcolorbox}

$ $

注意，将两个多项式 $f$ 与 $x^{-h}$ 相乘，会得到在 $\mathcal{R}_{\langle n, p \rangle}$ 中的一个同余多项式。因此，旋转后的多项式，即 $f \cdot x^{-h}$ 的结果，也会在 $\mathcal{R}_{\langle n, p \rangle}$ 中有一个同余多项式。

由于计算在多项式环 $\mathbb{Z}_p[x] / (x^n + 1)$ 中进行，系数在旋转越过 $x^n (= -1)$ 的边界时，其符号会发生变化。




\subsubsection{示例}
\label{subsec:coeff-rotation-ex}

假设有如下多项式 $f \in \mathbb{Z}_8[x] / (x^4 + 1)$：

对 $\mathbb{Z}_8$ 使用中心化剩余表示，即 $\{-4,-3,-2,-1,0,1,2,3\}$。

$f = 2 + 3x - 4x^2 -x^3$




$ $

多项式环 $\mathbb{Z}_8[x] / (x^4 + 1)$ 具有以下 4 个同余关系：

{\boldmath{$x^4 \equiv -1$}}

$x^4 \cdot x^{-1} \equiv -1 \cdot x^{-1}$

{\boldmath{$x^3 \equiv -x^{-1}$}}

$ $

$x^4 \equiv -1$

$x^4 \cdot x^{-3} \equiv -1 \cdot x^{-3}$

{\boldmath{$x \equiv -x^{-3}$}}

$ $

$x^4 \equiv -1$

$x^4 \cdot x^{-2} \equiv -1 \cdot x^{-2}$

{\boldmath{$x^2 \equiv -x^{-2}$}}

$ $



根据总结~\ref*{subsec:coeff-rotation} 中的系数旋转方法，向左旋转 1 个位置等价于计算 $f \cdot x^{-1}$，如下所示：

$f\cdot x^{-1} = 2\cdot(x^{-1}) + 3x\cdot(x^{-1}) - 4x^2\cdot(x^{-1}) - x^3\cdot(x^{-1})$

$\equiv -2x^{3} + 3 - 4x^1 - x^2$

$= 3 - 4x^1 - x^2 -2x^{3}$

$ $

类似地，向左旋转 3 个位置等价于计算 $f \cdot x^{-3}$，如下所示：

$f\cdot x^{-3} = 2\cdot(x^{-3}) + 3x\cdot(x^{-3}) - 4x^2\cdot(x^{-3}) - x^3\cdot(x^{-3})$

$\equiv -2x - 3x^2 + 4x^3 - 1$

$= -1 - 2x - 3x^2 + 4x^3$

$= -1 - 2x - 3x^2 + (4 \equiv -4 \bmod 8) x^3$

$\equiv -1 - 2x - 3x^2 -4x^3$

\clearpage

\section{分解}
\label{sec:decomp}
分解是一种在保持数值不变的同时，以较小的底数（基数）表示一个数的数学技术。本节介绍数的分解与多项式分解。

\subsection{数的分解}
\label{subsec:number-decomp}
固定模数 $q\ge 2$，并记 $\mathbb{Z}_q=\mathbb{Z}/q\mathbb{Z}$。设 $\gamma\in\mathbb{Z}_q$。数的分解将 $\gamma$ 表示为以 $\beta$ 为基数的多个数之和：

$ $

$\gamma = \gamma_1 \dfrac{q}{\beta^1} + \gamma_2 \dfrac{q}{\beta^2} + \cdots + \gamma_\ell \dfrac{q}{\beta^\ell}  $

$ $

\noindent 其中，$\beta\ge 2$ 和 $\ell\ge1$ 均为整数。若要求每个规范代表元 $0\le\gamma<q$ 都具有满足 $\gamma_i\in\{0,1,\ldots,\beta-1\}$ 的唯一展开，应假设 $q=\beta^\ell$。更一般地，$\beta^\ell\mid q$ 仅保证权重为整数；若 $d=q/\beta^\ell>1$，则精确展开存在当且仅当 $d\mid\gamma$。图~\ref{fig:decomp} 展示完整分解；截断情形见 \autoref{subsec:approx-decomp}。
当 $q=\beta^\ell$ 时，每个 $\gamma_i$ 是 $\gamma$ 的基 $\beta$ 展开的一个数位，其中 $i=1$ 为最高位。当 $q$ 和 $\beta$ 均为 $2$ 的幂时，才可用二进制位移解释这些权重。

\begin{figure}[h!]
    \centering
  \includegraphics[width=0.8\linewidth]{figures/decomp.pdf}
  \caption{数的分解示意图。}
  \label{fig:decomp}
\end{figure}

将数 $\gamma$ 按基数 $\beta$、层数 $\ell$ 进行分解，定义如下：


$ $

$\textsf{Decomp}^{\beta, \ell}(\gamma) = (\gamma_1, \gamma_2, \text{ } \cdots , \text{ } \gamma_\ell)$.
 
$ $

数的分解也称为基数分解，其中的底数 $\beta$ 称为基数。


\subsubsection{示例}

假设在 $\mathbb{Z}_{16}$ 中取 $\gamma=13$，希望以基数 $\beta = 2$、层数 $\ell = 4$ 分解 13。其分解如下：

$ $

$13 = 1 \cdot \dfrac{16}{2^1} + 1 \cdot \dfrac{16}{2^2} + 0 \cdot \dfrac{16}{2^3} + 1 \cdot \dfrac{16}{2^4}$

$ $

$\textsf{Decomp}^{2, 4}(13) = (1, 1, 0, 1)$


\subsection{多项式分解}
\label{subsec:poly-decomp}

现在，假设在多项式环 ${\mathbb{Z}_q[x] / (x^n + 1)}$ 中有一个多项式 $f$。因此，$f$ 的每个系数 $c_i$ 都是模 $q$ 的整数。多项式分解将 $f$ 表示为以 $\beta$ 为基数、$\ell$ 为层数的多个多项式之和：


\begin{tcolorbox}[title={\textbf{\tboxlabel{\ref*{subsec:poly-decomp}} 多项式分解}}]

给定 $f\in \mathbb{Z}_q[x]/(x^n+1)$，选取满足 $q=\beta^\ell$ 的整数 $\beta\ge2$、$\ell\ge1$，写为

$ $

$
f=\sum\limits_{i=1}^{\ell} f_i\,\dfrac{q}{\beta^i}, \qquad f_i\in \mathbb{Z}_q[x]/(x^n+1)
$

$ $

其中，每个 $f_i$ 都通过将 $f$ 的每个系数按基数 $\beta$ 分解而得到。若 $f=\sum\limits_j c_j x^j$ 且 $c_j\in\mathbb{Z}_q$，则
$c_j=\sum\limits_{i=1}^{\ell} c_{j,i}\,\dfrac{q}{\beta^i}$，其中 $c_{j,i}\in\{0,\ldots,\beta-1\}$，
并且 $f_i=\sum_j c_{j,i} x^j$。
将多项式 $f$ 按基数 $\beta$、层数 $\ell$ 分解，记为：

$ $

$\textsf{Decomp}^{\beta, \ell}(f) = (f_1, f_2, \text{ } \cdots , \text{ } f_\ell)$
 $ $
\end{tcolorbox}




\subsubsection{示例}

假设在多项式环 $\mathbb{Z}_{16}[x] / (x^4 + 1)$ 中有如下多项式：

$ $

$f = 6x^3 + 3x^2 + 14x + 7 \in \mathbb{Z}_{16}[x] / (x^4 + 1)$

$ $

希望按基数 $\beta = 4$、层数 $\ell = 2$ 分解上述多项式。则多项式 $f$ 的各次项系数分别分解如下：

$ $

${\bm{x}^{\bm{3}}}$: $6 = \color{blue}{1 \cdot \dfrac{16}{4^1}} \color{black}+ \color{red}{2 \cdot \dfrac{16}{4^2}}$

${\bm{x}^{\bm{2}}}$: $3 = \color{blue}{0 \cdot \dfrac{16}{4^1}} \color{black}+ \color{red}{3 \cdot \dfrac{16}{4^2}}$

${\bm{x}^{\bm{1}}}$: $14 = \color{blue}{3 \cdot \dfrac{16}{4^1}} \color{black}+ \color{red}{2 \cdot \dfrac{16}{4^2}}$

${\bm{x}^{\bm{0}}}$: $7 = \color{blue}{1 \cdot \dfrac{16}{4^1}} \color{black}+ \color{red}{3 \cdot \dfrac{16}{4^2}}$

$ $

分解得到的多项式如下：

$f = 6x^3 + 3x^2 + 14x + 7 = \color{blue}{(1x^3 + 0x^2 + 3x + 1) \cdot \dfrac{16}{4^1}} \color{black}+ \color{red}{(2x^3 + 3x^2 + 2x + 3) \cdot \dfrac{16}{4^2}} \color{black}$

$ $

$\textsf{Decomp}^{4, 2}(6x^3 + 3x^2 + 14x + 7) = (1x^3 + 0x^2 + 3x + 1, 2x^3 + 3x^2 + 2x + 3)$

\subsubsection{讨论}

分解之后，多项式的原始系数被拆成了更小的数。这一性质在全同态加密密文的多项式乘法中非常有用，可以降低噪声的增长速度。通常，密文的多项式系数较大，因为它们是均匀随机选取的数。减小这些较大的常数对于控制同态乘法中的噪声增长很重要。\autoref{subsec:tfhe-mult-cipher} 将进一步讨论这一点。


\subsection{近似分解}
\label{subsec:approx-decomp}

\begin{figure}[h!]
    \centering
  \includegraphics[width=0.7\linewidth]{figures/decomp3.pdf}
  \caption{近似分解示意图。}
  \label{fig:decomp3}
\end{figure}

选取满足 $\beta^\ell\mid q$ 的整数 $\beta\ge2$、$\ell\ge1$，并令 $d=q/\beta^\ell$。即使整除条件成立，只要 $d>1$ 且输入不是 $d$ 的倍数，就仍须截断。将 $\gamma$ 的规范代表元舍入到最接近的 $d$ 的整数倍，再将该整数倍模 $q$ 约减，并按上述权重分解。固定最近整数舍入的平局规则后，结果满足
$$
\gamma\equiv\sum_{i=1}^{\ell}\gamma_i\frac{q}{\beta^i}+\varepsilon\pmod q,
\qquad \gamma_i\in\{0,\ldots,\beta-1\},\qquad |\varepsilon|\le\frac{q}{2\beta^\ell}.
$$
残差取圆周舍入误差的一个整数提升。这里必须使用同余：例如 $q=16$、$\beta=2$、$\ell=2$ 时，输入 $15$ 舍入到 $16\equiv0\pmod{16}$，此时 $\varepsilon=-1$。多项式情形逐系数处理。若 $\beta^\ell\nmid q$，上述有理数权重不能定义这里的整数环 Gadget；可以改用 \autoref{subsec:gadget-decomposition} 中的整数 Gadget。


\subsection{Gadget 分解}
\label{subsec:gadget-decomposition}

Gadget 权重和允许的数位集合都是定义的一部分；任意权重不一定能表示每个剩余类。对于任意 $q\ge2$，一种精确的整数选择是 $g_i=\beta^{i-1}$，取 $\ell=\lceil\log_\beta q\rceil$，并使用规范代表元的普通进制数位。它在整数中满足 $\gamma=\sum_{i=1}^{\ell}\gamma_i g_i$。这与下面的降序权重不同，两种权重的数位向量不可混用。对于截断 Gadget，应使用带残差的模 $q$ 重构式，并在后续噪声分析中计入该残差。

Gadget 分解是数的分解（\autoref{subsec:number-decomp}）的一般化形式。在数的分解中，数 $\gamma$ 被分解如下：

$\gamma = \gamma_1 \dfrac{q}{\beta^1} + \gamma_2 \dfrac{q}{\beta^2} + \cdots + \gamma_\ell \dfrac{q}{\beta^\ell} $

$ $

在 Gadget 分解中，按如下方式分解 $\gamma$：

$\gamma = \gamma_1 g_1 + \gamma_2 g_2 + \cdots + \gamma_\ell g_\ell $

$ $

将 $\vec{g} = (g_1, g_2, \cdots, g_\ell)$ 称为 Gadget 向量，并记 $\textsf{Decomp}^{\vec{g}}(\gamma) = (\gamma_1, \gamma_2, \text{ } \cdots , \text{ } \gamma_\ell) $。

$ $

于是，$\gamma = \langle \textsf{Decomp}^{\vec{g}}(\gamma), \vec{g} \rangle $。

$ $

对于数的分解（\autoref{subsec:number-decomp}），其 Gadget 向量为 $\vec{g} = \Bigg(\dfrac{q}{\beta}, \dfrac{q}{\beta^2}, \cdots, \dfrac{q}{\beta^\ell}\Bigg)$。


\clearpage

\section{单位根}
\label{sec:roots}
\textbf{- 参考资料：}
\href{https://e.math.cornell.edu/people/belk/numbertheory/CyclotomicPolynomials.pdf}{域与分圆多项式}~\cite{cyclotomic-polynomial}


\subsection{定义}
\label{subsec:roots-def}
\begin{tcolorbox}[title={\textbf{\tboxdef{\ref*{subsec:roots-def}} 单位根的定义}}]
\begin{itemize}
\item \textbf{$\bm{n}$ 次单位根：}满足方程 $\zeta^n = 1$ 的复数 $\zeta$。
\item \textbf{$\bm{n}$ 次本原单位根：}满足 $\textsf{ord}_{\mathbb{C}}(\zeta) = n$ 的任意 $n$ 次单位根 $\zeta$。将 $n$ 次本原单位根的集合记为 $P(n)$。
\end{itemize}

$ $

一个 $n$ 次本原单位根可以视为全部 $n$ 个 $n$ 次单位根的\textit{生成元}。
\end{tcolorbox}

\subsection{定理}
\label{subsec:roots-theorem}


\begin{tcolorbox}[title={\textbf{\tboxtheorem{\ref*{subsec:roots-theorem}.1} $\bm{n}$ 次单位根的公式}}]
给定 $\zeta^n = 1$，恰有 $n$ 个不同的 $n$ 次单位根：

$ $

$\zeta = e^{2k\pi i/n} = \cos\left(\dfrac{2k\pi}{n}\right) + i\cdot\sin\left(\dfrac{2k\pi}{n}\right)$, 

$ $

\noindent 对应于 $n$ 个不同的 $k$ 值，其中 $k = \{0, 1, \cdots, n-1\}$。
\end{tcolorbox}
\begin{myproof} 
    \begin{enumerate}
    \item 假设 $\zeta = e^{2k\pi i/n}$。则有 $\zeta^n = (e^{2k\pi i/n})^n = e^{2k\pi i}$。由于 $\zeta^n = 1$，需要寻找使 $e^{2k\pi i} = 1$ 成立的 $k$ 值。
    \item 欧拉公式给出 $e^{i\cdot x} = \text{cos}(x) + i \cdot \text{sin}(x)$。因此，若 $x = 2k\pi$，则有 $e^{2k\pi i} = \text{cos}(2k\pi) + i \cdot \text{sin}(2k\pi)$。当 $k = 0, 1, 2, ...$ 时，该表达式的值为 1。所以，对于任意整数 $k \geq 0$，都有 $e^{2k\pi i} = 1$。
    \item 若 $\zeta = e^{2k\pi i/n} = \text{cos}(\frac{2k\pi}{n}) + i\cdot\text{sin}(\frac{2k\pi}{n})$，则 $k = 0, 1, ... \text{ } n-1$ 所对应的前 $n$ 个根互不相同，因为它们位于复平面的圆上：$x$ 轴表示实部，$y$ 轴表示虚部的系数；对于 $k = \{0, 1, \cdots, n-1\}$，各点的角度分别为 $2k\pi/n$。
    \item 注意，$\zeta^n = 1$ 是 $n$ 次多项式，因此至多有 $n$ 个根。所以，可以将 $k = \{0, 1, \cdots, n-1\}$ 所对应的前 $n$ 个根 $e^{2k\pi i/n}$ 视为这 $n$ 个互异的根，而忽略其余的根（即 $k \geq n$），因为后者只是复平面圆上前 $n$ 个根的重复（见 \autoref{fig:complex-plane}）。
    \end{enumerate}
\end{myproof}

\begin{figure}[h!]
    \centering
  \includegraphics[width=0.4\linewidth]{figures/euler-formula.png}
  \caption{复平面中由欧拉公式描述的圆。\href{https://en.wikipedia.org/wiki/Euler's_formula}{（来源）}}
  \label{fig:complex-plane}
\end{figure}


\begin{tcolorbox}[title={\textbf{\tboxtheorem{\ref*{subsec:roots-theorem}.2} 单位根的阶}}]
给定 $\zeta \in \mathbb{C}$（复数域）以及满足 $n \geq 1$ 的 $\zeta^n = 1$，$\zeta$ 是 $n$ 次单位根，当且仅当 $\textsf{ord}_{\mathbb{C}}(\zeta) \text{ } | \text{ } n$。
\end{tcolorbox}
\begin{proof}
    使用
    定理~\ref*{subsec:order-theorem}.1：
    \begin{enumerate}
    \item \textit{正向证明：}由于 $\textsf{ord}_{\mathbb{C}}(\zeta)=k$ 是使 $\zeta^k = 1$ 成立的最小整数，对于任意满足 $\zeta^n = 1$ 的 $n$，$n$ 必须是 $k$ 的倍数。这意味着 $k \mid n$。
    \item \textit{反向证明：}若 $k \mid n$，则 $n$ 是 $k$ 的倍数，因此 $\zeta^n = 1$。
    \end{enumerate}
    %$(Here, note that our \b\bmodulo computation wraps around a circle of Euler's formula on the complex plane instead of a finite integer field).
\end{proof}

\begin{tcolorbox}[title={\textbf{\tboxtheorem{\ref*{subsec:roots-theorem}.3} 全部 $\bm{n}$ 次单位根的集合}}]
全部 $n$ 次单位根的集合等于并集 $\bigcup_{d|n} P(d)$，即所有满足 $d \text{ } | \text{ } n$ 的 $d$ 次本原单位根集合之并。
\end{tcolorbox}
\begin{myproof}
    \begin{enumerate}
    \item 设 $\omega = e^{2\pi i/n}$。
    给定 $\zeta^n = 1$，对于每个 $n$ 次单位根 $\zeta$，在 $k_i = \{0, 1, \cdots, n-1\}$ 时有 $\zeta = \omega^{k_i}$。根据定理~\ref*{subsec:order-theorem}.1，$(\omega^{k_i})^n = 1$ 当且仅当 $\textsf{ord}_{\mathbb{C}}(\omega^{k_i}) \text{ } | \text{ } n$。
    \item 设 $\textsf{ord}_{\mathbb{C}}(\omega^{k_i}) = d_i$，则 $(\omega^{k_i})^{d_i} = 1$。综合这两个事实可知，每个 $n$ 次单位根 $\omega^{k_i}$ 也都是 $d_i$ 次本原单位根（即 $\zeta^{d_i} = 1$ 的一个解），也就是 $\omega^{k_i} \in P(d_i)$。
    \item 对于每个 $\textsf{ord}_{\mathbb{C}}(\omega^{k_i}) = d_i$，都有 $d_i \text{ } | \text{ } n$。对于每个整除 $n$ 的 $d_i$，所有（本原）$d_i$ 次单位根也都是 $n$ 次单位根。这是因为，满足 $\zeta^{d_i} = 1$ 的（本原）$d_i$ 次单位根也满足 $\zeta^{n} = 1$，因为 $n$ 是 $d_i$ 的倍数。
    \item 步骤 2 表明，每个 $n$ 次单位根，都是某个整除 $n$ 的 $d_i$ 所对应的 $d_i$ 次本原单位根。步骤 3 表明，当 $d_i$ 整除 $n$ 时，每个 $d_i$ 次单位根也都是 $n$ 次单位根。综合这两个结论，全部 $n$ 次单位根的集合，等于所有满足 $d_i$ 整除 $n$ 的 $d_i$ 次本原单位根集合之并（即 $\bigcup_{d|n} P(d)$）。
    \end{enumerate}
\end{myproof}

\begin{tcolorbox}[title={\textbf{\tboxtheorem{\ref*{subsec:roots-theorem}.4} $\bm{n}$ 次本原单位根的判定条件}}]
给定 $k = \{0, 1, \cdots, n-1\}$ 所对应的 $n$ 次单位根 $\zeta = \omega^k$，其中 $\omega = e^{2\pi i/n}$。则 $\zeta$ 是 $n$ 次本原单位根，当且仅当 $\text{gcd}(n, k) = 1$，即 $k$ 与 $n$ 互素。
\end{tcolorbox}
\begin{myproof}
    \begin{enumerate}
    \item 注意，对于 $k = 1$，有 $\zeta^n = 1$ 和 $\zeta = \omega$。因此，$\textsf{ord}_{\mathbb{C}}(\omega) = n$。
    \item 定理~\ref*{subsec:order-theorem}.2 表明，若 $ord_{\mathbb{F}}(a) = k$，则对于任意 $n \geq 1$，都有 $ord_{\mathbb{F}}(a^n) = \dfrac{k}{\text{gcd}(k,n)}$。类似地，若 $\textsf{ord}_{\mathbb{C}}(\omega) = n$，则对于任意 $k \geq 1$，都有 $\textsf{ord}_{\mathbb{C}}(\omega^k) = \dfrac{n}{\text{gcd}(k, n)}$。 当 $k=0$ 时，也可直接验证同一公式：$\operatorname{ord}(\omega^0)=1=n/\gcd(n,0)$；这包括 $n=1$ 的情况。
    \item 由步骤 2 可知，$\textsf{ord}_{\mathbb{C}}(\omega^k) = n$（即 $\omega^k$ 是 $n$ 次本原单位根），当且仅当 $\text{gcd}(k, n) = 1$。
    \end{enumerate}
\end{myproof}

\begin{tcolorbox}[title={\textbf{\tboxtheorem{\ref*{subsec:roots-theorem}.5} $\bm{n}$ 次本原单位根的个数}}]
$n$ 次本原单位根的个数为 $\phi(n)$，即 $\{1, \cdots, n \}$ 中与 $n$ 互素的元素个数。
\end{tcolorbox}

\begin{proof}
$ $
\begin{enumerate}
\item 给定 $\zeta^n = 1$，单位根为 $\zeta = \omega^k$，其中 $\omega = e^{2\pi i/n}$ 且 $k = \{0, 1, \cdots, n-1\}$。
\item 根据定义，$\omega^k$ 是 $n$ 次本原单位根，当且仅当 $\textsf{ord}_{\mathbb{C}}(\omega^k) = n$。
\item $\omega$ 是 $n$ 次本原单位根，因为 $\textsf{ord}_{\mathbb{C}}(\omega) = n$。
\item 根据定理~\ref*{subsec:order-theorem}.2，若 $\textsf{ord}_{\mathbb{C}}(\omega) = n$，则有 $\textsf{ord}_{\mathbb{C}}(\omega^k) = \dfrac{n}{\text{gcd}(k,n)}$。因此，要使 $\textsf{ord}_{\mathbb{C}}(\omega^k) = n$ 成立，$\text{gcd}(k,n)$ 必须为 1。换言之，$k$ 与 $n$ 必须互素。
当 $n>1$ 时，允许的指数为满足 $1\le k<n$ 且与 $n$ 互素的整数，共有 $\phi(n)$ 个。当 $n=1$ 时，唯一的根为 $1$，其阶为 $1$，且 $\phi(1)=1$；此时允许指数 $k=0$。
\end{enumerate}
\end{proof}


\clearpage

\section{分圆多项式}
\label{sec:cyclotomic}
\textbf{- 参考资料：}
\href{https://e.math.cornell.edu/people/belk/numbertheory/CyclotomicPolynomials.pdf}{域与分圆多项式}~\cite{cyclotomic-polynomial}

\subsection{定义}
\label{subsec:cyclotomic-def}
\begin{tcolorbox}[title={\textbf{\tboxdef{\ref*{subsec:cyclotomic-def}} 分圆多项式}}]
 \textbf{第 $\bm{n}$ 个分圆多项式：}以所有 $n$ 次本原单位根为根的多项式，即：
 
 \[ \Phi_n(x) = \prod_{\zeta \in P(n)} (x - \zeta)  = \prod_{\substack{0 \leq k \leq n-1,\\ 
 \text{gcd}(k, n) = 1}} (x - \omega^k) \textcolor{white}{......} \text{，其中 } \omega = e^{2\pi i/n} \]

 回顾欧拉公式：$e^{2k\pi i/n} = \cos\left(\dfrac{2k\pi}{n}\right) + i\cdot\sin\left(\dfrac{2k\pi}{n}\right)$

$ $

一些预先计算出的分圆多项式如下：
 
\begin{multicols}{2}
$\Phi_1(x) = x - 1$
\newline $\Phi_2(x) = x + 1$
\newline $\Phi_3(x) = x^2 + x + 1$
\newline $\Phi_4(x) = x^2 + 1$
\newline $\Phi_5(x) = x^4 + x^3 + x^2 + x + 1$
\newline $\Phi_6(x) = x^2 - x + 1$
\newline $\Phi_7(x) = x^6 + x^5 + x^4 + x^3 + x^2 + x + 1$
\newline $\Phi_8(x) = x^4 + 1$
\newline $\Phi_9(x) = x^6 + x^3 + 1$
\newline $\Phi_{10}(x) = x^4 - x^3 + x^2 - x + 1$
\end{multicols}
\end{tcolorbox}


例如，

\[\Phi_4(x) = \prod_{\substack{0 \leq k \leq 3,\\ 
 \text{gcd}(k, 4) = 1}} (x - \omega^k) = (x - \omega^1)(x - \omega^3) = (x - e^{2\pi i / 4})(x - e^{2\cdot3\pi i / 4}) = (x - e^{\pi i / 2})(x - e^{3\pi i / 2})\]

$= \left(x - \left(\cos\left(\dfrac{\pi}{2}\right) + i \cdot \sin\left(\dfrac{\pi}{2}\right) \right) \right) \cdot \left(x - \left(\cos\left(\dfrac{3\pi}{2}\right) + i \cdot \sin\left(\dfrac{3\pi}{2}\right) \right) \right)$

$= (x - i)(x + i) = x^2 + 1$



\subsection{定理}
\label{subsec:cyclotomic-theorem}


\begin{tcolorbox}[title={\textbf{\tboxtheorem{\ref*{subsec:cyclotomic-theorem}.1} 第 $M$ 个分圆多项式的根}}]

假设 $M$ 是 2 的幂，并考虑第 $M$ 个分圆多项式 $\Phi_M(x) = x^n + 1$，其中 $M = 2n$。则第 $M$ 个分圆多项式的根为 $\omega, \omega^3, \omega^5, \cdots, \omega^{2n-1}$，其中 $\omega = e^{i\pi/n}$。

\end{tcolorbox}

\begin{proof}

根据 \autoref{subsec:cyclotomic-def} 中的定义~\ref*{subsec:cyclotomic-def}，$\Phi_M(x)$ 的根为 $e^{2k\pi i/M} = e^{2k\pi i/(2n)} = e^{k\pi i/n}$，其中 $0 \leq k < M = 2n$ 且 $\textsf{gcd}(k, M = 2n) = 1$，因此 $k = \{1, 3, 5, \cdots, 2n-1\}$。若令 $\omega = e^{i\pi/n}$，则 $\Phi_M(x)$ 的根为 $\omega, \omega^3, \omega^5, \cdots, \omega^{2n-1}$。

\end{proof}

\begin{tcolorbox}[title={\textbf{\tboxtheorem{\ref*{subsec:cyclotomic-theorem}.2} 将多项式分解为分圆多项式的乘积}}]
对于任意正整数 $n$，有 \[ x^n - 1 = \prod_{d \mid n} \Phi_d(x) \]
\end{tcolorbox}
\begin{myproof}
\begin{enumerate}
    \item $x^n - 1$ 的根是全部 $n$ 次单位根。因此，$x^n - 1 = (x - \omega^0)(x - \omega^1)...(x - \omega^{n-1})$，其中 $\zeta = \omega^k$。
    \item 定理~\ref*{subsec:roots-theorem}.3 表明，每个 $n$ 次单位根（$\omega^k$）都是某个 $d$ 次本原单位根，其中 $d$ 整除 $n$。换言之，每个 $n$ 次单位根都属于某个满足 $d \mid n$ 的 $P(d)$。另一方面，由定义有 $\Phi_d(x) = \Pi_{\zeta \in P(d)} (x - \zeta)$。因此，$x^n - 1$ 是所有满足 $d \mid n$ 的 $\Phi_d(x)$ 的乘积。
\end{enumerate}
\end{myproof}
\begin{tcolorbox}[title={\textbf{\tboxtheorem{\ref*{subsec:cyclotomic-theorem}.3} 分圆多项式的整数系数性质}}]
分圆多项式的所有系数都是整数。
\end{tcolorbox}
\begin{myproof}
对 $n$ 使用强归纳法。当 $n=1$ 时，$\Phi_1(x)=x-1$ 是首一整系数多项式。设 $n>1$，并假设结论对每个正整数 $d<n$ 成立，令
\[
F(x)=\prod_{\substack{d\mid n\\d<n}}\Phi_d(x).
\]
则 $F(x)$ 是 $\mathbb{Z}[x]$ 中的首一多项式。由定理~\ref*{subsec:cyclotomic-theorem}.2，在 $\mathbb{C}[x]$ 中有 $x^n-1=F(x)\Phi_n(x)$。以首一多项式 $F$ 作长除法只涉及整数运算，因此商和余式均属于 $\mathbb{Z}[x]$。由 $\mathbb{C}$ 上多项式除法的唯一性，商必为 $\Phi_n$，余式为零。因此 $\Phi_n\in\mathbb{Z}[x]$，归纳完成。
\end{myproof}
\begin{tcolorbox}[title={\textbf{\tboxtheorem{\ref*{subsec:cyclotomic-theorem}.4} $\bm{\Phi_{nk}(x)}$ 的公式}}]
若 $k \mid n$，则 $\Phi_{nk}(x) = \Phi_n(x^k)$。
\end{tcolorbox}
\begin{myproof}
对于满足 $k\mid n$ 的正整数，由定理~\ref*{subsec:order-theorem}.3，$\operatorname{ord}(\zeta)=nk$ 当且仅当 $\operatorname{ord}(\zeta^k)=n$。因此，$\Phi_n(x^k)$ 与 $\Phi_{nk}(x)$ 具有完全相同的根。这些根均为单根：$\Phi_n$ 的每个根 $\eta$ 都非零且为单根，而多项式 $x^k-\eta$ 的导数 $kx^{k-1}$ 在根处非零，所以它有 $k$ 个互异复根。两个多项式都首一且具有相同的单根，因此相等。映射 $\zeta\mapsto\zeta^k$ 在这些根集合上是 $k$ 对一，而非一对一；若对每个原像都乘入一次因子，就会错误地重复计算 $\Phi_n$ 的根。
\end{myproof}

\clearpage

\section{环上的单位根与分圆多项式}
\label{sec:cyclotomic-polynomial-integer-ring}
在 \autoref{sec:roots} 和 \autoref{sec:cyclotomic} 中，已经学习了复数上（即 $X \in \mathbb{C}$）的 $\mu$ 次单位根与第 $\mu$ 个分圆多项式的定义和性质，概括如下：

\begin{itemize}
\item \textbf{$\bm \mu$ 次单位根}是方程 $X^\mu = 1$ 在 $\mathbb{C}$ 中的解。换言之，所有 $\mu$ 次单位根都可以写成 $X = e^{2 \pi i k / \mu}$，其中整数 $k$ 满足 $0 \le k < \mu$。
\item \textbf{$\bm \mu$ 次本原单位根（记为 $\bm \omega$）}是阶（见 \autoref{subsec:order-def}）等于 $\mu$ 的 $\mu$ 次单位根，即满足 $\omega^{\mu} = 1$，且对任意满足 $d \mid \mu$ 的 $1 \le d < \mu$ 都有 $\omega^d \neq 1$。
\item 给定任意一个 $\mu$ 次本原单位根 $\omega$，计算 $\omega^{k'}$ 即可生成全部 $\mu$ 次本原单位根，其中整数 $k'$ 满足 $0 < k' < \mu$ 且 $\textsf{gcd}(k', \mu) = 1$（见 \autoref{subsec:order-theorem} 中的定理~\ref*{subsec:roots-theorem}.4）。
\item \textbf{第 $\bm \mu$ 个分圆多项式}定义为以全部 $\mu$ 次本原单位根为根的多项式，即 \[ \Phi_{\mu}(x) = \prod_{\omega \in P({\mu})} (x - \omega) = \prod_{\substack{0 \leq k \leq {\mu}-1,\\ \textsf{gcd}(k, {\mu}) = 1}} (x - \omega^k) \]
\end{itemize}

本节介绍 $\mathbb{Z}_p$ 上的第 $\mu$ 个分圆多项式，其中 $p$ 为素数，其结构如下：

\begin{tcolorbox}[title={\textbf{\tboxdef{\ref*{sec:cyclotomic-polynomial-integer-ring}} 环 $\mathbb{Z}_p$ 上的单位根与分圆多项式}}]


\begin{itemize}
\item \textbf{$\bm \mu$ 次单位根（记为 $\bm \omega$）}是 $X^\mu \equiv 1 \bmod p$ 的解。与复数情形不同，这些解不能写成 $X = e^{2 \pi i k / \mu}$。
\item \textbf{$\bm \mu$ 次本原单位根}定义为阶等于 $\mu$ 的 $\mu$ 次单位根，即满足 $\omega^{\mu} \equiv 1 \bmod p$，且对于任意满足 $d \mid \mu$ 的 $1 \le d < \mu$，都有 $\omega^d \not\equiv 1 \bmod p$。
\item 给定任意一个 $\mu$ 次本原单位根 $\omega$，计算 $\omega^{k'}$ 即可生成全部 $\mu$ 次本原单位根，其中整数 $k'$ 满足 $0 < k' < \mu$ 且 $\textsf{gcd}(k', \mu) = 1$。
\item \textbf{第 $\bm \mu$ 个分圆多项式}定义为以全部 $\mu$ 次本原单位根为根的多项式，即 \[ \Phi_{\mu}(x) = \prod_{\omega \in P({\mu})} (x - \omega) = \prod_{\substack{0 \leq k \leq {\mu}-1,\\ \textsf{gcd}(k, {\mu}) = 1}} (x - \omega^k) \]
\end{itemize}

\end{tcolorbox}

\begin{table}[h] %usepackage{array} 
\begin{tabular}{|c||c||c|}
\hline \hline
& \textbf{$\bm{\mathbb{C}}$ 上的多项式} & \textbf{$\bm{\mathbb{Z}}_{\bm{p}}$ 上的多项式} \\
& \textbf{（复数）} & \textbf{（环）} \\ \hline \hline
\textbf{单位根}&所有满足 $X^\mu = 1$ 的 $X \in \mathbb{C}$（可通过&所有满足 $X^\mu \equiv 1 \bmod p$ 的 $X \in \mathbb{Z}_p$\\
\textbf{的定义}&$X = e^{2 \pi i k / \mu}$ 计算，其中整数 $k$&\\
\textbf{（$\bm \mu$ 次）}&满足 $0 \leq k \leq \mu - 1$）&\\
\textbf{}&&\\\hline
\textbf{本原单位根}&满足以下条件的 $\mu$ 次单位根 $\omega$：&满足以下条件的 $\mu$ 次单位根 $\omega$：\\
\textbf{的定义}&$\omega^{\mu} = 1$，且 $\omega^{d} \neq 1$&$\omega^{\mu} \equiv 1 \bmod p$，且 $\omega^{d} \not\equiv 1 \bmod p$\\
\textbf{}&对任意满足 $d \mid \mu$ 的 $1 \le d < \mu$ 成立&对任意满足 $d \mid \mu$ 的 $1 \le d < \mu$ 成立\\
\textbf{（$\bm \mu$ 次）}&&\\
\textbf{}&&\\
\textbf{}&&\\\hline
\textbf{分圆多项式}&\multicolumn{2}{|c|}{以全部 $\mu$ 次本原单位根为根的多项式，如下所示：}\\
\textbf{的定义}&\multicolumn{2}{|c|}{$ \Phi_{\mu}(x) = \prod_{\omega \in P(\mu)} (x - \omega) $ \text{ }（见 \autoref{subsec:cyclotomic-def} 中的定义~\ref*{subsec:cyclotomic-def}）}\\
\textbf{（第 $\bm \mu$ 个）}&\multicolumn{2}{|c|}{}\\
\textbf{}&\multicolumn{2}{|c|}{}\\
\textbf{}&\multicolumn{2}{|c|}{}\\\hline
\textbf{求本原}&对于 $\omega = e^{2 \pi i/ \mu}$，计算满足以下条件的所有 ${\omega}^k$：&先找到第 $\mu$ 个分圆多项式的一个\\
\textbf{单位根}&$0 < k < \mu $ 且 $\textsf{gcd}(k, \mu) = 1$&本原根 $\omega$，然后\\
\textbf{（$\bm \mu$ 次）}&（见 \autoref{subsec:roots-theorem} 中的定理~\ref*{subsec:roots-theorem}.4）&计算满足以下条件的所有 $\omega^k \bmod p$：\\
\textbf{}&&$0 < k < \mu $ 且 $\textsf{gcd}(k, \mu) = 1$\\
\textbf{}&&\\\hline\hline
\end{tabular}
\caption{$X \in \mathbb{C}$ 上与 $X \in \mathbb{Z}_p$ 上的单位根和分圆多项式的比较。}
\label{tab:cyclotomic-polynomial-comparison}
\end{table}


对于第 $\mu$ 个分圆多项式，无论在 $X \in \mathbb{C}$ 上还是在 $X \in \mathbb{Z}_p$ 上，它的每个根 $\omega$（即 $\mu$ 次本原单位根）的阶均为 $\mu$，也就是说，在 $X \in \mathbb{C}$ 上有 $\omega^{\mu} = 1$，在 $X \in \mathbb{Z}_p$ 上有 $\omega^{\mu} \equiv 1 \bmod p$。此外，通过计算满足 $\textsf{gcd}(k', \mu) = 1$ 的 $\omega^{k'}$，可以由任意一个根 $\omega$ 生成第 $\mu$ 个分圆多项式的全部根。

\autoref{tab:cyclotomic-polynomial-comparison} 比较了 $\mathbb{C}$（复数）上与 $\mathbb{Z}_p$（环）上的单位根和第 $\mu$ 个分圆多项式的性质。

\clearpage

\section{拉格朗日多项式插值}
\label{sec:polynomial-interpolation}
假设给定 $n+1$ 个二维坐标点 $(x_0, y_0), (x_1, y_1), \cdots, (x_n, y_n)$，其中所有 $x$ 值互不相同，但 $y$ 值不一定互异。拉格朗日多项式插值是一种求出经过这 $n+1$ 个坐标点、次数至多为 $n$ 的唯一多项式的方法。给定的点 $(x_i,y_i)$ 可以位于 $\mathbb{C}^2$（包含实数的复平面），也可以位于某个素数 $p$ 所对应的 $\mathbb{Z}_p^2$ 中。


\begin{tcolorbox}[title={\textbf{\tboxtheorem{\ref*{sec:polynomial-interpolation}} 拉格朗日多项式插值}}]

假设给定 $n+1$ 个二维坐标点 $(x_0, y_0), (x_1, y_1), \cdots, (x_n, y_n)$，其中所有 $X$ 值互不相同，而 $Y$ 值可以重复。$(X, Y)$ 所在的域可以是 $(x_i, y_i) \in \mathbb{C}^2$（其中包含实数域），也可以是 $(x_i, y_i) \in \mathbb{Z}_p^2$（其中 $p$ 为素数）。则存在唯一的、次数至多为 $n$ 的多项式 $f(X)$ 经过这 $n+1$ 个坐标点。该多项式 $f(X)$ 的计算方式如下：

$f(X) = \sum\limits_{j=0}^{n}\dfrac{(X-x_0)\cdot(X-x_1)\cdots(X-x_{j-1})\cdot(X-x_{j+1})\cdots(X-x_{n})}{(x_j-x_0)\cdot(x_j-x_1)\cdots(x_j-x_{j-1})\cdot(x_j-x_{j+1})\cdots(x_j-x_{n})}\cdot y_j$

$\textcolor{white}{f(X) }= \sum\limits_{j=0}^{n} \left( \prod\limits_{\substack{0 \le k \le n\\ k \ne j}} \dfrac{X - x_k}{x_j - x_k} \cdot y_j\right)$

\end{tcolorbox}

\begin{myproof}
\begin{enumerate}


\item 首先，证明存在一个次数为 $n$ 或更低的多项式 $f(X)$，经过 $n+1$ 个互异坐标点 $(x_0, y_0), (x_1, y_1), \ldots, (x_n, y_n)$。构造多项式 $f(X)$ 如下：

$f(X) = \sum\limits_{j=0}^{n}\dfrac{(X-x_0)\cdot(X-x_1)\cdots(X-x_{j-1})\cdot(X-x_{j+1})\cdots(X-x_{n})}{(x_j-x_0)\cdot(x_j-x_1)\cdots(x_j-x_{j-1})\cdot(x_j-x_{j+1})\cdots(x_j-x_{n})}\cdot y_j$

$ \textcolor{white}{f(X) }= \sum\limits_{j=0}^{n} \left( \prod\limits_{\substack{0 \le k \le n\\ k \ne j}} \dfrac{X - x_k}{x_j - x_k} \cdot y_j\right)$

$ \textcolor{white}{f(X) }= \sum\limits_{j=0}^{n} \ell_j(X) \cdot y_j$ \textcolor{red}{ $\rhd$ 其中 $\ell_j(X) =  \prod\limits_{\substack{0 \le k \le n\\ k \ne j}} \dfrac{X - x_k}{x_j - x_k}$}

$ $

将 $\{\ell_0(X), \ell_1(X), \ldots, \ell_{n}(X) \}$ 称为 $\leq n$ 次多项式的拉格朗日基。根据 $f(X)$ 的构造，对于每个 $(x_i, y_i) \in \{(x_0, y_0), (x_1, y_1), \ldots, (x_n, y_n)\}$，\text{ }当 $i' = i$ 时有 $\ell_i(x_{i'}) = 1$，当 $i' \neq i$ 时有 $\ell_i(x_{i'}) = 0$。因此，对于 $0 \leq i \leq n$，有 $f(x_i) = \sum\limits_{j=0}^{n} \ell_j(x_i) \cdot y_j = 1 \cdot y_i = y_i$。换言之，$f(X)$ 经过 $n+1$ 个互异坐标点 $\{(x_0, y_0), (x_1, y_1), \ldots, (x_n, y_n)\}$。当 $(X, Y$ 所在的域为 $(x_i, y_i) \in \mathbb{C}^2$（即实数和复数）或 $(x_i, y_i) \in \mathbb{Z}_p^2$（其中 $p$ 为素数）时，均可计算出满足要求的 $f(X)$。特别地，在 $\bmod \, p$ 中也可以计算出有效的 $f(X)$，因为由定理~\ref*{subsec:order-theorem}.4（\autoref{subsec:order-theorem}）中的费马小定理可知，$a^{p - 1} \equiv 1 \bmod p$ 当且仅当 $a$ 与 $p$ 互素。这意味着，当 $p$ 为素数时，对所有 $a \in \mathbb{Z}_p^{\times}$（即去掉 $\{0\}$ 后的 $\mathbb{Z}_p$），都有 $a^{p - 1} \equiv 1 \bmod p$。由于 $\mathbb{Z}_p^{\times}$ 中的每个元素都有逆元，
所以在 $f(X) = \sum\limits_{j=0}^{n}\dfrac{(X-x_0)\cdot(X-x_1)\cdots(X-x_{j-1})\cdot(X-x_{j+1})\cdots(X-x_{n})}{(x_j-x_0)\cdot(x_j-x_1)\cdots(x_j-x_{j-1})\cdot(x_j-x_{j+1})\cdots(x_j-x_{n})}\cdot y_j$ 的公式中，每次除法都可以通过乘以相应分母的逆元来完成。
 
\item 接下来，证明不可能有两个互不相同、次数为 $n$ 或更低的多项式 $f_1(X)$ 与 $f_2(X)$，同时经过相同的 $n+1$ 个互异 $(X, Y)$ 坐标点。假设存在这样的两个多项式 $f_1(X)$ 和 $f_2(X)$，则 $f_{1-2}(X) = f_1(X) - f_2(X)$ 是一个次数为 $n$ 或更低、经过 $(x_0, 0), (x_1, 0), \cdots, (x_n, 0)$ 的新多项式。这意味着 $f_{1-2}(X)$ 有 $n+1$ 个互异的根，换言之，$f_{1-2}(X)$ 的次数至少为 $n+1$。这与 $f_{1-2}(X)$ 的次数至多为 $n$ 的假设矛盾。因此，不存在两个这样的多项式 $f_1(X)$ 与 $f_2(X)$，同时经过相同的 $n+1$ 个互异 $(X, Y)$ 坐标点。

\item 已经证明，存在一个次数为 $n$ 或更低的多项式 $f(X)$，经过 $n+1$ 个互异 $(X, Y)$ 坐标点，而且不存在两个或更多互异的多项式满足要求。因此，满足这一要求的多项式唯一。

\end{enumerate}
\end{myproof}



\subsection{插值唯一性的域条件}
\label{subsec:polynomial-interpolation-field-condition}


\begin{tcolorbox}[title={\textbf{\tboxtheorem{\ref*{subsec:polynomial-interpolation-field-condition}} 插值唯一性的域条件}}]

如果一个多项式定义在域 $\mathbb{F}$ 上
（即它的 $x$、$y$
坐标和系数均属于 $\mathbb{F}$），那么对于 $\mathbb{F}$ 中任意 $n+1$ 个互异的 $x$ 值和任意 $n+1$ 个可重复的 $y$ 值，都存在唯一一个次数至多为 $n$ 的多项式对它们进行插值。

$ $

这样的域包括 $\mathbb{R}$（实数域）、$\mathbb{C}$（复数域），以及任意有限域 $\mathbb{Z}_t$（其中 $t$ 为素数）。

\end{tcolorbox}


\begin{myproof}
\begin{enumerate}

\item 回顾定理~\ref*{sec:polynomial-interpolation}，拉格朗日多项式插值定义为：

$f(X) = \sum\limits_{j=0}^{n} \left( \prod\limits_{\substack{0 \le k \le n\\ k \ne j}} \dfrac{X - x_k}{x_j - x_k} \cdot y_j\right)$

$ $

如果多项式 $f(X)$ 的所有 $x$、$y$ 值和系数都属于某个域，那么 $(x_j - x_k)^{-1}$ 一定存在且仍属于该域。这是因为域中的每个非零元素都有逆元，而 $(x_j - x_k)$ 不为零（因为 $x_j \neq x_k$）。这说明 $\prod\limits_{\substack{0 \le k \le n\\ k \ne j}} \dfrac{X - x_k}{x_j - x_k} \cdot y_j$ 也属于同一个域，因此保证了 $f(X)$ 存在且定义良好。$f(X)$ 的次数至多为 $n$，因为其可能出现的最高次项为 $X^{n}$。

\item 对于 $f(X)$ 的唯一性，假设另一个次数至多为 $n$ 的多项式 $g(X)$ 也经过相同的 $n+1$ 个点。令 $h(X) = f(X) - g(X)$，则 $h(X)$ 的次数至多为 $n$，并且对于全部 $(n+1)$ 个互异节点 $x_0,\dots,x_{n}$，都有 $h(x_i) = 0$，因此 $h$ 有 $(n+1)$ 个互异的根。在域上，次数为 $d$ 的非零多项式至多有 $d$ 个根；次数至多为 $n$ 的非零多项式 $h(X)$ 至多有 $n$ 个根。由于 $h(X)$ 有 $(n+1)$ 个根，$h(X)$ 必须是零多项式，即 $f(X) = g(X)$。因此，$f(X)$ 是唯一的。


\end{enumerate}
\end{myproof}

\clearpage

\section{向量与矩阵}
\label{sec:matrix}


\subsection{向量运算}
\label{subsec:vector-arithmetic}

本节介绍向量和矩阵计算的基本运算，以及向量投影、平面投影、平面或空间的基等进一步的概念。


\begin{tcolorbox}[title={\textbf{\tboxdef{\ref*{subsec:vector-arithmetic}} 向量运算}}]

\begin{itemize}
\item \textbf{加法：}给定两个 $n \times 1$ 向量（即 $n$ 维向量），每个向量由 $n$ 个数组成：

$\vec{a} = (a_0, a_1,\cdots, a_{n-1})$, $\vec{b} = (b_0, b_1,\cdots, b_{n-1})$

$ $

向量加法定义为：

$ \vec{a} + \vec{b} = (a_0 + b_0, \text{ } a_1  +  b_1, \text{ }\cdots, \text{ } a_{n-1}  +  b_{n-1})$

$ $

\item \textbf{点积：}给定两个 $n$ 维向量：

$\vec{a} = (a_0, a_1,\ldots, a_{n-1})$, $\vec{b} = (b_0, b_1,\ldots, b_{n-1})$

$ $

点积定义为：

$\langle \vec{a}, \vec{b} \rangle =  \vec{a} \cdot \vec{b} = a_0 b_0 + a_1 b_1 +\cdots + a_{n-1} b_{n-1}$.

$\vec{a} \cdot \vec{b}$ 也可以表示为 $\|\vec{a}\|\,\|\vec{b}\|\cos\theta$，其中 $\theta$ 是 $\vec{a}$ 与 $\vec{b}$ 的夹角。当 $\vec{a}$ 和 $\vec{b}$ 指向同一方向时，$\vec{a}\cdot\vec{b}$ 取得最大值（$\|\vec{a}\|\,\|\vec{b}\|$）。当 $\vec{a}$ 与 $\vec{b}$ 正交时，$\vec{a}\cdot\vec{b} = 0$。

$ $


\item \textbf{Hadamard 积：}给定两个 $n$ 维向量：

$\vec{a} = (a_0, a_1,\ldots, a_{n-1})$, $\vec{b} = (b_0, b_1,\ldots, b_{n-1})$

$ $

Hadamard 积定义为：

$\vec{a} \odot \vec{b} = (a_0 b_0, \text{ } a_1 b_1, \text{ }\cdots, \text{ } a_{n-1} b_{n-1})$

$ $

\item \textbf{Hermitian 内积：}给定两个 $n$ 维复向量：

$\vec{a} = (a_0 + i \cdot a'_0, \text{ } a_1+ i\cdot a'_1, \text{ }\ldots,\text{ }  a_{n-1}+ i\cdot a'_{n-1})$ 

$\vec{b} = (b_0 + i\cdot b'_0, \text{ } b_1 + i\cdot b'_1, \text{ }\ldots, \text{ } b_{n-1} + i\cdot b'_{n-1})$

$ $

Hermitian 内积是对第二个操作数取共轭（见 \autoref{subsec:imaginary}）后的点积：

$\langle\langle \vec{a}, \vec{b} \rangle\rangle =  \vec{a} \cdot \overline{\vec{b}} $

$ = (a_0 + i \cdot a'_0, \text{ } a_1+ i\cdot a'_1, \text{ }\cdots,\text{ } a_{n-1}+ i\cdot a'_{n-1})  \cdot (b_0 - i\cdot b'_0, \text{ } b_1 - i\cdot b'_1, \text{ }\cdots, \text{ } b_{n-1} - i\cdot b'_{n-1})$




\end{itemize}
\end{tcolorbox}




\subsection{不同类型的矩阵}
\label{subsec:vandermonde}


\begin{tcolorbox}[title={\textbf{\tboxdef{\ref*{subsec:vandermonde}} 矩阵}}]


\begin{itemize}

\item $n \times n$ 的\textbf{单位矩阵}和\textbf{反序单位矩阵}定义为：

$I_n = \begin{bmatrix}
1 & 0 & 0 & \cdots & 0\\
0 & 1 & 0 & \cdots & 0 \\
0 & 0 & 1 & \cdots & 0 \\
\vdots & \vdots & \vdots & \ddots & \vdots \\
0 & 0 & 0 & \cdots & 1 \\
\end{bmatrix}$, $I^R_n = \begin{bmatrix}
0 & \cdots & 0 & 0 & 1\\
0 & \cdots & 0 & 1 & 0 \\
0 & \cdots & 1 & 0 & 0 \\
\vdots & \iddots & \vdots & \vdots & \vdots \\
1 & 0 & 0 & \cdots & 0 \\
\end{bmatrix}$

$ $

\item 矩阵 $X$ 的\textbf{转置}是将各元素关于主对角线对调，记为 $X^T$，即：


$X = \begin{bmatrix}
a_1 & a_2 & a_3 & \cdots & a_n\\
b_1 & b_2 & b_3 & \cdots & b_n \\
c_1 & c_2 & c_3 & \cdots & c_n \\
\vdots & \vdots & \vdots & \ddots & \vdots \\
m_1 & m_2 & m_3 & \cdots & m_n \\
\end{bmatrix}$, $X^T = \begin{bmatrix}
a_1 & b_1 & c_1 & \cdots & m_1\\
a_2 & b_2 & c_2 & \cdots & m_2\\
a_3 & b_3 & c_3 & \cdots & m_3\\
\vdots & \vdots & \vdots & \ddots & \vdots \\
a_n & b_n & c_n & \cdots & m_n\\
\end{bmatrix}$

$ $


\item \textbf{Vandermonde 矩阵}是一个 $(m + 1) \times (n + 1)$ 矩阵，定义为：

$\mathit{Vander}(x_0, x_1, \cdots, x_{m}) = \begin{bmatrix}
1 & x_0 & x_0^2 & \cdots & x_0^n\\
1 & x_1 & x_1^2 & \cdots & x_1^n \\
1 & x_2 & x_2^2 & \cdots & x_2^n \\
\vdots & \vdots & \vdots & \ddots & \vdots \\
1 & x_m & x_m^2 & \cdots & x_m^{n} \\
\end{bmatrix}$

%A Vandermonde matrix is often used for representing $m$ different $n$-degree polynomials.

\end{itemize}

\end{tcolorbox}


 





\subsection{矩阵运算}
\label{subsec:matrix-arithmetic}


矩阵与向量相乘、矩阵与矩阵相乘分别定义如下：




\begin{tcolorbox}[title={\textbf{\tboxdef{\ref*{subsec:matrix-arithmetic}} 矩阵运算}}]

\begin{itemize}
\item \textbf{矩阵与向量相乘：}给定 $m \times n$ 矩阵 $A$ 和 $n$ 维向量 $x$：


$A = \begin{bmatrix}
 a_{1, 1} & a_{1, 2} & a_{1, 3} & \cdots & a_{1, n}\\
 a_{2, 1} & a_{2, 2} & a_{2, 3} & \cdots & a_{2, n} \\
 a_{3, 1} & a_{3, 2} & a_{3, 3} & \cdots & a_{3, n} \\
\vdots & \vdots & \vdots & \ddots & \vdots \\
 a_{m, 1} & a_{m, 2} & a_{m, 3} & \cdots & a_{m, n} \\
\end{bmatrix} = \begin{bmatrix} 
\vec{a}_{1, * } \\ 
\vec{a}_{2, * } \\ 
\vec{a}_{3, * } \\ 
\vdots\\
\vec{a}_{m, * } \\ 
\end{bmatrix}, \text{ } \vec{x} = (x_1, x_2, \cdots, x_{n})$


$ $

$A \cdot \vec{x}$ 的结果是一个 $m$ 维向量，计算如下：

$A \cdot \vec{x} = \Big(\vec{a}_{1, * } \cdot \vec{x}, \text{ } \vec{a}_{2, * } \cdot \vec{x}, \text{ }\cdots,\text{ } \vec{a}_{m, * } \cdot \vec{x} \Big) = \left(\sum\limits_{i=1}^{n} a_{1, i} \cdot  x_i, \text{ } \sum\limits_{i=1}^{n}  a_{2, i} \cdot x_i,\cdots, \text{ } \sum\limits_{i=1}^{n} a_{m, i} \cdot x_i \right)$


$ $


\item \textbf{矩阵与矩阵相乘：}给定 $m \times n$ 矩阵 $A$ 和 $n \times k$ 矩阵 $B$：



$A = \begin{bmatrix}
 a_{1, 1} & a_{1, 2} & a_{1, 3} & \cdots & a_{1, n}\\
 a_{2, 1} & a_{2, 2} & a_{2, 3} & \cdots & a_{2, n} \\
 a_{3, 1} & a_{3, 2} & a_{3, 3} & \cdots & a_{3, n} \\
\vdots & \vdots & \vdots & \ddots & \vdots \\
 a_{m, 1} & a_{m, 2} & a_{m, 3} & \cdots & a_{m, n} \\
\end{bmatrix}$, $B = \begin{bmatrix}
 b_{1, 1} & b_{1, 2} & b_{1, 3} & \cdots & b_{1, k}\\
 b_{2, 1} & b_{2, 2} & b_{2, 3} & \cdots & b_{2, k} \\
 b_{3, 1} & b_{3, 2} & b_{3, 3} & \cdots & b_{3, k} \\
\vdots & \vdots & \vdots & \ddots & \vdots \\
 b_{n, 1} & b_{n, 2} & b_{n, 3} & \cdots & b_{n, k} \\
\end{bmatrix}$

$ $

$A \cdot B$ 的结果是一个 $m \times k$ 矩阵，计算如下：

$A \cdot B = \begin{bmatrix}
\sum\limits_{i=1}^{n}a_{1, i}b_{i, 1} & \sum\limits_{i=1}^{n}a_{1, i}b_{i, 2} & \sum\limits_{i=1}^{n}a_{1, i}b_{i, 3} & \cdots & \sum\limits_{i=1}^{n}a_{1, i}b_{i, k} \\
\sum\limits_{i=1}^{n}a_{2, i}b_{i, 1} & \sum\limits_{i=1}^{n}a_{2, i}b_{i, 2} & \sum\limits_{i=1}^{n}a_{2, i}b_{i, 3} & \cdots & \sum\limits_{i=1}^{n}a_{2, i}b_{i, k} \\
\sum\limits_{i=1}^{n}a_{3, i}b_{i, 1} & \sum\limits_{i=1}^{n}a_{3, i}b_{i, 2} & \sum\limits_{i=1}^{n}a_{3, i}b_{i, 3} & \cdots & \sum\limits_{i=1}^{n}a_{3, i}b_{i, k} \\
\vdots & \vdots & \vdots & \ddots & \vdots \\
\sum\limits_{i=1}^{n}a_{m, i}b_{i, 1}
& \sum\limits_{i=1}^{n}a_{m, i}b_{i, 2} & \sum\limits_{i=1}^{n}a_{m, i}b_{i, 3} & \cdots & \sum\limits_{i=1}^{n}a_{m, i}b_{i, k} \\
\end{bmatrix}$


\end{itemize}
\end{tcolorbox}

根据上述矩阵与向量运算的定义，可以得到以下代数性质：

\begin{tcolorbox}[title={\textbf{\tboxtheorem{\ref*{subsec:matrix-arithmetic}} 矩阵运算的性质}}]

\begin{itemize}

\item \textbf{结合律：}

$(AB)C = A(BC)$

$A(Bx) = (AB)x$


$ $

 
\item \textbf{分配律：}

$A(x + y) = Ax + Ay$

%$A(x \odot y) = Ax \odot Ay $

%$A\langle x, y\rangle = \langle Ax, Ay\rangle$

%$A\langle \langle x, y\rangle\rangle = \langle\langle Ax, Ay\rangle\rangle$

$ $

但是，$Ax \cdot Ay \neq A(x \cdot y)$，因为结果的维数不匹配。另外，$A(x \odot y) \neq Ax \odot Ay $。还需注意 $A\langle x, y\rangle \neq \langle Ax, Ay\rangle$，以及 $A\langle \langle x, y\rangle\rangle \neq \langle\langle Ax, Ay\rangle\rangle$。

$ $

\item \textbf{不满足交换律：}

$Ax \neq xA$

$AB \neq BA$


\end{itemize}
\end{tcolorbox}

\begin{proof}

$ $

可以为向量或矩阵的每个元素引入变量，按照运算定义展开各等式的左右两边，再比较得到的表达式，从而验证 \ref*{subsec:matrix-arithmetic} 所述的性质。具体展开留给读者作为练习。

\end{proof}


\subsection{投影}
\label{subsec:projection}

这里介绍两类投影：向量投影和正交投影（即平面投影）。



\begin{figure}[!h]
    \centering
    \subfloat[向量投影 \href{https://upload.wikimedia.org/wikipedia/commons/thumb/9/98/Projection_and_rejection.png/200px-Projection_and_rejection.png}{（来源）}]{
  \includegraphics[width=0.25\linewidth]{figures/vector-projection.png}

  \label{fig:vector-projection}
}
    \subfloat[正交投影（平面／子空间投影）\href{https://www.researchgate.net/publication/333616378/figure/fig2/AS:766235117617152@1559696096437/Geometric-illustration-of-the-orthogonal-projection-operator-P-A-vector-x-R-m-is.png}{（来源）}]{
  \includegraphics[width=0.4\linewidth]{figures/orthogonal-projection.png}

  \label{fig:orthogonal-projection}
}
  \caption{}
\end{figure}


\para{向量投影：}给定同一 $n$ 维向量空间中的两个向量 $\vec{a}$ 和 $\vec{b}$，向量投影 $\textsf{Proj}_{\vec{b}}(\vec{a})$ 描述 $\vec{a}$ 在 $\vec{b}$ 方向上的分量，即 $\vec{a}$ 中平行于 $\vec{b}$ 的部分。在 \autoref{fig:vector-projection} 的例子中，$\vec{a}$ 在 $\vec{b}$ 上的投影是 $\vec{a}_1$，其长度 $\vec{a}_1$ 在几何上等于 $||a_1|| = ||a|| \cos \theta = ||a||\dfrac{a \cdot b}{||a||\cdot||b||} = \dfrac{a \cdot b}{||b||}$。设 $\vec{b'}$ 为 $\vec{b}$ 方向上的单位向量，即 $\vec{b'} = \dfrac{\vec{b}}{||b||}$。于是 $\vec{a}_1 = ||a_1||\cdot\vec{b'} = \dfrac{a \cdot b}{||b||} \cdot \dfrac{\vec{b}}{||b||} = \dfrac{a \cdot b}{||b||^2}\vec{b}$，从而 $\textsf{Proj}_{\vec{b}}(\vec{a}) = \dfrac{a \cdot b}{||b||^2}\vec{b}$。

\para{正交投影：}给定向量 $\vec{x}$，以及张成子空间 $P$ 的一组两两正交向量 $\vec{p}_0, \vec{p}_1, \cdots, \vec{p}_{n-1}$，正交投影 $\textsf{Proj}_P(\vec{x})$ 描述 $\vec{x}$ 落在 $P$ 内的部分，即 $\vec{x}$ 在 $P$ 中的分量。在图示例子中，$\vec{x}$ 在 $P$ 上的投影（红色箭头），等于 $\vec{x}$ 在张成 $P$ 的各正交基向量 $\vec{p}_i$ 上的投影之和：$\textsf{Proj}_P(\vec{x}) = \sum\limits_{i=0}^{n-1}\textsf{Proj}_{\vec{p}_i}(\vec{x})$。这一计算可以理解为在由 $n$ 个正交向量定义的坐标系中表示 $\vec{x}$。

\begin{tcolorbox}[title={\textbf{\tboxdef{\ref*{subsec:projection}} 向量投影与正交投影}}]

\begin{itemize}

\item \textbf{向量投影：}给定同一向量空间中的两个向量 $\vec{a}$ 和 $\vec{b}$，$\vec{a}$ 在 $\vec{b}$ 上的向量投影为：

$\textsf{Proj}_{\vec{b}}(\vec{a}) = \vec{a}_p  = \dfrac{\vec{a} \cdot \vec{b}}{||b||^2} \vec{b}$ \textcolor{red}{ $\rhd$ 其中 $||\vec{a}_p|| = ||\vec{a}|| \cos \theta$}

$ $

\item \textbf{正交基：}如果 $n$ 维平面（或子空间）$P$ 由两两正交的 $n$ 维向量 $\vec{p}_0, \vec{p}_1, \cdots, \vec{p}_{n-1}$ 张成，

则将矩阵 $P = \begin{bmatrix} 
\vec{p}_0\\
\vec{p}_1\\
\vdots\\
\vec{p}_{n-1}\\
\end{bmatrix}
$
定义为平面 $P$ 的一组正交基。


\item \textbf{正交投影：}给定正交基矩阵 $P = \begin{bmatrix} 
\vec{p}_0\\
\vec{p}_1\\
\vdots\\
\vec{p}_{n-1}\\
\end{bmatrix}
$，

向量 $\vec{a}$ 在 $P$ 上的正交投影为：

$\textsf{Proj}_P(\vec{a}) = \sum\limits_{i=0}^{n-1} \textsf{Proj}_{\vec{p}_i}(\vec{a})$



\end{itemize}


\end{tcolorbox}


根据正交投影的定义，可以推导出以下性质：

\para{正交基：}在 $n$ 维向量空间中，任意 $n$ 个两两正交的向量所张成的子空间 $P$ 就是整个向量空间。此外，该向量空间中的任意向量在 $P$ 上的正交投影都是唯一确定的向量。

\para{非正交基：}在 $n$ 维向量空间中，假设有 $n$ 个非正交向量满足两个条件：(i) 它们张成整个向量空间；(ii) 它们线性无关，即其中任何一个向量都不能表示为其他向量的线性组合。那么，由这 $n$ 个向量构成的 $n \times n$ 矩阵 $P$ 组成整个向量空间 $V$ 的一组基。对于空间中的每个 $\vec{v}$，矩阵与向量的乘法 $P\vec{v}$ 都得到唯一确定的向量。不过，由于这 $n$ 个基向量并不正交，表达式 $\textsf{Proj}_P(\vec{v})$ 并不是向量 $\vec{v}$ 在 $P$ 上的有效几何投影。尽管如此，计算 $P\vec{v}$ 仍可理解为将 $\vec{v}$ 唯一地变换到另一个坐标系中：该坐标系使用 $P$ 中的 $n$ 个非正交向量描述这个向量空间。

\begin{tcolorbox}[title={\textbf{\tboxtheorem{\ref*{subsec:projection}} 变换后向量的唯一性}}]

\begin{itemize}
\item \textbf{正交基：}如果 $n$ 个向量是 $n$ 维向量空间中平面 $P$ 的一组正交基，则 $P$ 等于整个向量空间，且对于空间中的每个向量 $\vec{v}$，$\textsf{Proj}_P(v)$ 都是唯一确定的向量。
\item \textbf{非正交基：}如果 $n$ 个向量是 $n$ 维向量空间中平面 $P$ 的一组非正交基，即这些向量线性无关并张成 $P$，则 $P$ 等于整个向量空间，并且 $P\vec{v}$ 的结果是唯一确定的向量。
\end{itemize}
\end{tcolorbox}




\subsection{多项式环的基}
\label{subsec:polynomial-ring-basis}

给定 $(n-1)$ 次多项式环 $\mathbb{Z}[X] / (X^n + 1)$，该环的一组基定义为满足以下两个要求的多项式集合：

\begin{itemize}

\item \textbf{线性无关：}基中的任何一个多项式，都不能表示为同一组基中其他多项式的线性组合。
\item \textbf{张成多项式环：}基中多项式的线性组合可以表示该环中的任意多项式。
\end{itemize}

对于 $(n-1)$ 次多项式环，构成其一组基的多项式个数恰为 $n$。


\subsection{整数（环）上多项式与向量之间的同构}
\label{subsec:poly-vector-transformation}



现在，定义从 $(n-1)$ 次多项式环到 $n$ 维向量空间的映射 $\sigma$：将输入多项式在 $n$ 个互异的 $X \in \mathbb{Z}_t$ 坐标处（例如 $x_0, x_1, \cdots, x_{n-1}$）求得的 $y$ 值依次排列，构成输出向量。这里要求 $\mathbb{Z}_t$ 为域，即 $t$ 为素数，以保证经过 $(n+1)$ 个互异 $X$ 点的插值多项式 $f(X)$ 定义良好，如 \autoref{subsec:polynomial-interpolation-field-condition} 所述。具体地，$\sigma$ 定义为：

$\sigma: f(X) \in \mathbb{Z}_t[X] / (X^n + 1) \text{ } \longrightarrow \text{ } (f(x_0), f(x_1), f(x_2), \cdots, f(x_{n-1})) \in \mathbb{Z}^n_t$

$ $

下面说明映射 $\sigma$ 为何是同构，即 $\sigma$ 是从 $\mathbb{Z}_t[X] / (X^n + 1)$ 到 $\mathbb{Z}^n_t$ 的一一对应双射，并且保持代数运算 $(+, \cdot)$，也就是说，$\sigma$ 对加法和乘法均为同态。

$ $

\para{双射：}在 $(n-1)$ 次多项式环中，将多项式在预先固定的 $n$ 个互异 $x$ 坐标处求值得到的一组 $y$ 值，可以唯一确定该多项式。其原因是，从代数上看，经过给定的 $n$ 个互异 $(x, y)$ 坐标点，只有一个次数为 $(n-1)$ 或更低的多项式。这一点已由拉格朗日多项式插值证明（见 \autoref{sec:polynomial-interpolation} 中的定理~\ref*{sec:polynomial-interpolation}）。

$ $

\para{同态：}映射 $\sigma$ 对运算 $(+, \cdot)$ 的同态性，意味着以下两个关系成立：

$\sigma(f_a(X) + f_b(X)) = \sigma(f_a(X)) + \sigma(f_b(X))$

$\sigma(f_a(X) \cdot f_b(X)) = \sigma(f_a(X)) \odot \sigma(f_b(X))$ \textcolor{red}{\text{ } \# $\odot$ 是向量的 Hadamard 积（见总结~\ref*{subsec:vector-arithmetic}）}

$ $

为证明映射 $\sigma$ 的同态性，将输入多项式 $f_a(X)$、$f_b(X)$ 以及经 $\sigma$ 映射得到的输出向量记为：


$f_a(X) = a_0 + a_1X + a_2X^2 + \cdots + a_{n-1}X^{n-1}$

$ \sigma(f_a(X)) = \left({f_a(x_0)}, {f_a(x_1)}, {f_a(x_2)}, \cdots, {f_a(x_{n-1}}) \right)  = \left(\sum\limits_{i=0}^{n-1}a_i(x_0)^i, \sum\limits_{i=0}^{n-1}a_i(x_1)^i, \sum\limits_{i=0}^{n-1}a_i(x_2)^i, \cdots, \sum\limits_{i=0}^{n-1}a_i(x_{n-1})^i\right)$


$ $

$f_b(X) = b_0 + b_1X + b_2X^2 + \cdots + b_{n-1}X^{n-1}$

$ \sigma(f_b(X)) = \left({f_b(x_0)}, {f_b(x_1)}, {f_b(x_2)}, \cdots, {f_b(x_{n-1}}) \right) = \left(\sum\limits_{i=0}^{n-1}b_i(x_0)^i, \sum\limits_{i=0}^{n-1}b_i(x_1)^i, \sum\limits_{i=0}^{n-1}b_i(x_2)^i, \cdots, \sum\limits_{i=0}^{n-1}b_i(x_{n-1})^i\right)$

$ $

在上述设定下，可以看出 $\sigma$ 对运算 $(+)$ 保持同态，如下所示：

$\bm{\sigma}\bm{(}f_a(X) + f_b(X)\bm{)} = \bm{\sigma\bigg(}(a_0 + b_0) + (a_1 + b_1)X + (a_2 + b_2)X^2 + \cdots + (a_{n-1} + b_{n-1})X^{n-1}\bm{\bigg)}$

$= \left(\sum\limits_{i=0}^{n-1}(a_i + b_i)(x_0)^i, \text{ }\sum\limits_{i=0}^{n-1}(a_i + b_i)(x_1)^i, \text{ }\sum\limits_{i=0}^{n-1}(a_i + b_i)(x_2)^i, \cdots, \text{ }\sum\limits_{i=0}^{n-1}(a_i + b_i)(x_{n-1})^i\right)$

$= \left(\sum\limits_{i=0}^{n-1}a_i(x_0)^i, \text{ }\sum\limits_{i=0}^{n-1}a_i(x_1)^i, \text{ }\sum\limits_{i=0}^{n-1}a_i(x_2)^i, \cdots, \text{ } \sum\limits_{i=0}^{n-1}a_i(x_{n-1})^i\right)$

\text{ } $+ \left(\sum\limits_{i=0}^{n-1}b_i(x_0)^i, \text{ } \sum\limits_{i=0}^{n-1}b_i(x_1)^i, \text{ } \sum\limits_{i=0}^{n-1}b_i(x_2)^i, \cdots, \text{ } \sum\limits_{i=0}^{n-1}b_i(x_{n-1})^i\right)$

$= \bm{\sigma(}f_a(X)\bm{)} + \bm{\sigma(}f_b(X)\bm{)}$

$ $

同样，$\sigma$ 对运算 $(\cdot)$ 也保持同态，如下所示：


$\bm{\sigma(}f_a(X)\cdot f_b(X)\bm{)} = \bm{\sigma\bigg(}\left(\sum\limits_{i=0}^{n-1}a_iX^i\right)\cdot\left( \sum\limits_{i=0}^{n-1}b_iX^i\right)\bm{\bigg)}$

$ $

$= \Bigg( \left(\sum\limits_{i=0}^{n-1}a_ix_0^i\right)\left( \sum\limits_{i=0}^{n-1}b_ix_0^i\right), \text{ }
\left(\sum\limits_{i=0}^{n-1}a_ix_1^i\right)\left( \sum\limits_{i=0}^{n-1}b_ix_1^i\right), \text{ }
\left(\sum\limits_{i=0}^{n-1}a_ix_2^i\right)\left( \sum\limits_{i=0}^{n-1}b_ix_2^i\right),$

$  \text{ } \text{ } \text{ } \text{ } \text{ } \text{ }
\cdots,
\left(\sum\limits_{i=0}^{n-1}a_ix_{n-1}^i\right)\left( \sum\limits_{i=0}^{n-1}b_ix_{n-1}^i\right)
\Bigg)$

$ $

$ = \Bigg(\sum\limits_{i=0}^{n-1}a_i(x_0)^i, \text{ } \sum\limits_{i=0}^{n-1}a_i(x_1)^i, \text{ } \cdots, \text{ } \sum\limits_{i=0}^{n-1}a_i(x_{n-1})^i\Bigg)$

$  \text{ } \text{ } \text{ } \text{ }  \odot  \Bigg(\sum\limits_{i=0}^{n-1}b_i(x_0)^i, \text{ } 
\sum\limits_{i=0}^{n-1}b_i(x_1)^i, \text{ }
\cdots, \text{ } \sum\limits_{i=0}^{n-1}b_i(x_{n-1})^i\Bigg)$

$ $

$= \bm{\sigma(}f_a(X)\bm{)} \odot \bm{\sigma(}f_b(X)\bm{)}$

$ $

因此，$\sigma$ 保持以下同态关系：

$\sigma(f_a(X) + f_b(X)) = \sigma(f_a(X)) + \sigma(f_b(X))$

$\sigma(f_a(X) \cdot f_b(X)) = \sigma(f_a(X)) \odot \sigma(f_b(X))$

$ $

不过，对于 $\sigma(f_a(X) \cdot f_b(X)) = \sigma(f_a(X)) \odot \sigma(f_b(X))$，还需要进一步说明为什么该关系在多项式环中成立，具体如下。


\para{多项式环中的约减：}假设没有多项式环 $X^n + 1$ 的设定，那么将 $f_a(X)$ 和 $f_b(X)$ 相乘，得到的 $f_a(X) \cdot f_b(X)$ 可能是一个次数高于 $n-1$ 的新多项式。这个较高次的多项式仍可解码为预期的正确向量。考虑以下情形：

$\sigma(f_a(X)) = (f_a(x_0), f_a(x_1), \cdots, f_a(x_{n-1})) = (v_0, v_1, \cdots, v_{n-1})$

$\sigma(f_b(X)) = (f_b(x_0), f_b(x_1), \cdots, f_b(x_{n-1})) = (u_0, u_1, \cdots, u_{n-1})$

$ $

此时，有：

$\sigma(f_a(X) \cdot f_b(X)) = (f_a(x_0)\cdot f_b(x_0), f_a(x_1)\cdot f_b(x_1), \cdots, f_a(x_{n-1})\cdot f_b(x_{n-1})) = (v_0u_0, v_1u_1, \cdots, v_{n-1}u_{n-1})$

$= (v_0, v_1, \cdots, v_{n-1}) \odot (u_0, u_1, \cdots, u_{n-1}) $

$ $

如上所示，即使 $f_a(X) \cdot f_b(X)$ 得到的多项式次数高于 $n-1$，仍可将其解码为预期的正确向量。然而，映射 $\sigma$ 会失去多项式与向量之间的同构性，因为当多项式次数高于 $n-1$ 时，经过给定的 $n$ 个互异 $X$ 坐标点 $\{x_0, x_1, \cdots, x_{n-1}\}$ 的多项式可能不止一个。这会带来问题：如果 $\sigma$ 只支持从多项式到向量的映射，而不支持从向量到多项式的映射，就无法在一开始将向量转换为多项式并进行同构计算。另一个问题是，多项式次数高于 $n-1$ 时，解码（即多项式求值）的计算开销也会增大。

为解决这两个问题，可以将 $n$ 个互异求值坐标 $X$ 选为多项式环模数 $X^n + 1$ 的根，其中 $n$ 是 2 的某个幂；再将 $f_a(X) \cdot f_b(X)$ 模 $X^n + 1$ 约减为次数至多为 $n - 1$ 的新多项式。设 $f_{ab}(X) = f_a(X) \cdot f_b(X)$，并令 $f'_{ab}(X)$ 为约减后的多项式，使得对于某个商多项式 $Q(X)$ 有 $f_{ab}(X) = Q(X)\cdot (X^n + 1) + f'_{ab}(X)$。根据总结~\ref*{subsubsec:polynomial-ring-discuss}（\autoref{subsubsec:polynomial-ring-discuss}），如果在 $X^n + 1$ 的根处求值，$Q(X)$ 项为零，因此 $f_{ab}(X)$ 和 $f'_{ab}(X)$ 的求值结果相同。所以，只要将 $n$ 个互异求值点 $\{x_0, x_1, \cdots, x_{n-1}\}$ 选为 $X^n + 1$ 的根，就能保证 $f'_{ab}(X)$ 的解码向量与预期的 $f_{ab}(X)$ 的解码向量完全相同。于是可以用约减后的多项式 $f'_{ab}(X)$ 替代高次多项式 $f_{ab}(X)$，并继续使用 $f'_{ab}(X)$ 进行后续多项式加法或乘法。多项式环中的约减既能提高多项式加法和乘法的效率，也能保持映射 $\sigma$ 的同构性，因此可以自由地在向量和多项式之间转换，并进行加法和乘法运算。

要应用上述多项式环约减，模多项式可以选择为任何具有至少 $n$ 个互异根的多项式。在实际应用中，常选择 $X^n + 1$ 作为多项式环的模数，它是第 $(\mu=2n)$ 个分圆多项式（见 \autoref{subsec:cyclotomic-def}）。选择分圆多项式，特别是第 $(\mu=2n)$ 个分圆多项式 $X^n + 1$ 作为模数，是因为其 $n$ 个互异的根（即 $(\mu=2n)$ 次本原单位根）有明确的表达，即使 $n$ 很大也能快速找到。

$ $

\para{多项式系数的模约减：}此外，为提高计算效率，常对多项式系数按某个模数 $t$ 约减，将系数大小限制在给定范围内。假设两个多项式 $f_c(X)$ 和 $f_d(X)$ 的系数模 $t$ 同余：

$f_c(X) = w_0 + w_1\cdot X + w_2\cdot X^2 + \cdots + w_{n-1}\cdot X^{n-1}$

$f_d(X) = w'_0 + w'_1\cdot X + w'_2\cdot X^2 + \cdots + w'_{n-1}\cdot X^{n-1} $

$w_i \equiv w'_i \bmod t$

$ $

则对于任意 $x \in \mathbb{Z}$，其求值结果 $f_c(x)$ 与 $f_d(x)$ 一定模 $t$ 同余，如下所示：

$f_c(x) = w_0 + w_1\cdot x + w_2\cdot x^2 + \cdots + w_{n-1}\cdot x^{n-1}$

$\equiv (w_0 + w_1\cdot x + w_2\cdot x^2 + \cdots + w_{n-1}\cdot x^{n-1}) \bmod t $

$\equiv (w_0  \bmod t) + (w_1  \bmod t)\cdot x + (w_2  \bmod t)\cdot x^2 + \cdots + (w_{n-1} \bmod t)\cdot x^{n-1} $

$\equiv w'_0 + w'_1\cdot x + w'_2\cdot x^2 + \cdots + w'_{n-1}\cdot x^{n-1} $

$= f_d(x)$

\para{总结：}由于 $\sigma$ 既是双射又是同态，$\sigma$ 就是 $(n-1)$ 次多项式环 $\mathbb{Z}_t[X]/(X^n + 1)$ 与 $n$ 维向量空间 $\mathbb{Z}_t^n$ 之间的同构映射。


$ $


\subsubsection{寻找合适的模数 $t$}
\label{subsubsec:poly-vector-transformation-modulus}

要通过同构映射，将 $\mathbb{Z}_t[X]/(X^n + 1)$ 中的多项式求值为一个 $n$ 维向量，需要在 $X^n + 1 \bmod t$ 的 $n$ 个互异根处求值。然而，并非所有（次数，模数）组合 $ = (n, t)$ 都能使 $X^n + 1 \bmod t$ 具有 $n$ 个互异的根。例如，若 $n = 2$ 且 $t = 3$，则对于 $X = \{0, 1, 2\}$ 的任意取值，均有 $X^2+ 1 \not\equiv 0 \bmod 3$。因此，目标是在固定 $n$ 后，找到满足要求的 $t$，使 $X^n + 1 \pmod t$ 存在 $n$ 个互异的根，从而能够使用同构映射 $\sigma$。

先提出两个约束：(1) 选择 $t$ 为素数；(2) 保证 $t-1$ 是 $2n$ 的倍数。

由定理~\ref*{subsec:order-theorem}.4（\autoref{subsec:order-theorem}）中的费马小定理，已知 $a^{t - 1} \equiv 1 \bmod t$ 当且仅当 $a$ 与 $t$ 互素。这意味着，若 $t$ 为素数，则对于所有 $a \in \mathbb{Z}_t^{\times}$（即 $\mathbb{Z}_t$ 中除去 $\{0\}$ 的元素），都有 $a^{t - 1} \equiv 1 \bmod t$。假设 $g$ 是 $\mathbb{Z}_t^{\times}$ 的生成元，其各次幂能够生成 $\mathbb{Z}_t^{\times}$ 的全部元素。则有 $\textsf{Ord}_{\mathbb{Z}_t}(g) = t - 1$ 和 $g^{t - 1} \equiv 1 \bmod t$。由于对于某个 $k$ 有 $t - 1 = k \cdot 2n$，所以 $g^{k\cdot2n} \equiv (g^{k})^{2n} \equiv 1 \bmod t$，从而 $\textsf{Ord}_{\mathbb{Z}_t}(g^{k}) \leq 2n$。但是，由于 $\textsf{Ord}_{\mathbb{Z}_t}(g) = t - 1$，对于所有满足 $a < t - 1 = k \cdot 2n$ 的 $a$，都有 $g^a \not\equiv 1 \bmod t$。换言之，对于所有满足 $b < 2n$ 的 $b$，都有 $(g^k)^b \not\equiv 1 \bmod t$。因此，$\textsf{Ord}_{\mathbb{Z}_t}(g^{k}) = 2n$。

令 $c = g^k$。由于 $\textsf{Ord}_{\mathbb{Z}_t}(c) = 2n$，有 $c^{2n} \equiv 1 \bmod t$，换言之，$(c^n)^2 \equiv 1 \bmod t$。
此时，$c^n$ 只能为 1 或 -1。原因是，在关系 $X^2 \equiv 1 \bmod t$ 中，$X$ 只能取 $1$ 或 $-1 \equiv t - 1 \bmod t$。代入 $X = c^n$ 后，$c^n$ 只能取 $1$ 或 $-1 \equiv t - 1 \bmod t$。但由于 $\textsf{Ord}_{\mathbb{Z}_t}(c) = 2n$，$c^n$ 不可能等于 1，因为 $1^1 = 1$，而 1 小于 $c$ 的阶 $2n > 1$。因此，$c^n$ 只能等于 $-1 \equiv t - 1 \bmod t$。如果 $c^n = -1 \equiv t - 1 \bmod t$，那么 $c$ 就是 $X^n + 1 \bmod t$ 的根，因为 $X^n + 1 = c^n + 1 \equiv (t - 1) + 1 \equiv 0 \bmod t$。

因此，给定分圆多项式 $X^n + 1$，如果选择素数 $t$，使其对于某个整数 $k$ 满足 $t - 1 = k \cdot 2n$，则 $X^n + 1 \bmod t$ 的一个根为 $X = c = g^k$。

找到 $X^n + 1$ 的一个根之后，就可以推导出 $X^n + 1$ 的全部 $n$ 个互异的根。假设 $\omega$ 是 $X^n + 1 \bmod t$ 的一个根，则有：

$\omega^n + 1 \equiv 0 \bmod t$

$\omega^n \equiv t - 1 \bmod t$

$\omega^{2n} \equiv (t - 1)\cdot(t - 1) \equiv  t^2 - 2t + 1 \equiv   1 \bmod t$

虽然 $\omega^{2n} \equiv 1 \bmod t$，但 $\omega^{n} \not\equiv 1 \bmod t$；否则就有 $X^n + 1 = \omega^n + 1 = 1 + 1 = 2 \neq 0$，这与 $\omega$ 是 $X^n + 1$ 的根相矛盾。因此，$\textsf{Ord}_{\mathbb{Z}_t}(c) = 2n$。

接下来，按如下方式推导出 $X^n + 1 \bmod t$ 的其余 $n-1$ 个互异的根：

$(\omega^3)^n + 1 \equiv (\omega^{2n})\cdot \omega^n + 1 \equiv \omega^n + 1 \equiv t - 1 + 1 \equiv 0 \bmod t$

$(\omega^5)^n + 1 \equiv (\omega^{4n})\cdot \omega^n + 1 \equiv \omega^n + 1 \equiv t - 1 + 1 \equiv 0 \bmod t$

$\vdots$

$(\omega^{2n - 1})^n + 1 \equiv 0 \bmod t$ \\ \textcolor{red}{ $\rhd$ 对任意奇数 $k = 2j+1$，有 $(\omega^{2j+1})^n = (\omega^{2n})^j \cdot \omega^n = 1^j \cdot (-1) = -1$，因此 $(\omega^k)^n + 1 = 0$}

$ $

由于 $\textsf{Ord}_{\mathbb{Z}_t}(c) = 2n$，$\omega, \omega^3, \omega^5, \cdots, \omega^{2n - 1}$ 在 $\mathbb{Z}_t^{\times}$ 中互不相同。因此，$\{\omega^{2i + 1}\}_{i=0}^{n-1}$ 是 $X^n + 1$ 的 $n$ 个互异的根。同时，由于它们是分圆多项式 $X^n + 1$ 的根，所以也是 $n$ 个互异的 $(\mu=2n)$ 次本原单位根。


$ $


将上述结论归纳如下：

\begin{tcolorbox}[title={\textbf{\tboxtheorem{\ref*{subsubsec:poly-vector-transformation-modulus}} 整数（环）上多项式与向量之间的同构}}]

\begin{itemize}
\item 假设有一个 $(n-1)$ 次多项式环 $\mathbb{Z}_t[X] / F(X)$，其中 $F(X)$ 是具有 $n$ 个互异根 $\{x_0, \cdots, x_{n-1}\}$ 的 $n$ 次多项式；另有一个 $n$ 维向量空间 $\mathbb{Z}_t^n$（模 $t$ 的整数向量），其中的向量为 $\vec{v}$。

$\sigma: f(X) \in \mathbb{Z}_t[X] / F(X) \text{ } \longrightarrow \text{ } (f(x_0), f(x_1), f(x_2), \cdots, f(x_{n-1})) \in \mathbb{Z}_t^n$

$ $

则 $\sigma$ 对运算 $(+, \cdot)$ 保持同构：


$\sigma(f_a(X) + f_b(X)) = \sigma(f_a(X)) + \sigma(f_b(X))$

$\sigma(f_a(X) \cdot f_b(X)) = \sigma(f_a(X)) \odot \sigma(f_b(X))$

$ $

%\item If the $n$-degree polynomial $F(x)$ has a fewer number of roots than $n$, then isomorphism between the vectors and polynomials over the $(+, \cdot)$ operations holds without modulo $t$ as follows:

%$\sigma': f(X) \in \mathbb{Z}[X] / F(X) \text{ } \longrightarrow \text{ } (f(x_0), f(x_1), f(x_2), \cdots, f(x_{n-1})) \in \mathbb{Z}^n$

%$ $

\item 假设第 $(\mu=2n)$ 个分圆多项式为 $X^n + 1 \bmod t$，其中 $t$ 是素数，$t - 1$ 是 $2n$ 的某个倍数，且 $g$ 是 $\mathbb{Z}_t^{\times}$ 的生成元。则 $X^n + 1$ 的 $n$ 个互异的根（即 $(\mu=2n)$ 次本原单位根）可以高效地计算为 $\{\omega^{2i + 1}\}_{i=0}^{n-1}$，其中 $\omega = g^{\frac{t-1}{2n}} \bmod t$。


\end{itemize}
\end{tcolorbox}





\subsection{复数上多项式与向量之间的同构}
\label{subsec:poly-vector-transformation-complex}


在定理~\ref*{subsec:polynomial-ring-basis}（\autoref{subsec:polynomial-ring-basis}）中，已经学习了同构映射 $\sigma: f(X) \in \mathbb{Z}_t[X]/(X^n + 1) \longrightarrow (f(x_0),f(x_1),f(x_2), \cdots, f(x_{n-1})) \in \mathbb{Z}_t^n$，其中 $x_0, x_1, x_2, \cdots, x_{n-1} \in \mathbb{Z}$ 是分圆多项式 $X^n + 1$ 的根（即 $(\mu=2n)$ 次本原单位根），具体为 $\omega, \omega^3, \omega^5, \cdots, \omega^{2n-1}$。这里，$\omega$ 可以是 $X^n + 1$ 的任意一个根，因为每个 $\omega$ 都能生成全部根。本小节将说明复数上的向量空间与多项式环之间的同构，具体如下：

$\sigma_c: f(X) \in \mathbb{R}[X]/(X^n + 1) \longrightarrow (f(\omega),f(\omega^3),f(\omega^5), \cdots, f(\omega^{2n-1})) \in \mathbb{\hat{C}}^{n}$

$ $

其中，$X \in \mathbb{C}$，而 $\omega = e^{i\pi/n}$ 是复数上分圆多项式 $X^n + 1$ 的根，即 $(\mu=2n)$ 次本原单位根（见 \autoref{subsec:cyclotomic-theorem} 中的定理~\ref*{subsec:cyclotomic-theorem}.1）。将 $\mathbb{\hat{C}}^{n}$ 定义为一个 $n$ 维的\textit{特殊}向量空间，其中每个向量的后半部分元素，都是其前半部分元素按逆序排列后的共轭，例如 $(v_0, v_1, \cdots, v_{\frac{n}{2}-1}, \overline{v}_{\frac{n}{2}-1}, \cdots, \overline{v}_1, \overline{v}_0 )$。

\subsubsection{$\mathbb{C}^{\frac{n}{2}}$ 与 $\hat{\mathbb{C}}^{n}$ 之间的同构}
\label{subsec:poly-vector-transformation-complex-isomorphism1}

本节将 $\mathbb{C}^{\frac{n}{2}}$ 和 $\hat{\mathbb{C}}^{n}$ 都视为由复向量组成的向量空间，它们之间的映射定义为 $\phi: \mathbb{C}^{\frac{n}{2}} \rightarrow \hat{\mathbb{C}}^{n}$，如下所示：

$\phi(z_0, z_1, \dots, z_{\frac{n}{2}-1}) = (z_0, z_1, \dots, z_{\frac{n}{2}-1}, \overline{z_{\frac{n}{2}-1}}, \dots, \overline{z_1}, \overline{z_0})$


$ $

换言之，$\hat{\mathbb{C}}^{n}$ 是 $\mathbb{C}^{\frac{n}{2}}$ 的共轭扩展：$\hat{\mathbb{C}}^{n}$ 的前半部分元素与 $\mathbb{C}^{\frac{n}{2}}$ 相同，而 $\hat{\mathbb{C}}^{n}$ 的后半部分元素是 $\mathbb{C}^{\frac{n}{2}}$ 逆序排列后的共轭。

下面证明，$\mathbb{C}^{\frac{n}{2}}$ 与 $\hat{\mathbb{C}}^{n}$ 之间的映射 $\phi$ 是同构。


$ $

\para{双射：}映射 $\phi$ 的双射性是直接的。对于 $\mathbb{C}^{\frac{n}{2}}$ 中的每个向量，将其逆序共轭附加到末尾，即可唯一地映射到 $\hat{\mathbb{C}}^{n}$ 中的一个向量。反过来，对于 $\hat{\mathbb{C}}^{n}$ 中的任意向量，只需去掉后半部分元素，就能唯一地映射回 $\mathbb{C}^{\frac{n}{2}}$，从而建立一一对应关系。

\para{同态：}下面证明，$\phi$ 对加法 $(+)$ 和向量逐元素乘法 $(\odot)$ 保持同态。给定两个向量 $\vec{u}, \vec{v} \in \mathbb{C}^{\frac{n}{2}}$：
\begin{itemize}
    \item \textbf{加法：}$\phi(\vec{u} + \vec{v}) = \phi(\vec{u}) + \phi(\vec{v})$ 成立，因为复共轭对加法满足分配性，即 $\overline{u + v} = \overline{u} + \overline{v}$。
    \item \textbf{逐元素乘法：}$\phi(\vec{u} \odot \vec{v}) = \phi(\vec{u}) \odot \phi(\vec{v})$ 成立，因为乘积的共轭等于各因子共轭的乘积，即 $\overline{u \cdot v} = \overline{u} \cdot \overline{v}$。
\end{itemize}

$ $

因此，映射 $\phi: \mathbb{C}^{\frac{n}{2}} \rightarrow \hat{\mathbb{C}}^{n}$ 是同构。

\begin{tcolorbox}[title={\textbf{以实数或复数为标量的复向量空间}}]

根据数学定义，向量空间必须满足\textit{数乘封闭性}：对于空间中的任意向量 $\vec{v}$，以及标量域中的任意标量 $c$，乘积 $c \cdot \vec{v}$ 仍必须属于同一向量空间。在这一要求下，$\mathbb{C}^{\frac{n}{2}}$ 自然构成以复数为标量的复向量空间；由于复数包含实数，它也可以视为以实数为标量的向量空间。另一方面，$\hat{\mathbb{C}}^{n}$ 是一个\textit{只能以实数为标量}的复向量空间。

$ $

原因如下：向量 $\vec{v} = (v_0, v_1, \dots, v_{n-1}) \in \hat{\mathbb{C}}^n$ 必须满足 Hermitian 对称约束，即对于所有 $k$，都有 $v_{n-1-k} = \overline{v_k}$。考虑 $n = 4$ 的例子：

$$\vec{v} = (1+i,\ 2+2i,\ 2-2i,\ 1-i) \in \hat{\mathbb{C}}^4$$

将 $\vec{v}$ 乘以复标量 $c = i$：

$$i \cdot \vec{v} = (-1+i,\ -2+2i,\ 2+2i,\ 1+i)$$

要使结果仍属于 $\hat{\mathbb{C}}^4$，第三个分量必须是第二个分量的共轭。但 $\overline{-2+2i} = -2-2i$，并不等于 $2+2i$。因此，$i \cdot \vec{v} \notin \hat{\mathbb{C}}^4$，不满足封闭性。

$ $

更一般地，对于任意 $\vec{v} \in \hat{\mathbb{C}}^n$ 和任意标量 $c$，乘积 $c \cdot \vec{v}$ 保持 Hermitian 对称，当且仅当 $\overline{c} = c$，即 $c$ 为实数。因此，只有将标量限制为实数时，$\hat{\mathbb{C}}^n$ 才构成向量空间。这就是 $\hat{\mathbb{C}}^n$ 必须视为以实数为标量的复向量空间的原因。

$ $

本书始终将 $\mathbb{C}^{\frac{n}{2}}$ 与 $\hat{\mathbb{C}}^{n}$ 都视为以实数为标量、由复向量组成的向量空间。

$ $

\textbf{向量空间的同构：}向量空间之间的一个映射要成为同构，除了是双射并且对运算 $(+, \cdot)$ 保持同态以外，还必须对数乘保持同态。因此，还需证明 $\phi(c \cdot \vec{u}) = c \cdot \phi(\vec{u})$。在这里，这一点总是成立，因为标量 $c$ 为实数，取复共轭时不变，即 $\overline{c} = c$，因此 $\overline{c \cdot u_k} = c \cdot \overline{u_k}$。

$ $

\textbf{向量空间的维数：}向量空间的维数由完整描述其中向量所需的独立标量个数确定。因此，$\mathbb{C}^{\frac{n}{2}}$ 的维数为 $n$：它包含 $\dfrac{n}{2}$ 个元素，每个元素需要 2 个实标量描述，分别对应实部和虚部。$\hat{\mathbb{C}}^{n}$ 的维数也为 $n$，因为前 $\dfrac{n}{2}$ 个元素需要 $\dfrac{n}{2}\cdot 2$ 个标量描述，而后 $\dfrac{n}{2}$ 个元素只是前 $\dfrac{n}{2}$ 个元素按逆序排列后的共轭。

\end{tcolorbox}


$ $



\subsubsection{$\hat{\mathbb{C}}^{n}$ 与 $\mathbb{R}[X] / (X^n + 1)$ 之间的同构}
\label{subsec:poly-vector-transformation-complex-isomorphism2}

下面采用与 \autoref{subsec:poly-vector-transformation} 开头相同的推理，证明 $\sigma_c$ 是 $\hat{\mathbb{C}}^{n}$ 与 $\mathbb{R}[X] / (X^n + 1)$ 之间的同构，即它既是双射又是同态。

\begin{figure}[h!]
    \centering
  \includegraphics[width=0.3\linewidth]{figures/cyclotomic-polynomial.pdf}
  \caption{第 8 个分圆多项式 $x^4 + 1$ 的四个根的示意图。}
  \label{fig:cyclotomic-polynomial}
\end{figure}

$ $

\para{双射：}根据欧拉公式 $e^{i\theta} = \cos\theta + i\cdot\sin\theta$（见 \autoref{subsec:euler}），可以得到算术关系 $\omega = \overline{\omega^{2n-1}}, \text{ } \omega^3 = \overline{\omega^{2n-3}}, \cdots, \omega^{n-1} = \overline{\omega^{n+1}}$。换言之，一半的根分别是另一半根的共轭。\autoref{fig:cyclotomic-polynomial} 中的复平面图形直观地展示了这一点：红色箭头表示第 8 个分圆多项式 $X^4 + 1$ 的根，每个根都包含实部和虚部。如图所示，一半的红色箭头（即根）是另一半关于 $x$ 轴（实轴）的镜像。因此，可以将这些根表示为一个 $n$ 维向量，其元素为 $X^n + 1$ 的根，且后半部分元素是前半部分元素逆序排列后的共轭。基于这种向量结构，映射 $\sigma_c$ 可以重写为：

$\sigma_c(f(X)) = (f(\omega),f(\omega^3),f(\omega^5), \cdots, f(\omega^{n-1}), f(\overline{\omega^{n-1}}), f(\overline{\omega^{n-3}}), \cdots, f(\overline{\omega^3}), f(\overline{\omega}))$

$ $

由于 $f(\overline{X}) = \overline{f(X)}$（因为 $f(X)$ 的系数均为实数），可以将 $\sigma_c$ 重写为：

$\sigma_c(f(X)) = (f(\omega),f(\omega^3),f(\omega^5), \cdots, f(\omega^{n-3}), f(\omega^{n-1}), \overline{f(\omega^{n-1})}, \overline{f(\omega^{n-3})}, \cdots, \overline{f(\omega^3)}, \overline{f(\omega)})$

$ $

这种向量结构恰好符合 $\hat{\mathbb{C}}^n$ 的定义：$n$ 维向量 $\vec{\hat v}$ 的后半部分元素，是其前半部分元素逆序排列后的共轭。

为了证明双射性，还需说明：每个 $f(X) \in \mathbb{R}[X] / (X^n + 1)$ 都映射到某个 $\vec{\hat v} \in \hat{\mathbb{C}}^{n}$，并且不存在两个不同的 $f_1(X), f_2(X) \in \mathbb{R}[X] / (X^n + 1)$ 映射到同一个 $\vec{\hat v} \in \hat{\mathbb{C}}^{n}$。第一个要求成立，因为每个多项式 $f(X) \in \mathbb{R}[X] / (X^n + 1)$ 都可以在 $X^n + 1$ 的 $n$ 个互异根处求值，得到有效的数值。第二个要求也成立，因为在 $(n-1)$ 次多项式环中，每组 $n$ 个互异的 $(x, y)$ 坐标点只能对应一个次数不超过 $(n-1)$ 的多项式；这里将 $X$ 坐标固定为 $X^n + 1$ 的 $n$ 个互异根 $\{\omega, \omega^3, \cdots, \omega^{2n - 1}\}$。这一点已由拉格朗日多项式插值证明（见 \autoref{sec:polynomial-interpolation} 中的定理~\ref*{sec:polynomial-interpolation}）。

$ $

\para{同态：}$\sigma_c$ 是同态。根据 \autoref{subsec:poly-vector-transformation} 中的推理，无论 $X$ 是模整数还是复数，关系 $\sigma_c(f_a(X) + f_b(X)) = \sigma_c(f_a(X)) + \sigma_c(f_b(X))$ 和 $\sigma_c(f_a(X) \cdot f_b(X)) = \sigma_c(f_a(X)) \odot \sigma_c(f_b(X))$ 都成立。此外，对于任意实标量 $c \in \mathbb{R}$，有 $\sigma_c(c \cdot f(X)) = c \cdot \sigma_c(f(X))$，因为先对多项式作数乘再求值，与先求值再对结果作数乘是相同的。

$ $

由于 $\sigma_c$ 既是双射，又对运算 $(+, \cdot)$ 保持同态，因此它是同构。


\begin{tcolorbox}[title={\textbf{\tboxtheorem{\ref*{subsec:poly-vector-transformation-complex}} 复数上多项式与向量之间的同构}}]

下面的映射 $\sigma_c$ 是复数上多项式与向量之间的同构：

$\sigma_c: f(X) \in \mathbb{R}[X]/(X^n + 1) \longrightarrow (f(\omega),f(\omega^3),f(\omega^5), \cdots, f(\omega^{2n-1})) \in \hat{\mathbb{C}}^{n} \text{ } (\cong \mathbb{C}^{\frac{n}{2}})$

$ $

其中，$\omega = e^{i\pi/n}$ 是复数上分圆多项式 $X^n + 1$ 的根，即 $(\mu=2n)$ 次本原单位根；$\hat{\mathbb{C}}^{n}$ 是一个以实数为标量、由复向量组成的 $n$ 维向量空间，其中后半部分元素是前半部分元素逆序排列后的共轭。

\end{tcolorbox}





\subsection{多项式环与实标量复向量空间之间的基变换}
\label{subsec:poly-vector-basis-transfer}


假设多项式 $f_0(X), f_1(X), \cdots, f_{n-1}(X)$ 构成 $(n-1)$ 次多项式环的一组基，且 $\sigma$ 是从 $(n-1)$ 次多项式环到 $n$ 维向量空间的同构映射。那么，$(\sigma(f_0(X)), \sigma(f_1(X)), \cdots, \sigma(f_{n-1}(X)))$ 就构成该 $n$ 维向量空间的一组基。这是因为，经 $\sigma$ 映射得到的向量保持关于运算 $(+, \cdot)$ 的代数关系，而基向量与子空间之间的基关系，也就是线性无关性与张成性，可以用由运算 $(+, \cdot)$ 构成的线性代数表达式描述。因此，如果一组多项式构成基，那么它们经 $\sigma$ 映射得到的向量也保持相应的基关系。

同样的原理适用于多项式环与以实数为标量的复向量空间之间的对应。给定多项式环 $\mathbb{R}[X]/(x^n + 1)$，构造 $\mathbb{R}[X]/(x^n + 1)$ 的一组基最直接的方法如下：

$f_0(X) = 1$

$f_1(X) = X$

$f_2(X) = X^2$

$\vdots$

$f_{n-1}(X) = X^{n-1}$

$ $

这 $n$ 个多项式线性无关，因为每个多项式都只含有其独有次数的单项式，而这个单项式无法由其他单项式的线性组合表示。此外，这 $n$ 个多项式张成整个多项式环 $\mathbb{R}[X]/(x^n + 1)$，因为对每个多项式作数乘可以赋予对应次数项任意系数，将它们相加就可以表示多项式环 $\mathbb{R}[X]/(X^n + 1)$ 中的任意多项式。


现在，对上述构成 $(n-1)$ 次多项式环 $\mathbb{R}[X]/(x^n + 1)$ 的一组基的 $n$ 个多项式应用映射 $\sigma_c$。根据多项式基到向量基的转换原理（见 \autoref{subsec:polynomial-ring-basis} 中的定理~\ref*{subsec:polynomial-ring-basis}），可以利用这 $n$ 个多项式，即 $(n-1)$ 次多项式环的一组基，以及从多项式到向量的同构映射 $\sigma_c$，计算 $n$ 维特殊向量空间 $\hat{\mathbb C}^{n}$ 的一组基，如下所示：

$W = \begin{bmatrix}
\sigma_c(f_0(X)) \\
\sigma_c(f_1(X)) \\
\sigma_c(f_2(X)) \\
\vdots \\
\sigma_c(f_{n-1}(X)) \\
\end{bmatrix}
= \begin{bmatrix}
\sigma_c(1) \\
\sigma_c(X) \\
\sigma_c(X^2) \\
\vdots \\
\sigma_c(X^{n-1}) \\
\end{bmatrix}$
$= \begin{bmatrix}
1 & 1 & \cdots & 1 & 1\\
(\omega) & (\omega^3) & (\omega^5) & \cdots & (\omega^{2n-1})\\
(\omega)^2 & (\omega^3)^2 & (\omega^5)^2 & \cdots & (\omega^{2n-1})^2\\
\vdots & \vdots & \vdots & \ddots & \vdots \\
(\omega)^{n-1} & (\omega^3)^{n-1} & (\omega^5)^{n-1} & \cdots & (\omega^{2n-1})^{n-1}\\
\end{bmatrix}$

\text{ } \text{ }  $= \begin{bmatrix}
1 & 1 & 1 & \cdots & 1\\
(\omega) & (\omega^3) & \cdots & (\overline{\omega})^3 & (\overline{\omega})\\
(\omega)^2 & (\omega^3)^2 & \cdots & (\overline{\omega}^3)^2 & (\overline{\omega})^2\\
\vdots & \vdots & \vdots & \ddots & \vdots \\
(\omega)^{n-1} & (\omega^3)^{n-1} & \cdots & (\overline{\omega}^3)^{n-1} & (\overline{\omega})^{n-1}\\
\end{bmatrix}$

$ $

$W$ 是 $n$ 维特殊向量空间 $\hat{\mathbb C}^{n}$ 的一组有效的基。



\begin{tcolorbox}[title={\textbf{\tboxtheorem{\ref*{subsec:poly-vector-transformation-complex}} 多项式环与向量空间之间的基变换}}]

如果 $n$ 个多项式构成一个 $(n-1)$ 次多项式环的一组基，并通过从 $(n-1)$ 次多项式环到 $n$ 维向量空间的同构映射 $\sigma$（在复数情形下为 $\sigma_c$）转换为 $n$ 个互异向量，那么转换得到的 $n$ 个向量就构成该 $n$ 维向量空间的一组基。

\end{tcolorbox}

\clearpage

\section{欧拉公式}
\label{sec:euler}

\subsection{勾股恒等式}
\label{subsec:pythagorean}



\begin{figure}[!h]
    \centering
    \subfloat[勾股定理 \href{https://socratic.org/questions/58a20042b72cff17846a6feb}{（来源）}]{
  \includegraphics[width=0.3\linewidth]{figures/pythagorean-theorem.png}

  \label{fig:pythagorean-theorem}
}
    \subfloat[勾股恒等式 \href{https://socratic.org/questions/58a20042b72cff17846a6feb}{（来源）}]{
  \includegraphics[width=0.3\linewidth]{figures/pythagorean-identity.png}

  \label{fig:pythagorean-identity}
}
  \caption{}
\end{figure}

\begin{tcolorbox}[title={\textbf{\tboxtheorem{\ref*{subsec:pythagorean}} 勾股恒等式}}]
$\sin^2 \theta + \cos^2 \theta = 1$
\end{tcolorbox}

\begin{proof}
$ $
\begin{enumerate}

\item 在 \autoref{fig:pythagorean-theorem} 中，根据勾股定理，有 $a^2 + b^2 = c^2$。

\item 在 \autoref{fig:pythagorean-identity} 中，根据正弦与余弦的定义，

$\sin\theta = \dfrac{a}{c}, \text{ } \cos\theta = \dfrac{b}{c}$

\item 结合步骤 1 和步骤 2：

$\sin^2 \theta + \cos^2 \theta = \dfrac{a^2}{c^2} + \dfrac{b^2}{c^2} = \dfrac{a^2 + b^2}{c^2} = 1$


\end{enumerate}
\end{proof}

\subsection{虚数}
\label{subsec:imaginary}


\begin{tcolorbox}[title={\textbf{\tboxdef{\ref*{subsec:imaginary}} 虚数}}]

\begin{itemize}

\item 将 \boldmath{$i$} 定义为满足 $i^2 =-1$ 的虚数单位。
\item \textbf{复数}是形如 $a + b\,i$ 的数，其中 $a$ 和 $b$ 均为实数（例如，$a = 5.6 + 4.3i)$

\item \textbf{实数}不包含虚部（例如，$a = 13.4$）。

\item 若 $a$ 和 $\overline{a}$ 具有相同的实部以及符号相反的虚部，则 $\overline{a}$ 是 $a$ 的\textbf{共轭}。

(e.g., $a = 3 + i\cdot3.4$, \text{ } $\overline{a} = 3 - i\cdot3.4$)

\item \textbf{Hermitian 向量}指后半部分元素为前半部分元素逆序排列后复共轭的向量，如下面的 $n$ 维向量所示：

$\vec{v} = (v_1, v_2, v_3, \cdots, v_{\frac{n}{2}-1}, v_{\frac{n}{2}}, \overline{v}_{\frac{n}{2}}, \overline{v}_{\frac{n}{2} - 1}, \cdots, \overline{v}_3, \overline{v}_2, \overline{v}_1)$


\end{itemize}
\end{tcolorbox}



\subsection{欧拉公式}
\label{subsec:euler}




\begin{tcolorbox}[title={\textbf{\tboxdef{\ref*{subsec:euler}} 欧拉公式}}]

$e^{i\theta} = \cos\theta + i\cdot\sin\theta$

\end{tcolorbox}

\begin{figure}[!h]
    \centering
  \includegraphics[width=0.4\linewidth]{figures/euler-formula.png}
  \caption{复平面单位圆上的欧拉公式示意图。\href{https://en.wikipedia.org/wiki/Euler's_formula}{（来源）}}
\end{figure}

在 \autoref{fig:complex-plane} 中，$e^{i\theta}$ 的值表示为复平面圆上的一个坐标点，其中 $x$ 轴表示实部，$y$ 轴表示虚部。随着 $\theta$ 增大，虚部（$\sin\theta$）在 $1$ 与 $-1$ 之间振荡，在虚轴上分别达到 $i$ 和 $-i$；实部（$\cos\theta$）在 $1$ 与 $-1$ 之间振荡，周期为 $2\pi$。


\subsection{以复数上分圆多项式的根构造 Vandermonde 矩阵}
\label{subsec:vandermonde-euler}

 
 
本小节利用复数上第 $\mu$ 个分圆多项式的 $n$ 个互异根，构造 Vandermonde 矩阵（见 \autoref{subsec:vandermonde}），其中 $\mu$ 是 2 的幂，具体如下：



\begin{tcolorbox}[title={\textbf{\tboxtheorem{\ref*{subsec:vandermonde-euler}} 以复数上指标为 2 的幂的分圆多项式的根构造 Vandermonde 矩阵}}]

假设有一个 $n \times n$ 的 Vandermonde 矩阵（其中 $n$ 是 2 的幂），由第 $\mu$ 个分圆多项式的 $n$ 个互异根组成（见 \autoref{subsec:cyclotomic-theorem} 中的定理~\ref*{subsec:cyclotomic-theorem}.1），其中 $\mu$ 是 2 的幂，且 $n = \dfrac{\mu}{2}$。换言之，$V = \mathit{Vander}(x_0, x_1, \cdots, x_{n-1})$，并且对于 $1 \leq j \leq n$，每个根均为 $x_j = (e^{i\pi/n})^{2j-1}$，即 $\mu$ 次本原单位根。则有：

$V \cdot V^T = \begin{bmatrix}
0 & \cdots & 0 & 0 & n\\
0 & \cdots & 0 & n & 0\\
0 & \cdots & n & 0 & 0\\
\vdots & \iddots & \vdots & \vdots & \vdots \\
n & 0 & 0 & \cdots & 0\\
\end{bmatrix} = n \cdot I^R_n$

$ $

并且 $V^{-1} = \dfrac{V^T \cdot I_n^R}{n}$


\end{tcolorbox}

\begin{proof}

$ $

\begin{enumerate}

\item 给定 $\omega = e^{i\pi/n}$，其中每个 $x_j = (\omega)^{2j-1}$。因此，可以展开如下：

$V \cdot V^T = \begin{bmatrix}
1 & (\omega) & (\omega)^2 & \cdots & (\omega)^{n-1}\\
1 & (\omega^3) & (\omega^3)^2 & \cdots & (\omega^3)^{n-1}\\
1 & (\omega^5) & (\omega^5)^2 & \cdots & (\omega^5)^{n-1}\\
\vdots & \vdots & \vdots & \ddots & \vdots \\
1 & (\omega^{2n-1}) & (\omega^{2n-1})^2 & \cdots & (\omega^{2n-1})^{n-1}\\
\end{bmatrix} 
\cdot 
\begin{bmatrix}
1 & 1 & 1 & \cdots & 1\\
(\omega) & (\omega^3) & (\omega^5) & \cdots & (\omega^{2n-1})\\
(\omega)^2 & (\omega^3)^2 & (\omega^5)^2 & \cdots & (\omega^{2n-1})^2\\
\vdots & \vdots & \vdots & \ddots & \vdots \\
(\omega)^{n-1} & (\omega^3)^{n-1} & (\omega^5)^{n-1} & \cdots & (\omega^{2n-1})^{n-1}\\
\end{bmatrix} $

$ $


$=
\begin{bmatrix}
\sum\limits_{k=0}^{n-1} \omega^{2k} & \sum\limits_{k=0}^{n-1} \omega^{4k}  & \sum\limits_{k=0}^{n-1} \omega^{6k} & \cdots & \sum\limits_{k=0}^{n-1} \omega^{2nk} \\

\sum\limits_{k=0}^{n-1} \omega^{4k} & \sum\limits_{k=0}^{n-1} \omega^{6k}  & \sum\limits_{k=0}^{n-1} \omega^{8k} & \cdots & \sum\limits_{k=0}^{n-1} \omega^{2k(n+1)} \\

\sum\limits_{k=0}^{n-1} \omega^{6k} & \sum\limits_{k=0}^{n-1} \omega^{8k} & \sum\limits_{k=0}^{n-1} \omega^{10k} & \cdots & \sum\limits_{k=0}^{n-1} \omega^{2k(n+2)} \\

\vdots & \vdots & \vdots & \ddots & \vdots \\
\sum\limits_{k=0}^{n-1} \omega^{2nk} & \sum\limits_{k=0}^{n-1} \omega^{2(n+1)k} & \sum\limits_{k=0}^{n-1} \omega^{2(n+2)k} & \cdots & \sum\limits_{k=0}^{n-1} \omega^{2(n+n-1)k} \\

\end{bmatrix}$

$ $

\item 矩阵 $V \cdot V^T$ 的副对角线元素为 $\sum\limits_{k=0}^{n-1} \omega^{2nk}$。可以推导：

$\sum\limits_{k=0}^{n-1} \omega^{2nk} = \sum\limits_{k=0}^{n-1} (e^{i\pi/n})^{2nk} = \sum\limits_{k=0}^{n-1}e^{2\pi k i} = \sum\limits_{k=0}^{n-1}(\cos(2\pi k) + i\sin(2\pi k)) = \sum\limits_{k=0}^{n-1} (1 + 0) = n$

这说明，矩阵 $V \cdot V^T$ 的副对角线元素为 $n$。

$ $

\item 接下来，证明矩阵 $V \cdot V^T$ 除副对角线外的所有元素均为 0。也就是说，证明：

$\sum\limits_{k=0}^{n-1} \omega^{2k} = \sum\limits_{k=0}^{n-1} \omega^{4k} = \sum\limits_{k=0}^{n-1} \omega^{6k} = \cdots = \sum\limits_{k=0}^{n-1} \omega^{2(n-1)k} = \sum\limits_{k=0}^{n-1} \omega^{2(n+1)k} = \cdots = \sum\limits_{k=0}^{n-1} \omega^{2(2n-1)k} = 0$

证明中需要使用等比数列求和公式 $\sum\limits_{i=0}^{n-1}x^i = \dfrac{x^n - 1}{x - 1}$：

\begin{tcolorbox}[title={\textbf{\tboxtheorem{\ref*{subsec:vandermonde-euler}.1} 等比数列求和公式}}]

设等比数列之和为 $S_n = 1 + x + x^2 + \cdots + x^{n - 1}$

则 $x \cdot S_n = x + x^2 + x^3 + \cdots + x^{n}$

$x \cdot S_n - S_n = (x + x^2 + x^3 + \cdots + x^{n}) - (1 + x + x^2 + \cdots + x^{n - 1}) = x^n - 1$

$S_n\cdot (x - 1) = x^n - 1$

$S_n = \dfrac{x^n - 1}{x - 1}$ \textcolor{red}{ $\rhd$ 其中要求 $x \neq 1$}

\end{tcolorbox}

利用等比数列求和公式 $\sum\limits_{i=0}^{n-1}x^i = \dfrac{x^n - 1}{x - 1}$，

对于 $1 \leq m \leq (2n - 1)$，有 $\sum\limits_{k=0}^{n-1} \omega^{2mk} = \dfrac{(\omega^{2m})^n - 1}{\omega^{2m} - 1} = \dfrac{(\omega^{2n})^{m} - 1}{\omega^{2m} - 1} = \dfrac{1 - 1}{\omega^{2m} - 1} = 0$ \textcolor{red}{ $\rhd$ 因为 $\textsf{Ord}(\omega) = 2n$}

$ $

因此，

$\sum\limits_{k=0}^{n-1} \omega^{2k} = \sum\limits_{k=0}^{n-1} \omega^{4k} = \sum\limits_{k=0}^{n-1} \omega^{6k} = \cdots = \sum\limits_{k=0}^{n-1} \omega^{2(n-1)k} = \sum\limits_{k=0}^{n-1} \omega^{2(n+1)k} = \cdots = \sum\limits_{k=0}^{n-1} \omega^{2(2n-1)k} = 0$

$ $

\item 根据步骤 2 和步骤 3 的证明，$V \cdot V^T = 
\begin{bmatrix}
0 & \cdots & 0 & 0 & n\\
0 & \cdots & 0 & n & 0\\
0 & \cdots & n & 0 & 0\\
\vdots & \iddots & \vdots & \vdots & \vdots \\
n & 0 & 0 & \cdots & 0\\
\end{bmatrix} = n \cdot I_n^R$

\item 给定 $V \cdot V^T = n \cdot I_n^R$，

$V^{-1} \cdot V \cdot V^T = V^{-1} \cdot n \cdot I_n^R$

$V^T = V^{-1} \cdot n \cdot I_n^R$

$V^T \cdot I_n^R = V^{-1} \cdot n \cdot I_n^R  \cdot I_n^R$

$V^T \cdot I_n^R = V^{-1} \cdot n$ \textcolor{red}{\text{ } \# 因为 $I_n^R  \cdot I_n^R = I_n$}

$V^{-1} = \dfrac{V^T \cdot I_n^R}{n}$

\end{enumerate}
\end{proof}

后续在 CKKS 方案（\autoref{sec:ckks}）中，将使用 $V^{-1}$ 将复向量编码为实向量，并使用 $V^T$ 将实向量解码为复向量（见 \autoref{subsec:ckks-enc-dec}）。

$ $

%\para{Condition for $\bm \mu$:} It's worthwhile to note that the property $V\cdot V^T = n\cdot I_n^R$ does not hold if $\mu$ (denoting the $\mu$-th cyclotomic polynomial) is not a power of 2. In particular, step 3 of the proof does not hold anymore if $\mu$ is not a power of 2:

%$\sum\limits_{k=0}^{n-1} \omega^{2k} \neq \sum\limits_{k=0}^{n-1} \omega^{4k} \neq \sum\limits_{k=0}^{n-1} \omega^{6k} \neq \cdots \neq \sum\limits_{k=0}^{n-1} \omega^{2(n-1)k} \neq \sum\limits_{k=0}^{n-1} \omega^{2(n+1)k} \neq \cdots \neq \sum\limits_{k=0}^{n-1} \omega^{2(2n-1)k} \neq 0$


\para{关于 $\bm \mu$ 的条件：}恒等式 $V \cdot V^T = n \cdot I_n^R$ 只依赖于
$\textsf{Ord}(\omega) = 2n$（步骤 3），因此对任意 $\mu = 2n$（即任意偶数）都成立。
同样，节点 $\omega, \omega^3, \ldots, \omega^{2n-1}$ 是 $X^n + 1$ 的 $n$ 个根，
这对任意 $n$ 都成立，因为 $(\omega^{2j-1})^n = \omega^{2nj}\cdot\omega^{-n} = 1\cdot(-1) = -1$。

不过，将 $\mu$ 取为 2 的幂的原因并非来自本定理，
而是来自全同态加密方案本身。在实际应用中，CKKS、BFV 和 BGV 通常采用
分圆环 $\mathbb{Z}[X]/(X^n + 1)$，其中 $n$ 是 2 的幂。这种选择使 $X^n + 1 = \Phi_{2n}(X)$ 在有理数域上不可约，而这一性质成立
当且仅当 $n$ 是 2 的幂；同时，它还支持高效的负循环 NTT 乘法，
并具有清晰的 Ring-LWE 安全性归约。一旦将
$n$（从而也将 $\mu = 2n$）固定为 2 的幂，这些节点就自动成为本原的
$\mu$ 次单位根：根据定理~\ref*{subsec:roots-theorem}.4（\autoref{subsec:roots-theorem}），
$\omega^k$ 是 $\mu$ 次本原单位根，当且仅当 $\gcd(\mu, k) = 1$；而每个奇数指数
$1, 3, \ldots, 2n-1$ 都与 $\mu = 2n$ 互素，当且仅当 $2n$ 不含奇素因子，即 $2n$
是 2 的幂。因此，$\{\omega, \omega^3, \ldots, \omega^{2n-1}\}$ 恰好为本原的
$(\mu=2n)$ 次单位根，是选择二幂分圆环的结果。同时，只要 $\mu$ 为偶数，$V \cdot V^T = n \cdot I_n^R$ 仍然成立。



\subsection{以环（$\mathbb{Z}_p$）上分圆多项式的根构造 Vandermonde 矩阵}
\label{subsec:vandermonde-euler-integer-ring}

定理~\ref*{subsec:vandermonde-euler}（见 \autoref{subsec:vandermonde-euler}）证明了 $V \cdot V^T = n \cdot I^R_n$，其中 $V$ 为 Vandermonde 矩阵 $V = \mathit{Vander}(x_0, x_1, \cdots, x_{n-1})$，每个 $x_i$ 是 $X \in \mathbb{C}$（即复数）上的 $\mu$ 次本原单位根，且 $\mu$ 为 2 的幂。本小节将证明，即使每个 $x_i$ 是 $X \in \mathbb{Z}_p$（即环）上的 $\mu$ 次本原单位根，关系 $V \cdot V^T = n \cdot I^R_n$ 也成立。具体地，将证明定理~\ref*{subsec:vandermonde-euler}：


\begin{tcolorbox}[title={\textbf{\tboxtheorem{\ref*{subsec:vandermonde-euler-integer-ring}} 以环（$\mathbb{Z}_p$）上指标为 2 的幂的分圆多项式的根构造 Vandermonde 矩阵}}]

证明方法与定理~\ref*{subsec:vandermonde-euler}（见 \autoref{subsec:vandermonde-euler}）相同。假设有一个 $n \times n$ 的 Vandermonde 矩阵（其中 $n$ 是 2 的幂），由环 $X \in \mathbb{Z}_p$ 上第 $\mu$ 个分圆多项式的 $n$ 个互异根组成，其中 $\mu$ 是 2 的幂且 $n = \dfrac{\mu}{2}$。换言之，$V = \mathit{Vander}(x_0, x_1, \cdots, x_{n-1})$，其中每个 $x_i$ 都是 $X^n + 1$ 的根，即 $(\mu=2n)$ 次本原单位根。则有：

$V \cdot V^T = \begin{bmatrix}
0 & \cdots & 0 & 0 & n\\
0 & \cdots & 0 & n & 0\\
0 & \cdots & n & 0 & 0\\
\vdots & \iddots & \vdots & \vdots & \vdots \\
n & 0 & 0 & \cdots & 0\\
\end{bmatrix} = n \cdot I^R_n$

$ $

并且 $V^{-1} = n^{-1}\cdot V^T \cdot I_n^R$


\end{tcolorbox}
\begin{proof}
$ $
\begin{enumerate}
\item 将 $V \cdot V^T$ 展开如下：

$V \cdot V^T = \begin{bmatrix}
1 & (\omega) & (\omega)^2 & \cdots & (\omega)^{n-1}\\
1 & (\omega^3) & (\omega^3)^2 & \cdots & (\omega^3)^{n-1}\\
1 & (\omega^5) & (\omega^5)^2 & \cdots & (\omega^5)^{n-1}\\
\vdots & \vdots & \vdots & \ddots & \vdots \\
1 & (\omega^{2n-1}) & (\omega^{2n-1})^2 & \cdots & (\omega^{2n-1})^{n-1}\\
\end{bmatrix} 
\cdot 
\begin{bmatrix}
1 & 1 & 1 & \cdots & 1\\
(\omega) & (\omega^3) & (\omega^5) & \cdots & (\omega^{2n-1})\\
(\omega)^2 & (\omega^3)^2 & (\omega^5)^2 & \cdots & (\omega^{2n-1})^2\\
\vdots & \vdots & \vdots & \ddots & \vdots \\
(\omega)^{n-1} & (\omega^3)^{n-1} & (\omega^5)^{n-1} & \cdots & (\omega^{2n-1})^{n-1}\\
\end{bmatrix} $

$ $


$=
\begin{bmatrix}
\sum\limits_{k=0}^{n-1} \omega^{2k} & \sum\limits_{k=0}^{n-1} \omega^{4k}  & \sum\limits_{k=0}^{n-1} \omega^{6k} & \cdots & \sum\limits_{k=0}^{n-1} \omega^{2nk} \\

\sum\limits_{k=0}^{n-1} \omega^{4k} & \sum\limits_{k=0}^{n-1} \omega^{6k}  & \sum\limits_{k=0}^{n-1} \omega^{8k} & \cdots & \sum\limits_{k=0}^{n-1} \omega^{2k(n+1)} \\

\sum\limits_{k=0}^{n-1} \omega^{6k} & \sum\limits_{k=0}^{n-1} \omega^{8k} & \sum\limits_{k=0}^{n-1} \omega^{10k} & \cdots & \sum\limits_{k=0}^{n-1} \omega^{2k(n+2)} \\

\vdots & \vdots & \vdots & \ddots & \vdots \\
\sum\limits_{k=0}^{n-1} \omega^{2nk} & \sum\limits_{k=0}^{n-1} \omega^{2(n+1)k} & \sum\limits_{k=0}^{n-1} \omega^{2(n+2)k} & \cdots & \sum\limits_{k=0}^{n-1} \omega^{2(n+n-1)k} \\

\end{bmatrix}$

$ $

其中，$\omega$（即 $(\mu=2n)$ 次本原单位根）的阶为 $2n$。

$ $

\item 矩阵 $V \cdot V^T$ 的副对角线元素为 $\sum\limits_{k=0}^{n-1} \omega^{2nk}$。由于 $\textsf{Ord}_p(\omega) = 2n$，可知 $\omega^{2n} \equiv 1 \bmod p$。因此，矩阵 $V \cdot V^T$ 的每个副对角线元素都为 $\sum\limits_{k=0}^{n-1} 1 = n$。

$ $

\item 接下来，证明矩阵 $V \cdot V^T$ 除副对角线外的其他元素均为 $0$。也就是说，证明：

$\sum\limits_{k=0}^{n-1} \omega^{2k} = \sum\limits_{k=0}^{n-1} \omega^{4k} = \sum\limits_{k=0}^{n-1} \omega^{6k} = \cdots = \sum\limits_{k=0}^{n-1} \omega^{2(n-1)k} = \sum\limits_{k=0}^{n-1} \omega^{2(n+1)k} = \cdots = \sum\limits_{k=0}^{n-1} \omega^{2(2n-1)k} = 0$

$ $

由等比数列求和公式（定理~\ref*{subsec:vandermonde-euler}.1），上式成立。如定理~\ref*{subsec:vandermonde-euler} 的证明步骤 3 所示，每个元素都可写成 $\sum_{k=0}^{n-1} (\omega^{2m})^k$，其中 $m$ 为某个整数，且 $2m$ 不是 $2n$ 的倍数。令 $r = \omega^{2m}$，则有：

$\sum\limits_{k=0}^{n-1} r^k = \dfrac{r^n - 1}{r - 1} = \dfrac{(\omega^{2m})^n - 1}{\omega^{2m} - 1} = \dfrac{(\omega^n)^{2m} - 1}{\omega^{2m} - 1}$

$ $

由于 $\omega$ 是 $X^n + 1$ 的根，所以 $\omega^n \equiv -1 \bmod p$。因此：

$\dfrac{(-1)^{2m} - 1}{\omega^{2m} - 1} = \dfrac{1 - 1}{\omega^{2m} - 1} = 0$

$ $

由此，所有非副对角线位置上的和都为 0。

$ $

\item 根据步骤 1 和步骤 2，矩阵 $V \cdot V^T$ 的副对角线元素为 $n$，其他位置的元素都为 $0$。

\item 现在推导 $V^{-1}$ 的公式。给定 $V \cdot V^T = n \cdot I_n^R$，

$V^{-1} \cdot V \cdot V^T = V^{-1} \cdot n \cdot I_n^R$

$V^T = V^{-1} \cdot n \cdot I_n^R$

$V^T \cdot I_n^R = V^{-1} \cdot n \cdot I_n^R  \cdot I_n^R$

$V^T \cdot I_n^R = V^{-1} \cdot n$ \textcolor{red}{\text{ } \# 因为 $I_n^R  \cdot I_n^R = I_n$}

$ $

这里有一点需要注意：模运算不支持直接进行通常的数值除法（如 \autoref{subsec:modulo-division} 所述）。因此，定理~\ref*{subsec:vandermonde-euler}（见 \autoref{subsec:vandermonde-euler}）中的公式 $V^{-1} = \dfrac{V^T \cdot I_n^R}{n}$ 不能直接用于当前情形，因为模 $p$ 的运算不允许直接用 $n$ 除 $V^T \cdot I_n^R$。所以，改为将 $V^T \cdot I_n^R$ 乘以 $n$ 的逆元，即 $n^{-1}$。继续推导如下：

$V^T \cdot I_n^R = V^{-1} \cdot n$

$V^T \cdot I_n^R \cdot n^{-1}= V^{-1} \cdot n \cdot n^{-1}$

$V^{-1} = n^{-1}\cdot V^T \cdot I_n^R$

\end{enumerate}
\end{proof}

最终证明了 $V\cdot V^T = n \cdot I_n^R$ 和 $V^{-1} = n^{-1}\cdot V^T \cdot I_n^R$。后续在 BFV 方案（\autoref{sec:bfv}）中，将使用 $V^{-1}$ 将整数向量编码为多项式系数向量，再使用 $V^T$ 将其解码回整数向量（见 \autoref{subsec:bfv-batch-encoding}）。

$ $

\para{关于 $\bm \mu$ 的条件：}与 CKKS 中的情况类似，如果 $\mu$（表示第 $\mu$ 个分圆多项式的指标）不是 2 的幂，性质 $V\cdot V^T = n\cdot I_n^R$ 就不成立。具体来说，当 $\mu$ 不是 2 的幂时，证明的步骤 3 不再成立：

$\sum\limits_{k=0}^{n-1} \omega^{2k} \neq \sum\limits_{k=0}^{n-1} \omega^{4k} \neq \sum\limits_{k=0}^{n-1} \omega^{6k} \neq \cdots \neq \sum\limits_{k=0}^{n-1} \omega^{2(n-1)k} \neq \sum\limits_{k=0}^{n-1} \omega^{2(n+1)k} \neq \cdots \neq \sum\limits_{k=0}^{n-1} \omega^{2(2n-1)k} \neq 0$

\clearpage

\section{模数重缩放}
\label{sec:modulus-rescaling}

\subsection{同余关系中的模数重缩放}
\label{subsec:modulo-rescaling}

设 $q,q'$ 为正整数模数，$r=q'/q$，并固定具有整数平移不变性的最近整数舍入规则 $R(x)=\lfloor x+1/2\rfloor$。对于整数 $h$，有 $R(x+h)=R(x)+h$。记 $\delta_x=R(rx)-rx$，则 $|\delta_x|\le1/2$。

\begin{table}[h!]
\centering
\begin{tabular}{|l|l|}
\hline
原关系 & 模 $q'$ 下的缩放关系 \\
\hline
$a\equiv b\pmod q$ & $R(ra)\equiv R(rb)\pmod{q'}$ \\
$a+c\equiv b+d\pmod q$ & $R(ra)+R(rc)\equiv R(rb)+R(rd)+\eta\pmod{q'}$ \\
$ac\equiv bd\pmod q$ & $R(rac)\equiv R(rbd)\pmod{q'}$ \\
\hline
\end{tabular}
\caption{同余关系的缩放，其中 $\eta=\delta_a+\delta_c-\delta_b-\delta_d$ 且 $|\eta|\le2$。}
\label{tab:rescaling}
\end{table}

\begin{proof}
若 $a=b+kq$，则 $ra=rb+kq'$，并且 $R(ra)=R(rb)+kq'$。因此，对同余式两侧的整体分别舍入，会精确保持同余，无论 $ra$ 是否为整数。相同论证也适用于整个乘积 $ac$ 和 $bd$。

对于分别舍入的加数，由 $a+c=b+d+kq$ 得到
\[
R(ra)+R(rc)-R(rb)-R(rd)=kq'+\delta_a+\delta_c-\delta_b-\delta_d.
\]
最后四项组成整数 $\eta$，满足 $|\eta|\le2$。当它们抵消时，特别是所有缩放后的加数均为整数时，关系精确成立；否则必须保留这个显式误差。
\end{proof}

这些结论是单向蕴含，一般不是等价关系。例如，取 $q=8$、$q'=2$，则 $a=0$ 与 $b=1$ 均舍入为零，但 $a\not\equiv b\pmod8$。舍入会丢失信息。此外，分别缩放乘积的两个因子会将整体尺度变为 $r^2$，不能替代将整个乘积只乘一次 $r$。

对于固定整数乘数 $s$，有用的恒等式为 $R(ra)s=ras+\delta_a s$，误差至多为 $|s|/2$。因此，在 LWE 相位 $b-\sum_i a_i s_i$ 中，分别舍入 $b$ 和各个 $a_i$ 会引入 $\delta_b-\sum_i\delta_{a_i}s_i$。后续模数切换正确性需要约束这个误差，而不能只使用没有误差定义的近似同余符号。

\subsubsection{示例}
\label{subsec:modulo-rescaling-ex}

假设有如下同余关系：

$b \equiv a\cdot s + \Delta \cdot m + e \bmod q$，\text{ }其中：$q = 30$，\text{ }$s = 5$，\text{ }$a = 10$，\text{ }$\Delta = 10$，\text{ }$m = 1$，\text{ }$e = 10$，\text{ }$b = 40$

$ $

首先，代入给定的示例值，验证上述同余关系是否成立：

$b \equiv a\cdot s + \Delta \cdot m + e \bmod 30$

$40 \equiv 10 \cdot 5 + 10 \cdot 1 + 10 \bmod 30$

$40 \equiv 70 \bmod 30$ 

$ $

该同余关系成立。

$ $

现在，希望将模数按 $30 \rightarrow 3$ 重缩放。根据 \autoref{tab:rescaling} 中介绍的重缩放原理，计算缩放后的值如下：

$q'= 3$, \text{ } $s = 5$, \text{ } $m = 1$

$\hat{a} = \Bigg\lceil a\cdot\dfrac{3}{30} \Bigg\rfloor = \Bigg\lceil 10\cdot\dfrac{3}{30} \Bigg\rfloor = 1$

$\hat{\Delta} = \Bigg\lceil \Delta\cdot\dfrac{3}{30} \Bigg\rfloor = \Bigg\lceil 10\cdot\dfrac{3}{30} \Bigg\rfloor = 1$



$\hat{e} = \Bigg\lceil e\cdot\dfrac{3}{30} \Bigg\rfloor = \Bigg\lceil 10\cdot\dfrac{3}{30} \Bigg\rfloor = 1$

$\hat{b} = \Bigg\lceil b\cdot\dfrac{3}{30} \Bigg\rfloor = \Bigg\lceil 40\cdot\dfrac{3}{30} \Bigg\rfloor = 4$

$ $

将模数按 $30 \rightarrow 3$ 重缩放后，得到的同余关系推导如下：

$\Bigg\lceil b\dfrac{3}{30} \Bigg\rfloor \equiv \Bigg\lceil s\cdot a \dfrac{3}{30} \Bigg\rfloor + \Bigg\lceil m \cdot \Delta \dfrac{3}{30} \Bigg\rfloor + \Bigg\lceil e \dfrac{3}{30} \Bigg\rfloor \bmod 3$

$\hat{b} \equiv \hat{a} \cdot s + \hat{\Delta} \cdot m + \hat{e}  \bmod 3$
 \text{ }（这是精确同余关系，因为所有缩放后的值均无小数部分。）

$4 \equiv 1 \cdot 5 + 1 \cdot 1 + 1 \bmod 3$

$4 \equiv 7 \bmod 3$

$ $

该例精确成立，因为 $ra$、$r\Delta$ 和 $re$ 都是整数，所有舍入误差均为零。若部分缩放后的项不是整数，误差抵消后仍可能精确成立；一般情况下，应保留 \autoref{tab:rescaling} 中相应的误差项。











\clearpage

\section{中国剩余定理}
\label{sec:chinese-remainder}
\textbf{- 参考资料 1：}
\href{https://brilliant.org/wiki/chinese-remainder-theorem/}{Brilliant -- 中国剩余定理}~\cite{crt}

\noindent \textbf{- 参考资料 2：}
\href{https://www.youtube.com/watch?v=fz1vxq5ts5I}{YouTube -- 扩展欧几里得算法教程}



\begin{tcolorbox}[title={\textbf{\tboxtheorem{\ref*{sec:chinese-remainder}.1} 中国剩余定理}}]

假设有两两互素且大于一的整数 $n_0, n_1, n_2, \cdots, n_k$，令 $N = n_0 n_1 \cdots n_k$。分别在各模数 $n_0, n_1, n_2, \dots, n_k$ 下选取共 $k + 1$ 个随机整数 $a_0, a_1, a_2, \cdots, a_k$，即 $a_0 \in \mathbb{Z}_{n_0}$、$a_1 \in \mathbb{Z}_{n_1}$、$\cdots$、$a_k \in \mathbb{Z}_{n_k}$。则存在唯一的解 $x \bmod N$，使得对于每个 $0 \le i \le k$，都有 $x \equiv a_i \pmod{n_i}$。也就是说：

\text{ } $x \equiv a_0 \bmod n_0$

\text{ } $x \equiv a_1 \bmod n_1$

\text{ } $x \equiv a_2 \bmod n_2$

\text{ } \text{ } $\vdots$

\text{ } $x \equiv a_k \bmod n_k$

$ $

为了计算 $x$，首先对每个 $0 \leq i \leq k$，按如下方式计算 $y_i$ 和 $z_i$：

$y_i = \dfrac{N}{n_i}, \text{ } z_i = y_i^{-1} \bmod n_i$

$ $

每个 $y_i$ 的逆元，即 $y_i^{-1}$，都可以通过扩展欧几里得算法计算（参见 \href{https://www.youtube.com/watch?v=fz1vxq5ts5I}{YouTube 教程}）。于是，唯一解 $x$ 可以计算为：

$x = \sum\limits_{i=0}^k a_i y_i z_i $ \textcolor{red}{ $\rhd$ 也可以计算 $ x = \sum\limits_{i=0}^k |a_i z_i|_{n_i} y_i$，其中 $|a_i z_i|_{n_i} = a_i z_i \bmod n_i$}

$ $

由于这样的 $x$ 在 $\bmod \text{ } N$ 中唯一，因此 $x \bmod N$ 与 $(a_0, a_1, a_2, \cdots, a_k)$ 之间存在同构映射。

$ $

此外，对于所有 $0 \leq i \leq k$，都有 $y_iz_i \equiv (y_iz_i)^2 \bmod N$ \text{ }
\end{tcolorbox}

\begin{proof}

$ $

\begin{enumerate}
\item 给定 $x = \sum\limits_{i=0}^k a_i y_i z_i$，对于每个满足 $0 \leq i \leq k$ 的 $i$，计算 $x \bmod n_i$：

$x \bmod n_i = \sum\limits_{j=0}^k a_j y_j z_j \bmod n_i$

$= a_0 y_0 z_0 + a_1 y_1 z_1 + a_2 y_2 z_2 + \cdots + a_k y_k z_k \bmod n_i$

$= a_i y_i z_i \bmod n_i$ \textcolor{red}{ $\rhd$ 因为对于所有 $j \neq i$，都有 $y_j \equiv 0 \bmod n_i$，这些项都是 $n_i$ 的倍数}

$= a_i$ \textcolor{red}{ $\rhd$ 因为 $y_i z_i \equiv y_i y_i^{-1} \equiv 1 \bmod n_i$}

$ $

因此，$x$ 在各模数 $n_0, n_1, n_2, \cdots, n_k$ 下都与 $a_0, a_1, a_2, \cdots, a_k$ 同余。


$ $

此外，以下关系也成立：

$x \bmod n_i = \sum\limits_{j=0}^k |a_j z_j|_{n_j} y_j \text{ } (\bmod n_i)$

$= |a_0 z_0|_{n_0} y_0 + |a_1 z_1|_{n_1} y_1 + |a_2 z_2|_{n_2} y_2 + \cdots + |a_k z_k|_{n_k} y_k \bmod n_i$

$= |a_i z_i|_{n_i} y_i \bmod n_i$ 

$= a_i$ 

$ $

\item 接下来，证明 $x$ 在模 $N$ 下是唯一解。假设存在两个解 $x$ 和 $x'$，满足：

$x \equiv x' \equiv a_0 \bmod n_0$

$x \equiv x' \equiv a_1 \bmod n_1$

$x \equiv x' \equiv a_2 \bmod n_2$

\text{ } $\vdots$

$x \equiv x' \equiv a_k \bmod n_k$

$ $

根据模同余的定义，有 $ n_0 \mid (x - x') , n_1 \mid (x - x') , \text{ } n_2 \mid (x - x'), \cdots, \text{ } n_k \mid (x - x')$。

由于 $n_0, n_1, n_2, \cdots, n_k$ 两两互素，必有 $n_0n_1n_2n_3\cdots n_k \mid (x - x')$，或等价地，$N \mid (x - x')$。这意味着 $x \equiv x' \bmod N$。因此，$x$ 在模 $N$ 下是唯一解。

$ $

\item 下面证明，对于所有 $0 \leq i \leq k$，均有 $y_iz_i \equiv (y_iz_i)^2 \bmod N$ \text{ }。

$ $

在模 $n_i$ 下，$y_iz_i \equiv 1 \bmod n_i$，因为 $z_i$ 是 $y_i$ 模 $n_i$ 的逆元。对于所有满足 $i \neq j$ 的其他模数 $n_j$，有 $y_iz_i \equiv 0 \bmod n_j$，因为 $y_i = \dfrac{N}{n_i}$，所以 $n_j$ 整除 $y_i$。

$ $

将 $(y_iz_i) \equiv 1 \bmod n_i$ 两边平方，得到 $(y_iz_i)^2 \equiv 1 \bmod n_i$。同样，将 $(y_iz_i) \equiv 0 \bmod n_j$ 两边平方，得到 $(y_iz_i)^2 \equiv 0 \bmod n_j$。

$ $

因此，$y_iz_i - (y_iz_i)^2 \equiv 0 \bmod n_i$ 且 $y_iz_i - (y_iz_i)^2 \equiv 0 \bmod n_j$。换言之，对于所有 $0 \leq j \leq k$，都有 $y_iz_i - (y_iz_i)^2 \equiv 0 \bmod n_j$。

$ $

再进行与步骤 2 相同的推理：由于各个两两互素的 $n_j$ 都整除 $y_iz_i - (y_iz_i)^2$，所以 $n_0n_1\cdots n_k = N$ 整除 $y_iz_i - (y_iz_i)^2$。因此，$y_iz_i - (y_iz_i)^2 \equiv 0 \bmod N$，即 $y_iz_i \equiv (y_iz_i)^2 \bmod N$。这对所有 $0 \leq i \leq k$ 都成立。

\end{enumerate}

\end{proof}


\subsection{应用：剩余数系统（RNS）}
\label{subsec:crt-application}

现代处理器通常以至多 64 位的数据字进行运算。对于超过 64 位的数据，可以使用中国剩余定理高效地处理：令每个两两互素的模数 $n_i$ 都满足 $\log_2n_i \leq 64$，其中 $N = n_0\cdot n_1\cdots\cdot n_k$，从而将大数 $a \bmod N$ 表示为 $\vec{a}_{\mathit{crt}} = (a_0, a_1, \cdots , a_k)$，其中 $a \equiv a_i \bmod n_i$。这样，对于任意两个大数 $a$ 和 $b \bmod N$，可以按如下方式计算 $a + b \bmod N$ 和 $a \cdot b \bmod N$：

\begin{itemize}
\item $a + b \equiv \sum\limits_{i=0}^k a_i y_i z_i + \sum\limits_{i=0}^k b_i y_i z_i \equiv \sum\limits_{i=0}^k (a_i y_i z_i + b_i y_i z_i) \equiv \sum\limits_{i=0}^k (a_i + b_i) y_i z_i \bmod N $

$ $

\item $a \cdot b \equiv \sum\limits_{i=0}^k a_i y_i z_i \cdot \sum\limits_{i=0}^k b_i y_i z_i \equiv \sum\limits_{i=0}^k (a_i \cdot b_i) (y_i z_i)^2  + \sum\limits_{i\neq j}^k (a_i \cdot b_j) y_i z_iy_j z_j \equiv \sum\limits_{i=0}^k (a_i \cdot b_i) (y_i z_i)^2$

\textcolor{red}{ $\rhd$ 对于满足 $i \neq j$ 的所有项 $y_iz_iy_jz_j$，其模 $N$ 的值都为 0，因为 $y_iy_j \bmod N \equiv 0$。\\
这是因为 $y_i = n_0n_1\cdots n_{i-1}n_{i+1}\cdots$ 且 $y_j = n_0n_1\cdots n_{j-1}n_{j+1}\cdots$。\\
因此，$y_iy_j$ 是 $N$ 的倍数。}

$ $

$\equiv \sum\limits_{i=0}^k (a_i \cdot b_i) (y_i z_i) \bmod N$

\textcolor{red}{ $\rhd$ 这是因为 $(y_i z_i) \equiv (y_i z_i)^2$，如定理~\ref*{sec:chinese-remainder}.1 的证明步骤 3 所示。}


\end{itemize}

$ $

因此，中国剩余定理给出如下有用的公式：

\begin{tcolorbox}[title={\textbf{\tboxtheorem{\ref*{sec:chinese-remainder}.2} 中国剩余定理的应用}}]

假设有两个大数 $a = \sum\limits_{i=0}^k a_i y_i z_i \bmod N$ 和 $b = \sum\limits_{i=0}^k b_i y_i z_i \bmod N$，其中 $N$ 是两两互素模数 $n_0 \cdot n_1 \cdots n_k$ 的乘积，则存在如下同构：

$a \xrightarrow{\sigma} \vec{a}_{\mathit{crt}} = (a_0, a_1, \cdots, a_k)$

$b \xrightarrow{\sigma} \vec{b}_{\mathit{crt}} = (b_0, b_1, \cdots, b_k)$

$ $

根据上述同构，有：

\begin{itemize}

\item $a + b \equiv \sum\limits_{i=0}^k (a_i + b_i) y_i z_i \bmod N \Longleftrightarrow \vec{a}_{\mathit{crt}} + \vec{b}_{\mathit{crt}} = ((a_i+b_i)\bmod n_i)_{i=0}^{k}$

\item $a \cdot b \equiv \sum\limits_{i=0}^k (a_i \cdot b_i) y_i z_i \bmod N \Longleftrightarrow \vec{a}_{\mathit{crt}} \odot \vec{b}_{\mathit{crt}} = ((a_i b_i)\bmod n_i)_{i=0}^{k}$


\end{itemize}



其中，各分量的加法或乘法都可以在模 $n_i$ 下独立完成。

\end{tcolorbox}


\clearpage

\section{泰勒级数}
\label{sec:taylor-series}
泰勒级数将解析函数在某一点附近表示为无穷幂级数，前提是该级数在相应邻域内收敛到这个函数。截断该级数，就得到函数的多项式近似。形式化地，函数的泰勒级数由其各阶导数在同一点处的取值构成无穷级数。给定函数 $f(X)$，其以 $a$ 为中心的泰勒级数表示为：

$
f(a) + \dfrac{f'(a)}{1!}(X - a) + \dfrac{f''(a)}{2!}(X - a)^2 + \dfrac{f'''(a)}{3!}(X - a)^3 + \cdots = \sum\limits_{d=0}^{\infty}\dfrac{f^{(d)}(a)}{d!}(X - a)^d
$

截断到次数 $D$ 得到第 $D$ 个泰勒多项式 $T_D(X)$。在泰勒级数表示 $f$ 的邻域内，误差在 $D\to\infty$ 时趋于零，但它不必在每次增加次数时都减小，也不必随到 $a$ 的距离单调变化。若实函数在连接 $a$ 与 $X$ 的区间上具有 $D+1$ 阶连续导数，泰勒定理给出
\[
|f(X)-T_D(X)|\le\frac{\sup_{u}|f^{(D+1)}(u)|}{(D+1)!}|X-a|^{D+1},
\]
其中上确界取遍该区间。在级数的收敛区域之外，提高次数一般不能保证精度改善。

\clearpage



\section{利用 FFT 与 NTT 高效计算多项式乘法}
\label{sec:ntt}
\textbf{- 参考资料：}
\href{http://web.cecs.pdx.edu/~maier/cs584/Lectures/lect07b-11-MG.pdf}{多项式与快速傅里叶变换（FFT）}~\cite{ntt}

\subsection{背景与动机}
\label{subsec:ntt-motivation}

本章计算模 $X^n+1$ 的负循环乘积，而非未经约减的 $2n-2$ 次乘积。使用商环约减时，求值节点必须为 $X^n+1$ 的根。卷积公式中，指标超出 $0,\ldots,n-1$ 的系数视为零。复杂度计算的是域上的算术运算次数，而非位操作次数或浮点误差。

给定两个 $(n-1)$ 次多项式：

$A(X) = \sum\limits_{i=0}^{n-1}a_iX^i$, \textcolor{white}{...} $B(X) = \sum\limits_{i=0}^{n-1}b_iX^i$

多项式乘积 $C(X) = A(X)\cdot B(X)$ 的计算方式如下：

$C(X) = \sum\limits_{i=0}^{2n-2}c_iX^{i}$，其中 $c_i = \sum\limits_{k=0}^{i}a_kb_{i-k}$

计算 $\vec{c} = (c_0, c_1, \cdots, c_{2n-2})$ 的这一运算也称为 $\vec a$ 与 $\vec b$ 的卷积，记为 $\vec{c} = \vec{a} \otimes \vec{b}$。其时间复杂度，即两个数之间乘法运算的总次数，为 $O(n^2)$。

另一种多项式乘法方法基于\textbf{点值表示}。次数为 $(n-1)$ 或更低的多项式 $A(X)$ 的点值表示，是由 $n$ 个坐标点 $\{(x_0, y_0), (x_1, y_1), \cdots (x_{n-1}, y_{n-1})\}$ 组成的集合，其中各个 $x_i$ 是互不相同的 $X$ 坐标，而各个 $y_i$ 不必是互不相同的 $Y$ 坐标。给定次数为 $(n-1)$ 或更低的多项式的点值表示，可以通过多项式插值（见 \autoref{sec:polynomial-interpolation}）恢复该多项式。
将次数为 $(n-1)$ 或更低的多项式 $A(X)$ 和 $B(X)$ 的点值表示记为：

$A(X)$ : $\bm ( ({x}_0, {y}_0^{\langle a \rangle}), ({x}_1, {y}_1^{\langle a \rangle}), \cdots ({x}_{n-1}, {y}_{n-1}^{\langle a \rangle}) \bm )$

$B(X)$ : $\bm ( ({x}_0, {y}_0^{\langle b \rangle}), ({x}_1, {y}_1^{\langle b \rangle}), \cdots ({x}_{n-1}, {y}_{n-1}^{\langle b \rangle}) \bm )$


于是，多项式 $C(X) = A(X) \cdot B(X)$ 的点值表示，可以通过对 $A(X)$ 和 $B(X)$ 的点值表示中的 $y$ 值作 Hadamard 积得到（见 \autoref{subsec:vector-arithmetic} 中的定义~\ref*{subsec:vector-arithmetic}），如下所示：

$C(X)$：$\bm ( ({x}_0, {y}_0^{\langle c \rangle}), ({x}_1, {y}_1^{\langle c \rangle}), \cdots ({x}_{n-1}, {y}_{n-1}^{\langle c \rangle}) \bm )$，其中 ${y}_i^{\langle c \rangle} = {y}_i^{\langle a \rangle} \cdot {y}_i^{\langle b \rangle}$

然而，仅凭这 $n$ 个坐标点，无法恢复多项式 $C(X)$，因为 $C(X)$ 的次数为 $2n-2$，或者低于 $2n-2$。但如果将所有多项式，包括 $A(X), B(X)$ 和 $C(X)$，都视为多项式环 $\mathbb{R}[X]/ (X^n + 1)$（或 $\mathbb{Z}_p[X]/(X^n + 1)$）中的元素，就可以把 $(2n-2)$ 次多项式 $C(X)$ 约减为环中与之同余、次数为 $(n-1)$ 或更低的多项式。此时，$C(X)$ 的 $n$ 个坐标点就足以恢复 $C(X)$。

不过，这种新方法的时间复杂度仍为 $O(n^2)$。两个多项式点值表示之间的 Hadamard 积仅需 $O(n)$，但在 $n$ 个互异 $x$ 值处对多项式求值需要 $O(n^2)$，因为每个多项式有 $n$ 项，需要对 $n$ 个互异 $x$ 值分别计算各项。用于恢复 $C(X)$ 的多项式插值也需要 $O(n^2)$。

为解决这一效率问题，本节介绍一种高效的多项式求值技术：它能在 $O(n \log n)$ 时间内，求出一个多项式在 $n$ 个互异单位根处的值。这类技术分为快速傅里叶变换（FFT）和数论变换（NTT）。两者在方法上几乎相同，区别在于 FFT 使用复数上的多项式环（见 \autoref{sec:roots}），而 NTT 使用有限域上的多项式环，例如素数模整数域（见 \autoref{sec:cyclotomic-polynomial-integer-ring}）。基于 FFT 或 NTT 的多项式乘法包含三个步骤：(1) 正向 FFT 或 NTT；(2) 点值乘法；(3) 逆 FFT 或 NTT。

\subsection{正向 FFT 或 NTT}
\label{subsec:ntt-forward}

基二递归取 $n=2^L$。对于有限域 NTT，还要求素数 $p$ 满足 $2n\mid p-1$，并在 $\mathbb Z_p^\times$ 中选取阶为 $2n$ 的 $\omega$。这些条件保证节点互异，且 $n$ 及几何级数中的所有非零分母可逆。复数变换则取 $\omega=e^{\pi i/n}$。

对于 FFT，使用多项式环 $\mathbb{R}[X]/ (X^n + 1)$；对于 NTT，使用 $\mathbb{Z}_p[X]/ (X^n + 1)$，其中 $X^n + 1$ 为分圆多项式。目标多项式的求值坐标 $x$ 选为 $X^n + 1$ 的 $n$ 个互异根，即 $\omega^1$、$\omega^3, \ldots, \omega^{2n-1}$，其中 $\omega$ 为 $2n$ 次本原单位根。于是，多项式 $A(X)$ 的点值表示为 $\bm ( ({x}_0, {y}_0^{\langle a \rangle}), ({x}_1, {y}_1^{\langle a \rangle}), \cdots ({x}_{n-1}, {y}_{n-1}^{\langle a \rangle}) \bm )$，其中：

${y}_i^{\langle a \rangle} = A(\omega^{2i + 1}) = \sum\limits_{j=0}^{n-1}a_{j}\cdot (\omega^{2i + 1})^j = \sum\limits_{j=0}^{n-1}a_{j}\cdot \omega^{(2i + 1)\cdot j}$

$ $

将向量 $\vec{y}^{\langle a \rangle} = (y_0^{\langle a \rangle}, y_1^{\langle a \rangle}, \cdots, y_{n-1}^{\langle a \rangle})$ 称为系数向量 $\vec{a} = (a_0, a_1, \cdots, a_{n-1})$ 的离散傅里叶变换（DFT），记为 $\vec{y}^{\langle a \rangle} = \textsf{DFT}(\vec{a})$。如 \autoref{subsec:ntt-motivation} 所述，直接计算 DFT 需要 $O(n^2)$ 时间，因为对于一个含 $n$ 项的多项式，需要在 $n$ 个互异的 $X$ 值处求值。

\subsubsection{基本思路}
\label{subsec:ntt-forward-overview}

FFT 或 NTT 是计算 DFT 的改进方法，将时间复杂度从 $O(n^2)$ 降低到 $O(n \log n)$。FFT 的基本思路是：将次数为 $(n-1)$ 或更低的待求值多项式 $A(X)$，拆成两个次数约为原来一半的多项式 $A_0(X)$ 和 $A_1(X)$，如下所示：


$A(X) = a_0 + a_1X + a_2X^2 + \cdots + a_{n-1}X^{n - 1}$

$\textcolor{white}{A(X) }= A_0(X^2) + X \cdot A_1(X^2)$
$A_0(X) = a_0 + a_2X + a_4X^2 + \cdots + a_{n-2}X^{\frac{n}{2} - 1}$

$A_1(X) = a_1 + a_3X + a_5X^2 + \cdots + a_{n-1}X^{\frac{n}{2} - 1}$

将一个多项式拆成两个次数减半的多项式的上述过程，称为 Cooley-Tukey 步骤。将 $A(X)$ 拆成低次多项式 $A_0(X)$ 和 $A_1(X)$ 后，在 $2n$ 次本原单位根的奇数次幂 $\{\omega^1, \omega^3, \omega^5, \cdots, \omega^{2n-1}\}$ 处求 $A(X)$ 的值，等价于在所列的 $n$ 个平方节点（其中仅有 $n/2$ 个互异节点） $\{(\omega^2)^1, (\omega^2)^3, (\omega^2)^5, \ldots, (\omega^2)^{2n-1}\}$ 处求 $A_0(X)$ 和 $A_1(X)$ 的值，再计算 $A_0(X^2) + X\cdot A_1(X^2)$。但是，$2n$ 次本原单位根 $\omega$ 的阶为 $2n$，即 $\omega^{2n} = 1$，并且对于所有整数 $1\le m < 2n$，都有 $\omega^m \neq 1$。因此，$\{(\omega^2)^1, (\omega^2)^3, (\omega^2)^5, \ldots, (\omega^2)^{2n-1}\}$ 的后半部分只是前半部分的重复。这意味着，只需分别在 $\dfrac{n}{2}$ 个互异的 $x$ 坐标处求 $A_0(X)$ 和 $A_1(X)$ 的值，而无需在 $n$ 个互异坐标处求值，因为另一半的求值结果与前半部分相同：代入多项式的 $x$ 本身就相同。

继续递归地将 $A_0(X)$ 和 $A_1(X)$ 拆为次数减半的多项式，并在数量减半，即 $n/2$ 个 $n$ 次单位根处求值。拆分总轮数为 $\log n$，每轮单位根与系数相乘的次数至多为 $n$，总计为 $O(n \log n)$。

\subsubsection{具体过程}
\label{subsec:ntt-forward-details}

假设多项式环为 $\mathbb{Z}_{p}[X] / (X^8 + 1)$，即素数 $p$ 对应的有限域上的多项式环，或为复数上的多项式环 $\mathbb{R}[X] / (X^8 + 1)$。
用 $\omega$ 表示一个 $(2n=16)$ 次本原单位根，
则 $X^8 + 1$ 的 8 个互异根为 $\{\omega^1, \omega^3, \omega^5, \omega^7, \omega^9, \omega^{11}, \omega^{13}, \omega^{15}\}$。

$ $

将待求值的目标多项式定义为：

$A(X) = a_0 + a_1X + a_2X^2 + a_3X^3 + a_4X^4 + a_5X^5 + a_6X^6 + a_7X^7$

$ $

利用 Cooley-Tukey 步骤，将这个 7 次多项式拆成以下两个 3 次多项式：

$A_0(X) = a_0 + a_2X + a_4X^2 + a_6X^3$

$A_1(X) = a_1 + a_3X + a_5X^2 + a_7X^3$

$A(X) = A_0(X^2) + X \cdot A_1(X^2)$

$ $

再递归地将上述两个 3 次多项式拆成 1 次多项式：


$A_{0,0}(X) = a_0 + a_4X$, \textcolor{white}{...} $A_{0,1}(X) = a_2 + a_6X$

$A_0(X) = A_{0,0}(X^2) + X\cdot A_{0,1}(X^2)$

$ $

$A_{1,0}(X) = a_1 + a_5X$, \textcolor{white}{...} $A_{1,1}(X) = a_3 + a_7X$

$A_1(X) = A_{1,0}(X^2) + X\cdot A_{1,1}(X^2)$

$ $

$A(X) = A_0(X^2) + X \cdot A_1(X^2)$

$\mathcolor{white}{A(X)} = \underbrace{\underbrace{(\underbrace{A_{0,0}(X^4)}_{\text{FFT 第 1 层}} + X^2\cdot \underbrace{A_{0,1}(X^4)}_{\text{FFT 第 1 层}})}_{\text{FFT 第 2 层}} + X \cdot \underbrace{(\underbrace{A_{1,0}(X^4)}_{\text{FFT 第 1 层}} + X^2\cdot \underbrace{A_{1,1}(X^4)}_{\text{FFT 第 1 层}})}_{\text{FFT 第 2 层}}}_{\text{FFT 第 3 层}}$

$ $

为了在 $n$ 个互异根 $X = \{\omega^1, \omega^3, \ldots, \omega^{15}\}$ 处求出 $A(X)$ 的值，从第 $1 \leq l \leq 3$ 层开始，对上述公式的每一层 FFT 在 $X = \{\omega^1, \omega^3, \cdots, \omega^{15}\}$ 处求值。

$ $

\para{FFT 第 $\bm{l = 1}$ 层：}在 $X = \{\omega^1, \omega^3, \cdots, \omega^{15}\}$ 处求出 $A_{0,0}(X^4)$、$A_{0,1}(X^4)$、$A_{1,0}(X^4)$ 和 $A_{1,1}(X^4)$ 的值。但是，将 $X = \{\omega^1, \omega^3, \cdots, \omega^{15}\}$ 代入 $X^4$ 后，只会得到两个不同的值 $\omega^4$ 和 $\omega^{12}$，它们对应于 $X^2+1$ 的根。这是因为 $\omega$ 的阶为 $2n$，即 $\omega^{2n} = 1$，因此 $(\omega^1)^4 = (\omega^5)^4 = (\omega^9)^4 = (\omega^{13})^4 = \omega^4$ 且 $(\omega^3)^4 = (\omega^7)^4 = (\omega^{11})^4 = (\omega^{15})^4 = \omega^{12}$。
所以，只需在 2 个而非 8 个互异的 $x$ 值处，分别求出 $A_{0,0}(X^4)$、$A_{0,1}(X^4)$、$A_{1,0}(X^4)$ 和 $A_{1,1}(X^4)$ 的值。每次求值只需要常数次算术运算：1 次乘法和 1 次加法。由于共有 4 个待求值的多项式，即 $A_{0, 0}(X^4), A_{0, 1}(X^4), A_{1, 0}(X^4), A_{1, 1}(X^4)$，这一层共进行 $4 \cdot 2 = 8$ 次 FFT 基本计算。

$ $

\para{FFT 第 $\bm{l = 2}$ 层：}利用 FFT 第 1 层的求值结果作为基础，在 $X = \{\omega^1, \omega^3, \cdots, \omega^{15}\}$ 处求出 $A_{0}(X^2)$ 和 $A_{1}(X^2)$ 的值。但是，将 $X = \{\omega^1, \omega^3, \cdots, \omega^{15}\}$ 代入 $X^2$ 后，只会得到 4 个不同的值 $\omega^2$、$\omega^6$、$\omega^{10}$ 和 $\omega^{14}$，它们对应于 $X^4+1$ 的根。这是因为 $\omega$ 的阶为 $2n$，即 $\omega^{2n} = 1$，因此 $(\omega^1)^2 = (\omega^9)^2$、$(\omega^3)^2 = (\omega^{11})^2$、$(\omega^5)^2 = (\omega^{13})^2$ 和 $(\omega^7)^2 = (\omega^{15})^2$。所以，只需在 4 个而非 8 个互异的 $x$ 值处，分别求出 $A_{0}(X^2)$ 和 $A_{1}(X^2)$ 的值。每次求值只需常数次算术运算：1 次乘法和 1 次加法。这里复用了 FFT 第 1 层的结果，并且第 2 层的计算结构与第 1 层相同。待求值的多项式共有 2 个，即 $A_{0}(X^2), A_{1}(X^2)$，因此这一层共进行 $2 \cdot 4 = 8$ 次 FFT 基本计算。

$ $

\para{FFT 第 $\bm{l = 3}$ 层：}利用 FFT 第 2 层的求值结果作为基础，在 $X = \{\omega^1, \omega^3, \cdots, \omega^{15}\}$ 处求出 $A(X)$ 的值。在最后一层，所有 8 个 $X$ 值都互不相同，因此都需要求值；每次求值仍只需要常数次算术运算：1 次乘法和 1 次加法。待求值的多项式仅有 1 个，即 $A(X)$，因此这一层共进行 $1 \cdot 8 = 8$ 次 FFT 基本计算。

$ $

\para{一般化：}假设目标多项式的次数至多为 $n-1$，包含 $n$ 项，定义 $L = \log_2 n$ 为 FFT 的总层数。正向 FFT 共需 $L$ 层，其中第 $l$ 层需要对 $2^{L - l}$ 个多项式，分别在 $2^l$ 个互异的 $X$ 值处求值。因此，正向 FFT 的基本计算总次数为 $\sum_{l=1}^{L}(2^{L-l} \cdot 2^l) = L \cdot 2^L = n \log n$
%$\log (n) \cdot (2^{L - l} \cdot 2^l) = 2^L \log n = n \log n$
。所以，正向 FFT 的时间复杂度为 $O(n \log n)$。



使用 FFT 时，层数每降低一级，需要求值的 $x$ 点数就减半，而待求值的多项式数量加倍；两者相互抵消，使每层的计算量保持为 $O(n)$。由于共有 $\log n$ 层，总时间复杂度为 $O(n \log n)$。这一优化的关键在于求值坐标 $x$ 的特殊性质：它的幂，即 $\omega^i$，具有周期性。为保证这种周期性，本章的负循环 FFT 将求值点选为 $2n$ 次本原单位根的奇数次幂；通常的循环 DFT 则使用 $n$ 次本原单位根的所有幂。


\subsection{逐点乘法}
\label{subsec:pointwise-multiplication}

对多项式 $A(X)$ 和 $B(X)$ 分别应用正向 FFT（见 \autoref{subsec:ntt-forward}），得到 $\vec{y}^{\langle a \rangle}$ 和 $\vec{y}^{\langle b \rangle}$ 后，就可以通过 Hadamard 积 $ \vec{y}^{\langle c \rangle} = \vec{y}^{\langle a \rangle} \odot \vec{y}^{\langle b \rangle}$，在 $O(n)$ 时间内计算出 $C(X) = A(X) \cdot B(X)$ 的点值表示，如 \autoref{subsec:ntt-motivation} 所述。

\subsection{逆 FFT 或 NTT}
\label{subsec:ntt-backward}

至此，已经计算出：

$C(X)$：$\bm ( ({x}_0, {y}_0^{\langle c \rangle}), ({x}_1, {y}_1^{\langle c \rangle}), \cdots ({x}_{n-1}, {y}_{n-1}^{\langle c \rangle}) \bm )$，其中 ${y}_i^{\langle c \rangle} = {y}_i^{\langle a \rangle} \cdot {y}_i^{\langle b \rangle}$

$ $

最后一步是将 ${y}_i^{\langle c \rangle}$ 转回 $C(X)$ 的多项式系数 $\vec{c}$。这一逆向操作称为逆 FFT。

给定 $(n-1)$ 次多项式
$C(X) = \sum\limits_{i=0}^{n-1}c_iX^i$，正向 FFT 在计算上等价于在 $X^n + 1$ 的 $n$ 个互异根（即 $2n$ 次单位根）处对多项式求值，如下所示：

${y}_i^{\langle c \rangle} = C(\omega^{2i+1}) = \sum\limits_{j=0}^{n-1}c_{j}\cdot (\omega^{2i + 1})^j = \sum\limits_{j=0}^{n-1}c_{j}\cdot \omega^{(2i+1)j}$

上述求值等价于以下矩阵与向量的乘法：

$ $



$\vec y^{\langle c \rangle} = W \cdot \vec{c}$，\text{ }其中 $W = \begin{bmatrix}
(\omega^1)^0 & (\omega^1)^1 & \cdots & (\omega^1)^{n-1}\\
(\omega^3)^0 & (\omega^3)^1 & \cdots & (\omega^3)^{n-1}\\
\vdots & \vdots & \ddots & \vdots \\
(\omega^{2n-1})^0 & (\omega^{2n-1})^1 & \cdots & (\omega^{2n-1})^{n-1}\\
\end{bmatrix} $，\text{ }$\vec{c} = (c_0, c_1, \cdots, c_{n-1})$

$ $

将 $W$ 的各元素记为 $(W)_{i,j} = \omega^{(2i+1)j}$。逆 FFT 就是将上述计算反过来。为实现这种反演，需要找到逆矩阵 $W^{-1}$，使得 $W^{-1} \cdot \vec y^{\langle c \rangle} = W^{-1} \cdot (W \cdot \vec{c}) = (W^{-1} \cdot W) \cdot \vec{c} = I_n \cdot \vec{c} =  \vec{c}$。
可以构造如下逆矩阵 $W^{-1}$：

$(W^{-1})_{j,k} = n^{-1}\cdot \omega^{-(2k+1)j}$

$ $

下面说明为何 $W^{-1} \cdot W = I_n$ 成立。$W^{-1} \cdot W$ 的每个元素计算为：

$(W^{-1} \cdot W)_{j, i} = \sum\limits_{k=0}^{n-1}(n^{-1}\cdot \omega^{-(2k+1)j} \cdot \omega^{(2k+1)i}) = n^{-1}\cdot\sum\limits_{k=0}^{n-1}\omega^{(2k+1)(i-j)}$

要使 $W^{-1} \cdot W$ 等于 $I_n$，需要满足：

\[
(W^{-1} \cdot W)_{j, i} = \begin{cases}
\text{1，若 } j = i \\
\text{0，若 } j \neq i \\
\end{cases}
\]

如果 $j = i$，则上述条件成立，因为：

$(W^{-1} \cdot W)_{j, i} = n^{-1}\cdot\sum\limits_{k=0}^{n-1}\omega^{(2k+1)(0)} =  n^{-1}\cdot\sum\limits_{k=0}^{n-1}1 = n^{-1} \cdot n = 1$

$ $

对于 $j \neq i$ 的情形，使用等比数列求和公式 $\sum\limits_{i=0}^{n-1}x^i = \dfrac{x^n - 1}{x - 1}$，证明如下：

\begin{tcolorbox}[title={\textbf{\tboxtheorem{\ref*{subsec:ntt-backward}} 等比数列求和公式}}]

设等比数列之和为 $S_n = 1 + x + x^2 + \cdots + x^{n - 1}$

则 $x \cdot S_n = x + x^2 + x^3 + \cdots + x^{n}$

$x \cdot S_n - S_n = (x + x^2 + x^3 + \cdots + x^{n}) - (1 + x + x^2 + \cdots + x^{n - 1}) = x^n - 1$

$S_n\cdot (x - 1) = x^n - 1$

$S_n = \dfrac{x^n - 1}{x - 1}$ \textcolor{red}{ $\rhd$ 其中要求 $x \neq 1$}

\end{tcolorbox}

目标是计算 $\sum\limits_{k=0}^{n-1}\omega^{(2k+1)(i-j)} = \omega^{i-j} \sum\limits_{k=0}^{n-1}(\omega^{2(i-j)})^k$。在等比数列求和公式中取 $x = \omega^{2(i-j)}$：

$\omega^{i-j} \sum\limits_{k=0}^{n-1}(\omega^{2(i-j)})^k = \omega^{i-j} \dfrac{(\omega^{2(i-j)})^n - 1}{\omega^{2(i-j)} - 1}$ 

$= \omega^{i-j} \dfrac{(\omega^{2n})^{i - j} - 1}{\omega^{2(i-j)} - 1}$

$= \omega^{i-j} \dfrac{(1)^{i - j} - 1}{\omega^{2(i-j)} - 1}$ \textcolor{red}{ $\rhd$ 因为 $\omega$ 的阶为 $2n$，所以 $\omega^{2n}=1$}

\textcolor{red}{ $\rhd$ 这里分母不能为 0。由于 $i \neq j$ 且 $|i-j| < n$，指数 $2(i-j)$ 不是 $2n$ 的倍数。}

$= 0$

$ $

因此，$(W^{-1} \cdot W)_{j, i}$ 在 $j = i$ 时为 1，在 $j \neq i$ 时为 0。所以，逆 FFT 可以计算为 $\vec{c} = W^{-1} \cdot y^{\langle c \rangle}$，其中：

$c_i = \sum_{j=0}^{n-1} (W^{-1})_{i,j} \cdot y_j^{\langle c \rangle}$

$ = \sum\limits_{j=0}^{n-1} \left( n^{-1} \cdot \omega^{-(2j+1)i} \right) \cdot y_j^{\langle c \rangle}$ \textcolor{red}{ $\rhd$ 因为 $(W^{-1})_{j,k} = n^{-1} \cdot \omega^{-(2k+1)j}$，且 $(W^{-1})_{i,j} = n^{-1} \cdot \omega^{-(2j+1)i}$}

$ = n^{-1}\cdot \sum\limits_{j=0}^{n-1}{y}_j^{\langle c \rangle} \cdot \omega^{-(2j+1)i}$

$ = n^{-1} \sum\limits_{j=0}^{n-1} y_j^{\langle c \rangle} \cdot (\omega^{-2ji} \cdot \omega^{-i}) $ \textcolor{red}{ $\rhd$ 因为 $\omega^{-(2j+1)i} = \omega^{-2ji - i} = \omega^{-2ji} \cdot \omega^{-i}$}

$= n^{-1} \omega^{-i} \sum\limits_{j=0}^{n-1} {y}_j^{\langle c \rangle} (\omega^{-2})^{ji}$

$ $


复用正向过程中介绍的递归拆分技术（见 \autoref{subsec:ntt-forward-details}），逆变换同样可以达到 $O(n \log n)$ 的时间复杂度。


\begin{comment}
\begin{algorithm}[t]
\algrenewcommand\algorithmicindent{1em} % Adjust indentation herei.e., 
\begin{algorithmic}[1]
\State \textsf{FFT}$(\vec{a}):$ %\textcolor{red}{ $\rhd$ an RLWE ciphertext storing $k$ zeros and non-zero random numbers, whose ratio varies}
\State $n' \gets$ $\vec{a}$ 的长度
\If{$n' = 1$}
    \State \textbf{返回} $a$
\EndIf
\State $\omega = e^{\frac{2\pi}{k}}$ \textcolor{red}{ $\rhd$ 在 NTT 情形下为 $\omega = g^{\frac{p-1}{n}}$}
\State $\omega' = 1$
\State $\vec{a}^{\langle 0 \rangle} = (a_0, a_2, \cdots, a_{n'-2})$
\State $\vec{a}^{\langle 1 \rangle} = (a_1, a_3, \cdots, a_{n'-1})$
\State $\vec{y}^{\langle 0 \rangle} = \textsf{FFT}(\vec{a}^{\langle 0 \rangle})$
\State $\vec{y}^{\langle 1 \rangle} = \textsf{FFT}(\vec{a}^{\langle 1 \rangle})$
\For{$i \in \{0, 1, \cdots \frac{n'}{2} - 1\}$}
    \State $y_i \gets y_i^{\langle 0 \rangle} + \omega \cdot y_i^{\langle 1 \rangle}$
    \State $y_{i+\frac{n}{2}} \gets y_i^{\langle 0 \rangle} + \omega \cdot y_i^{\langle 1 \rangle}$
    \State $\omega' = \omega' \cdot \omega$
\EndFor
\State \textbf{返回} $\vec{y}$
\end{algorithmic}
\caption{快速傅里叶变换（FFT）}
\label{alg:fft}
\end{algorithm}
\end{comment}

\clearpage





\part{后量子密码}
\label{part:pqc}

\renewcommand{\thesection}{B-\arabic{section}}
\setcounter{section}{0}

本部分介绍基于格的密码体制：LWE、RLWE、GLWE、GLev 和 GGSW。它们是构建 FHE 方案的重要基础。


$ $

\begin{tcolorbox}[
    title = \textbf{所需背景知识},    % box title
    colback = white,    % light background; tweak to taste
    colframe = black,  % frame colour
    boxrule = 0.8pt,     % line thickness
    left = 1mm, right = 1mm, top = 1mm, bottom = 1mm % inner padding
]



\begin{itemize}
\item \autoref{sec:group}: \nameref{sec:group}
\item \autoref{sec:field}: \nameref{sec:field}
\item \autoref{sec:polynomial-ring}: \nameref{sec:polynomial-ring}
\end{itemize}

\end{tcolorbox}

$ $

\clearpage

\section{格密码}
\label{sec:lattice}
\textbf{- 参考资料：}
\href{https://mysite.science.uottawa.ca/mnevins/papers/StephenHarrigan2017LWE.pdf}{格密码与
带误差学习问题}~\cite{lattice-crypto}


$ $

格密码在明确的计算困难假设及参数选择下，为后量子安全提供候选方案；不能无条件地声称它能够抵抗所有量子攻击。本节介绍后续使用的 LWE（带误差学习）、RLWE（环带误差学习）、GLWE、GLev 和 GGSW 构造。





\subsection{概述}
\label{subsec:lattice-overview}

本概述中的所有等式均在 $\mathbb Z_q$ 上成立。掩码独立均匀采样；误差先从指定分布采样为小整数，仅在运算时模 $q$ 约减。

假设有一个未知的 $k$ 维向量 $\vec{s}$ 作为秘密密钥，以及许多公开已知的 $k$ 维向量 $\vec{a}^{\langle i \rangle}$。

再假设有大量如下点积 $\vec{s} \cdot \vec{a}^{\langle i \rangle}$：

$\vec{s} \cdot \vec{a}^{\langle 0 \rangle} = s_0\cdot a_{0}^{\langle 0 \rangle} + s_1\cdot a_{1}^{\langle 0 \rangle} + \cdots + s_{k-1}\cdot a_{k-1}^{\langle 0 \rangle} = b^{\langle 0 \rangle}$

$\vec{s} \cdot \vec{a}^{\langle 1 \rangle} = s_0\cdot a_{0}^{\langle 1 \rangle} + s_1\cdot a_{1}^{\langle 1 \rangle} + \cdots + s_{k-1}\cdot a_{k-1}^{\langle 1 \rangle} = b^{\langle 1 \rangle}$

$\vec{s} \cdot \vec{a}^{\langle 2 \rangle} = s_0\cdot a_{0}^{\langle 2 \rangle} + s_1\cdot a_{1}^{\langle 2 \rangle} + \cdots + s_{k-1}\cdot a_{k-1}^{\langle 2 \rangle} = b^{\langle 2 \rangle}$

\text{ } $\vdots$

$ $

在无噪声情形下，若 $k$ 个样本的系数矩阵在 $\mathbb Z_q$ 上可逆，它们就能确定 $\vec s$。当 $q$ 为素数时，这要求矩阵满秩；当模数为合数时，行列式必须为单位。仅有 $k$ 个方程，并不能保证方程独立或唯一恢复。

但是，如果在上述每个方程中都随机加入一个未知的小噪声 $e^{\langle i \rangle}$，即误差，得到：

$b^{\langle 0 \rangle} = s_0\cdot a_{0}^{\langle 0 \rangle} + s_1\cdot a_{1}^{\langle 0 \rangle} + \cdots + s_{k-1}\cdot a_{k-1}^{\langle 0 \rangle} + e^{\langle 0 \rangle} \pmod q$

$b^{\langle 1 \rangle} = s_0\cdot a_{0}^{\langle 1 \rangle} + s_1\cdot a_{1}^{\langle 1 \rangle} + \cdots + s_{k-1}\cdot a_{k-1}^{\langle 1 \rangle} + e^{\langle 1 \rangle} \pmod q$

$b^{\langle 2 \rangle} = s_0\cdot a_{0}^{\langle 2 \rangle} + s_1\cdot a_{1}^{\langle 2 \rangle} + \cdots + s_{k-1}\cdot a_{k-1}^{\langle 2 \rangle} + e^{\langle 2 \rangle} \pmod q$

\text{ } $\vdots$


$ $

噪声使得这些方程不能直接作为关于 $\vec s$ 的精确线性方程求解。若 $k$ 个误差各有 $r$ 种可能值，枚举全部 $r^k$ 个误差向量是一种可行攻击。这只是该枚举方法的工作量上界，并非所有算法的下界，也不是安全性证明。额外样本可能帮助攻击者，而且攻击者始终可以忽略它们。LWE 困难性是针对明确维数、模数、秘密与误差分布以及样本数量的计算假设；过小或具有特殊结构的误差可能使某些实例易解。

基于这一直观认识，格密码所依赖的数学困难问题可描述如下：

$ $

\begin{tcolorbox}[title={\textbf{\tboxlabel{\ref*{subsec:lattice-overview}} LWE（带误差学习）与 RLWE 问题}}]
\textbf{\underline{LWE 问题}}

考虑如下形式的样本：$b = \vec{s} \cdot\vec{a} + e$，其中 $e$ 是后面将介绍的小噪声。

对每个样本，从 $\mathbb Z_q^k$ 中均匀选取 $\vec a$，并独立地从明确指定的分布 $\chi$ 采样小整数误差 $e$，随后将 $e$ 模 $q$ 约减。秘密 $\vec s$ 从其指定分布采样一次并重复使用。误差并非在 $\mathbb Z_q$ 上均匀分布。安全实验允许获得关于安全参数的多项式数量的这类样本。

$(\vec{a}^{\langle 1 \rangle}, b^{\langle 1 \rangle})$，其中 $b^{\langle 1 \rangle} = \vec{s} \cdot \vec{a}^{\langle 1 \rangle} + e^{\langle 1 \rangle}$

$(\vec{a}^{\langle 2 \rangle}, b^{\langle 2 \rangle})$，其中 $b^{\langle 2 \rangle} = \vec{s} \cdot \vec{a}^{\langle 2 \rangle} + e^{\langle 2 \rangle}$

$(\vec{a}^{\langle 3 \rangle}, b^{\langle 3 \rangle})$，其中 $b^{\langle 3 \rangle} = \vec{s} \cdot \vec{a}^{\langle 3 \rangle} + e^{\langle 3 \rangle}$

\text{ } $\vdots$

假设攻击者拥有大量二元组 $(\vec{a}^{\langle i \rangle}, b^{\langle i \rangle})$。在这一设定下，以下困难问题构成 LWE（带误差学习）密码体制的安全基础：

$ $

\begin{itemize}
\item \textbf{搜索问题：}给定这些样本，恢复秘密 $\vec s$。搜索 LWE 假设认为，对于所选参数族，不存在以不可忽略概率成功的高效算法。
\item \textbf{判定问题：}区分 LWE 样本 $(\vec a,\langle\vec a,\vec s\rangle+e)$ 与 $\mathbb Z_q^k\times\mathbb Z_q$ 中独立均匀的二元组。判定 LWE 假设认为，对于所选参数族，每个高效区分器的优势都可忽略。
\end{itemize}

$ $

搜索到判定的归约需要前提；这两种表述并不对所有模数和分布无条件等价。例如，Regev 的引理 4.2 要求模数为关于维数多项式有界的素数，并满足其中的误差及预言机条件。其最坏情形格问题归约也有明确的噪声和模数要求，而且是量子归约。参见 \href{https://cims.nyu.edu/~regev/papers/qcrypto.pdf}{Regev, On Lattices, Learning with Errors, Random Linear Codes, and Cryptography} 的定理 1.1 及第 4 节。

$ $

\textbf{\underline{RLWE 问题}}

在 RLWE 的代数形式中，使用 $\mathbb Z_q[X]/(X^n+1)$ 上的 $B=AS+E$，多项式代表元的次数小于 $n$。搜索问题要求恢复 $S$，判定问题要求将这些二元组与均匀环元素对区分开。仅替换记号并不能转移所有 LWE 安全性定理：环、模数以及误差和秘密分布必须满足相应环 LWE 结论的前提。

\end{tcolorbox}




\noindent\textbf{小秘密的适用范围：}正规形归约可以把均匀 LWE 秘密改成从\emph{误差分布}中采样的秘密。对于素数幂模数，仅根据掩码选择可逆矩阵 $A_0$，对应 $\vec b_0=A_0\vec s+\vec e_0$。对新鲜行向量样本 $(\vec a,b)$，令
\[
\vec a'=-\vec aA_0^{-1},\qquad b'=b-\vec aA_0^{-1}\vec b_0
=\vec a'\vec e_0+e.
\]
新的掩码仍均匀，新秘密的分布为 $\chi^k$；均匀输入对也映射成均匀对。预处理材料是此归约的内部状态。在环版本中，同一恒等式要求初始环掩码可逆。参见 \href{https://web.cs.ucla.edu/~sahai/work/web/2009\%20Publications/ACPS09.pdf}{Applebaum--Cash--Peikert--Sahai，第 3.1 节引理 2}。该论证不会把高斯误差变成均匀三元秘密。因此，后续三元秘密构造使用针对实际分布的判定假设，除非另给分布匹配的归约。维数计数或“小”这个形容词都不能代替该归约。

\subsection{LWE 密码体制}
\label{subsec:lattice-scheme}

下面的正确性公式假设 $t\mid q$，取 $\Delta=q/t$，并以 $[x]_q$ 表示 $x\bmod q$ 的中心化整数代表元。舍入作用于该代表元，解码结果还需模 $t$ 约减。所有位移描述仅适用于 $\Delta$ 为二幂的情形；一般的缩放是乘以或除以 $\Delta$。完整密文是 $(\vec a,b)$，并非只有 $b$。

LWE 密码体制使用加密公式 $b = \vec{s} \cdot \vec{a} + \Delta \cdot m + e$，其中 $\vec{s}$ 为秘密密钥，$\vec{a}$ 为每次加密时选取的公开随机向量，$m$ 为明文，$e$ 为每次加密时从正态分布中随机选取的小噪声，$b$ 为密文。$\Delta$ 是明文 $M$ 的尺度因子，相当于将 $m$ 向左移动 $\text{log}_2\Delta$ 个二进制位。加密前，先将明文左移若干位，即 $\text{log}_2\Delta$ 位，为低位中的误差预留足够空间。

\begin{figure}[h!]
    \centering
  \includegraphics[width=0.8\linewidth]{figures/TFHE-fig1.pdf}
  \caption{LWE 中的明文缩放与加噪示意图：$\Delta \cdot m + e \in \mathbb{Z}_q$}
  \label{fig:scaling}
\end{figure}

\autoref{fig:scaling} 直观展示了 $\Delta \cdot m + e$：将明文 $m$ 左移 $\text{log}_2\Delta$ 位，并加入噪声 $e$。实际的加密和解密公式如下：

\begin{tcolorbox}[title={\textbf{\tboxlabel{\ref*{subsec:lattice-scheme}} 基于格的 LWE 密码体制}}]
\begin{itemize}
    \item \textbf{\underline{加密}：}$b^{\langle i \rangle} = \vec{s} \cdot \vec{a}^{\langle i \rangle} + \Delta \cdot m^{\langle i \rangle} + e^{\langle i \rangle}$，其中 $b^{\langle i \rangle}$ 和 $\vec{a}^{\langle i \rangle}$ 对攻击者也是公开已知的，而 $\vec{s}, m^{\langle i \rangle}, e^{\langle i \rangle}$ 则未知，仅秘密密钥持有者知晓。

    $ $
    
    \item \textbf{\underline{解密}：}$\left\lceil [b^{\langle i\rangle}-\vec s\cdot\vec a^{\langle i\rangle}]_q/\Delta\right\rfloor\bmod t=m^{\langle i\rangle}$ $\Big($，前提是 $|e^{\langle i \rangle}| < \dfrac{\Delta}{2}\Big)$。
\end{itemize}
\end{tcolorbox}

$\lfloor \rceil_\Delta$ 表示将数舍入到最接近的 $\Delta$ 的整数倍。例如，$\lfloor 16 \rceil_{10} = 20$ 表示将 16 舍入到最接近的 10 的整数倍。再如，$\lfloor 17 \rceil_{8} = 16$ 表示将 17 舍入到最接近的 8 的整数倍；由于 17 更接近 16 而非 24，因此结果为 16。


$ $

\noindent \textbf{正确性：}在解密过程中，计算 $b^{\langle i \rangle} - \vec{s} \cdot \vec{a}^{\langle i \rangle}$ 得到 $\Delta \cdot m^{\langle i \rangle}+ e^{\langle i \rangle}$，如 \autoref{fig:scaling} 所示。随后计算 $\lceil \Delta \cdot m^{\langle i \rangle} + e^{\langle i \rangle} \rfloor_{\Delta}$，即将数值舍入到最接近的 $\Delta$ 的整数倍；只要所加噪声满足 $|e^{\langle i \rangle}| < \frac{\Delta}{2}$，就会得到 $\Delta \cdot m^{\langle i \rangle}$。也就是说，如果噪声小于 $\frac{\Delta}{2}$，它就会在舍入过程中被消去。最后，将 $\Delta \cdot m^{\langle i \rangle}$ 右移 $\text{log}_2 \Delta$ 位，即可得到 $m^{\langle i \rangle}$。因此，只要保证 $|e^{\langle i \rangle}| < \frac{\Delta}{2}$，也就是噪声 $e^{\langle i \rangle}$ 必须小于相应阈值，就能够在解密的舍入过程中消去 $e^{\langle i \rangle}$，恢复原始的 $\Delta \cdot m^{\langle i \rangle}$。将 $m^{\langle i \rangle}$ 按 $\Delta$ 缩放有两个目的：(i) 加密时在低位为 $e^{\langle i \rangle}$ 留出空间，使噪声位与明文位不相互干扰，避免破坏明文；(ii) 解密时消去低位中的噪声 $e^{\langle i \rangle}$，同时保持明文 $m^{\langle i \rangle}$ 不受损。

$ $

\noindent \textbf{安全性：}固定实际秘密分布、误差分布、模数、维数和样本上限 $Q$。令 $\delta(T,Q)$ 表示时间为 $T$ 的区分器在真实 LWE 样本预言机与独立均匀二元组预言机之间的最大接受概率差。秘密重复使用，而每次查询的均匀掩码与误差均独立新鲜。这是判定假设，不仅是密钥恢复困难的断言。

在自适应左右加密游戏中，隐藏位 $\vartheta$ 从每次提交的消息对 $(m_0,m_1)$ 中选择一个消息。定义猜测偏差 $\epsilon=|\Pr[\widehat\vartheta=\vartheta]-1/2|$。归约自行选择 $\vartheta$，对每次查询从样本预言机取得 $(\vec a,z)$，返回 $(\vec a,z+\Delta m_\vartheta)$。真实样本下这恰好是加密；均匀样本下，即使给定此前全部记录，新鲜的 $z$ 加上任意消息偏移仍然均匀。因此逐步归纳可知，完整视图与 $\vartheta$ 独立。归约输出攻击者是否猜中，得到
\[
\epsilon\le\delta(T',Q),
\]
其中 $T'$ 包含模拟开销。由于假设中的优势已经允许 $Q$ 条样本，不需要再乘 $Q$。如果将 IND 优势定义成两个消息实验的输出概率差，对应上界是 $2\delta(T',Q)$。该结论涵盖经典加密接口中的重复和自适应选择明文，前提是不额外公开依赖秘密的材料。面对量子攻击者时，基础判定假设也必须涵盖量子区分器；此陈述不允许叠加态加密查询。

解密先减去依赖秘密的项，再通过舍入去除足够小的误差，除以 $\Delta$，最后模 $t$ 约减。噪声界与安全假设是两个独立要求：足够小而能正确解码，并不自动意味着足够大而能保证安全。

\subsection{RLWE 密码体制}
\label{subsec:lattice-scheme2}

RLWE 密码体制使用与 \tboxlabel{\ref*{subsec:lattice-scheme}} 相同的公式，但将 $\vec{s}, \vec{a}^{\langle i \rangle}, b^{\langle i \rangle}, m^{\langle i \rangle}, e^{\langle i \rangle}$ 替换为多项式 $S, A^{\langle i \rangle}, B^{\langle i \rangle}, M^{\langle i \rangle}, E^{\langle i \rangle}$，如下所示：

\begin{tcolorbox}[title={\textbf{\tboxlabel{\ref*{subsec:lattice-scheme2}} 基于格的 RLWE 密码体制}}]
\begin{itemize}
    \item \textbf{\underline{加密}：}$B^{\langle i \rangle} = S \cdot A^{\langle i \rangle} + \Delta \cdot M^{\langle i \rangle} + E^{\langle i \rangle}$，其中 $B^{\langle i \rangle}$ 和 $A^{\langle i \rangle}$ 对攻击者也是公开已知的，而 $S, M^{\langle i \rangle}, E^{\langle i \rangle}$ 则未知，仅秘密密钥持有者知晓。

    $ $
    
    \item \textbf{\underline{解密}：}$\left\lceil [B^{\langle i\rangle}-S\cdot A^{\langle i\rangle}]_q/\Delta\right\rfloor\bmod t=M^{\langle i\rangle}$ \\$\Big($前提是 $||E^{\langle i \rangle}||_{\infty} < \dfrac{\Delta}{2}$，即 $E^{\langle i \rangle}$ 的每个系数的绝对值都小于 $\dfrac{\Delta}{2}\Big)$

    $ $

    $\lceil  \rfloor_{\Delta}$ 表示对多项式每一项的系数分别进行舍入。

\end{itemize}



\end{tcolorbox}

上述安全归约在相应判定 RLWE 假设下适用于环元素对，也在实际 GLWE 假设下适用于 $k$ 个环掩码。将环乘法展开成系数方程，得到的是具有负循环结构的矩阵，并非独立均匀的 LWE 行。因此可以沿用证明模板，但标量 LWE 困难性定理并不会自动认证某个环或模结构参数族。


\clearpage

\section{LWE 密码体制}
\label{sec:lwe}

\textbf{- 参考资料：}
\href{https://www.zama.ai/post/tfhe-deep-dive-part-1}{TFHE 深入解析：第一部分 -- 密文类型}~\cite{tfhe-1}



$ $




\subsection{参数设置}
设 $[0, t-1]$ 为明文范围，$[0, q-1]$ 为密文范围，其中 $t \ll q$，即 $t$ 远小于 $q$。随机选取长度为 $k$ 的向量 $\vec{s}$ 作为秘密密钥，其中包含 $k$ 个从 $\{-1, 0, 1\}$ 中采样的三元数，记为 $\vec{s} \xleftarrow{\$} \{-1, 0, 1\}^k$。令明文尺度因子为 $\Delta = \dfrac{q}{t}$。

\subsection{加密}
\label{subsec:lwe-enc}

\begin{enumerate}
\item 假设要加密明文 $m \in \mathbb{Z}_t$。\item 随机选取长度为 $k$ 的向量 $\vec{a} \in \mathbb{Z}_q^k$，作为一次性的公开随机掩码，记为 $\vec{a} \xleftarrow{\$} \mathbb{Z}_q^k$。
\item 从高斯分布 $\chi_\sigma$ 中采样一个小的、一次性随机噪声 $e \in \mathbb{Z}_q$，记为 $e \xleftarrow{\chi_\sigma} \mathbb{Z}_q$。
\item 将 $m$ 按 $\Delta$ 缩放，即计算 $\Delta \cdot m$。这将 $m \in \mathbb{Z}_t$ 转换为 $\Delta \cdot m \in \mathbb{Z}_q$。
\item 计算 $b = \vec{a} \cdot \vec{s} + \Delta \cdot m + e \in \mathbb{Z}_q$。
\end{enumerate}

$ $

LWE 加密公式汇总如下：

$ $

\begin{tcolorbox}[title={\textbf{\tboxlabel{\ref*{subsec:lwe-enc}} LWE 加密}}]
\textbf{\underline{初始设置}：}$\Delta = \dfrac{q}{t}$，$\vec{s} \xleftarrow{\$} \{-1, 0, 1\}^k$，其中 $t$ 整除 $q$。

$ $

$ $

\textbf{\underline{加密输入}：}$m \in \mathbb{Z}_t$，$\vec{a} \xleftarrow{\$} \mathbb{Z}_q^k$，$e \xleftarrow{\chi_\sigma} \mathbb{Z}_q$

$ $

%\textbf{Encryption}

\begin{enumerate}
\item 向上缩放 $m \longrightarrow \Delta \cdot m \text{ } \in \mathbb{Z}_q$

\item 计算 $b = \vec{a} \cdot \vec{s} + \Delta  m + e \pmod q$
\item $\textsf{LWE}_{\vec{s},\sigma}(\Delta  m + e) = (\vec{a}, b) \text{ } \in \mathbb{Z}_q^{k + 1}$ 
\end{enumerate}

\end{tcolorbox}

\subsection{解密}
\label{subsec:lwe-dec} 

\begin{enumerate}
\item 给定密文 $(\vec{a}, b)$，其中 $b = \vec{a} \cdot \vec{s} + \Delta \cdot m + e \in \mathbb{Z}_q$，计算 $b - \vec{a} \cdot \vec{s}$，得到与 $\Delta \cdot m + e \in \mathbb{Z}_q$ 相同的值。
\item 将 $\Delta \cdot m + e \in \mathbb{Z}_q$ 舍入到最接近的 $\Delta$ 的整数倍，即按基数 $\Delta$ 舍入，记为 $\lceil \Delta \cdot m + e \rfloor_{\Delta}$。只要 $e$ 足够小并满足 $e < \dfrac{\Delta}{2}$，舍入就能成功消去 $e$，得到 $\Delta m$。如果 $e \geq \dfrac{\Delta}{2}$，则噪声 $e$ 的部分高位会与明文 $m$ 重叠，无法被消去，从而破坏明文 $m$ 的部分低位。
\item 计算 $\dfrac{\Delta m} {\Delta}$，这等价于将 $\lceil \Delta \cdot m + e \rfloor_{\Delta}$ 向右移动 $\text{log}_2 \Delta$ 位。这里假设 $\Delta$ 为 2 的幂；若 $\Delta$ 不是 2 的幂，则将 $m$ 按 $\Delta$ 向上或向下缩放，分别意味着将该数乘以或除以 $\Delta$。
\end{enumerate}

$ $

LWE 解密公式汇总如下：

$ $

\begin{tcolorbox}[title={\textbf{\tboxlabel{\ref*{subsec:lwe-dec}} LWE 解密}}]

\textbf{\underline{解密输入}：}$\textsf{ct} = (\vec{a}, b) \text{ } \in \mathbb{Z}_q^{k+1}$

$ $

%\textbf{Decryption}
\begin{enumerate}
\item $\textsf{LWE}^{-1}_{S,\sigma}(\textsf{ct}) = b - \vec{a}\cdot \vec{s} = \Delta  m + e \pmod q$

\item 向下缩放 $\Bigg\lceil\dfrac{ \Delta  m + e } {\Delta}\Bigg\rfloor \bmod t = m \text{ } \in \mathbb{Z}_t$
\end{enumerate}

为了正确解密，噪声 $e$ 应满足 $e < \dfrac{\Delta}{2}$。

\end{tcolorbox}

解密时，秘密密钥持有者可以利用秘密密钥 $\vec{s}$ 直接计算 $\vec{a} \cdot \vec{s}$，因此能够从 $b$ 中减去 $\vec{a} \cdot \vec{s}$。

将明文 $m$ 按 $\Delta$ 缩放有两个目的：(i) 加密时，将 $m$ 左移 $\text{log}_2\Delta$ 位，与低位中的噪声 $e$ 分离；而噪声 $e$ 是使攻击者难以猜测 $m$ 或 $\vec{s}$ 的必要成分；(ii) 解密时，通过右移 $\text{log}_2\Delta$ 位消去低位噪声 $e$，同时保持高位的 $m$ 不受损。将明文 $m$ 右移（即缩放）$\text{log}_2\Delta$ 位，再加入噪声 $e$ 的过程如 \autoref{fig:scaling} 所示。

\subsubsection{$t$ 不整除 $q$ 的情形}
\label{subsubsec:lwe-noise-bound}

在总结~\ref*{subsec:lwe-enc}（\autoref{subsec:lwe-enc}）中，假设 $t$ 整除 $q$。此时，明文 $m$ 的大小没有额外的上下界限制：其数值可以反复越过模 $t$ 的边界，而解密仍然正确。这是因为 $m$
在模 $t$ 下的回绕，可以由解密时缩放后明文的模 $q$ 约减对应处理。

另一方面，假设 $t$ 不整除 $q$。此时，将尺度因子设为 $\Delta = \left\lfloor\dfrac{q}{t}\right\rfloor$。
只要 $q \gg t$，即使 $m$ 很大并越过模 $t$ 的边界，仍然可以正确解密。下面说明其原因。

在本小节的分析中，选择 $t$ 为奇数（素数）。出于计算效率的考虑，这是全同态加密中常见的做法（见 \autoref{subsubsec:poly-vector-transformation-modulus}）。
假设采用中心化剩余表示，明文域为 $\left[-\dfrac{t-1}{2}, \dfrac{t-1}{2}\right]$，密文域为 $\left[-\dfrac{q}{2}, \dfrac{q}{2} - 1\right]$。
将明文 $m \bmod t$ 表示为 $m = m' + vt$，其中 $m' \in \mathbb{Z}_t$，整数 $v$ 表示 $m$ 越过模 $t$ 边界的次数。将明文尺度因子设为 $\Delta = \left\lfloor\dfrac{q}{t}\right\rfloor$，则按 $\Delta$ 缩放并加入噪声后的明文可以表示为：

$\left\lfloor\dfrac{q}{t}\right\rfloor\cdot m + e $

$= \left\lfloor\dfrac{q}{t}\right\rfloor\cdot m' + \left\lfloor\dfrac{q}{t}\right\rfloor\cdot vt + e$ \textcolor{red}{ $\rhd$ 使用 $m = m' + vt$}

$= \left\lfloor\dfrac{q}{t}\right\rfloor\cdot m' + \dfrac{q}{t}\cdot vt - \left(\dfrac{q}{t} - \left\lfloor\dfrac{q}{t}\right\rfloor\right)\cdot vt + e$  

$= \left\lfloor\dfrac{q}{t}\right\rfloor\cdot m' + qv - \left(\dfrac{q}{t} - \left\lfloor\dfrac{q}{t}\right\rfloor\right)\cdot vt + e$

$= \left\lfloor\dfrac{q}{t}\right\rfloor\cdot m' + qv - \epsilon \cdot vt + e$ \textcolor{red}{ $\rhd$ 其中 $\epsilon = \dfrac{q}{t} - \left\lfloor\dfrac{q}{t}\right\rfloor$，是 $[0, 1)$ 之间的小数}

$ $

回顾 LWE 的解密关系：$\dfrac{b - \vec{a}\cdot \vec{s} \bmod q}{\Delta} \bmod t = \dfrac{\Delta m + e \bmod q}{\Delta} \bmod t$。根据上述表达式，可以通过以下计算解密消息：



$\left\lceil\dfrac{1}{\Delta} \cdot \left(b - \vec{a}\cdot\vec{s} \bmod q\right)\right\rfloor \bmod t$

$= \left\lceil\dfrac{1}{\Delta} \cdot \left(\Delta m + e \bmod q\right)\right\rfloor \bmod t$

$= \left\lceil\dfrac{1}{\left\lfloor\frac{q}{t}\right\rfloor} \cdot \left(\left\lfloor\dfrac{q}{t}\right\rfloor \cdot (m' + vt) + e \bmod q\right)\right\rfloor \bmod t$

$= \left\lceil\dfrac{1}{\left\lfloor\frac{q}{t}\right\rfloor} \cdot \left(\left\lfloor\dfrac{q}{t}\right\rfloor \cdot m' + \left\lfloor\dfrac{q}{t}\right\rfloor\cdot vt + e \bmod q\right)\right\rfloor \bmod t$

$= \left\lceil\dfrac{1}{\left\lfloor\frac{q}{t}\right\rfloor} \cdot \left(\left\lfloor\dfrac{q}{t}\right\rfloor \cdot m' + \left(\dfrac{q}{t} - \epsilon \right)\cdot vt + e \bmod q\right)\right\rfloor \bmod t$

$= \left\lceil\dfrac{1}{\left\lfloor\frac{q}{t}\right\rfloor} \cdot \left(\left\lfloor\dfrac{q}{t}\right\rfloor \cdot m' + vq - \epsilon vt + e \bmod q\right)\right\rfloor \bmod t$


$= \left\lceil\dfrac{1}{\left\lfloor\frac{q}{t}\right\rfloor} \cdot \left(\left\lfloor\dfrac{q}{t}\right\rfloor\cdot m'  - \epsilon vt + e \bmod q \right)\right\rfloor \bmod t$ 

$ $

为了正确解密，希望能够省去内层的模 $q$ 约减。也就是说，在中心化剩余表示 $\left[-\dfrac{q}{2}, \dfrac{q}{2}-1\right]$ 下，希望满足 $-\dfrac{q}{2} \leq \left\lfloor\dfrac{q}{t}\right\rfloor\cdot m'  - \epsilon vt + e < \dfrac{q}{2}$，或者采用更严格的条件 $\left|\left\lfloor\dfrac{q}{t}\right\rfloor\cdot m'  - \epsilon vt + e\right| < \dfrac{q}{2}$。

$ $

为保证这些条件成立，引入一个特殊假设：$|-\epsilon vt + e| < \dfrac{\Delta}{2}$。将这一假设用于上述表达式，可推导出：

$\left|\left\lfloor\dfrac{q}{t}\right\rfloor\cdot m'  - \epsilon vt + e\right| $

$ \leq \left|\left\lfloor\dfrac{q}{t}\right\rfloor\cdot m'\right| + \left|- \epsilon vt + e\right|$

$ \leq \left|\left\lfloor\dfrac{q}{t}\right\rfloor\cdot \dfrac{t-1}{2}\right| + \left|- \epsilon vt + e\right|$ \textcolor{red}{ $\rhd$ 假设 $t$ 为奇素数，这是全同态加密中的常见做法}

$ \leq \left|\dfrac{q}{t}\cdot \dfrac{t-1}{2}\right| + \left|- \epsilon vt + e\right|$

$ = \dfrac{q}{2} - \dfrac{q}{2t} + \left|- \epsilon vt + e\right|$

$ < \dfrac{q}{2} - \dfrac{q}{2t} + \dfrac{\Delta}{2}$ \textcolor{red}{ $\rhd$ 使用特殊假设 $|-\epsilon vt + e| < \dfrac{\Delta}{2}$}

$ = \dfrac{q}{2} + \dfrac{1}{2}\cdot \left(\Delta - \dfrac{q}{t}\right)$

$< \dfrac{q}{2}$ \textcolor{red}{ $\rhd$ 因为 $\Delta < \dfrac{q}{t}$}


$ $

因此，若假设 $|-\epsilon vt + e| < \dfrac{\Delta}{2}$，便可将 $\left\lfloor\dfrac{q}{t}\right\rfloor\cdot m'  - \epsilon vt + e \bmod q$ 简化为 $\left\lfloor\dfrac{q}{t}\right\rfloor\cdot m'  - \epsilon vt + e$。继续推导：


$\left\lceil\dfrac{1}{\left\lfloor\frac{q}{t}\right\rfloor} \cdot \left(\left\lfloor\dfrac{q}{t}\right\rfloor\cdot m'  - \epsilon vt + e \bmod q \right)\right\rfloor \bmod t$ 

$= \left\lceil\dfrac{1}{\left\lfloor\frac{q}{t}\right\rfloor} \cdot \left(\left\lfloor\dfrac{q}{t}\right\rfloor\cdot m'  - \epsilon vt + e \right)\right\rfloor \bmod t$ 



$= \left\lceil m' + \dfrac{-\epsilon  vt + e  }{\lfloor\frac{q}{t}\rfloor} \right\rfloor \bmod t$ 

$= m' + \left\lceil \dfrac{-\epsilon vt + e }{\lfloor\frac{q}{t}\rfloor} \right\rfloor \bmod t$ \textcolor{red}{ $\rhd$ 因为 $\lceil m' \rfloor = m'$}

$= m' \bmod t$ \textcolor{red}{ $\rhd$ 使用特殊假设 $|-\epsilon vt + e| < \dfrac{\Delta}{2} = \dfrac{\lfloor\frac{q}{t}\rfloor}{2}$}

因此，当明文尺度因子设为 $\Delta=\left\lfloor\dfrac{q}{t}\right\rfloor$，且 $t$ 为奇数（素数）时，只要满足误差界条件 $|-\epsilon vt + e| < \dfrac{\Delta}{2} = \dfrac{\lfloor\frac{q}{t}\rfloor}{2}$，就能正确解密。以下情况会破坏这一条件并导致解密失败：(1) 噪声 $e$ 相对于 $q$ 过大；(2) 明文模数 $t$ 相对于 $q$ 过大；(3) 明文越过模 $t$ 边界的次数过多，即 $v$ 过大。满足这些误差界条件的一种通用方法，是将密文模数 $q$ 设得足够大。换言之，如果 $q \gg t$ 且 $q \gg e$，上述误差界就成立。

可以将总结~\ref*{subsec:lwe-enc}（见 \autoref{subsec:lwe-enc}）中的明文尺度因子公式推广为 $\left\lfloor\dfrac{q}{t}\right\rfloor$，其中 $t$ 为奇数（素数）。


\begin{tcolorbox}[title={\textbf{\tboxlabel{\ref*{subsubsec:lwe-noise-bound}} 奇数明文模数 $\bm{t}$ 下的噪声预算}}]


给定明文尺度因子 $\Delta=\left\lfloor\dfrac{q}{t}\right\rfloor$，且 $t$ 为奇数（素数），只要满足以下误差界条件，LWE 解密就能正确进行：

$|-\epsilon vt + e| < \dfrac{\Delta}{2}$

$ $

其中，$\epsilon = \dfrac{q}{t} - \left\lfloor\dfrac{q}{t}\right\rfloor$ 是 $[0, 1)$ 之间的小数，$v$ 表示明文 $m$ 越过模 $t$ 边界的次数。

\end{tcolorbox}



\clearpage

\section{RLWE 密码体制}
\label{sec:rlwe}
RLWE 密码体制的密文是二元组 $(A, B)$，其中 $B = S \cdot A + \Delta \cdot M + E$。公开随机掩码 $A$ 与秘密密钥 $S$ 都是 $(n-1)$ 次多项式；消息 $M$ 和噪声 $E$ 都是 $(n-1)$ 次多项式。与 LWE 类似，每个密文都使用新的公开随机掩码 $A$，而所有密文使用同一个秘密密钥 $S$。为简化记号，本节将每个密文实例记为 $(A, B)$，而不使用 $(A^{\langle i \rangle}, b^{\langle i \rangle})$。

在 RLWE 中，所有多项式运算都在多项式环 $\mathbb{Z}_q[x] / (x^n + 1)$ 中进行，其中 $x^n + 1$ 为分圆多项式，对于某个整数 $f$ 满足 $n = 2^f$，且多项式系数属于 $\mathbb{Z}_q$。因此，RLWE 中所有多项式的系数范围为 $\mathbb{Z}_q$，最高次数为 $n -1$。为简化记号，记 $\mathcal{R}_{\langle n,q \rangle} = \mathbb{Z}_q[x] / (x^n + 1)$。

\subsection{参数设置}

设 $t$ 为明文空间大小，$q$ 为密文空间大小，其中 $t < q$，即 $t$ 远小于 $q$，并且 $t | q$，即 $t$ 整除 $q$。随机选取 $(n-1)$ 次多项式 $S \in \mathcal{R}_{\langle n, q \rangle}$ 作为秘密密钥，其系数取自 $\{-1, 0, 1\}$。令明文尺度因子为 $\Delta = \left\lfloor\dfrac{q}{t}\right\rfloor$。

RLWE 的设置参数与 LWE 类似。区别在于，$S$ 不再是从 $\{-1, 0, 1\}$ 中采样的长度为 $k$ 的向量，而是包含 $n$ 个秘密系数的 $(n-1)$ 次多项式。每个系数都是从 $\{-1, 0, 1\}$ 中随机采样的三元数，记为 $S \xleftarrow{\$} \mathcal{R}_{\langle n, \textit{tern} \rangle}$。


\subsection{加密}
\label{subsec:rlwe-enc}

多项式代表元的次数至多为 $n-1$，允许高次系数为零。噪声系数是采样后模 $q$ 约减的整数，采用 \autoref{subsec:lwe-dec} 中的约定。正确性不等式约束这些提升后的系数，而不是对无界高斯分布给出确定性保证。

\begin{enumerate}
\item 假设有一个 $(n-1)$ 次多项式 $M \in \mathcal{R}_{\langle n, t \rangle}$，其系数表示要加密的明文数值。
\item 随机选取 $(n-1)$ 次多项式 $A \in \mathcal{R}_{\langle n,q \rangle}$ 作为一次性的公开随机掩码，记为 $A \xleftarrow{\$} \mathcal{R}_{\langle n, q \rangle}$。
\item 随机选取小系数多项式 $E \in \mathcal{R}_{\langle n,q \rangle}$ 作为一次性噪声，其 $n$ 个系数都是 $\mathbb{Z}_q$ 中从高斯分布 $\chi_\sigma$ 随机采样的小数值，记为 $E \xleftarrow{\chi_\sigma} \mathcal{R}_{\langle n, q \rangle}$。
\item 将 $M$ 按 $\Delta$ 缩放，即计算 $\Delta \cdot M$。这将 $M \in \mathcal{R}_{\langle n,t \rangle}$ 转换为 $\Delta \cdot M \in \mathcal{R}_{\langle n,q \rangle}$。
\item 计算 $B = A \cdot S + \Delta \cdot M + E \bmod \mathcal{R}_{\langle n,q \rangle}$，即按 $n$ 约减多项式次数，并将系数模 $q$ 约减。
\item 最终密文为 $(A, B)$。
\end{enumerate}

$ $

RLWE 加密公式汇总如下：

$ $
\begin{tcolorbox}[title={\textbf{\tboxlabel{\ref*{subsec:rlwe-enc}} RLWE 加密}}]
\textbf{\underline{初始设置}：}$\Delta = \left\lfloor\dfrac{q}{t}\right\rfloor$，$S \xleftarrow{\$} \mathcal{R}_{\langle n, \textit{tern} \rangle}$

$ $

$ $

\textbf{\underline{加密输入}：}$M \in \mathcal{R}_{\langle n, t \rangle}$，$A \xleftarrow{\$} \mathcal{R}_{\langle n, q \rangle}$，$E \xleftarrow{\chi_\sigma} \mathcal{R}_{\langle n, q \rangle}$

$ $

\begin{enumerate}
\item 向上缩放 $M \longrightarrow \Delta  M \text{ } \in \mathcal{R}_{\langle n, q \rangle}$
\item 计算 $B = A \cdot S + \Delta  M + E  \text{ } \bmod \mathcal{R}_{\langle n,q \rangle}$
\item $\textsf{RLWE}_{S,\sigma}(\Delta  M + E) = (A, B) \text{ } \in \mathcal{R}_{\langle n,q \rangle}^2$
\end{enumerate}

\end{tcolorbox}

\subsection{解密}
\label{subsec:rlwe-dec}

\begin{enumerate}
\item 给定密文 $(A, B)$，其中 $B = A \cdot S + \Delta \cdot M + E \in \mathcal{R}_{\langle n,q \rangle}$，计算 $B - A \cdot S = \Delta \cdot M + E$。
\item 将多项式 $\Delta \cdot M + E \in \mathcal{R}_{\langle n,q \rangle}$ 的每个系数舍入到最接近的 $\Delta$ 的整数倍，即按基数 $\Delta$ 舍入，记为 $\lceil \Delta \cdot M + E \rfloor_{\Delta}$。这一操作可以成功消去 $E$，得到 $\Delta \cdot M$。需要注意的是，噪声多项式 $E$ 的每个系数 $e_i$ 都必须足够小，满足 $|e_i| < \dfrac{\Delta}{2}$，才能在舍入过程中被消去。否则，$e_i$ 的部分高位会与明文 $m_i$ 系数的低位重叠并将其破坏，而且无法通过舍入消去。
\item 计算 $\dfrac{\Delta \cdot M}{\Delta}$，等价于将 $\Delta \cdot M$ 的每个多项式系数按 $\Delta$ 向下缩放；若 $\Delta$ 是 2 的幂，也就等价于将每个系数右移 $\text{log}_2 \Delta$ 位。
\end{enumerate}

$ $

RLWE 解密公式汇总如下：


\begin{tcolorbox}[title={\textbf{\tboxlabel{\ref*{subsec:rlwe-dec}} RLWE 解密}}]

\textbf{\underline{解密输入}：}$\textsf{ct} = (A, B) \text{ } \in \mathcal{R}_{\langle n, q \rangle}^{2}$

$ $

\begin{enumerate}

\item $\textsf{RLWE}^{-1}_{S,\sigma}(\textsf{ct}) = B - A \cdot S = \Delta  M + E \text{ } \in \mathcal{R}_{\langle n,q \rangle}$ 
\item 向下缩放 $\Bigg\lceil\dfrac{ \Delta M + E}{\Delta}\Bigg\rfloor \bmod t = M \text{ } \in \mathcal{R}_{\langle n,t \rangle}$

\end{enumerate}

当 $t\mid q$ 时，要求 $\|E\|_\infty<\Delta/2$。若 $t\nmid q$，应逐系数使用 \autoref{subsubsec:lwe-noise-bound} 的有效误差界：对于未约减系数 $m_i=m_i'+v_it$，其中 $m_i'$ 为中心化代表元，在该小节的奇数 $t$ 假设下要求 $|e_i-(q\bmod t)v_i|<\Delta/2$。仅有 $q\gg t$ 并不充分。

\end{tcolorbox}

\clearpage

\section{GLWE 密码体制}
\label{sec:glwe}
GLWE 密码体制是同时涵盖 LWE 与 RLWE 的一般化形式。GLWE 密文是元组 $(\{A_i\}_{i=0}^{k-1}, B)$，其中 $B = \sum\limits_{i=0}^{k-1}{(A_i \cdot S_i)} + \Delta \cdot M + E$。公开量 $\{A_i\}_{i=0}^{k-1}$ 和秘密密钥 $\{S_i\}_{i=0}^{k-1}$ 都是由 $k$ 个 $(n-1)$ 次多项式组成的列表；消息 $M$ 与噪声 $E$ 各为一个 $(n-1)$ 次多项式。与 LWE 和 RLWE 类似，每个密文都使用新的公开随机掩码 $A$，而所有密文使用相同的秘密密钥 $S$。为简化记号，本节将每个密文实例记为 $(\{A_i\}_{i=0}^{k-1}, B)$，而不使用 $(\{A_i\}_{i=0}^{k-1, \langle j \rangle}, B^{\langle j \rangle})$。




\subsection{参数设置}
设 $t$ 为明文空间大小，$q$ 为密文空间大小，其中 $t < q$，即 $t$ 远小于 $q$，且 $t | q$，即 $t$ 整除 $q$。随机选取由 $k$ 个 $(n-1)$ 次多项式构成的列表作为秘密密钥，其中每个系数都是从 $\{-1, 0, 1\}$ 中随机选取的三元数，即 $\{S_i\}_{i=0}^{k-1} \xleftarrow{\$} \mathcal{R}_{\langle n, \textit{tern} \rangle}^k$。令明文尺度因子为 $\Delta = \left\lfloor\dfrac{q}{t}\right\rfloor$。

GLWE 的设置参数与 RLWE 类似。区别在于，$S$ 不再是包含 $n$ 个秘密系数的单个 $(n-1)$ 次多项式，而是由 $k$ 个这样的 $(n-1)$ 次多项式组成的列表，总共包含 $n \cdot k$ 个秘密系数。


\subsection{加密}
\label{subsec:glwe-enc}

多项式代表元的次数至多为 $n-1$，允许高次系数为零。噪声系数是采样后模 $q$ 约减的整数，采用 \autoref{subsec:lwe-dec} 中的约定。正确性不等式约束这些提升后的系数，而不是对无界高斯分布给出确定性保证。

假设有一个 $(n-1)$ 次多项式 $M \in \mathcal{R}_{\langle n, t \rangle}$，其系数表示要加密的明文数值。


\begin{enumerate}
\item 随机选取由 $k$ 个 $(n-1)$ 次多项式组成的列表 $\{A_i\}_{i=0}^{k-1} \xleftarrow{\$} \mathcal{R}_{\langle n,q \rangle}^{k}$，作为一次性的公开随机量。
\item 随机选取小系数多项式 $E \xleftarrow{\chi_\sigma} \mathcal{R}_{\langle n,q \rangle}$ 作为一次性噪声，其 $n$ 个系数都是 $\mathbb{Z}_q$ 中从高斯分布 $\chi_\sigma$ 随机采样的小数值。
\item 将 $M$ 按 $\Delta$ 缩放，即计算 $\Delta \cdot M$。这将 $M \in \mathcal{R}_{\langle n, t \rangle}$ 转换为 $\Delta \cdot M \in \mathcal{R}_{\langle n,q \rangle}$。
\item 计算 $B = \sum\limits_{i=0}^{k-1}{(A_i \cdot S_i)} + \Delta \cdot M + E \in \mathcal{R}_{\langle n,q \rangle}$。
\end{enumerate}

$ $

GLWE 加密公式汇总如下：

$ $

\begin{tcolorbox}[title={\textbf{\tboxlabel{\ref*{subsec:glwe-enc}} GLWE 加密}}]

\textbf{\underline{初始设置}：}$\Delta = \left\lfloor\dfrac{q}{t}\right\rfloor$，$\{S_i\}_{i=0}^{k-1} \xleftarrow{\$} \mathcal{R}_{\langle n, \textit{tern} \rangle}^k$

$ $

$ $

\textbf{\underline{加密输入}：}$M \in \mathcal{R}_{\langle n, t \rangle}$，$\{A_i\}_{i=0}^{k-1} \xleftarrow{\$} \mathcal{R}_{\langle n,q \rangle}^{k}$，$E \xleftarrow{\chi_\sigma} \mathcal{R}_{\langle n,q \rangle}$

$ $

\begin{enumerate}

\item 向上缩放 $M \longrightarrow \Delta M \text { } \in \mathcal{R}_{\langle n, q\rangle}$

\item 计算 $B = \sum\limits_{i=0}^{k-1}{(A_i \cdot S_i)} + \Delta  M + E \text{ } \in \mathcal{R}_{\langle n,q \rangle}$

\item $\textsf{GLWE}_{S,\sigma}(\Delta M + E) = (\{A_i\}_{i=0}^{k-1}, B) \text{ } \in \mathcal{R}_{\langle n,q \rangle}^{k + 1}$ 



\end{enumerate}

\end{tcolorbox}

\subsection{解密}
\label{subsec:glwe-dec}

\begin{enumerate}
\item 给定密文 $(\{A_i\}_{i=0}^{k-1}, B)$，其中 $B = \sum\limits_{i=0}^{k-1}{(A_i \cdot S_i)} + \Delta \cdot M + E \in \mathcal{R}_{\langle n,q \rangle}$，计算 $B - \sum\limits_{i=0}^{k-1}{(A_i \cdot S_i)} = \Delta \cdot M + E$。
\item 将多项式 $\Delta \cdot M + E \in \mathcal{R}_{\langle n,q \rangle}$ 的每个系数舍入到最接近的 $\Delta$ 的整数倍，即按基数 $\Delta$ 舍入，记为 $\lceil \Delta \cdot M + E \rfloor_{\Delta}$。这一操作可以成功消去 $E$，得到 $\Delta \cdot M$。需要注意的是，$E$ 的每个系数 $e_i$ 都必须满足 $|e_i| < \dfrac{\Delta}{2}$，才能在舍入时被消去。否则，$e_i$ 的部分高位会与明文 $m_i$ 系数的低位重叠，无法在解密时消去，从而破坏明文 $m_i$。
\item 计算 $\dfrac{\Delta \cdot M} {\Delta}$，等价于将 $\Delta \cdot M$ 的每个多项式系数按 $\Delta$ 向下缩放。
\end{enumerate}

$ $

GLWE 解密公式汇总如下：

$ $

\begin{tcolorbox}[title={\textbf{\tboxlabel{\ref*{subsec:glwe-dec}} GLWE 解密}}]
\textbf{\underline{解密输入}：}$\textsf{ct} = (\{A_i\}_{i=0}^{k-1}, B) \text{ } \in \mathcal{R}_{\langle n,q \rangle}^{k + 1}$

\begin{enumerate}
\item $\textsf{GLWE}^{-1}_{S,\sigma}(\textsf{ct}) = B - \sum\limits_{i=0}^{k-1}{(A_i \cdot S_i)} = \Delta  M + E \text{ } \in \mathcal{R}_{\langle n,q \rangle}$

\item 向下缩放
$\Bigl\lceil \dfrac{ \Delta  M + E }{\Delta}\Bigr\rfloor \bmod t = M \text{ } \in \mathcal{R}_{\langle n, t \rangle}$

\end{enumerate}

为了正确解密，噪声多项式 $E$ 的每个系数 $e_i$ 都应满足 $|e_i| < \dfrac{\Delta}{2}$。

\end{tcolorbox}


\subsubsection{讨论}

下面的大小比较仅描述表示方式与暴力搜索空间。具体安全性共同依赖于 $n,k,q$、误差和秘密分布及样本数量；单独增大某个维数并不构成安全归约。
\begin{enumerate}
\item \textbf{LWE} 是多项式环次数参数为 $n = 1$ 时的 GLWE 特例。此时，$\{A_i\}_{i=0}^{k-1}, \{S_i\}_{i=0}^{k-1}, E$ 和 $M$ 中的所有多项式都是零次常数多项式，而 $A_i$ 和 $S_i$ 分别含有 $k$ 个这样的常数，因此各自构成一个向量。
\item \textbf{RLWE} 是 $k = 1$ 时的 GLWE 特例。此时，秘密密钥 $S$ 是单个多项式 $S_0$，每次加密也只使用单个多项式 $A_0$ 作为公开量。
\item \textbf{$\bm{n}$ 的大小：}多项式 $S_i = s_{i,0} + s_{i,1}X + \cdots + s_{i, n-1}X^{n-1} $ 有 $n$ 个系数。若各系数独立均匀取三元值，则完整秘密有 $3^{nk}$ 种可能取值。这是穷举候选数量，并非最佳攻击的下界，也不是搜索假设（\autoref{subsec:lattice-overview}）的证明。更高次数的明文多项式可以表示更多系数，但实际吞吐量还取决于运算成本。
%\item \textbf{Size of $\bm{q}$:} Increasing $q$ increases the ciphertext molulus space, making it hard to attack the cryptosystem. 
%Also, a large modulo space creates an enough separation between the plaintext bits and the noise bits in the ciphertext, which reduces the growth rate of the noise bits across ciphertext multiplications (will be discussed in \autoref{subsubsec:glwe-mult-plain-discussion}) and thus we can do a more number of ciphertext multiplications before the noise bits overflow to the plaintext bit area. 
\item \textbf{$\bm{k}$ 的大小：}增大 $k$ 会增加秘密多项式 $(S_0, S_1, \ldots, S_{k-1})$ 和公开掩码多项式 $(A_0, A_1,\ldots, A_{k-1})$。每个密文仍只有一个消息多项式 $M$ 和一个误差多项式 $E$。比较两个 $n,k$ 选择时，需要考察其实际环或模结构假设与分布，而不能只数秘密系数。
\item \textbf{减小密文大小：}种子压缩密文可以保存新鲜独立的公开种子 $d$ 与 $B$。明确指定的公开展开函数 $G$ 生成全部 $kn$ 个掩码系数，并使用参数/索引域分离及无偏系数映射。其输入绝不是秘密。由于种子公开，普通隐藏种子的 PRG 安全性并不能说明 $(d,G(d))$ 可替换为带独立均匀 $A$ 的 $(d,A)$：检验 $A=G(d)$ 就能区分两者。对固定展开函数，所需联合假设比较的是 $(d,\langle G(d),S\rangle+E)$ 与均匀 $V$ 对应的 $(d,V)$，并须使用实际秘密/误差分布及样本数量。
\end{enumerate}

\noindent\textbf{明确的种子模型：}若展开是理想随机预言机 $H:\{0,1\}^h\to\mathcal R_{\langle n,q\rangle}^k$，经典攻击者作 $Q$ 次加密查询和 $Q_H$ 次预言机查询时，其猜测偏差满足
\[
\epsilon\le\delta(T',Q)+\min\left\{1,\frac{QQ_H+Q(Q-1)/2}{2^h}\right\}.
\]
这里 $\delta$ 是实际分布的判定 GLWE 优势，采用 \autoref{subsec:lattice-scheme} 的约定。模拟密文时独立选择种子，把挑战样本的均匀掩码编程到此前未赋值的 $H(d)$，再向身体加入选定消息。仅当种子此前已被查询，或重复了更早的加密种子时，模拟才失败。第 $j$ 次加密时至多有 $Q_H+j-1$ 个点已经占用，因此得到所列联合界。在该事件之外，真实样本产生真实加密，均匀样本身体则隐藏消息；失败时中止并作独立公平猜测。这证明了上述经典随机预言机结论，既没有不加假设地把具体 XOF 替换为理想预言机，也没有为量子预言机查询证明同一个界。关于公开 XOF 展开与秘密伪随机生成的区别，参见 \href{https://pq-crystals.org/kyber/data/kyber-specification-round3-20210804.pdf}{Kyber v3.02 说明第 4.3.1 与 5.4 节}。重复掩码还会暴露两个消息相位之差加上误差差值，因此种子的新鲜性至关重要。

\subsection{GLWE 的另一种形式}
\label{subsec:glwe-alternative}

下面给出 \tboxlabel{\ref*{subsec:glwe-enc}} 的另一种形式，在加密与解密公式中翻转每个 $A_iS_i$ 的符号：

\begin{tcolorbox}[title={\textbf{\tboxlabel{\ref*{subsec:glwe-alternative}} GLWE 密码体制的另一种形式}}]

\textbf{\underline{初始设置}：}$\Delta = \left\lfloor\dfrac{q}{t}\right\rfloor$，$\{S_i\}_{i=0}^{k-1} \xleftarrow{\$} \mathcal{R}_{\langle n, \textit{tern} \rangle}^k$

$ $

\par\noindent\rule{\textwidth}{0.4pt}


\textbf{\underline{加密输入}：}$M \in \mathcal{R}_{\langle n, t \rangle}$，$\{A_i\}_{i=0}^{k-1} \xleftarrow{\$} \mathcal{R}_{\langle n,q \rangle}^{k}$，$E \xleftarrow{\chi_\sigma} \mathcal{R}_{\langle n,q \rangle}$

$ $

\begin{enumerate}
\item 向上缩放 $M \longrightarrow \Delta M \text { } \in \mathcal{R}_{\langle n, q\rangle}$

\item 计算 $B = \textcolor{red}{-\sum\limits_{i=0}^{k-1}{(A_i \cdot S_i)}} + \Delta  M + E \text{ } \in \mathcal{R}_{\langle n,q \rangle}$

\item $\textsf{GLWE}_{S,\sigma}(\Delta M + E) = (\{A_i\}_{i=0}^{k-1}, B) \text{ } \in \mathcal{R}_{\langle n,q \rangle}^{k + 1}$ 

\end{enumerate}

\par\noindent\rule{\textwidth}{0.4pt}


\textbf{\underline{解密输入}：}$\textsf{ct} = (\{A_i\}_{i=0}^{k-1}, B) \text{ } \in \mathcal{R}_{\langle n,q \rangle}^{k + 1}$

\begin{enumerate}
\item $\textsf{GLWE}^{-1}_{S,\sigma}(\textsf{ct}) = B \text{ } \textcolor{red}{ + \sum\limits_{i=0}^{k-1}{(A_i \cdot S_i)}} = \Delta  M + E \text{ } \in \mathcal{R}_{\langle n,q \rangle}$

\item 向下缩放
$\Bigl\lceil \dfrac{ \Delta  M + E }{\Delta}\Bigr\rfloor \bmod t = M  \text{ } \in \mathcal{R}_{\langle n, t \rangle}$

\end{enumerate}

为了正确解密，噪声多项式 $E$ 的每个系数 $e_i$ 都应满足 $|e_i| < \dfrac{\Delta}{2}$。

\end{tcolorbox}

即使 $A_iS_i$ 项的符号发生翻转，解密阶段也会通过加上相反符号的对应项将其抵消，因此仍保持与原形式相同的解密正确性。





\subsection{公钥加密}
\label{subsec:glwe-public-key-enc}

\textbf{适用范围：}本小节限定 $k=1$，即标准 RLWE 公钥结构，并要求足够大的二幂环维数 $n$，以及能够论证所需判定 RLWE 假设的参数和分布族。下面保留只有一个分量的向量记号。这不是适用于任意 GLWE 维数或 $n=1$ 的 LWE 特例的通用公钥构造；这些推广需要合适的公开矩阵构造及各自的安全论证。与前面的设置相同，展示的解密界假设 $t\mid q$。


\autoref{subsec:glwe-enc} 中的加密方案假设，每个明文都由秘密密钥持有者加密。本节介绍公钥加密方案：生成与秘密密钥相对应的公钥后，任何知道公钥的人都可以加密明文，而只有秘密密钥持有者能够解密。其基本思路是在初始化阶段预先计算部分加密所需的分量，并作为公钥公开。在实际加密阶段，将公钥乘以额外的随机量 $U$，再加入额外噪声 $E_1, \vec{E}_2$，为每个密文产生不可预测的随机性。具体方案如下：

\begin{tcolorbox}[title={\textbf{\tboxlabel{\ref*{subsec:glwe-public-key-enc}} RLWE 公钥加密（$k=1$）}}]

\textbf{\underline{初始设置}：}
\begin{itemize}
\item 尺度因子 $\Delta = \left\lfloor\dfrac{q}{t}\right\rfloor$
\item 秘密密钥 $\vec{S} = \{S_i\}_{i=0}^{k-1} \xleftarrow{\$} \mathcal{R}_{\langle n, \textit{tern} \rangle}^k$
\item 按如下方式生成公钥对 $(\mathit{PK}_1, \vv{\mathit{PK}}_2) \in \mathcal{R}_{\langle n, q \rangle}^{k+1}$：

$\vec{A} = \{A_i\}_{i=0}^{k-1} \xleftarrow{\$} \mathcal{R}_{\langle n, q \rangle}^k$, \text{ } $E \xleftarrow{\chi_{\sigma}} \mathcal{R}_{\langle n, q \rangle}$

$\mathit{PK}_1 = \vec{A} \cdot \vec{S} + E \in \mathcal{R}_{\langle n, q \rangle}$

$\vv{\mathit{PK}}_2 = \vec{A} \in \mathcal{R}_{\langle n, q \rangle}^{k}$

\end{itemize}

\par\noindent\rule{\textwidth}{0.4pt}

\textbf{\underline{加密输入}：}$M \in \mathcal{R}_{\langle n, t \rangle}$，\text{ }$U \xleftarrow{\$} \mathcal{R}_{\langle n,\textit{tern} \rangle}, \text{ } E_1 \xleftarrow{\chi_{\sigma}} \mathcal{R}_{\langle n,q \rangle}, \text{ } \vec{E}_2 \xleftarrow{\chi_{\sigma}} \mathcal{R}_{\langle n,q \rangle}^k$

$ $

\begin{enumerate}

\item 向上缩放 $M \longrightarrow \Delta M \text { } \in \mathcal{R}_{\langle n, q\rangle}$

\item 计算：

$B = \mathit{PK}_1\cdot U + \Delta  M + E_1 \text{ } \in \mathcal{R}_{\langle n,q \rangle}$

$\vec{D} = \vv{{\mathit{PK}}}_2 \cdot U + \vec{E}_2 \in \mathcal{R}_{\langle n,q \rangle}^{k}$ \textcolor{red}{ $\rhd$ $\vv{\mathit{PK}}_2 \cdot U$ 将 $\vv{\mathit{PK}}_2$ 的每个元素乘以 $U$}

\item $\textsf{GLWE}_{S,\sigma}(\Delta M + E_{\mathit{all}}) = (\vec{D}, B) \text{ } \in \mathcal{R}_{\langle n,q \rangle}^{k+1}$ \textcolor{red}{ $\rhd$ 其中 $E_{\mathit{all}} = E \cdot U + E_1 - \vec{E}_2\cdot\vec{S}$}



\end{enumerate}


\par\noindent\rule{\textwidth}{0.4pt}

\textbf{\underline{解密输入}：}$\textsf{ct} = (\vec{D}, B) \text{ } \in \mathcal{R}_{\langle n,q \rangle}^{k + 1}$

\begin{enumerate}
\item $\textsf{GLWE}^{-1}_{S,\sigma}(\textsf{ct}) = B - \vec{D} \cdot \vec{S} = \Delta  M + E_{\mathit{all}} \text{ } \in \mathcal{R}_{\langle n,q \rangle}$

\item 向下缩放
$\Bigl\lceil \dfrac{ \Delta  M + E_{\mathit{all}} }{\Delta}\Bigr\rfloor \bmod t = M  \text{ } \in \mathcal{R}_{\langle n, t \rangle}$

\end{enumerate}

为了正确解密，噪声多项式 $E_{\mathit{all}}$ 的每个系数 $e_i$ 都应满足 $|e_i| < \dfrac{\Delta}{2}$。

\end{tcolorbox}


解密过程第一步中的等式推导如下：

$\textsf{GLWE}^{-1}_{S,\sigma}\bm{(}\text{ } \textsf{ct} = (\vec{D}, B) \text{ }\bm{)} = B - \vec{D}\cdot\vec{S}$

$ = (\mathit{PK}_1\cdot U + \Delta  M + E_1) - (\vv{\mathit{PK}}_2 \cdot U + \vec{E}_2)\cdot \vec{S} $

$=  (\vec{A} \cdot \vec{S} + E)\cdot U + \Delta M + E_1 - (\vec{A}\cdot U)\cdot \vec{S} - \vec{E}_2\cdot\vec{S}$

$=  (U\cdot\vec{A}) \cdot \vec{S} + E \cdot U + \Delta M + E_1 - (U\cdot \vec{A}) \cdot \vec{S} - \vec{E}_2\cdot\vec{S}$

$= \Delta M + E \cdot U + E_1 - \vec{E}_2\cdot\vec{S}$

$= \Delta M + E_{\mathit{all}}$ \textcolor{red}{ $\rhd$ 其中 $E_{\mathit{all}} = E \cdot U + E_1 - \vec{E}_2\cdot\vec{S}$}


$ $

\para{安全性：}在 $k=1$ 时，将公钥写成 $(A,P)$，其中 $P=AS+E$；密文为 $(D,B)=(AU+E_2,PU+E_1+\Delta M)$。假设针对实际 $S$ 分布及误差 $E$ 的判定 RLWE 成立，也假设针对实际 $U$ 分布的两样本判定 RLWE 成立，其两个误差分别服从 $E_2,E_1$ 的分布。本小节中两个秘密均为三元分布，三种误差分布相同，因此同一个相应的两样本 RLWE 假设可覆盖两次使用。仅有针对高斯秘密的定理并不能推出这一点。

单挑战 IND-CPA 证明使用两个混合。首先，将 $P=AS+E$ 替换为独立均匀环元素。归约将一条 RLWE 样本嵌入公钥，自行选择新鲜 $U,E_1,E_2$ 就能执行全部加密，不需要知道 $S$。随后，在均匀 $(A,P)$ 下，将挑战的 $(D,B-\Delta M)$ 替换为独立均匀环元素。归约取得共享同一个未知秘密 $U$ 的两条样本 $(A,Z_2),(P,Z_1)$，以其掩码作为公钥，等挑战消息选定后返回 $(Z_2,Z_1+\Delta M)$。真实样本给出 $Z_2=AU+E_2$ 和 $Z_1=PU+E_1$；均匀样本则完全隐藏选中的消息。其他加密都可用新鲜、已知的随机性公开执行。因此，采用 \autoref{subsec:lattice-scheme} 的猜测偏差约定，有
\[
\epsilon_{\mathrm{PKE}}\le\delta_S(T'_1,1)+\delta_U(T'_2,2).
\]
这两项使用所述实际分布，并计入归约开销。这是限定构造的完整条件性 CPA 归约，并非某个数值参数族的认证，也不涵盖依赖秘密的公钥加密查询。参见 \href{https://web.eecs.umich.edu/~cpeikert/pubs/ideal-lwe.pdf}{Lyubashevsky--Peikert--Regev，第 1.1 节 Applications}。对于未加限制的向量版本，第二个混合需要的是以 $U$ 为秘密、具有 $k+1$ 条样本的秩一 RLWE，而不仅是第一个混合所用的秩 $k$ GLWE 假设。

抵消运算另外给出了正确性恒等式，其中 $E_{\mathit{all}}=EU+E_1-E_2S$。

在负循环环中，对提升后的系数有
\[
\|E_{\mathit{all}}\|_\infty\le\|E\|_\infty\|U\|_1+\|E_1\|_\infty+\|E_2\|_\infty\|S\|_1.
\]
当 $t\mid q$ 时，右端小于 $\Delta/2$ 即足以保证解密。

如果不限制适用范围，那么在 $n=1$ 且 $U$ 为三元数时，即使 LWE 向量维数 $k$ 很大，这个构造也不安全：攻击者只需尝试 $U\in\{-1,0,1\}$，检验所有 $D_i-A_iU$ 是否很小，再解码 $B-\mathit{PK}_1U$。特别地，$U=0$ 以 $1/3$ 的概率发生，此时直接得到 $B=\Delta m+E_1$。因此，不能从展示的代数公式推出一般 GLWE 公钥安全性。


\clearpage

\section{GLev}
\label{sec:glev}
GLev 密文是由多个 GLWE 密文组成的列表，分别加密明文列表 $\dfrac{q}{\beta^1}M, \dfrac{q}{\beta^2}M, \ldots, \dfrac{q}{\beta^l}M$，其中 $M$ 是编码为多项式的明文。GLev 密文中的第 $i$ 个 GLWE 密文使用不同的明文尺度因子 $\Delta_i = \dfrac{q}{\beta^i}$。\autoref{fig:glev} 直观展示了 GLev 密文的结构。


对于这里的整数环权重，要求整数 $\beta\ge2$ 且 $\beta^l\mid q$。只有在期望第一行恢复取自 $\{0,\ldots,t-1\}$ 的系数代表元时，才需要 $t\le\beta$。GLev 也可作为一般环消息的 gadget 密文使用；在这种用法中，并不要求每一行都能单独解码原消息。若要用降幂 gadget 精确重构每个剩余，还须满足 $q=\beta^l$，如 \autoref{sec:decomp} 所述。


\subsection{加密}
\label{subsec:glev-enc}

\begin{tcolorbox}[title={\textbf{\tboxlabel{\ref*{subsec:glev-enc}} GLev 加密}}]


$\textsf{GLev}_{S, \sigma}^{\beta, l}(M) = \Bigl \{\textsf{GLWE}_{S, \sigma}\left(\dfrac{q}{\beta^i} M + E_i\right)  \Bigr \}_{i=1}^{l} \in \mathcal{R}_{\langle n, q \rangle }^{(k+1) \cdot l}$
\end{tcolorbox}

\begin{figure}[h!]
    \centering
  \includegraphics[width=1.0\linewidth]{figures/TFHE-fig2.pdf}
  \caption{GLev 密文结构示意图。}
  \label{fig:glev}
\end{figure}



\subsection{解密}

第一行的相位为 $(q/\beta)M+E_1\pmod q$。若 $\|E_1\|_\infty<q/(2\beta)$，最近整数解码逐系数恢复 $M\bmod\beta$。只有当原代表元能由其模 $\beta$ 的类唯一确定时，才能恢复它，例如系数属于 $\{0,\ldots,t-1\}$ 且 $t\le\beta$。类似地，若 $\|E_i\|_\infty<q/(2\beta^i)$，第 $i$ 行能解码模 $\beta^i$ 的结果。后面的行对误差的容忍度更低，但不一定无法解码；gadget 运算也不要求各行单独解码。

\subsection{Lev 与 RLev}

Lev 是 $n=1$ 时的 GLev，RLev 是 $k=1$ 时的 GLev。



\clearpage

\section{GGSW}
\label{sec:ggsw}
GGSW 密文由一组 GLev 密文组成。在 GGSW 密码体制中，秘密密钥 $S$ 是由 $k$ 个多项式组成的列表，即 $S_0, S_1, ... \text{ } S_{k-1}$。GGSW 密文中的各个 GLev 密文（以 $i$ 索引）分别加密明文 $-S_0 \cdot M, -S_1 \cdot M, \ldots, -S_{k-1} \cdot M$，以及 $M$。\autoref{fig:ggsw} 直观展示了这一结构。


\subsection{加密}
\label{subsec:ggsw-enc}

\begin{tcolorbox}[title={\textbf{\tboxlabel{\ref*{subsec:ggsw-enc}} GGSW 加密}}]
$\textsf{GGSW}_{S, \sigma}^{\beta, l}(M) = \Bigl \{ \{ \textsf{GLev}_{S, \sigma}^{\beta, l}(-S_i \cdot M)  \}_{i=0}^{k-1}, \textsf{GLev}_{S, \sigma}^{\beta, l}(M) \Bigr \} \in \mathcal{R}_{\langle n, q \rangle }^{(k+1) \cdot l \cdot (k+1)}$
\end{tcolorbox}

\begin{figure}[h!]
    \centering
  \includegraphics[width=1.0\linewidth]{figures/TFHE-fig3.pdf}
  \caption{GGSW 密文结构示意图。}
  \label{fig:ggsw}
\end{figure}


\subsection{解密}

在 GLev 的消息范围与噪声条件下，解密最后一个 GLev 块即可恢复 $M$。其他块的第一行在相应误差界下解码得到 $-S_iM\bmod\beta$。要从该乘积恢复 $M\bmod\beta$，要求 $S_i$ 是 $\mathbb Z_\beta[X]/(X^n+1)$ 中的单位。非零并不充分：在 $\mathbb Z_5[X]/(X^2+1)$ 中，$X-2$ 和 $X+2$ 均非零，但 $(X-2)(X+2)=0$。使用最后一个块时，无需对 $S_i$ 求逆。



\subsection{GSW 与 RGSW}

GSW 是 $n=1$ 时的 GGSW，RGSW 是 $k=1$ 时的 GGSW。

\clearpage


\part{全同态加密的通用技术}
\label{part:generic-fhe}

\renewcommand{\thesection}{C-\arabic{section}}
\setcounter{section}{0}

本部分介绍 TFHE、CKKS、BGV 和 BFV 等多种 FHE 方案共同使用的同态计算技术。

根据 \autoref{sec:glwe}，$\textsf{GLWE}_{S,\sigma}(\Delta M + E) = (A_0, A_1, \gap{$\cdots$} , A_{k-1}, B) \in \mathcal{R}_{\langle n, q \rangle }^{k + 1}$，其中 $\mathcal{R}_{\langle n,q \rangle} = \mathbb{Z}_q[x] / (x^n + 1)$，$B$ 按 $B = \sum\limits_{i=0}^{k-1}{(A_i \cdot S_i)} + \Delta \cdot M + E$ 计算。每个 $A_i$ 是作为公钥分量的 $(n-1)$ 次多项式，其各个系数在 $\mathcal{R}_{\langle n, q \rangle }$ 中均匀随机采样。$E$ 是作为噪声的 $(n-1)$ 次多项式，其各个系数依据高斯分布 $\chi_\sigma$ 从 $\mathcal{R}_{\langle n, q \rangle }$ 中采样。私钥 $S$ 是由 $k$ 个 $(n-1)$ 次多项式构成的列表，满足 $S = (S_0, S_1, \gap{$\cdots$} S_{k-1} ) \in \mathcal{R}_{\langle n, q \rangle }^k$；每个多项式 $S_i$ 的各个系数都是从 $\mathbb{Z}_2$（即 $\{0, 1\}$）中随机采样得到的二元数。

基于这一 GLWE 设置，本部分介绍同态运算的通用构件，包括密文与密文加法、密文与明文加法、密文与明文乘法、模数切换和密钥切换。


\clearpage


\section{GLWE 密文与密文加法}
\label{sec:glwe-add-cipher}
两个输入必须具有相同的秘密密钥、环、模数 $q$ 和尺度 $\Delta=q/t$；所述模 $t$ 的明文运算假设 $t\mid q$。在这些运算章节中，$\textsf{GLWE}_{S,\sigma}(\text{phase})$ 是具有所列相位的密文的简写；求值后的密文不一定具有新鲜加密时相同的高斯分布或标准差 $\sigma$。所有噪声界均针对系数约减前的整数提升。

\textbf{- 参考资料：}
\href{https://www.zama.ai/post/tfhe-deep-dive-part-2}{深入理解 TFHE（二）：编码与线性分层运算}~\cite{tfhe-2}


$ $

假设有两个 GLWE 密文，分别加密不同的明文 $M^{\langle 1 \rangle}, M^{\langle 2 \rangle}$：

$\textsf{GLWE}_{S, \sigma}(\Delta M^{\langle 1 \rangle} + E^{\langle 1 \rangle}) = \textsf{ct}^{\langle 1 \rangle} = ( A_0^{\langle 1 \rangle}, A_1^{\langle 1 \rangle}, \ldots, A_{k-1}^{\langle 1 \rangle}, B^{\langle 1 \rangle}) \in \mathcal{R}_{\langle n,q \rangle}^{k + 1}$

$\textsf{GLWE}_{S, \sigma}(\Delta M^{\langle 2 \rangle} + E^{\langle 2 \rangle}) = \textsf{ct}^{\langle 2 \rangle} = ( A_0^{\langle 2 \rangle}, A_1^{\langle 2 \rangle}, \ldots, A_{k-1}^{\langle 2 \rangle}, B^{\langle 2 \rangle}) \in \mathcal{R}_{\langle n,q \rangle}^{k + 1}$

$ $

\noindent 定义如下密文加法运算：

$\textsf{ct}^{\langle 1 \rangle} + \textsf{ct}^{\langle 2 \rangle} = ( A_0^{\langle 1 \rangle} + A_0^{\langle 2 \rangle}, \text{ } A_1^{\langle 1 \rangle} + A_1^{\langle 2 \rangle}, \ldots, A_{k-1}^{\langle 1 \rangle} + A_{k-1}^{\langle 2 \rangle}, \text{ } B^{\langle 1 \rangle} + B^{\langle 2 \rangle} )$

$ $

\noindent 则有以下结论：

\begin{tcolorbox}[title={\textbf{\tboxlabel{\ref*{sec:glwe-add-cipher}} GLWE 同态加法}}]
$\textsf{GLWE}_{S, \sigma}(\Delta M^{\langle 1 \rangle} + E^{\langle 1 \rangle} ) + \textsf{GLWE}_{S, \sigma}(\Delta M^{\langle 2 \rangle} + E^{\langle 2 \rangle}) $

$ = ( \{A_i^{\langle 1 \rangle}\}_{i=0}^{k-1}, \text{ } B^{\langle 1 \rangle}) + (\{A_i^{\langle 2 \rangle}\}_{i=0}^{k-1}, \text{ } B^{\langle 2 \rangle}) $

$ = ( \{A_i^{\langle 1 \rangle} + A_i^{\langle 2 \rangle}\}_{i=0}^{k-1}, \text{ } B^{\langle 1 \rangle} + B^{\langle 2 \rangle} ) $

$= \textsf{GLWE}_{S, \sigma}(\Delta(M^{\langle 1 \rangle} + M^{\langle 2 \rangle}) + E^{\langle 3 \rangle})$ \textcolor{red}{ $\rhd$ 其中 $E^{\langle 3 \rangle} = E^{\langle 1  \rangle} + E^{\langle 2 \rangle}$}
\end{tcolorbox}


相位恒等式在代数上无条件成立。若 $\|E^{\langle1\rangle}+E^{\langle2\rangle}\|_\infty<\Delta/2$，解码才会返回 $M^{\langle1\rangle}+M^{\langle2\rangle}\bmod t$；例如，两个输入的噪声界之和小于 $\Delta/2$ 就是充分条件。

$ $

%\noindent \textbf{\underline{Proof}}
\begin{myproof}
\begin{enumerate}
\item 定义以下记号： \\
$A_0^{\langle 3 \rangle} = A_0^{\langle 1 \rangle} + A_0^{\langle 2 \rangle}$ \\
$A_1^{\langle 3 \rangle} = A_1^{\langle 1 \rangle} + A_1^{\langle 2 \rangle}$ \\
$\vdots$ \\
$A_{k-1}^{\langle 3 \rangle} = A_{k-1}^{\langle 1 \rangle} + A_{k-1}^{\langle 2 \rangle}$ \\
$E^{\langle 3 \rangle} = E^{\langle 1 \rangle} + E^{\langle 2 \rangle}$ \\
$B^{\langle 3 \rangle} = B^{\langle 1 \rangle} + B^{\langle 2 \rangle}$
\item 可推导得： \\
$B^{\langle 3 \rangle} = B^{\langle 1 \rangle} + B^{\langle 2 \rangle}$ \\
$ = \sum\limits_{i=0}^{k-1}{(A_i^{\langle 1 \rangle} \cdot S_i)} + \Delta \cdot M^{\langle 1 \rangle} + E^{\langle 1 \rangle} + \sum\limits_{i=0}^{k-1}{(A_i^{\langle 2 \rangle} \cdot S_i)} + \Delta \cdot M^{\langle 2 \rangle} + E^{\langle 2 \rangle}$ \\ 
$= \sum\limits_{i=0}^{k-1}{((A_i^{\langle 1 \rangle} + A_i^{\langle 2 \rangle}) \cdot S_i)} + \Delta \cdot (M^{\langle 1 \rangle} + M^{\langle 2 \rangle}) + (E^{\langle 1 \rangle} + E^{\langle 2 \rangle})$ \textcolor{red}{ $\rhd$ 由交换律与分配律} \\
$= \sum\limits_{i=0}^{k-1}{(A_i^{\langle 3 \rangle} \cdot S_i)} + \Delta \cdot (M^{\langle 1 \rangle} + M^{\langle 2 \rangle}) + E^{\langle 3 \rangle}$ \\
\item 由于 $B^{\langle 3 \rangle} = \sum\limits_{i=0}^{k-1}{(A_i^{\langle 3 \rangle} \cdot S_i)} + \Delta \cdot (M^{\langle 1 \rangle} + M^{\langle 2 \rangle}) + E^{\langle 3 \rangle}$，

因此 $(A_0^{\langle 3 \rangle}, A_1^{\langle 3 \rangle}, A_2^{\langle 3 \rangle}, \ldots, A_{k-1}^{\langle 3 \rangle}, B^{\langle 3 \rangle})$ 构成密文 $\textsf{GLWE}_{S, \sigma}(\Delta \cdot (M^{\langle 1 \rangle} + M^{\langle 2 \rangle}) + E^{\langle 3 \rangle})$。 \\
\item 因此， \\
$\textsf{GLWE}_{S, \sigma}(\Delta M^{\langle 1 \rangle} + E^{\langle 1 \rangle}) + \textsf{GLWE}_{S, \sigma}(\Delta M^{\langle 2 \rangle} + E^{\langle 2 \rangle})$ \\
$ = ( A_0^{\langle 1 \rangle} + A_0^{\langle 2 \rangle}, \text{ } A_1^{\langle 1 \rangle} + A_1^{\langle 2 \rangle},\ldots, A_{k-1}^{\langle 1 \rangle} + A_{k-1}^{\langle 2 \rangle}, \text{ } B^{\langle 1 \rangle} + B^{\langle 2 \rangle} )$ \\
$ = ( A_0^{\langle 3 \rangle}, A_1^{\langle 3 \rangle}, A_2^{\langle 3 \rangle}, \ldots,  A_{k-1}^{\langle 3 \rangle}, B^{\langle 3 \rangle})$ \\
$ = ( \{A_i^{\langle 3 \rangle}\}_{i=0}^{k-1}, B^{\langle 3 \rangle})$ \\
$= \textsf{GLWE}_{S, \sigma}(\Delta (M^{\langle 1 \rangle} + M^{\langle 2 \rangle}) + E^{\langle 3 \rangle})$

%\begin{flushright}
%\qedsymbol{} 
%\end{flushright}

\end{enumerate}
\end{myproof}

\subsection{讨论}
\label{subsubsec:glwe-add-cipher-discuss} 

\para{噪声消除：}使用秘密密钥作减法会精确消去掩码乘积；只有所得误差各系数的绝对值仍小于 $\Delta/2$ 时，舍入才会去除误差。若合并后的误差超出此界，两个有效输入也可能相加得到无法正确解码的输出。

$ $

\para{噪声增长：}新误差为 $E^{\langle3\rangle}=E^{\langle1\rangle}+E^{\langle2\rangle}$。独立零均值误差的方差相加；若各自方差为 $\sigma^2$，则 $h$ 个独立样本之和的方差为 $h\sigma^2$。仅有零均值并不充分。一般有 $\operatorname{Var}(e_1+e_2)=\operatorname{Var}(e_1)+\operatorname{Var}(e_2)+2\operatorname{Cov}(e_1,e_2)$。密文与自身相加得到 $2e$，其方差为 $4\sigma^2$，并非 $2\sigma^2$。确定性的系数界可由三角不等式给出，无需独立性假设。

$ $

\para{$t\nmid q$ 时的明文溢出：}掩码乘积和密文分量始终按模 $q$ 解释。当 $t\mid q$ 时，明文模 $t$ 的回绕与乘以 $\Delta=q/t$ 相容。当 $t\nmid q$ 且 $\Delta=\lfloor q/t\rfloor$ 时，将未约减明文的每个系数写为 $m=m'+vt$，其中 $m'$ 为中心化代表元，则有效误差是 $e-(q\bmod t)v$。在 \autoref{subsubsec:lwe-noise-bound} 的假设下，要求 $|e-(q\bmod t)v|<\Delta/2$。这一条件依赖溢出次数 $v$，并非仅由 $q\gg t$ 保证。最后模 $t$ 的约减也是消息解密的一部分。

 


\clearpage

\section{GLWE 密文与明文加法}
\label{sec:glwe-add-plain}
密文与 $\Lambda$ 的整数提升使用相同的环和尺度。为得到所述模 $t$ 的解码明文，假设 $t\mid q$、$\Delta=q/t$ 且 $\|E\|_\infty<\Delta/2$。

假设有一个 GLWE 密文 \textsf{ct} 和一个新的明文多项式 $\Lambda$，如下所示：


$\textsf{ct} = \textsf{GLWE}_{S, \sigma}(\Delta M + E) = ( A_0, A_1, \ldots, A_{k-1}, B) \in \mathcal{R}_{\langle n,q \rangle}^{k + 1}$

$\Lambda$：新的明文多项式

$\Delta \Lambda$：按 $\Delta$ 放大后的新明文多项式

$ $

\noindent 定义如下密文与明文加法运算：

$\textsf{ct} + \Delta\Lambda = ( A_0, \text{ } A_1, \ldots, A_{k-1}, \text{ } B + \Delta \Lambda)$

$ $

\noindent 则有以下结论：

\begin{tcolorbox}[title={\textbf{\tboxlabel{\ref*{sec:glwe-add-plain}} GLWE 密文与明文的同态加法}}]
$\textsf{GLWE}_{S, \sigma}(\Delta M + E) + \Delta\Lambda $

$=  (\{A_i\}_{i=0}^{k-1}, \text{ } B) + \Delta\Lambda$

$=  (\{A_i\}_{i=0}^{k-1}, \text{ } B + \Delta\Lambda)$

$= \textsf{GLWE}_{S, \sigma}(\Delta (M + \Lambda) + E)$
\end{tcolorbox}


这意味着，将按 $\Delta$ 放大后的明文多项式 $\Lambda$ 加到加密 $M$ 的 GLWE 密文上，再解密所得密文，便得到 $M + \Lambda$。




\begin{myproof}

\begin{enumerate}
\item 由于 $B = \sum\limits_{i=0}^{k-1}{(A_i \cdot S_i)} + \Delta \cdot M + E$，

$B + \Delta \cdot \Lambda = \sum\limits_{i=0}^{k-1}{(A_i \cdot S_i)} + \Delta \cdot M + E + \Delta \cdot \Lambda = \sum\limits_{i=0}^{k-1}{(A_i \cdot S_i)} + \Delta \cdot (M + \Lambda) + E$  \\

这意味着 $(A_0, A_1, ... \text{ } A_{k-1}
, B + \Delta \Lambda)$ 构成密文 $\textsf{GLWE}_{S, \sigma}(\Delta (\Lambda + M) + E)$。
\item 因此， \\
$\textsf{GLWE}_{S, \sigma}(\Delta M + E) + \Delta\Lambda$ \\
$ = (A_0, \text { } A_1, ... \text{ } A_{k-1}, \text { } B + \Delta \Lambda)$ \\
$ = (\{A_i\}_{i=0}^{k-1}, \text { } B + \Delta \Lambda)$ \\
$= \textsf{GLWE}_{S, \sigma}(\Delta (M + \Lambda) + E)$

%\begin{flushright}
%\qedsymbol{} 
%\end{flushright}
\end{enumerate}

\end{myproof}


\para{噪声增长：}相对于未约减消息 $M+\Lambda$，相位误差仍为 $E$。当 $t\mid q$ 时，将消息模 $t$ 约减不会改变有效误差。若 $t\nmid q$，消息回绕可能引入 \autoref{subsubsec:lwe-noise-bound} 描述的溢出误差；此时“噪声不变”只能理解为相对于未约减消息而言。

\clearpage

\section{GLWE 密文与明文乘法}
\label{sec:glwe-mult-plain}
消息解码假设 $t\mid q$ 且 $\Delta=q/t$。约束噪声时，对公开乘数 $\Lambda$ 采用固定的系数整数提升。下面的密文相位等式并不声称其分布与新鲜加密相同。

\textbf{- 参考资料：}
\href{https://www.zama.ai/post/tfhe-deep-dive-part-2}{深入理解 TFHE（二）：编码与线性分层运算}~\cite{tfhe-2}

$ $

假设有一个 GLWE 密文 \textsf{ct}：

$\textsf{ct} = \textsf{GLWE}_{S, \sigma}(\Delta M + E) = ( A_0, A_1, \ldots, A_{k-1}, B) \in \mathcal{R}_{\langle n,q \rangle}^{k + 1}$

$ $

\noindent 以及如下新的明文多项式 $\Lambda$：

$\Lambda = \sum\limits_{i=0}^{n-1}(\Lambda_i \cdot X^i) \in \mathcal{R}_{\langle n, q \rangle}$

$ $

\noindent 定义如下密文与明文乘法运算：

$\Lambda \cdot \textsf{ct} = (\Lambda \cdot A_0, \Lambda \cdot A_1, \ldots, \Lambda \cdot A_{k-1}, \Lambda \cdot B)$

$ $

\noindent 当满足变换所需的单位根与可逆性条件时，NTT 乘法使用 $O(n\log n)$ 次域运算（\autoref{sec:ntt}）；例如，满足 $2n\mid q-1$ 的素数 $q$ 支持负循环变换。任意模数并不自动具备这样的 NTT。下面的代数恒等式与所选乘法算法无关：

\begin{tcolorbox}[title={\textbf{\tboxlabel{\ref*{sec:glwe-mult-plain}} GLWE 密文与明文的乘法}}]
$\Lambda \cdot \textsf{GLWE}_{S, \sigma}(\Delta M + E)$

$= \Lambda \cdot (\{A_i^{\langle 1 \rangle}\}_{i=0}^{k-1}, \text{ } B^{\langle 1 \rangle})$

$= (\{\Lambda\cdot A_i^{\langle 1 \rangle}\}_{i=0}^{k-1}, \text{ } \Lambda \cdot B^{\langle 1 \rangle})$

$= \textsf{GLWE}_{S, \sigma}(\Delta (M \cdot \Lambda) + \Lambda\cdot E )$
\end{tcolorbox}

所得相位为 $\Delta M\Lambda+\Lambda E$。只有当 $\|\Lambda E\|_\infty<\Delta/2$ 时，它才解码得到 $M\Lambda\bmod t$。

$ $

%\noindent \textbf{\underline{Proof}}
\begin{myproof}
\begin{enumerate}
\item 定义以下记号： \\
$A_0' = \Lambda \cdot A_0$ \\
$A_1' = \Lambda \cdot A_1$ \\
$\vdots$ \\
$A_{k-1}' = \Lambda \cdot A_{k-1}$ \\
$E' = \Lambda \cdot E$ \\
$B' = \Lambda \cdot B$ \\
\item 可推导得： \\
$B' = \Lambda \cdot B$ \\
$= \Lambda \cdot (\sum\limits_{i=0}^{k-1}{(A_i \cdot S_i)} + \Delta \cdot M + E)$ 
$= \sum\limits_{i=0}^{k-1}{(\Lambda \cdot A_i \cdot S_i)} + \Delta \cdot \Lambda \cdot M + \Lambda \cdot E$ \\ \textcolor{red}{ $\rhd$ 由多项式环的分配律} \\
$= \sum\limits_{i=0}^{k-1}{((\Lambda \cdot A_i) \cdot S_i)} + \Delta \cdot (\Lambda \cdot M) + (\Lambda \cdot E)$  \\
$= \sum\limits_{i=0}^{k-1}{(A_i' \cdot S_i)} + \Delta \cdot (\Lambda \cdot M) + (E')$ \\
\item 由于 $B' = \sum\limits_{i=0}^{k-1}{(A_i' \cdot S_i)} + \Delta \cdot (\Lambda \cdot M) + (E')$，

$(A_0', A_1', \ldots, A_{k-1}'
, B')$ 构成密文 $\textsf{GLWE}_{S, \sigma}(\Delta \cdot \Lambda \cdot M+E^{\prime})$。
\item 因此， \\
$\Lambda \cdot \textsf{GLWE}_{S, \sigma}(\Delta M + E)$ \\
$ = (\Lambda \cdot A_0, \text { } \Lambda \cdot A_1, \ldots, \Lambda \cdot A_{k-1}, \text { } \Lambda \cdot B)$ \\
$ = ( \{A'_{i}\}_{i=0}^{k-1}, \text { } \Lambda \cdot B)$ \\
$= \textsf{GLWE}_{S, \sigma}(\Delta (M \cdot \Lambda) + \Lambda \cdot E)$

%\begin{flushright}
%\qedsymbol{} 
%\end{flushright}
\end{enumerate}
\end{myproof}

使用秘密密钥作减法会精确消去掩码乘积，但舍入只能在解码界内去除 $E'=\Lambda E$。在负循环环中，每个输出系数都是至多 $n$ 个乘积的带符号和，因此
\[
\|\Lambda E\|_\infty\le\|\Lambda\|_1\|E\|_\infty\le n\|\Lambda\|_\infty\|E\|_\infty.
\]
这里 $\|\Lambda\|_1$ 表示所选提升的系数绝对值之和。

重复乘法可能使噪声界增大，直至无法正确解码。解密用于读取结果，并不会刷新已有密文。对于较大的乘数，预先计算的 GLev 表示可以通过下面的 gadget 分解（\autoref{subsec:gadget-decomposition}）给出更好的误差界。

\subsection{利用 gadget 分解抑制噪声}
\label{subsubsec:gadget-decomposition-noise-suppression}

这里使用精确的降幂 gadget $g_i=q/\beta^i$，要求 $q=\beta^l$，并采用 $\{0,\ldots,\beta-1\}$ 中的无符号数位。预先计算的输入是 $\textsf{GLev}^{\beta,l}_{S,\sigma}(\Delta M)$，而非任意单个 GLWE 密文。这样的低噪声 GLev 表示是额外的加密输入，不能假设经过先前运算的密文会自动免费获得它。

在密文与明文乘法 $\Lambda \cdot \textsf{GLWE}_{S, \sigma}(\Delta M)$ 中，噪声 $E$ 增长为 $E' = \Lambda \cdot E$。为限制这种增长，可以对 $\Lambda$ 进行分解（\autoref{subsec:number-decomp}），并使用 $M$ 的 GLev 加密（\autoref{subsec:glev-enc}），具体如下：


$\Lambda = \Lambda_1 \dfrac{q}{\beta^1} + \Lambda_2 \dfrac{q}{\beta^2} + \cdots + \Lambda_l \dfrac{q}{\beta^l} \longrightarrow \textsf{Decomp}^{\beta, l}(\Lambda) = (\Lambda_1, \Lambda_2, \cdots, \Lambda_l)$

$ $

$\textsf{GLev}_{S, \sigma}^{\beta, l}(\Delta M) = \Bigg\{ \textsf{GLWE}_{S, \sigma}\left(\Delta M \dfrac{q}{\beta^1} + E_1\right), \textsf{GLWE}_{S, \sigma}\left(\Delta M \dfrac{q}{\beta^2} + E_2\right), \cdots \textsf{GLWE}_{S, \sigma}\left(\Delta M \dfrac{q}{\beta^l} + E_l\right) \Bigg\}$

$ $


将明文 $M$ 加密为 $\textsf{GLev}_{S, \sigma}^{\beta, l}(\Delta M)$，而非 $\textsf{GLWE}_{S, \sigma}(\Delta M)$；相应地，计算 $\textsf{Decomp}^{\beta, l}(\Lambda) \cdot \textsf{GLev}_{S, \sigma}^{\beta, l}(\Delta M)$，而非 $\Lambda \cdot \textsf{GLWE}_{S, \sigma}(\Delta M)$。注意，两种计算得到的明文结果相同，如下所示：

$\textsf{Decomp}^{\beta, l}(\Lambda) \cdot \textsf{GLev}_{S, \sigma}^{\beta, l}(\Delta M)$

$= (\Lambda_1, \Lambda_2, \cdots, \Lambda_l) \cdot \left (\textsf{GLWE}_{S, \sigma}\left(\dfrac{q}{\beta} \Delta M + E_1\right), \text{ } \textsf{GLWE}_{S, \sigma}\left(\dfrac{q}{\beta^2} \Delta M + E_2\right), \text{ } \cdots, \text{ } \textsf{GLWE}_{S, \sigma}\left(\dfrac{q}{\beta^l} \Delta M + E_l\right) \right )$

$= \Lambda_1\cdot\textsf{GLWE}_{S, \sigma}\left(\dfrac{q}{\beta} \Delta M + E_1\right) +  \Lambda_2\cdot\textsf{GLWE}_{S, \sigma}\left(\dfrac{q}{\beta^2} \Delta M + E_2\right) + \cdots + \Lambda_l\cdot\textsf{GLWE}_{S, \sigma}\left(\dfrac{q}{\beta^l} \Delta M + E_l\right)$

$= \textsf{GLWE}_{S, \sigma}\left(\Lambda_1\cdot\dfrac{q}{\beta}\Delta M + \Lambda_1 E_1\right) +\textsf{GLWE}_{S, \sigma}\left(\Lambda_2\cdot\dfrac{q}{\beta^2}\Delta M + \Lambda_2 E_2\right)+ \cdots + \textsf{GLWE}_{S, \sigma}\left(\Lambda_l\cdot\dfrac{q}{\beta^l}\Delta M + \Lambda_l E_l\right)$

$= \textsf{GLWE}_{S, \sigma}\left(\Lambda_1\cdot\dfrac{q}{\beta} \Delta M + \Lambda_2\cdot\dfrac{q}{\beta^2} \Delta M + \cdots + \Lambda_l\cdot\dfrac{q}{\beta^l} \Delta M+E_{\textit{all}}\right)$

$= \textsf{GLWE}_{S, \sigma}\left(\left(\Lambda_1\cdot\dfrac{q}{\beta} + \Lambda_2\cdot\dfrac{q}{\beta^2} + \cdots + \Lambda_l\cdot\dfrac{q}{\beta^l}\right)\cdot \Delta M  + E_{\textit{all}} \right)$ \textcolor{red}{ $\rhd$ 其中 $E_{\textit{all}} = \sum\limits_{i=1}^l \Lambda_iE_i$}

$= \textsf{GLWE}_{S, \sigma}\left(\Lambda \cdot \Delta M + E_{\textit{all}}\right)$ \textcolor{red}{ $\rhd$ 解密结果为 $\Lambda\cdot M$}

$ $

精确分解时，最终误差为 $E_{\textit{all}}=\sum_{i=1}^l\Lambda_iE_i$。若 $\|E_i\|_\infty\le B_{\mathrm{g}}$，无符号数位给出
\[
\|E_{\textit{all}}\|_\infty\le\sum_{i=1}^l\|\Lambda_i\|_1\|E_i\|_\infty\le nl(\beta-1)B_{\mathrm{g}}.
\]
若要使用 $\beta/2$ 的数位界，必须明确改用平衡分解。这可能改善大乘数对应的最坏情形误差界，但不能保证每个输入的实际误差都更小。只有各自满足误差界时，两种计算才会解码得到相同消息。对于截断分解 $\Lambda\equiv\sum_i\Lambda_i g_i+\rho\pmod q$，最终相位的误差变为 $\sum_i\Lambda_iE_i-\rho\Delta M$；还必须约束这一残差。图~\ref{fig:decomp2} 展示了预期的误差界改善。

\begin{figure}[h!]
    \centering
  \includegraphics[width=0.8\linewidth]{figures/decomp2.pdf}
  \caption{利用 gadget 分解降低密文与明文乘法中的噪声。}
  \label{fig:decomp2}
\end{figure}

\subsubsection{讨论}
\label{subsubsec:glwe-mult-plain-discuss}


分解技术还具有重要用途：在 GLWE 密钥切换（\autoref{sec:glwe-key-switching}）中，我们将介绍如何将分解后的掩码多项式 $\textsf{Decomp}^{\beta,l}(A_i)$
与预先计算的密钥切换密钥
$\textsf{KSK}_i=\textsf{GLev}^{\beta,l}_{S',\sigma}(S_i)$,
结合起来实现密钥切换。因此，尽管每次操作输出的都是普通 GLWE 密文，gadget 分解仍可在多次密钥切换调用中反复使用。

对于反复重新初始化普通密文噪声 $E$ 的技术，我们将在 \autoref{subsec:tfhe-noise-bootstrapping} 介绍 TFHE 的噪声自举。


\clearpage

\section{GLWE 模数切换}
\label{sec:modulus-switching}
在 GLWE 密码体制中，模数切换是将密文的模数域改为更小（或更大）的模数域，同时保证密文仍能解密得到相同明文的过程。例如，设有密文 $\textsf{LWE}_{S, \sigma}(\Delta m) \in \mathbb{Z}_q^{k+1}$。如果将密文模数按 $q \rightarrow q'$ 切换，则密文变为 $\textsf{LWE}_{S, \sigma}\left(\Delta\dfrac{q'}{q} m + e'\right) \in \mathbb{Z}_{q'}^{k+1}$。密文的其他分量，例如噪声（$e$）和公钥分量（$a_0, a_1, \cdots, a_{k-1}$、$b$），均按 $\dfrac{q'}{q}$ 缩放，变为 $\left\lceil e\dfrac{q'}{q} \right\rfloor$ 和 $\left(\left\lceil a_0\dfrac{q'}{q}\right\rfloor, \left\lceil a_1\dfrac{q'}{q}\right\rfloor, \cdots, \left\lceil a_{k-1}\dfrac{q'}{q}\right\rfloor, \left\lceil b\dfrac{q'}{q}\right\rfloor\right)$。为切换 LWE 密文的模数，我们使用 \autoref{sec:modulus-rescaling} 介绍的模数重缩放技术。
相同的模数切换技术也适用于 RLWE 密文。本节将介绍如何切换（即重缩放）LWE 和 RLWE 密文的模数，并证明其正确性。
%That is, after  LWE/RLWE ciphertexts, we will prove why they are still valid LWE/RLWE ciphertexts.


\subsection{LWE 模数切换}
\label{subsec:modulus-switch-lwe}

\textbf{- 参考资料：}
\href{https://www.jeremykun.com/2022/07/16/modulus-switching-in-lwe/}{LWE 中的模数切换}

$ $

回顾 LWE 密码体制（\autoref{subsec:lwe-enc}），其组成如下：

\begin{itemize}
\item \textbf{\underline{参数设置}：} $\Delta = \dfrac{q}{t}$，\text{ } $S = (s_0, s_1, \ldots, s_{k-1}) \xleftarrow{\$} \mathbb{Z}_2^{k}$

\item \textbf{\underline{加密输入}：} $m \in \mathbb{Z}_t$，\text{ } $\vv{a} = (a_0, a_1, \ldots, a_{k-1}) \xleftarrow{\$} \mathbb{Z}_q^{k}$，\text{ } $e \xleftarrow{\chi_\sigma} \mathbb{Z}_q$

$ $

\item \textbf{\underline{加密}：} $\textsf{LWE}_{\vv{s},\sigma}(\Delta  m + e) = (\vv{a}, b) \in \mathbb{Z}_q^{k + 1}$ \text{ }（其中 $b = \vv{a} \cdot \vv{s} + \Delta  m + e \in \mathbb{Z}_q$）

$ $

\item \textbf{\underline{解密}：} $\textsf{LWE}^{-1}_{S,\sigma}(\textsf{ct}) = b - \vv{a}\cdot \vv{s} = \Delta  m + e \text{ } \text{ } \in \mathbb{Z}_q$

\end{itemize}

$ $


在 LWE 密码体制中，模数切换将原 LWE 密文的模数域转换成较小的模数域。可以将这一过程看作：除明文 $m$ 和私钥 $S$ 外，将原密文的所有分量缩小到更小的域中。该操作保持原明文 $m$ 的数值和完整性，同时将尺度因子 $\Delta$ 缩小为 $\hat{\Delta}$，将噪声 $e$ 缩小为 $\hat{e}$（注意，噪声的变化不影响原明文 $m$，因为噪声最终会在解密后的舍入过程中被消去），并将 $\vv{a}$ 缩小为 $\vv{\hat{a}}$。TFHE 在自举过程中使用模数切换来提高计算效率（详见 \autoref{subsec:tfhe-noise-bootstrapping}）。此外，BGV（\autoref{subsec:bfv-mult-cipher}）和 CKKS（\autoref{subsec:ckks-mult-cipher}）的密文与密文乘法算法也会使用模数切换。

LWE 模数切换的核心思想是对 LWE 方案中的同余关系进行重缩放。LWE 的同态运算包括密文与密文加法、密文与明文加法、密文与明文乘法以及密文与密文乘法。这些算法中使用的同余关系本质上都可由同一个基本关系改写得到：$b = \vv{a}\cdot \vv{s} + \Delta m + e \bmod q$。因此，将 LWE 密文的模数按 $q \rightarrow q'$ 切换，等价于将上述同余关系的模数按 $q \rightarrow q'$ 重缩放。

由此可知，在 LWE 密码体制中按 $q \rightarrow \hat{q}$（其中 $q > \hat{q}$）进行模数切换，就是按如下方式转换原 LWE 密文 $\textsf{LWE}_{S, \sigma}(\Delta m + e)$：

\begin{tcolorbox}[title={\textbf{\tboxlabel{\ref*{subsec:modulus-switch-lwe}} LWE 模数切换}}]

给定满足 $b = \vv{a}\cdot \vv{s} + \Delta m + e \bmod q$ 且 $m \in \mathbb{Z}_t$ 的 LWE 密文 $(\vv{a}, b)$，将其模数从 $q$ 切换到 $\hat q$，等价于将 $(A, b)$ 更新为 $(\vv{\hat{a}}, \hat b)$，具体如下：


$\vv{\hat{a}} = (\hat{a}_0, \hat{a}_1, \ldots, \hat{a}_{k-1})$，其中各分量满足 $\hat{a}_i = \Big\lceil a\dfrac{\hat{q}}{q} \Big\rfloor \in \mathbb{Z}_{\hat{q}}$ \textcolor{red}{ $\rhd$ $\lceil \rfloor$ 表示舍入到最近的整数}

 $\hat{b} = \Big\lceil b \dfrac{\hat{q}}{q} \Big\rfloor \in \mathbb{Z}_{\hat{q}}$

$\textsf{LWE}_{{S},\sigma}(\hat{\Delta}  m + \hat{e} + \epsilon_{\textit{all}}) = (\hat{a}_0, \hat{a}_1, \ldots, \hat{b}) \in \mathbb{Z}_{\hat{q}}^{k + 1}$ \textcolor{red}{ $\rhd$ 其中 $\epsilon_{\textit{all}}$ 是一个较小的舍入误差}

$ $


上述更新使 $\hat e$ 和 $\hat \Delta$ 发生如下变化：

$\hat{e} = \Big\lceil e\dfrac{\hat{q}}{q} \Big\rfloor \in \mathbb{Z}_{\hat{q}}$, \text{ }

$\hat{\Delta} = \Delta\dfrac{\hat{q}}{q}$ \textcolor{red}{ $\rhd$ 该值应为整数}

$ $


与此同时，$S$ 和 $m$ 保持不变。

\end{tcolorbox}

$ $

注意，要说明 $(\hat{a}_0, \hat{a}_1, \ldots, \hat{b}) \in \mathbb{Z}_{\hat{q}}^{k + 1}$ 是 $\hat{\Delta}m$ 的有效 LWE 密文，需要证明以下关系成立：

$\hat{b} = \sum\limits_{i=0}^{k-1}\hat{a}_i \cdot s_i + \hat{\Delta} m + \hat{e} \in \mathbb{Z}_{\hat{q}}$

$ $

%\textbf{\underline{Proof}}
\begin{myproof}
\begin{enumerate}
\item 注意以下关系：

$\hat{b} = \Big\lceil b \dfrac{\hat{q}}{q} \Big\rfloor = b\dfrac{\hat{q}}{q} + \epsilon_b$（其中 $-0.5 < \epsilon_b < 0.5$，表示舍入产生的偏差）

$\hat{a_i} = \Big\lceil a_i \dfrac{\hat{q}}{q} \Big\rfloor = a_i\dfrac{\hat{q}}{q} + \epsilon_{a_i}$（其中 $-0.5 < \epsilon_{a_i} < 0.5$）

$\hat{e} = \Big\lceil e \dfrac{\hat{q}}{q} \Big\rfloor = e\dfrac{\hat{q}}{q} + \epsilon_e$（其中 $-0.5 < \epsilon_e < 0.5$）
\item 注意以下关系：

$b = \vv{a} \cdot \vv{s} + \Delta  m + e = \sum\limits_{i=0}^{k-1}(a_is_i) + \Delta m + e  \in \mathbb{Z}_q$

$b = \sum\limits_{i=0}^{k-1}(a_is_i) + \Delta m + e + H \cdot q$（这里将模 $q$ 的关系改写为加上 $q$ 的某个倍数 $H \cdot q$）

\item 根据步骤 1 和步骤 2：

$\hat{b} = b\dfrac{\hat{q}}{q} + \epsilon_b \in \mathbb{Z}_{\hat{q}}$

$= \left(\sum\limits_{i=0}^{k-1}(a_is_i) + \Delta m + e + H \cdot q\right) \cdot \dfrac{\hat{q}}{q}  + \epsilon_b$

$= \dfrac{\hat{q}}{q} \cdot \sum\limits_{i=0}^{k-1}(a_is_i) + \dfrac{\hat{q}}{q} \cdot \Delta m + \dfrac{\hat{q}}{q} \cdot e + \dfrac{\hat{q}}{q} \cdot H \cdot q + \epsilon_b$

$= \sum\limits_{i=0}^{k-1}\left(\dfrac{\hat{q}}{q} \cdot a_is_i\right) + \hat{\Delta} m + (\hat{e} - \epsilon_e) + \hat{q}\cdot H + \epsilon_b$

$= \sum\limits_{i=0}^{k-1}\left((\hat{a}_i - \epsilon_{a_i}) \cdot s_i\right) + \hat{\Delta} m + \hat{e} - \epsilon_e + \hat{q}\cdot H + \epsilon_b$

$= \sum\limits_{i=0}^{k-1}(\hat{a}_is_i - \epsilon_{a_i}s_i) + \hat{\Delta} m + \hat{e} - \epsilon_e + \epsilon_b \in \mathbb{Z}_{\hat{q}}$

$= \sum\limits_{i=0}^{k-1}\hat{a}_is_i + \hat{\Delta} m + \left( \hat{e} - \epsilon_e + \epsilon_b - \sum\limits_{i=0}^{k-1}\epsilon_{a_i}s_i \right) \in \mathbb{Z}_{\hat{q}}$

$= \sum\limits_{i=0}^{k-1}\hat{a}_is_i + \hat{\Delta} m + \hat{e} + \epsilon_{\textit{all}} \in \mathbb{Z}_{\hat{q}}$ \textcolor{red}{ $\rhd$ 其中 $\epsilon_{\textit{all}} = \left(- \epsilon_e + \epsilon_b - \sum\limits_{i=0}^{k-1}\epsilon_{a_i}s_i \right)$}

$ $

$\epsilon_{\textit{all}}$ 的最大可能值为：

$\epsilon_{\textit{all}} = |-0.5| + |0.5| + |-0.5 \cdot k| = 1 + 0.5k$ 

因此，LWE 模数切换得到的是一个近似同余关系（\autoref{sec:modulus-rescaling}）。不过，只要 $\hat \Delta$ 足够大，$\epsilon_{\textit{all}} = 1 + 0.5k$ 就会在 LWE 解密的右移过程中被消去，最终恢复原明文 $m$。此外，在实际应用中，当 $k$ 足够大时，$\sum\limits_{i=0}^{k-1}\epsilon_{a_i}s_i$
相对于密文模数仍较小，因为每个 $a_i$ 和 $s_i$ 都由均匀采样产生。

\para{注意：} 如果 $\hat \Delta$ 不够大，$\epsilon_{\textit{all}}$ 就可能无法在解密时被消去，从而破坏明文 $m$。此外，如果 $\Delta \rightarrow \hat\Delta$ 缩小得过多，$\hat\Delta m$ 与 $\hat e$ 之间的距离会变得过小；此时 $\hat e = \Big\lceil e \dfrac{\hat{q}}{q} \Big\rfloor$ 的舍入过程可能影响 $\hat \Delta m$ 的最低有效位，从而破坏明文。

$ $

\item 综上，$\hat{b}$ 近似满足：

$\hat{b} = \sum\limits_{i=0}^{k-1}\hat{a}_is_i + \hat{\Delta} m + \hat{e} + \epsilon_{\textit{all}}  \approx \sum\limits_{i=0}^{k-1}\hat{a}_i  s_i + \hat{\Delta} m + \hat{e} \in \mathbb{Z}_{\hat{q}}$

因此，$(\hat{a}_0, \hat{a}_1, \ldots, \hat{b}) = \textsf{LWE}_{{S},\sigma}(\hat{\Delta} m + \hat{e} + \epsilon_{\textit{all}})$，对其解密即可得到 $m$。

%\begin{flushright}
%\qedsymbol{} 
%\end{flushright}

\end{enumerate}
\end{myproof}

\subsection{示例}

假设 LWE 的参数设置如下：

$ $

$t = 4$

$q = 64$

$n = 4$

$\Delta = \dfrac{q}{t} = 16$

$m = 1 \in \mathbb{Z}_t$

$S = (s_0, s_1, s_2, s_3) = (0, 1, 1, 0) \in \{-1, 0, 1\}^4$

$A = (a_0, a_1, a_2, a_3) = (-25, 12, -3, 7) \in \mathbb{Z}_q^4$

$e = 1 \in \mathbb{Z}_q$

$b = a_0s_0 + a_1s_1 + a_2s_2 + a_3s_3 + \Delta m + e = 26 \in \mathbb{Z}_q$

$\textsf{LWE}_{S, \sigma}(\Delta m + e) = \textsf{ct} = (a_0, a_1, a_2, a_3, b) = (-25, 12, -3, 7, 26) \in \mathbb{Z}_q^{n+1}$

$ $

现在希望将模数从 $q=64$ 切换到 $\hat{q} = 32$，可得：

$\hat{\Delta} = \Delta\cdot\dfrac{32}{64} = 8$

$\hat{e} = \Big\lceil 1 \cdot \dfrac{32}{64} \Big\rfloor = 1$

$\textsf{LWE}_{S, \sigma}(\hat{\Delta} m + \hat{e} + \epsilon_{\textit{all}}) = \hat{\textsf{ct}} = (\hat{a_0}, \hat{a_1}, \hat{a_2}, \hat{a_3}, \hat{b})$

$ = \left(\Big\lceil -25\cdot\dfrac{32}{64}\Big\rfloor, \Big\lceil12\cdot\dfrac{32}{64}\Big\rfloor, \Big\lceil-3\cdot\dfrac{32}{64}\Big\rfloor, \Big\lceil7\cdot\dfrac{32}{64}\Big\rfloor, \Big\lceil26\cdot\dfrac{32}{64}\Big\rfloor\right)$

$ = (-12, 6, -1, 4, 13) \in \mathbb{Z}_{\hat{q}}^{n+1}$

$ $

接下来验证以下 LWE 约束是否成立：

$\hat{b} = \hat{a}_0s_0 + \hat{a}_1s_1 + \hat{a}_2s_2 + \hat{a}_3s_3 + \hat{\Delta}m + \hat{e} \in \mathbb{Z}_{32}$

$13 = 0 + 6 - 1 + 0 + 8 \cdot 1 + 1 \in \mathbb{Z}_{32}$

$13 \approx 14 \in \mathbb{Z}_{32}$

此处出现了大小为 1 的微小差异，来源于以下舍入偏差：

$\hat {a_0} = \lceil -12.5 \rfloor = -12$、$\hat {a_2} = \lceil -1.5 \rfloor = -1$、$\hat{a_3} =  \lceil 3.5 \rfloor = 4$ 和 $\hat{e} = \lceil 0.5 \rfloor = 1 $

$ $

代入 LWE 解密公式计算：

$\hat b - (\hat{a}_0s_0 + \hat{a}_1s_1 + \hat{a}_2s_2 + \hat{a}_3s_3) = 13 - (0(-12) + 1(6) + 1(-1) + 0(4)) = 13 - (6 - 1) = 8 = \hat m + \hat e \in \mathbb{Z}_{32}$

$ $

$m = \left \lceil \dfrac{8}{\hat \Delta} \right \rfloor = \left \lceil \dfrac{8}{8} \right \rfloor = 1$，结果正确。


\subsection{讨论}
\label{subsubsec:modulus-switch-lwe-discuss}

\begin{figure}[h!]
    \centering
  \includegraphics[width=0.7\linewidth]{figures/modulus-switching.pdf}
  \caption{含噪声的缩放明文示意图：$\Delta \cdot m + e \in \mathbb{Z}_q$}
  \label{fig:modulus-switch}
\end{figure}

\para{$m$ 与 $e$ 之间的距离缩小：} 将 LWE 密文的模数按 $q \rightarrow \hat{q}$ 切换后，底层含噪明文 $\Delta m + e$ 缩小为 $\hat{\Delta}m + \hat{e}$，如 \autoref{fig:modulus-switch} 所示。注意，在按 $q \rightarrow \hat{q}$ 切换模数后，$\Delta m$ 缩小为 $\hat{\Delta} m$，但不会丢失明文比特信息。具体而言，明文值 $m$ 保持不变，而尺度因子 $\Delta$ 缩小为 $\hat{\Delta}$，噪声 $e$ 缩小为 $\hat e$。不过，与 $e$ 的最高有效位和 $\Delta m$ 的最低有效位之间的距离相比，模数切换后 $\hat e$ 的最高有效位与 $\hat \Delta m$ 的最低有效位之间的距离更小。




%\para{Hard Threshold on the $\Delta m$ Term's Value:} In \autoref{subsubsec:glwe-add-cipher-discuss}, we learned that the value of $\Delta m$ should not overflow or underflow its MSB $q$ during homomorphic computation (i.e., $\Delta m$ should be within the range of modulo $q$ without wrapping around), because if it wraps around the MSB, then our derived congruence relation of modulus switch loses its correctness. This is necessary in order to get the correct functionality of modulus switch on LWE ciphertexts. In particular, if we need to handle the case of $\Delta m$'s overflow, then our congruence relation of modulus switch should include the additional term $H_2 \cdot \Delta$ to take into account the added multiples of $\Delta$ caused by the overflows of $\Delta m$. But problematically, this term does not get eliminated during the modulus switch as $H_2 \cdot \Delta$ is not divisible by $\hat{q}$, which breaks our derived congruence relation above. Therefore, the applications of GLWE homomorphic encryption should ensure to manage their scaled plaintext $\Delta m$ to be always within than its maximum upper and minimum lower boundaries of the ciphertext modulus domain (i.e., $[0, q - 1]$ or $[-\dfrac{q}{2}, \dfrac{q}{2} - 1]$) throughout all computations. Meanwhile, the masking terms $A_i\cdot S_i$ are allowed to overflow or underflow ciphertext indefinitely during homomorphic computations, because during the decryption procedure they will be exactly subtracted and become zeroed out modulo $q$.  


\subsection{RLWE 模数切换}
\label{subsec:modulus-switch-rlwe}




RLWE 模数切换与 LWE 模数切换类似。回顾 RLWE 密码体制（\autoref{subsec:rlwe-enc}），其组成如下：

\begin{itemize}
\item \textbf{\underline{参数设置}：}
$\Delta = \dfrac{q}{t}$, \text{ } $S = s_0 + s_1X + s_2X^2 + \cdots + s_{n-1}X^{n-1} \xleftarrow{\$} \mathcal{R}_{\langle n, \textit{tern} \rangle}$

$ $

\item \textbf{\underline{加密输入}：}

$M \in \mathcal{R}_{\langle n, t \rangle}$

$A = a_0 + a_1X + a_2X^2 + \cdots + a_{n-1}X^{n-1} \xleftarrow{\$} \mathcal{R}_{\langle n, q \rangle}$

$E = e_0 + e_1X + e_2X^2 + \cdots + e_{n-1}X^{n-1} \xleftarrow{\chi_\sigma} \mathcal{R}_{\langle n, q \rangle}$ 

$ $

\item \textbf{\underline{加密}：}
$\textsf{RLWE}_{S,\sigma}(\Delta \cdot M + E) = (A, B) \in \mathcal{R}_{\langle n,q \rangle}^2$ 

其中 $B = A \cdot S + \Delta \cdot M + E = b_0 + b_1X + b_2X^2 + \cdots + b_{n-1}X^{n-1}$

$ $

\item \textbf{\underline{解密}：} $\textsf{RLWE}^{-1}_{S,\sigma}(\textsf{ct}) = B - A \cdot S = \Delta  M + E \text{ } \in \mathcal{R}_{\langle n,q \rangle}$
\end{itemize}

$ $

RLWE 模数切换过程如下：

\begin{tcolorbox}[title={\textbf{\tboxlabel{\ref*{subsec:modulus-switch-rlwe}} RLWE 模数切换}}]

对于满足 $B = AS + \Delta M + E$ 且 $M \in \mathcal{R}_{\langle n, q \rangle}$ 的 RLWE 密文 $(A, B)$，将模数从 $q$ 切换到 $\hat q$，等价于将 $(A, B)$ 更新为 $(\hat A, \hat B)$，具体如下：


$\hat{A} = \hat{a}_0 + \hat{a}_1X + \hat{a}_2X^2 + \cdots + \hat{a}_{n-1}X^{n-1}$，其中每个系数满足 $\hat{a}_i = \Big\lceil a_i\dfrac{\hat{q}}{q} \Big\rfloor \in \mathbb{Z}_{\hat{q}}$

$\hat{B} = \hat{b}_0 + \hat{b}_1X + \hat{b}_2X^2 + \cdots + \hat{b}_{n-1}X^{n-1}$，其中每个系数满足 $\hat{b}_i = \Big\lceil b_i\dfrac{\hat{q}}{q} \Big\rfloor \in \mathbb{Z}_{\hat{q}}$

$\textsf{RLWE}_{{S},\sigma}(\hat{\Delta}  M + \hat{E} + E^{\langle \epsilon_{\textit{all}} \rangle}) = (\hat{A}, \hat{B}) \in \mathcal{R}_{\langle n, \hat{q} \rangle}^{2}$ 

$ $

上述更新使 $\Delta$ 和 $E$ 发生如下变化：

$\hat{\Delta} = \Delta\dfrac{\hat{q}}{q}$ \textcolor{red}{ $\rhd$ 该值应为整数}

$\hat{E} = \hat{e}_0 + \hat{e}_1X + \hat{e}_2X^2 + \cdots + \hat{e}_{n-1}X^{n-1}$，其中每个系数满足 $\hat{e}_i = \Big\lceil e_i\dfrac{\hat{q}}{q} \Big\rfloor \in \mathbb{Z}_{\hat{q}}$

$ $

与此同时，$S$ 和 $M$ 保持不变。

\end{tcolorbox}

证明与 LWE 模数切换类似。

$ $

\textbf{\underline{证明}}

\begin{enumerate}
\item 注意以下关系：

$\hat{b}_i = \Big\lceil b_i \dfrac{\hat{q}}{q} \Big\rfloor = b_i\dfrac{\hat{q}}{q} + \epsilon_{b_i}$（其中 $-0.5 < \epsilon_{b_i} < 0.5$，表示舍入产生的偏差）

$\hat{a}_i = \Big\lceil a_i \dfrac{\hat{q}}{q} \Big\rfloor = a_i\dfrac{\hat{q}}{q} + \epsilon_{a_i}$（其中 $-0.5 < \epsilon_{a_i} < 0.5$）

$\hat{e}_i = \Big\lceil e_i \dfrac{\hat{q}}{q} \Big\rfloor = e_i\dfrac{\hat{q}}{q} + \epsilon_{e_i}$（其中 $-0.5 < \epsilon_{e_i} < 0.5$）

\item 注意以下关系：


$B - A\cdot S$ 

$ = (b_0 + b_1X + \cdots + b_{n-1}X^{n-1} ) - (a_{0} + a_{1}X + \cdots + a_{n-1}X^{n-1})(s_{0} + s_{1}X + \cdots + s_{n-1}X^{n-1})$ 


$ = \left(b_0 - \left( \sum\limits_{i=0}^{0}(a_{0-i}s_{i}) - \sum\limits_{i=1}^{n-1}(a_{n-i}s_{i}) \right)\right)$

$ + \left(b_1 - \left( \sum\limits_{i=0}^{1}(a_{1-i}s_{i}) - \sum\limits_{i=2}^{n-1}(a_{n+1-i}s_{i})   \right) \right)\cdot X$ 



$ + \left(b_2 - \left( \sum\limits_{i=0}^{2}(a_{2-i}s_{i}) - \sum\limits_{i=3}^{n-1}(a_{n+2-i}s_{i})   \right) \right)\cdot X^2$ 


$\vdots$ 


$ + \left(b_{n-1} - \left(  \sum\limits_{i=0}^{n-1}(a_{n-1-i}s_{i}) -  \sum\limits_{i=n}^{n-1}(a_{n+n-1-i}s_{i})  \right) \right)\cdot X^{n-1}$ \textcolor{red}{ $\rhd$ 按相同幂次合并各项}

$ $

$= \sum\limits_{h=0}^{n-1}  \left(b_h - \left( \sum\limits_{i=0}^{h}(a_{h-i}s_{i}) -  \sum\limits_{i=h+1}^{n-1}(a_{n+h-i}s_{i})  \right) \right)\cdot X^{h}  $

%$ $
$ $

因此，

$ $

$B = \sum\limits_{h=0}^{n-1}  b_h  X^{h}  $

$A\cdot S = \sum\limits_{h=0}^{n-1}  \left(\sum\limits_{i=0}^{h}(a_{h-i}s_{i}) - \sum\limits_{i=h+1}^{n-1}(a_{n+h-i}s_{i}) \right)\cdot X^{h}  $

%$= \sum\limits_{h=0}^{n-1}  C_h \cdot X^{h}  $, where $C_h = b_h - \left(  \sum\limits_{i=0}^{h}(a_{h-i}s_{i}) - \sum\limits_{i=h+1}^{n-1}(a_{n+h-i}s_{i})  \right)$

$ $




\item 根据步骤 2，

$B = A \cdot S + \Delta  M + E$ 


$\sum\limits_{h=0}^{n-1}  b_h  X^{h}  = \sum\limits_{h=0}^{n-1}  \left(\sum\limits_{i=0}^{h}(a_{h-i}s_{i}) - \sum\limits_{i=h+1}^{n-1}(a_{n+h-i}s_{i}) \right)\cdot X^{h}  + \Delta \sum\limits_{h=0}^{n-1}  m_h  X^{h} + \sum\limits_{h=0}^{n-1}  e_h  X^{h}  \in \mathbb{Z}_q$

$\sum\limits_{h=0}^{n-1}  b_h  X^{h}  = \sum\limits_{h=0}^{n-1}  \left(\sum\limits_{i=0}^{h}(a_{h-i}s_{i}) - \sum\limits_{i=h+1}^{n-1}(a_{n+h-i}s_{i}) \right)\cdot X^{h}  + \Delta \sum\limits_{h=0}^{n-1}  m_h  X^{h} + \sum\limits_{h=0}^{n-1}  e_h  X^{h} + H \cdot q$

（这里将模 $q$ 的关系改写为加上 $H \cdot q$；它是一个 $(n-1)$ 次多项式，其中每个系数 $c_i$ 都是 $q$ 的某个倍数。）

$ $

\item 根据步骤 1 和步骤 3，对于 $0 \leq j \leq n-1$ 中的每个 $j$：

$\hat{b}_j = b_j\dfrac{\hat{q}}{q} + \epsilon_{b_j} \in \mathbb{Z}_{\hat{q}}$

$= \left(\sum\limits_{i=0}^{j}(a_{j-i}s_{i}) - \sum\limits_{i=j+1}^{n-1}(a_{n+j-i}s_{i}) + \Delta m_j + e_j + c_j \cdot q\right) \cdot \dfrac{\hat{q}}{q} + \epsilon_{b_j}$

$= \dfrac{\hat{q}}{q} \sum\limits_{i=0}^{j}(a_{j-i}s_{i}) - \dfrac{\hat{q}}{q} \sum\limits_{i=j+1}^{n-1}(a_{n+j-i}s_{i}) + \dfrac{\hat{q}}{q} \cdot \Delta m_j + \dfrac{\hat{q}}{q} e_j + \dfrac{\hat{q}}{q} \cdot c_j \cdot q + \epsilon_{b_j}$

$=  \sum\limits_{i=0}^{j}\left(\dfrac{\hat{q}}{q}\cdot a_{j-i}s_{i}\right) - \sum\limits_{i=j+1}^{n-1}\left(\dfrac{\hat{q}}{q} \cdot a_{n+j-i}s_{i}\right) + \hat{\Delta} m_j + (\hat{e}_j - \epsilon_{e_j}) + \hat{q} \cdot c_j + \epsilon_{b_j}$

$=  \sum\limits_{i=0}^{j}((\hat{a}_{j-i}-\epsilon_{a_{j-i}})\cdot s_{i}) - \sum\limits_{i=j+1}^{n-1}((\hat{a}_{n+j-i}-\epsilon_{a_{n+j-i}})\cdot s_{i}) + \hat{\Delta} m_j + (\hat{e}_j - \epsilon_{e_j}) + \hat{q} \cdot c_j + \epsilon_{b_j}$

$=  \sum\limits_{i=0}^{j}(\hat{a}_{j-i}s_i) 
- \sum\limits_{i=j+1}^{n-1}(\hat{a}_{n+j-i}s_i)
- \sum\limits_{i=0}^{j}(\epsilon_{a_{j-i}} s_i)
+ \sum\limits_{i=j+1}^{n-1}(\epsilon_{a_{n+j-i}} s_i)
+ \hat{\Delta} m_j + (\hat{e}_j - \epsilon_{e_j}) + \hat{q} \cdot c_j + \epsilon_{b_j}$


$=  \left(\sum\limits_{i=0}^{j}(\hat{a}_{j-i}s_i) 
- \sum\limits_{i=j+1}^{n-1}(\hat{a}_{n+j-i}s_i)\right)
+ \hat{\Delta} m_j + \hat{e}_j 
+ \left(\epsilon_{b_j} - \epsilon_{e_j}
- \sum\limits_{i=0}^{j}(\epsilon_{a_{j-i}} s_i)
+ \sum\limits_{i=j+1}^{n-1}(\epsilon_{a_{n+j-i}} s_i)\right) + \hat{q} \cdot c_j
$

$=  \left(\sum\limits_{i=0}^{j}(\hat{a}_{j-i}s_i) 
- \sum\limits_{i=j+1}^{n-1}(\hat{a}_{n+j-i}s_i)\right)
+ \hat{\Delta} m_j + \hat{e}_j 
+ \epsilon_{\textit{all}} \in \mathbb{Z}_{\hat{q}}
$

\textcolor{red}{ $\rhd$ 其中 $\epsilon_{\textit{all}} = \epsilon_{b_j} - \epsilon_{e_j}
- \sum\limits_{i=0}^{j}(\epsilon_{a_{j-i}} s_i)
+ \sum\limits_{i=j+1}^{n-1}(\epsilon_{a_{n+j-i}} s_i) \approx 0$}


$ $

\item 综上，对于每个 $0 \leq j \leq n-1 $，多项式系数 $\hat{b_j}$ 近似满足：

$\hat{b}_j =  \left(\sum\limits_{i=0}^{j}(\hat{a}_{j-i}s_i) 
- \sum\limits_{i=j+1}^{n-1}(\hat{a}_{n+j-i}s_i)\right)
+ \hat{\Delta} m_j + \hat{e}_j 
+ \epsilon_{\textit{all}}$

$ \approx \left(\sum\limits_{i=0}^{j}(\hat{a}_{j-i}s_i) 
- \sum\limits_{i=j+1}^{n-1}(\hat{a}_{n+j-i}s_i)\right)
+ \hat{\Delta} m_j + \hat{e}_j$

因此，$(\hat{A}, \hat{B}) = \textsf{RLWE}_{{S},\sigma}(\hat{\Delta} M + \hat{E} + E^{\langle \epsilon_{\textit{all}} \rangle})$，对其解密即可得到 $M$。

\begin{flushright}
\qedsymbol{} 
\end{flushright}

\end{enumerate}




\subsection{GLWE 模数切换}
\label{subsec:modulus-switch-glwe}

GLWE 模数切换是 RLWE 模数切换的推广。唯一的区别在于：RLWE 的 $A$ 和 $S$ 各为一个多项式，而 GLWE 的 $A$ 和 $S$ 各为包含 $k$ 个多项式的列表。因此，可以对 GLWE 中的 $k$ 个多项式应用与 RLWE 相同的模数切换技术。

回顾 GLWE 密码体制（\autoref{subsec:glwe-enc}），其组成如下：

\begin{itemize}
\item \textbf{\underline{初始设置}：} $\Delta = \dfrac{q}{t}$，\text{ } $\{S_i\}_{i=0}^{k-1} \xleftarrow{\$} \mathcal{R}_{\langle n, \textit{tern} \rangle}^k$

$ $

\item \textbf{\underline{加密输入}：} $M \in \mathcal{R}_{\langle n, t \rangle}$，\text{ } $\{A_i\}_{i=0}^{k-1} \xleftarrow{\$} \mathcal{R}_{\langle n,q \rangle}^{k}$，\text{ } $E \xleftarrow{\chi_\sigma} \mathcal{R}_{\langle n,q \rangle}$

$ $

\item \textbf{\underline{加密}：} $\textsf{GLWE}_{S,\sigma}(\Delta M + E) = (\{A_i\}_{i=0}^{k-1}, B) \text{ } \in \mathcal{R}_{\langle n,q \rangle}^{k + 1}$

其中 $B = \sum\limits_{i=0}^{k-1}{(A_i \cdot S_i)} + \Delta  M + E \text{ } \in \mathcal{R}_{\langle n,q \rangle}$

$ $

\item \textbf{\underline{解密}：} $\textsf{GLWE}^{-1}_{S,\sigma}(\textsf{ct}) = B - \sum\limits_{i=0}^{k-1}{(A_i \cdot S_i)} = \Delta  M + E \text{ } \in \mathcal{R}_{\langle n,q \rangle}$

\end{itemize}

$ $

GLWE 模数切换过程如下：

\begin{tcolorbox}[title={\textbf{\tboxlabel{\ref*{subsec:modulus-switch-glwe}} GLWE 模数切换}}]

给定满足 $B = \vv A\cdot \vv S + \Delta M + E \bmod q$ 且 $M \in \mathcal{R}_{\langle n, q \rangle}$ 的 GLWE 密文 $(\{A_i\}_{i=0}^{k-1}, B)$，将其模数从 $q$ 切换到 $\hat q$，等价于将 $(\{A\}_{i=0}^{k-1}, B)$ 更新为 $(\{\hat{A}_i\}_{i=0}^{k-1}, \hat B)$，具体如下：


$\hat{A_i} = \hat{a}_{i,0} + \hat{a}_{i,1}X + \hat{a}_{i,2}X^2 + \cdots + \hat{a}_{i, {n-1}}X^{n-1}$，其中每个系数满足 $\hat{a}_{i,j} = \Big\lceil a_{i,j}\dfrac{\hat{q}}{q} \Big\rfloor \in \mathbb{Z}_{\hat{q}}$


$\hat{B} = \hat{b}_0 + \hat{b}_1X + \hat{b}_2X^2 + \cdots + \hat{b}_{n-1}X^{n-1}$，其中每个系数满足 $\hat{b}_j = \Big\lceil b_j\dfrac{\hat{q}}{q} \Big\rfloor \in \mathbb{Z}_{\hat{q}}$

$\textsf{GLWE}_{{S},\sigma}(\hat{\Delta}  M + \hat{E} + E^{\langle \epsilon_{\textit{all}} \rangle}) = (\{\hat{A}_i\}_{i=0}^{k-1}, \hat{B}) \in \mathcal{R}_{\langle n, \hat{q} \rangle}^{k+1}$ 

$ $

上述更新使 $E$ 和 $\Delta$ 发生如下变化：

$\hat{E} = \hat{e}_0 + \hat{e}_1X + \hat{e}_2X^2 + \cdots + \hat{e}_{n-1}X^{n-1}$，其中每个系数满足 $\hat{e}_j = \Big\lceil e_j\dfrac{\hat{q}}{q} \Big\rfloor \in \mathbb{Z}_{\hat{q}}$


$\hat{\Delta} = \Delta\dfrac{\hat{q}}{q}$ \textcolor{red}{ $\rhd$ 该值应为整数}

$ $


与此同时，$\vv{S}$ 和 $M$ 保持不变。



\end{tcolorbox}

证明与 RLWE 模数切换类似。对 GLWE 密文进行模数切换后，其同余关系（即 $B = \sum\limits_{i=0}^{k-1}A_i\cdot S_i + \Delta M + E$）中第 $j$ 个多项式系数的累计舍入误差如下：

$\epsilon_{j, all} = \epsilon_{b_j} - \epsilon_{e_j}
- \sum\limits_{l=0}^{k-1}\sum\limits_{i=0}^{j}(\epsilon_{a_{l,j-i}} \cdot s_{l,i})
+ \sum\limits_{l=0}^{k-1}\sum\limits_{i=j+1}^{n-1}(\epsilon_{a_{l,n+j-i}} \cdot s_{l,i})$

\textcolor{red}{ $\rhd$ 由总结~\ref*{subsec:modulus-switch-rlwe} 的证明步骤 4 推导得到：$\epsilon_{\textit{all}} = \epsilon_{b_j} - \epsilon_{e_j}
- \sum\limits_{i=0}^{j}(\epsilon_{a_{j-i}} s_i)
+ \sum\limits_{i=j+1}^{n-1}(\epsilon_{a_{n+j-i}} s_i)$}

$ $

注意，GLWE 模数切换产生的舍入误差可能比 RLWE 模数切换大约 $k$ 倍。不过，从长期来看，该误差相对于密文模数仍较小，因为各舍入误差独立且均匀分布，其和相对于模数增长较慢（中心极限定理）。


\clearpage


\section{GLWE 密钥切换}
\label{sec:glwe-key-switching}
下面的精确降幂 gadget 假设 $q=\beta^l$、整数 $\beta\ge2$，并采用 \autoref{subsec:poly-decomp} 的无符号系数数位。两个密钥使用相同的环和模数，但分量数量 $k$ 与 $k'$ 可以不同。密钥持有者预先计算旧秘密各分量在新密钥下的加密。解码消息的误差界假设 $t\mid q$。公开求值密钥还需要整体方案的安全论证，尤其是在密钥形成环路时；下面的算术证明并不是这种安全证明。

\textbf{- 参考资料：}
\href{https://www.zama.ai/post/tfhe-deep-dive-part-3}{深入理解 TFHE（三）：密钥切换与分层乘法}~\cite{tfhe-3}

$ $

密钥切换通过公开求值和预先生成的密钥切换密钥，将 GLWE 密文的秘密密钥从 $S$ 换为 $S'$，而不公开明文。它保持相位的消息项，但将 $E$ 换为一般不同的误差 $E'$。能否正确解码取决于新误差界；密钥切换并不是噪声刷新。

$ $

\noindent 回顾 $\textsf{GLWE}_{S,\sigma}(\Delta M + E) = (A_0, A_1, \ldots , A_{k-1}, B) \in \mathcal{R}_{\langle n, q \rangle }^{k + 1}$

其中 $B = \sum\limits_{i=0}^{k-1}{(A_i \cdot S_i)} + \Delta \cdot M + E$

私钥 $S$ 是包含 $k$ 个多项式的列表：$S = (S_0, S_1, \ldots, S_{k-1})$

$ $

\noindent 同时回顾 GLev（\autoref{sec:glev}）的定义：

$\textsf{GLev}_{S, \sigma}^{\beta, l}(M) = \Bigl \{\text{GLWE}_{S, \sigma}\Bigl (\dfrac{q}{\beta^i} \cdot M + E_i \Bigr ) \Bigr \}_{i=1}^{l} \in \mathcal{R}_{\langle n, q \rangle }^{(k+1) \cdot l}$

$ $

\noindent 现在，将 $k$ 个密钥切换密钥分别记为：

$\mathit{KSK}_i = \textsf{GLev}_{S', \sigma}^{\beta, l}(S_i)$ \\
$ = \left ( \textsf{GLWE}_{S', \sigma} \left ( 
 \dfrac{q}{\beta^1}S_i + E_{i, 1}\right ), \textsf{GLWE}_{S', \sigma} \left ( 
 \dfrac{q}{\beta^2}S_i + E_{i, 2}\right ), \ldots, \textsf{GLWE}_{S', \sigma} \left ( 
 \dfrac{q}{\beta^l}S_i  + E_{i, l} \right) \right ) \in \mathcal{R}_{\langle n, q \rangle }^{(k' + 1) \cdot l}$

$ $

\noindent 它是使用 $S'$ 对私钥 $S$ 进行 GLWE 加密得到的列表；新私钥 $S'$ 则是包含 $k'$ 个多项式的列表 $(S'_0, S'_1, \ldots, S'_{k'-1})$。于是，可以按以下方式完成 GLWE 密文的密钥切换 $S \rightarrow S'$：


\begin{tcolorbox}[title={\textbf{\tboxlabel{\ref*{sec:glwe-key-switching}} GLWE 密钥切换}}]

给定 $\textsf{GLWE}_{S,\sigma}(\Delta M + E) = (A, B)$，

$\textsf{GLWE}_{S',\sigma}(\Delta M + E') = (\overbrace{0, 0, \ldots, 0}^{k'}, B) - \sum \limits_{i=0}^{k-1} \langle \textsf{Decomp}^{\beta, l}(A_i), \mathit{KSK}_i \rangle$
\end{tcolorbox}



\begin{myproof}

\begin{enumerate}
\item 根据多项式分解原理（\autoref{subsec:poly-decomp}），给定多项式 $A_i \in \mathbb{Z}_q[x] / (x^n + 1)$，其分解如下：

$ $

$\textsf{Decomp}^{\beta, l}(A_i) = (A_{ i,1 }, A_{ i,2 }, \ldots, A_{ i,l })$，其中

$ $

$A_i = A_{ i, 1 } \dfrac{q}{\beta^1} + A _{ i, 2 }\dfrac{q}{\beta^2} + \cdots + A_{ i, l } \dfrac{q}{\beta^l}  $

$ $

\item $\langle \textsf{Decomp}^{\beta, l}(A_i), \mathit{KSK}_i \rangle $ \\
$ =  A_{ i,1 } \cdot \textsf{GLWE}_{S', \sigma} \left ( 
 \dfrac{q}{\beta^1}S_i + E_{i, 1} \right ) + A_{ i,2 } \cdot \textsf{GLWE}_{S', \sigma} \left ( 
 \dfrac{q}{\beta^2}S_i  + E_{i, 2} \right ) + \cdots + A_{ i,l } \cdot \textsf{GLWE}_{S', \sigma} \left ( 
 \dfrac{q}{\beta^l}S_i  + E_{i, l} \right ) $ 

\textcolor{red}{ $\rhd$ 各个 \textsf{GLWE} 密文分别加密明文 $S_i\dfrac{q}{\beta}, S_i\dfrac{q}{\beta^2}, \cdots, S_i\dfrac{q}{\beta^l}$，明文尺度因子均为 1}

$ $
 
$= \textsf{GLWE}_{S', \sigma} \left ( 
 \left ( A_{ i,1 } \dfrac{q}{\beta^1} + A_{ i,2 } \dfrac{q}{\beta^2} + \cdots + A_{ i,l } \dfrac{q}{\beta^l} \right ) \cdot S_i + E_i \right )$ \textcolor{red}{ $\rhd$ 其中 $E_i = \sum\limits_{j=1}^{l} A_{ i, j } E_{i, j}$}

 
$= \textsf{GLWE}_{S', \sigma} \left ( 
 \left ( A_i \right ) \cdot S_i + E_i \right ) = \textsf{GLWE}_{S', \sigma} ( 
 A_i  S_i  + E_i) $ 

$ $
 
 \item $\sum \limits_{i=0}^{k-1} \langle \textsf{Decomp}^{\beta, l}(A_i), \mathit{KSK}_i \rangle = \sum \limits_{i=0}^{k-1} (\textsf{GLWE}_{S', \sigma} ( 
 A_i  S_i  + E_i)) $ 
 
 $= \textsf{GLWE}_{S', \sigma} \left ( \sum \limits_{i=0}^{k-1}  
 (A_i  S_i   + E_i) \right ) $

$ $
 
\item $B$ 等价于在私钥 $S'$ 下的平凡 GLWE 加密，其中各个 $A_0, A_1, \ldots, A_{k'-1}$ 都为 0。也就是说，$\textsf{GLWE}_{S', \sigma}(B) = (\overbrace{0, 0, \ldots}^{k'}, B)$。注意，即使不知道 $S'$，也能构造 $\textsf{GLWE}_{S', \sigma}(B)$，因为所有 $A_iS_i$ 项均为 0。

$ $
 
 \item $\textsf{GLWE}_{S', \sigma}(B) - \textsf{GLWE}_{S', \sigma} \left( \sum \limits_{i=0}^{k-1} (A_i  S_i + E_i )\right) $
 
 $= \textsf{GLWE}_{S', \sigma}(B - \sum \limits_{i=0}^{k-1} (A_i  S_i + E_i)) $
 
$= \textsf{GLWE}_{S', \sigma}(\Delta M + E')$ \textcolor{red}{ $\rhd$ 其中 $E' = E - \sum\limits_{i=0}^{k-1}E_i$}

$ $

\item 将上述关系显式展开，可得：

$\textsf{GLWE}_{S', \sigma}(B) - \textsf{GLWE}_{S', \sigma} \left(\sum \limits_{i=0}^{k-1} (A_i  S_i + E_i)\right)$

$ = (\overbrace{0, 0, \ldots}^{k'}, B) - (\overbrace{A'_0, A'_1, \ldots, A'_{k'-1}}^{k'}, B')  $ \text { } \textcolor{red}{ $\rhd$ 其中 $B' = \sum\limits_{i=0}^{k'-1} A'_iS'_i + \sum\limits_{i=0}^{k-1} A_iS_i + \sum\limits_{i=0}^{k-1}E_i$}

$(\overbrace{-A'_0, -A'_1, \ldots, -A'_{k'-1}}^{k'}, \text{ } B - B')$

$ $

对上述密文进行解密可得：

$ \textsf{GLWE}_{S', \sigma}^{-1} \bm{\Big(} \text{ } (\overbrace{-A'_0, -A'_1, \ldots, -A'_{k'-1}}^{k'}, \text{ } B - B') \text{ } \bm{\Big)}$

$ = B -  B' - \sum\limits_{i=0}^{k'-1} -A'_iS_i'$

$= B - \sum\limits_{i=0}^{k'-1} A'_iS'_i - \sum\limits_{i=0}^{k-1} A_iS_i - \sum\limits_{i=0}^{k-1}E_i + \sum\limits_{i=0}^{k'-1} A'_iS_i'$

$= B - \sum\limits_{i=0}^{k-1} A_iS_i - \sum\limits_{i=0}^{k-1}E_i$

$= \Delta M + E - \sum\limits_{i=0}^{k-1}E_i $

$= \Delta M + E'$ \textcolor{red}{ $\rhd$ 其中 $E' = E - \sum\limits_{i=0}^{k-1}E_i$}

严格来说，$B - \sum\limits_{i=0}^{k-1} A_iS_i = \Delta M + E + Kq$，其中 $K$ 是一个多项式，以 $q$ 的倍数表示系数回绕的量。不过，上述所有运算均在模 $q$ 下进行，因此 $Kq$ 项会被消去。

$ $

\item 因此，$ (0, 0, \ldots , B) - \sum \limits_{i=0}^{k-1} \langle \textsf{Decomp}^{\beta, l}(A_i), \mathit{KSK}_i \rangle = \textsf{GLWE}_{S',\sigma}(\Delta M + E') $

$\textsf{GLWE}_{S',\sigma}(\Delta M + E')$ 是对明文 $\Delta M$ 的加密，其明文尺度因子为 1。不过，也可以将该密文重新解释为对明文 $M$ 的加密，其尺度因子为 $\Delta$。这样，无需任何额外计算，就能恢复密文原来的尺度因子 $\Delta$。
\end{enumerate}
\end{myproof}


\para{噪声与截断：}精确分解时，
\[
E'=E-\sum_{i=0}^{k-1}\sum_{j=1}^l A_{i,j}E_{i,j},\qquad
\|E'\|_\infty\le\|E\|_\infty+\sum_{i,j}\|A_{i,j}\|_1\|E_{i,j}\|_\infty.
\]
若所有密钥切换误差的范数至多为 $B_{\mathrm{ks}}$，无符号数位给出的上界为 $\|E'\|_\infty\le\|E\|_\infty+knl(\beta-1)B_{\mathrm{ks}}$。为了保证解码，它必须仍小于 $\Delta/2$。代数上的相位恒等式并不意味着新密文具有原先的误差分布。

对于权重为整数的截断降幂 gadget，写成 $A_i\equiv\sum_jA_{i,j}(q/\beta^j)+\rho_i\pmod q$，同一公开算法得到
\[
E'=E-\sum_{i,j}A_{i,j}E_{i,j}+\sum_i\rho_iS_i,
\]
因此上界中还必须包含 $\sum_i\|\rho_i\|_\infty\|S_i\|_1$。省略这一残差会使证明失效。一般 $q$ 也可采用精确整数 gadget，但必须在密文数位与密钥切换密钥中一致地定义其权重和分解。




\clearpage


%\section{GLWE Input Vector Rotation}
%\label{sec:glwe-input-rotation}
%\input{zh-CN/c06-glwe-input-rotation}
%\clearpage


\part{全同态加密方案}
\label{part:fhe-schemes}

\renewcommand{\thesection}{D-\arabic{section}}
\setcounter{section}{0}

本部分介绍四种最常见的 FHE 方案：TFHE、CKKS、BGV 和 BFV，以及它们的 RNS 变体。

\newpage

\section{TFHE 方案}
\label{sec:tfhe}

TFHE 方案用于整数的同态加法与乘法，尤其适合逻辑电路这样的按位计算。与 BFV、BGV 或 CKKS 不同，TFHE 的特点是噪声自举速度快，因此能够高效处理乘法深度较大的计算。TFHE 的噪声自举技术还可以进一步用于函数加密。

在 TFHE 中，每个明文都被加密为一个 LWE 密文。因此，TFHE 的密文与密文加法、密文与明文加法，以及密文与明文乘法，均基于 \autoref{part:generic-fhe} 中介绍的 GLWE 同态加法和乘法实现；取 $n = 1$ 即可将 GLWE 化为 LWE。

本节介绍 TFHE 的特有组成部分：密钥切换、密文与密文乘法、系数提取，以及噪声自举。

$ $

\begin{tcolorbox}[
    title = \textbf{所需背景知识},    % box title
    colback = white,    % light background; tweak to taste
    colframe = black,  % frame colour
    boxrule = 0.8pt,     % line thickness
    left = 1mm, right = 1mm, top = 1mm, bottom = 1mm % inner padding
]
\begin{itemize}
  \item \autoref{sec:modulo}: \nameref{sec:modulo}
  \item \autoref{sec:group}: \nameref{sec:group}
  \item \autoref{sec:field}: \nameref{sec:field}
  \item \autoref{sec:order}: \nameref{sec:order}
  \item \autoref{sec:polynomial-ring}: \nameref{sec:polynomial-ring}
  \item \autoref{sec:decomp}: \nameref{sec:decomp}
  \item \autoref{sec:modulus-rescaling}: \nameref{sec:modulus-rescaling}
  \item \autoref{sec:lattice}: \nameref{sec:lattice}
  \item \autoref{sec:lwe}: \nameref{sec:lwe}
  \item \autoref{sec:rlwe}: \nameref{sec:rlwe}
  \item \autoref{sec:glwe}: \nameref{sec:glwe}
  \item \autoref{sec:glev}: \nameref{sec:glev}
  \item \autoref{sec:ggsw}: \nameref{sec:ggsw}
  \item \autoref{sec:glwe-add-cipher}: \nameref{sec:glwe-add-cipher}
  \item \autoref{sec:glwe-add-plain}: \nameref{sec:glwe-add-plain}
  \item \autoref{sec:glwe-mult-plain}: \nameref{sec:glwe-mult-plain}
  \item \autoref{subsec:modulus-switch-lwe}: \nameref{subsec:modulus-switch-lwe}
  \item \autoref{sec:glwe-key-switching}: \nameref{sec:glwe-key-switching}
\end{itemize}
\end{tcolorbox}

\clearpage

\subsection{加密与解密}
\label{subsec:tfhe-enc-dec}

TFHE 基于 LWE 密码系统（\autoref{sec:lwe}）进行加密和解密，它等价于取 $n = 1$ 的 GLWE 密码系统（\autoref{sec:glwe}）。不过，与前面介绍的 LWE 密码系统不同，TFHE 从二元集合 $\{0, 1\}$ 中采样秘密密钥的元素，而不是从三元集合 $\{-1, 0, 1\}$ 中采样。


\begin{tcolorbox}[title={\textbf{\tboxlabel{\ref*{subsec:tfhe-enc-dec}} TFHE 加密与解密}}]
\textbf{\underline{初始设置}：} $\Delta = \dfrac{q}{t}$, $\vec{s} \xleftarrow{\$} \mathbb{Z}_2^k$  \textcolor{red}{ $\rhd$ 其中 $t$ 整除 $q$，且 $\vec{s}$ 的每个元素都是零次多项式}

$ $

\par\noindent\rule{\textwidth}{0.4pt}

\textbf{\underline{加密输入}：} $m \in \mathbb{Z}_t$, $\vec{a} \xleftarrow{\$} \mathbb{Z}_q^k$, $e \xleftarrow{\chi_\sigma} \mathbb{Z}_q$ \textcolor{red}{ $\rhd$ $\vec{a}$ 的每个元素都是零次多项式}
\begin{enumerate}
\item 放大尺度：$m \longrightarrow \Delta \cdot m \text{ } \in \mathbb{Z}_q$

\item 计算 $b = \vec{a} \cdot \vec{s} + \Delta  m + e \text{ } \bmod q$
\item $\textsf{LWE}_{\vec{s},\sigma}(\Delta  m + e) = (\vec{a}, b) \text{ } \in \mathbb{Z}_q^{k + 1}$ 
\end{enumerate}

\par\noindent\rule{\textwidth}{0.4pt}

\textbf{\underline{解密输入}：} $\textsf{ct} = (\vec{a}, b) \text{ } \in \mathbb{Z}_q^{k+1}$
\begin{enumerate}
\item $\textsf{LWE}^{-1}_{\vec{s},\sigma}(\textsf{ct}) = b - \vec{a}\cdot \vec{s} = \Delta  m + e  \pmod q$

\item 缩小尺度：$\Bigg\lceil\dfrac{ \Delta  m + e } {\Delta}\Bigg\rfloor = m \text{ } \in \mathbb{Z}_t$ \textcolor{red}{ $\rhd$ 即执行 $q \rightarrow t$ 的模数切换}
\end{enumerate}

$ $

\textbf{{正确解密的条件}：}
\begin{itemize}
\item 同态运算过程中增长的噪声 $e$ 应满足：$e < \dfrac{\Delta}{2}$。
\end{itemize}

\end{tcolorbox}

为简化记号，本节经常将 $\textsf{LWE}_{\vec{s},\sigma}(\Delta  m + e)$ 写成 $\textsf{LWE}_{\vec{s},\sigma}(\Delta  m)$，因为 $\textsf{LWE}_{\vec{s},\sigma}(\Delta  m + e) \approx \textsf{LWE}_{\vec{s},\sigma}(\Delta  m)$（即二者解密得到同一消息）。即使我们用 $\textsf{LWE}_{\vec{s},\sigma}(\Delta  m)$ 代替 $\textsf{LWE}_{\vec{s},\sigma}(\Delta  m + e)$，也应将其理解为对 $\Delta m + e$ 的加密，即放大后的消息中仍然包含噪声。

\subsection{密文与密文的同态加法}
\label{subsec:tfhe-add-cipher}

TFHE 的密文与密文加法采用 LWE 的密文加法方案，它等价于取 $n = 1$ 的 GLWE 密文加法方案（\autoref{sec:glwe-add-cipher}）。


\begin{tcolorbox}[title={\textbf{\tboxlabel{\ref*{subsec:tfhe-add-cipher}} TFHE 密文与密文加法}}]
$\textsf{LWE}_{\vec{s}, \sigma}(\Delta m^{\langle 1 \rangle} + e^{\langle 1 \rangle} ) + \textsf{LWE}_{\vec{s}, \sigma}(\Delta m^{\langle 2 \rangle} + e^{\langle 2 \rangle}) $

$ = ( \vec{a}^{\langle 1 \rangle}, \text{ } b^{\langle 1 \rangle}) + (\vec{a}^{\langle 2 \rangle}, \text{ } b^{\langle 2 \rangle}) $

$ = ( \vec{a}^{\langle 1 \rangle} + \vec{a}^{\langle 2 \rangle}, \text{ } b^{\langle 1 \rangle} + b^{\langle 2 \rangle} ) $

$= \textsf{LWE}_{\vec{s}, \sigma}(\Delta(m^{\langle 1 \rangle} + m^{\langle 2 \rangle}) + e^{\langle 1 + 2 \rangle} )$
\label{Here}
\end{tcolorbox}


\subsection{密文与明文的同态加法}
\label{subsec:tfhe-add-plain}

TFHE 的密文与明文加法（其中 $\lambda$ 是待加的常数）采用 LWE 的密文与明文加法方案，它等价于取 $n = 1$ 的 GLWE 密文与明文加法方案（\autoref{sec:glwe-add-plain}）。

\begin{tcolorbox}[title={\textbf{\tboxlabel{\ref*{subsec:tfhe-add-plain}} TFHE 密文与明文加法}}]
$\textsf{LWE}_{\vec{s}, \sigma}(\Delta m + e) + \Delta \lambda $

$=  (\vec{a}, \text{ } b) + \Delta \lambda$

$=  (\vec{a}, \text{ } b + \Delta\lambda)$

$= \textsf{LWE}_{\vec{s}, \sigma}(\Delta (m + \lambda) + e )$
\end{tcolorbox}


\subsection{密文与明文的同态乘法}
\label{subsec:tfhe-mult-plain}

TFHE 的密文与明文乘法采用 LWE 的密文与明文乘法方案，它等价于取 $n = 1$ 的 GLWE 密文与明文乘法方案（\autoref{sec:glwe-mult-plain}）。

\begin{tcolorbox}[title={\textbf{\tboxlabel{\ref*{subsec:tfhe-mult-plain}} TFHE 密文与明文乘法}}]
$\textsf{LWE}_{\vec{s}, \sigma}(\Delta m + e) \cdot \lambda$

$= (\vec{a}, \text{ } b) \cdot \lambda$

$= (\lambda\cdot \vec{a}, \text{ } \lambda \cdot b)$

$= \textsf{LWE}_{\vec{s}, \sigma}(\Delta (m \cdot \lambda) + \lambda \cdot e )$
\end{tcolorbox}


\subsection{同态密钥切换}
\label{subsec:tfhe-key-switching}

\textbf{- 参考资料：}
\href{https://www.zama.ai/post/tfhe-deep-dive-part-3}{深入理解 TFHE：第三部分——密钥切换与分层乘法}~\cite{tfhe-3}

TFHE 的密钥切换方案将 LWE 密文对应的秘密密钥从 $\vec{s}$ 变为 $\vec{s}_{'}$，两个密钥向量的维数可以相同，也可以不同。该方案本质上就是 LWE 密钥切换。具体而言，它等价于取 $n = 1$ 的另一种 GLWE 表示（\autoref{subsec:glwe-alternative}）所采用的密钥切换方案（\autoref{sec:glwe-key-switching}），如下所示：

\begin{tcolorbox}[title={\textbf{\tboxlabel{\ref*{subsec:tfhe-key-switching}} TFHE 密钥切换}}]

给定 $\textsf{LWE}_{\vec{s},\sigma}(\Delta m + e) = (\vec{a}, b)$，

$\textsf{LWE}_{\vec{s}_{'},\sigma}(\Delta m + e') = (0, b) - \bm{\langle} \textsf{Decomp}^{\beta, l}(\vec{a}), \text{ } \textsf{Lev}_{\vec{s}_{'}, \sigma}^{\beta, l}(\vec{s}) \bm{\rangle}$
\end{tcolorbox}


\subsection{密文与密文的同态乘法}
\label{subsec:tfhe-mult-cipher}

\textbf{- 参考资料：}
\href{https://www.zama.ai/post/tfhe-deep-dive-part-3}{深入理解 TFHE：第三部分——密钥切换与分层乘法}~\cite{tfhe-3}

$ $

TFHE 支持以下形式的两个密文相乘：$\textsf{LWE}_{\vec{s}, \sigma}(\Delta m_1) \cdot \textsf{GSW}_{\vec{s}, \sigma}^{\beta, l}(m_2)$。

$ $

\noindent 第一项 $\textsf{LWE}_{\vec{s}, \sigma}(\Delta m_1 + e_1)$ 可由以下任一种方式得到：
\begin{itemize}
\item 对明文 $m_1$ 进行一次新的 LWE 加密（\autoref{subsec:glwe-enc}）。
\item 对两个 LWE 密文执行同态加法所得的结果（\autoref{sec:glwe-add-cipher}）。
\item 对一个 LWE 密文和一个明文执行同态乘法所得的结果（\autoref{sec:glwe-mult-plain}）。
\end{itemize}

$ $

\noindent 第二项 $\textsf{GSW}_{\vec{s}, \sigma}^{\beta, l}(m_2)$ 可由以下任一种方式得到：
\begin{itemize}
\item 对明文 $m_2$ 进行一次新的 GSW 加密（\autoref{subsec:ggsw-enc}）。
\item 通过\textit{电路自举}将 $\textsf{LWE}_{\vec{s}, \sigma}(\Delta m_2 + e_2)$ 转换为 $\textsf{GSW}_{\vec{s}, \sigma}^{\beta, l}(m_2)$（将在未来版本中介绍）。
\end{itemize}

$ $

\noindent 回顾以下关系：

$\textsf{LWE}_{\vec{s}, \sigma}(\Delta m_1 + e_1) = (\vec{a}, b) \in \mathbb{Z}_{q}^{k + 1}$，其中 $b = \vec{a} \cdot \vec{s} + \Delta m_1 + e_1$

$ $

$\textsf{GSW}_{\vec{s}, \sigma}^{\beta, l}(m_2) = \Bigr \{ \{ \textsf{Lev}_{\vec{s}, \sigma}^{\beta, l} (-s_i \cdot m_2)  \}_{i=0}^{k-1}, \textsf{Lev}_{\vec{s}, \sigma}^{\beta, l}(m_2) \Bigl \} \in \mathbb{Z}_{q }^{(k+1) \cdot l \cdot (k'+1)}$ \textcolor{red}{ $\rhd$ 由 \autoref{subsec:ggsw-enc} 可得}

$ $

\noindent 使用以下记号：

$\textsf{GSW}_{\vec{s}, \sigma}^{\beta, l}(m_2) = {\bar{\textsf{ct}}} = (\bar{\textsf{ct}}_0,. \bar{\textsf{ct}}_1, \gap{$\cdots$} \bar{\textsf{ct}}_k)$ 

$\bar{\textsf{ct}}_i = \textsf{Lev}_{\vec{s}, \sigma}^{\beta, l}(-s_i \cdot m_2)$，其中 $0 \leq i \leq (k-1)$

$\bar{\textsf{ct}}_k = \textsf{Lev}_{\vec{s}, \sigma}^{\beta, l}(m_2)$

$\textsf{ct} = \textsf{LWE}_{\vec{s}, \sigma}(\Delta m_1 + e_1) = (\vec{a}, b) = (a_0, a_1, \gap{$\cdots$}, a_{k-1}, b) = (\textsf{ct}_0, \textsf{ct}_1, \cdots, \textsf{ct}_k)$

$ $

\noindent 定义以下 TFHE 密文乘法运算：

$\textsf{ct} \cdot {\bar{\textsf{ct}}} = \sum\limits_{i=0}^{k}\langle \textsf{Decomp}^{\beta, l}(\textsf{ct}_i), \bar{\textsf{ct}}_i \rangle$

$ $

\noindent 于是下式成立：

\begin{tcolorbox}[title={\textbf{\tboxlabel{\ref*{subsec:tfhe-mult-cipher}} TFHE 密文与密文乘法}}]
$\textsf{ct} = \textsf{LWE}_{\vec{s}, \sigma}(\Delta m_1 + e_1) = (a_0, a_1, \cdots, a_{k-1}, b)$

$\bar{\textsf{ct}} = \textsf{GSW}_{\vec{s}, \sigma}^{\beta, l}(m_2) = \bm( \textsf{Lev}_{\vec{s}, \sigma}^{\beta, l}(-s_0\cdot m_2), \textsf{Lev}_{\vec{s}, \sigma}^{\beta, l}(-s_1\cdot m_2), \cdots, \textsf{Lev}_{\vec{s}, \sigma}^{\beta, l}(-s_{k-1}\cdot m_2), \textsf{Lev}_{\vec{s}, \sigma}^{\beta, l}(m_2)  \bm)$

$\textsf{LWE}_{\vec{s}, \sigma}(\Delta m_1 + e_1) \cdot \textsf{GSW}_{\vec{s}, \sigma}^{\beta, l}(m_2) $

$= \sum\limits_{i=0}^{k}\langle \textsf{Decomp}^{\beta, l}(\textsf{ct}_i), \bar{\textsf{ct}}_i \rangle $

$\approx \textsf{LWE}_{\vec{s}, \sigma}(\Delta m_1 m_2)$
\end{tcolorbox}

这意味着，将两个 TFHE 密文（一个为 LWE，另一个为 GSW）相乘，并解密所得的 LWE 密文，结果与将两个原始明文相乘相同。


%\textbf{\underline{Proof}}
\begin{myproof}
\begin{enumerate}
\item $\sum\limits_{i=0}^{k}\langle \textsf{Decomp}^{\beta, l}(\textsf{ct}_i), \bar{\textsf{ct}}_i \rangle$ \\
$= \langle \textsf{Decomp}^{\beta, l}(a_0), \bar{\textsf{ct}}_0 \rangle + \langle \textsf{Decomp}^{\beta, l}(a_1), \bar{\textsf{ct}}_1 \rangle + \gap{$\cdots$} + \langle \textsf{Decomp}^{\beta, l}(a_{k-1}), \bar{\textsf{ct}}_{k-1} \rangle + \langle \textsf{Decomp}^{\beta, l}(b), \bar{\textsf{ct}}_k \rangle$ \\
\textcolor{red}{ $\rhd$ 展开两个向量的点积}
\item 对于 $i = k$： \\
$\textsf{Decomp}^{\beta, l}(b) = (b_1, b_2, \ldots , b_l)$，其中 $b = b_1\dfrac{q}{\beta^1} + b_2\dfrac{q}{\beta^2} + \cdots + b_l\dfrac{q}{\beta^l}$ \textcolor{red}{ $\rhd$ 由 \autoref{subsec:poly-decomp} 可得}\\
$\bar{\textsf{ct}}_k = \textsf{Lev}_{\vec{s}, \sigma}^{\beta, l}(m_2) = \left(\textsf{LWE}_{\vec{s}, \sigma}\left(m_{2}\dfrac{q}{\beta^1} + e_{2, 1}\right), \textsf{LWE}_{\vec{s}, \sigma}\left(m_{2}\dfrac{q}{\beta^2} + e_{2, 2}\right), \ldots, \textsf{LWE}_{\vec{s}, \sigma}\left(m_{2}\dfrac{q}{\beta^l} + e_{2, l}\right) \right)$ \\
$ $

因此： \\
$\langle \textsf{Decomp}^{\beta, l}(b), \bar{\textsf{ct}}_k \rangle$\\
$= b_1 \cdot \textsf{LWE}_{\vec{s}, \sigma} \left (m_{2}\dfrac{q}{\beta^1}  + e_{2, 1}\right ) + b_2 \cdot \textsf{LWE}_{\vec{s}, \sigma} \left (m_{2}\dfrac{q}{\beta^2}  + e_{2, 2}\right ) + \gap{$\cdots$} + b_l \cdot \textsf{LWE}_{\vec{s}, \sigma} \left (m_{2}\dfrac{q}{\beta^l}  + e_{2, l}\right)$ \\
$= \textsf{LWE}_{\vec{s}, \sigma} \left (b_1m_{2}\dfrac{q}{\beta^1}  + b_1e_{2, 1}\right ) + \textsf{LWE}_{\vec{s}, \sigma} \left (b_2m_{2}\dfrac{q}{\beta^2}  + b_2e_{2,2}\right ) + \gap{$\cdots$} + \textsf{LWE}_{\vec{s}, \sigma} \left (b_lm_{2}\dfrac{q}{\beta^l} + b_le_{2, l}\right)$ \textcolor{red}{ $\rhd$ 由 \autoref{sec:glwe-mult-plain} 可得} \\
$= \textsf{LWE}_{\vec{s}, \sigma} \left (b_1m_{2}\dfrac{q}{\beta^1} + b_2m_{2}\dfrac{q}{\beta^2} + \gap{$\cdots$} + b_lm_{2}\dfrac{q}{\beta^l} + b_1e_{2,1} + \cdots + b_le_{2,l}\right)$ \textcolor{red}{ $\rhd$ 由 \autoref{sec:glwe-add-cipher} 可得}\\
$= \textsf{LWE}_{\vec{s}, \sigma} \left (m_{2} \cdot \left ( b_1\dfrac{q}{\beta^1} + b_2\dfrac{q}{\beta^2} + \gap{$\cdots$} + b_l\dfrac{q}{\beta^l} + e^{\langle k \rangle} \right)\right)$ \textcolor{red}{ $\rhd$ 其中 $e^{\langle k \rangle} = \sum\limits_{i=1}^{l}b_ie_{2, i}$}\\
$= \textsf{LWE}_{\vec{s}, \sigma} (m_{2}b + e^{\langle k \rangle})$ \textcolor{red}{ $\rhd$ 由 \autoref{subsec:poly-decomp} 可得}

\item 对于 $0 \leq i \leq (k - 1)$： \\
$\textsf{Decomp}^{\beta, l}(a_i) = (a_{\langle i, 1 \rangle}, a_{\langle i, 2 \rangle}, \gap{$\cdots$}, a_{\langle i, l \rangle})$，其中 $a_i = a_{\langle i, 1 \rangle}\dfrac{q}{\beta^1} + a_{\langle i, 2 \rangle}\dfrac{q}{\beta^2} + \gap{$\cdots$} + a_{\langle i, l \rangle}\dfrac{q}{\beta^l}$ \\
$\bar{\textsf{ct}}_i =  \textsf{Lev}_{\vec{s}, \sigma}^{\beta, l}(-s_im_2) $

$= \left(\textsf{LWE}_{\vec{s}, \sigma}\left(-s_im_2\dfrac{q}{\beta^1} + e_{i, 1}\right), \textsf{LWE}_{\vec{s}, \sigma}\left(-s_im_2\dfrac{q}{\beta^2} + e_{i, 2}\right), \gap{$\cdots$}, \textsf{LWE}_{\vec{s}, \sigma}\left(-s_im_2\dfrac{q}{\beta^l} + e_{i, l}\right) \right)$ \\

$ $
因此： \\
$\langle \textsf{Decomp}^{\beta, l}(a_0), \bar{\textsf{ct}}_0 \rangle + \langle \textsf{Decomp}^{\beta, l}(a_1), \bar{\textsf{ct}}_1 \rangle + \gap{$\cdots$} + \langle \textsf{Decomp}^{\beta, l}(a_{k-1}), \bar{\textsf{ct}}_{k-1} \rangle$ \\
$= \sum\limits_{i=0}^{k-1}\langle \textsf{Decomp}^{\beta, l}(a_i), \bar{\textsf{ct}}_i \rangle$ \\
$= \sum\limits_{i=0}^{k-1}\Bigg(a_{\langle i, 1\rangle} \cdot \textsf{LWE}_{\vec{s}, \sigma}\left(-s_im_2\dfrac{q}{\beta^1} + e_{i,1}\right) + a_{\langle i, 2\rangle} \cdot \textsf{LWE}_{\vec{s}, \sigma}\left(-s_im_2\dfrac{q}{\beta^2} + e_{i,2}\right) + $

$ \cdots +  a_{\langle i, l\rangle} \cdot \textsf{LWE}_{\vec{s}, \sigma}\left(-s_im_2\dfrac{q}{\beta^l} + e_{i,l}\right)\Bigg)$ \\
$= \sum\limits_{i=0}^{k-1}\Bigg(\textsf{LWE}_{\vec{s}, \sigma}\left(-a_{\langle i, 1\rangle}s_im_2\dfrac{q}{\beta^1} + a_{\langle i, 1 \rangle}e_{i,1}\right) + \textsf{LWE}_{\vec{s}, \sigma}\left(-a_{\langle i, 2\rangle}s_im_2\dfrac{q}{\beta^2} + a_{\langle i, 2 \rangle}e_{i,2}\right) + $

$\gap{$\cdots$} + \textsf{LWE}_{\vec{s}, \sigma}\left(-a_{\langle i, l\rangle}s_im_2\dfrac{q}{\beta^l} + a_{\langle i, l \rangle}e_{i,l}\right)\Bigg)$ \\
$= \sum\limits_{i=0}^{k-1}\textsf{LWE}_{\vec{s}, \sigma}\left(-a_{\langle i, 1\rangle}s_im_2\dfrac{q}{\beta^1} -a_{\langle i, 2\rangle}s_im_2\dfrac{q}{\beta^2} + \gap{$\cdots$} -a_{\langle i, l\rangle}s_im_2\dfrac{q}{\beta^l} + a_{\langle i, 1 \rangle}e_{i,1} + \cdots + a_{\langle i, l \rangle}e_{i,l}\right)$ \\
$= \sum\limits_{i=0}^{k-1}\textsf{LWE}_{\vec{s}, \sigma}\left(-s_im_2 \cdot \left(a_{\langle i, 1\rangle}\dfrac{q}{\beta^1} + a_{\langle i, 2\rangle}\dfrac{q}{\beta^2} + \gap{$\cdots$} + a_{\langle i, l\rangle}\dfrac{q}{\beta^l}\right) + a_{\langle i, 1 \rangle}e_{i,1} + \cdots + a_{\langle i, l \rangle}e_{i,l}\right)$ \\
$= \sum\limits_{i=0}^{k-1}\textsf{LWE}_{\vec{s}, \sigma}(-s_im_2a_i + e^{\langle i \rangle})$ \textcolor{red}{ $\rhd$ 其中 $e^{\langle i \rangle} = a_{\langle i, 1 \rangle}e_{i,1} + \cdots + a_{\langle i, l \rangle}e_{i,l}$}
\item 根据步骤 2 和 3， \\
$\sum\limits_{i=0}^{k}\langle \textsf{Decomp}^{\beta, l}(\textsf{ct}_i), \bar{\textsf{ct}}_i \rangle$ \\
$= \sum\limits_{i=0}^{k-1}\textsf{LWE}_{\vec{s}, \sigma}(-s_im_2a_i + e^{\langle i \rangle}) + \textsf{LWE}_{\vec{s}, \sigma} (m_{2}b + e^{\langle k \rangle})$ \\ 
$= \textsf{LWE}_{\vec{s}, \sigma}\left(\sum\limits_{i=0}^{k-1}(-s_im_2a_i) + m_{2}b + \sum\limits_{i=0}^ke^{\langle i \rangle}\right)$ \textcolor{red}{ $\rhd$ 两个 LWE 密文相加} \\
$= \textsf{LWE}_{\vec{s}, \sigma}\left(m_2b - \sum\limits_{i=0}^{k-1}m_2a_is_i + \sum\limits_{i=0}^ke^{\langle i \rangle}\right)$  \\ 
$= \textsf{LWE}_{\vec{s}, \sigma}\left(m_2(b - \sum\limits_{i=0}^{k-1}a_is_i) + \sum\limits_{i=0}^ke^{\langle i \rangle}\right)$\\
$= \textsf{LWE}_{\vec{s}, \sigma}\left(m_2(\Delta m_1 + e_1) + \sum\limits_{i=0}^ke^{\langle i \rangle}\right)$ \\
$= \textsf{LWE}_{\vec{s}, \sigma}\left(\Delta m_1m_2 + m_2e_1 + \sum\limits_{i=0}^ke^{\langle i \rangle}\right)$ \\
$\approx \textsf{LWE}_{\vec{s}, \sigma}(\Delta m_1m_2)$ \textcolor{red}{ $\rhd$ 因为 $m_2e_1 + \sum\limits_{i=0}^ke^{\langle i \rangle}$ 很小，所以 $m_2e_1$ 也很小}
\end{enumerate}

\end{myproof}

%This issue of noise growth is similar to that in ciphertext-to-plaintext multiplication (\autoref{subsubsec:glwe-mult-plain-discussion}). 
为了抑制噪声增长，需要进行噪声自举（将在 \autoref{subsec:tfhe-noise-bootstrapping} 中介绍）。

\subsubsection{推广到 GLWE 与 GGSW 的乘法}
\label{subsubsec:tfhe-glwe-to-ggsw-multiplication}

可以将 TFHE 的 LWE 与 GSW 乘法进一步推广为以下两个密文之间的 GLWE 与 GGSW 乘法：$\textsf{GLWE}_{\vec{S}, \sigma}(\Delta M_1) \cdot \textsf{GGSW}_{\vec{S}, \sigma}^{\beta, l}(M_2)$，其中 $M_1$、$M_2$ 和 $S$ 都是 $(n-1)$ 次多项式。


$ $

\noindent 第一项 $\textsf{GLWE}_{\vec{S}, \sigma}(\Delta M_1)$ 可由以下任一种方式得到：
\begin{itemize}
\item 对明文 $M_1$ 进行一次新的 GLWE 加密（\autoref{subsec:glwe-enc}）。
\item 对两个 GLWE 密文执行同态加法所得的结果（\autoref{sec:glwe-add-cipher}）。
\item 对一个 GLWE 密文和一个明文执行同态乘法所得的结果（\autoref{sec:glwe-mult-plain}）。
\end{itemize}

$ $

\noindent 第二项 $\textsf{GGSW}_{\vec{S}, \sigma}^{\beta, l}(M_2)$ 可由以下任一种方式得到：
\begin{itemize}
\item 对明文 $M_2$ 进行一次新的 GGSW 加密（\autoref{subsec:ggsw-enc}）。
\item 通过\textit{电路自举}将 $\textsf{GLWE}_{\vec{S}, \sigma}(\Delta M_2)$ 转换为 $\textsf{GGSW}_{\vec{S}, \sigma}^{\beta, l}(M_2)$（将在未来版本中介绍）。
\end{itemize}

$ $

\noindent 回顾以下关系：

$\textsf{GLWE}_{\vec{S}, \sigma}(\Delta M_1) = (A_0, A_1, \gap{$\cdots$}, A_{k-1}, B) \in \mathcal{R}_{n, q}^{k + 1}$，其中 $B = \sum\limits_{i=0}^{k-1}(A_i \cdot S_i) + \Delta M_1 + E$ \\ \textcolor{red}{ $\rhd$ 由 \autoref{subsec:glwe-enc} 可得}

$ $

$\textsf{GGSW}_{\vec{S}, \sigma}^{\beta, l}(M_2) = \Bigr \{ \{ \textsf{GLev}_{\vec{S}, \sigma}^{\beta, l} (-S_i \cdot M_2)  \}_{i=0}^{k-1}, \textsf{GLev}_{\vec{S}, \sigma}^{\beta, l}(M_2) \Bigl \} \in \mathcal{R}_{\langle n, q \rangle }^{(k+1) \cdot l \cdot (k'+1)}$ \textcolor{red}{ $\rhd$ 由 \autoref{subsec:ggsw-enc} 可得}

$ $

\noindent 使用以下记号：

$\textsf{GGSW}_{\vec{S}, \sigma}^{\beta, l}(M_2) = {\bar{C}} = (\bar{C_0},. \bar{C_1}, \gap{$\cdots$} \bar{C_k})$ 

$\bar{C_i} = \textsf{GLev}_{\vec{S}, \sigma}^{\beta, l}(-S_i \cdot M_2)$，其中 $0 \leq i \leq (k-1)$

$\bar{C_k} = \textsf{GLev}_{\vec{S}, \sigma}^{\beta, l}(M_2)$

$\textsf{ct} = \textsf{GLWE}_{\vec{S}, \sigma}(\Delta M_1) = (C_0, C_1, \gap{$\cdots$}, C_k) = (A_0, A_1, \gap{$\cdots$}, A_{k-1}, B)$

$ $

\noindent 定义以下 TFHE 密文乘法运算：

$\textsf{ct} \cdot {\bar{C}} = \sum\limits_{i=0}^{k}\langle \textsf{Decomp}^{\beta, l}(C_i), \bar{C_i} \rangle$

$ $

\noindent 于是下式成立：

\begin{tcolorbox}[title={\textbf{\tboxlabel{\ref*{subsubsec:tfhe-glwe-to-ggsw-multiplication}} 推广到 GLWE 与 GGSW 的乘法}}]
$\textsf{ct} = \textsf{GLWE}_{\vec{S}, \sigma}(\Delta M_1) = (A_0, A_1, \cdots, A_{k-1}, B)$

$\bar{C} = \textsf{GGSW}_{\vec{S}, \sigma}^{\beta, l}(M_2)$

$ = \bm(\textsf{GLev}_{\vec{S}, \sigma}^{\beta, l}(-S_0\cdot M_2), \textsf{GLev}_{\vec{S}, \sigma}^{\beta, l}(-S_1\cdot M_2), \cdots, \textsf{GLev}_{\vec{S}, \sigma}^{\beta, l}(-S_{k-1}\cdot M_2), \textsf{GLev}_{\vec{S}, \sigma}^{\beta, l}(M_2)\bm)$

$ $

$\textsf{GLWE}_{\vec{S}, \sigma}(\Delta M_1) \cdot \textsf{GGSW}_{\vec{S}, \sigma}^{\beta, l}(M_2) = \sum\limits_{i=0}^{k}\langle \textsf{Decomp}^{\beta, l}(C_i), \bar{C_i} \rangle \approx \textsf{GLWE}_{\vec{S}, \sigma}(\Delta M_1 M_2)$
\end{tcolorbox}

这意味着，将两个 TFHE 密文（一个为 GLWE，另一个为 GGSW）相乘，并解密所得的 GLWE 密文，结果与将两个原始明文相乘相同。


%\textbf{\underline{Proof}}
\begin{myproof}
\begin{enumerate}
\item $\sum\limits_{i=0}^{k}\langle \textsf{Decomp}^{\beta, l}(C_i), \bar{C_i} \rangle$ \\
$= \langle \textsf{Decomp}^{\beta, l}(A_0), \bar{C_0} \rangle + \langle \textsf{Decomp}^{\beta, l}(A_1), \bar{C_1} \rangle + \gap{$\cdots$} + \langle \textsf{Decomp}^{\beta, l}(A_{k-1}), \bar{C}_{k-1} \rangle + \langle \textsf{Decomp}^{\beta, l}(B), \bar{C_k} \rangle$ \\
\textcolor{red}{ $\rhd$ 展开两个向量的点积}
\item 对于 $i = k$： \\
$\textsf{Decomp}^{\beta, l}(B) = (B_1, B_2, \gap{$\cdots$}, B_l)$，其中 $B = B_1\dfrac{q}{\beta^1} + B_2\dfrac{q}{\beta^2} + \gap{$\cdots$} + B_l\dfrac{q}{\beta^l}$ \textcolor{red}{ $\rhd$ 由 \autoref{subsec:poly-decomp} 可得}\\
$\bar{C}_k = \textsf{GLev}_{\vec{S}, \sigma}^{\beta, l}(M_2) = \left(\textsf{GLWE}_{\vec{S}, \sigma}\left(M_{2}\dfrac{q}{\beta^1}\right), \textsf{GLWE}_{\vec{S}, \sigma}\left(M_{2}\dfrac{q}{\beta^2}\right), \gap{$\cdots$}, \textsf{GLWE}_{\vec{S}, \sigma}\left(M_{2}\dfrac{q}{\beta^l}\right) \right)$ 

\textcolor{red}{ $\rhd$ 为简化表达，省略每个 GLWE 密文中的噪声项 $E_{2,1}, \ldots, E_{2,l}$} \\
$ $

因此： \\
$\langle \textsf{Decomp}^{\beta, l}(B), \bar{C_k} \rangle$\\
$= B_1 \cdot \textsf{GLWE}_{\vec{S}, \sigma} \left (M_{2}\dfrac{q}{\beta^1} \right ) + B_2 \cdot \textsf{GLWE}_{\vec{S}, \sigma} \left (M_{2}\dfrac{q}{\beta^2} \right ) + \gap{$\cdots$} + B_l \cdot \textsf{GLWE}_{\vec{S}, \sigma} \left (M_{2}\dfrac{q}{\beta^l}\right)$ \\
$= \textsf{GLWE}_{\vec{S}, \sigma} \left (B_1M_{2}\dfrac{q}{\beta^1} \right ) + \textsf{GLWE}_{\vec{S}, \sigma} \left (B_2M_{2}\dfrac{q}{\beta^2} \right ) + \gap{$\cdots$} + \textsf{GLWE}_{\vec{S}, \sigma} \left (B_lM_{2}\dfrac{q}{\beta^l}\right)$ \textcolor{red}{ $\rhd$ 由 \autoref{sec:glwe-mult-plain} 可得} \\
$= \textsf{GLWE}_{\vec{S}, \sigma} \left (B_1M_{2}\dfrac{q}{\beta^1} + B_2M_{2}\dfrac{q}{\beta^2} + \gap{$\cdots$} + B_lM_{2}\dfrac{q}{\beta^l}\right)$ \textcolor{red}{ $\rhd$ 由 \autoref{sec:glwe-add-cipher} 可得}\\
$= \textsf{GLWE}_{\vec{S}, \sigma} \left (M_{2} \cdot \left ( B_1\dfrac{q}{\beta^1} + B_2\dfrac{q}{\beta^2} + \gap{$\cdots$} + B_l\dfrac{q}{\beta^l} \right)\right)$ \\
$= \textsf{GLWE}_{\vec{S}, \sigma} (M_{2}B)$ \textcolor{red}{ $\rhd$ 由 \autoref{subsec:poly-decomp} 可得}
\item 对于 $0 \leq i \leq (k - 1)$： \\
$\textsf{Decomp}^{\beta, l}(A_i) = (A_{\langle i, 1 \rangle}, A_{\langle i, 2 \rangle}, \gap{$\cdots$}, A_{\langle i, l \rangle})$，其中 $A_i = A_{\langle i, 1 \rangle}\dfrac{q}{\beta^1} + A_{\langle i, 2 \rangle}\dfrac{q}{\beta^2} + \gap{$\cdots$} + A_{\langle i, l \rangle}\dfrac{q}{\beta^l}$ \\
$\bar{C_i} =  \textsf{GLev}_{\vec{S}, \sigma}^{\beta, l}(-S_iM_2) = \left(\textsf{GLWE}_{\vec{S}, \sigma}\left(-S_iM_2\dfrac{q}{\beta^1}\right), \textsf{GLWE}_{\vec{S}, \sigma}\left(-S_iM_2\dfrac{q}{\beta^2}\right), \gap{$\cdots$}, \textsf{GLWE}_{\vec{S}, \sigma}\left(-S_iM_2\dfrac{q}{\beta^l}\right) \right)$ \\

$ $
因此： \\
$\langle \textsf{Decomp}^{\beta, l}(A_0), \bar{C_0} \rangle + \langle \textsf{Decomp}^{\beta, l}(A_1), \bar{C_1} \rangle + \gap{$\cdots$} + \langle \textsf{Decomp}^{\beta, l}(A_{k-1}), \bar{C}_{k-1} \rangle$ \\
$= \sum\limits_{i=0}^{k-1}\langle \textsf{Decomp}^{\beta, l}(A_i), \bar{C_i} \rangle$ \\
$= \sum\limits_{i=0}^{k-1}\left(A_{\langle i, 1\rangle} \cdot \textsf{GLWE}_{\vec{S}, \sigma}\left(-S_iM_2\dfrac{q}{\beta^1}\right) + A_{\langle i, 2\rangle} \cdot \textsf{GLWE}_{\vec{S}, \sigma}\left(-S_iM_2\dfrac{q}{\beta^2}\right) + \gap{$\cdots$} + A_{\langle i, l\rangle} \cdot \textsf{GLWE}_{\vec{S}, \sigma}\left(-S_iM_2\dfrac{q}{\beta^l}\right)\right)$ \\
$= \sum\limits_{i=0}^{k-1}\left(\textsf{GLWE}_{\vec{S}, \sigma}\left(-A_{\langle i, 1\rangle}S_iM_2\dfrac{q}{\beta^1}\right) + \textsf{GLWE}_{\vec{S}, \sigma}\left(-A_{\langle i, 2\rangle}S_iM_2\dfrac{q}{\beta^2}\right) + \gap{$\cdots$} + \textsf{GLWE}_{\vec{S}, \sigma}\left(-A_{\langle i, l\rangle}S_iM_2\dfrac{q}{\beta^l}\right)\right)$ \\
$= \sum\limits_{i=0}^{k-1}\textsf{GLWE}_{\vec{S}, \sigma}\left(-A_{\langle i, 1\rangle}S_iM_2\dfrac{q}{\beta^1} + -A_{\langle i, 2\rangle}S_iM_2\dfrac{q}{\beta^2} + \gap{$\cdots$} + -A_{\langle i, l\rangle}S_iM_2\dfrac{q}{\beta^l}\right)$ \\
$= \sum\limits_{i=0}^{k-1}\textsf{GLWE}_{\vec{S}, \sigma}\left(-S_iM_2 \cdot \left(A_{\langle i, 1\rangle}\dfrac{q}{\beta^1} + A_{\langle i, 2\rangle}\dfrac{q}{\beta^2} + \gap{$\cdots$} + A_{\langle i, l\rangle}\dfrac{q}{\beta^l}\right)\right)$ \\
$= \sum\limits_{i=0}^{k-1}\textsf{GLWE}_{\vec{S}, \sigma}(-S_iM_2A_i)$
\item 根据步骤 2 和 3， \\
$\sum\limits_{i=0}^{k}\langle \textsf{Decomp}^{\beta, l}(C_i), \bar{C_i} \rangle$ \\
$= \sum\limits_{i=0}^{k-1}\textsf{GLWE}_{\vec{S}, \sigma}(-S_iM_2A_i) + \textsf{GLWE}_{\vec{S}, \sigma} (M_{2}B)$ \\ 
$= \textsf{GLWE}_{\vec{S}, \sigma}\Big(\sum\limits_{i=0}^{k-1}(-S_iM_2A_i) + M_{2}B\Big)$ \textcolor{red}{ $\rhd$ 两个 GLWE 密文相加} \\
$= \textsf{GLWE}_{\vec{S}, \sigma}\Big(BM_2 - \sum\limits_{i=0}^{k-1}M_2A_iS_i\Big)$  \\ 
$= \textsf{GLWE}_{\vec{S}, \sigma}\Big(M_2(B - \sum\limits_{i=0}^{k-1}A_iS_i)\Big)$\\
$= \textsf{GLWE}_{\vec{S}, \sigma}(M_2(\Delta M_1 + E))$ \\
$= \textsf{GLWE}_{\vec{S}, \sigma}(\Delta M_1M_2 + M_2E)$ \\
$\approx \textsf{GLWE}_{\vec{S}, \sigma}(\Delta M_1M_2)$ \textcolor{red}{ $\rhd$ 因为 $E$ 很小，所以 $M_2E$ 也很小}
\end{enumerate}

\end{myproof}

\subsection{系数提取}
\label{subsec:tfhe-extraction}

\textbf{- 参考资料：}
\href{https://www.zama.ai/post/tfhe-deep-dive-part-4}{深入理解 TFHE：第四部分——可编程自举}~\cite{tfhe-4}

$ $

在 TFHE 中，系数提取是指从以 GLWE 密文形式加密的多项式中提取某个系数的过程。提取出的系数仍以 LWE 密文的形式表示（\autoref{sec:lwe}）。 %We will explain how to extract the coefficient of a plaintext polynomial $M$ from a GLWE ciphertext and RLWE ciphertext, respectively.


%\subsubsection{Coefficient Extraction from a GLWE Ciphertext}
%\label{subsubsec:tfhe-extraction-glwe}

注意，在 GLWE 密码系统中，明文 $M$ 被编码为一个多项式，其中每个系数都编码一个明文值 $m_0, m_1, \cdots, m_{n-1}$。

%\subsubsection{Overview}
%\label{subsec:tfhe-extraction-overview}

假设 GLWE 密文的设置如下： \\
$M = \sum\limits_{j=0}^{n-1}m_jX^j \in \mathcal{R}_{\langle n, q \rangle}$ 

$S = \left(S_0 = \sum\limits_{j=0}^{n-1}s_{0,j}X^j, S_1 = \sum\limits_{j=0}^{n-1}s_{1,j}X^j, \gap{$\cdots$}, S_{k-1} = \sum\limits_{j=0}^{n-1}s_{k-1,j}X^j \right)$ 

$\textsf{GLWE}_{\vec{S}, \sigma}(\Delta M) = \left(A_0 = \sum\limits_{j=0}^{n-1}a_{0,j}X^j, A_1 = \sum\limits_{j=0}^{n-1}a_{1,j}X^j, \gap{$\cdots$}, A_{k-1} = \sum\limits_{j=0}^{n-1}a_{k-1,j}X^j, B = \sum\limits_{j=0}^{n-1}b_{j}X^j\right)$ 

$B = \sum\limits_{i=0}^{k-1}A_iS_i + \Delta M + E$ 

$E = \sum\limits_{i=0}^{n-1}e_iX^i$ 

$ $

\noindent 注意：

$\Delta M + E = B - \sum\limits_{i=0}^{k-1}A_iS_i$ 

$ = (\Delta m_0 + \Delta m_1X + \gap{$\cdots$} + \Delta m_{n-1}X^{n-1}) + (e_0 + e_1X + \gap{$\cdots$} + e_{n-1}X^{n-1})$

$= (\Delta m_0 + e_0) + (\Delta m_1 + e_1)X + \gap{$\cdots$} + (\Delta m_{n-1} + e_{n-1})X^{n-1}$

$ $

\noindent 该公式还可以写为：

$B - \sum\limits_{i=0}^{k-1}A_iS_i$ 

$ = (b_0 + b_1X + \gap{$\cdots$} + b_{n-1}X^{n-1} )$ 

$ - (a_{0,0} + a_{0,1}X + \gap{$\cdots$} + a_{0, n-1}X^{n-1})(s_{0,0} + s_{0,1}X + \gap{$\cdots$} + s_{0, n-1}X^{n-1})$ 

$ - (a_{1,0} + a_{1,1}X + \gap{$\cdots$} + a_{1, n-1}X^{n-1})(s_{1,0} + s_{1,1}X + \gap{$\cdots$} + s_{1, n-1}X^{n-1})$

$ - \gap{$\cdots$} $ 

$ - (a_{k-1,0} + a_{k-1,1}X + \gap{$\cdots$} + a_{k-1, n-1}X^{n-1})(s_{k-1,0} + s_{k-1,1}X + \gap{$\cdots$} + s_{k-1, n-1}X^{n-1})$ 

$ $

$ = \left(b_0 - \left( \sum\limits_{i=0}^{k-1} \sum\limits_{j=0}^{0}(a_{i,0-j}s_{i,j}) - \sum\limits_{i=0}^{k-1} \sum\limits_{j=1}^{n-1}(a_{i,n-j}s_{i,j}) \right)\right)$

$ + \left(b_1 - \left( \sum\limits_{i=0}^{k-1} \sum\limits_{j=0}^{1}(a_{i,1-j}s_{i,j}) - \sum\limits_{i=0}^{k-1} \sum\limits_{j=2}^{n-1}(a_{i,n+1-j}s_{i,j})   \right) \right)\cdot X$ 



$ + \left(b_2 - \left( \sum\limits_{i=0}^{k-1} \sum\limits_{j=0}^{2}(a_{i,2-j}s_{i,j}) - \sum\limits_{i=0}^{k-1} \sum\limits_{j=3}^{n-1}(a_{i,n+2-j}s_{i,j})   \right) \right)\cdot X^2$ 

$ $

$\gap{$\cdots$}$ 

$ $

$ + \left(b_{n-1} - \left(  \sum\limits_{i=0}^{k-1} \sum\limits_{j=0}^{n-1}(a_{i,n-1-j}s_{i,j}) - \sum\limits_{i=0}^{k-1} \sum\limits_{j=n}^{n-1}(a_{i,n+(n-1)-j}s_{i,j})  \right) \right)\cdot X^{n-1}$ 

\textcolor{red}{ $\rhd$ 按相同指数合并各项}

$ $


$ $

$= \sum\limits_{h=0}^{n-1}  \left(b_h - \left(  \sum\limits_{i=0}^{k-1} \sum\limits_{j=0}^{h}(a_{i,h-j}s_{i,j}) - \sum\limits_{i=0}^{k-1} \sum\limits_{j=h+1}^{n-1}(a_{i,n+h-j}s_{i,j})  \right) \right)\cdot X^{h}  $


$ $

$= \sum\limits_{h=0}^{n-1}  C_h \cdot X^{h}  $，其中 $C_h = b_h - \left(  \sum\limits_{i=0}^{k-1} \sum\limits_{j=0}^{h}(a_{i,h-j}s_{i,j}) - \sum\limits_{i=0}^{k-1} \sum\limits_{j=h+1}^{n-1}(a_{i,n+h-j}s_{i,j})  \right)$

$ $

\noindent 注意，在上述 $(n-1)$ 次多项式中，每个 $X^h$ 项的系数 $C_h$ 都可以表示为如下 LWE 密文 $\textsf{ct}_h$：

$S' = (s_{0,0}, s_{0,1}, \gap{$\cdots$}, s_{0,n-1}, s_{1,0}, s_{1,1}, \gap{$\cdots$}, s_{1, n-1}, \gap{$\cdots$}, s_{k-1, n-1}) = (s'_0, s'_1, \gap{$\cdots$}, s'_{nk-1} ) \in \mathbb{Z}_q^{nk}$


$\textsf{ct}_h = (a'_0, a'_1, \gap{$\cdots$}, a'_{nk-1}, b_h) \in \mathbb{Z}_q^{nk + 1}$


\[
    \text{，其中 } a'_{n \cdot i + j} =   
\begin{cases}
    a_{i,h - j} \text{（若 } 0 \leq j \leq h\text{）}\\
    -a_{i,n + h - j} \text{（若 } h+1 \leq j \leq n-1\text{）}\\
\end{cases}
\centering , b_h \text{ 直接取自多项式 } B
\]


\noindent 注意 $b_h - \sum\limits_{i=0}^{nk-1}s'_ia'_i = \Delta m_h + e_h$。这意味着可以用 $C_h$ 的加密形式 $\textsf{LWE}_{\vec{s}_{'}, \sigma}(\Delta m_h)$ 来替代它，即用加密了 $M$ 的第 $h$ 个系数的 LWE 密文 $\textsf{ct}_h$ 来表示。因此，我们从 $\textsf{GLWE}_{\vec{S}, \sigma}(\Delta M)$ 中提取出了 $\textsf{LWE}_{\vec{s}_{'}, \sigma}(\Delta m_h)$。这一操作称为系数提取；它只是重新排列 GLWE 密文多项式的系数以提取一个 LWE 密文，因此不会增加任何噪声。

得到 $\textsf{LWE}_{\vec{s}_{'}, \sigma}(\Delta m_h)$ 后，即可执行 $\vec{s}_{'} \rightarrow \vec{s}$ 的密钥切换（\autoref{subsec:tfhe-key-switching}）。

\begin{tcolorbox}[title={\textbf{\tboxlabel{\ref*{subsec:tfhe-extraction}} GLWE 密文的系数提取}}]
给定以下 GLWE 密文：

$M = \sum\limits_{j=0}^{n-1}m_jX^j \in \mathcal{R}_{\langle n, t \rangle}$

$\vec{S} = \left(S_0 = \sum\limits_{j=0}^{n-1}s_{0,j}X^j, S_1 = \sum\limits_{j=0}^{n-1}s_{1,j}X^j, \gap{$\cdots$}, S_{k-1} = \sum\limits_{j=0}^{n-1}s_{k-1,j}X^j \right)$

$\textsf{GLWE}_{\vec{S}, \sigma}(\Delta M) = \left(A_0 = \sum\limits_{j=0}^{n-1}a_{0,j}X^j, A_1 = \sum\limits_{j=0}^{n-1}a_{1,j}X^j, \gap{$\cdots$}, A_{k-1} = \sum\limits_{j=0}^{n-1}a_{k-1,j}X^j, B = \sum\limits_{j=0}^{n-1}b_{j}X^j\right)$

$B = \sum\limits_{i=0}^{k-1}A_iS_i + \Delta M + E \bmod q$, \text{ } $E = \sum\limits_{i=0}^{n-1}e_iX^i$

$ $

$\textsf{LWE}_{\vec{s}_{'}, \sigma}(\Delta m_h)$ 是一个 LWE 密文，它加密了 $\Delta M$ 的第 $h$ 个系数（即 $\Delta m_h$）。可以按如下方式从 $\textsf{GLWE}_{\vec{S}, \sigma}(\Delta M)$ 中提取 $\textsf{LWE}_{\vec{s}_{'}, \sigma}(\Delta m_h)$：

$ $

$\vec{s}_{'} = (s_{0,0}, s_{0,1}, \gap{$\cdots$}, s_{0,n-1}, s_{1,0}, s_{1,1}, \gap{$\cdots$}, s_{1, n-1}, \gap{$\cdots$}, s_{k-1, n-1}) = (s'_0, s'_1, \gap{$\cdots$}, s'_{nk-1} ) \in \mathbb{Z}_q^{nk}$

$\textsf{LWE}_{\vec{s}_{'}, \sigma}(\Delta m_h) = (a_0', a_1', \gap{$\cdots$} , a_{nk-1}', b_h) \in \mathbb{Z}_q^{nk + 1}$


\[
    \text{，其中 } a'_{n \cdot i + j} =   
\begin{cases}
    a_{i,h - j} \text{（若 } 0 \leq j \leq h\text{）}\\
    -a_{i,n + h - j} \text{（若 } h+1 \leq j \leq n-1\text{）}\\
\end{cases}
, b_h \text{ 取自多项式 } B
\]

得到 $\textsf{LWE}_{\vec{s}_{'}, \sigma}(\Delta m_h)$ 后，执行 $\vec{s}_{'} \rightarrow \vec{s}$ 的密钥切换（\autoref{subsec:tfhe-key-switching}）。

\end{tcolorbox}



\subsection{噪声自举}
\label{subsec:tfhe-noise-bootstrapping}


\textbf{- 参考资料：}
\href{https://www.zama.ai/post/tfhe-deep-dive-part-4}{深入理解 TFHE：第四部分——可编程自举}~\cite{tfhe-4}

$ $

连续对 TFHE 密文进行同态加法不一定会使噪声 $e$ 增大，因为 $e$ 是按照高斯分布随机生成的，许多噪声相加后的均值为 0。另一方面，连续进行同态乘法会使噪声增长，因为噪声项也会相乘，从而增大其幅度。因此，需要在噪声侵入存放明文 $m$ 的高位之前，以某种方式将噪声\textit{重置}（即防止 \autoref{fig:scaling} 中红色噪声位溢出到蓝色明文位）。将噪声重新初始化为较小值的过程称为噪声自举。

如本节开头所述，TFHE 使用 LWE（即多项式次数为 0 的 GLWE）对明文进行加密与解密。也就是说，每个明文都是 $m$（一个数），被编码为零次多项式。此外，加密各个 $m$ 的秘密密钥 S 是向量 $ \{s_0, s_1, \text{ } \cdots \text{ }, s_{k-1} \}$，而不是一个多项式。另一方面，TFHE 的噪声自举使用 GLWE 密文之间的同态加法，以及 GLWE 与 GGSW 密文之间的同态乘法。

%Nonetheless, the reason why we learned TFHE ciphertext addition (\autoref{subsec:glwe-add} and \autoref{subsec:glwe-add-plain}), multiplication (\autoref{subsec:tfhe-mult-cipher} and \autoref{subsec:glwe-mult-plain}), and polynomial rotation (\autoref{subsec:coeff-rotation}) in terms of GLWE is because the bootstrapping procedure we will explain in this subsection requires homomorphic addition and multiplication of GLWE ciphertexts.

假设有如下 TFHE 密文：

$ $

$\textsf{LWE}_{\vec{s}, \sigma}(\Delta m) = (a_0, a_1, \gap{$\cdots$} a_{k-1}, b)$

$b = \sum\limits_{i=0}^{k-1} a_is_i + \Delta m + e_b$

$\vec{s} = (s_0, s_1, \gap{$\cdots$} s_{k-1})$

$ $

，其中 $e_b$ 是经过许多次密文（或明文）乘法累积而成的较大噪声。噪声自举的目标是将 $(a_0, a_1, \gap{$\cdots$} a_{k-1}, b)$ 转换为 $(a_0', a_1', \gap{$\cdots$} a_{k-1}', b')$，使其满足：

$ $

$b' = \sum\limits_{i=0}^{k-1} a_i's_i + \Delta m + e_s$

$ $

，其中 $e_s$ 是重新初始化后的噪声。

$ $
\subsubsection{概述}
\label{subsec:bootstrapping-overview}

为了实现噪声自举，需要构造一个专门设计的 $(n-1)$ 次多项式 $V(X)$，称为查找表（Lookup Table，LUT）。在介绍 $V(X)$ 之前，先通过一个初步的 LUT 多项式 $V_q(X)$ 说明其基本思想。设想多项式 $V_q(X)$ 中每个项 $X^{j}$ 的指数 $j = \Delta m_i + e_*$ 是带有噪声 $e_* \in \mathbb{Z}_{\Delta}$ 的明文 $m_i$，而对应的系数 $v_{j} = m_i$ 则是不含噪声的明文。因此，将 $(q-1)$ 次多项式 $V_q(X)$ 定义如下：

$V_q(X) = v_0 + v_1X^1 + v_2X^2 + \gap{$\cdots$} + v_{q-1}X^{q-1}$

\text{ } $= m_0X^{\Delta m_0 + e_0} + m_0X^{\Delta m_0 + e_1} + m_0X^{\Delta m_0 + e_2} + \gap{$\cdots$} + m_0X^{\Delta m_0 + e_{\Delta - 1}}$ \textcolor{red}{ $\rhd$ 共 $\Delta$ 项}

\text{ } $ + \text{ } m_1X^{\Delta m_1 + e_0} + m_1X^{\Delta m_1 + e_1} + m_1X^{\Delta m_1 + e_2} + \gap{$\cdots$} + m_1X^{\Delta m_1 + e_{\Delta - 1}}$ \textcolor{red}{ $\rhd$ 共 $\Delta$ 项}

\text{ } $ + \gap{$\cdots$} $

\text{ } $ + \text{ } m_{t - 1}X^{\Delta m_{t - 1} + e_0} + m_{t - 1}X^{\Delta m_{t - 1} + e_1} + m_{t - 1}X^{\Delta m_{t - 1} + e_2} + \gap{$\cdots$} + m_{t - 1}X^{\Delta m_{t - 1} + e_{\Delta - 1}}$ \textcolor{red}{ $\rhd$ 共 $\Delta$ 项}

$ $

在上式中，$m_i$ 和 $e_k$ 分别遍历所有可能的明文消息和误差值（其中 $m_i \in \mathbb{Z}_t$ 且 $e_k \in \mathbb{Z}_{\Delta}$）。 %Note that $\mathbb{Z}_{q} = \{0, 1, \cdots, q-1\} = \{\Delta m_0 + e_0, \Delta m_0 + e_1, \cdots, \Delta m_{t-1} + e_{\Delta - 1}\}$.

我们将 $V_q(X)$ 设计为具有如下性质：每个项 $v_{j}X^{j}$ 都表示特殊的映射（指数，系数）$= (\Delta m_i + e_{*}, m_i)$，其中 $e_*$ 可以取 $\mathbb{Z}_{\Delta}$ 中的任意值。在 TFHE 的设置阶段，使用新定义的 GLWE 密钥 $\vec{S}_{bk}$ 对 $V_q(X)$ 进行 GLWE 加密；这一\textit{自举密钥}与 LWE 秘密密钥 $\vec{s}$ 不同。$\vec{S}_{bk}$ 是由系数为二进制值的 $(n-1)$ 次多项式构成的列表。随后，在噪声自举阶段，利用多项式系数旋转方法 1（\autoref{subsec:coeff-rotation} 中的小结~\ref*{subsec:coeff-rotation}.1），计算 $V \cdot X^{-(\Delta m + e)} = V'$，将 $V$ 的系数向左旋转 $\Delta m + e$ 个位置。接着，利用系数提取技术（\autoref{subsec:tfhe-extraction}）提取该多项式常数项的系数，即旋转后的多项式 $V'$ 中最左侧零次项的系数。此外，$V_q(X)$ 在 TFHE 设置阶段已被加密为 GLWE 密文，因此从旋转后的 $V'_q(X)$ 中提取出的常数项系数，是对 $m$ 的 LWE 加密（即 $\textsf{LWE}_{\vec{s}, \sigma}(\Delta m)$），其噪声已重新初始化（即充分降低）。

因此，至少在当前阶段，可以从概念上将噪声自举过程理解为：

$ $

\begin{enumerate}
\item \textbf{\underline{输入}：} 加密 $m$ 的带噪声密文 $\textsf{LWE}_{\vec{s}, \sigma}(\Delta m + e)$
\item 将输入转换为 $X^{-(\Delta m + e)}$ 的形式，作为 $V_q(X)$（查找表）的旋转因子。
\item 计算 $V_q \cdot X^{-(\Delta m + e)} = V'_q$，将 $V_q$ 向左旋转 $\Delta m + e$ 个位置。
\item 以 LWE 密文形式提取旋转后的 $V'_q(X)$ 的常数项系数 $m$，即得到 $\textsf{LWE}_{\vec{s}, \sigma}(\Delta m)$。
\item \textbf{\underline{输出}：} 明文 $m$ 的 LWE 加密 $\textsf{LWE}_{\vec{s}, \sigma}(\Delta m)$，其噪声已重新初始化
\end{enumerate}

$ $

输出 $\textsf{LWE}_{\vec{s}, \sigma}(\Delta m)$ 与输入密文加密相同的明文消息，但噪声已充分降低。因此，可以将输出 $\textsf{LWE}_{\vec{s}, \sigma}(\Delta m)$ 用于后续的 TFHE 同态运算，例如加法或乘法。在该噪声自举过程中，多项式 $V_q$ 充当\textit{字典}，保存从带噪声明文 $\Delta m + e$（即以 $\Delta m + e = j$ 表示，其中 $v_jX^j$）到无噪声明文 $m$（即以 $m = v_j$ 表示，其中 $v_jX^j$）的映射。因此，$V_q$ 被称为查找表（LUT）。

那么，$V_q(X)$ 的次数应取多大？为了让 $V_q(X)$ 编码从 $\Delta m + e \in \mathbb{Z}_q$ 到 $m \in \mathbb{Z}_t$ 的所有可能映射，$V_q(X)$ 应为 $(q-1)$ 次多项式。然而，$q$ 非常大，处理一个 $(q-1)$ 次多项式在计算上不可行。因此，实际使用的是一个小得多的多项式 $V(X)$，其次数只有 $n-1$。回顾 TFHE 的设置，有 $n \ll q$。因此，需要以\textit{比例正确}的方式，将较大的密文空间 $\Delta m + e \in \mathbb{Z}_q$ \textit{压缩}到小得多的空间 $\mathbb{Z}_n$，并将压缩后的值编码为 $X^j$ 的指数。为实现这种 $\mathbb{Z}_q \rightarrow \mathbb{Z}_n$ 的按比例压缩，将使用 \autoref{subsec:modulus-switch-lwe} 中介绍的 LWE 模数切换技术。

\subsubsection{噪声自举中的模数切换}
\label{subsec:bootstrapping-modulus-switch}

为了避免使用庞大的 $(q-1)$ 次多项式 $V_q$，我们将 $q$ 个可能的密文元素 $\Delta m + e \in \mathbb{Z}_q$ 压缩为 $(n-1)$ 次多项式 $V$ 的 $n$ 个不同指数；其中，$V$ 中的每个项 $v_jX^j$ 都表示映射 $j \rightarrow v_j$，即从带噪声明文映射到无噪声明文。但是，将 $(n-1)$ 次多项式 $V$ 的系数向左旋转时，$v_jX^j$ 一旦越过 $X^0$ 次项与 $X^{n-1}$ 次项之间的边界，$v_j$ 的符号就会翻转为 $-v_j$（见 \autoref{subsec:coeff-rotation-ex} 中的例子）。由于这种系数符号翻转，$(n-1)$ 次多项式 $V$ 理论上可以编码共 $2n$ 种不同的系数状态：$(v_0, v_1, v_2, \gap{$\cdots$}, v_{n-1}, -v_0, -v_1, \gap{$\cdots$}, -v_{n-1})$。为了将这 $2n$ 个不同的系数中的每一个移到 $V$ 的常数项位置（即将系数 $v_j$ 移到 $V$ 的最左侧项），旋转运算 $V \cdot X^{-j}$ 可以使用 $2n$ 个不同的 $j$ 值，即 $\{0, 1, 2, \gap{$\cdots$}, n-1, n, \gap{$\cdots$}, 2n-1\}$，从而将 $(v_0, v_1, v_2, \gap{$\cdots$}, v_{n-1}, -v_0, -v_1, \gap{$\cdots$}, -v_{n-1})$ 中的每一个系数移到常数项位置。这说明 $V\cdot X^{-j}$ 中的指数 $j$ 可以取 $2n$ 个不同值中的任意一个，以覆盖 $V$ 的全部 $2n$ 种可能的（包括符号翻转后的）系数状态。此外，记住 $j = \Delta m + e$。因此，将对 $\Delta m + e$ 执行 $q \rightarrow 2n$ 的模数切换。利用 LWE 模数切换技术（\autoref{subsec:modulus-switch-lwe}），原始 LWE 密文 $\textsf{LWE}_{\vec{s}, \sigma}(\Delta m + e) = ({a}_0, {a}_1, \gap{$\cdots$} {a}_{k-1}, {b}) \in \mathbb{Z}_q^{k+1}$，即噪声自举过程的初始输入，将转换为：

$ $

$\textsf{LWE}_{\vec{s}, \sigma}(\hat{\Delta} m) = (\hat{a}_0, \hat{a}_1, \gap{$\cdots$} \hat{a}_{k-1}, \hat{b}) \in \mathbb{Z}_{2n}^{k+1}$

$\vec{s} = (s_0, s_1, \gap{$\cdots$} s_{k-1}) \in \mathbb{Z}_2^{k}$  \textcolor{red}{ $\rhd$ 因为每个 $s_i$ 都是二进制值，所以秘密密钥保持不变}

$\hat{\Delta} = \Delta\dfrac{2n}{q} = \dfrac{2n}{t} \in \mathbb{Z}_{2n}$

$\hat{a}_i = \left \lceil a_i\dfrac{2n}{q} \right \rfloor \in \mathbb{Z}_{2n}$

$\hat{e} = \left \lceil e\dfrac{2n}{q} \right \rfloor \in \mathbb{Z}_{2n}$

$\hat{b} = \left \lceil b\dfrac{2n}{q} \right \rfloor \approx \sum\limits_{i=0}^{k-1} \hat{a}_is_i + \hat{\Delta} m + \hat{e}  \in \mathbb{Z}_{2n}$



$ $

\para{$V(X)$ 的次数与安全性}：如果目标只是设计次数最低的多项式 $V$，使其系数能够映射明文 $m$ 的所有可能取值，那么构造一个 ${t}$ 次多项式 $V$ 就足够了。尽管如此，将 $V$ 的次数选为 $2n$ 而不是 ${t}$，是为了保证足够的安全级别：多项式次数 $n$ 越高，方案对攻击的抵抗能力就越强。



\subsubsection{将可用明文空间减半}
\label{subsec:tfhe-zero-padding} 

一个问题是，LUT 多项式 $V(X)$ 的旋转具有\textit{负循环}性质，即 $V(X)\cdot X^{n} = -V(X)$（也就是说，系数每旋转 $n$ 个位置就会翻转符号）。更一般地：


$V(X)\cdot X^{-j}= V(X)\cdot X^{2n - j} = V(X)\cdot X^{4n - j} = \cdots$


$
= V(X)\cdot X^{-(j \bmod 2n)} =  
\begin{cases}
    v_j  + v_{j+1}X + \cdots, \text{当 } 0\leq i<n\\
    -v_{j} - v_{j-1}X - \cdots, \text{当 } n\leq j<2n
\end{cases}
$

$ $

\noindent ，其中 $v_j$ 表示将多项式 $V(X)$ 向左旋转 $j$ 个位置后得到的常数项系数。问题在于，每当旋转越过 $X^0$ 与 $X^{n-1}$ 之间的边界时，$v_j$ 的符号就会翻转。给定经过模数切换的值 $\hat{a}_j$、$\hat{e}$ 和 $\hat{b}$，构造如下 LUT 多项式 $V(X)$：

$ $

$V(X) = v_0 + v_1X^1 + v_2X^2 + \gap{$\cdots$} + v_{n-1}X^{n-1}$

\text{ } $= m_0X^{\hat{\Delta} m_0 + \hat{e}_0} + m_0X^{\hat{\Delta} m_0 + \hat{e}_1} + m_0X^{\hat{\Delta} m_0 + \hat{e}_2} + \gap{$\cdots$} +  m_0X^{\hat{\Delta} m_0 + \hat{e}_{\hat{\Delta} - 1}}$ \textcolor{red}{ $\rhd$ 共 $\hat{\Delta}$ 项}

\text{ } $ + \text{ } m_1X^{\hat{\Delta} m_1 + \hat{e}_0} + m_1X^{\hat{\Delta} m_1 + \hat{e}_1} + m_1X^{\hat{\Delta} m_1 + \hat{e}_2} + \gap{$\cdots$} + m_1X^{\hat{\Delta} m_1 + \hat{e}_{\hat{\Delta} - 1}}$ \textcolor{red}{ $\rhd$ 共 $\hat{\Delta}$ 项}

\text{ } $ + \gap{$\cdots$} $

\text{ } $ + \text{ } m_{t/2 - 1}X^{\hat{\Delta} m_{t/2 - 1} + \hat{e}_0} + m_{t/2 - 1}X^{\hat{\Delta} m_{t/2 - 1} + \hat{e}_1} + m_{t/2 - 1}X^{\hat{\Delta} m_{t/2 - 1} + \hat{e}_2} + \gap{$\cdots$} + m_{t/2 - 1}X^{\hat{\Delta} m_{t/2 - 1} + \hat{e}_{\hat{\Delta} - 1}}$ \textcolor{red}{ $\rhd$ 共 $\hat{\Delta}$ 项}

$ $


回顾，当 $j= \{0, 1, \cdots , n-1\}$ 时，计算 $V(X)\cdot X^{-j}$ 可以将 $V(X)$ 的系数向左旋转 $\{0, 1, \cdots n-1 \}$ 个位置。对 $V(X)$ 执行每一次 $j$ 个位置的旋转，其中 $j= \{0, 1, \cdots , n-1\}$，旋转后的多项式 $V'(X)$ 的常数项系数依次为：

$\overbrace{\overbrace{\underbrace{\Delta m_0, \Delta m_0, \gap{$\cdots$}}_{\text{系数对应项为 } X^{\hat{\Delta}m_0 + \hat{e}_*}}}^{\hat{\Delta} \text{ 次重复}} \overbrace{\underbrace{\Delta m_1, \Delta m_1, \gap{$\cdots$}}_{\text{系数对应项为 } X^{\hat{\Delta}m_1 + \hat{e}_*}}}^{\hat{\Delta} \text{ 次重复}} \gap{$\cdots$} \overbrace{\underbrace{\Delta m_{{t}/{2}-1}, \Delta m_{{t}/{2}-1}, \gap{$\cdots$}}_{\text{系数对应项为 } X^{\hat{\Delta}m_{{t}/{2}-1}+ \hat{e}_*} }}^{\hat{\Delta} \text{ 次重复}}}^{\text{旋转 $j = \{0, 1, \cdots n-1\}$ 个位置后 $V'(X)$ 常数项的系数}}$ 



$ $

在上式中，噪声 $\hat{e}_*$ 的取值范围为 $[0, \hat{\Delta})$。注意，我们将所有指数 $\hat\Delta m_i + \hat{e}_0, \hat\Delta m_i + \hat{e}_1, \cdots,\hat\Delta m_i + \hat{e}_{\hat\Delta - 1}$ 设计为映射到同一个系数值 $m_i$；这是因为，在对应噪声 $\hat{e}_*$ 被消除后，它们解密得到的底层明文 $m_i$ 相同。因此，每个 $m_i$ 都在 $V(X)$ 的系数中连续重复使用 $\hat\Delta$ 次。可以将这种连续重复理解为一种鲁棒性：对于任意噪声 $\hat{e}_* \in \mathbb{Z}_{\hat{\Delta}}$，都能将相应的 $\hat{\Delta}m_i + \hat{e}_*$ 映射到 $m_i$。

到目前为止，上述 $m_0, m_1, \cdots, m_{{t}/{2} - 1}$ 个系数组成的序列，是我们期望 $V'(X)$（即旋转后的多项式）在执行各次 $0, 1, \cdots, n-1$ 旋转时返回的常数项系数（其中每个 $m_i + 1 = m_{i+1}$）。然而，当旋转次数处于 $[n, 2n -1)$ 范围时，系数映射的正确性就会被破坏，因为多项式环的负循环性质使得系数在从 $X^0$ 次项越过边界到达 $X^{n-1}$ 次项时发生符号翻转。具体而言，每旋转 $n, n+1, \cdots, 2n-1$ 个位置后，多项式 $V'(X)$ 的常数项系数如下：

$ $

$\overbrace{\overbrace{\underbrace{-m_{0}, -m_{0}, \gap{$\cdots$}}_{-\text{系数对应项为 } X^{\hat{\Delta}m_{0} + \hat{e}_*}}}^{\hat{\Delta} \text{ 次重复}} \overbrace{\underbrace{-m_{1}, -m_{1}, \gap{$\cdots$}}_{-\text{系数对应项为 } X^{\hat{\Delta}m_{1} + \hat{e}_*}}}^{\hat{\Delta} \text{ 次重复}} \gap{$\cdots$} \overbrace{\underbrace{-m_{t/2 - 1}, -m_{t/2 - 1}, \gap{$\cdots$}}_{-\text{系数对应项为 } X^{\hat{\Delta}m_{t/2 - 1} + \hat{e}_*} }}^{\hat{\Delta} \text{ 次重复}}}^{\text{分别旋转 $n, n+1, \cdots 2n-1$ 个位置后 $V'(X)$ 常数项的系数}}$ 

$ $

由上式可见，旋转后的 $V'(X)$ 的常数项系数呈现旋转周期为 $n$ 的负循环模式：第二组 $n$ 次旋转得到的系数，恰好是第一组 $n$ 次旋转所得系数的相反数。下面说明这种负循环行为为什么会破坏（指数，系数）= $(\hat{\Delta} m + \hat{e}, m)$ 的映射。由于 TFHE 的明文和密文取值均定义在环上，旋转 LUT 多项式 $V(X)$ 时，理想情况下，旋转后的多项式 $V'(X)$ 的常数项值（即映射后的明文值）应按循环方式回绕，呈现环状模式；每个值连续重复 $\hat{\Delta}$ 次，以抵抗最大为 $\hat{e}_* \in Z_{\hat{\Delta}}$ 的噪声。然而，多项式环的负循环性质使后半组旋转的常数项值出现问题，因为它们恰好是前半组对应值的相反数，从而破坏了第一组与第二组数值之间按环回绕的模式。
%Specifically, as we gradually rotate $V$ with the rotation count $[0, 2n)$, the constant term's coefficient (i.e., $m_i$) should be ideally a value that is either the same as before or greater than the previous value ($m_{i-1}$), all the way to $2n$ rotations (to represent a ring pattern). This constraint holds for the first-half $i = [0, n)$, as the rotated $V'(X)$'s constant term's coefficient is $m_0, \cdots, m_1, \cdots, m_{t/2 - 1}$. However, for the second-half $i = [n, 2n)$, the rotated $V'(X)$'s constant term's coefficient is $-m_0, \cdots, -m_1, \cdots, -m_{t/2 - 1}$, which violates our constraint that $m_i \leq m_{i+1}$. Therefore, the combination of the first-half and the second-half domains of $i$ cannot mathematically implement a continuous modulo value range $\mathbb{Z}_t$. 

因此，$V(X)$ 作为完美 LUT 存在局限：它只能在 $i \in \mathbb{Z}_t$ 的一半范围内保持（指数，系数）= $(\hat{\Delta} m + \hat{e}, m)$ 的正确映射，而无法同时覆盖另一半。

\para{解决方法：}
我们已经看到，当 $V(X)$ 的旋转不超过 $n-1$ 个位置（即第一组旋转）时，$V(X)$ 的（指数，系数）= $(\hat{\Delta} m + \hat e, \Delta m)$ 映射仍保持环状一致性。因此，避免 LUT 多项式旋转的负循环问题，最简单的方法是在 TFHE 应用的\textit{设计中}限制噪声自举期间 $V(X)$ 的旋转不超过 $n-1$ 个位置。为此，当 TFHE 应用的计算流程处理明文值时（即在考虑同态运算之前的原始明文计算逻辑中），应确保每个计算步骤的可能输入和输出，仅涉及 $\mathbb{Z}_t$ 内预先指定的连续 $\dfrac{t}{2}$ 个模剩余值。该约束使 $\hat\Delta m + \hat e$ 的可能取值被限制在 $\mathbb{Z}_{2n}$ 内连续的 $n$ 个值中。由于 LUT 多项式 $V(X)$ 通过计算 $V(X)\cdot X^{-(\hat\Delta m + \hat e)}$ 进行旋转，应用设计将 $\hat\Delta m + \hat e$ 限制为至多 $n-1$（在 $2n - 1$ 范围内），便使 $V(X)$ 在噪声自举期间最多旋转 $n-1$ 个位置。这样就能避免出现使系数符号翻转的、有问题的 $\{n, n+1, \gap{$\cdots$}, 2n -1 \}$ 次旋转。

因此，以将应用可用明文值减半、限制为 $\mathbb{Z}_t$ 内连续的 $\dfrac{t}{2}$ 个值为代价，可以避免 $V(X)$ 的负循环旋转问题，从而保持 $V(X)$ 的（指数，系数）= $(\hat{\Delta} m + \hat e, m)$ 映射正确。

\begin{comment}
与此同时，密文空间 $q$ 保持不变。例如，取简单情形 $t = 8$，TFHE 应用可以设计如下编码方案：

%Note that during the blind rotation (\autoref{subsec:bootstrapping-blind-rotation}), each step's cumulative rotation positions of the homomorphic MUX gate may end up rotating $V(X)$ more than $n$ positions. However, enforcing $0 \leq \hat\Delta m + e < n$ guarantees the rotation-finished $V'(X)$ to have the cumulative rotation count to be some integer within $[0, n) \bmod 2n$.

\begin{multicols}{2}
\begin{itemize}
 \item $0000_2 \rightarrow 000_2 (= 0)$
 \item $0001_2 \rightarrow 001_2 (= 1)$
 \item $0010_2 \rightarrow 010_2 (= 2)$
 \item $0011_2 \rightarrow 011_2 (= 3)$
 \item $0100_2 \rightarrow 100_2 (= -4)$
 \item $0101_2 \rightarrow 101_2 (= -3)$
 \item $0110_2 \rightarrow 110_2 (= 2)$
 \item $0111_2 \rightarrow 111_2 (= 1)$
\end{itemize}
\end{multicols}

通过这种二进制编码，可以在 $t$ 个明文状态中选取 $\dfrac{t}{2}$ 个状态来编码有符号整数，此时最高有效位（MSB）始终为 0。
\end{comment}

综合上述考虑，最终的 LUT 多项式 $V(X)$ 如下：

\begin{tcolorbox}[title={\textbf{\tboxlabel{\ref*{subsec:tfhe-zero-padding}} 查找表多项式 $\bm{V(X)}$ 的结构}}]

$ $

$V(X) = v_0 + v_1X^1 + v_2X^2 + \gap{$\cdots$} + v_{n-1}X^{n-1}$

\text{ } $= m_0X^{\hat{\Delta} m_0 + \hat{e}_0} + m_0X^{\hat{\Delta} m_0 + \hat{e}_1} + m_0X^{\hat{\Delta} m_0 + \hat{e}_2} + \gap{$\cdots$} + m_0X^{\hat{\Delta} m_0 + \hat{e}_{\hat{\Delta} - 1}}$

\text{ } $ + \text{ } m_1X^{\hat{\Delta} m_1 + \hat{e}_0} + m_1X^{\hat{\Delta} m_1 + \hat{e}_1} + m_1X^{\hat{\Delta} m_1 + \hat{e}_2} + \gap{$\cdots$} + m_1X^{\hat{\Delta} m_1 + \hat{e}_{\hat{\Delta} - 1}}$

\text{ } $ + \gap{$\cdots$} $

\text{ } $ + m_{{t}/{2} - 1}X^{\hat{\Delta} m_{{t}/{2} - 1} + \hat{e}_0} + m_{{t}/{2} - 1}X^{\hat{\Delta} m_{{t}/{2} - 1} + \hat{e}_1} + m_{{t}/{2} - 1}X^{\hat{\Delta} m_{{t}/{2} - 1} + \hat{e}_2}$

\text{ } \text{ } $ + \gap{$\cdots$} + m_{{t}/{2} - 1}X^{\hat{\Delta} m_{{t}/{2} - 1} + \hat{e}_{\hat{\Delta} - 1}}$

$ $

\noindent ，其中 $\hat\Delta = \Delta \cdot \dfrac{2n}{q} = \dfrac{2n}{t}$。应用应确保 $m_0, m_1, \cdots, m_{t/2 - 1}$ 是 $\mathbb{Z}_t$ 内一些连续的模 $\dfrac{t}{2}$ 剩余值。该约束使 $V(X)$ 的旋转位置 $\hat{\Delta} m_i + \hat {e}_{*}$（其中 $e_* \in \mathbb{Z}_{\hat{\Delta}}$）至多有 $n$ 种连续的可能值，从而防止 $V(X)$ 旋转超过一个完整周期并触发负循环问题。


\end{tcolorbox}

LUT 多项式 $V(X)$ 也称为累加器。

\subsubsection{盲旋转}
\label{subsec:bootstrapping-blind-rotation} 

盲旋转是指旋转一个\textit{加密}多项式的系数，使观察者无法得知这些系数旋转了多少个位置。旋转后，也无法看出哪个系数移到了哪个次数的项。盲旋转使用基本的多项式旋转方法 1（\autoref{subsec:coeff-rotation} 中的小结~\ref*{subsec:coeff-rotation}），区别在于计算 $V(X)\cdot X^{-i}$ 是以同态方式完成的。

注意 $X^{\hat{\Delta} m + \hat{e}_b} = X^{\hat{b} - \sum_{i=0}^{k-1}{\hat{a}_is_i}}$，因此可以通过计算 $V \cdot X^{-(\hat{b} - \sum_{i=0}^{k-1}{\hat{a}_is_i})} = V \cdot X^{-\hat{b} + \sum_{i=0}^{k-1}{\hat{a}_is_i}}$ 来旋转 $V$。实际上，在不知道 LWE 秘密密钥 $S$ 的情况下，无法直接计算 LWE 解密公式 $-\hat{b} + \sum_{i=0}^{k-1}{\hat{a}_is_i}$（或 $\hat{b} - \sum_{i=0}^{k-1}{\hat{a}_is_i}$）。不过，只要在 TFHE 设置阶段给定 $\{GGSW_{\vec{S}_{bk}, \sigma}^{\beta, l}(s_i)\}_{i=0}^{k-1}$，便存在一种数学方法，使我们无须知道秘密密钥 $S$ 就能计算 $V \cdot X^{-\hat{b} + \sum_{i=0}^{k-1}{\hat{a}_is_i}}$。注意，$\{GGSW_{\vec{S}_{bk}, \sigma}^{\beta, l}(s_i)\}_{i=0}^{k-1}$ 是使用 GLWE 秘密密钥 $\vec{S}_{bk}$（即\textit{自举}密钥）对 LWE 秘密密钥 $S$ 进行的 GGSW 加密。按照以下流程，使用 $\vec{S}_{bk}$ 同态计算 $V \cdot X^{-\hat{b} + \sum_{i=0}^{k-1}{\hat{a}_is_i}}$，即将 $V$ 的系数向左盲旋转 $\hat{b} + \sum_{i=0}^{k-1}{\hat{a}_is_i}$ 个位置：

\begin{enumerate}
\item 在 TFHE 设置阶段，使用自举密钥 $\vec{S}_{bk}$ 对多项式 $V$ 进行 GLWE 加密，从而使 $V$ 的每个系数都受到加密保护。
\item 计算 $V_0 = V \cdot X^{-\hat{b}}$，本质上就是将 $V$ 的各多项式向左旋转 $\hat{b}$ 个位置。由于 $\hat{b}$ 是 LWE 密文中公开可见的已知值，可以直接计算 $V_0 = V \cdot X^{-\hat{b}}$ 的旋转。
\item 计算 $V_1 = V_0 \cdot X^{\hat{a}_0s_0} = s_0 \cdot (V_0 \cdot X^{\hat{a}_0} - V_0) + V_0$。
该公式适用于 $s_0 \in \{0, 1\}$ 这一特殊情形：若 $s_0 = 0$，则 $V_1 = V_0$；若 $s_0 = 1$，则 $V_1 = V_0 \cdot X^{\hat{a}_0}$。对 $s_0 \cdot (V_0 \cdot X^{\hat{a}_0} - V_0) + V_0$ 的计算通过 TFHE 同态加法和乘法实现。我们将其称为 $V_0$ 的盲旋转，其中选择位 $s_0$（即秘密密钥向量 $S$ 的第一个元素）使用 $\vec{S}_{bk}$（即自举密钥）加密为 GGSW 密文，而 $V_0$ 则是以 GLWE 密文形式加密的多项式。将 GLWE 加密的 $V_0$ 与 $x^{\hat{a}_0}$ 相乘，采用 GLWE 密文与明文乘法（\autoref{sec:glwe-mult-plain}）；从所得结果中减去 GLWE 加密的 $V_0$，采用 GLWE 密文与密文加减法（\autoref{sec:glwe-add-cipher}）；将结果乘以 GGSW 加密的 $s_0$，采用 GLWE 与 GGSW 乘法（\autoref{subsubsec:tfhe-glwe-to-ggsw-multiplication}）；最后加上 GLWE 加密的 $V_0$，采用 GLWE 密文与密文加法。若 $s_0 = 1$，公式中的 $X^{\hat{a}_0}$ 项会乘到 $V_0$ 上，等效于将 $V_0$ 的系数向右旋转 $\hat{a}_0$ 个位置。若 $s_0 = 0$，则 $V_0$ 不发生旋转，保持原样。无论哪种情况，所得 $V_1$ 都以新的 GLWE 密文形式加密。在这一盲旋转过程中，除非知道 $s_0$ 和 $\vec{S}_{bk}$，否则无法判断 $V_0$ 是否发生旋转，也无法知道旋转了多少个位置。
\item 使用与上一步相同的盲旋转方法，计算 $V_2 = V_1 \cdot X^{\hat{a}_1s_1} = s_1 \cdot (V_1 \cdot X^{\hat{a}}_1 - V_1) + V_1$ 的 GLWE 加密。注意，已知的公开信息包括 $\hat{a}_1$、$V_1$ 的 GLWE 加密，以及使用 $\vec{S}_{bk}$ 加密 $s_1$ 得到的 GGSW 密文。
\item 继续以相同方式计算 $V_3, V_4, \gap{$\cdots$}, V_k$ 的 GLWE 加密，最终得到 $V_k$ 的 GLWE 加密；其计算结果等价于：

$V' = V_k$

$\text{ } = V_{k-1} \cdot X^{\hat{a}_{k-1}s_{k-1}}$

$\text{ } = V_0 \cdot X^{\hat{a}_0s_0}X^{\hat{a}_1s_1}\cdots X^{\hat{a}_{k-1}s_{k-1}}$

$\text{ } = V \cdot X^{- \hat{b} + \sum_{i=0}^{k-1}{\hat{a}_is_i}} $

$\text{ } = V \cdot X^{-(\hat{\Delta} m + \hat{e}_b)}$

\end{enumerate}

这意味着，最终多项式 $V_k$ 的 GLWE 加密所对应的常数项系数为 $m$，因为 $V(X)$ 被设计为具有映射（$\hat\Delta m + \hat e_{*}, m$）。

注意，虽然应用所使用的明文空间被限制为 $\mathbb{Z}_t$ 内连续的 $\dfrac{t}{2}$ 个模剩余值（\autoref{subsec:tfhe-zero-padding}），但密文空间没有这一限制。也就是说，盲旋转的中间步骤允许将 $V(X)$ 旋转超过 $n$ 个位置，因为后续步骤可以将无效位置重新转回有效位置。因此，噪声自举的正确性只取决于完成后的状态 $V_k(X)$。旋转完成的 $V_k(X)$ 必须相当于向左旋转了 $\hat\Delta m + \hat e$ 个位置。因此，只须保证 $\hat\Delta m + \hat e$ 落在 $\mathbb{Z}_t$ 域内预先指定的连续 $\dfrac{t}{2}$ 个模剩余值范围中；这等价于确保累计旋转次数至多为 $n-1$ 个位置，以避免提取到符号相互矛盾的系数。

\begin{figure}[h!]
    \centering
  \includegraphics[width=0.2\linewidth]{figures/mux.pdf}
  \caption{MUX 逻辑门示意图}
  \label{fig:mux}
\end{figure}

\para{同态 MUX 逻辑门：} 在步骤 3 中，公式 $s_0 \cdot (V_0 \cdot X^{\hat{a}_0} - V_0) + V_0$ 实现了 \autoref{fig:mux} 所示的 MUX 逻辑门。本例中，$s_0$ 是选择位，用于在两个输入 $V_0$ 和 $V_0 \cdot X^{\hat{a}_0}$ 之间进行选择。如果 $s_0$ 为 1，则输出为 $V_0 \cdot X^{\hat{a}_0}$；否则输出为 $V_0$。在本设计中，同态计算 $s_0 \cdot (V_0 \cdot X^{\hat{a}_0} - V_0) + V_0$ 实际上实现了一个同态 MUX 门，其中两个输入采用 LWE 加密，选择位采用 GGSW 加密。

\para{GLWE 与 GGSW 密文的维数：} 加密 LUT 多项式 $V(X)$ 的 GLWE 密文维数为 $k' \times n$，其中 $k'$ 是 $\vec{A}$ 的长度，即公开掩码多项式的数量，$n$ 是 $A, B,$ 和 $V$ 的多项式次数。加密 $\vec{S}_{\textit{bk}}$ 的每个元素（即自举密钥多项式的每个系数）的 GGSW 密文维数为 $(k' + 1) \times l \times (k' + 1)$。实际中令 $k' = 1$，使这些密文简化为 RLWE 与 RGSW 密文。设置 $k' = 1$ 是出于计算效率的考虑。

\subsubsection{系数提取}

接下来，利用系数提取技术（\autoref{subsec:tfhe-extraction}），提取旋转后的多项式 $V'(X)$' 的常数项系数，得到 $m$ 的加密：$\textsf{LWE}_{\vec{s}', \sigma}(\Delta m)$，其中 $\vec{s}'$ 是长度为 $k'\cdot n$ 的向量。此时，原始密文 $\textsf{LWE}$ 的旧噪声 $\hat{e}_b$ 已经消失，而自举后的新密文 $\textsf{LWE}_{\vec{s}', \sigma}(\Delta m)$ 含有新生成的较小噪声 $e_s$。

事实上，盲旋转过程（\autoref{subsec:bootstrapping-blind-rotation}）中的同态 MUX 逻辑涉及大量密文乘法和加法，直到系数提取之前，这些运算都会累积额外噪声。为限制一系列 MUX 逻辑运算中累积的噪声，可以仔细调整安全参数。例如，可用较窄的高斯分布采样噪声 $e$，同时选取足够大的 $n$，以补偿减小噪声对安全性的影响。另一方面，增大 $q$ 会使系统更容易受到格规约攻击，从而降低安全性。


\subsubsection{密钥切换}
\label{subsec:bootstrapping-key-switch} 

系数提取的输出是 $\textsf{LWE}_{\vec{s}', \sigma}(\Delta m)$，它使用长度为 $k'\cdot n$ 的秘密密钥 $\vec{s}'$ 加密。为保持一致，需要执行 $\vec{s}' \rightarrow \vec{s}$ 的密钥切换，其中 $\vec{s}$ 是原始的长度为 $k$ 的密钥。可采用 \autoref{subsec:tfhe-key-switching} 中小结~\ref*{subsec:tfhe-key-switching} 所介绍的技术完成这一密钥切换。



\subsubsection{噪声自举小结}
\label{subsec:tfhe-summary}

TFHE 的噪声自举过程归纳如下：

\begin{tcolorbox}[title={\textbf{\tboxlabel{\ref*{subsec:tfhe-summary}} TFHE 噪声自举流程}}]

\textbf{\underline{查找表加密}：} 使用长度为 $k'$ 的自举密钥 $\vec{S}_{bk}$，将 LUT 多项式 $V(X)$ 加密为 GLWE 密文。

  
\begin{enumerate}
\item \textbf{\underline{模数切换}：} 对 TFHE 密文 $\textsf{LWE}_{\vec{s}, \sigma}(\Delta m + e_b)$ 执行 $q \rightarrow 2n$ 的模数切换，得到 $\textsf{LWE}_{\vec{s}, \sigma}(\hat{\Delta} m + \hat e_b)$，其中 $\hat\Delta = \Delta \cdot \dfrac{2n}{q} = \dfrac{2n}{t}$。

$ $

\item \textbf{\underline{盲旋转}：} 通过递归计算下式，将 GLWE 加密的多项式 $V$ 向左旋转 $(\hat{b} - \sum\limits_{i=0}^{k-1}{\hat{a}_is_i}) = (\hat{\Delta}m + \hat{e}_b)$ 个位置：

$V_0 = V \cdot X^{-\hat{b}}$

$V_1 = V_0 \cdot X^{\hat{a}_0s_0} = s_0 \cdot (V_0 \cdot X^{\hat{a}_0} - V_0) + V_0$

\text{ } $\vdots$

$V_k = V_{k-1} \cdot X^{\hat{a}_{k-1}s_{k-1}} = s_{k-1} \cdot (V_{k-1} \cdot X^{\hat{a}_{k-1}} - V_{k-1}) + V_{k-1}$

$= V_0 \cdot X^{\hat{a}_0s_0}X^{\hat{a}_1s_1}\cdots X^{\hat{a}_{k-1}s_{k-1}}$

$= V \cdot X^{- \hat{b} + \sum_{i=0}^{k-1}{\hat{a}_is_i}} $

$ = V \cdot X^{-(\hat{\Delta}m + \hat{e}_b)}$

$ $

\noindent 上述盲旋转的每个实际步骤，都通过以下 TFHE 密文与密文乘法和加法计算：

$\textsf{GLWE}_{\vec{S}_{bk}, \sigma}(V_{i+1}) = \textsf{GGSW}_{\vec{S}_{bk}, \sigma}^{\beta, l}(s_i) \cdot (\textsf{GLWE}_{\vec{S}_{bk}, \sigma}(V_i) \cdot X^{\hat{a}_i} - \textsf{GLWE}_{\vec{S}_{bk}, \sigma}(V_i)) + \textsf{GLWE}_{\vec{S}_{bk}, \sigma}(V_i)$

$ $

\item \textbf{\underline{系数提取}：} 从旋转后的多项式 $V_k$ 中同态提取常数项系数 $m$，得到 $\textsf{LWE}_{\vec{s}', \sigma}(\Delta m)$，其中 $\vec{s}'$ 是长度为 $k' \cdot n$ 的向量。

$ $

\item \textbf{\underline{密钥切换}：} 同态执行 LWE 密文的 $\vec{s}' \rightarrow \vec{s}$ 密钥切换。

%\item \textbf{Rescaling:} Multiply $\textsf{LWE}_{\vec{s}, \sigma}(\hat{\Delta}m)$ by $\dfrac{q}{2n}$ to get:  

%$\dfrac{q}{2n} \cdot LWE_{S, \sigma}(\hat{\Delta}m)$

%$\textsf{LWE}_{\vec{s}, \sigma}(\dfrac{q}{2n} \cdot \hat{\Delta}m)$

%$ = LWE_{S, \sigma}(\Delta m)$

%, whose noise get reinitialized to $e_s$. 
\end{enumerate}

$ $

\para{\underline{将可用明文范围减半}：} LUT 多项式 V 的负循环旋转会带来问题。为避免该问题，要求应用确保明文 $m$ 只使用 $\mathbb{Z}_t$ 内连续的 $\dfrac{t}{2}$ 个模剩余值。这样可以防止 $V(X)$ 旋转超过 $n-1$ 个位置，从而避免系数提取得到符号相互矛盾的系数。

$ $

\para{\underline{$k'$ 的选择}：} 为提高计算效率，将 $k'$ 设置为 1，使 GLWE 与 GGSW 密文简化为 RLWE 与 RGSW 密文。
%\item Piping the LUT output to a \textbf{negacyclic function}.
%\item Multiplying the LUT output with \textbf{another negacyclically rotating polynomial}. 
%\end{itemize}
\end{tcolorbox}


\subsubsection{示例：噪声自举}
\label{subsec:tfhe-noise-bootstrapping-ex}

假设 GLWE 的安全参数设置如下：$n = 16$, $t = 8$, $q = 64$, $k = 8$

$\mathbb{Z}_{t=8} = \{-4, -3, -2, -1, 0 , 1, 2, 3\}$

$\mathbb{Z}_{q=64} = \{ -32, -31, -30, \gap{$\cdots$}, 29, 30, 31 \}$

$\Delta = \dfrac{q}{t} = \dfrac{64}{8} = 8$

$ $

并假设有如下 LWE 密文：

$\vec{s} = (1, 0, 0, 1, 1, 1, 0, 1) = \mathbb{Z}_2^{k=8}$

$m = 1  \in \mathbb{Z}_{t=8}$

$\Delta m = 1 \cdot 8 = 8 \in \mathbb{Z}_{q=64}$



$\textsf{LWE}_{\vec{s}, \sigma}(\Delta m) = (a_0, a_1, a_2, a_3, a_4, a_5, a_6, a_7, b) 
= (8, -28, 4, -32, 0, 31, -6, 7, 24) \in \mathbb{Z}_{q=64}^{k+1=9}$

$e = 2 \in \mathbb{Z}_{q=64}$（为保证正确解密，应满足 $|e| < \dfrac{\Delta}{2} = 4$）

$b = \sum\limits_{i=0}^{7}a_is_i + \Delta m + e = (8 -32 + 31 +7) + 8 + 2 = 24 \in \mathbb{Z}_{q=64}$

$ $

此时，TFHE 噪声自举过程如下：

$ $

\begin{enumerate}
\item \textbf{模数切换：} 对 $\textsf{LWE}_{\vec{s}, \sigma}(\Delta m)$ 执行 $q \rightarrow 2n$ 的模数切换，即 $64 \rightarrow 32$。模数切换后，原始 LWE 密文变为：

$\mathbb{Z}_{2n=32} = \{-16, -15, -14, \gap{$\cdots$}, 13, 14, 15\}$

$\vec{s} = (1, 0, 0, 1, 1, 1, 0, 1) = \mathbb{Z}_2^{k=8}$


$\hat{\Delta} = \Delta\dfrac{2n}{64} = 8\dfrac{32}{64} = 4$

$\hat{\Delta}m = 4 \cdot 1 = 4 \in \mathbb{Z}_{2n=32}$

$\hat{e} =  \left\lceil e \dfrac{2n}{q} \right\rfloor = \left\lceil 2\dfrac{32}{64} \right\rfloor = 1 \in \mathbb{Z}_{2n=32}$

$ $

$\textsf{LWE}_{\vec{s}, \sigma}(\hat{\Delta}m) = (\hat{a}_0, \hat{a}_1, \hat{a}_2, \hat{a}_3, \hat{a}_4, \hat{a}_5, \hat{a}_6, \hat{a}_7, \hat{b}) \in \mathbb{Z}_{2n=32}^{k+1=9}$


$ = (\Big\lceil 8\dfrac{32}{64} \Big\rfloor, \Big\lceil -28\dfrac{32}{64}\Big\rfloor, \Big\lceil 4\dfrac{32}{64}\Big\rfloor, \Big\lceil -32\dfrac{32}{64}\Big\rfloor, \Big\lceil 0\dfrac{32}{64}\Big\rfloor, \Big\lceil 31\dfrac{32}{64}\Big\rfloor, \Big\lceil -6\dfrac{32}{64}\Big\rfloor, \Big\lceil 7\dfrac{32}{64}\Big\rfloor, \Big\lceil 24\dfrac{32}{64}\Big\rfloor)$

$ $

$= (4, {-14}, 2, -16, 0, 16, {-3}, 4, 12)$

$ $

注意 $\sum\limits_{i=0}^{7}(\hat{a}_is_i) + \hat{\Delta}m + \hat{e} = (4 - 16 + 16 + 4) + 4 + 1 = 13 \in \mathbb{Z}_{2n=32}$

$ $

$\hat b = 12 \approx 13 = \sum\limits_{i=0}^{7}(\hat{a}_is_i) + \hat{\Delta}m + \hat{e}$

$\hat b$ 中的这一微小差异，来自模数切换时对 $\hat a_0, \hat a_1, \cdots , \hat e$ 取整产生的累积噪声。

$ $

\item \textbf{盲旋转：} 假设应用通过限制可用明文值来避免多项式负循环旋转的问题：仅使用 $\mathbb{Z}_8 = \{-4, -3, -2, -1, 0, 1, 2, 3\}$ 内以下连续的 $\dfrac{8}{2}$ 个模剩余值，即 $\{-2, -1, 0, 1\}$。这意味着 $V(X)\cdot X^{i}$ 中的 $i=\hat\Delta m + \hat e$ 只能取以下值：$i = \{-8, -7, \cdots, 6, 7\}$。根据这些要求，\autoref{tab:lut} 给出了将 $\hat{\Delta} m + \hat{e}$ 映射到 $\Delta m$ 的查找表多项式 $V(X)$。


\begin{table}[h]
\centering
\footnotesize
%\noindent\adjustbox{max width=\columnwidth}{
\begin{tabular}{|c||c|c|c|c|c|c|c|c|c|c|} % left align
\hline
% remove line in a table, \usepackage{adjustbox}
\multicolumn{10}{|l|}{$V(X) = v_0 + v_1X + v_2X^2 + v_3X^3 + v_4X^4 + v_5X^5 + v_6X^6 + v_7X^7$}  \\ 
\multicolumn{10}{|l|}{\textcolor{white}{......} $ + v_8X^8 + v_9X^9 + v_{10}X^{10} + v_{11}X^{11} + v_{12}X^{12} + v_{13}X^{13} + v_{14}X^{14} + v_{15}X^{15}$}  \\ 
\multicolumn{10}{|l|}{\textcolor{white}{....}$= 0 + 0X +0X^2 +0X^3 +1X^4 +1X^5 +1X^6 +1X^7$}  \\ 
\multicolumn{10}{|l|}{\textcolor{white}{......}$+ 2X^8  + 2X^9 + 2X^{10} + 2X^{11} + 1X^{12} + 1X^{13} + 1X^{14} + 1X^{15}$} \\ 
\hline
\hline
\textbf{\boldmath$i = \hat{\Delta}m + \hat{e}$} & $-8$ & $-7$ & $-6$ & $-5$ & $-4$ & $-3$ & $-2$ & $-1$ \\
(in $V \cdot X^{-i}$) & ($\textcolor{orange}{110}\textcolor{green}{00}_2$) & ($\textcolor{orange}{110}\textcolor{green}{01}_2$) & ($\textcolor{orange}{110}\textcolor{green}{10}_2$) & ($\textcolor{orange}{110}\textcolor{green}{11}_2$) & ($\textcolor{orange}{111}\textcolor{green}{00}_2$) & ($\textcolor{orange}{111}\textcolor{green}{01}_2$) & ($\textcolor{orange}{111}\textcolor{green}{10}_2$) & ($\textcolor{orange}{111}\textcolor{green}{11}_2$)\\
\hline
\textbf{常数项的} & $-2$ & $-2$ & $-2$ & $-2$ & $-1$ & $-1$ & $-1$ & $-1$ \\
 \textbf{系数（$V\cdot X^{-i}$）}& $\textcolor{orange}{110}_2$ & $\textcolor{orange}{110}_2$ & $\textcolor{orange}{110}_2$ & $\textcolor{orange}{110}_2$ & $\textcolor{orange}{111}_2$ & $\textcolor{orange}{111}_2$ & $\textcolor{orange}{111}_2$ & $\textcolor{orange}{111}\textcolor{green}{00}_2$ \\
\hline
\textbf{$\bm{m}$（明文）} & $-2$ & $-2$ & $-2$ & $-2$ & $-1$ & $-1$ & $-1$ & $-1$ \\
\hline
\hline
\textbf{\boldmath$i = \hat{\Delta}m + \hat{e}$} & $0$ & $1$ & $2$ & $3$ & $4$ & $5$ & $6$ & $7$ \\
(in $V \cdot X^{-i}$) & ($\textcolor{orange}{000}_2$)& ($\textcolor{orange}{000}_2$)& ($\textcolor{orange}{000}_2$)& ($\textcolor{orange}{000}_2$)& ($\textcolor{orange}{001}_2$)& ($\textcolor{orange}{001}_2$)& ($\textcolor{orange}{001}_2$)&($\textcolor{orange}{001}_2$)\\
\hline
\textbf{常数项的} & $0$ & $0$ & $0$ & $0$ & $1$ & $1$ & $1$ & $1$ \\
\textbf{系数（$V\cdot X^{-i}$）}& $\textcolor{orange}{000}\textcolor{green}{00}_2$ & $\textcolor{orange}{000}\textcolor{green}{00}_2$ & $\textcolor{orange}{000}\textcolor{green}{00}_2$ & $\textcolor{orange}{000}\textcolor{green}{00}_2$ & $\textcolor{orange}{001}\textcolor{green}{00}_2$ & $\textcolor{orange}{001}\textcolor{green}{00}_2$ & $\textcolor{orange}{001}\textcolor{green}{00}_2$ & $\textcolor{orange}{001}\textcolor{green}{00}_2$ \\
\hline
\textbf{$\bm{m}$（明文）} & $0$ & $0$ & $0$ & $0$ & $1$ & $1$ & $1$ & $1$ \\
\hline
\end{tabular}
\centering
\caption{LWE 参数设置为 $n=16, q=64, t=8$ 时的查找表。
\textcolor{orange}{橙色}表示明文 $m$ 的比特。\textcolor{green}{绿色}表示噪声 $e$ 的比特。 %\textcolor{cyan}{Cyan} is the 0-padding bit at the MSB of $\hat\Delta m + e$ (and thus the MSB of $m$), ensured by the application's usage of plaintext numbers.
}
\label{tab:lut}
\end{table}

注意，$V(X)$ 中 $X^8 \sim X^{15}$ 项的系数取 $\{2, 1\}$，而不是 $\{-2, -1\}$。这样，当 $V$ 向左旋转 $\{-8,-7,-6,-5,4,5,6,7\}$ 个位置时，由于越过指数 $n$ 的边界而发生回绕，常数项系数的符号会翻转为 $\{-2, -1\}$。

实际自举时，将 \autoref{tab:lut} 中的 $V(X)$（它已被 GLWE 加密）向左盲旋转 $\hat{b} - \sum\limits_{i=0}^{7}\hat{a}_is_i = 4$ 个位置，计算如下：

$\hat\Delta m + \hat e = \hat{b} - \sum\limits_{i=0}^{7}\hat{a}_is_i = 12 - (4 - 16 + 16 + 4) = 4 \text{ mod 32} \in \mathbb{Z}_{2n=32}$ 


在 \autoref{tab:lut} 中，如果旋转次数 $i = 4$，则对应的常数项系数为 $v_4 = 1 = m$。由于 $\Delta = 4$，最终得到 $\textsf{LWE}_{\vec{s},\sigma}(\Delta m) = 1$。



$ $

实际的盲旋转计算如下：

$\vec{s} = (1, 0, 0, 1, 1, 1, 0, 1)$

$\textsf{LWE}_{\vec{s}, \sigma}(\hat{\Delta}m) = (\hat{a}_0, \hat{a}_1, \hat{a}_2, \hat{a}_3, \hat{a}_4, \hat{a}_5, \hat{a}_6, \hat{a}_7, \hat{b}) =  (4, -14, 2, -16, 0, 16, -3, 4, 12)$

$V_0 = V \cdot X^{-\hat{b}} = V \cdot X^{-12} = v_{12} + v_{13}X + v_{14}X^2 + \cdots $

$V_1 = V_0 \cdot X^{\hat{a}_0s_0} = s_0 \cdot (V_0 \cdot X^{\hat{a}_0} - V_0) + V_0 = V_0 \cdot X^{4} = v_{8} + v_{9}X + v_{10}X^2 + \cdots$ 

$V_2 = V_1 \cdot X^{\hat{a}_1s_1} = s_1 \cdot (V_1 \cdot X^{\hat{a}_1} - V_1) + V_1 = V_1 = v_{8} + v_{9}X + v_{10}X^2 + \cdots$

$V_3 = V_2 \cdot X^{\hat{a}_2s_2} = s_2 \cdot (V_2 \cdot X^{\hat{a}_2} - V_2) + V_2 = V_2 = v_{8} + v_{9}X + v_{10}X^2 + \cdots$

$V_4 = V_3 \cdot X^{\hat{a}_3s_3} = s_3 \cdot (V_3 \cdot X^{\hat{a}_3} - V_3) + V_3 = V_3 \cdot X^{-16} = -v_{8} - v_{9}X - v_{10}X^2 - \cdots$

$V_5 = V_4 \cdot X^{\hat{a}_4s_4} = s_4 \cdot (V_4 \cdot X^{\hat{a}_4} - V_4) + V_4 = V_4 \cdot X^{0} = -v_{8} - v_{9}X - v_{10}X^2 - \cdots$

$V_6 = V_5 \cdot X^{\hat{a}_5s_5} = s_5 \cdot (V_5 \cdot X^{\hat{a}_5} - V_5) + V_5 = V_5 \cdot X^{16} = v_{8} + v_{9}X + v_{10}X^2 + \cdots$

$V_7 = V_6 \cdot X^{\hat{a}_6s_6} = s_6 \cdot (V_6 \cdot X^{\hat{a}_6} - V_6) + V_6 = V_6 = v_{8} + v_{9}X + v_{10}X^2 + \cdots$

$V_8 = V_7 \cdot X^{\hat{a}_7s_7} = s_7 \cdot (V_7 \cdot X^{\hat{a}_7} - V_7) + V_7 = V_7 \cdot X^{4} = v_{4} + v_5X + v_6X^2 + \cdots $

$ $

盲旋转的最终输出是 $V_8$ 的 GLWE 密文 $\textsf{GLWE}_{\vec{S}, \sigma}(V_8)$，其常数项系数为 $v_4 = m = 1$。

$ $

上述盲旋转的每一步都通过以下 TFHE 密文与密文乘法计算：

$\textsf{GLWE}_{\vec{S}, \sigma}(V_{i+1}) = \textsf{GGSW}_{\vec{S}, \sigma}^{\beta, l}(s_i) \cdot (\textsf{GLWE}_{\vec{S}, \sigma}(V_i) \cdot X^{\hat{a}_i} - \textsf{GLWE}_{\vec{S}, \sigma}(V_i)) + \textsf{GLWE}_{\vec{S}, \sigma}(V_i)$

$ $

这一计算留给读者作为练习。

$ $

\item \textbf{系数提取：} 盲旋转结束时，最终得到如下 GLWE 密文：

$\textsf{GLWE}_{\vec{S}, \sigma}(V_8)$

$ = \textsf{GLWE}_{\vec{S}, \sigma}\bm(\hat{\Delta} \cdot (v_4 + v_5X + v_6X^2 + v_7X^3 + v_8X^{4} + v_9X^{5} + v_{10}X^{6} + v_{11}X^{7} + v_{12}X^{8} + v_{13}X^{9} + v_{14}X^{10} + v_{15}X^{11} - v_{0}X^{12} - v_{1}X^{13} - v_{2}X^{14} - v_{3}X^{15})\bm)$

$ = \textsf{GLWE}_{\vec{S}, \sigma}\bm(\hat{\Delta}\cdot(1 + 1X + 1X^2 + 1X^3 + 2X^{4} + 2X^{5} + 2X^{6} + 2X^{7} + 1X^{8} + 1X^{9} + 1X^{10} + 1X^{11} - 0X^{12} - 0X^{13} - 0X^{14} - 0X^{15})\bm)$

$ = \left(A_0 = \sum\limits_{j=0}^{15}(a_{0,0} + a_{0,1}X + \cdots), A_1 = \cdots, A_{k-1} = \cdots, B = \sum\limits_{j=0}^{15}b_{j}X^j\right)$

现在，利用系数提取公式（小结~\ref{subsec:tfhe-extraction}），提取加密多项式 $\textsf{GLWE}_{\vec{S}, \sigma}(\hat{\Delta} \cdot (1 + 1X + 1X^2 + \cdots))$ 的常数项系数。具体而言，提取的是与 $\textsf{LWE}_{\vec{s}, \sigma}(\Delta m_0)$ 对应的常数项系数。通过以下计算提取 $\textsf{LWE}_{\vec{s}, \sigma}(\Delta m_0)$：

$\textsf{LWE}_{\vec{s}, \sigma}(\Delta m_0) = (a_0', a_1', \gap{$\cdots$} , a_{nk-1}', b_h)$ \textcolor{red}{ $\rhd$ 其中 $h = 0$}

\[
    \text{，其中 } a'_{n \cdot i + j} =   
    \begin{cases}
      a_{i,0 - j} \text{（若 } 0 \leq j \leq 0\text{）}\\
      -a_{i,n + 0 - j} \text{（若 } 0+1 \leq j \leq n-1\text{）}\\
  \end{cases}
, b_0 \text{ 取自多项式 } B
\]


\end{enumerate}


\subsubsection{讨论}
\label{subsec:tfhe-noise-bootstrapping-discussion}

\begin{itemize}
\item \textbf{可编程自举}：前述自举（\autoref{subsec:tfhe-noise-bootstrapping}）使用一个简单查找表 $V(X)$，将 $\Delta m + e$ 映射到 $\Delta m$；而通过修改 $V(X)$ 的系数，可以让 $\Delta m + e$ 映射到不同的值。例如，改变 LUT 输入与输出之间的映射，可以实现 AND、OR、XOR、CMUX 等逻辑门，相关内容见 \autoref{subsec:tfhe-noise-bootstrapping-gate}。这种修改 $V(X)$ 中指数与系数之间映射的自举称为可编程自举。如果将 $V(X)$ 加密为 GLWE 密文，就能同时隐藏映射关系及每个输入实例，从而实现\textit{函数加密}。注意，普通自举（\autoref{subsec:tfhe-noise-bootstrapping}）和可编程自举（\autoref{subsec:tfhe-noise-bootstrapping-gate}）产生的噪声量相同。
\item \textbf{自举噪声}：在自举过程中旋转 LUT 多项式 $V(X)$ 时，需要在同态 MUX 门中执行许多 TFHE 乘法以得到 $V_0 \cdots V_k$，因此，在最终完成噪声重新初始化之前，不可避免地会产生额外噪声。不过，仔细选择参数可以限制模数切换和盲旋转期间这些额外噪声的增长。
\end{itemize}


\subsubsection{应用：门自举}
\label{subsec:tfhe-noise-bootstrapping-gate}

除了实现盲旋转中使用的同态 MUX 逻辑门（\autoref{subsec:bootstrapping-blind-rotation}），还可以利用 LUT 多项式 $V(X)$ 实现 AND、NAND、OR、XOR 等其他同态逻辑门。实现这些门时，每个密文加密一个单比特明文；也可以将若干比特通过线性组合公式合并，再用 LUT 同时处理。假设 $q = 32$、$t = 8$、$m \in \mathbb{Z}_8 = \{-4,-3, -2,-1,0,1,2,3\}$、$\Delta = \dfrac{q}{t} = 4$、$\hat\Delta = \dfrac{\Delta\cdot 2n}{q} = 2$，并将逻辑门输入编码为 LWE 明文 $0 \rightarrow -1$ 和 $1 \rightarrow 1$，且最大累积噪声为 $e = [-1, 1]$。


\begin{table}[ht]
  \centering
  \resizebox{\columnwidth}{!}{%
    \begin{tabular}{c c | r r | c | c | c | c}
      \hline
        \multicolumn{8}{c}{查找表多项式 $V(X) = 1 + 1X + 1X^2 + 1X^3 + 1X^4 + 1X^5 + 1X^6 + 1X^7$}\\
      \hline
      \multicolumn{2}{c|}{\textbf{解码值}} & 
      \multicolumn{2}{c|}{\textbf{编码值}} & 
      \multicolumn{1}{c|}{\begin{tabular}{c}
        \textbf{编码的} \\ \textbf{线性} \\
        \textbf{组合}
      \end{tabular}} & 
      \multicolumn{1}{c|}{\begin{tabular}{c}
        \textbf{放大后的} \textbf{编码} \\
         \textbf{组合}
      \end{tabular}} &
      \textbf{自举} & 
      \begin{tabular}{c}
        \textbf{解码} \\ \textbf{结果} \\
      \end{tabular} \\
        \hline
      $d_1$ & $d_2$ & $m_1$ & $m_2$ & $m_1 + m_2 - 1$ & $\hat\Delta\cdot(m_1+m_2-1)$ &
      $V(X)\cdot X^{-\hat\Delta\cdot(m_1 + m_2 - 1)+e}$ &   $d_1 \land d_2$ \\
      \hline
      0 & 0 &
      $-1$ & $-1$ &
      $-3$ &
      $-6$ &
      常数项系数为 $-1$ &
      0 \\
      0 & 1 & 
      $-1$ & $1$ &
      $-1$ &
      $-2$ &
      常数项系数为 $-1$ &
      0 \\
      1 & 0 & 
      $1$ & $-1$ &
      $-1$ &
      $-2$ &
      常数项系数为 $-1$ &
      0 \\
      1 & 1 & 
      $1$ & $1$ &
      $1$ &
      $2$ &
      常数项系数为 $1$ &
      1 \\
      \hline
    \end{tabular}
  }
  \caption{采用额外编码的 AND 运算真值表示例。}
  \label{tab:gate-and}
\end{table}


%\begin{figure}[h!]
%    \centering
%  \includegraphics[width=\linewidth]{figures/torus-gate.png}
%  \caption{An illustration of an AND gate in Torus \href{https://assets-global.website-files.com/622ef9de9152c97467eac748/6298a348a0c75850f0688c83_NAND_bootstrap.png}{(Source)}}
%  \label{fig:torus-gate}
%\end{figure}


\autoref{tab:gate-and} 展示了 AND 逻辑门的可编程自举设计。在这一应用中，将 LUT 多项式 $V$ 定义为 $V(X) = \sum\limits_{i=0}^{7}X^i$。LUT 多项式 $V(X)$ 将一半明文域映射为 $1$，另一半映射为 $-1$（各项越过 $X^7$ 边界时会发生回绕）。在此设计中，每个比特都单独加密为一个独立的 TFHE 密文。门输入 0 和 1 分别编码为 $-1$ 和 $1$。AND 门的线性组合，即同态计算公式，为 $\textsf{LWE}_{\vec{s}, \sigma}(\Delta m_1) + \textsf{LWE}_{\vec{s}, \sigma}(\Delta m_2) - 1$。只有两个输入均为正时，输出才为正（即 $1$）；此时盲旋转将 $V$ 向左旋转 $\hat\Delta\cdot 1 + e$ 个位置，常数项系数为 $1$。因此，盲旋转和系数提取将输出噪声降低后的 $\textsf{LWE}_{\vec{s}, \sigma}(\Delta \cdot 1)$，它是 $1$ 的编码。该设计能够容忍的最大噪声为 $|e| = 1$。若要容忍更大的噪声，应增大 $q$ 和 $n$。

注意，AND 门的 LUT 布局本身具有负循环性质，这是一个特殊情形，因此可以利用 $V(X)$ 固有的负循环旋转性质，将 $V(X)$ 中全部 $2n=16$ 种系数状态用于 AND 门映射函数的输出。然而，在许多应用中，LUT 布局并不一定像该 AND 门示例一样具有负循环性质。前述噪声自举的 LUT 布局（\autoref{subsec:bootstrapping-overview}）也不是负循环函数，而是恒等函数（因为它仅仅移除噪声）。因此，对于大多数应用，需要只使用 $t$ 明文空间中的 $\dfrac{t}{2}$ 部分，以避免 $V(X)$ 的旋转超过 $n-1$ 个位置（\autoref{subsec:tfhe-zero-padding}）。

除 AND 门之外，其他逻辑门也可以按类似方式构造，每种门采用不同的线性组合公式和 LUT 布局。

\para{除法：} TFHE 不支持任意大小明文数的直接除法。这是因为 TFHE 的 LWE 向量元素属于环 $Z_q$，其中每个元素 $g$ 不一定存在乘法逆元 $g^{-1}$，因而难以通过将目标数乘以 $g^{-1}$ 来实现除法。可以改用门自举所实现的逻辑门来完成二进制除法。为支持二进制除法，每个明文必须是一个比特，并被加密为独立的密文。也可以设计线性组合公式，将多个比特合并并同时处理，类似于之前为处理 AND 门的两个输入比特所设计的线性组合。

\subsubsection{应用：神经网络中的自举}
\label{subsec:tfhe-neural-network}

\begin{figure}[h!]
    \centering
  \includegraphics[width=0.8\linewidth]{figures/neural-network.pdf}
  \caption{神经网络示意图}
  \label{fig:neural-network}
\end{figure}

\begin{figure}[h!]
    \centering
  \includegraphics[width=0.8\linewidth]{figures/nn-homomorphic.pdf}
  \caption{神经网络中的可编程自举示意图}
  \label{fig:neural-network2}
\end{figure}

同态加密可以应用于深度神经网络中的神经元。通常，每个神经元的计算包含两个步骤：

\begin{enumerate}
  \item \textbf{输入值的线性组合：} 给定输入特征值（或中间值）集合 $(x_1, x_2, \gap{$\cdots$}, x_n)$、权重集合 $(w_1, w_2, \gap{$\cdots$}, w_n)$ 和偏置 $b$，计算：$y = \sum\limits_{i=1}^{n}x_iw_i + b$。
\item \textbf{激活函数：} 计算 $f(y)$，其中 $f$ 是非线性激活函数，例如 $\sin$ 函数、ReLU、sigmoid、双曲正切等。
\end{enumerate}

TFHE 可以将第一步的线性组合公式 $y = \sum\limits_{i=1}^{n}x_iw_i + b$ 同态计算为 $\sum\limits_{i=1}^{n}\textsf{LWE}_{\vec{s}, \sigma}(x_i) \cdot w_i + b$，其实现使用密文加法（\autoref{sec:glwe-add-cipher}）以及密文与明文乘法（\autoref{sec:glwe-mult-plain}）。

然而，第二步的非线性函数无法直接表示为密文的加法和乘法。为解决这一问题，可以通过可编程自举求值激活函数，使自举的输出与激活函数的输出相同或接近。如果在第二步使用自举，则可以在每个神经元计算结束时刷新噪声，因此原则上能够处理任意深度的神经网络，而无须担心噪声持续增长。



\subsection{离散环面上的 TFHE}
\label{subsec:torus}

\textbf{- 参考资料：}
\href{https://eprint.iacr.org/2021/1402.pdf}{（离散化）环面上的
全同态加密
指南}~\cite{torus}

$ $

环面 $\mathbb{T}$ 是介于 0 和 1 之间并按模 1 回绕的连续实数域，即 $[0, 1)$。

离散环面 $\mathbb{T}_t$ 是如下有限实数集合：$\Big(0, \dfrac{1}{t}, \dfrac{2}{t}, \gap{$\cdots$}, \dfrac{t - 1}{t}\Big)$


前面各小节介绍的 TFHE 方案采用如下设置：

$ $

$m \in \mathbb{Z}_t$ 
$\vec{s} = \{0, 1\}^{k}$
$e \in \mathbb{Z}_q$

$\textsf{LWE}_{\vec{s}, \sigma}(\Delta m) = (a_0, a_1, \gap{$\cdots$}, a_k, b) \in \mathbb{Z}_t^{k + 1}$

$ $

而原始 TFHE 方案是基于离散环面设计的：


$ $

$m \in \mathbb{T}_{t}, \text{ } \vec{s} \in \{0, 1\}^{k}, \text{ } e \in \mathbb{T}_{q}$

$ $

$\textsf{LWE}_{\vec{s}, \sigma}(m) = \textsf{ct} = (a_0, a_1, \gap{$\cdots$}, a_k, b) \in \mathbb{T}_{q}^{k+1}$

$b = \sum\limits_{i=0}^{k}(a_is_i) + m + e \in \mathbb{T}_q$

$\textsf{LWE}_{\vec{s}, \sigma}^{-1}(\textsf{ct}) =  \Big\lceil b - \sum\limits_{i=0}^{k-1}(a_is_i) \Big\rfloor_{\frac{1}{t}} = \Big\lceil m + e \Big\rfloor_{\frac{1}{t}} = m$，其中给定 $e < \dfrac{1}{2t}$

\textcolor{red}{ $\rhd$ 其中 $\lceil x \rfloor_{\frac{1}{t}}$ 表示将 $x$ 舍入为最接近的 $\dfrac{1}{t}$ 的整数倍}

$ $

原始 TFHE 的区别在于，所有值（多项式系数或向量元素）都在模 1 的浮点数域（即 $[0, 1)$）中计算，而不是在大整数域（即 $[0, q)$）中计算。这意味着明文也必须编码为 $[0, 1)$ 中的值，而不是 $[0, q)$ 中的整数。注意，原始 TFHE 方案不需要尺度因子 $\Delta$，因为环面 $[0, 1)$ 的连续域提供了可区分 $q$ 个离散小数值的浮点精度；只要 $E$ 中的每个系数（或向量元素）$e_i$ 小于 $\dfrac{1}{2t}$，解密过程就可以成功消除噪声 $E$。

基于环面和基于整数环的 TFHE 方案均建立在相同的基本原理之上。



\newpage

\section{BFV 方案}
\label{sec:bfv}
BFV 支持在整数明文模数下对明文多项式进行精确运算。与 CKKS 的近似运算不同，只要每一步均满足 BFV 正确性界，其解码结果就是精确的；仅仅采用整数编码，并不能保证任意噪声或电路深度下都能正确解密。应用还必须选择能够保持预期整数语义的明文模数与编码。本章以 Fan 和 Vercauteren 的原始方案~\cite{cryptoeprint:2012/144} 为参考。

在 BFV 中，每个明文都被加密为一个 RLWE 密文。因此，BFV 的密文与密文加法、密文与明文加法，以及密文与明文乘法，均基于 GLWE 的同态加法与乘法（见 \autoref{part:generic-fhe}）实现；取 $k = 1$ 即可将 GLWE 化为 RLWE。




\begin{tcolorbox}[
    title = \textbf{所需背景知识},    % box title
    colback = white,    % light background; tweak to taste
    colframe = black,  % frame colour
    boxrule = 0.8pt,     % line thickness
    left = 1mm, right = 1mm, top = 1mm, bottom = 1mm % inner padding
]

\begin{itemize}
\item \autoref{sec:modulo}: \nameref{sec:modulo}
\item \autoref{sec:group}: \nameref{sec:group}
\item \autoref{sec:field}: \nameref{sec:field}
\item \autoref{sec:order}: \nameref{sec:order}
\item \autoref{sec:polynomial-ring}: \nameref{sec:polynomial-ring}
\item \autoref{sec:decomp}: \nameref{sec:decomp}
\item \autoref{sec:roots}: \nameref{sec:roots}
\item \autoref{sec:cyclotomic}: \nameref{sec:cyclotomic}
\item \autoref{sec:cyclotomic-polynomial-integer-ring}: \nameref{sec:cyclotomic-polynomial-integer-ring}
\item \autoref{sec:matrix}: \nameref{sec:matrix}
\item \autoref{sec:euler}: \nameref{sec:euler}
\item \autoref{sec:modulus-rescaling}: \nameref{sec:modulus-rescaling}
\item \autoref{sec:chinese-remainder}: \nameref{sec:chinese-remainder}
\item \autoref{sec:polynomial-interpolation}: \nameref{sec:polynomial-interpolation}
\item \autoref{sec:ntt}: \nameref{sec:ntt}
\item \autoref{sec:lattice}: \nameref{sec:lattice}
\item \autoref{sec:rlwe}: \nameref{sec:rlwe}
\item \autoref{sec:glwe}: \nameref{sec:glwe}
\item \autoref{sec:glwe-add-cipher}: \nameref{sec:glwe-add-cipher}
\item \autoref{sec:glwe-add-plain}: \nameref{sec:glwe-add-plain}
\item \autoref{sec:glwe-mult-plain}: \nameref{sec:glwe-mult-plain}
\item \autoref{subsec:modulus-switch-rlwe}: \nameref{subsec:modulus-switch-rlwe}
\item \autoref{sec:glwe-key-switching}: \nameref{sec:glwe-key-switching}
\end{itemize}
\end{tcolorbox}

\clearpage

\subsection{单值编码}
\label{subsec:bfv-single-encoding}

这种二进制多项式编码要求 $0\leq m<2^n$。只有选定的整数系数代表没有模 $t$ 回绕，且乘积次数严格小于 $n$ 时，在 $X=2$ 处求值才保持整数加法和乘法。它不是从 $\mathbb Z_t[X]/(X^n+1)$ 到无界整数的环同态。例如，取 $n=2$ 时，整数 $2$ 的编码 $X$ 在商环中平方得到 $-1$，而不是整数 $4$。

BFV 支持两种编码方案：单值编码和批量编码。本小节介绍单值编码。



\begin{tcolorbox}[title={\textbf{\tboxlabel{\ref*{subsec:bfv-single-encoding}} BFV 编码}}]

\textbf{\underline{输入整数}：}将满足 $0\leq m<2^n$ 的整数 $m$ 分解如下：

$m = b_{n-1}\cdot 2^{n-1} + b_{n-2}\cdot 2^{n-2} + \cdots + b_1\cdot 2^1 + b_0\cdot 2^0$  \text { } ，其中每个 $b_i \in \{0, 1\}$

$ $


\textbf{\underline{编码后的多项式}：} $M(X) = b_0 + b_1X + b_2X^2 + \cdots + b_{n-1}X^{n-1} \in \mathcal{R}_{\langle n, t\rangle}$

$ $

\textbf{\underline{解码}：} $M(X=2) = m$

\end{tcolorbox}


下面分析小结~\ref*{subsec:bfv-single-encoding} 中的编码方案，是否能保证加法和乘法之后仍能正确解码。这等价于证明：

$\text{Decode}(\sigma(m_1) + \sigma(m_2)) = m_1 + m_2$ \textcolor{red}{$\rhd$ 其中 $\sigma$ 为编码函数}

$\text{Decode}(\sigma(m_1) \cdot \sigma(m_2)) = m_1 \cdot m_2$

设 $\sigma(m_1) = M_1(X)$ 且 $\sigma(m_2) = M_2(X)$，则

$\text{Decode}(\sigma(m_1) + \sigma(m_2)) = \text{Decode}(M_1(X) + M_2(X))$


$= \text{Decode}(M_{1+2}(X))$ \textcolor{red}{$\rhd$ 其中 $M_{1+2}(X) = M_1(X) + M_2(X)$}

$= M_{1+2}(2)$ \textcolor{red}{$\rhd$ 因为解码就是在 $X = 2$ 处对多项式求值}

$= M_{1}(2) + M_{2}(2)$ \textcolor{red}{$\rhd$ 因为在 $X=2$ 处对 $M_{1+2}(X)$ 求值，等价于将 $M_{1+2}(X)$ 拆分为 $M_{1}(X)$ 和 $M_{2}(X)$，分别对 $M_1(2)$ 与 $M_2(2)$ 求值后相加}

$ $

类似地，乘法之后的解码也保持正确：

$\text{Decode}(\sigma(m_1) \cdot \sigma(m_2)) = \text{Decode}(M_1(X) \cdot M_2(X))$

$= \text{Decode}(M_{1\cdot2}(X))$ \textcolor{red}{$\rhd$ 其中 $M_{1\cdot2}(X) = M_1(X) \cdot M_2(X)$}

$= M_{1\cdot2}(2)$ \textcolor{red}{$\rhd$ 因为解码就是在 $X = 2$ 处对多项式求值}

$= M_{1}(2) \cdot M_{2}(2)$ \textcolor{red}{$\rhd$ 因为在 $X=2$ 处对 $M_{1\cdot2}(X)$ 求值，等价于将 $M_{1\cdot2}(X)$ 拆分为 $M_{1}(X)$ 和 $M_{2}(X)$，分别对 $M_1(2)$ 与 $M_2(2)$ 求值后相乘}




$ $

因此，这些整数恒等式在所述无回绕及次数条件下成立，并不是商环中无条件成立的恒等式。

$ $

总结~\ref*{subsec:bfv-single-encoding} 的编码只有在所有结果多项式的次数严格小于 $n$，且所选整数系数不发生模 $t$ 回绕时，才保持整数语义。次数达到 $n$ 时，模 $X^n+1$ 约减就已经改变在 $2$ 处的求值。
由于这些约束，单值编码并不是全同态加密的理想选择。而且，尽管每个多项式包含 $n$ 个系数，却只能编码一个值，因此计算效率也较低。


\subsection{批量编码}
\label{subsec:bfv-batch-encoding}

下述标量槽位公式假设 $n\geq2$ 是二次幂，且 $t$ 是满足 $2n\mid(t-1)$ 的素数。此时 $\omega=g^{(t-1)/(2n)}$ 是本原 $2n$ 次单位根，范德蒙矩阵可逆。若 $t=p^r$ 是素数幂，且 $p\equiv1\pmod{2n}$，将根提升到 $\mathbb Z_{p^r}$ 后也能得到这些标量槽位；直接使用 $t-1$ 的指数公式一般不适用于素数幂。其他参数可能只能得到数量较少的扩环槽位。本章的 $W^T$ 是按行排列求值点的范德蒙矩阵，$W$ 是其转置；这里显式规定矩阵朝向。

单值编码方案（\autoref{subsec:bfv-single-encoding}）每次只对一个值进行编码与解码，而批量编码利用高维向量，同时对大量数值执行这些操作。因此，批量编码比单值编码更高效。此外，通过向量加法和 Hadamard 积，可以同时对批量编码的值逐元素执行同态加法与乘法。因此，批量编码能够采用 SIMD（单指令多数据）方式，比单值编码更高效地完成同态运算。

BFV 编码将模 $t$ 的、含 $n$ 个槽位的整数输入向量 $\vec{v} = (v_0, v_1, v_2, \cdots v_{n-1})$，转换为模 $t$ 的另一个含 $n$ 个槽位的向量 $\vec{m} = (m_0, m_1, m_2, \cdots m_{n-1})$。后者的各元素作为编码所得 $(n-1)$ 次或更低次数多项式 $M(X) \in \mathbb{Z}_t[X] / (X^n + 1)$ 的系数。

$ $

\subsubsection{编码第一阶段（\textsf{Encoding\textsubscript{1}}）}
\label{subsubsec:bfv-encoding-1} 

在 \autoref{subsec:poly-vector-transformation} 中，我们已经知道，通过映射 $\sigma$，可以将 $(n-1)$ 次或更低次数的多项式同构地映射为一个含 $n$ 个槽位的向量（为简化记号，将 \autoref{subsec:poly-vector-transformation} 中的 $\sigma_c$ 记为 $\sigma$）：

$\sigma: M(X) \in \mathbb{Z}_t[X]/(X^n + 1) \longrightarrow  (M(\omega),M(\omega^3),M(\omega^5), \cdots, M(\omega^{2n-1})) \in \mathbb{Z}_t^n$

$ $

，即在 $n$ 个不同的 $(\mu=2n)$ 次本原单位根 $\omega, \omega^3, \omega^5, \cdots, \omega^{2n-1}$ 处对多项式 $M(X)$ 求值。令 $\vec{m}$ 为包含多项式 $M(X)$ 的 $n$ 个系数的向量，则可将映射 $\sigma$ 表示为：

$\vec{v} = W^T \cdot \vec{m}$

$ $

，其中 $W^T$ 如下：

$W^T = \begin{bmatrix}
1 & (\omega) & (\omega)^2 & \cdots & (\omega)^{n-1}\\
1 & (\omega^3) & (\omega^3)^2 & \cdots & (\omega^3)^{n-1}\\
1 & (\omega^5) & (\omega^5)^2 & \cdots & (\omega^5)^{n-1}\\
\vdots & \vdots & \vdots & \ddots & \vdots \\
1 & (\omega^{2n-1}) & (\omega^{2n-1})^2 & \cdots & (\omega^{2n-1})^{n-1}\\
\end{bmatrix}$  \textcolor{red}{ \text{ } \# $W^T$ 是 \autoref{subsec:poly-vector-transformation} 中所述 $W$ 的转置}

$ $

注意，$W^T$ 的每一行与 $\vec{m}$ 的点积，给出多项式 $M(X)$ 在相应 $X = \{w, w^3, w^5, \cdots, w^{2n-1}\}$ 处的求值结果。在 BFV 编码方案中，\textsf{Encoding\textsubscript{1}} 将含 $n$ 个槽位的输入向量 $\vec{v}$ $\in \mathbb{Z}_t^n$ 编码为明文多项式 $M(X) \in \mathbb{Z}_t[X] / (X^n + 1)$，而 \textsf{Decoding\textsubscript{2}} 则将 $M(X)$ 解码回 $\vec{v}$。由于 $W^T \cdot \vec{m}$ 得到的 $\vec{v}$ 是 $M(X)$ 的解码，因此将 $W^T$ 称为解码矩阵。另一方面，\textsf{Encoding\textsubscript{1}} 的目标是将 $\vec{v}$ 编码为 $M(X)$，以便基于 $M(X)$ 进行同态计算。由关系 $\vec{v} = W^T \cdot \vec{m}$，可以推导出编码公式：

$(W^T)^{-1} \cdot \vec{v} = (W^T)^{-1} W^T \cdot \vec{m}$

$\vec{m} = (W^T)^{-1} \cdot \vec{v}$

$ $

因此，需要求出 $W^T$ 的逆矩阵 $(W^T)^{-1}$，作为编码矩阵。根据 \autoref{subsec:vandermonde-euler} 中的定理~\ref*{subsec:vandermonde-euler}，已经知道 $V^{-1} = \dfrac{V^T \cdot I_n^R}{n}$，其中 $V = W^T$ 且 $V^T = W$。换言之，$(W^T)^{-1} = \dfrac{W \cdot I_n^R}{n}$。因此，BFV 的编码公式可以写为：

$\vec{m} = (W^T)^{-1} \cdot \vec{v} = \dfrac{W \cdot I_n^R \cdot \vec{v}}{n}$，\text{ } 其中

$ $

$W =  \begin{bmatrix}
1 & 1 & 1 & \cdots & 1\\
(\omega) & (\omega^3) & (\omega^5) & \cdots & (\omega^{2n-1})\\
(\omega)^2 & (\omega^3)^2 & (\omega^5)^2 & \cdots & (\omega^{2n-1})^2\\
\vdots & \vdots & \vdots & \ddots & \vdots \\
(\omega)^{n-1} & (\omega^3)^{n-1} & (\omega^5)^{n-1} & \cdots & (\omega^{2n-1})^{n-1}\\
\end{bmatrix}$

$= \begin{bmatrix}
1 & 1 & \cdots & 1 & 1 & \cdots & 1 & 1\\
(\omega) & (\omega^3) & \cdots & (\omega^{n - 1}) & (\omega^{-(n - 1)}) & \cdots & (\omega^{-3}) & (\omega^{-1})\\
(\omega)^2 & (\omega^3)^2 & \cdots & (\omega^{n - 1})^2 & (\omega^{-(n - 1)})^2 & \cdots & (\omega^{-3})^2 & (\omega^{-1})^2\\
\vdots & \vdots & \vdots & \ddots & \vdots \\
(\omega)^{n-1} & (\omega^3)^{n-1} & \cdots & (\omega^{n - 1})^{n-1} & (\omega^{-(n - 1)})^{n-1} & \cdots & (\omega^{-3})^{n-1} & (\omega^{-1})^{n-1}\\
\end{bmatrix}$

$ $

，其中 $\omega$ 是模 $t$ 的 $2n$ 次本原单位根（这意味着 $t \equiv 1 \bmod 2n$）。因此有 $\omega = g^{\frac{t - 1}{2n}} \bmod t$（$g$ 是 $\mathbb{Z}_t^{\times}$ 的生成元，见 \autoref{subsubsec:poly-vector-transformation-modulus}）。

在 \autoref{subsec:poly-vector-transformation} 中，我们已经知道 $W$ 是 $n$ 维向量空间的一组基。因此，对于 $n$ 维向量空间 $\mathbb{Z}_t^{n}$ 中的每个 $\vec{v}$，都存在唯一对应的向量 $\dfrac{W \cdot I_n^R\cdot\vec{v}}{n} = \vec{m}$（参见 \autoref{subsec:projection} 中的定理~\ref*{subsec:projection}）。由此，以 $\vec{m}$ 的 $n$ 个元素为系数构成的多项式 $M(X)$，也是与 $\vec{v}$ 双射对应的唯一多项式。

注意，通过计算 $\dfrac{W \cdot I_n^R \cdot \vec{v}}{n}$
，这给出线性双射。该环同构把 $\vec v$ 的逐槽乘法对应到 $M$ 的负循环多项式乘法，而不是系数向量 $\vec m$ 的逐分量乘法。

\subsubsection{编码第二阶段（\textsf{Encoding\textsubscript{2}}）}

得到含 $n$ 个槽位的向量 $\vec{m}$ 后，使用尺度因子 $\Delta = \left\lfloor \dfrac{q}{t} \right\rfloor$ 对其放大（即相乘），其中 $q$ 是密文模数。将 $\vec{m}$ 乘以 $\Delta$，得到 $\Delta \vec{m}$。
$\Delta\vec{m}$ 中的 $n$ 个整数将作为 RLWE 加密中明文多项式的 $n$ 个系数。最终编码后的明文多项式为 $\Delta M = \sum\limits_{i=0}^{n-1}\Delta m_iX^i$。


\subsubsection{解码第一阶段（\textsf{Decoding\textsubscript{1}}）}
\label{subsubsec:bfv-enc-dec-decoding1} 

编码与解码示意图中的除以 $\Delta$，针对的是理想、未约减且无噪声的编码多项式。对于实际密文，应先使用 \autoref{subsec:bfv-enc-dec} 的缩放舍入解码器，再对 $\mathcal R_{\langle n,t\rangle}$ 中恢复出的多项式应用槽位求值矩阵。

RLWE 密文完成解密的前半部分、得到 $\Delta M = \sum\limits_{i=0}^{n-1}\Delta m_iX^i$ 后，计算 $\dfrac{\Delta\vec{m}}{\Delta} = \vec{m}$。




\subsubsection{解码第二阶段（\textsf{Decoding\textsubscript{2}}）}


在 \autoref{subsubsec:bfv-encoding-1} 中，我们已经推导出以下解码公式，用于将一个系数为模 $t$ 整数、次数为 $(n-1)$ 的多项式转换成包含 $n$ 个槽位的输入向量：

$\vec{v} = W^T \cdot \vec{m}$


\subsubsection{小结}
\label{subsubsec:bfv-encoding-summary}

\begin{tcolorbox}[title={\textbf{\tboxlabel{\ref*{subsubsec:bfv-encoding-summary}} BFV 的编码与解码}}]



\textbf{\underline{输入}：} 包含 $n$ 个槽位的模 $t$ 整数向量 $\vec{v} = (v_0, v_1, \cdots, v_{n-1}) \in \mathbb{Z}_t^n$

\par\noindent\rule{\textwidth}{0.4pt}

\textbf{\underline{编码}}
\begin{enumerate}

\item 应用变换 $\vec{m} = n^{-1}\cdot W \cdot I_n^R \cdot \vec{v}$，将 $\vec{v} \in \mathbb{Z}_t^n$ 转换为 $\vec{m} \in \mathbb{Z}_t^n$

，其中，$W$ 是按如下方式构造的 $n$ 维向量空间的一组基：

$W =  \begin{bmatrix}
1 & 1 & 1 & \cdots & 1\\
(\omega) & (\omega^3) & (\omega^5) & \cdots & (\omega^{2n-1})\\
(\omega)^2 & (\omega^3)^2 & (\omega^5)^2 & \cdots & (\omega^{2n-1})^2\\
\vdots & \vdots & \vdots & \ddots & \vdots \\
(\omega)^{n-1} & (\omega^3)^{n-1} & (\omega^5)^{n-1} & \cdots & (\omega^{2n-1})^{n-1}\\
\end{bmatrix}$

$= \begin{bmatrix}
1 & 1 & \cdots & 1 & 1 & \cdots & 1 & 1\\
(\omega) & (\omega^3) & \cdots & (\omega^{n - 1}) & (\omega^{-(n - 1)}) & \cdots & (\omega^{-3}) & (\omega^{-1})\\
(\omega)^2 & (\omega^3)^2 & \cdots & (\omega^{n - 1})^2 & (\omega^{-(n - 1)})^2 & \cdots & (\omega^{-3})^2 & (\omega^{-1})^2\\
\vdots & \vdots & \vdots & \ddots & \vdots \\
(\omega)^{n-1} & (\omega^3)^{n-1} & \cdots & (\omega^{n - 1})^{n-1} & (\omega^{-(n - 1)})^{n-1} & \cdots & (\omega^{-3})^{n-1} & (\omega^{-1})^{n-1}\\
\end{bmatrix}$

$ $

，其中，$\omega$ 是模 $t$ 的一个本原 $2n$ 次单位根。这意味着 $\omega = g^{\frac{t - 1}{2n}} \bmod t$（$g$ 是 $\mathbb{Z}_t^{\times}$ 的一个生成元，见 \autoref{subsubsec:poly-vector-transformation-modulus}）。

$ $

\item 将 $\vec{m}$ 转换为缩放后的整数向量 $\Delta\vec{m}$，其中，$\Delta=\lfloor \dfrac{q}{t}\rfloor$ 是尺度因子。如果 $\Delta$ 过于接近 1（即 $t$ 太大），则\textbf{噪声预算}会过小（容易导致解密失败）。如果 $\Delta$ 过于接近 $q$（即 $t$ 太小），则\textbf{消息容量}会过小（即明文模数 $t$ 限制了可编码数值的范围）。最终编码得到的明文多项式为 $\Delta M = \sum\limits_{i=0}^{n-1} \Delta m_i X^i \text{ } \in \mathbb{Z}_q[X] / (X^n + 1)$。 %The rounding of $\lceil \Delta \vec{m} \rfloor$ during the encoding process makes CKKS an approximate FHE scheme.

\end{enumerate}

\par\noindent\rule{\textwidth}{0.4pt}

\textbf{\underline{解码}：} 从明文多项式 $\Delta M = \sum\limits_{i=0}^{n-1}\Delta m_iX^i$ 中恢复 $\vec{m} = \dfrac{\Delta \vec{m}}{\Delta}$。然后，
计算 $\vec{v} = W^T \cdot \vec{m}$。

\end{tcolorbox}


%\para{BFV's Encoding Range:} \autoref{tab:polynomial-encoding-capacity-integer} depicts the range of input values that can be encoded by the polynomials $M(X) \in \mathcal{R}_{\langle n, t \rangle}$ over $X \in \mathbb{Z}_t$. As illustrated in \autoref{tab:polynomial-encoding-capacity-integer}, the encoding range of integer inputs is determined by the size of the $n$ and $t$ parameters. Therefore, these two parameters should be chosen sufficiently large enough to encode all values that can be used by an application. 

不过，小结~\ref*{subsubsec:bfv-encoding-summary} 还不是 BFV 批量编码的最终版本。在 \autoref{subsec:bfv-rotation} 中，我们将介绍如何在不解密封装输入向量的密文的情况下，对其槽位进行同态旋转。为了支持这种同态旋转，需要对小结~\ref*{subsubsec:bfv-encoding-summary} 中介绍的编码方案稍作调整。我们将在 \autoref{subsec:bfv-rotation} 中说明具体方法，并在 \autoref{subsubsec:bfv-rotation-summary} 的小结~\ref*{subsubsec:bfv-rotation-summary} 中给出 BFV 的最终编码方案。


\subsection{加密与解密}
\label{subsec:bfv-enc-dec}

这里使用 RLWE 相位 $B+AS$ 描述对称 BFV 加密，即 \autoref{subsec:glwe-alternative} 中取 $k=1$ 的另一种符号约定（\autoref{sec:rlwe}）。用 $[U]_q$ 表示逐系数中心模约减，并设 $r_q=q-t\Delta$。完整解码器 $\textsf{Dec}$ 使用比例 $t/q$。本章其他相位推导中的 $\textsf{RLWE}^{-1}$ 表示最终缩放舍入之前的带噪相位。任何一种记号都不表示同态运算后的噪声仍是新鲜噪声。



\begin{tcolorbox}[title={\textbf{\tboxlabel{\ref*{subsec:bfv-enc-dec}} BFV 加密与解密}}]

\textbf{\underline{初始设置}：} $\Delta=\left\lfloor\dfrac{q}{t}\right\rfloor \text{ 是多项式编码的明文尺度因子}, \text{ } S \xleftarrow{\$} \mathcal{R}_{\langle n, \textit{tern} \rangle}$

基本方案的明文模数可以是任意满足 $q\gg t$ 的整数 $t\geq2$。批编码和素数幂自举方法还需满足各自小节所列的附加条件。这里 $S$ 的系数取自三元集合 $\{-1,0,1\}$。

\par\noindent\rule{\textwidth}{0.4pt}

\textbf{\underline{加密输入}：} $\Delta M \in \mathcal{R}_{\langle n, q \rangle}$、$A \xleftarrow{\$} \mathcal{R}_{\langle n, q \rangle}$、$E \xleftarrow{\chi_\sigma} \mathcal{R}_{\langle n, q \rangle}$


$ $

\begin{enumerate}


\item 计算 $B = -A \cdot S + \Delta M + E \text{ } \in \mathcal{R}_{\langle n,q \rangle}$

\item $\textsf{RLWE}_{S,\sigma}(\Delta M + E) = (A, B) \text{ } \in \mathcal{R}_{\langle n,q \rangle}^2$ 

\end{enumerate}

\par\noindent\rule{\textwidth}{0.4pt}

\textbf{\underline{解密输入}：} $\textsf{ct} = (A, B) \text{ } \in \mathcal{R}_{\langle n,q \rangle}^2$

$\textsf{Dec}_S(A,B)=\left[\left\lceil\dfrac{t}{q}[B+AS]_q\right\rfloor\right]_t=M$。

（噪声 $E = \sum\limits_{i=0}^{n-1}e_iX^i$ 由舍入过程消除。）



$ $

\textbf{\underline{正确解密的条件}：}

对于相位 $\Delta M+E\pmod q$，一个逐系数充分条件是 $\|tE-r_qM\|_\infty<q/2$。特别地，当 $M$ 使用满足 $\|M\|_\infty\leq t/2$ 的中心代表，且右侧为正时，$\|E\|_\infty<q/(2t)-r_q/2$ 是充分条件。同态运算后对明文模约减还会产生余数误差，推导见 \autoref{subsubsec:bfv-noise-analysis}。这是 BFV 的缩放舍入条件；不能在不改变假设的情况下，用 \autoref{subsubsec:lwe-noise-bound} 讨论的另一种 $1/\Delta$ 约定替代它。

\end{tcolorbox}


在本节中，为简洁起见，我们经常将 $\textsf{RLWE}_{S,\sigma}(\Delta  M + E)$ 写作 $\textsf{RLWE}_{S,\sigma}(\Delta  M)$，因为 $\textsf{RLWE}_{S,\sigma}(\Delta M + E) \approx \textsf{RLWE}_{S,\sigma}(\Delta M)$（即二者解密得到相同的消息）。即使我们用 $\textsf{RLWE}_{S,\sigma}(\Delta  M)$ 代替 $\textsf{RLWE}_{S,\sigma}(\Delta  M + E)$，也应将其理解为对 $\Delta  M + E$ 的加密（即噪声包含在缩放后的消息中）。


我们将在 \autoref{subsubsec:bfv-noise-analysis} 中更详细地说明 BFV 正确解密的条件。



\subsection{密文与密文相加}
\label{subsec:bfv-add-cipher}

BFV 的密文与密文相加采用 RLWE 的密文与密文相加方案，但将加密和解密公式中 $A\cdot S$ 项的符号取反。具体而言，这等价于 $k = 1$ 时的另一种 GLWE 版本（\autoref{subsec:glwe-alternative}）的密文与密文相加方案。


\begin{tcolorbox}[title={\textbf{\tboxlabel{\ref*{subsec:bfv-add-cipher}} BFV 密文与密文相加}}]
$\textsf{RLWE}_{S, \sigma}(\Delta M^{\langle 1 \rangle} + E^{\langle 1 \rangle} ) + \textsf{RLWE}_{S, \sigma}(\Delta M^{\langle 2 \rangle} + E^{\langle 2 \rangle}) $

$ = ( A^{\langle 1 \rangle}, \text{ } B^{\langle 1 \rangle}) + (A^{\langle 2 \rangle}, \text{ } B^{\langle 2 \rangle}) $

$ = ( A^{\langle 1 \rangle} + A^{\langle 2 \rangle}, \text{ } B^{\langle 1 \rangle} + B^{\langle 2 \rangle} ) $

$= \textsf{RLWE}_{S, \sigma}(\Delta(M^{\langle 1 \rangle} + M^{\langle 2 \rangle}) + E^{\langle 1 \rangle} + E^{\langle 2 \rangle})$
\end{tcolorbox}


\subsubsection{噪声界分析}
\label{subsubsec:bfv-noise-analysis}

令 $r_q=q-t\Delta$，中心相位为 $x=[\Delta M+E]_q=\Delta M+E-qK$。乘以 $t/q$ 得到精确恒等式
\[
\frac{t}{q}x=M+\frac{tE-r_qM}{q}-tK.
\]
最后模 $t$ 约减时，整数倍 $tK$ 消失。因此，\autoref{subsec:bfv-enc-dec} 中的总结~\ref*{subsec:bfv-enc-dec} 在以下条件下正确解密：
\[
\|tE-r_qM\|_\infty<q/2.
\]
该证明没有删去内部模 $q$ 约减。当 $r_q\ne0$ 时，用 $1/\Delta$ 代替 $t/q$ 并不能得到相同恒等式，而需要另行指定解码器和代表元约定。

对于密文加法，写成 $M^{\langle1\rangle}+M^{\langle2\rangle}=M^{\langle1+2\rangle}+tH$，其中 $M^{\langle1+2\rangle}$ 是所选中心明文代表。因为 $t\Delta=q-r_q$，相加密文的相位为
\[
\Delta M^{\langle1+2\rangle}+E_{\rm add}\pmod q,
\qquad E_{\rm add}=E^{\langle1\rangle}+E^{\langle2\rangle}-r_qH.
\]
其正确性的充分条件为
\[
\bigl\|t(E^{\langle1\rangle}+E^{\langle2\rangle}-r_qH)-r_qM^{\langle1+2\rangle}\bigr\|_\infty<q/2.
\]
项 $-r_qH$ 记录明文回绕，不能省略。乘法或旋转后，只要明确有效明文和完整误差，也可应用相同推理。这与 \autoref{subsubsec:lwe-noise-bound} 中需要显式保留余数项的原则一致。增大 $q/t$ 可能增加正确性预算，但安全性仍需根据完整参数集评估。

\subsection{密文与明文相加}
\label{subsec:bfv-add-plain}

相加的明文多项式使用尺度 $\Delta$。若 $M+\Lambda=M_{\rm out}+tH$，则明文约减后的有效误差为 $E-r_qH$，其中 $r_q=q-t\Delta$，与 \autoref{subsubsec:bfv-noise-analysis} 中的约定相同。

BFV 的密文与明文相加采用 RLWE 的密文与明文相加方案，但将加密和解密公式中 $A\cdot S$ 项的符号取反。具体而言，这等价于 $k = 1$ 时的另一种 GLWE 版本（\autoref{subsec:glwe-alternative}）的密文与明文相加方案（\autoref{sec:glwe-add-plain}）。


\begin{tcolorbox}[title={\textbf{\tboxlabel{\ref*{subsec:bfv-add-plain}} BFV 密文与明文相加}}]
$\textsf{RLWE}_{S, \sigma}(\Delta M + E) + \Delta\Lambda $

$=  (A, \text{ } B) + \Delta\Lambda$

$=  (A, \text{ } B + \Delta\cdot\Lambda)$

$= \textsf{RLWE}_{S, \sigma}(\Delta (M + \Lambda) + E)$
\end{tcolorbox}



\subsection{密文与明文相乘}
\label{subsec:bfv-mult-plain}

这里 $\Lambda$ 是未缩放的明文乘数。若 $M\Lambda=M_{\rm out}+tH$，则有效误差为 $\Lambda E-r_qH$。对于多项式乘法，$\|\Lambda E\|_\infty\leq\|\Lambda\|_1\|E\|_\infty$；这个卷积界及余数项都必须计入解码条件。


BFV 的密文与明文相乘采用 RLWE 的密文与明文相乘方案，但将加密和解密公式中 $A\cdot S$ 项的符号取反。具体而言，这等价于 $k = 1$ 时的另一种 GLWE 版本（\autoref{subsec:glwe-alternative}）的密文与明文相乘方案（\autoref{sec:glwe-mult-plain}）。


\begin{tcolorbox}[title={\textbf{\tboxlabel{\ref*{subsec:bfv-mult-plain}} BFV 密文与明文相乘}}]
$\textsf{RLWE}_{S, \sigma}(\Delta M + E) \cdot \Lambda$

$= (A, \text{ } B) \cdot \Lambda$

$= (A \cdot \Lambda, \text{ }  B \cdot \Lambda )$

$= \textsf{RLWE}_{S, \sigma}(\Delta (M \cdot \Lambda) + \Lambda E)$
\end{tcolorbox}




\subsection{密文与密文相乘}
\label{subsec:bfv-mult-cipher}

以下四阶段表述是一个需要显式误差界的教学变体：它使用辅助模数 $Q=q\Delta$，并以 $1/\Delta$ 重缩放。原始 FV 乘法在重线性化之前，以 $t/q$ 缩放整数提升后的张量乘积分量；两种比例只有在 $t\mid q$ 时才相同。这里保留该变体以解释模提升，但完整保留其附加余数和舍入误差，并不声称所有 BFV 实现都按此顺序执行这四个阶段。

\noindent \textbf{- 参考资料 1：}
\href{https://www.inferati.com/blog/fhe-schemes-bfv}{BFV 加密方案简介}~\cite{inferati-bfv}

\noindent \textbf{- 参考资料 2：}
\href{https://eprint.iacr.org/2012/144.pdf}{某种程度上、部分与全同态加密}~\cite{cryptoeprint:2012/144}

%\noindent \textbf{- Reference 3:} 
%\href{https://eprint.iacr.org/2016/510.pdf}{A Full RNS Variant $of FV like Somewhat Homomorphic Encryption Schemes}~\cite{cryptoeprint:2016/510}


给定两个密文 $\textsf{RLWE}_{S, \sigma}(\Delta M^{\langle 1\rangle })$ 和 $\textsf{RLWE}_{S, \sigma}(\Delta M^{\langle 2\rangle})$，密文与密文相乘的目标是构造一个新密文，使其解密结果为 $\Delta M^{\langle 1 \rangle} M^{\langle 2 \rangle}$。密文与密文相乘比密文与明文相乘更复杂，包括四个步骤：（1）\textsf{ModRaise}；（2）多项式乘法；（3）重线性化；（4）重缩放。

为便于理解，我们将基于 RLWE 的另一种版本（\autoref{subsec:glwe-alternative} 中的定理~\ref*{subsec:glwe-alternative}）来解释 BFV 的密文与密文相乘；在该版本中，加密和解密公式中 $AS$ 项的符号取反。

\subsubsection{模数提升（\textsf{ModRaise}）}
\label{subsubsec:bfv-mult-cipher-modraise}

选定 $(A,B)$ 的整数系数代表，在 $\mathbb Z[X]/(X^n+1)$ 中将未约减相位写为
\[
AS+B=\Delta M+E+qK.
\]
把相同系数代表嵌入模数 $Q=q\Delta$ 后，相位变为 $\Delta M+E+qK\pmod Q$，并不是只有误差 $E$ 的、原明文 $\Delta M$ 的加密。一般不能断言 $(q-1)n+(q-1)<Q$；这还要求 $n$ 与 $\Delta$ 满足额外关系。上述同余本身也不需要无回绕条件。后续两个阶段的乘积和相位都按模 $Q$ 解释；重缩放时会显式计入整数提升及其中的 $Q$ 倍数。

BFV 密文与密文相乘的第一步是对两个输入密文 $(A^{\langle 1 \rangle}, B^{\langle 1 \rangle}) \bmod q$ 和 $(A^{\langle 2 \rangle}, B^{\langle 2 \rangle}) \bmod q$ 进行\textit{模数提升}，
按如下方式将模数变为 $q \rightarrow Q$（其中 $Q = q\cdot \Delta$）：

$(A^{\langle 1 \rangle}, B^{\langle 1 \rangle}) \bmod Q$

$(A^{\langle 2 \rangle}, B^{\langle 2 \rangle}) \bmod Q$

$ $

经过 \textsf{ModRaise} 后，这两个密文的解密结果为：

$A^{\langle 1 \rangle} \cdot S + B^{\langle 1 \rangle} = \Delta M^{\langle 1 \rangle} + E^{\langle 1 \rangle} + K_1q \pmod Q$

$A^{\langle 2 \rangle} \cdot S + B^{\langle 2 \rangle} = \Delta M^{\langle 2 \rangle} + E^{\langle 2 \rangle} + K_2q \pmod Q$


$ $

因此，\textit{模数提升后}的密文具有以下形式：

$\textsf{RLWE}_{S, \sigma}(\Delta M^{\langle 1 \rangle} + K_1q) = (A^{\langle 1 \rangle}, B^{\langle 1 \rangle}) \bmod Q$

$\textsf{RLWE}_{S, \sigma}(\Delta M^{\langle 2 \rangle} + K_2q) = (A^{\langle 2 \rangle}, B^{\langle 2 \rangle}) \bmod Q$


\subsubsection{多项式乘法}
\label{subsubsec:bfv-mult-cipher-multiplication}

接下来的目标是构造一个新密文，它加密的内容为
$(\Delta M^{\langle 1 \rangle} + E^{\langle 1 \rangle} + K_1q) \cdot (\Delta M^{\langle 2 \rangle} + E^{\langle 2 \rangle} + K_2q)$. 

首先，可以推导出以下关系：

$(\Delta M^{\langle 1 \rangle} + E^{\langle 1 \rangle} + K_1q) \cdot (\Delta M^{\langle 2 \rangle} + E^{\langle 2 \rangle} + K_2q)$

$ = (A^{\langle 1 \rangle} \cdot S + B^{\langle 1 \rangle}) \cdot (A^{\langle 2 \rangle} \cdot S + B^{\langle 2 \rangle})$

$ = \underbrace{B^{\langle 1 \rangle}B^{\langle 2 \rangle}}_{D_0}  + \underbrace{(B^{\langle 2 \rangle}A^{\langle 1 \rangle} + B^{\langle 1 \rangle}A^{\langle 2 \rangle})}_{D_1} \cdot S + \underbrace{(A^{\langle 1 \rangle} \cdot A^{\langle 2 \rangle})}_{D_2} \cdot S \cdot S $



$= D_0 + D_1\cdot S + D_2\cdot S^2$



$ $

同时，我们还具有以下关系：

$\textsf{RLWE}_{S, \sigma}^{-1}(\Delta M^{\langle 1 \rangle} + K_1q) = \Delta M^{\langle 1 \rangle} + E^{\langle 1 \rangle} + K_1q$

$\textsf{RLWE}_{S, \sigma}^{-1}(\Delta M^{\langle 2 \rangle} + K_2q) = \Delta M^{\langle 2 \rangle} + E^{\langle 2 \rangle} + K_2q$

$ $

将这些关系结合起来，得到：

$\textsf{RLWE}_{S, \sigma}^{-1}(\Delta M^{\langle 1 \rangle} + K_1q) \cdot \textsf{RLWE}_{S, \sigma}^{-1}(\Delta M^{\langle 2 \rangle} + K_2q) = D_0 + D_1\cdot S + D_2\cdot S^2$

$ $

注意，作为密文分量的 $D_0, D_1,$ 和 $D_2$ 是已知值，而 $S$ 只有私钥持有者知道。因此，可以将 $D_0 + D_1\cdot S + D_2\cdot S^2$ 视为解密公式：给定密文分量 $D_0, D_1, D_2$ 和私钥 $S$，便可得到 $(\Delta M^{\langle 1 \rangle} + E^{\langle 1 \rangle} + K_1q) \cdot (\Delta M^{\langle 2 \rangle} + E^{\langle 2 \rangle} + K_2q)$。换言之，可以令 $(D_0, D_1, D_2)$ 为一种新的密文形式，它能够由 $S$ 解密为 $(\Delta M^{\langle 1 \rangle} + E^{\langle 1 \rangle} + K_1q) \cdot (\Delta M^{\langle 2 \rangle} + E^{\langle 2 \rangle} + K_2q)$。

三元组 $(D_0,D_1,D_2)$ 是关于秘密密钥 $S$ 的二次密文，因此解密使用的是平方 $S^2$，不是平方根。继续相乘会使秘密密钥次数相加；重复平方则在每步将次数翻倍。重线性化以引入求值密钥误差为代价，将密文恢复为两个分量及一次形式。



\subsubsection{重线性化}
\label{subsubsec:bfv-mult-cipher-relinearization}

对于任意 $Q$，本小节使用整数分解基 $g_i=\beta^{i-1}$，取 $l=\lceil\log_\beta Q\rceil$，逐系数精确分解满足 $D_2\equiv\sum_{i=1}^lD_{2,i}g_i\pmod Q$。这是递增权重约定，不是递减的 $Q/\beta^i$ 约定；相匹配的 RLev 各行加密 $g_iS^2$。令 $E_{\rm rl}=\sum_iD_{2,i}E_i'$ 表示累计误差。重线性化后的相位等于原相位加 $E_{\rm rl}$，并非与原相位完全相等。

重线性化将三元组 $(D_0,D_1,D_2)\in\mathcal R_{\langle n,Q\rangle}^3$ 转换为两个分量。取 $\textsf{ct}_\alpha=(D_1,D_0)$，并通过分解构造 $\textsf{ct}_\beta$，则有
\[
\phi_S(\textsf{ct}_\alpha+\textsf{ct}_\beta)\equiv D_0+D_1S+D_2S^2+E_{\rm rl}\pmod Q.
\]

在 $D_0+D_1S+D_2S^2$ 中，前两项构成合成密文 $\textsf{ct}_\alpha=(D_1,D_0)$。分解构造产生的 $\textsf{ct}_\beta$ 满足相位关系 $\phi_S(\textsf{ct}_\beta)\equiv D_2S^2+E_{\rm rl}\pmod Q$。公开秘密密钥函数（如 $S^2$）的加密，还需要所选构造的求值密钥安全性假设；下列代数并不构成对该假设的证明。

构造加密 $D_2 \cdot S^2$ 的密文有一种直接的方法：将 $S^2$ 加密为 RLWE 密文 $\textsf{RLWE}_{S, \sigma}(S^2) = (A^{\langle s \rangle}, B^{\langle s \rangle})$，使其满足 $A^{\langle s \rangle}\cdot S + B^{\langle s \rangle} = S^2 + E^{\langle s \rangle} \bmod Q$（其中，密文模数为 $Q$，明文尺度因子为 $\Delta = 1$）。然后，将 $D_2$ 看作模 $Q$ 的明文多项式，与 $D_2$ 进行密文与明文相乘（\autoref{sec:glwe-mult-plain}）。但这种方法在实际中不可行，因为计算 $D_2 \cdot \textsf{RLWE}_{S, \sigma}(S^2)$ 会产生如下巨大的噪声：

$D_2 \cdot (A^{\langle s \rangle}, B^{\langle s \rangle}) = (D_2\cdot A^{\langle s \rangle}, D_2 \cdot B^{\langle s \rangle})$

$ $

，其解密结果为：

$D_2\cdot A^{\langle s \rangle} \cdot S + D_2\cdot B^{\langle s \rangle} = D_2 \cdot S^2 + D_2\cdot E^{\langle s \rangle} \pmod Q$ 

$ $

。在上述解密表达式 $D_2S^2 + D_2E^{\langle s \rangle} \bmod Q$ 中，$D_2S^2$ 项可以模 $Q$ 约减，因为在最终解密公式 $D_0 + D_1S + D_2S^2 \bmod Q$ 中，这一项本来就允许模 $Q$ 约减。然而，问题出在噪声 $D_2\cdot E^{\langle s \rangle}$：其系数可以取 $[0, Q - 1]$ 中的任意值（因为多项式 $D_2 = A^{\langle 1 \rangle} A^{\langle 2 \rangle}$ 的每个系数都可以取 $[0, Q - 1]$ 中的任意值）。如此大的噪声无法保证最终的正确解密。

为避免这一噪声问题，改进的方法是使用Gadget 分解技术（\autoref{subsec:gadget-decomposition}），将加密 $D_2 S^2$ 的 RLWE 密文表示为多个噪声较小的 RLWE 密文之和。为此，我们使用加密 $S^2$ 的 RLev 密文（\autoref{sec:glev}）。假设Gadget 向量为 $\vec{g} = \Bigg(\beta^{0}, \beta^{1}, \beta^{2}, \cdots, \beta^{l-1}\Bigg)$。回顾我们的目标：在已知 $D_2$、未知 $S$ 且已知 $\textsf{RLev}_{S, \sigma}^{\beta, l}(S^2) = \left\{\textsf{RLWE}_{S, \sigma}\left(\beta^{i-1}\cdot S^2\right)\right\}_{i=1}^{l}$ 的情况下，求得 $\textsf{ct}_\beta = \textsf{RLWE}_{S, \sigma}( S^2 \cdot D_2)$。于是，可以按如下方式构造 $\textsf{ct}_\beta$：

$\textsf{ct}_\beta = \textsf{RLWE}_{S, \sigma}(S^2 \cdot D_2)$

$= \textsf{RLWE}_{S, \sigma}\left( S^2 \cdot \left(D_{2,1}\beta^{0} + D_{2,2}\beta^{1} + \cdots D_{2,l}\beta^{l-1}\right)\right)$ \textcolor{red}{ $\rhd$ 对 $D_2$ 进行分解}

$= \textsf{RLWE}_{S, \sigma}\left( S^2 \cdot D_{2,1} \cdot \beta^{0}\right) + \textsf{RLWE}_{S, \sigma}\left( S^2 \cdot D_{2,2} \cdot \beta^{1}\right) + \cdots + \textsf{RLWE}_{S, \sigma}\left( S^2 \cdot D_{2,l} \cdot \beta^{l-1}\right)$

$= D_{2,1}\cdot\textsf{RLWE}_{S, \sigma}\left( S^2 \cdot \beta^{0}\right) + D_{2,2}\cdot\textsf{RLWE}_{S, \sigma}\left( S^2 \cdot \beta^{1}\right) + \cdots + D_{2,l}\cdot\textsf{RLWE}_{S, \sigma}\left( S^2 \cdot \beta^{l-1}\right)$
\textcolor{red}{ $\rhd$ 其中，每个 \textsf{RLWE} 密文都以 $S^2\beta^{0}, S^2\beta^{1}, \cdots, S^2\beta^{l-1}$ 为明文，明文尺度因子为 $\Delta = 1$}

$ $

$= \bm{\langle} \textsf{Decomp}^{\beta, l}(D_2), \text{ } \textsf{RLev}_{S, \sigma}^{\beta, l}( S^2) \bm{\rangle}$ \textcolor{red}{ $\rhd$ 将 \textsf{Decomp} 和 \textsf{RLev} 视为向量并求内积}

$ $

对上述结果解密，得到：

$\textsf{RLWE}_{S, \sigma}^{-1}(\textsf{ct}_\beta = \bm{\langle} \textsf{Decomp}^{\beta, l}(D_2), \text{ } \textsf{RLev}_{S, \sigma}^{\beta, l}( S^2) \bm{\rangle} \bm{)}$ \textcolor{red}{ $\rhd$ $\textsf{RLev}_{S, \sigma}^{\beta, l}( S^2)$ 的尺度因子均为 1}

$= D_{2,1}\cdot\left(E_1' +  S^2\beta^{0}\right) + D_{2,2}\cdot\left(E_2' +  S^2\beta^{1}\right) + \cdots + D_{2,l}\cdot\left(E_l' +  S^2\beta^{l-1}\right)$ \textcolor{red}{ \text{ } \# 其中，每个 $E_i'$ 都是嵌入 $\textsf{RLWE}_{S, \sigma}\left(S^2\cdot\beta^{i-1}\right)$ 中的噪声}

$ $

$= \sum\limits_{i=1}^{l} (E_i'\cdot D_{2,i}) +  S^2\cdot\left(D_{2,1}\beta^{0} + D_{2,2}\beta^{1} + \cdots + D_{2, l}\beta^{l-1}\right)$

$= \sum\limits_{i=1}^{l} \epsilon_{i} + D_2\cdot S^2$ \textcolor{red}{ \text{ } \# 其中，每个 $\epsilon_i = E_i'\cdot D_{2,i}$}

$\equiv D_2S^2+E_{\rm rl}\pmod Q$，其中 $\|E_{\rm rl}\|_\infty\leq\sum_{i=1}^l\|D_{2,i}\|_1\|E_i'\|_\infty$。小数字给出有用的界，但不能与每一个可能的单密钥噪声实例作严格大小比较。

$ $

因此，我们得到以下完整关系：


$\textsf{RLWE}_{S, \sigma}^{-1}(\Delta M^{\langle 1 \rangle} + K_1q) \cdot \textsf{RLWE}_{S, \sigma}^{-1}(\Delta M^{\langle 2 \rangle} + K_2q) \bmod Q$

$= (\Delta M^{\langle 1 \rangle} + E^{\langle 1 \rangle} + K_1q) \cdot (\Delta M^{\langle 2 \rangle} + E^{\langle 2 \rangle} + K_2q) \bmod Q$

$ = (A^{\langle 1 \rangle} \cdot S + B^{\langle 1 \rangle}) \cdot (A^{\langle 2 \rangle} \cdot S + B^{\langle 2 \rangle})  \bmod Q$

$ = D_0 + D_1\cdot S + D_2\cdot S^2  \bmod Q$ \textcolor{red}{ $\rhd$ $D_0 = B^{\langle 1 \rangle} B^{\langle 2 \rangle}$, \text{ } $D_1 = A^{\langle 1 \rangle} B^{\langle 2 \rangle} + A^{\langle 2 \rangle} B^{\langle 1 \rangle}$, \text{ } $D_2 = A^{\langle 1 \rangle} A^{\langle 2 \rangle}$}

$ = \textsf{RLWE}_{S, \sigma}^{-1}(\textsf{ct}_\alpha) + \textsf{RLWE}_{S, \sigma}^{-1}(\textsf{ct}_\beta) - \sum\limits_{i=1}^{l} (E_i'\cdot D_{2,i})  \bmod Q$ 

\textcolor{red}{ $\rhd$ $\textsf{ct}_\alpha = (D_1, D_0) = (A_\alpha, B_\alpha)$, \text{ } $\textsf{ct}_\beta = \langle \textsf{Decomp}^{\beta, l}(D_2), \textsf{RLev}_{S, \sigma}^{\beta, l}(S^2) \rangle = (A_\beta, B_\beta)$}

$ $

$ = \textsf{RLWE}_{S, \sigma}^{-1}(\textsf{ct}_\alpha + \textsf{ct}_\beta) - \sum\limits_{i=1}^{l} (E_i'\cdot D_{2,i})  \bmod Q$

$ = \textsf{RLWE}_{S, \sigma}^{-1}\bm((A_{\alpha+\beta}, B_{\alpha+\beta})\bm) - \sum\limits_{i=1}^{l} (E_i'\cdot D_{2,i}) \bmod Q$  \textcolor{red}{ $\rhd$ $A_{\alpha+\beta} = A_{\alpha} + A_\beta$, \text{ } $B_{\alpha + \beta} = B_\alpha + B_\beta$}


$ $

从上述推导中提取以下关系：

$(\Delta M^{\langle 1 \rangle} + E^{\langle 1 \rangle} + K_1q) \cdot (\Delta M^{\langle 2 \rangle} + E^{\langle 2 \rangle} + K_2q) \bmod Q$

$=  \Delta^2M^{\langle 1 \rangle}M^{\langle 2 \rangle} + \Delta\cdot (M^{\langle 1 \rangle}E^{\langle 2 \rangle} + M^{\langle 2 \rangle}E^{\langle 1 \rangle}) + q\cdot(\Delta M^{\langle 1 \rangle}K_2 + \Delta M^{\langle 2 \rangle}K_1 + E^{\langle 1 \rangle}K_2 + E^{\langle 2 \rangle}K_1) + K_1 K_2 q^2 + E^{\langle 1 \rangle} E^{\langle 2 \rangle} \bmod Q$

$= \textsf{RLWE}_{S, \sigma}^{-1}\bm((A_{\alpha+\beta}, B_{\alpha+\beta})\bm) - \sum\limits_{i=1}^{l} (E_i'\cdot D_{2,i}) \bmod Q$  

$ $

可以将上述关系改写为：


$\textsf{RLWE}_{S, \sigma}^{-1}\bm((A_{\alpha+\beta}, B_{\alpha+\beta})\bm) = A_{\alpha+\beta}\cdot S + B_{\alpha+\beta} \bmod Q$

$ = \Delta^2M^{\langle 1 \rangle}M^{\langle 2 \rangle} + \Delta\cdot (M^{\langle 1 \rangle}E^{\langle 2 \rangle} + M^{\langle 2 \rangle}E^{\langle 1 \rangle}) + q\cdot(\Delta M^{\langle 1 \rangle}K_2 + \Delta M^{\langle 2 \rangle}K_1 + E^{\langle 1 \rangle}K_2 + E^{\langle 2 \rangle}K_1) + K_1 K_2 q^2 + E^{\langle 1 \rangle} E^{\langle 2 \rangle} + \sum\limits_{i=1}^{l} (E_i'\cdot D_{2,i}) \bmod Q$ 

$ $

该表达式包含乘积明文、提升项 $qK_i$、输入误差乘积及重线性化误差。下一小节会跟踪这些项经过教学变体的 $1/\Delta$ 重缩放后的变化，并应用完整 BFV 解码条件，而不预先假设它与无噪声加密近似相等。


\subsubsection{重缩放}
\label{subsubsec:bfv-mult-cipher-rescaling}

令 $(A_*,B_*)=(A_{\alpha+\beta},B_{\alpha+\beta})$ 为模数 $Q=q\Delta$ 下重线性化后的密文，并选定整数系数代表。定义
\[
C_A=\Bigl\lceil A_*/\Delta\Bigr\rfloor,\quad C_B=\Bigl\lceil B_*/\Delta\Bigr\rfloor,
\quad R_A=C_A-A_*/\Delta,\quad R_B=C_B-B_*/\Delta.
\]
$R_A,R_B$ 的每个系数绝对值至多为 $1/2$。用 $L_i=\Delta M^{\langle i\rangle}+E^{\langle i\rangle}+qK_i$ 表示整数输入相位，并保留 \autoref{subsubsec:bfv-mult-cipher-relinearization} 中的 $E_{\rm rl}$，则
\[
A_*S+B_*=L_1L_2+E_{\rm rl}+QK_*,
\qquad C_AS+C_B\equiv\frac{L_1L_2+E_{\rm rl}}{\Delta}+R_AS+R_B\pmod q.
\]
这里使用 $Q/\Delta=q$，所以被约减的 $Q$ 倍数在除法后成为整数倍 $q$。设 $r_q=q-t\Delta$，展开乘积并去除整数倍 $q$ 后得到
\[
C_AS+C_B\equiv\Delta M^{\langle1\rangle}M^{\langle2\rangle}+\eta\pmod q,
\]
\[
\begin{split}
\eta={}&M^{\langle1\rangle}E^{\langle2\rangle}+M^{\langle2\rangle}E^{\langle1\rangle}
+\frac q\Delta(E^{\langle1\rangle}K_2+E^{\langle2\rangle}K_1)\\
&+\frac{q r_q}{\Delta}K_1K_2
+\frac{E^{\langle1\rangle}E^{\langle2\rangle}+E_{\rm rl}}{\Delta}+R_AS+R_B.
\end{split}
\]
特别地，$q^2/\Delta=qt+q r_q/\Delta$。后一项一般不可忽略，不能将其视为近似零的误差而删去。$\eta$ 的表达式可以含有有理数加项，但完整表达式在模 $q$ 下与整数输出相位一致。

解码时，将 $M^{\langle1\rangle}M^{\langle2\rangle}=M_{\rm out}+tH$ 写成所选中心明文代表。有效输出误差为 $E_{\rm out}=\eta-r_qH$。标准 BFV 解码器正确性的一个充分条件是
\[
\|t(\eta-r_qH)-r_qM_{\rm out}\|_\infty<q/2.
\]
当 $t\mid q$ 时，所有 $r_q$ 项消失，但明文与噪声的乘积、求值密钥误差和各分量舍入误差仍然存在。

\para{噪声分析：}上述乘积都是负循环多项式乘积。应使用 $\|UV\|_\infty\leq\|U\|_1\|V\|_\infty\leq n\|U\|_\infty\|V\|_\infty$，不能套用逐系数标量界。另有 $\|R_AS+R_B\|_\infty\leq(\|S\|_1+1)/2$。用这些不等式约束 $\eta$ 的其余各项，再应用三角不等式与最终阈值比较。仅凭 $q\gg t$ 或 $Q=q\Delta$，不能无条件保证乘法成功。

\para{模数切换与重缩放：}在这一特定构造中，以 $1/\Delta=q/Q$ 缩放分量，会将辅助模数从 $Q$ 改回 $q$，同时降低明文尺度。两个术语描述的是不同目的，并不总能互换。标准 FV 乘法可以通过对未约减整数乘积执行缩放舍入，在恢复明文尺度的同时保持对外密文模数不变。

\subsubsection{小结}
\label{subsubsec:bfv-mult-cipher-summary}

将以上内容综合起来，BFV 的密文与密文相乘可总结如下：

\begin{tcolorbox}[title={\textbf{\tboxlabel{\ref*{subsubsec:bfv-mult-cipher-summary}} BFV 的密文与密文相乘}}]


假设有以下两个 RLWE 密文：

$\textsf{RLWE}_{S, \sigma}(\Delta M^{\langle 1 \rangle}) = (A^{\langle 1 \rangle}, B^{\langle 1 \rangle}) \bmod q$，\text{ } 其中 $B^{\langle 1 \rangle} = -A^{\langle 1 \rangle} \cdot S + \Delta M^{\langle 1 \rangle} + E^{\langle 1 \rangle} \bmod q$

$\textsf{RLWE}_{S, \sigma}(\Delta M^{\langle 2 \rangle}) = (A^{\langle 2 \rangle}, B^{\langle 2 \rangle})  \bmod q$，\text{ } 其中 $B^{\langle 2 \rangle} = -A^{\langle 2 \rangle} \cdot S + \Delta M^{\langle 2 \rangle} + E^{\langle 2 \rangle} \bmod q$

$ $

这两个密文的乘法按如下步骤执行：

$ $

\begin{enumerate}

\item \textbf{\underline{模数提升（\textsf{ModRaise}）}}

按如下方式将密文 $(A^{\langle 1 \rangle}, B^{\langle 1 \rangle}) \bmod q$ 和 $(A^{\langle 2 \rangle}, B^{\langle 2 \rangle}) \bmod q$ 的模数强制提升至 $Q$（其中 $Q = q \cdot \Delta$）：

$(A^{\langle 1 \rangle}, B^{\langle 1 \rangle}) \bmod Q$

$(A^{\langle 2 \rangle}, B^{\langle 2 \rangle}) \bmod Q$


$ $

\item \textbf{\underline{乘法}}

在模 $Q$ 下计算以下多项式乘法：

$D_0 = B^{\langle 1 \rangle}B^{\langle 2 \rangle} \bmod Q$

$D_1 = B^{\langle 2 \rangle}A^{\langle 1 \rangle} +  B^{\langle 1 \rangle}A^{\langle 2 \rangle} \bmod Q$

$D_2 = A^{\langle 1 \rangle} \cdot A^{\langle 2 \rangle} \bmod Q$


$ $

\item \textbf{\underline{重线性化}}


计算：

$ \textsf{ct}_\alpha = (D
_1, D_0)$

$ \textsf{ct}_\beta = \bm{\langle} \textsf{Decomp}^{\beta, l}(D_2), \text{ } \textsf{RLev}_{S, \sigma}^{\beta, l}( S^2) \bm{\rangle}$.

$ \textsf{ct}_{\alpha + \beta} = \textsf{ct}_{\alpha} + \textsf{ct}_{\beta}$

$ $

于是，以下性质成立：

$\phi_S(\textsf{ct}_{\alpha+\beta})\equiv(\Delta M^{\langle1\rangle}+E^{\langle1\rangle}+qK_1)(\Delta M^{\langle2\rangle}+E^{\langle2\rangle}+qK_2)+E_{\rm rl}\pmod Q$。

$ $




\item \textbf{\underline{重缩放}}

将 $\textsf{ct}_{\alpha + \beta} = (A_{\alpha + \beta}, B_{\alpha + \beta}) \bmod Q$ 更新为 $\textsf{ct'}_{\alpha + \beta} = \left(\left\lceil\dfrac{A_{\alpha + \beta}}{\Delta}\right\rfloor, \left\lceil\dfrac{B_{\alpha + \beta}}{\Delta}\right\rfloor\right) \bmod q$。

$ $

这一明文重缩放过程也可以视为密文 $\textsf{ct}_{\alpha + \beta}$ 按 $Q \rightarrow q$ 进行模数切换。


$ $

\end{enumerate}

注意，密文与密文相乘后，明文尺度因子 $\Delta=\left\lfloor\dfrac{q}{t}\right\rfloor$、密文模数 $q$ 和私钥 $S$ 均保持不变。



\end{tcolorbox}


\para{\textsf{ModRaise} 的目的：}在缩放之前把张量乘积模 $q$ 约减，可能丢失缩放舍入所需的信息。此教学变体改为在 $Q=q\Delta$ 下计算，因为被约减的 $Q$ 倍数在除以 $\Delta$ 后成为 $q$ 倍数。它并不要求每个明文乘积都小于 $Q$；对于多项式卷积及较大明文，这种断言一般不成立。仍需保留提升项 $qK_i$，并计入 \autoref{subsubsec:bfv-mult-cipher-rescaling} 推导的完整误差。

$ $

\para{交换\textsf{重线性化}与\textsf{重缩放}的顺序：}两种顺序都可以构造，但其求值密钥、模数和舍入误差界不同，不能直接使用同一组密钥和误差证明交换步骤。标准 FV 通常先缩放二次整数乘积，再在目标模数下重线性化。



\subsubsection{应用于 TFHE 的密文与密文相乘}
\label{subsubsec:bfv-mult-cipher-tfhe}

本小节有意采用另一种 LWE 相位 $b+\langle\vec a,\vec s\rangle$。若用于 \autoref{subsec:tfhe-enc-dec} 中相位为 $b-\langle\vec a,\vec s\rangle$ 的 TFHE 约定，须先将公开掩码取负，输出时再取负还原。张量乘积构造是一种分层 LWE 乘法替代方案，并非 TFHE 自举所使用的外积。它需要对乘积 $s_i s_j$ 的求值密钥、适当参数及完整重缩放误差界，并不能消除任意 LWE 乘法的噪声。


在 \autoref{subsec:tfhe-mult-cipher} 中，外积的输入是一个 LWE 密文和一个 GSW 密文。从任意 LWE 密文获得后者，可能需要电路自举。下述张量构造改为使用另一种二次规模的求值密钥，并要求相应的重缩放预算。这给出一种替代运算，并不宣称其性能优于 TFHE 外积或可编程自举。与本书 TFHE 的编码一样，要求 $t\mid q$ 且 $\Delta=q/t$。给定两个 LWE 密文，其整数提升满足：

$\textsf{LWE}_{\vec{s}, \sigma}(\Delta m^{\langle 1 \rangle} + e^{\langle 1 \rangle}) = (\vec{a}^{\langle 1 \rangle}, b^{\langle 1 \rangle }) \Rightarrow \,\,\, \vec{a}^{\langle 1 \rangle} \cdot \vec{s} + b^{\langle 1 \rangle} = \Delta m^{\langle 1 \rangle} + e^{\langle 1 \rangle} + k_1q$

$\textsf{LWE}_{\vec{s}, \sigma}(\Delta m^{\langle 2 \rangle} + e^{\langle 2 \rangle})  = (\vec{a}^{\langle 2 \rangle}, b^{\langle 2 \rangle}) \Rightarrow \,\,\, \vec{a}^{\langle 2 \rangle} \cdot \vec{s} + b^{\langle 2 \rangle} = \Delta m^{\langle 2 \rangle} + e^{\langle 2 \rangle} + k_2q$

$ $

将这两个等式相乘，得到以下关系：


$(\Delta m^{\langle 1 \rangle} + e^{\langle 1 \rangle} + k_1q) \cdot (\Delta m^{\langle 2 \rangle} + e^{\langle 2 \rangle} + k_2q)$

$ = (\vec{a}^{\langle 1 \rangle} \cdot \vec{s} + b^{\langle 1 \rangle}) \cdot (\vec{a}^{\langle 2 \rangle} \cdot \vec{s} + b^{\langle 2 \rangle})$ 

$ = b^{\langle 1 \rangle}b^{\langle 2 \rangle}  + (b^{\langle 2 \rangle}\vec{a}^{\langle 1 \rangle} + b^{\langle 1 \rangle}\vec{a}^{\langle 2 \rangle}) \cdot \vec{s} + (\vec{a}^{\langle 1 \rangle} \cdot \vec{s}) \cdot (\vec{a}^{\langle 2 \rangle} \cdot \vec{s}) $



$ = \underbrace{b^{\langle 1 \rangle}b^{\langle 2 \rangle}}_{d_0}  + \underbrace{(b^{\langle 2 \rangle}\vec{a}^{\langle 1 \rangle} + b^{\langle 1 \rangle}\vec{a}^{\langle 2 \rangle})}_{d_1} \cdot \vec{s} + \underbrace{(\vec{a}^{\langle 1 \rangle} \otimes \vec{a}^{\langle 2 \rangle})}_{\vec{d}_2} \cdot (\vec{s} \otimes \vec{s}) $  


$= d_0 + d_1\cdot \vec{s} + \vec{d}_2\cdot \vec{s}^{\otimes}$ \textcolor{red}{ $\rhd$ 其中 $\vec{s}^{\otimes} = \vec{s} \otimes \vec{s}$}

$ $

，其中，$\otimes$ 表示两个向量的外积。例如，给定两个长度为 $n$ 的向量 $\vec{v}$ 和 $\vec{u}$，$\vec{v} \otimes \vec{u}$ 等价于将以下 $n$ 个不同的、长度为 $n$ 的向量连接而成的长度为 $n^2$ 的向量：$v_0\cdot\vec{u}, \,\, v_1\cdot\vec{u}, \,\, \cdots v_{n-1}\cdot\vec{u}$。注意，上述关系与 BFV 密文与密文相乘中推导出的关系类似。因此，类似于 BFV 的其余步骤，可以对这一 LWE 项按如下方式进行重线性化与重缩放：


\begin{tcolorbox}[title={\textbf{\tboxlabel{\ref*{subsubsec:bfv-mult-cipher-tfhe}} 分层 LWE 张量乘法}}]


假设有以下两个 LWE 密文：

$\textsf{LWE}_{\vec{s}, \sigma}(\Delta M^{\langle 1 \rangle}) = (\vec{a}^{\langle 1 \rangle}, b^{\langle 1 \rangle}) \bmod q$，\text{ } 其中 $b^{\langle 1 \rangle} = -\vec{a}^{\langle 1 \rangle} \cdot \vec{s} + \Delta m^{\langle 1 \rangle} + e^{\langle 1 \rangle} \bmod q$

$\textsf{LWE}_{\vec{s}, \sigma}(\Delta M^{\langle 2 \rangle}) = (\vec{a}^{\langle 2 \rangle}, b^{\langle 2 \rangle}) \bmod q$，\text{ } 其中 $b^{\langle 2 \rangle} = -\vec{a}^{\langle 2 \rangle} \cdot \vec{s} + \Delta m^{\langle 2 \rangle} + e^{\langle 2 \rangle} \bmod q$

$ $

这两个 LWE 密文的乘法按如下步骤执行：

$ $

\begin{enumerate}

\item \textbf{\underline{模数提升（\textsf{ModRaise}）}}

将 $(\vec{a}^{\langle 1 \rangle}, b^{\langle 1 \rangle})$ 和 $(\vec{a}^{\langle 2 \rangle}, b^{\langle 2 \rangle})$ 各分量选定的整数提升视为模 $Q=q\Delta$ 的元素。其相位仍包含 $qk_1$ 和 $qk_2$，并不是在较大模数下对相同缩放明文的新鲜加密：

$(\vec{a}^{\langle 1 \rangle}, b^{\langle 1 \rangle}) \bmod Q$

$(\vec{a}^{\langle 2 \rangle}, b^{\langle 2 \rangle}) \bmod Q$


$ $

\item \textbf{\underline{乘法}}

在模 $Q$ 下计算以下多项式乘法：

$d_0 = b^{\langle 1 \rangle}b^{\langle 2 \rangle} \bmod Q$

$d_1 = b^{\langle 2 \rangle}\vec{a}^{\langle 1 \rangle} +  b^{\langle 1 \rangle}\vec{a}^{\langle 2 \rangle} \bmod Q$

$\vec{d}_2 = \vec{a}^{\langle 1 \rangle} \otimes \vec{a}^{\langle 2 \rangle} \bmod Q$


$ $

\item \textbf{\underline{重线性化}}


计算：

$ \textsf{ct}_\alpha = (d
_1, d_0)$

$ \textsf{ct}_\beta = \bm{\langle} \textsf{Decomp}^{\beta, l}(\vec{d}_2), \text{ } \textsf{Lev}_{\vec{s}, \sigma}^{\beta, l}(\vec{s}^{\otimes}) \bm{\rangle}$.

$ \textsf{ct}_{\alpha + \beta} = \textsf{ct}_{\alpha} + \textsf{ct}_{\beta}$

$ $

于是，以下性质成立：

$\phi_{+}(\textsf{ct}_{\alpha+\beta})\equiv(\Delta m^{\langle1\rangle}+e^{\langle1\rangle}+qk_1)(\Delta m^{\langle2\rangle}+e^{\langle2\rangle}+qk_2)+e_{\rm rl}\pmod Q$，其中 $\phi_+(\vec a,b)=b+\langle\vec a,\vec s\rangle$。

$ $




\item \textbf{\underline{重缩放}}

将 $\textsf{ct}_{\alpha + \beta} = (\vec{a}_{\alpha + \beta}, b_{\alpha + \beta}) \bmod Q$ 更新为 $\textsf{ct'}_{\alpha + \beta} = \left(\left\lceil\dfrac{\vec{a}_{\alpha + \beta}}{\Delta}\right\rfloor, \left\lceil\dfrac{b_{\alpha + \beta}}{\Delta}\right\rfloor\right) \bmod q$。

$ $

这一明文重缩放过程也可以视为密文 $\textsf{ct}_{\alpha + \beta}$ 按 $Q \rightarrow q$ 进行模数切换。


$ $

\end{enumerate}

注意，密文与密文相乘后，明文尺度因子 $\Delta=\dfrac{q}{t}$、密文模数 $q$ 和私钥 $\vec{s}$ 均保持不变。



\end{tcolorbox}



\textbf{求值密钥与条件性正确性。}设 LWE 维数为 $k$，选择整数分解权重 $g_j=\beta^j$，其中 $0\leq j<L$，并要求 $\beta^L\geq Q$。对每个有序下标对 $(u,v)$，公开相位为 $g_js_us_v+e_{u,v,j}^{\rm key}\pmod Q$ 的行密文。将 $\vec d_2$ 的每个分量在模 $Q$ 下精确分解为绝对值不超过 $D_g$ 的数字。若 $|e_{u,v,j}^{\rm key}|\leq B_{\rm key}$，则所得重线性化误差满足
\[
|e_{\rm rl}|\leq k^2LD_gB_{\rm key}.
\]
该密钥不同于 TFHE 的 GGSW 自举密钥。输入分量与明文均取中心代表，记 $h=\|\vec s\|_1$，并设 $|e_i|\leq B_i$，则整数相位中的进位满足
\[
|k_i|\leq\kappa_i:=\frac{h+2}{2}+\frac{B_i}{q}.
\]
若 $R_a,R_b$ 是最后除以 $\Delta$ 时各分量的舍入误差，直接展开可得
\[
\begin{aligned}
\phi_+(\textsf{ct}_{\rm out})&\equiv\Delta m_1m_2+\eta\pmod q,\\
\eta&=m_1e_2+m_2e_1+t(e_1k_2+e_2k_1)
      +\frac{e_1e_2+e_{\rm rl}}\Delta
      +\langle R_a,\vec s\rangle+R_b.
\end{aligned}
\]
其中 $qk_1m_2$、$qk_2m_1$ 和 $qt k_1k_2$ 在模 $q$ 下消失。特别地，由于 $t\mid q$，一般 BFV 教学变体中危险的模数余数项在此消失。解码的一个充分条件是
\[
\frac t2(B_1+B_2)+t(B_1\kappa_2+B_2\kappa_1)
+\frac{B_1B_2+k^2LD_gB_{\rm key}}\Delta+\frac{h+1}{2}
<\frac\Delta2.
\]
明文模 $t$ 约减只增加 $q$ 的倍数，不再增加误差。这给出了上述运算顺序的条件性证明。若改为先用 $1/\Delta=t/q$ 缩放整数张量，再在模 $q$ 下重线性化，则密钥模数为 $q$，求值密钥误差不再除以 $\Delta$，张量分量舍入的界为 $(1+h+h^2)/2$。这两种顺序各有独立论证，密钥不能互换。附带的精确检查使用非零输入误差和求值密钥误差验证了两种顺序，但不据此断言参数达到安全要求或性能更快。

\subsection{同态密钥切换}
\label{subsec:bfv-key-switching}

分解基及其加密各行必须使用相同约定。对于一般模数 $q$，可使用整数权重 $\beta^{i-1}$，并采用足够多的位数精确分解；若采用截断分解，残差与旧秘密密钥的乘积也必须计入密钥切换误差。

BFV 的密钥切换方案将 RLWE 密文的秘密密钥从 $S$ 改为 $S'$。其本质上是 RLWE 的密钥切换方案，只是在加密和解密公式中将 $A\cdot S$ 项的符号取反。具体而言，它等价于 $k = 1$ 时的另一种 GLWE 版本（\autoref{subsec:glwe-alternative}）的密钥切换方案（\autoref{sec:glwe-key-switching}），如下所示：

\begin{tcolorbox}[title={\textbf{\tboxlabel{\ref*{subsec:bfv-key-switching}} BFV 的密钥切换}}]
$\textsf{RLWE}_{S',\sigma}(\Delta M) = (0, B) + \bm{\langle} \textsf{Decomp}^{\beta, l}(A), \text{ } \textsf{RLev}_{S', \sigma}^{\beta, l}(S) \bm{\rangle}$
\end{tcolorbox}




\subsection{输入向量槽位的同态旋转}
\label{subsec:bfv-rotation}


本节将介绍：当输入向量 $\vec{v}$ 已被编码为多项式并加密为 RLWE 密文后，如何对其元素进行同态旋转。在 \autoref{subsec:coeff-rotation} 中，我们学习了如何旋转多项式的系数。然而，旋转明文多项式 $M(X)$ 或 RLWE 密文多项式 $(A(X), B(X))$，不一定会旋转生成它们的输入向量。

对输入向量槽位（即输入向量）进行同态旋转的关键要求是：该运算必须作用于 RLWE 密文，并且在运算结束后，对 RLWE 密文进行解密和解码所恢复出的输入向量，应处于预期的旋转状态。我们将这一任务分为以下两个子问题：

\begin{enumerate}[leftmargin=3\parindent]
\item 如何通过将明文多项式 $M$ 更新为 $M'$，间接旋转输入向量？

\item 如何通过将 RLWE 密文多项式 $(A, B)$ 更新为 $(A', B')$，间接将明文多项式 $M$ 更新为 $M'$？
\end{enumerate}



\subsubsection{通过更新明文多项式旋转输入向量槽位}

在这一任务中，目标是修改明文多项式 $M(X)$，使输入向量 $\vec{v}$ 的前半部分元素在其内部向左循环移动 $h$ 个位置，同时使 $\vec{v}$ 的后半部分元素也在其内部向左循环移动 $h$ 个位置。具体而言，如果 $\vec{v}$ 定义如下：

$\vec{v} = (v_0, v_1, \cdots, v_{n-1})$

$ $

那么，将移动 $h$ 个位置后的向量 $\vec{v}^{\langle h \rangle}$ 记为：

$\vec{v}^{\langle h \rangle} = (\underbrace{v_h, v_{h+1}, \cdots, v_0, v_1, \cdots, v_{h-2}, v_{h-1},}_{\text{前半部分 $\dfrac{n}{2}$ 个元素向左旋转 $h$ 个位置}} \text{ } \underbrace{v_{\frac{n}{2}+h}, v_{\frac{n}{2}+h+1}, \cdots, v_{\frac{n}{2}+h-2}, v_{\frac{n}{2}+h-1}}_{\text{后半部分 $\dfrac{n}{2}$ 个元素向左旋转 $h$ 个位置}})$

$ $

回顾 \autoref{subsec:bfv-batch-encoding}，BFV 编码方案的组成部分如下：

$\vec{v} = (v_0, v_1, v_2, \cdots, v_{n-1})$ \textcolor{red}{ $\rhd$ 包含 $n$ 个槽位的输入向量}

$ $

$W =  \begin{bmatrix}
1 & 1 & 1 & \cdots & 1\\
(\omega) & (\omega^3) & (\omega^5) & \cdots & (\omega^{2n-1})\\
(\omega)^2 & (\omega^3)^2 & (\omega^5)^2 & \cdots & (\omega^{2n-1})^2\\
\vdots & \vdots & \vdots & \ddots & \vdots \\
(\omega)^{n-1} & (\omega^3)^{n-1} & (\omega^5)^{n-1} & \cdots & (\omega^{2n-1})^{n-1}\\
\end{bmatrix}$

$= \begin{bmatrix}
1 & 1 & \cdots & 1 & 1 & \cdots & 1 & 1\\
(\omega) & (\omega^3) & \cdots & (\omega^{n - 1}) & (\omega^{-(n - 1)}) & \cdots & (\omega^{-3}) & (\omega^{-1})\\
(\omega)^2 & (\omega^3)^2 & \cdots & (\omega^{n - 1})^2 & (\omega^{-(n - 1)})^2 & \cdots & (\omega^{-3})^2 & (\omega^{-1})^2\\
\vdots & \vdots & \vdots & \ddots & \vdots \\
(\omega)^{n-1} & (\omega^3)^{n-1} & \cdots & (\omega^{n - 1})^{n-1} & (\omega^{-(n - 1)})^{n-1} & \cdots & (\omega^{-3})^{n-1} & (\omega^{-1})^{n-1}\\
\end{bmatrix}$

，其中 $\omega = g^{\frac{t - 1}{2n}} \bmod t$（$g$ 是 $\mathbb{Z}_t^{\times}$ 的一个生成元）

\textcolor{red}{ $\rhd$ 将 $\vec{v}$ 转换为 $\vec{m}$ 的编码矩阵（即明文多项式 $M(X)$ 的 $n$ 个系数）}
 


$ $

$\Delta\vec{m}  = n^{-1}\cdot\Delta W\cdot I_n^R\cdot\vec{v}$  

\textcolor{red}{ $\rhd$ 包含明文多项式经缩放后的 $n$ 个整数系数的向量}

$ $

$\Delta M = \sum\limits_{i=0}^{n-1} ( \Delta m_i  X^i)$ 

\textcolor{red}{ $\rhd$ 对输入向量 $\vec{v}$ 进行同构编码得到的整数多项式}


$ $

由 \autoref{subsec:poly-vector-transformation} 可知，对多项式 $M(X)$ 进行解码 $\vec{v}$，等价于在以下 $n$ 个不同的本原 $(\mu=2n)$ 次单位根处计算 $M(X)$ 的值：$\{\omega, \omega^3, \omega^5, \cdots, \omega^{2n-3}, \omega^{2n-1} \}$。因此，上述解码过程等价于：

$\vec{v} = \dfrac{W^T \Delta \vec{m}}{\Delta} = \bm{(}W^T_0\cdot \vec{m}, \text{ } W^T_1\cdot \vec{m}, \text{ } W^T_2\cdot \vec{m}, \text{ } \cdots, \text{ } W^T_{n-1}\cdot \vec{m}\bm{)}$  

$= \bm{(} \text{ } M(\omega), \text{ } M(\omega^3), \text{ } M(\omega^5), \cdots, M(\omega^{2n-3}), \text{ } M(\omega^{2n-1}) \text{ } \bm{)}$


$= \bm{(} \text{ } M(\omega), \text{ } M(\omega^3), \text{ } M(\omega^5), \cdots, M(\omega^{n-3}), \text{ } M(\omega^{n-1}), \text{ } M(\omega^{-(n-1)}), \text{ } M(\omega^{-(n-3)}), \cdots, M(\omega^{-3}), \text{ } M(\omega^{-1}) \text{ } \bm{)}$


%%$\textsf{LWE}_{S, \sigma}(\Delta M) = (A, B) \in \mathcal{R}_{\langle n, q \rangle }^{2}$ \textcolor{red}{  $\rhd$ An LWE ciphertext that encrypts $\Delta M$}

%$A\cdot S + B = \Delta M + E$


$ $

接下来的任务是将 $M(X)$ 修改为 $M'(X)$，使解码 $M'(X)$ 得到经过修改的输入向量 $\vec{v}^{\langle h \rangle}$：它将 $\vec{v}$ 的前半部分元素在其内部向左循环旋转 $h$ 个位置，并将后半部分元素也在其内部向左循环旋转 $h$ 个位置。为实现这一旋转，我们采用以下两步方案：

\begin{enumerate}[leftmargin=3\parindent]
\item 为将 $M(X)$ 转换为 $M'(X)$，定义新的映射 $\sigma_M$ 如下：

$\sigma_M : (M(X), h) \in (\mathcal{R}_{\langle n, t \rangle}, \mathbb{Z}_{n}) \longrightarrow M'(X) \in \mathcal{R}_{\langle n, t \rangle}$

，其中，$h$ 是要对 $\vec{v}$ 施加的旋转位置数。
\item 为将 $M'(X)$ 解码为旋转后的输入向量 $\vec{v}^{\langle h \rangle}$，需要修改 \textsf{Encoding\textsubscript{1}}（\autoref{subsubsec:bfv-encoding-1}）中的同构映射 $\sigma : M(X) \in \mathcal{R}_{\langle n, t \rangle} \longrightarrow \vec{v} \in \mathbb{Z}^n$，重新设计解码方案。
\end{enumerate}

$ $

\para{将 $\bm{M(X)}$ 转换为 $\bm{M'(X)}$：} 第一个任务是将 $M(X)$ 转换为 $M'(X)$，等价于应用新的映射 $\sigma_M : (M(X), h) \in (\mathcal{R}_{\langle n, t \rangle}, \mathbb{Z}_{n}) \longrightarrow M'(X) \in \mathcal{R}_{\langle n, t \rangle}$，使解码 $M'(X)$ 得到旋转后的输入向量 $\vec{v}^{\langle h \rangle}$：$\vec{v}$ 的前半部分元素在其内部向左循环旋转 $h$ 个位置，后半部分元素也在其内部向左循环旋转 $h$ 个位置。为了设计满足这一要求的 $\sigma_M$，我们将使用具有以下两个特殊性质的数 $5^j$（这些性质来自数论）：
\begin{itemize}
\item 当整数 $j$ 满足 $0 \leq j < \dfrac{n}{2}$ 时，$(5^j \bmod 2n)$ 和 $(-5^j \bmod 2n)$ 生成 $[0, 2n)$ 之间的所有奇数。
\item 对每个满足 $0 \leq j < \dfrac{n}{2}$ 的整数 $j$，均有 $(5^j \bmod 2n) + (-5^j \bmod 2n) \equiv 0 \bmod 2n$。
\end{itemize}
例如，假设模数 $2n = 16$。那么，

\begin{multicols}{2}
\noindent $5^0 \bmod 16 = 1$
\newline $5^1 \bmod 16 = 5$
\newline $5^2 \bmod 16 = 9$
\newline $5^3 \bmod 16 = 13$
\newline $-(5)^0 \bmod 16 = 15$
\newline $-(5)^1 \bmod 16 = 11$
\newline $-(5)^2 \bmod 16 = 7$
\newline $-(5)^3 \bmod 16 = 3$
\end{multicols}


如上所示，$0 \leq j < 4$ 生成 $[0, 16)$ 之间的所有奇数。此外，对 $0 \leq j < 4$ 中的每个 $j$，均有 $(5^j \bmod 16) + (-5^j \bmod 16) = 16$。

$ $

定义 $J(h) = 5^h \bmod 2n$，\text{ } 以及 \text{ } $J_*(h) = -5^h \bmod 2n$。基于 $J(h)$ 和 $J_*(h)$，将映射 $\sigma_M : M(X) \rightarrow M'(X)$ 定义如下：

$\sigma_{M} : M(X) \rightarrow M(X^{J(h)})$

$ $

给定明文多项式 $M(X)$，为了使其解码得到的输入向量 $\vec{v}$ 的前半部分元素向左循环旋转 $h$ 个位置，后半部分元素也向左循环旋转 $h$ 个位置，我们通过应用 $\sigma_M$ 映射，将当前明文多项式 $M(X)$ 更新为新的多项式 $M'(X) = M(X^{J(h)}) = M(X^{5^h})$，其中，$h$ 是 $\vec{v}$ 的前、后半部分元素各自向左旋转的位置数。


$ $

\para{将 $\bm{M'(X)}$ 解码为 \boldmath$\vec{v}^{\langle h \rangle}$}：第二个任务是修改原始解码方案，使其能够成功地将 $M'(X)$ 解码为旋转后的输入向量 $\vec{v}^{\langle h \rangle}$。为此，我们将原来的同构映射 $\sigma : M(X) \longrightarrow \vec{v}$ 从：

$\sigma: M(X) \in \mathcal{R}_{\langle n, q \rangle} \longrightarrow  \bm{(}M(\omega), M(\omega^3),M(\omega^5), \cdots, M(\omega^{2n-1})\bm{)} \in \mathbb{Z}^n$ \textcolor{red}{ \text{ } \# 该映射在 \autoref{subsec:poly-vector-transformation} 中设计}

$ $

，改为：

$\sigma_J: M(X) \in \mathcal{R}_{\langle n, t \rangle} \longrightarrow  \bm{(}M(\omega^{J(0)}),M(\omega^{J(1)}),M(\omega^{J(2)}), \cdots, M(\omega^{J(\frac{n}{2} - 1)}), $

\textcolor{white}{$\sigma_J: M(X) \in \mathcal{R}_{\langle n, q \rangle} \longrightarrow  \bm{(}$} $M(\omega^{J_*(0)}), \text{ } M(\omega^{J_*(1)}), \text{ } M(\omega^{J_*(2)}), \cdots, M(\omega^{J_*(\frac{n}{2} - 1)}) \text{ } \bm{)} \in \mathbb{Z}^{n}$


$ $

$\sigma$ 与 $\sigma_J$ 的共同点在于：它们都在 $n$ 个不同的本原 $(\mu=2n)$ 次单位根处计算多项式 $M(X)$ 的值（即 $w^i$，其中 $i$ 取 $[0,2n]$ 之间的所有奇数）。对于 $\sigma_J$ 映射，注意当 $j$ 遍历 $0 \leq j < \dfrac{n}{2} $ \text{ } 时，$J(j) = 5^j \bmod 2n$ \text{ } 和 \text{ } $J_*(j) = -5^j \bmod 2n$ \text{ } 覆盖了 $[0, 2n]$ 之间的所有奇数。因此，在 $0 \leq j < \frac{n}{2}$ 范围内，$\omega^{J(j)}$ 和 $\omega^{J_*(j)}$ 覆盖了全部 $n$ 个不同的本原 $(\mu=2n)$ 次单位根。

同时，$\sigma$ 和 $\sigma_J$ 的差别在于输出向量元素的顺序。在 $\sigma$ 映射中，对 $M(X)$ 求值所用坐标的顺序为 $\omega, \omega^3, \cdots, \omega^{2n - 1}$；而在 $\sigma_J$ 映射中，求值所用坐标的顺序为 $\omega^{J(0)}, \omega^{J(1)}, \cdots, \omega^{J(\frac{n}{2} - 1)}, \omega^{J_*(0)}, \omega^{J_*(1)}, \cdots, \omega^{J_*(\frac{n}{2} - 1)}$。稍后将说明为何这样修改顺序。

在原始的 \textsf{Decoding\textsubscript{2}} 过程（\autoref{subsec:bfv-batch-encoding}）中，对明文多项式 $M(X)$ 应用 $\sigma$ 映射，等价于进行以下计算：

$\vec{v}_{'} = \bm{(} \text{ } M(\omega), \text{ } M(\omega^3), \text{ } M(\omega^5), \cdots, M(\omega^{2n-3}), \text{ } M(\omega^{2n-1}) \bm{)}$ 

$ $

$ = \bm{(} \text{ } M(\omega), \text{ } M(\omega^3), \text{ } M(\omega^5), \cdots, M(\omega^{n-1}), \text{ } M(\omega^{-(n-1)}), \cdots, \text{ } M(\omega^{-3}), \text{ } M(\omega^{-1}) \text{ } \bm{)}$ 

$ $

$= \bm{(} \text{ } W^T_0\cdot \vec{m}, \text{ } W^T_1\cdot \vec{m}, \text{ } W^T_2\cdot \vec{m}, \text{ } \cdots, \text{ } W^T_{n-1}\cdot \vec{m} \text{ } \bm{)} $ 
$ $ \textcolor{red}{ $\rhd$ 其中，$W_i^T$ 是 $W^T$ 的第 $(i+1)$ 行}

$ = W^T \vec{m} $
%$= \dfrac{I_n^R}{n}\cdot \bm{(} \text{ } M(\omega), \text{ } M(\omega^3), \text{ } M(\omega^5), \cdots, M(\omega^{2n-3}), \text{ } M(\omega^{2n-1}) \text{ } )$




$ $

类似地，对明文多项式 $M(X)$ 应用修改后的 $\sigma_J$ 映射，等价于进行以下计算：

$\vec{v}_{'} =  \bm{(} \text{ } 
M(\omega^{J(0)}), \text{ } M(\omega^{J(1)}), \text{ } M(\omega^{J(2)}), \cdots,  M(\omega^{J(\frac{n}{2}-1)}), \text{ } M(\omega^{J_*(0)}), \text{ } M(\omega^{J_*(1)}), \cdots,  M(\omega^{J_*(\frac{n}{2}-1)}) \text{ } \bm{)}$


$ $

$= \bm{(} \text{ } \hathat{W}^*_0\cdot \vec{m}, \text{ } \hathat{W}^*_1\cdot \vec{m}, \text{ } \hathat{W}^*_2\cdot \vec{m}, \text{ } \cdots, \text{ } \hathat{W}^*_{n-1}\cdot \vec{m} \text{ } \bm{)} $ 

$ $

$ = \hathat{W}^* \vec{m} $，其中

$ $

\noindent $\hathat{W}^* = \begin{bmatrix}
1 & (\omega^{J(0)}) & (\omega^{J(0)})^2 & \cdots & (\omega^{J(0)})^{n-1}\\
1 & (\omega^{J(1)}) & (\omega^{J(1)})^2 & \cdots & (\omega^{J(1)})^{n-1}\\
1 & (\omega^{J(2)}) & (\omega^{J(2)})^2 & \cdots & (\omega^{J(2)})^{n-1}\\
\vdots & \vdots & \vdots & \ddots & \vdots \\
1 & (\omega^{J(\frac{n}{2}-1)}) & (\omega^{J(\frac{n}{2}-1)})^2 & \cdots & (\omega^{J(\frac{n}{2}-1)})^{n-1}\\
1 & (\omega^{J_*(0)}) & (\omega^{J_*(0)})^2 & \cdots & (\omega^{J_*(0)})^{n-1}\\
1 & (\omega^{J_*(1)}) & (\omega^{J_*(1)})^2 & \cdots & (\omega^{J_*(1)})^{n-1}\\
1 & (\omega^{J_*(2)}) & (\omega^{J_*(2)})^2 & \cdots & (\omega^{J_*(2)})^{n-1}\\
\vdots & \vdots & \vdots & \ddots & \vdots \\
1 & (\omega^{J_*(\frac{n}{2}-1)}) & (\omega^{J_*(\frac{n}{2}-1)})^2 & \cdots & (\omega^{J_*(\frac{n}{2}-1)})^{n-1}\\
\end{bmatrix}$

$ $

$ $

\noindent $\hathat{W} = \begin{bmatrix}
1 & 1 & \cdots & 1 & 1 & 1 & \cdots & 1\\
(\omega^{J(\frac{n}{2} - 1)}) & (\omega^{J(\frac{n}{2} - 2)}) & \cdots & (\omega^{J(0)}) & (\omega^{J_*(\frac{n}{2} - 1)}) & (\omega^{J_*(\frac{n}{2} - 2)}) & \cdots & (\omega^{J_*(0)})\\
(\omega^{J(\frac{n}{2} - 1)})^2 & (\omega^{J(\frac{n}{2} - 2)})^2 & \cdots & (\omega^{J(0)})^2 & (\omega^{J_*(\frac{n}{2} - 1)})^2 & (\omega^{J_*(\frac{n}{2} - 2)})^2 & \cdots & (\omega^{J_*(0)})^2 \\
\vdots & \vdots & \ddots & \vdots & \vdots & \vdots & \vdots & \vdots \\
(\omega^{J(\frac{n}{2} - 1)})^{n-1} & (\omega^{J(\frac{n}{2} - 2)})^{n-1} & \cdots & (\omega^{J(0)})^{n-1} & (\omega^{J_*(\frac{n}{2} - 1)})^{n-1} & (\omega^{J_*(\frac{n}{2} - 2)})^{n-1} & \cdots  & (\omega^{J_*(0)})^{n-1}
\end{bmatrix}$

$ $

$ $

从现在开始，将 \textsf{Encoding\textsubscript{1}} 过程（\autoref{subsec:bfv-batch-encoding}）中的 $W$ 替换为 $\hathat{W}$，并将 \textsf{Decoding\textsubscript{2}} 过程中的 $W^T$ 替换为 $\hathat{W}^*$。

$ $

为证明 $\hathat{W}$ 和 $W$ 一样是有效的编码矩阵，且 $\hathat{W}^*$ 和 $W^T$ 一样是有效的解码矩阵，需要证明以下两点：
\begin{itemize}[leftmargin=3\parindent]

\item \textbf{${\hathat{\bm W}}$ 是 $\bm{n}$ 维向量空间的一组基：} 这一点成立，因为 $\hathat{W}$ 只是对 $W$ 的列重新排序，仍然是 $\bm{n}$ 维向量空间的一组基。

\item \boldmath{$\hathat{W}^{*} \cdot \hathat{W} = n \cdot I_n^R$（满足 \autoref{subsec:vandermonde-euler} 中的定理~\ref*{subsec:vandermonde-euler}）}：这一证明分为两个子证明：
\begin{enumerate}[leftmargin=2\parindent]
\item \boldmath{${\hathat{ W}^{ *} \cdot \hathat{ W}}$ 的反对角线上各值为 $\bm{n}$：} \boldmath$\hathat{ W}^{ *} \cdot \hathat{ W}$ 的反对角线上每个元素的计算结果均为 $\sum\limits_{i=0}^{n-1} \omega^{2nik} = n$，其中 $k$ 为某个整数。
\item \boldmath{${\hathat{ W}^{ *} \cdot \hathat{ W}}$ 的其余位置各值均为 0：} 除反对角线上的元素外，其余所有元素均为 $\sum\limits_{i=0}^{n-1}r^i=\dfrac{r^n-1}{r-1}=0\quad(r=\omega^{a+b}\ne1)$（由等比数列求和公式可得）。
\end{enumerate}

我们提供了 \href{https://github.com/wang2000star/fhe-textbook-cn/blob/chinese-edition/source%20code/bfv_j_matrix_inverse_proof.py}{\underline{Python 脚本}}，用于通过实验验证这一点。

\end{itemize}




因此，$\hathat{W}^*$ 和 $\hathat{W}$ 是将 $\vec{v}$ 转换为 $M(X)$ 的有效编码与解码矩阵。

$ $

现在考虑 $\vec{v}^{\langle h \rangle}$ 的结构（即将 $\vec{v}$ 的前半部分元素在其内部向左循环旋转 $h$ 个位置，同时将后半部分元素也在其内部向左循环旋转 $h$ 个位置）。回顾 $\vec{v}$ 的表达式：

$\vec{v} = \bm{(} \text{ } 
M(\omega^{J(0)}), \text{ } M(\omega^{J(1)}), \cdots,  M(\omega^{J(\frac{n}{2}-1)}), \text{ } 
M(\omega^{J_*(0)}), \text{ } M(\omega^{J_*(1)}),  \cdots,  M(\omega^{J_*(\frac{n}{2}-1)}) \bm{)}$


\text{ } \text{ } $= \bm{(} \text{ } \hathat{W}^*_0\cdot \vec{m}, \text{ } \hathat{W}^*_1\cdot \vec{m}, \text{ } \hathat{W}^*_2\cdot \vec{m}, \text{ } \cdots, \text{ } \hathat{W}^*_{n - 1}\cdot \vec{m} \text{ } \bm{)} $ 

$ $

因此，分别将 $\vec{v}$ 的前、后半部分元素向左旋转 $h$ 个位置，所得 $\vec{v}^{\langle h \rangle}$ 的状态如下：

$\vec{v}^{\langle h \rangle} = \bm{(} \text{ } M(\omega^{J(h)}), \text{ } M(\omega^{J(h+1)}), \cdots, M(\omega^{J(\frac{n}{2}-1)}), \text{ } M(\omega^{J(0)}), \text{ }M(\omega^{J(1)}), \cdots, \text{ } M(\omega^{J(h-2)}), \text{ } M(\omega^{J(h-1)}),$

\text{ }\text{ }\text{ } $ M(\omega^{J_*(h)}), \text{ } M(\omega^{J_*(h+1)}), \cdots, M(\omega^{J_*(\frac{n}{2}-1)}), \text{ } M(\omega^{J_*(0)}), \text{ }M(\omega^{J_*(1)}), \cdots, \text{ } M(\omega^{J_*(h-2)}), \text{ } M(\omega^{J_*(h-1)}) \text{ } \bm{)}$

$ $

注意，上述 $\vec{v}^{\langle h \rangle}$ 的计算等价于：将 $\hathat{W}^*$ 的上方 $\frac{n}{2}$ 行在其内部向上循环旋转 $h$ 个位置，将 $\hathat{W}^{*}$ 的下方 $\frac{n}{2}$ 行在其内部向上循环旋转 $h$ 个位置，再将得到的矩阵与 $\vec{m}$ 相乘。然而，实际中不希望如此直接修改解码矩阵 $\hathat{W}^*$，因为这样解码矩阵就无法保持一致。因此，我们将通过把 $\vec{m}$ 修改为 $\vec{m}_{'}$（即把 $M(X)$ 修改为 $M'(X)$）来保持关系 $\vec{v}^{\langle h \rangle} = \hathat{W}^* \cdot \vec{m}_{'}$ 成立，而不直接修改 $\hathat{W}^*$。提取 $\hathat{W}^*$ 的上半部分行，将这一 $\dfrac{n}{2} \times n$ 矩阵记为 $\hathat{H}_1^*$。那么，$\hathat{H}_1^*$ 等价于形式为 $V(\omega^{J(0)}, \omega^{J(1)}, \cdots, \omega^{J(\frac{n}{2}-1)})$ 的范德蒙德矩阵（\autoref{subsec:vandermonde} 中的定义~\ref*{subsec:vandermonde}）。类似地，提取 $\hathat{W}^*$ 的下半部分行，将这一 $\dfrac{n}{2} \times n$ 矩阵记为 $\hathat{H}_2^*$。那么，$\hathat{H}_2^*$ 等价于形式为 $V(\omega^{J_*(0)}, \omega^{J_*(1)}, \cdots, \omega^{J_*(\frac{n}{2}-1)})$ 的范德蒙德矩阵（\autoref{subsec:vandermonde} 中的定义~\ref*{subsec:vandermonde}）。

现在，将 $\hathat{H}_1^*$ 的各行向上旋转 $h$ 个位置，并记为 $\hathat{H}_1^{*\langle h \rangle}$；将 $\hathat{H}_2^*$ 的各行向上旋转 $h$ 个位置，并记为 $\hathat{H}_2^{*\langle h \rangle}$。将由 $\vec{v}^{\langle h \rangle}$ 的前半部分元素构成的 $\dfrac{n}{2}$ 维向量记为 $\vec{v}_1^{\langle h \rangle}$，将由 $\vec{v}^{\langle h \rangle}$ 的后半部分元素构成的 $\dfrac{n}{2}$ 维向量记为 $\vec{v}_2^{\langle h \rangle}$。
那么，计算（即解码）$\vec{v}_1^{\langle h \rangle} = \hathat{H}_1^{*\langle h \rangle}\cdot \vec{m}$，等价于将 $M(X)$ 修改为 $M'(X) = M(X^{J(h)})$（其系数向量为 $\vec{m}^{\langle h \rangle}$）后，再计算（即解码）$\vec{v}_1^{\langle h \rangle} = \hathat{H}_1^{*}\cdot \vec{m}^{\langle h \rangle}$。原因如下：

$\vec{v}_1^{\langle h \rangle} =\hathat{H}_1^{*\langle h \rangle}\cdot \vec{m}$

$ = \bm{(} \text{ } \hathat{W}^*_{h}\cdot\vec{m}, \text{ } \hathat{W}^*_{h+1}\cdot\vec{m}, \text{ } \hathat{W}^*_{h+2}\cdot\vec{m}, \cdots, \hathat{W}^*_{\frac{n}{2}-1}\cdot\vec{m}, \text{ } \hathat{W}^*_{0}\cdot\vec{m}, \text{ }\hathat{W}^*_{1}\cdot\vec{m}, \cdots, \text{ } \hathat{W}^*_{h-2}\cdot\vec{m}, \text{ } \hathat{W}^*_{h-1}\cdot\vec{m}  \bm{)}$

$ $

$= \bm{(} \text{ } M((\omega^{J(h)})^{J(0)}), \text{ } M((\omega^{J(h)})^{J(1)}), \text{ } M((\omega^{J(h)})^{J(2)}), \cdots, M((\omega^{J(h)})^{J(\frac{n}{2}-1)}) \text{ } \bm{)}$


\textcolor{red}{ $\rhd$ 这等价于在以下 $\dfrac{n}{2}$ 个不同的 $(\mu=2n)$ 次单位根处计算多项式 $M(X^{J(h)})$ 的值：$\omega^{J(0)}, \omega^{J(1)},  \omega^{J(2)}, \cdots, \omega^{J(\frac{n}{2} - 1)}$}

$ $

$= \bm{(} \text{ } M(\omega^{J(h)\cdot J(0)}), \text{ } M(\omega^{J(h)\cdot J(1)}), \text{ } M(\omega^{J(h)\cdot J(2)}), \cdots, M(\omega^{J(h)\cdot J(\frac{n}{2}-1)}) \text{ } \bm{)}$

$= \bm{(} \text{ } M(\omega^{5^h\cdot 5^0}), \text{ } M(\omega^{5^h\cdot 5^1}), \text{ } M(\omega^{5^h\cdot 5^2}), \cdots, M(\omega^{5^h\cdot 5^{n/2-1}}) \text{ } \bm{)}$

$= \bm{(} \text{ } M(\omega^{5^{h}}), \text{ } M(\omega^{5^{h+1}}), \text{ } M(\omega^{5^{h+2}}), \cdots, M(\omega^{5^{h+n/2-1}}) \text{ } \bm{)}$ \textcolor{red}{ $\rhd$ 注意 $5^{\frac{n}{2}} \bmod 2n = 1$}

$ $

$= \bm{(} \text{ } M(\omega^{J(h)}), \text{ } M(\omega^{J(h+1)}), \text{ } M(\omega^{J(h+2)}), \cdots, M(\omega^{J(\frac{n}{2}-1)}), \text{ } M(\omega^{J(0)}), \text{ }M(\omega^{J(1)}),$ 
                                    
$\cdots, \text{ } M(\omega^{J(h-2)}), \text{ } M(\omega^{J(h-1)})  \bm{)}$ 

$ $

$= \hathat{H}_1^{*}\cdot \vec{m}^{\langle h \rangle}$ \textcolor{red}{ \text{ } \# 其中，$\vec{m}^{\langle h \rangle}$ 包含多项式 $M(X^{J(h)})$ 的 $n$ 个系数 }


$ $

类似地，计算（即解码）$\vec{v}_2^{\langle h \rangle} = \hathat{H}_2^{*\langle h \rangle}\cdot \vec{m}$，等价于将 $M(X)$ 修改为 $M'(X) = M(X^{J(h)})$（其系数向量为 $\vec{m}^{\langle h \rangle}$）后，再计算（即解码）$\vec{v}_2^{\langle h \rangle} = \hathat{H}_2^{*}\cdot \vec{m}^{\langle h \rangle}$。原因如下：

$\vec{v}_2^{\langle h \rangle} =\hathat{H}_2^{*\langle h \rangle}\cdot \vec{m}$

$ $


$ = \bm{(} \text{ } \hathat{W}^*_{\frac{n}{2} +h}\cdot\vec{m}, \text{ } \hathat{W}^*_{\frac{n}{2} + h+1}\cdot\vec{m}, \text{ } \hathat{W}^*_{\frac{n}{2} +h+2}\cdot\vec{m}, \cdots, \hathat{W}^*_{n -1}\cdot\vec{m}, \text{ } \hathat{W}^*_{\frac{n}{2}}\cdot\vec{m}, \text{ }\hathat{W}^*_{\frac{n}{2} + 1}\cdot\vec{m}, \cdots,$

\text{ }\text{ }\text{ }\text{ }$ \text{ } \hathat{W}^*_{\frac{n}{2} + h-2}\cdot\vec{m}, \text{ } \hathat{W}^*_{\frac{n}{2} + h-1}\cdot\vec{m}  \bm{)}$

$ $

$= \bm{(} \text{ } M((\omega^{J(h)})^{J_*(0)}), \text{ } M((\omega^{J(h)})^{J_*(1)}), \text{ } M((\omega^{J(h)})^{J_*(2)}), \cdots, M((\omega^{J(h)})^{J_*(\frac{n}{2}-1)}) \text{ } \bm{)}$


\textcolor{red}{ $\rhd$ 这等价于在以下 $\dfrac{n}{2}$ 个不同的 $(\mu=2n)$ 次单位根处计算多项式 $M(X^{J(h)})$ 的值：$\omega^{J_*(0)}, \omega^{J_*(1)},  \omega^{J_*(2)}, \cdots, \omega^{J_*(\frac{n}{2} - 1)}$}

$ $

$= \bm{(} \text{ } M(\omega^{J(h)\cdot J_*(0)}), \text{ } M(\omega^{J(h)\cdot J_*(1)}), \text{ } M(\omega^{J(h)\cdot J_*(2)}), \cdots, M(\omega^{J(h)\cdot J_*(\frac{n}{2}-1)}) \text{ } \bm{)}$

$= \bm{(} \text{ } M(\omega^{5^h\cdot -5^0}), \text{ } M(\omega^{5^h\cdot -5^1}), \text{ } M(\omega^{5^h\cdot -5^2}), \cdots, M(\omega^{5^h\cdot -5^{n/2-1}}) \text{ } \bm{)}$

$= \bm{(} \text{ } M(\omega^{-5^{h}}), \text{ } M(\omega^{-5^{h+1}}), \text{ } M(\omega^{-5^{h+2}}), \cdots, M(\omega^{-5^{h+n/2-1}}) \text{ } \bm{)}$ \textcolor{red}{ $\rhd$ 注意 $-(5^{\frac{n}{2}} \bmod 2n) = -1$}

$ $

$= \bm{(} \text{ } M(\omega^{J_*(h)}), \text{ } M(\omega^{J_*(h+1)}), \text{ } M(\omega^{J_*(h+2)}), \cdots, M(\omega^{J_*(\frac{n}{2}-1)}), \text{ } M(\omega^{J_*(0)}), \text{ }M(\omega^{J_*(1)}),$ 
                                    
$\cdots, \text{ } M(\omega^{J_*(h-2)}), \text{ } M(\omega^{J_*(h-1)})  \bm{)}$ 

$ $

$= \hathat{H}_2^{*}\cdot \vec{m}^{\langle h \rangle}$


$ $

上述推导表明，$\vec{v}_1^{\langle h \rangle} = \hathat{H}_1^{*}\cdot \vec{m}^{\langle h \rangle}$，且 $\vec{v}_2^{\langle h \rangle} = \hathat{H}_2^{*}\cdot \vec{m}^{\langle h \rangle}$。结合这两点，得到结论 $\vec{v}^{\langle h \rangle} = \hathat{W}^* \cdot \vec{m}^{\langle h \rangle}$：将输入向量 $\vec{v}$ 的前半部分元素向左旋转 $h$ 个位置、后半部分元素也向左旋转 $h$ 个位置，等价于将明文多项式 $M(X)$ 更新为 $M(X^{J(h)})$，再使用解码矩阵 $\hathat{W}^*$ 进行解码。



然而，现在出现了一个新问题：无法直接将明文 $M(X)$ 更新为 $M(X^{J(h)})$，因为 $M(X)$ 已被加密为 RLWE 密文。因此，需要更新 RLWE 密文分量 $(A, B)$，从而\textit{间接地}将 $M(X)$ 更新为 $M(X^{J(h)})$。下一小节将介绍具体方法。


\subsubsection{通过更新密文多项式来更新明文多项式}

给定 RLWE 密文 $\textsf{ct} = (A, B)$，目标是将 $\textsf{ct} = (A, B)$ 更新为 $C^{\langle h \rangle} = (A^{\langle h \rangle}, B^{\langle h \rangle})$，使其解密后得到输入向量 $\vec{v}^{\langle h \rangle}$。也就是说，应满足以下关系：

$\textsf{RLWE}^{-1}_{S, \sigma}\bm{(} \text{ } C^{\langle h \rangle} = (A^{\langle h \rangle}, B^{\langle h \rangle}) \text{ }\bm{)} = \Delta M(X^{J(h)}) + E'$

$ $

回顾 RLWE 密码系统（\autoref{sec:rlwe}）的另一种版本（\autoref{subsec:glwe-alternative}），明文与密文对满足以下关系：

$\Delta M(X) + E(X) = A(X)\cdot S(X) + B(X) \approx \Delta M(X)$

$ $

在上述关系中应用 $X = X^{J(h)}$，可推导出以下关系：

$\Delta M(X^{J(h)}) + E(X^{J(h)}) = A(X^{J(h)}) \cdot S(X^{J(h)}) + B(X^{J(h)}) \approx \Delta M(X^{J(h)})$

$ $

这一关系意味着，以 $S(X^{J(h)})$ 为秘密密钥解密密文 $C^{\langle h \rangle} = (A(X^{J(h)}), B(X^{J(h)}))$，即可得到 $\Delta M(X^{J(h)})$。因此，$C^{\langle h \rangle} = (A(X^{J(h)}), B(X^{J(h)}))$ 就是我们要寻找的 RLWE 密文，因为对它解密，再对其明文 $M(X^{J(h)})$ 解码，就能得到输入向量 $\vec{v}^{\langle h \rangle}$。

逐分量应用自同构，得到密文 $C^{\langle h\rangle}=(A(X^{J(h)}),B(X^{J(h)}))$，其秘密密钥为 $S(X^{J(h)})$。随后，使用 \autoref{subsec:bfv-key-switching} 中的小结~\ref*{subsec:bfv-key-switching} 将密钥切换回原密钥 $S(X)$。这会产生新的密文分量，并增加密钥切换误差；未经改变的同一分量对并不同时对应两个秘密密钥。下式只是在明文记号中省略该误差：

$ $

\noindent $\underbrace{\textsf{RLWE}_{S(X),\sigma}\bm{(}\Delta M(X^{J(h)})\bm{)}}_{\substack{\text{以下运算的结果：} \\ \text{明文与密文相加}}} = \underbrace{\bm{(} \text{ } 0, B(X^{J(h)}) \text{ } \bm{)}}_{\substack{\text{明文 } B(X^{J(h)}) \\ \text{（平凡密文）}}} + \bm{\langle} \text{ } \underbrace{\textsf{Decomp}^{\beta, l}\bm{(}A(X^{J(h)})\bm{)}, \text{ } \textsf{RLev}_{S(X), \sigma}^{\beta, l}\bm{(}S(X^{J(h)})\bm{)} \text{ } \bm{\rangle}}_{\substack{\substack{\text{加密 $A(X^{J(h)}) \cdot S(X^{J(h)})$ 的 RLWE 密文}\\ \text{其密钥已按 $S(X^{J(h)})\rightarrow S(X)$ 切换}}}}$



\subsubsection{小结}
\label{subsubsec:bfv-rotation-summary}

BFV 输入向量的旋转过程总结如下：

\begin{tcolorbox}[title={\textbf{\tboxlabel{\ref*{subsec:bfv-rotation}} BFV 输入向量槽位的同态旋转}}]

为支持输入向量槽位旋转，将 \textsf{Encoding\textsubscript{1}} 中的原始编码矩阵更新如下：

$ $

{\footnotesize{\noindent $\hathat{W} = \begin{bmatrix}
1 & 1 & \cdots & 1 & 1 & 1 & \cdots & 1\\
(\omega^{J(\frac{n}{2} - 1)}) & (\omega^{J(\frac{n}{2} - 2)}) & \cdots & (\omega^{J(0)}) & (\omega^{J_*(\frac{n}{2} - 1)}) & (\omega^{J_*(\frac{n}{2} - 2)}) & \cdots & (\omega^{J_*(0)})\\
(\omega^{J(\frac{n}{2} - 1)})^2 & (\omega^{J(\frac{n}{2} - 2)})^2 & \cdots & (\omega^{J(0)})^2 & (\omega^{J_*(\frac{n}{2} - 1)})^2 & (\omega^{J_*(\frac{n}{2} - 2)})^2 & \cdots & (\omega^{J_*(0)})^2 \\
\vdots & \vdots & \ddots & \vdots & \vdots & \vdots & \ddots & \vdots \\
(\omega^{J(\frac{n}{2} - 1)})^{n-1} & (\omega^{J(\frac{n}{2} - 2)})^{n-1} & \cdots & (\omega^{J(0)})^{n-1} & (\omega^{J_*(\frac{n}{2} - 1)})^{n-1} & (\omega^{J_*(\frac{n}{2} - 2)})^{n-1} & \cdots  & (\omega^{J_*(0)})^{n-1}
\end{bmatrix}$}}

$ $

$ $


，并将 \textsf{Decoding\textsubscript{2}} 中的原始解码矩阵更新如下：

$\hathat{W}^* = \begin{bmatrix}
1 & (\omega^{J(0)}) & (\omega^{J(0)})^2 & \cdots & (\omega^{J(0)})^{n-1}\\
1 & (\omega^{J(1)}) & (\omega^{J(1)})^2 & \cdots & (\omega^{J(1)})^{n-1}\\
1 & (\omega^{J(2)}) & (\omega^{J(2)})^2 & \cdots & (\omega^{J(2)})^{n-1}\\
\vdots & \vdots & \vdots & \ddots & \vdots \\
1 & (\omega^{J(\frac{n}{2}-1)}) & (\omega^{J(\frac{n}{2}-1)})^2 & \cdots & (\omega^{J(\frac{n}{2}-1)})^{n-1}\\
1 & (\omega^{J_*(0)}) & (\omega^{J_*(0)})^2 & \cdots & (\omega^{J_*(0)})^{n-1}\\
1 & (\omega^{J_*(1)}) & (\omega^{J_*(1)})^2 & \cdots & (\omega^{J_*(1)})^{n-1}\\
1 & (\omega^{J_*(2)}) & (\omega^{J_*(2)})^2 & \cdots & (\omega^{J_*(2)})^{n-1}\\
\vdots & \vdots & \vdots & \ddots & \vdots \\
1 & (\omega^{J_*(\frac{n}{2}-1)}) & (\omega^{J_*(\frac{n}{2}-1)})^2 & \cdots & (\omega^{J_*(\frac{n}{2}-1)})^{n-1}\\
\end{bmatrix}$

$ $

$ $

，其中，$J(h)$ 是\textit{旋转辅助公式}：$J(h) = 5^h \bmod 2n$，
\text{ } $J_*(h) = -5^h \bmod 2n$. 

$ $

使用 $\hathat{W}$ 进行编码：$\vec{m} = n^{-1}\cdot \hathat{W} \cdot I_n^R \cdot \vec{v}$。使用 $\hathat{W}^*$ 进行解码：$\vec{v} = \hathat{W}^* \cdot \vec{m}$

$ $

假设有如下 RLWE 密文和密钥切换密钥：

$\textsf{RLWE}_{S, \sigma}(\Delta M) = (A, B)$, \text{ } $\textsf{RLev}_{S, \sigma}^{\beta, l}\bm(S(X^{J(h)})\bm)$

$ $

那么，要使密文所对应原始输入向量 $\vec{v}$ 的前半部分元素在其内部向左循环旋转 $h$ 个位置，同时使 $\vec{v}$ 的后半部分元素也在其内部向左循环旋转 $h$ 个位置，过程如下：

\begin{enumerate}
\item 将 $A(X)$、$B(X)$ 更新为 $A(X^{J(h)})$、$B(X^{J(h)})$。
\item 执行如下从 $S(X^{J(h)})$ 到 $S(X)$ 的密钥切换（\autoref{subsec:bfv-key-switching}）：

$\textsf{RLWE}_{S(X),\sigma}\bm{(}\Delta M(X^{J(h)})\bm{)} = \bm{(} 0, B(X^{J(h)}) \bm{)} \text{ } + \text{ } \bm{\langle}  \textsf{Decomp}^{\beta, l}\bm{(}A(X^{J(h)})\bm{)}, \text{ } \textsf{RLev}_{S(X), \sigma}^{\beta, l}\bm{(}S(X^{J(h)})\bm{)} \bm{\rangle}$
\end{enumerate}


\end{tcolorbox}


\subsubsection{编码示例}
\label{subsubsec:bfv-batch-encoding-ex}

%$bfv_vector = [10, 0, 5, 13]
%bfv_vector2 = [2, 4, 8, 6]


在上述示例中，计算 $\bmod t$ 和 $\bmod q$ 时采用无符号模表示（例如 $[0, t-1], [0, q]$）。不过，即使使用居中模表示（例如 $[-\dfrac{t}{2}, \dfrac{t}{2})$、$[-\dfrac{q}{2}, \dfrac{q}{2})$），结果的正确性仍然成立。在许多实际的全同态加密库中，通常使用居中模运算来更有效地管理噪声增长。

假设采用如下设置：

$\mu = 8, \, n = \dfrac{\mu}{2} = 4, \, t = 17, \, q = 2^8 = 256, \, n^{-1} = 13, \, \Delta = \left\lfloor\dfrac{q}{t}\right\rfloor = 15$

$ $

$\mathcal{R}_{\langle 4, 17 \rangle} = \mathbb{Z}_{17}[X] / (X^4 + 1)$

$ $

$X^4 + 1 \pmod{17}$ 的根为 $X = \{2, 8, 15, 9\}$，验证如下：

$2^4 \equiv 8^4 \equiv 15^4 \equiv 9^4 \equiv 16 \equiv -1 \bmod{17}$

$ $

定义~\ref*{subsec:cyclotomic-def}（位于 \autoref{subsec:cyclotomic-def}）指出，第 $\mu$ 个分圆多项式的根是本原 $\mu$ 次单位根。而定义~\ref*{subsec:roots-def}（位于 \autoref{subsec:roots-def}）指出，本原 $\mu$ 次单位根的阶为 $\mu$。这些定义同时适用于 $X\in\mathbb{C}$（复数）上的分圆多项式和 $X \in \mathbb{Z}_t$（环）上的分圆多项式。

由于 $\{2, 8, 15, 9\}$ 是环 $\mathbb{Z}_{17}$ 上第 $(\mu=8)$ 个分圆多项式 $X^4 + 1$ 的根，因此它们也是 $\mathbb{Z}_{17}$ 的本原 $(\mu=8)$ 次单位根。所以，它们的阶（\autoref{subsec:order-def}）为 $\mu=8$，验证如下：

$2^8 \equiv 8^8 \equiv 15^8 \equiv 9^8 \equiv 1 \bmod 17$

$2^4 \equiv 8^4 \equiv 15^4 \equiv 9^4 \equiv 16 \not\equiv 1 \bmod 17$

$ $

定义~\ref*{subsec:cyclotomic-def}（位于 \autoref{subsec:cyclotomic-def}）和定理~\ref*{subsec:cyclotomic-theorem}（\autoref{subsec:cyclotomic-theorem}）还指出，对于每个本原 $\mu$ 次单位根 $\omega$，$\{\omega^k\}_{\textsf{gcd}(k, \mu) = 1}$ 生成第 $\mu$ 个分圆多项式的全部根。注意，对于第 $(\mu=8)$ 个分圆多项式 $X^4 + 1$，它的根 $\{2, 8, 15, 9\}$ 按如下方式生成全部 $(\mu=8)$ 次单位根：

$\{2^1, 2^3, 2^5, 2^7\} \equiv \{8^1, 8^3, 8^5, 8^7\} \equiv \{15^1, 15^3, 15^5, 15^7\} \equiv \{9^1, 9^3, 9^5, 9^7\} \equiv \{2, 8, 15, 9\} \bmod 17$

$ $

在 $X^4 + 1$ 的根 $\{2, 8, 15, 9\}$ 中，选取 $\omega = 9$ 作为基础根，按如下方式构造编码矩阵 $\hathat{W}$ 和解码矩阵 $\hathat{W}^*$：

$\hathat{W} = \begin{bmatrix}
1 & 1 & 1 & 1\\
\omega^{J(1)} & \omega^{J(0)} & \omega^{J_*(1)} & \omega^{J_*(0)}\\
(\omega^{J(1)})^2 & (\omega^{J(0)})^2 & (\omega^{J_*(1)})^2 & (\omega^{J_*(0)})^2\\
(\omega^{J(1)})^3 & (\omega^{J(0)})^3 & (\omega^{J_*(1)})^3 & (\omega^{J_*(0)})^3\\
\end{bmatrix}$ \textcolor{red}{  $\rhd$ 其中 $J(h) = 5^h \bmod 8$}

$ = \begin{bmatrix}
1 & 1 & 1 & 1\\
9^{5} & 9^{1} & 9^{3} & 9^{7}\\
(9^{5})^2 & (9^{1})^2 & (9^{3})^2 & (9^{7})^2\\
(9^{5})^3 & (9^{1})^3 & (9^{3})^3 & (9^{7})^3\\
\end{bmatrix} \equiv \begin{bmatrix}
1 & 1 & 1 & 1\\
8 & 9 & 15 & 2\\
13 & 13 & 4 & 4\\
2 & 15 & 9 & 8\\
\end{bmatrix} \bmod{17}$ 

$ $ 

$\hathat{W}^* = \begin{bmatrix}
1 & \omega^{J(0)} & (\omega^{J(0)})^2 & (\omega^{J(0)})^3\\
1 & \omega^{J(1)} & (\omega^{J(1)})^2 & (\omega^{J(1)})^3\\
1 & \omega^{J_*(0)} & (\omega^{J_*(0)})^2 & (\omega^{J_*(0)})^3\\
1 & \omega^{J_*(1)} & (\omega^{J_*(1)})^2 & (\omega^{J_*(1)})^3\\
\end{bmatrix} \equiv \begin{bmatrix}
1 & 9 & 13 & 15\\
1 & 8 & 13 & 2\\
1 & 2 & 4 & 8\\
1 & 15 & 4 & 9\\
\end{bmatrix}  \bmod{17}$

$ $

注意，定理~\ref*{subsec:vandermonde-euler-integer-ring}（位于 \autoref{subsec:vandermonde-euler-integer-ring}）可以验证如下：

$\hathat{W}^* \cdot \hathat{W} = \begin{bmatrix}
1 & 9 & 13 & 15\\
1 & 8 & 13 & 2\\
1 & 2 & 4 & 8\\
1 & 15 & 4 & 9\\
\end{bmatrix} \cdot \begin{bmatrix}
1 & 1 & 1 & 1\\
8 & 9 & 15 & 2\\
13 & 13 & 4 & 4\\
2 & 15 & 9 & 8\\
\end{bmatrix} = \begin{bmatrix}
0 & 0 & 0 & 4\\
0 & 0 & 4 & 0\\
0 & 4 & 0 & 0\\
4 & 0 & 0 & 0\\
\end{bmatrix} = n I_n^{R} \pmod{17}$

$ $

现在，假设要对 $\mathbb{Z}_{17}$ 中的以下两个输入向量（即输入向量槽位）进行编码：

$\vec{v}_1 = (10, 3, 5, 13)$

$\vec{v}_2 = (2, 4, 3, 6)$

$\vec{v}_1 + \vec{v}_2 = (10, 3, 5, 13) + (2, 4, 3, 6) \equiv (12, 7, 8, 2) \bmod 17$

$ $

这两个向量的编码如下：

$ $

$\vec{m}_1 = n^{-1} \hathat{W} \cdot I_n^R \cdot \vec{v} = 13 \cdot \begin{bmatrix}
1 & 1 & 1 & 1\\
8 & 9 & 15 & 2\\
13 & 13 & 4 & 4\\
2 & 15 & 9 & 8\\
\end{bmatrix} \cdot \begin{bmatrix}
0 & 0 & 0 & 1\\
0 & 0 & 1 & 0\\
0 & 1 & 0 & 0\\
1 & 0 & 0 & 0\\
\end{bmatrix} \cdot \begin{bmatrix}
10\\
3\\
5\\
13\\
\end{bmatrix} \equiv (12, 11, 12, 1)  \bmod 17$

$ $

$\vec{m}_2 = n^{-1} \hathat{W} \cdot I_n^R \cdot \vec{v} = 13 \cdot \begin{bmatrix}
1 & 1 & 1 & 1\\
8 & 9 & 15 & 2\\
13 & 13 & 4 & 4\\
2 & 15 & 9 & 8\\
\end{bmatrix} \cdot \begin{bmatrix}
0 & 0 & 0 & 1\\
0 & 0 & 1 & 0\\
0 & 1 & 0 & 0\\
1 & 0 & 0 & 0\\
\end{bmatrix} \cdot \begin{bmatrix}
2\\
4\\
3\\
6\\
\end{bmatrix} \equiv (8, 5, 14, 6) \bmod 17$

$ $

$ $

$\Delta M_1(X) = 15\cdot(12 + 11X + 12X^2 + 1X^3) = 180 + 165X + 180X^2 + 15X^3 \pmod{q}$ \textcolor{red}{ $\rhd$ 其中 $q = 256$}

$\Delta M_2(X) = 15\cdot(8 + 5X + 14X^2 + 6X^3) = 120 + 75X + 210X^2 + 90X^3 \pmod{q}$

$\Delta M_{1+2}(X) = \Delta M_1(X) + \Delta M_2(X) = \Delta (M_1(X) + M_2(X)) = 44 + 240X + 134X^2 + 105X^3 \pmod{q}$

$M_{1+2}(X) =\left\lceil\dfrac{17}{256}(44 + 240X + 134X^2 + 105X^3)\right\rfloor \bmod 17$

$= \left\lceil2.922 + 15.938X + 8.898X^2 + 6.973X^3\right\rfloor \bmod 17$

$3 + 16X + 9X^2 + 7X^3 \bmod 17 $

$ $


这一多项式与 $\vec{m}_{1+2}$ 的值一致，如下所示：

$\vec{m}_1 + \vec{m}_2 = (12, 11, 12, 1) + (8, 5, 14, 6) = (20, 16, 26, 7) \equiv (3, 16, 9, 7) \bmod 17$

$ $

最后，按如下方式解码 $\vec{m}_{1+2}$：



$\vec{v}_{1+2} = \hathat{W}^* \cdot \vec{m}_{1+2} = \begin{bmatrix}
1 & 9 & 13 & 15\\
1 & 8 & 13 & 2\\
1 & 2 & 4 & 8\\
1 & 15 & 4 & 9\\
\end{bmatrix} \cdot \begin{bmatrix}3\\16\\9\\7\end{bmatrix} = (12, 7, 8, 2) \bmod{17}$

$ $

这一结果与预期的向量和 $\vec{v}_1 + \vec{v}_2$ 一致。




\subsubsection{旋转示例}
\label{subsubsec:bfv-rotation-ex}

同样，在上述示例中，计算 $\bmod t$ 和 $\bmod q$ 时使用无符号模表示（例如 $[0, t-1], [0, q]$）。不过，即使使用居中模表示（例如 $[-\dfrac{t}{2}, \dfrac{t}{2})$、$[-\dfrac{q}{2}, \dfrac{q}{2})$），结果的正确性仍然成立。

假设采用如下设置：

$\mu = 16, n = \dfrac{\mu}{2} = 8, \text{ } t = 17, \text{ } q = 2^8 = 256, \text{ } n^{-1} = 15, \text{ } \Delta = 15$

$ $

$\mathcal{R}_{\langle 8, 17 \rangle} = \mathbb{Z}_{17}[X] / (X^8 + 1)$

$ $

$X^8 + 1 \pmod{17}$ 的根为 $X = \{3, 5, 6, 7, 10, 11, 12, 14\}$，验证如下：

$3^8 \equiv 5^8 \equiv 6^8 \equiv 7^8 \equiv 10^8 \equiv 11^8 \equiv 12^8 \equiv 14^8 \equiv 16 \equiv -1 \bmod{17}$


$ $

在 $X^8 + 1$ 的根 $\{3, 5, 6, 7, 10, 11, 12, 14\}$ 中，选取 $\omega = 3$ 作为基础根，按如下方式构造编码矩阵 $\hathat{W}$ 和解码矩阵 $\hathat{W}^*$：

$\hathat{W} = \begin{bmatrix}
1 & 1 & 1 & 1 & 1 & 1 & 1 & 1\\
\omega^{J(3)} & \omega^{J(2)} & \omega^{J(1)} & \omega^{J(0)} & \omega^{J_*(3)} & \omega^{J_*(2)} & \omega^{J_*(1)} & \omega^{J_*(0)}\\
(\omega^{J(3)})^2 & (\omega^{J(2)})^2 & (\omega^{J(1)})^2 & (\omega^{J(0)})^2 & (\omega^{J_*(3)})^2 & (\omega^{J_*(2)})^2 & (\omega^{J_*(1)})^2 & (\omega^{J_*(0)})^2\\
(\omega^{J(3)})^3 & (\omega^{J(2)})^3 & (\omega^{J(1)})^3 & (\omega^{J(0)})^3 & (\omega^{J_*(3)})^3 & (\omega^{J_*(2)})^3 & (\omega^{J_*(1)})^3 & (\omega^{J_*(0)})^3\\
(\omega^{J(3)})^4 & (\omega^{J(2)})^4 & (\omega^{J(1)})^4 & (\omega^{J(0)})^4 & (\omega^{J_*(3)})^4 & (\omega^{J_*(2)})^4 & (\omega^{J_*(1)})^4 & (\omega^{J_*(0)})^4\\
(\omega^{J(3)})^5 & (\omega^{J(2)})^5 & (\omega^{J(1)})^5 & (\omega^{J(0)})^5 & (\omega^{J_*(3)})^5 & (\omega^{J_*(2)})^5 & (\omega^{J_*(1)})^5 & (\omega^{J_*(0)})^5\\
(\omega^{J(3)})^6 & (\omega^{J(2)})^6 & (\omega^{J(1)})^6 & (\omega^{J(0)})^6 & (\omega^{J_*(3)})^6 & (\omega^{J_*(2)})^6 & (\omega^{J_*(1)})^6 & (\omega^{J_*(0)})^6\\
(\omega^{J(3)})^7 & (\omega^{J(2)})^7 & (\omega^{J(1)})^7 & (\omega^{J(0)})^7 & (\omega^{J_*(3)})^7 & (\omega^{J_*(2)})^7 & (\omega^{J_*(1)})^7 & (\omega^{J_*(0)})^7\\
\end{bmatrix}$ 

$ \equiv \begin{bmatrix}
1 & 1 & 1 & 1 & 1 & 1 & 1 & 1\\
12&14&5&3&10&11&7&6\\
8&9&8&9&15&2&15&2\\
11&7&6&10&14&5&3&12\\
13&13&13&13&4&4&4&4\\
3&12&14&5&6&10&11&7\\
2&15&2&15&9&8&9&8\\
7&6&10&11&5&3&12&14\\
\end{bmatrix} \bmod{17}$ 

$ $ 

$\hathat{W}^* = \begin{bmatrix}
1 & \omega^{J(0)} & (\omega^{J(0)})^2 & (\omega^{J(0)})^3 & (\omega^{J(0)})^4 & (\omega^{J(0)})^5 & (\omega^{J(0)})^6 & (\omega^{J(0)})^7\\
1 & \omega^{J(1)} & (\omega^{J(1)})^2 & (\omega^{J(1)})^3 & (\omega^{J(1)})^4& (\omega^{J(1)})^5& (\omega^{J(1)})^6& (\omega^{J(1)})^7\\
1 & \omega^{J(2)} & (\omega^{J(2)})^2 & (\omega^{J(2)})^3 & (\omega^{J(2)})^4 & (\omega^{J(2)})^5 & (\omega^{J(2)})^6 & (\omega^{J(2)})^7\\
1 & \omega^{J(3)} & (\omega^{J(3)})^2 & (\omega^{J(3)})^3 & (\omega^{J(3)})^4& (\omega^{J(3)})^5& (\omega^{J(3)})^6& (\omega^{J(3)})^7\\
1 & \omega^{J_*(0)} & (\omega^{J_*(0)})^2 & (\omega^{J_*(0)})^3 & (\omega^{J_*(0)})^4 & (\omega^{J_*(0)})^5 & (\omega^{J_*(0)})^6 & (\omega^{J_*(0)})^7\\
1 & \omega^{J_*(1)} & (\omega^{J_*(1)})^2 & (\omega^{J_*(1)})^3 & (\omega^{J_*(1)})^4& (\omega^{J_*(1)})^5& (\omega^{J_*(1)})^6& (\omega^{J_*(1)})^7\\
1 & \omega^{J_*(2)} & (\omega^{J_*(2)})^2 & (\omega^{J_*(2)})^3 & (\omega^{J_*(2)})^4 & (\omega^{J_*(2)})^5 & (\omega^{J_*(2)})^6 & (\omega^{J_*(2)})^7\\
1 & \omega^{J_*(3)} & (\omega^{J_*(3)})^2 & (\omega^{J_*(3)})^3 & (\omega^{J_*(3)})^4& (\omega^{J_*(3)})^5& (\omega^{J_*(3)})^6& (\omega^{J_*(3)})^7\\
\end{bmatrix}$

$ \equiv \begin{bmatrix}
1&3&9&10&13&5&15&11\\
1&5&8&6&13&14&2&10\\
1&14&9&7&13&12&15&6\\
1&12&8&11&13&3&2&7\\
1&6&2&12&4&7&8&14\\
1&7&15&3&4&11&9&12\\
1&11&2&5&4&10&8&3\\
1&10&15&14&4&6&9&5\\
\end{bmatrix}  \bmod{17}$

$ $

$ $

注意，定理~\ref*{subsec:vandermonde-euler-integer-ring}（位于 \autoref{subsec:vandermonde-euler-integer-ring}）可以验证如下：

$\hathat{W}^* \cdot \hathat{W} = \begin{bmatrix}
1&3&9&10&13&5&15&11\\
1&5&8&6&13&14&2&10\\
1&14&9&7&13&12&15&6\\
1&12&8&11&13&3&2&7\\
1&6&2&12&4&7&8&14\\
1&7&15&3&4&11&9&12\\
1&11&2&5&4&10&8&3\\
1&10&15&14&4&6&9&5\\
\end{bmatrix} \cdot \begin{bmatrix}
1 & 1 & 1 & 1 & 1 & 1 & 1 & 1\\
12&14&5&3&10&11&7&6\\
8&9&8&9&15&2&15&2\\
11&7&6&10&14&5&3&12\\
13&13&13&13&4&4&4&4\\
3&12&14&5&6&10&11&7\\
2&15&2&15&9&8&9&8\\
7&6&10&11&5&3&12&14\\
\end{bmatrix}$

$ = \begin{bmatrix}
0 & 0 & 0 & 0 & 0 & 0 & 0 & 8\\
0 & 0 & 0 & 0 & 0 & 0 & 8 & 0\\
0 & 0 & 0 & 0 & 0 & 8 & 0 & 0\\
0 & 0 & 0 & 0 & 8 & 0 & 0 & 0\\
0 & 0 & 0 & 8 & 0 & 0 & 0 & 0\\
0 & 0 & 8 & 0 & 0 & 0 & 0 & 0\\
0 & 8 & 0 & 0 & 0 & 0 & 0 & 0\\
8 & 0 & 0 & 0 & 0 & 0 & 0 & 0\\
\end{bmatrix} = n I_n^{R} \pmod{17}$

$ $

现在，假设要对 $\mathbb{Z}_{17}$ 中的以下输入向量（即输入向量槽位）进行编码：

$\vec{v} = (1, 2, 3, 4, 5, 6, 7, 8)$

将这一向量向左旋转 3 个位置（即前半部分槽位与后半部分槽位分别在各自的组内循环回绕），得到一个新向量：

$\vec{v}_r = (4, 1, 2, 3, 8, 5, 6, 7) \bmod 17$

$ $

$\vec{v}$ 的编码如下：

$ $

$\vec{m} = n^{-1} \hathat{W} \cdot I_n^R \cdot \vec{v} $

$= 15 \cdot  \begin{bmatrix}
1 & 1 & 1 & 1 & 1 & 1 & 1 & 1\\
12&14&5&3&10&11&7&6\\
8&9&8&9&15&2&15&2\\
11&7&6&10&14&5&3&12\\
13&13&13&13&4&4&4&4\\
3&12&14&5&6&10&11&7\\
2&15&2&15&9&8&9&8\\
7&6&10&11&5&3&12&14\\
\end{bmatrix} \cdot \begin{bmatrix}
0 & 0 & 0 & 0 & 0 & 0 & 0 & 1\\
0 & 0 & 0 & 0 & 0 & 0 & 1 & 0\\
0 & 0 & 0 & 0 & 0 & 1 & 0 & 0\\
0 & 0 & 0 & 0 & 1 & 0 & 0 & 0\\
0 & 0 & 0 & 1 & 0 & 0 & 0 & 0\\
0 & 0 & 1 & 0 & 0 & 0 & 0 & 0\\
0 & 1 & 0 & 0 & 0 & 0 & 0 & 0\\
1 & 0 & 0 & 0 & 0 & 0 & 0 & 0\\
\end{bmatrix} \cdot \begin{bmatrix}
1\\
2\\
3\\
4\\
5\\
6\\
7\\
8\\
\end{bmatrix}$

$ \equiv (13, 16, 10, 5, 9, 12, 7, 1)  \bmod 17$


$ $

$\Delta M(X) = 15\cdot(13 + 16X + 10X^2 + 5X^3 + 9X^4 + 12X^5 + 7X^6 + X^7)$


$ = 195 + 240X + 150X^2 + 75X^3 + 135X^4 + 180X^5 + 105X^6 + 15X^7 \pmod{q}$ \textcolor{red}{ $\rhd$ 其中 $q = 256$}

$ $

这一多项式与 $\Delta M(X^{J(3)})$ 的值一致，如下所示：

$\Delta M(X^{J(3)}) = \Delta M(X^{13}) $

$= 195 + 180X - 150X^2 - 15X^3 + 135X^4 - 240X^5 - 105X^6 + 75X^7 \pmod{q}$

$ $



%Note that the coefficients of the scaled polynomials $\Delta M(X^{J(3)})$ are still within the range of the ciphertext domain $q=64$ (which must hold throughout all homomorphic computations to preserve correctness). 

现在，按如下方式解码 $M(X^{J(3)})$：

$\vec{m}_{J(3)} = \dfrac{\Delta \vec{m}_{J(3)}}{\Delta} = \dfrac{(195, 180, -150, -15, 135, -240, -105, 75)}{15} = (13, 12, -10, -1, 9, -16, -7, 5) \pmod{17}$

$ $


$\vec{v}_{J(3)} = \hathat{W}^* \cdot \vec{m}_{J(3)} = \begin{bmatrix}
1&3&9&10&13&5&15&11\\
1&5&8&6&13&14&2&10\\
1&14&9&7&13&12&15&6\\
1&12&8&11&13&3&2&7\\
1&6&2&12&4&7&8&14\\
1&7&15&3&4&11&9&12\\
1&11&2&5&4&10&8&3\\
1&10&15&14&4&6&9&5\\
\end{bmatrix} \cdot \begin{bmatrix}
13\\
12\\
-10\\
-1\\
9\\
-16\\
-7\\
5\end{bmatrix}$

$= (4, 1, 2, 3, 8, 5, 6, 7) \pmod{17}$

$= \vec{v}_r$

$ $

解码得到的 $\vec{v}_{J(3)}$ 与预期的旋转后输入向量 $\vec{v}_r$ 一致。


$ $

实际中，对加密分量应用自同构，得到
\[
\textsf{ct}^{\langle h=3\rangle}=(A(X^{J(3)}),B(X^{J(3)}))
\]
。其解密秘密密钥为 $S(X^{J(3)})$。在执行通常的解密之前，使用相应旋转密钥将秘密密钥切换回 $S(X)$；只要变换后的误差与密钥切换误差仍满足界限，解码就得到 $\vec v_r$。

$ $




\para{源代码：} 运行 \href{https://github.com/wang2000star/fhe-textbook-cn/blob/chinese-edition/source%20code/bfv.py}{\underline{这个 Python 脚本}}（BFV 加法）和 \href{https://github.com/wang2000star/fhe-textbook-cn/blob/chinese-edition/source%20code/bfv_rotation_only.py}{\underline{这个脚本}}（仅 BFV 旋转），可以执行 BFV 批量编码和同态输入向量旋转的示例。



\subsection{应用：矩阵乘法}
\label{subsec:bfv-matrix-multiplication}

BFV 可以先逐槽相乘，再通过旋转和加法求和，计算打包内积。在槽位数为二次幂的单个循环块内，倍增求和只需对数数量的旋转；若要合并两行，还需要行交换。下述矩阵与向量乘法的对角方法假设 $m\leq n$，且 $\rho$ 表示对全部 $n$ 个有效槽位的逻辑循环旋转。前节只实现两个半行各自旋转，因此仅凭该操作不能实现完整的 $\rho$。可以将逻辑向量放在更大环的一个循环行中，也可组合行内旋转、行交换自同构 $X\mapsto X^{-1}$ 和明文掩码。给定矩阵 $A$ 和向量 $\vec x$（\autoref{subsec:matrix-arithmetic} 中的定义~\ref*{subsec:matrix-arithmetic}）：

$A = \begin{bmatrix}
 a_{\langle 0, 0\rangle} & a_{\langle 0, 1\rangle} & a_{\langle 0, 2\rangle} & \cdots & a_{\langle 0, n-1\rangle}\\
 a_{\langle 1, 0\rangle} & a_{\langle 1, 1\rangle} & a_{\langle 1, 2\rangle} & \cdots & a_{\langle 1, n-1\rangle} \\
 a_{\langle 2, 0\rangle} & a_{\langle 2, 1\rangle} & a_{\langle 2, 2\rangle} & \cdots & a_{\langle 2, n-1\rangle} \\
\vdots & \vdots & \vdots & \ddots & \vdots \\
 a_{\langle m-1, 0\rangle} & a_{\langle m-1, 1\rangle} & a_{\langle m-1, 2\rangle} & \cdots & a_{\langle m-1, n-1\rangle} \\
\end{bmatrix} = \begin{bmatrix} 
\vec{a}_{\langle 0, * \rangle} \\ 
\vec{a}_{\langle 1, * \rangle} \\ 
\vec{a}_{\langle 2, * \rangle} \\ 
\vdots\\
\vec{a}_{\langle m-1, * \rangle} \\ 
\end{bmatrix}, \text{ } \vec{x} = (x_0, x_1, \cdots, x_{n-1})$

$ $

$A \cdot \vec{x}$ 的结果是一个 $m$ 维向量，计算方式如下：

$A \cdot \vec{x} = \Big(\vec{a}_{\langle 0, * \rangle} \cdot \vec{x}, \text{ } \vec{a}_{\langle 1, * \rangle} \cdot \vec{x}, \text{ }\cdots,\text{ } \vec{a}_{\langle m-1, * \rangle} \cdot \vec{x} \Big) = \left(\sum\limits_{i=0}^{n-1} a_{0, i} \cdot  x_i, \text{ } \sum\limits_{i=0}^{n-1}  a_{1, i} \cdot x_i, \cdots, \text{ } \sum\limits_{i=0}^{n-1} a_{m-1, i} \cdot x_i \right)$

$ $

将 $\vec{v}$ 向左旋转 $h$ 个位置后的结果定义为 $\rho(\vec{v}, h)$。回顾 Hadamard 乘积（\autoref{subsec:vector-arithmetic} 中的定义~\ref*{subsec:vector-arithmetic}），它定义为两个向量逐槽位相乘：

$\vec{a} \odot \vec{b} = (a_0 b_0, \text{ } a_1 b_1, \text{ } \cdots, \text{ } a_{n-1} b_{n-1})$

$ $

从矩阵 $A$ 中提取 $n$ 个不同的对角向量 $\vec{u}_i$，定义如下：

$\vec{u}_i = \{a_{\langle j \bmod m, \text{ } (i + j)\bmod n \rangle}\}_{j = 0}^{n-1}$

$ $

那么，原始矩阵与向量相乘的公式可以等价地构造如下：

$A\cdot\vec x=\left(\sum_{i=0}^{n-1}\vec u_i\odot\rho(\vec x,i)\right)_{0:m}$

，其计算结果等价于 $A \cdot \vec{x}$。
上述公式可以用于同态计算，因为 BFV 支持通过输入向量的多项式编码形式之间的密文与明文相乘，计算输入向量的 Hadamard 乘积（\autoref{subsec:bfv-mult-plain}）；BFV 也支持将 $\rho(\vec{v}, h)$ 作为同态输入向量槽位旋转来计算（\autoref{subsec:bfv-rotation}）。在同态计算上述公式后，只需使用所得输入向量的前 $m$ 个槽位（共 $n$ 个槽位）来存储 $A \cdot \vec{x}$ 的结果。



\subsection{噪声自举}
\label{subsec:bfv-bootstrapping}

\noindent \textbf{- 参考资料 1：}
\href{https://eprint.iacr.org/2022/1363.pdf}{重新审视 BGV 与 BFV 的自举}~\cite{cryptoeprint:2022/1363}


\noindent \textbf{- 参考资料 2：}
\href{https://eprint.iacr.org/2014/873.pdf}{HELib 的自举}~\cite{10.1007/s00145-020-09368-7}

\noindent \textbf{- 参考资料 3：}
\href{https://arxiv.org/pdf/1906.02867}{关于自举中低位提取多项式的说明}~\cite{huo2019notelowerdigitsextraction}

\noindent \textbf{- 参考资料 4：}
\href{https://eprint.iacr.org/2024/1587.pdf}{采用分圆素数模数的全同态加密}~\cite{cryptoeprint:2024/1587}

%\noindent \textbf{- Reference 4:} 
%\href{https://s-space.snu.ac.kr/bitstream/10371/152893/1/000000155273.pdf}{Bootstrapping Methods for Homomorphic Encryption}

%\noindent \textbf{- Reference 5:} 
%\href{https://www.esat.kuleuven.be/cosic/publications/thesis-418.pdf}{Bootstrapping Algorithms for BGV and FV}

$ $

在 BFV 中，连续进行密文与密文相乘会使噪声以乘法方式增长；一旦噪声溢出到消息位，消息就会遭到破坏。自举是重置已增长噪声的过程。

\subsubsection{总体思路}
\label{subsubsec:bfv-bootstrapping-high-level}

令 $t=p^r$、$\varepsilon>r$ 且 $P=p^\varepsilon$。下面分解解密函数，给出一般系数与槽位变换构造，并在 \autoref{subsubsec:bfv-bootstrapping-summary} 给出显式的逐门正确性充分条件。该条件中的同态多项式求值采用标准 FV：先计算整数张量乘积，再按当前明文模数 $T$ 以 $T/q$ 缩放；当 $T\nmid q$ 时，不能改用前述 $1/\Delta$ 教学变体。这是条件性正确性结论，并不是安全或优化参数的选择。使用秘密密钥在自身之下的加密，还需所引构造的求值密钥安全及循环安全性假设~\cite{cryptoeprint:2022/1363}。

对于当前明文模数 $T$，按精确的有理尺度 $\mu_T=q/T$ 衡量误差会更方便。令 $r_T=q-T\lfloor q/T\rfloor$，同一个密文满足
\[
\phi_S(C)\equiv\mu_TM_T+\mathcal E_T\pmod q,
\qquad \mathcal E_T=E_T-\frac{r_T}{T}M_T.
\]
这里 $M_T$ 的各系数取中心代表。若 $\|\mathcal E_T\|_\infty<q/(2T)$，解码就正确。改变误差坐标既不改变密文，也不把方案变成近似明文方案；这个条件正是 \autoref{subsec:bfv-enc-dec} 的解码条件。

对于相位满足 $AS+B=\Delta M+E\pmod q$ 的密文 $(A,B)$，目标明文恒等式为
\[
A'=\Bigl[\Bigl\lceil(P/q)A\Bigr\rfloor\Bigr]_P,\quad
B'=\Bigl[\Bigl\lceil(P/q)B\Bigr\rfloor\Bigr]_P,\quad
Z=[A'S+B']_P,
\qquad M=\Bigl[\Bigl\lceil Z/p^{\varepsilon-r}\Bigr\rfloor\Bigr]_{p^r}.
\]
\begin{enumerate}
\item \textbf{模数切换：}将输入密文模数从 $q$ 改为 $P$。这一步本身不改变原明文模数 $p^r$。参数适当时，切换后相位为 $p^{\varepsilon-r}M+E'\pmod P$，其中 $\|E'\|_\infty<p^{\varepsilon-r}/2$。
\item \textbf{同态解密：}使用在较大密文模数 $q$ 下加密 $S$ 的求值密文，此时明文模数为 $P$，尺度为 $\Delta'=\lfloor q/P\rfloor$。计算 $A'\cdot\textsf{RLWE}_{S,\sigma}(\Delta'S)+(0,\Delta'B')$。结果加密相位 $Z\pmod P$，并含有额外求值噪声。公开常数 $B'$ 必须以尺度 $\Delta'$ 编码。
\item \textbf{\textsf{CoeffToSlot}：}应用同态线性变换，将 $Z$ 的系数放入槽位。在满足 \autoref{subsubsec:bfv-rotation-summary} 的完全标量分裂条件时，对槽位向量应用的变换为 $n^{-1}\hathat W I_n^R$。
\item \textbf{数字提取：}在每个槽位中，通过低位提取多项式及精确除法步骤，计算 $[\lceil z_i/p^{\varepsilon-r}\rfloor]_{p^r}$。当 $p$ 为奇数时使用平衡数字；当 $p=2$ 时，删除无符号数字实现的是向下取整，因此删除之前需要加半单位偏移量。下文将说明这些数字恒等式；它们用于完整自举时，还须满足 \autoref{subsubsec:bfv-bootstrapping-coefftoslot} 的 CRT 适用域限制。
\item \textbf{\textsf{SlotToCoeff}：}应用逆线性变换 $\hathat W^*$，将恢复出的系数放回多项式。结果加密 $M$，并带有求值产生的新误差 $E^{\langle b\rangle}$。其界取决于参数和电路，并不自动小于任意输入误差 $E$，也不能直接与尺度不同的 $E'$ 比较。
\end{enumerate}

旧的 $E'$ 在求值舍入函数时被消除，而同态求值自身产生的误差仍然保留。可用的自举必须为后续操作留下足够的输出噪声预算。只有选择适当参数并控制总失败概率，重复自举才能支持任意预先给定的有限电路深度。

\subsubsection{模数切换}
\label{subsubsec:bfv-bootstrapping-modulus-switch}

令 $P=p^\varepsilon$，在约减前选用相容的整数提升，定义各分量舍入多项式 $R_A=A'-(P/q)A$ 和 $R_B=B'-(P/q)B$。其系数绝对值至多为 $1/2$。切换后的精确相位为
\[
A'S+B'\equiv p^{\varepsilon-r}M+E'\pmod P,
\quad E'=\frac PqE+\left(\frac Pq\Delta-p^{\varepsilon-r}\right)M+R_AS+R_B.
\]
因此，目标舍入的一个充分条件是 $\|E'\|_\infty<p^{\varepsilon-r}/2$。各分量舍入项不能用仅对旧误差舍入来替代。除缩放后的旧误差和尺度失配项外，该项的确定性界为 $(\|S\|_1+1)/2$。

\subsubsection{同态解密}
\label{subsubsec:bfv-bootstrapping-homomorphic-decryption}

写成 $Z=[A'S+B']_P=[p^{\varepsilon-r}M+E']_P$。取 $\Delta'=\lfloor q/P\rfloor$，令求值密文 $(A_s,B_s)$ 满足
\[
A_sS+B_s=\Delta'S+E_s\pmod q.
\]
同态内积得到密文
\[
C=(A'A_s,\ A'B_s+\Delta'B'),
\qquad \phi_S(C)=\Delta'(A'S+B')+A'E_s\pmod q.
\]
这明确区分了加密明文 $\Delta'S$ 与加密秘密密钥 $S$。若 $A'S+B'=Z+PK$，则
\[
\phi_S(C)\equiv\Delta'Z+A'E_s-r_PK\pmod q,
\qquad r_P=q-P\Delta'.
\]
因此，当 $P\nmid q$ 时，明文模 $P$ 约减表现为额外误差 $-r_PK$，并不需要单独执行密文侧的取模门。另有 $\|A'E_s\|_\infty\leq\|A'\|_1\|E_s\|_\infty$。仅有宽泛的不等式 $q\gg P$，并不能保证该误差落在量级为 $q/P$ 的解码预算内。

目标系数操作是 $[\lceil z_i/p^{\varepsilon-r}\rfloor]_{p^r}$。它通过一系列多项式求值和精确除法实现，而不是由一个 $p$ 次多项式完成整个舍入函数。若要逐系数求值，先执行 \textsf{CoeffToSlot}，再执行 \autoref{subsubsec:bfv-bootstrapping-digit-extraction} 的数字提取过程，最后执行 \textsf{SlotToCoeff}。同态求值误差与加密整数 $z_i$ 内部的旧误差 $E'$ 是不同量。

\subsubsection{\textsf{CoeffToSlot} 与 \textsf{SlotToCoeff}}
\label{subsubsec:bfv-bootstrapping-coefftoslot}

下列 $n$ 个标量槽位的矩阵假设 $p\equiv1\pmod{2n}$，并将一个本原单位根 Hensel 提升到 $\mathbb Z_{p^\varepsilon}$，随后在每个明文模数下相容地约减。一般地，当 $\gcd(p,2n)=1$ 且 $d=\operatorname{ord}_{2n}(p)$ 时，有 $n/d$ 个扩域环槽位。下述迹对偶基构造将全部系数打包进 $d$ 个密文的标量槽位，因此不必完全分裂为标量槽位。这里的矩阵乘法表示由自同构、密钥切换和明文乘法组成的同态线性电路。这条 CRT 路线在本书二次幂环中仍要求 $p$ 为奇数：当 $p=2$ 时，$X^n+1=(X+1)^n\pmod2$ 含重因子。二进制 CRT 自举必须改用与 $2$ 互素的分圆指标，或采用另有独立论证的构造。

\textbf{一般可实现的线性变换。}为包含替代的二进制环，暂记 $R_j=(\mathbb Z/p^j\mathbb Z)[X]/\Phi_{m_c}(X)$，其中 $p\nmid m_c$、$n=\varphi(m_c)$、$d=\operatorname{ord}_{m_c}(p)$ 且 $\ell=n/d$。本章原来的环取 $m_c=2n$。选择模 $p^\varepsilon$ 下相容的 $\Phi_{m_c}$ 首一因子，再约减到模 $p^j$；CRT 给出 $\ell$ 个次数为 $d$ 的非分歧扩域环槽位。

对 $a\in(\mathbb Z/m_c\mathbb Z)^\times$，令 $\tau_a(X)=X^a$。由于模 $p$ 的代数是有限域的直积，迹配对非退化，其 Gram 行列式在每个模 $p^j$ 下都是单位。因此存在迹对偶基 $b_0,\ldots,b_{n-1}$，满足
\[
\operatorname{Tr}(b_uX^v)=\delta_{uv},\qquad
\operatorname{Tr}(W)=\sum_a\tau_a(W),\qquad
z_u=\operatorname{Tr}(b_uZ).
\]
于是，每个基环线性映射 $\mathcal T:R_j\to R_j$ 都有显式实现
\[
\mathcal T(Z)=\sum_a\kappa_a\tau_a(Z),\qquad
\kappa_a=\sum_{u=0}^{n-1}\mathcal T(X^u)\tau_a(b_u).
\]
代入迹的展开式，所得结果就是 $\sum_u\mathcal T(X^u)z_u$。每一项只需一个公开多项式乘数，以及自同构后的密钥切换，不涉及解密，也不除以扩域次数。特别地，当 $p\mid d$ 时，不能把槽位迹除以 $d$。提升的 Frobenius 是在生成元上按 $X\mapsto X^p$ 定义的自同构；当 $j>1$ 时，它不是对模 $p^j$ 下任意元素直接应用 $W\mapsto W^p$。

对 $0\leq u<d$，通过 CRT 剩余定义 $\mathcal T_u$：槽位 $h$ 放入标量系数 $z_{u\ell+h}$，其余扩域坐标全为零。逆映射 $\mathcal U_u$ 读取各槽位的标量坐标，将其写入系数位置 $u\ell+h$，其他位置补零。CRT 矩阵求逆即可构造这些映射，迹对偶公式则将正逆映射都转化为同态操作。它们满足
\[
\sum_{u=0}^{d-1}\mathcal U_u\bigl(\mathcal T_u(Z)\bigr)=Z.
\]
打包映射在模 $p^\varepsilon$ 下使用，相容的解包映射在模 $p^r$ 下使用。数字多项式作用于标量项，而不是任意扩域元素。若其分子的每个槽位都能被 $p$ 整除，CRT 就推出整个多项式属于 $pR_j$，因此后续逐系数精确除法有良好定义。这给出了下述变换的一般但未经优化的构造~\cite{cryptoeprint:2022/1363}。

熟悉的完全分裂矩阵是 $d=1$ 的特殊情形：

\textsf{CoeffToSlot} 步骤的目标是以同态方式将多项式 $Z$ 的系数 $z_i$ 移入输入向量槽位。


在小结~\ref*{subsubsec:bfv-rotation-summary}（位于 \autoref{subsubsec:bfv-rotation-summary}）中，我们了解到，将输入槽位向量 $\vec{v}$ 转换为多项式系数向量 $\vec{m}$ 的编码公式为 $\vec{m} = n^{-1}\cdot\hathat{W} \cdot I_n^R \cdot \vec{v}$，其中，$\hathat{W}$ 是按如下方式构造的 $n$ 维向量空间的一组基：

$\hathat{W} = \begin{bmatrix}
1 & 1 & \cdots & 1 & 1 & 1 & \cdots & 1\\
(\omega^{J(\frac{n}{2} - 1)}) & (\omega^{J(\frac{n}{2} - 2)}) & \cdots & (\omega^{J(0)}) & (\omega^{J_*(\frac{n}{2} - 1)}) & (\omega^{J_*(\frac{n}{2} - 2)}) & \cdots & (\omega^{J_*(0)})\\
(\omega^{J(\frac{n}{2} - 1)})^2 & (\omega^{J(\frac{n}{2} - 2)})^2 & \cdots & (\omega^{J(0)})^2 & (\omega^{J_*(\frac{n}{2} - 1)})^2 & (\omega^{J_*(\frac{n}{2} - 2)})^2 & \cdots & (\omega^{J_*(0)})^2 \\
\vdots & \vdots & \ddots & \vdots & \vdots & \vdots & \ddots & \vdots \\
(\omega^{J(\frac{n}{2} - 1)})^{n-1} & (\omega^{J(\frac{n}{2} - 2)})^{n-1} & \cdots & (\omega^{J(0)})^{n-1} & (\omega^{J_*(\frac{n}{2} - 1)})^{n-1} & (\omega^{J_*(\frac{n}{2} - 2)})^{n-1} & \cdots  & (\omega^{J_*(0)})^{n-1}
\end{bmatrix}$

$ $

\textcolor{red}{ $\rhd$ 其中，旋转辅助函数为 $J(h) = 5^h \bmod 2n$}

$ $

因此，对于输入密文 $\textsf{ct} = \textsf{RLWE}_{S, \sigma}\bm(\Delta' Z\bm) \bmod q$，可以将其输入向量槽位理解为存储了某些值，对它们乘以 $n^{-1}\cdot\hathat{W} \cdot I_n^R$ 后就得到多项式 $Z$ 的系数 $z_i$。这意味着，如果对 $\textsf{RLWE}_{S, \sigma}\bm(\Delta' Z\bm)$ 的输入向量槽位\textit{同态}乘以 $n^{-1}\cdot\hathat{W} \cdot I_n^R$，那么所得密文的 $n$ 个输入向量槽位将包含 $Z$ 的 $n$ 个系数，等价于将 $Z$ 的系数移入输入向量槽位。因此，\textsf{CoeffToSlot} 步骤等价于同态计算 $n^{-1}\cdot\hathat{W} \cdot I_n^R \cdot \textsf{RLWE}_{S, \sigma}\bm(\Delta' Z\bm)$。可以使用 \autoref{subsec:bfv-matrix-multiplication} 中介绍的技术同态计算矩阵与向量相乘。

在 \textsf{CoeffToSlot} 步骤之后，可以通过对去噪多项式进行同态求值，消除每个 $z_i$ 的低 $\varepsilon-r$ 个（$p$ 进制）数位中的噪声；下一小节将介绍这一多项式。

在明文槽位中恢复 $M$ 的系数后，需要将它们移回原来的系数位置。密文仍含有有界的求值误差。这一步称为 \textsf{SlotToCoeff}，它恰好是 \textsf{CoeffToSlot} 的逆过程。由小结~\ref*{subsubsec:bfv-rotation-summary}（位于 \autoref{subsubsec:bfv-rotation-summary}）还可知，$n^{-1}\cdot\hathat{W} \cdot I_n^R$ 的逆矩阵为 $\hathat{W}^*$，其中：

$\hathat{W}^* = \begin{bmatrix}
1 & (\omega^{J(0)}) & (\omega^{J(0)})^2 & \cdots & (\omega^{J(0)})^{n-1}\\
1 & (\omega^{J(1)}) & (\omega^{J(1)})^2 & \cdots & (\omega^{J(1)})^{n-1}\\
1 & (\omega^{J(2)}) & (\omega^{J(2)})^2 & \cdots & (\omega^{J(2)})^{n-1}\\
\vdots & \vdots & \vdots & \ddots & \vdots \\
1 & (\omega^{J(\frac{n}{2}-1)}) & (\omega^{J(\frac{n}{2}-1)})^2 & \cdots & (\omega^{J(\frac{n}{2}-1)})^{n-1}\\
1 & (\omega^{J_*(0)}) & (\omega^{J_*(0)})^2 & \cdots & (\omega^{J_*(0)})^{n-1}\\
1 & (\omega^{J_*(1)}) & (\omega^{J_*(1)})^2 & \cdots & (\omega^{J_*(1)})^{n-1}\\
1 & (\omega^{J_*(2)}) & (\omega^{J_*(2)})^2 & \cdots & (\omega^{J_*(2)})^{n-1}\\
\vdots & \vdots & \vdots & \ddots & \vdots \\
1 & (\omega^{J_*(\frac{n}{2}-1)}) & (\omega^{J_*(\frac{n}{2}-1)})^2 & \cdots & (\omega^{J_*(\frac{n}{2}-1)})^{n-1}\\
\end{bmatrix}$

因此，\textsf{SlotToCoeff} 步骤等价于将去噪多项式求值的输出同态乘以 $\hathat{W}^*$。

下一小节将学习如何基于数位提取多项式 $G_\varepsilon(x)$，设计 BFV 的核心算法——去噪多项式。


\subsubsection{数位提取}
\label{subsubsec:bfv-bootstrapping-digit-extraction}


当 $p$ 为奇数时，使用平衡的 $p$ 进制数字 $d\in\{-(p-1)/2,\ldots,(p-1)/2\}$。将输入系数写成 $z=\sum_{j=0}^{\varepsilon-1}d_jp^j\pmod{p^\varepsilon}$。删除最低 $v=\varepsilon-r$ 个平衡数字，得到解密所需的 $\lceil z/p^v\rfloor\pmod{p^r}$。普通非负数字实现的则是向下取整：例如 $p=3,\varepsilon=3,m=1,e'=-1$ 时 $z=8$，虽然 $|e'|<9/2$，仍有 $\lfloor8/9\rfloor=0\ne1$。当 $p=2$ 时，使用非负数字，并在删除 $v$ 位之前加上 $2^{v-1}$，从而在非中点处实现 $\lceil z/2^v\rfloor=\lfloor(z+2^{v-1})/2^v\rfloor$。这里的 $p=2$ 公式仅是标量舍入恒等式，不能把所引用的 CRT 自举路线推广到本章的二次幂环；该路线要求分圆指标与 $p$ 互素，详见 \autoref{subsubsec:bfv-bootstrapping-coefftoslot}。

对 $j\geq1$，提升多项式 $F_j$ 满足
\[
F_j(d+p^ju)\equiv d\pmod{p^{j+1}}
\]
其中 $d$ 为允许的最低位数字，$u$ 为整数。下标 $j$ 表示算法中已知的精度，不必等于秘密输入中首个非零高位的位置。令 $L_\varepsilon=F_{\varepsilon-1}\circ\cdots\circ F_1$。

\para{第 1 次递归：}$F_1(z)\equiv d_0\pmod{p^2}$。

\para{第 2 次递归：}$F_2\circ F_1(z)\equiv d_0\pmod{p^3}$。

\para{第 3 次递归：}$F_3\circ F_2\circ F_1(z)\equiv d_0\pmod{p^4}$。

\para{第 $\bm{(\varepsilon-1)}$ 次递归：}$L_\varepsilon(z)\equiv d_0\pmod{p^\varepsilon}$。

因此，数字删除操作为
\[
G_\varepsilon(z)=\frac{z-L_\varepsilon(z)}p\pmod{p^{\varepsilon-1}}.
\]
这里对可整除的代表执行精确整数除法，输出模数降低一个 $p$ 的幂次。它不是在 $\mathbb Z_{p^\varepsilon}$ 中乘以 $p$ 的逆元，因为该逆元不存在。即使存在 $p^{-1}\bmod q$，用它乘密文分量也是不同操作，可能放大噪声。因此，$G_\varepsilon$ 表示多项式求值与除法组成的过程，而不是 $\mathbb Z_{p^\varepsilon}$ 上的单个多项式。重复执行
\[
G_{r+1}\circ\cdots\circ G_{\varepsilon-1}\circ G_\varepsilon(z)
\]
即可删除 $\varepsilon-r$ 位。每次调用 $G_j$ 都使用 $F_1,\ldots,F_{j-1}$ 提取自身输入的最低位；即使 $r>1$，也必须从 $F_1$ 开始。

\para{尺度因子的重新解释：}假设同态减法后，明文代表满足 $u=pw\pmod{p^j}$，密文相位为 $\Delta_j u+E_j$，其中 $\Delta_j=\lfloor q/p^j\rfloor$。选择一致且可整除的中心代表后，同一密文可在明文模数 $p^{j-1}$ 下重新解释为
\[
\Delta_jpw+E_j=\Delta_{j-1}w+(E_j-\delta_jw),\qquad
\delta_j=\Delta_{j-1}-p\Delta_j\in\{0,\ldots,p-1\}.
\]
这避免了将密文分量乘以 $p^{-1}\bmod q$。其中 $-\delta_jw$ 是按地板尺度衡量误差时的坐标变化；若按精确尺度衡量，它恰好抵消：

\[
\frac{q}{p^j}(pw)+\mathcal E_j
=\frac{q}{p^{j-1}}w+\mathcal E_j.
\]
因此，精确除法不增加自然误差，而且扩大了解码容差。两种描述等价：若误差界已经采用自然误差坐标，就不能再次累加地板尺度的漂移。使用相容且能被整除的中心系数代表时，这个恒等式精确成立，包括 $p^j\nmid q$ 的情形。

各地板尺度依次为 $\lfloor q/p^\varepsilon\rfloor,\lfloor q/p^{\varepsilon-1}\rfloor,\ldots,\lfloor q/p^r\rfloor$。与密文相加的明文必须采用当前尺度，而密文与明文相乘使用未缩放乘数，如 \autoref{subsec:bfv-mult-plain} 所述。多项式求值会改变自然误差；精确除法保持自然误差，同时改变明文模数。\autoref{subsubsec:bfv-bootstrapping-summary} 的显式递推可逐项检查所有中间解码条件。

剩余的局部代数任务是构造提升多项式。

\para{设计 \boldmath$F_{\varepsilon'}(z_i)$：}下面给出显式控制精度的插值构造。令 $j=\varepsilon'\geq1$，$D$ 为选定的连续 $p$ 个数字代表组成的集合。由费马小定理，对每个 $d\in D$，$(d^p-d)/p$ 都是整数。在 $\mathbb Z_{p^j}$ 上插值得到次数至多为 $p-1$ 的多项式 $H_j$，使得
\[
H_j(d)\equiv(d^p-d)/p\pmod{p^j}\quad(d\in D).
\]
不同数字节点之差在模 $p^j$ 下都是单位元，因此拉格朗日插值成立。选取系数的整数提升，并定义
\[
F_j(X)=X^p-pH_j(X).
\]
\begin{proof}
对于 $z=d+p^ju$，二项式定理给出 $z^p\equiv d^p\pmod{p^{j+1}}$。整数系数多项式的求值还满足 $H_j(z)\equiv H_j(d)\pmod{p^j}$。因此
\[
F_j(z)\equiv d^p-pH_j(d)\equiv d\pmod{p^{j+1}}.
\]
\end{proof}
这证明了各次递归所用的提升性质。其次数至多为 $p$，不是 $p-1$。若仅在 $\mathbb F_p$ 上分别插值各位数字函数，再重复使用任意整数系数提升，则不足以保证更高幂次下的正确性，因为进位会影响高位。系数必须按上述方法在所需精度下构造。这一局部存在性证明本身并没有给出优化同态实现的深度界或噪声界。

\subsubsection{小结}
\label{subsubsec:bfv-bootstrapping-summary}

下面总结推广后的明文变换。方框后的递推明确给出正确性条件，涵盖低模舍入、一般线性变换、全部同态门以及最终解码。安全性与优化仍是独立问题。

\begin{tcolorbox}[title={\textbf{\tboxlabel{\ref*{subsubsec:bfv-bootstrapping-summary}} BFV 自举}}]
令 $t=p^r$、$\varepsilon>r$、$P=p^\varepsilon$、$\Delta=\lfloor q/p^r\rfloor$ 和 $\Delta'=\lfloor q/P\rfloor$。 完整 CRT 流程须满足 \autoref{subsubsec:bfv-bootstrapping-coefftoslot} 中的相容环条件；在本章的二次幂环中，$p$ 为奇数。
\begin{enumerate}
\item \textbf{模数切换：}计算 $A'=[\lceil(P/q)A\rfloor]_P$ 和 $B'=[\lceil(P/q)B\rfloor]_P$。其相位满足 $A'S+B'\equiv p^{\varepsilon-r}M+E'\pmod P$，完整 $E'$ 见 \autoref{subsubsec:bfv-bootstrapping-modulus-switch}。要求 $\|E'\|_\infty<p^{\varepsilon-r}/2$。
\item \textbf{同态解密：}使用模数 $q$ 下 $\Delta'S$ 的加密，计算 $A'\cdot\textsf{RLWE}_{S,\sigma}(\Delta'S)+(0,\Delta'B')$。按 \autoref{subsubsec:bfv-bootstrapping-homomorphic-decryption} 同时跟踪密钥噪声项和明文约减余数项。
\item \textbf{\textsf{CoeffToSlot}：}同态地将 $Z=[p^{\varepsilon-r}M+E']_P$ 的系数放入槽位。在标量分裂情形下，槽位变换为 $n^{-1}\hathat W I_n^R$（\autoref{subsubsec:bfv-rotation-summary}）；对于扩域槽位，则使用 \autoref{subsubsec:bfv-bootstrapping-coefftoslot} 构造的 $d$ 个映射 $\mathcal T_u$。
\item \textbf{数字提取：}对当前精度 $j=\varepsilon,\varepsilon-1,\ldots,r+1$，用 $L_j=F_{j-1}\circ\cdots\circ F_1$ 提取最低位，再实现精确商 $G_j(z)=(z-L_j(z))/p\pmod{p^{j-1}}$。当 $p$ 为奇数时使用平衡数字；当 $p=2$ 时使用上文规定的舍入偏移量。最终槽位为 $m_i\pmod{p^r}$，所以未约减代表可以相差 $k_i'p^r$；当 $r>1$ 时，不能写成 $k_i'p$。
\para{\underline{尺度因子的重新解释}：}每次除法都把尺度从 $\lfloor q/p^j\rfloor$ 重新解释为 $\lfloor q/p^{j-1}\rfloor$，并计入上文推导的漂移 $-\delta_jw$。最终尺度为 $\lfloor q/p^r\rfloor=\Delta$。这里不假设 $p$ 存在模逆元。
\item \textbf{\textsf{SlotToCoeff}：}在标量分裂情形下应用逆线性变换 $\hathat W^*$，扩域槽位情形则对 $\mathcal U_u$ 的输出求和，恢复相位为 $\Delta M+E^{\langle b\rangle}\pmod q$ 的加密。完整输出误差包含全部求值误差、尺度漂移及明文约减项，并须满足 BFV 解码界。
\end{enumerate}
\end{tcolorbox}

\textbf{显式的条件性正确性证书。}采用 \autoref{subsubsec:bfv-bootstrapping-high-level} 定义的自然误差 $\mathcal E_T$。令 $C_R$ 满足 $\|UV\|_\infty\leq C_R\|U\|_\infty\|V\|_\infty$，并令 $h_S$ 是乘以秘密多项式这一线性算子的无穷范数上界。对于负循环环，可取 $C_R=n$ 和 $h_S=\|S\|_1$。对于公开多项式 $C$ 或自同构 $\tau_a$，分别以 $\lambda_C$、$\lambda_a$ 表示相应系数算子的无穷范数。这些都是可计算的矩阵范数，也适用于一般分圆环。

使用具有 $L$ 个数字、数字绝对值不超过 $D_g$ 的精确整数分解，求值密钥各行在输出秘密下的原始相位为 $g_iS_{\rm old}+E_i^{\rm key}$，且 $\|E_i^{\rm key}\|_\infty\leq B_{\rm key}$。一次密钥切换至多给自然误差上界增加
\[
B_{\rm ks}=C_RLD_gB_{\rm key}
\]
这里不需要独立性假设。若输入界为 $B_1,B_2$，可递推下列有效误差界；必要时可以将界向上取整：
\[
\begin{aligned}
B_{\rm add/sub}&=B_1+B_2,\qquad B_{\rm plain\,mult}=\lambda_CB_1,\\
B_{\rm plain\,add}&=B_1+(r_T/T)\|C\|_\infty,\\
B_{\rm auto}&=\lambda_aB_1+B_{\rm ks},\qquad B_{\rm exact\,div}=B_1,\\
\kappa_i&=(h_S+2)/2+B_i/q.
\end{aligned}
\]
明文加法使用 $\lfloor q/T\rfloor C$；精确除法要求明文能被 $p$ 整除，并将 $T$ 改为 $T/p$。标准 FV 乘法给出
\[
\begin{aligned}
B_{\rm mult}={}&C_R\left[\frac T2(B_1+B_2)
+T(\kappa_1B_2+\kappa_2B_1)+\frac TqB_1B_2\right]\\
&+\frac{1+h_S+h_S^2}{2}+B_{\rm ks}.
\end{aligned}
\]
为核实此界，将中心分量提升写成 $AS+B=(q/T)M+\mathcal E+qK$，由范数可得 $\|K\|_\infty\leq\kappa$。展开整数张量乘积，再乘以 $T/q$，保留下列自然误差项：
\[
\begin{aligned}
\mathcal E_{\rm mult}={}&M_1\mathcal E_2+M_2\mathcal E_1
+T(\mathcal E_1K_2+\mathcal E_2K_1)+\frac Tq\mathcal E_1\mathcal E_2\\
&+\rho_0+\rho_1S+\rho_2S^2+E_{\rm ks}\pmod q.
\end{aligned}
\]
省略的进位项都是 $q$ 的整数倍。各张量分量舍入误差 $\rho_i$ 的范数至多为 $1/2$，从而得到上述舍入界。明文模 $T$ 约减也只使精确缩放的明文改变 $q$ 的倍数。这就证明了递推式，而不是预先假定乘法正确。

对于最初的模数切换，输入自然误差界 $B_{\rm in}$ 满足下式即可：
\[
\frac PqB_{\rm in}+\frac{1+h_{\rm in}}2
<\frac{p^{\varepsilon-r}}2.
\]
这里 $h_{\rm in}$ 是乘以输入秘密这一算子的范数上界。旧误差此时成为整数明文 $Z$ 的一部分，并不是新求值密文的加密误差。若输入秘密的加密具有自然误差界 $B_s$，则同态解密后的求值误差账本从下式开始：
\[
B_Z\leq\lambda_{A'}B_s+(r_P/P)\|B'\|_\infty.
\]
若自举密钥使用地板尺度加密，且整数误差界为 $B_{\rm key}$，可取 $B_s=B_{\rm key}+(r_P/P)\|S_{\rm in}\|_\infty$。求值模数可以独立于输入模数选取；为简化记号，上式使用相同的 $q$。

要求显式电路中每个密文节点的界均满足 $B_C<q/(2T_C)$。模数切换不等式保证选择正确的舍入明文；迹对偶映射将其系数放入标量槽位；提升多项式和精确除法恢复模 $p^r$ 的各系数；逆映射重构 $M$。结合上述递推归纳，即可证明最终密文解密为 $M$。在地板尺度记号下，其输出误差为 $E^{\langle b\rangle}=\mathcal E_{p^r}+(r_{p^r}/p^r)M$。新误差界不必小于每一种输入误差，但必须满足后续计算的预算。

\texttt{audit/check\_bootstrap.py} 中的精确公开求值样例覆盖了完全标量分裂的 $\Phi_4,p=5$、在 $r=1$ 和 $r=2$ 下具有二次扩域槽位的 $\Phi_8,p=3$，以及互素二进制替代环 $\Phi_3,p=2$。各例的密文模数均不能被明文模数整除，输入误差和求值密钥误差也都非零。独立检查器在公开求值完成后审计每个中间密文；求值器不接收秘密或解密回调。这些小环样例验证了实现与代数，并不证明密码学安全性，也没有实现参考文献中优化的薄自举算法。

\para{同态解密的目的：}它将低模数密文中依赖秘密密钥的相位，转换为新求值模数及适当尺度下的加密明文，使后续舍入电路能够在不泄露秘密密钥和明文的情况下求值。尺度重新解释只是数字删除实现的一个组成部分，本身并不构成自举的完整解释或正确性证明。



\newpage

\section{CKKS 方案}
\label{sec:ckks}
CKKS 方案支持复数上的 SIMD 近似算术，实数是其中的特例，适用于机器学习等数值计算任务。BFV、BGV 和 TFHE 的原生明文空间是精确的有限集合，但通过适当的编码或电路也能实现定点或浮点运算；它们与 CKKS 的区别在于原生算术及误差的含义。~\cite{ckks-original}

对于二次幂 $n\geq 2$，应区分实代数 $\mathcal{R}_{\langle n\rangle}=\mathbb{R}[X]/(X^n+1)$、整数环 $R=\mathbb{Z}[X]/(X^n+1)$ 和密文环 $R_q=R/qR$。典范嵌入的逆映射将 $n/2$ 个复数槽位映为多项式 $M\in\mathcal{R}_{\langle n\rangle}$。随后，编码将其系数乘以 $\Delta>0$ 并舍入，得到 $R$ 中的整数多项式，再在模 $q$ 下加密。


$ $

\noindent 总体而言，CKKS 的加密过程如下：

\begin{enumerate}
\item \underline{\textsf{编码\textsubscript{1}}：}将目标输入复数编码为实数
\item \underline{\textsf{编码\textsubscript{2}}：}将实数编码为整数
\item \underline{\textsf{加密}：}使用 RLWE 加密整数
\end{enumerate}
得到的 RLWE 密文支持同态加法和乘法。

$ $

\noindent 所有同态运算完成后，CKKS 的解密过程如下：

\begin{enumerate}
\item \underline{\textsf{解密}：}将 RLWE 密文解密为明文整数
\item \underline{\textsf{解码\textsubscript{1}}：}将整数解码为实数
\item \underline{\textsf{解码\textsubscript{2}}：}将实数解码为复数
\end{enumerate}

$ $

CKKS 将放大后的逆嵌入舍入到整数格时，通常会产生编码误差。加密噪声以及后续算术、密钥切换、重缩放和函数近似产生的误差也会影响解码结果。这并非专门由平方根产生的误差，某些输入也能精确编码。与 BFV 的精确明文解密不同，一般的 CKKS 解密会保留数值误差。

CKKS 和 BFV 共用环算术与密钥切换技术，但编码、解码、尺度管理、乘法和自举必须遵循各自不同的明文含义。两者都使用环：$\mathbb{C}$ 本身是域，因此也是环。CKKS 的关键区别是特征为零的近似算术，而不是在明文模数下进行精确算术。



\begin{tcolorbox}[
    title = \textbf{所需背景知识},    % box title
    colback = white,    % light background; tweak to taste
    colframe = black,  % frame colour
    boxrule = 0.8pt,     % line thickness
    left = 1mm, right = 1mm, top = 1mm, bottom = 1mm % inner padding
]

\begin{itemize}
\item \autoref{sec:modulo}: \nameref{sec:modulo}
\item \autoref{sec:group}: \nameref{sec:group}
\item \autoref{sec:field}: \nameref{sec:field}
\item \autoref{sec:order}: \nameref{sec:order}
\item \autoref{sec:polynomial-ring}: \nameref{sec:polynomial-ring}
\item \autoref{sec:decomp}: \nameref{sec:decomp}
\item \autoref{sec:roots}: \nameref{sec:roots}
\item \autoref{sec:cyclotomic}: \nameref{sec:cyclotomic}
\item \autoref{sec:matrix}: \nameref{sec:matrix}
\item \autoref{sec:euler}: \nameref{sec:euler}
\item \autoref{sec:modulus-rescaling}: \nameref{sec:modulus-rescaling}
\item \autoref{sec:chinese-remainder}: \nameref{sec:chinese-remainder}
\item \autoref{sec:taylor-series}: \nameref{sec:taylor-series}
\item \autoref{sec:polynomial-interpolation}: \nameref{sec:polynomial-interpolation}
\item \autoref{sec:ntt}: \nameref{sec:ntt}
\item \autoref{sec:lattice}: \nameref{sec:lattice}
\item \autoref{sec:rlwe}: \nameref{sec:rlwe}
\item \autoref{sec:glwe}: \nameref{sec:glwe}
\item \autoref{sec:glev}: \nameref{sec:glev}
\item \autoref{sec:glwe-add-cipher}: \nameref{sec:glwe-add-cipher}
\item \autoref{sec:glwe-add-plain}: \nameref{sec:glwe-add-plain}
\item \autoref{sec:glwe-mult-plain}: \nameref{sec:glwe-mult-plain}
\item \autoref{subsec:modulus-switch-rlwe}: \nameref{subsec:modulus-switch-rlwe}
\item \autoref{sec:glwe-key-switching}: \nameref{sec:glwe-key-switching}
\item \autoref{sec:bfv}: \nameref{sec:bfv}
\end{itemize}
\end{tcolorbox}

\clearpage


\subsection{编码与解码}
\label{subsec:ckks-encoding-decoding}


CKKS 的编码与解码在本质上非常接近 BFV 的批量编码方案。BFV 基于更新后的 $\hathat W$ 和 $\hathat W^*$ 矩阵（\autoref{subsec:bfv-rotation} 中的总结~\ref*{subsubsec:bfv-rotation-summary}）设计批量编码方案（\autoref{subsec:bfv-enc-dec} 中的总结~\ref*{subsec:bfv-enc-dec}）。具体而言，BFV 在 $X^n+1$ 的各个根，即本原 $(\mu=2n)$ 次单位根处对多项式求值，从而将多项式解码为输入槽位向量（即 $\vec{v} = \hathat W^* \cdot \vec{m}$）；反过来，通过这一运算的逆运算将输入槽位向量编码为多项式（即 $\vec{m} = n^{-1} \cdot \hathat W \cdot I_n^R \cdot \vec{v}$）。这一编码与解码方案以总结~\ref*{subsec:poly-vector-transformation-complex}（\autoref{subsec:poly-vector-transformation-complex}）为基础，后者在环（有限域）上的 $n$ 槽位向量与 $(n-1)$ 次（或更低次）多项式之间构造了如下同构映射：

$\sigma: f(x) \in \mathbb{Z}_t[X] / F(X) \text{ } \longrightarrow \text{ } (f(\omega^1), f(\omega^3)), \cdots, f(\omega^{2n-1})) \in \mathbb{Z}_t^n$

其中，$\omega = g^{\frac{t - 1}{2n}}$ 是定义在模素数 $t$ 的环上的第 $(\mu=2n)$ 个分圆多项式 $X^n + 1$ 的根，即本原 $(\mu=2n)$ 次单位根。

$ $

CKKS 在复单位根上使用相同的求值与插值原理。复数也构成环；这里的完整求值向量属于具有共轭对称性的实代数 $\mathbb{\hat C}^n$。总结~\ref*{subsec:poly-vector-transformation-complex}（\autoref{subsec:poly-vector-transformation-complex}）给出了如下典范嵌入：

$\sigma_c: f(X) \in \mathbb{R}[X]/(X^n + 1) \longrightarrow (f(\omega),f(\omega^3),f(\omega^5), \cdots, f(\omega^{2n-1})) \in \mathbb{\hat{C}}^{n} \text{ } (\longrightarrow \mathbb{C}^{\frac{n}{2}})$


其中，$\omega = e^{i\pi/n}$ 是定义在复数域上的第 $(\mu=2n)$ 个分圆多项式 $X^n + 1$ 的根，即本原 $(\mu=2n)$ 次单位根；$\mathbb{\hat{C}}^{n}$ 是一个特殊的 $n$ 维复向量空间，其后半部分元素是前半部分元素按逆序排列的共轭。$\mathbb{\hat{C}}^{n}$ 与 $\mathbb{{C}}^{\frac{n}{2}}$ 同构，因为 $\mathbb{\hat{C}}^{n}$ 的后半部分元素完全由前半部分元素确定。因此，映射 $\sigma_c$ 本质上是在 $\dfrac{n}{2}$ 槽位复向量 $\vec{v} \in \mathbb{C}^{\frac{n}{2}}$ 与 $(n-1)$ 次（或更低次）实系数多项式 $\mathbb{R}[X] / (X^n + 1)$ 之间的同构。所以，CKKS 的批量编码方案将一个 $\dfrac{n}{2}$ 槽位的复数输入向量编码为 $(n-1)$ 次（或更低次）的实系数多项式，而解码过程则与此相反。

此外，回顾 BFV 中的做法：为了支持输入向量槽位的同态旋转，我们将 $W$ 和 $W^T$ 更新为 $\hathat W$ 和 $\hathat W^*$（见 \autoref{subsubsec:bfv-rotation-summary} 中的总结~\ref*{subsubsec:bfv-rotation-summary}）。同样，CKKS 批量编码方案使用 $\hathat W$ 和 $\hathat W^*$ 替代 $W$ 和 $W^T$，以支持同态旋转。因此，CKKS 批量编码方案中的同构映射更新为：


$\sigma_c: f(X) \in \mathbb{R}[X]/(X^n + 1) \longrightarrow \bm ( f(\omega^{J(0)}),f(\omega^{J(1)}),f(\omega^{J(2)}), \cdots, f(\omega^{J(\frac{n}{2} - 1)}), $

\textcolor{white}{$ f(X) \in \mathbb{R}[X]/(X^n + 1) \longrightarrow \bm ( $} $\cdots, f(\omega^{J_*(0)}),f(\omega^{J_*(1)}),f(\omega^{J_*(2)}), \cdots, f(\omega^{J_*(\frac{n}{2} - 1)}) \bm) \in \mathbb{\hat{C}}^{n} \text{ } (\longrightarrow \mathbb{C}^{\frac{n}{2}})$

其中，$J(h) = 5^h \bmod 2n$ 是用于辅助旋转的公式。


$ $

BFV 与 CKKS 的编码方案具有以下区别：

\begin{itemize}
\item \textbf{输入槽位值的类型：}BFV 的输入槽位值是 ${n}$ 个模 $t$ 的整数，它们被编码为 $n$ 个多项式系数（即模 $t$ 的整数）。而 CKKS 的输入槽位值是 $\dfrac{n}{2}$ 个复数，它们被编码为 $n$ 个多项式系数（即实数）。

\item \textbf{多项式系数的类型：} BFV 编码到 $\mathbb{Z}_t[X]/(X^n+1)$。CKKS 先插值得到实数系数，再将放大后的系数舍入为整数。两种方案实际加密的多项式都具有整数系数。

\item \textbf{尺度因子：} BFV 在密文模数 $q$ 下使用 $\Delta=\lfloor q/t\rfloor$。CKKS 根据目标精度、消息界和模数链选择 $\Delta$。正确解释密文要求包含全部累积误差的解密相位满足中心化不回绕界；系数可以为零或负数。在每一层都要求 $\|\Delta M+E\|_\infty<q_l/2$。

\item \textbf{编码精度：}对于 BFV，在解码过程中，将多项式系数缩小为 $\dfrac{ \Delta\cdot \vec{m} }{\Delta}$ 后能够保留输入值的精度
%as far as $\Delta$ is sufficiently large enough to round off the scaling error (as explained in \autoref{subsubsec:bfv-enc-dec-decoding1})
。而对于 CKKS，如果原始输入值的小数位数过多，尺度因子无法将所有小数位都左移到整数部分，那么缩小后的多项式系数就可能丢失精度。这意味着输入值的一些低位小数会被舍入，导致原始输入的精度损失。例如，设多项式系数为 $m_i = \dfrac{1}{3} = 0.33333\cdots$，尺度因子为 $\Delta = 100$。则缩放后的系数为 $\lceil \Delta m_i \rfloor = 33$，再缩小后得到 $\dfrac{33}{100} = 0.33$。由于 $0.33 \neq 0.33333\cdots$，CKKS 的编码与解码过程并不总能保证完全精确。正是由于这种编码误差，CKKS 被称为\textit{近似}加密方案。编码误差的影响可能随同态运算而增大，使误差幅度增加，解码结果逐渐偏离预期的精确值。减小 CKKS 编码误差的一种方式是增大 $\Delta$，使更多小数位左移到缩放后的整数部分。

\end{itemize}

$ $


\para{\boldmath$\vec{v}_{'} \in \mathbb{\hat C}^n$ 的结构：}注意，总结~\ref*{subsec:poly-vector-transformation-complex}（\autoref{subsec:poly-vector-transformation-complex}）中给出的 $\vec{v}_{'}$ 的原始解码方案为：

$\vec{v}_{'} = \bm{(} \text{ } M(\omega), \text{ } M(\omega^3), \text{ } M(\omega^5), \cdots, M(\omega^{2n-3}), \text{ } M(\omega^{2n-1}) \bm{)}$

它解码得到一个厄米向量：

$\vec{v}_{'} = (v_0, v_1, \cdots, v_{\frac{n}{2} - 1}, \overline v_{\frac{n}{2} - 1}, \cdots, \overline v_1, \overline v_0 )$

其后半部分元素是前半部分元素按逆序排列的共轭。

$ $

不过，将 $W$ 和 $W^T$ 替换为 $\hathat W$ 和 $\hathat W^*$ 后，上述解码方案改为以下支持同态旋转的形式：


$\vec{v}_{'} =  \bm{(} \text{ }
M(\omega^{J(0)}), \text{ } M(\omega^{J(1)}), \text{ } M(\omega^{J(2)}), \cdots,  M(\omega^{J(\frac{n}{2}-1)}), \text{ } M(\omega^{J_*(0)}), \text{ } M(\omega^{J_*(1)}), \cdots,  M(\omega^{J_*(\frac{n}{2}-1)}) \text{ } \bm{)}$

$\textcolor{white}{\vec{v}_{'}} =  \bm{(} \text{ }
M(\omega^{J(0)}), \text{ } M(\omega^{J(1)}), \text{ } M(\omega^{J(2)}), \cdots,  M(\omega^{J(\frac{n}{2}-1)}), \text{ } M(\overline\omega^{J(0)}), \text{ } M(\overline\omega^{J(1)}), \cdots,  M(\overline\omega^{J(\frac{n}{2}-1)}) \text{ } \bm{)}$

\textcolor{red}{ $\rhd$ 因为 $\omega^{-1} = (e^{\frac{i\pi}{n}})^{-1} = e^{\frac{-i\pi}{n}} = \overline\omega$，且给定 $J(h) = 5^h \bmod 2n$ 和 $J_*(h) = -5^h \bmod 2n$}

它解码得到如下\textit{正序}（而非逆序）厄米向量：

$\vec{v}_{'} = (v_0, v_1, \cdots, v_{\frac{n}{2} - 1}, \overline v_0, \overline v_1, \cdots, \overline v_{\frac{n}{2} - 1})$

其后半部分元素是前半部分元素按相同顺序排列的共轭。执行同态旋转时（将在 \autoref{subsec:ckks-rotation} 中说明），与 BFV 的同态旋转一样，$\vec{v}_{'}$ 的前半部分和后半部分元素分别在各自的组内循环旋转。



下面总结 CKKS 的编码与解码过程。它与 BFV 的编码与解码过程相似，后者见 \autoref{subsubsec:bfv-encoding-summary} 中的总结~\ref*{subsubsec:bfv-encoding-summary}：

\begin{tcolorbox}[title={\textbf{\tboxlabel{\ref*{subsec:ckks-encoding-decoding}} CKKS 的编码与解码}}]

\textbf{\underline{输入}：}一个 $\dfrac{n}{2}$ 槽位的复向量 $\vec{v} = (v_0, v_1, \cdots, v_{\frac{n}{2}-1}) \in \mathbb{C}^{\frac{n}{2}}$

\par\noindent\rule{\textwidth}{0.4pt}

\textbf{\underline{编码}：}

$ $

\begin{enumerate}
\item 将 $\vec{v}$ 转换（即同构变换）为如下 $n$ 槽位的\textit{正序}厄米向量 $\vec{v}_{'}$：

$\vec{v}_{'} = (v_0, v_1, \cdots, v_{\frac{n}{2}-1}, \overline v_0, \overline v_1, \cdots, \overline v_{\frac{n}{2}-1}) \in \mathbb{\hat C}^{n}$

\item 应用变换 $\vec{m} = \dfrac{\hathat W \cdot I_n^R \cdot \vec{v}_{'}}{n}$，将 $\vec{v}_{'}$ 转换为实数向量 $\vec{m}$

其中，$\hathat{W}$ 是按如下方式构造的 $n$ 维向量空间的一组基：

$ $

\noindent{\footnotesize{\noindent $\noindent\hathat W = \begin{bmatrix}
1 & 1 & \cdots & 1 & 1 & 1 & \cdots & 1\\
(\omega^{J(\frac{n}{2} - 1)}) & (\omega^{J(\frac{n}{2} - 2)}) & \cdots & (\omega^{J(0)}) & (\omega^{J_*(\frac{n}{2} - 1)}) & (\omega^{J_*(\frac{n}{2} - 2)}) & \cdots & (\omega^{J_*(0)})\\
(\omega^{J(\frac{n}{2} - 1)})^2 & (\omega^{J(\frac{n}{2} - 2)})^2 & \cdots & (\omega^{J(0)})^2 & (\omega^{J_*(\frac{n}{2} - 1)})^2 & (\omega^{J_*(\frac{n}{2} - 2)})^2 & \cdots & (\omega^{J_*(0)})^2 \\
\vdots & \vdots & \ddots & \vdots & \vdots & \ddots & \vdots & \vdots \\
(\omega^{J(\frac{n}{2} - 1)})^{n-1} & (\omega^{J(\frac{n}{2} - 2)})^{n-1} & \cdots & (\omega^{J(0)})^{n-1} & (\omega^{J_*(\frac{n}{2} - 1)})^{n-1} & (\omega^{J_*(\frac{n}{2} - 2)})^{n-1} & \vdots  & (\omega^{J_*(0)})^{n-1}
\end{bmatrix}$}}

\textcolor{red}{ $\rhd$ 其中 $\omega = e^{i\pi/n} = \cos \left(\dfrac{\pi}{n}\right) + i\sin \left(\dfrac{\pi}{n}\right)$、$J(h) = 5^h \bmod 2n$，且 $J_*(h) = -5^h \bmod 2n$}

\small{\noindent $ = \begin{bmatrix}
1 & 1 & \cdots & 1 & 1 & 1 & \cdots & 1\\
(\omega^{J(\frac{n}{2} - 1)}) & (\omega^{J(\frac{n}{2} - 2)}) & \cdots & (\omega^{J(0)}) & (\overline\omega^{J(\frac{n}{2} - 1)}) & (\overline\omega^{J(\frac{n}{2} - 2)}) & \cdots & (\overline\omega^{J(0)})\\
(\omega^{J(\frac{n}{2} - 1)})^2 & (\omega^{J(\frac{n}{2} - 2)})^2 & \cdots & (\omega^{J(0)})^2 & (\overline\omega^{J(\frac{n}{2} - 1)})^2 & (\overline\omega^{J(\frac{n}{2} - 2)})^2 & \cdots & (\overline\omega^{J(0)})^2 \\
\vdots & \vdots & \ddots & \vdots & \vdots & \ddots & \vdots & \vdots \\
(\omega^{J(\frac{n}{2} - 1)})^{n-1} & (\omega^{J(\frac{n}{2} - 2)})^{n-1} & \cdots & (\omega^{J(0)})^{n-1} & (\overline\omega^{J(\frac{n}{2} - 1)})^{n-1} & (\overline\omega^{J(\frac{n}{2} - 2)})^{n-1} & \vdots  & (\overline\omega^{J(0)})^{n-1}
\end{bmatrix}$}

\textcolor{red}{ $\rhd$ 因为 $\omega^{-1} = e^{\frac{-i\pi}{n}} = \overline{e^{\frac{i\pi}{n}}} = \overline\omega$}

$ $

\item 将 $\vec{m}$ 转换为缩放后的整数向量 $\lceil \Delta\vec{m} \rfloor \approx \Delta \vec{m}$，其中 $\Delta$ 是大于 1 的尺度因子，使得 $\Delta m_i$ 在所有情形下，包括整个同态运算过程中，都不会超过 $q_0$ 对应范围的上下界（即满足 $0 \leq \Delta m_i < q_0$ 或 $-\dfrac{q_0}{2} \leq \Delta m_i < \dfrac{q_0}{2}$）。最终编码得到的明文多项式为 $\Delta M = \sum\limits_{i=0}^{n-1} \lceil \Delta m_i \rfloor X^i \text{ } \in \mathbb{Z}_q[X] / (X^n + 1)$。编码过程中的 $\lceil \Delta \vec{m} \rfloor$ 舍入会产生编码误差，因此 CKKS 是一种近似加密方案。

\end{enumerate}

\par\noindent\rule{\textwidth}{0.4pt}

\textbf{\underline{解码}：}由明文多项式 $\Delta M = \sum\limits_{i=0}^{n-1}\Delta m_iX^i$ 恢复 $\vec{m} = \dfrac{\Delta \vec{m}}{\Delta}$。然后，
计算 $\vec{v}_{'} = \hathat W^* \cdot \vec{m}$，其中：

\noindent $\hathat{W}^* = \begin{bmatrix}
1 & (\omega^{J(0)}) & (\omega^{J(0)})^2 & \cdots & (\omega^{J(0)})^{n-1}\\
1 & (\omega^{J(1)}) & (\omega^{J(1)})^2 & \cdots & (\omega^{J(1)})^{n-1}\\
1 & (\omega^{J(2)}) & (\omega^{J(2)})^2 & \cdots & (\omega^{J(2)})^{n-1}\\
\vdots & \vdots & \vdots & \ddots & \vdots \\
1 & (\omega^{J(\frac{n}{2}-1)}) & (\omega^{J(\frac{n}{2}-1)})^2 & \cdots & (\omega^{J(\frac{n}{2}-1)})^{n-1}\\
1 & (\omega^{J_*(0)}) & (\omega^{J_*(0)})^2 & \cdots & (\omega^{J_*(0)})^{n-1}\\
1 & (\omega^{J_*(1)}) & (\omega^{J_*(1)})^2 & \cdots & (\omega^{J_*(1)})^{n-1}\\
1 & (\omega^{J_*(2)}) & (\omega^{J_*(2)})^2 & \cdots & (\omega^{J_*(2)})^{n-1}\\
\vdots & \vdots & \vdots & \ddots & \vdots \\
1 & (\omega^{J_*(\frac{n}{2}-1)}) & (\omega^{J_*(\frac{n}{2}-1)})^2 & \cdots & (\omega^{J_*(\frac{n}{2}-1)})^{n-1}\\
\end{bmatrix}$

$ $

\noindent $ = \begin{bmatrix}
1 & (\omega^{J(0)}) & (\omega^{J(0)})^2 & \cdots & (\omega^{J(0)})^{n-1}\\
1 & (\omega^{J(1)}) & (\omega^{J(1)})^2 & \cdots & (\omega^{J(1)})^{n-1}\\
1 & (\omega^{J(2)}) & (\omega^{J(2)})^2 & \cdots & (\omega^{J(2)})^{n-1}\\
\vdots & \vdots & \vdots & \ddots & \vdots \\
1 & (\omega^{J(\frac{n}{2}-1)}) & (\omega^{J(\frac{n}{2}-1)})^2 & \cdots & (\omega^{J(\frac{n}{2}-1)})^{n-1}\\
1 & (\overline\omega^{J(0)}) & (\overline\omega^{J(0)})^2 & \cdots & (\overline\omega^{J(0)})^{n-1}\\
1 & (\overline\omega^{J(1)}) & (\overline\omega^{J(1)})^2 & \cdots & (\overline\omega^{J(1)})^{n-1}\\
1 & (\overline\omega^{J(2)}) & (\overline\omega^{J(2)})^2 & \cdots & (\overline\omega^{J(2)})^{n-1}\\
\vdots & \vdots & \vdots & \ddots & \vdots \\
1 & (\overline\omega^{J(\frac{n}{2}-1)}) & (\overline\omega^{J(\frac{n}{2}-1)})^2 & \cdots & (\overline\omega^{J(\frac{n}{2}-1)})^{n-1}\\
\end{bmatrix}$

$ $

$ $

最后，只提取\textit{正序}厄米向量 $\vec{v}_{'}$ 的前 $\dfrac{n}{2}$ 个元素，即可恢复输入向量 $\vec{v}$。

\end{tcolorbox}

\begin{comment}
\begin{table}[h] %usepackage{array}
\centering
\begin{tabular}{|>{\centering\arraybackslash}p{0.2\columnwidth}||>{\centering\arraybackslash}p{0.75\columnwidth}|}
\hline \hline
& \textbf{Encoded Polynomial $M(X) \in \mathcal{R}_{\langle n \rangle}$ over $X \in \mathbb{C}$ (complex numbers)} \\ \hline \hline
\textbf{Range of Input Slot Value}& Depends on the range of $M(X)$'s coefficient\\ \hline
\textbf{Range of $\bm{M(X)}$'s coefficients}&Any real number $m_i$ with the custom precision scaling factor $\Delta$ under the constraint: $0 \leq \Delta m_i \leq q-1$ (or $-\dfrac{q}{2} \leq \Delta m_i \leq \dfrac{q}{2} - 1$) \\\hline
%\textbf{Range of $\bm Y$}& $\left[0, \text{ } (1+i)\cdot {t}\cdot 2^{ \lceil \frac{n+1}{2} \rceil}\right]$\\
%& or $\left[-(1+i)\cdot \dfrac{t}{2}\cdot 2^{ \lceil \frac{n + 1}{2} \rceil}, \text{ } (1+i)\cdot \dfrac{t}{2}\cdot 2^{ \lceil \frac{n+1}{2} \rceil}\right]$ (a rough estimate)\\
\hline
\end{tabular}
\caption{The encoding capacity of $M(X) \in \mathcal{R}_{\langle n, t \rangle}$.}
\label{tab:polynomial-encoding-capacity-complex}
\end{table}

\para{CKKS's Encoding Range:} \autoref{tab:polynomial-encoding-capacity-complex} depicts the range of input values that can be encoded by the polynomials $M(X) \in \mathcal{R}_{\langle n \rangle}$ over $X \in \mathbb{C}$. As illustrated in \autoref{tab:polynomial-encoding-capacity-complex}, the encoding range of complex number inputs is determined by the size of the $n$ and $q$ parameters. Therefore, these two parameters should be chosen sufficiently large to encode all values that can be used by an application.
\end{comment}

\para{CKKS 的近似性质：} 令 $P=\lceil\Delta M\rfloor$ 表示对系数逐项舍入到最近整数，则 $P=\Delta M+E_{\mathrm{enc}}$，且 $\|E_{\mathrm{enc}}\|_\infty\leq 1/2$。由单位模根上的三角不等式，每个解码槽位的编码误差至多为 $n/(2\Delta)$。这是保守的系数范数界，应与 CKKS 分析采用的典范嵌入范数区分。加密和后续计算还会产生额外误差，近似解码不会自动将它们消除。


\subsubsection{示例}
\label{subsubsec:ckks-encoding-ex}


设输入复向量的维数为 $\dfrac{n}{2} = 2$，界定多项式的次数为 $n$ = 4，尺度因子为 $\Delta = 1024$。

所使用的 $n$ 维向量空间的基为

$\hathat{W}= \begin{bmatrix}
1 & 1 & 1 & 1\\
\omega^{J(1)} & \omega^{J(0)} & \overline{\omega}^{J(1)} & \overline{\omega}^{J(0)}\\
(\omega^{J(1)})^2 & ({\omega^{J(0)}})^2 & (\overline{\omega}^{J(1)})^2 & (\overline{\omega}^{J(0)})^2\\
(\omega^{J(1)})^3 & (\omega^{J(0)})^3 & (\overline{\omega}^{J(1)})^3 & (\overline{\omega}^{J(0)})^3\\
\end{bmatrix}$
$= \begin{bmatrix}
1 & 1 & 1 & 1\\
\omega^5 & \omega & \overline{\omega}^5 & \overline{\omega}\\
\omega^2 & \omega^2 & \overline{\omega^2} & \overline{\omega}^2\\
\omega^7 & \omega^3 & \overline{\omega}^7 & \overline{\omega}^3\\
\end{bmatrix}$

其中 $\omega = e^{i\pi/n} = \cos \left(\dfrac{\pi}{n}\right) + i\sin \left(\dfrac{\pi}{n}\right)$

$ $


在上述设置下，假设需要编码的输入复向量为 $\vec{v} = (1.1 + 4.3i, \text{ } 3.5 - 1.4i)$。

$ $

首先，构造正序厄米向量 $\vec{v}_{'} = (1.1 + 4.3i, \text{ } 3.5 - 1.4i, \text{ } 1.1 - 4.3i, \text{ } 3.5 + 1.4i)$。

$ $

接下来，应用如下变换，将复向量 $\vec{v}_{'}$ 转换为实数向量 $\vec{m}$：


$\vec{m} = \dfrac{\hathat{W} \cdot I_n^R \cdot \vec{v}_{'}}{n} = \dfrac{1}{4} \cdot \begin{bmatrix}
1 & 1 & 1 & 1\\
\omega^5 & \omega & \overline{\omega}^5 & \overline{\omega}\\
\omega^2 & \omega^2 & \overline{\omega}^2 & \overline{\omega}^2\\
\omega^7 & \omega^3 & \overline{\omega}^7 & \overline{\omega}^3\\
\end{bmatrix} \cdot
\begin{bmatrix}
0 & 0 & 0 & 1 \\
0 & 0 & 1 & 0 \\
0 & 1 & 0 & 0 \\
1 & 0 & 0 & 0
\end{bmatrix}
\cdot
\begin{bmatrix} 1.1 + 4.3i\\3.5 - 1.4i\\1.1 - 4.3i\\3.5 + 1.4i \end{bmatrix}$

$ $

$= \dfrac{\hathat W \cdot I_n^R \cdot \vec{v}_{'}}{n} = \dfrac{1}{4} \cdot\begin{bmatrix}
1 & 1 & 1 & 1\\
\overline{\omega} & \overline{\omega}^5 & \omega  & \omega^5\\
\overline{\omega}^2 & \overline{\omega}^2 & \omega^2 & \omega^2\\
\overline{\omega}^3 & \overline{\omega}^7 & \omega^3 & \omega^7\\
\end{bmatrix}
\cdot
\begin{bmatrix} 1.1 + 4.3i\\3.5 - 1.4i\\1.1 - 4.3i\\3.5 + 1.4i \end{bmatrix}$

$ $

$= \dfrac{1}{4} \cdot
\begin{bmatrix}
(1.1 + 4.3i) + (3.5 - 1.4i) + (1.1 - 4.3i) + (3.5 + 1.4i)\\
(1.1 + 4.3i)\overline{\omega} + (3.5 - 1.4i)\overline{\omega}^5 + (1.1 - 4.3i)\omega+ (3.5 + 1.4i)\omega^5 \\
(1.1 + 4.3i)\overline{\omega^2} + (3.5 - 1.4i)\overline{\omega^2} + (1.1 - 4.3i)\omega^2+ (3.5 + 1.4i)\omega^2 \\
(1.1 + 4.3i)\overline{\omega}^3 + (3.5 - 1.4i)\overline{\omega}^7 + (1.1 - 4.3i)\omega^3 + (3.5 + 1.4i)\omega^7 \\
\end{bmatrix}$

$ $

$= \dfrac{1}{4} \cdot
\begin{bmatrix}
(1.1 + 4.3i) + (3.5 - 1.4i) + (1.1 - 4.3i) + (3.5 + 1.4i)\\
(1.1 + 4.3i)\overline{\omega} + (3.5 - 1.4i)\overline{\omega}^5 + (1.1 - 4.3i)\omega+ (3.5 + 1.4i)\omega^5 \\
(1.1 + 4.3i)\overline{\omega^2} + (3.5 - 1.4i)\overline{\omega^2} + (1.1 - 4.3i)\omega^2+ (3.5 + 1.4i)\omega^2 \\
(1.1 + 4.3i)\overline{\omega}^3 + (3.5 - 1.4i)\overline{\omega}^7 + (1.1 - 4.3i)\omega^3 + (3.5 + 1.4i)\omega^7 \\
\end{bmatrix}$

$ $

$= \dfrac{1}{4} \cdot
\begin{bmatrix}
9.2\\
1.1(\overline{\omega} + \omega) + 4.3i(\overline{\omega} - \omega) +
3.5(\overline{\omega}^5 + \omega^5) - 1.4i(\overline{\omega}^5 - \omega^5)\\
1.1(\overline{\omega}^2 + {\omega}^2) + 4.3i(\overline{\omega}^2 - {\omega}^2) +
3.5(\overline{\omega}^2 + {\omega}^2) - 1.4i(\overline{\omega}^2 - {\omega}^2)\\
1.1(\overline{\omega}^3 + {\omega}^3) + 4.3i(\overline{\omega}^3 - {\omega}^3) +
3.5(\overline{\omega}^7 + {\omega}^7) - 1.4i(\overline{\omega}^7 - {\omega}^7)\\
\end{bmatrix}$

$ $

$= \dfrac{1}{4} \cdot
\begin{bmatrix}
9.2\\
1.1\left(2\cos\dfrac{\pi}{4}\right) - 4.3i\left(2i\sin\dfrac{\pi}{4}\right) + 3.5\left(2\cos\dfrac{5\pi}{4}\right) + 1.4i\left(2i\sin\dfrac{5\pi}{4}\right)\\
1.1\left(2\cos\dfrac{\pi}{2}\right) - 4.3i\left(2i\sin\dfrac{\pi}{2}\right) + 3.5\left(2\cos\dfrac{\pi}{2}\right) + 1.4i\left(2i\sin\dfrac{\pi}{2}\right)\\
1.1\left(2\cos\dfrac{3\pi}{4}\right) - 4.3i\left(2i\sin\dfrac{3\pi}{4}\right) + 3.5 \left(2\cos\dfrac{7\pi}{4}\right) + 1.4i\left(2i\sin\dfrac{7\pi}{4}\right)\\
\end{bmatrix}$

$ $

$= \dfrac{1}{4} \cdot
\begin{bmatrix}
9.2\\
1.1\left(2\dfrac{\sqrt{2}}{2}\right) + 4.3\left(2\dfrac{\sqrt{2}}{2}\right) + 3.5\left(-2\dfrac{\sqrt{2}}{2}\right) - 1.4\left(-2\dfrac{\sqrt{2}}{2}\right)\\
1.1(2\cdot 0) + 4.3(2\cdot1) + 3.5(2\cdot 0) - 1.4(2\cdot 1)\\
1.1\left(-2\dfrac{\sqrt{2}}{2}\right) + 4.3\left(2\dfrac{\sqrt{2}}{2}\right) + 3.5 \left(2\dfrac{\sqrt{2}}{2}\right) - 1.4\left(-2\dfrac{\sqrt{2}}{2}\right)\\
\end{bmatrix}$

$ $

$= 0.25 \cdot
\begin{bmatrix}
9.2\\
1.1\sqrt{2} + 4.3{\sqrt{2}} - 3.5\sqrt{2} + 1.4\sqrt{2}\\
1.1(0) + 4.3(2) - 3.5(0) - 1.4(2)\\
-1.1\sqrt{2} + 4.3\sqrt{2} + 3.5\sqrt{2} + 1.4\sqrt{2}\\
\end{bmatrix} =
\begin{bmatrix}
2.3\\
0.825\sqrt{2}\\
1.45\\
2.025\sqrt{2}\\
\end{bmatrix} \approx (2.3, \text{ } 1.1667, \text{ } 1.45, \text{ } 2.8638)$

$ $

将实数向量 $\vec{m}$ 乘以尺度因子 $\Delta$ 并舍入，得到缩放后的整数向量 $\Delta\vec{m}$：

$\Delta\vec{m} \approx \lceil \Delta \vec{m} \rfloor = \lceil 1024 \cdot (2.3, \text{ } 1.1667, \text{ } 1.45, \text{ } 2.8638) \rfloor = (2355, \text{ } 1195, \text{ } 1485, \text{ } 2933)$

$ $

最终，$\vec{v} = (1.1 + 4.3i, \text{ } 3.5 - 1.4i)$ 被编码为如下明文多项式 $M(X)$：

$\Delta M(X) = 2355 +  1195X + 1485X^2 + 2933X^3 \in \mathcal{R}_{\langle 4 \rangle}$

$ $


为了解码 $\vec{m}$，计算：

$\vec{v}_{'} = \dfrac{\hathat W^*\cdot\Delta\vec{m}}{\Delta} = \begin{bmatrix}
1,\omega,\omega^2, \omega^3\\
1,\omega^5, \omega^2,\omega^7\\
1,\overline{\omega}, \overline{\omega^2}, \overline{\omega^3}\\
1,\overline{\omega^5}, \overline{\omega^2}, \overline{\omega^7}
\end{bmatrix}$
$\cdot \begin{bmatrix}
2355\\1195\\1485\\2933
\end{bmatrix}\cdot \dfrac{1}{1024}$

$=
\begin{bmatrix}
1,\dfrac{\sqrt{2}}{2} + \dfrac{i\sqrt{2}}{2},i, -\dfrac{\sqrt{2}}{2} + \dfrac{i\sqrt{2}}{2}\\
1,-\dfrac{\sqrt{2}}{2} - \dfrac{i\sqrt{2}}{2}, i,\dfrac{\sqrt{2}}{2} - \dfrac{i\sqrt{2}}{2}\\
1,\dfrac{\sqrt{2}}{2} - \dfrac{i\sqrt{2}}{2}, -i, -\dfrac{\sqrt{2}}{2} - \dfrac{i\sqrt{2}}{2}\\
1,-\dfrac{\sqrt{2}}{2} + \dfrac{i\sqrt{2}}{2}, -i, \dfrac{\sqrt{2}}{2} + \dfrac{i\sqrt{2}}{2}
\end{bmatrix}$
$\cdot \begin{bmatrix}
2.2998046875\\1.1669921875\\1.4501953125\\2.8642578125
\end{bmatrix}$

$ $

$ \approx (1.0997+4.3007i, \text{ } 3.5000-1.4003i, \text{ } 1.0997-4.3007i, \text{ } 3.5000+1.4003i)$

$ $

提取厄米向量 $\vec{v}_{'}$ 的前 $\dfrac{n}{2} = 2$ 个元素，即可恢复输入向量：

$(1.0997+4.3007i, \text{ } 3.5000-1.4003i)$

$\approx (1.1 + 4.3i, \text{ } 3.5 - 1.4i) = \vec{v}$ \textcolor{red}{  $\rhd$ 原始输入向量}

$ $

对放大后的逆嵌入作整数舍入，产生了原始槽位与解码槽位之间的小差异。所展示的中间小数系数是近似值，整数编码应根据含 $\sqrt2$ 的精确表达式计算。

$ $

\para{源代码：}运行以下文件可执行 CKKS 编码示例：\href{https://github.com/wang2000star/fhe-textbook-cn/blob/chinese-edition/source%20code/ckks.py}{\underline{此 Python 脚本}}.



\subsection{加密与解密}
\label{subsec:ckks-enc-dec}

CKKS 加密的是 $P=\lceil\Delta M\rfloor\in R$，而不是直接加密任意实数系数。下面给出对称密钥 RLWE 形式；公钥 CKKS 加密使用公开的 RLWE 样本，并具有相应的组合误差。

\begin{tcolorbox}[title={\textbf{\tboxlabel{\ref*{subsec:ckks-enc-dec}} CKKS 加密与解密}}]
选择小秘密 $S\in R$，均匀采样 $A\leftarrow R_q$ 及小整数误差 $E_0$，输出
\[
B=[-AS+P+E_0]_q,\qquad \textsf{ct}=(A,B)\in R_q^2.
\]
定义原始解密相位为
\[
\mathsf{Phase}_S(\textsf{ct})=[B+AS]_q=P+E_0=\Delta M+E,
\qquad E=E_{\mathrm{enc}}+E_0,
\]
其成立条件为 $\|\Delta M+E\|_\infty<q/2$。这里 $[\cdot]_q$ 选择中心化系数代表元。解码时，在选定的根上求相位的值，再除以 $\Delta$：
\[
\widehat v_j=\Delta^{-1}\mathsf{Phase}_S(\textsf{ct})(\omega^{J(j)})
=v_j+\Delta^{-1}E(\omega^{J(j)}).
\]
因此 $|\widehat v_j-v_j|\leq\|E\|_{\mathrm{can},\infty}/\Delta\leq n\|E\|_\infty/\Delta$。仅有 $|e_i|<\Delta/2$ 并不能推出槽位精度为 $1/\Delta$。舍入到 $1/\Delta$ 的整数倍也不能消除一般的 CKKS 噪声。
\end{tcolorbox}

下文在代数推导中使用 $\textsf{RLWE}^{-1}$ 时，均指这一原始相位；除以所记录的尺度和槽位解码是独立操作。$\textsf{RLWE}_{S,\sigma}(P)$ 表示加密关系，并不声称计算后的密文仍具有原来的独立误差分布 $\chi_\sigma$。

\subsection{密文与密文相加}
\label{subsec:ckks-add-cipher}

CKKS 与 BFV 一样逐分量相加密文（总结~\ref*{subsec:bfv-add-cipher}、\autoref{subsec:bfv-add-cipher}）。密文的密钥、模数和尺度必须相容。结果的相位误差为 $E^{\langle1\rangle}+E^{\langle2\rangle}$，相位之和必须满足不回绕界。




\begin{tcolorbox}[title={\textbf{\tboxlabel{\ref*{subsec:ckks-add-cipher}} CKKS 密文与密文相加}}]
$\textsf{RLWE}_{S, \sigma}(\Delta M^{\langle 1 \rangle} ) + \textsf{RLWE}_{S, \sigma}(\Delta M^{\langle 2 \rangle} ) $

$ = ( A^{\langle 1 \rangle}, \text{ } B^{\langle 1 \rangle}) + (A^{\langle 2 \rangle}, \text{ } B^{\langle 2 \rangle}) $

$ = ( A^{\langle 1 \rangle} + A^{\langle 2 \rangle}, \text{ } B^{\langle 1 \rangle} + B^{\langle 2 \rangle} ) $

$= \textsf{RLWE}_{S, \sigma}(\Delta(M^{\langle 1 \rangle} + M^{\langle 2 \rangle}) )$
\end{tcolorbox}


\subsection{密文与明文相加}
\label{subsec:ckks-add-plain}

CKKS 与 BFV 一样将明文加到 $B$ 分量（总结~\ref*{subsec:bfv-add-plain}、\autoref{subsec:bfv-add-plain}）。应按密文的实际尺度 $\Delta$ 编码明文，并将其编码误差计入误差预算。下述简写中的 $\Delta\Lambda$ 表示舍入后的整数编码。



\begin{tcolorbox}[title={\textbf{\tboxlabel{\ref*{subsec:ckks-add-plain}} CKKS 密文与明文相加}}]
$\textsf{RLWE}_{S, \sigma}(\Delta M) + \Delta\Lambda $

$=  (A, \text{ } B) + \Delta\Lambda$

$=  (A, \text{ } B + \Delta\cdot\Lambda)$

$= \textsf{RLWE}_{S, \sigma}(\Delta (M + \Lambda) )$
\end{tcolorbox}




\subsection{密文与密文相乘}
\label{subsec:ckks-mult-cipher}

CKKS 乘法先在当前模数下对密文作张量积，然后通常进行重线性化和重缩放。张量积将输入尺度 $\Delta_1,\Delta_2$ 变为 $\Delta_1\Delta_2$；除以 $w_l=q_l/q_{l-1}$ 的重缩放将其变为 $\Delta_1\Delta_2/w_l$，并消耗一层。下述推导采用定理~\ref*{subsec:glwe-alternative}（\autoref{subsec:glwe-alternative}）中的符号约定 $B=-AS+\Delta M+E\pmod q$。我们先推导包含误差的精确相位，再说明两种重线性化方法及重缩放。


\subsubsection{合成密文的推导}
\label{subsubsec:ckks-mult-cipher-relation}

对于模数和秘密密钥相同的两个密文，设其相位分别为 $\Delta M^{\langle 1\rangle}+E^{\langle 1\rangle}$ 和 $\Delta M^{\langle 2\rangle}+E^{\langle 2\rangle}$。在模 $q$ 下定义
\[
D_0=B^{\langle 1\rangle}B^{\langle 2\rangle},\quad
D_1=A^{\langle 1\rangle}B^{\langle 2\rangle}+A^{\langle 2\rangle}B^{\langle 1\rangle},\quad
D_2=A^{\langle 1\rangle}A^{\langle 2\rangle}.
\]
精确的相位恒等式为
\[
D_0+D_1S+D_2S^2\equiv\Delta^2M^{\langle 1\rangle}M^{\langle 2\rangle}+F\pmod q,
\]
\[
F=\Delta M^{\langle 1\rangle}E^{\langle 2\rangle}
+\Delta M^{\langle 2\rangle}E^{\langle 1\rangle}
+E^{\langle 1\rangle}E^{\langle 2\rangle}.
\]
CKKS 解码会保留这些误差项。重线性化构造 $C_\alpha=(D_1,D_0)$，以及相位为 $D_2S^2+E_{\mathrm{ks}}$ 的多项式对 $C_\beta$。于是 $C_\alpha+C_\beta\in R_q^2$ 的相位为 $\Delta^2M^{\langle 1\rangle}M^{\langle 2\rangle}+F+E_{\mathrm{ks}}$。这并未将尺度恢复为 $\Delta$；恢复尺度是随后重缩放步骤的作用。如果不进行重线性化，密文长度会随关于 $S$ 的次数增长；反复平方会在每一步将该次数翻倍。

\subsubsection{重线性化方法 1：密文分解}
\label{subsubsec:relinearization-gadget-decomposition}

对于任意 $q$，使用显式的整数 gadget：选择整数 $\beta\geq 2$，令 $l=\lceil\log_\beta q\rceil$，并取 $g_i=\beta^{i-1}$，其中 $1\leq i\leq l$。将 $D_2$ 各系数的无符号代表元分解为数位 $d_i\in\{0,\ldots,\beta-1\}$，使
\[
D_2\equiv\sum_{i=1}^l d_i g_i\pmod q.
\]
这里的升序 gadget 与降序权重 $q/\beta^i$ 是不同约定，同一个数位向量不能同时配用这两套权重。生成求值密钥对 $K_i=(A_i,B_i)$，使其满足
\[
B_i+A_iS\equiv g_iS^2+E_i\pmod q.
\]
于是
\[
C_\beta=\sum_{i=1}^l d_iK_i,\qquad
\mathsf{Phase}_S(C_\beta)\equiv D_2S^2+\sum_{i=1}^l d_iE_i\pmod q.
\]
对于负循环系数范数，有 $\|\sum_i d_iE_i\|_\infty\leq n(\beta-1)\sum_i\|E_i\|_\infty$。使用平衡数位或截断分解时，必须采用相应的数位界与残差界。求值密钥加密的是 $S^2$ 的 gadget 倍数，不附加 CKKS 消息尺度 $\Delta$。

\subsubsection{重线性化方法 2：密文模数切换}
\label{subsubsec:relinearization-modulus-switch}

另一种重线性化方法使用辅助正整数 $g$ 和模 $gq$ 下的求值密钥。应根据密钥切换的误差预算选择 $g$，它不必随机，也不必是二次幂。在 $R_{gq}$ 中均匀采样 $A'$，并采样小整数误差 $E'$，令
\[
\mathit{evk}_g=(A',B'),\qquad B'=[-A'S+gS^2+E']_{gq}.
\]
将 $D_2\in R_q$ 的系数提升为固定的中心化整数。计算多项式对 $Z=[D_2\mathit{evk}_g]_{gq}$，并令
\[
C_\beta=\left[\left\lceil Z/g\right\rfloor\right]_q.
\]
记 $R_A=\lceil Z_A/g\rfloor-Z_A/g$、$R_B=\lceil Z_B/g\rfloor-Z_B/g$，其每个系数的绝对值均不超过 $1/2$。在模 $q$ 下，输出相位满足
\[
\mathsf{Phase}_S(C_\beta)\equiv D_2S^2+D_2E'/g+R_B+R_AS\pmod q.
\]
因此，新增误差的系数范数满足
\[
\|E_{\mathrm{ks}}\|_\infty\leq
\frac{n\|D_2\|_\infty\|E'\|_\infty}{g}
+\frac{1+\|S\|_1}{2}.
\]
对 $Z$ 取模产生的 $gq$ 倍数，在除法后变为 $q$ 的倍数，因此在模 $q$ 下消失。舍入不对加法分配，也不能与乘以 $S$ 交换；其贡献应显式写成 $R_B+R_AS$。较大的 $g$ 会抑制求值密钥误差，但不能消除此舍入项。$g$ 的选择还会影响扩展模数求值密钥的安全参数。


\subsubsection{重缩放}
\label{subsubsec:ckks-mult-cipher-rescale}

设 $q_l=w_lq_{l-1}$，重缩放前的密文相位为 $\Delta_{\mathrm{in}}M+F$。重缩放将中心化密文代表元逐系数除以 $w_l$，再舍入到最近整数：
\[
\widehat A=[\lceil A/w_l\rfloor]_{q_{l-1}},\qquad
\widehat B=[\lceil B/w_l\rfloor]_{q_{l-1}}.
\]
令 $R_A=\lceil A/w_l\rfloor-A/w_l$、$R_B=\lceil B/w_l\rfloor-B/w_l$，则
\[
\widehat B+\widehat AS\equiv
\frac{\Delta_{\mathrm{in}}}{w_l}M+
\frac{F}{w_l}+R_B+R_AS\pmod {q_{l-1}}.
\]
因此实际输出尺度为 $\Delta_{\mathrm{out}}=\Delta_{\mathrm{in}}/w_l$，尺度元数据本身不作舍入。对于尺度分别为 $\Delta_1,\Delta_2$ 的密文乘积，有 $\Delta_{\mathrm{out}}=\Delta_1\Delta_2/w_l$。若 $\Delta_1=\Delta_2=\Delta$ 且 $w_l\approx\Delta$，结果接近 $\Delta$，但一般不与之相等。将尺度重新标为精确的 $\Delta$ 会引入依赖消息的误差，必须另行分析精度。

新的相位误差为 $\widehat F=F/w_l+R_B+R_AS$，不是逐系数的 $\lceil F/w_l\rfloor$。确定性界为
\[
\|\widehat F\|_\infty\leq\|F\|_\infty/w_l+(1+\|S\|_1)/2.
\]
正确解释结果还要求完整的输出相位位于模 $q_{l-1}$ 的中心化区间内。

对于 RNS 链 $q_l=\prod_{i=0}^l w_i$，应根据算术效率和目标尺度安排选择素数。RNS 改变的是表示形式，而不是 $q_l$ 的总信息量。RLWE 安全性取决于维数、总模数、秘密分布、误差分布和求值密钥参数，而不是仅由最小及最大的 RNS 素数决定。

\para{加密时的初始尺度因子 $\Delta$：} 应根据目标精度和可用模数空间选择 $\Delta$，并使 $q_0$ 足以容纳预期最终消息及误差。令 $\Delta$ 接近重缩放素数通常较为方便，但并非必须取 $\Delta=w_L$。

\para{噪声增长：} 乘法会传播消息与误差的交叉项及误差乘积，而重缩放会新增舍入误差。界取决于消息与误差的大小；并非所有输入的误差都会增长，例如乘以零时就未必如此。单独重缩放不能恢复已经损失的数值精度。

\subsubsection{BFV 与 CKKS 自举的比较}
\label{subsubsec:bfv-bootstrapping-ckks-comparison}

两种方案都通过求值适当的解密相关函数来刷新密文，但具体步骤与明文含义不同。在下文的 CKKS 流程中，提升密文系数会在较大模数下暴露原始相位 $\Delta M+E+q_0K$，无需另行计算包含加密秘密的内积。线性的系数与槽位变换以及近似模约减函数会近似消除 $q_0K$ 的贡献。BFV 自举则求值有限明文空间上的精确解密函数，并须满足相应参数和数位提取条件。

\para{乘法层级：} 即使保持固定密文模数，BFV 仍有有限的噪声预算。在 CKKS 标准的乘法、重线性化、重缩放流程中，重缩放会消耗一层模数链。到达 $q_0$ 意味着不能再按该安排进行重缩放，而不是在代数上禁止一切乘法。自举可补充可用层级，同时也会消耗若干层（\autoref{subsec:ckks-bootstrapping}）。

\para{噪声处理的局限：} 重缩放将原有相位误差和模数除以同一因子，并增加舍入误差。因此，它本身不会改善原有误差与模数的比值，也不能恢复损失的数值精度。CKKS 自举扩大可用模数，但会保留已有数值误差，并增加自身的近似与求值误差。


\subsubsection{总结}
\label{subsubsec:ckks-mult-cipher-summary}

\begin{tcolorbox}[title={\textbf{\tboxlabel{\ref*{subsubsec:ckks-mult-cipher-summary}} CKKS 密文与密文乘法}}]
对于密钥和模数 $q_l$ 相同、尺度分别为 $\Delta_1,\Delta_2$ 的密文：
\begin{enumerate}
\item 在模 $q_l$ 下计算 $D_0=B_1B_2$、$D_1=A_1B_2+A_2B_1$ 和 $D_2=A_1A_2$。它们的相位是两个输入相位的乘积。
\item 使用前述任一种方法重线性化，得到 $(A,B)=(D_1,D_0)+C_\beta$，其相位为 $\Delta_1\Delta_2 M_1M_2+F$；$F$ 包含乘法与密钥切换误差。
\item 若计算安排要求重缩放，令 $q_{l-1}=q_l/w_l$，并计算
\[
(\widehat A,\widehat B)=[(\lceil A/w_l\rfloor,\lceil B/w_l\rfloor)]_{q_{l-1}},
\quad \Delta_{\mathrm{out}}=\Delta_1\Delta_2/w_l.
\]
输出误差为 $F/w_l+R_B+R_AS$，并须满足上述不回绕与精度界。
\end{enumerate}
\para{交换\textsf{重线性化}与\textsf{重缩放}的顺序：} 可以先重缩放三分量密文，其舍入误差为 $R_0+R_1S+R_2S^2$，再使用适用于较低模数的密钥重线性化。两种安排并不是相同的密文计算，其误差和成本界也不同。
\end{tcolorbox}


\subsection{密文与明文相乘}
\label{subsec:ckks-mult-plain}

以选定的尺度 $\Delta_\Lambda$ 编码明文乘数 $\Lambda$，得到 $P_\Lambda=\lceil\Delta_\Lambda\Lambda\rfloor$。对于尺度为 $\Delta$ 的密文 $(A,B)$，计算
\[
(A',B')=(AP_\Lambda,BP_\Lambda)\pmod q,
\qquad \Delta'=\Delta\Delta_\Lambda.
\]
相位误差包含输入密文误差乘以 $P_\Lambda$ 的贡献，以及明文编码误差的贡献。若需要除以 $w_l$ 重缩放，则使用 \autoref{subsubsec:ckks-mult-cipher-summary} 的流程，输出尺度为 $\Delta\Delta_\Lambda/w_l$，并增加相应的舍入误差。这种常见安排消耗一层，但并非每次明文乘法都必须如此：精确表示的小整数常量可以采用 $\Delta_\Lambda=1$；模数和精度预算允许时，也可以延后重缩放。


\subsection{\textsf{ModDrop}}
\label{subsec:ckks-moddrop}

对于 $q_{l-1}\mid q_l$，定义
\[
(A',B')=([A]_{q_{l-1}},[B]_{q_{l-1}}).
\]
若 $A'=A-K_Aq_{l-1}$、$B'=B-K_Bq_{l-1}$，则
\[
A'S+B'=AS+B-(K_AS+K_B)q_{l-1}.
\]
因此新相位与原相位模 $q_{l-1}$ 同余。若 $\|\Delta M+E\|_\infty<q_{l-1}/2$，其中心化代表元仍精确等于 $\Delta M+E$；否则相位可能回绕，从而改变解码值。

\begin{tcolorbox}[title={\textbf{\tboxlabel{\ref*{subsec:ckks-moddrop}} CKKS \textsf{ModDrop}}}]
将每个密文分量对当前模数的一个因子取模。只有完整相位落在新的中心化区间内，所记录的尺度和已有相位误差才保持不变；可用模数空间会减少。
\end{tcolorbox}

\subsubsection{模数切换与 \textsf{ModDrop} 的区别}
\label{subsubsec:ckks-moddrop-vs-modswitch}

两种操作都会降低密文模数。重缩放还将记录的尺度和原有误差除以 $q_l/q_{l-1}$，并引入舍入误差。\textsf{ModDrop} 保持尺度不变，但只有满足其不回绕条件时才保持相位。它也可用于加法或乘法前对齐不同层级的密文，并非只能在不再进行乘法时使用。

\subsection{同态密钥切换}
\label{subsec:ckks-key-switching}

选取 \autoref{subsubsec:relinearization-gadget-decomposition} 中的精确整数 gadget $g_i$，并分解 $A\equiv\sum_i d_ig_i\pmod q$。新秘密密钥 $S'$ 下的求值密钥 $(a_i,b_i)$ 满足 $b_i+a_iS'\equiv g_iS+E_i\pmod q$。于是
\begin{tcolorbox}[title={\textbf{\tboxlabel{\ref*{subsec:ckks-key-switching}} CKKS 密钥切换}}]
\[
(A',B')=(0,B)+\sum_i d_i(a_i,b_i),\qquad
B'+A'S'\equiv B+AS+\sum_i d_iE_i\pmod q.
\]
\end{tcolorbox}
尺度不变，误差增加所列密钥切换项。正确解释要求新相位适合该模数。这与 \autoref{subsec:bfv-key-switching} 使用相同的基础 gadget 技术，但采用 CKKS 的近似解码含义。


\subsection{输入向量槽位的同态旋转}
\label{subsec:ckks-rotation}

BFV 与 CKKS 槽位旋转采用相同的根指数置换，但 BFV 在其明文环中的根上求值，CKKS 则在复数根上求值。在 \autoref{subsec:ckks-encoding-decoding} 的编码约定中，逆变换为 $n^{-1}\hathat W I_n^R$，解码变换为 $\hathat W^*$，矩阵顺序不可颠倒。

对于 $J(j)=5^j\bmod2n$，自同构 $M(X)\mapsto M(X^{J(h)})$ 将槽位 $j$ 变为原来的槽位 $j+h\bmod(n/2)$，因为 $5^j5^h=5^{j+h}$。在 $J_*(j)=-5^j\bmod2n$ 的约定下，共轭半部向相同方向旋转。
\begin{tcolorbox}[title={\textbf{\tboxlabel{\ref*{subsec:ckks-rotation}} CKKS 输入向量槽位的同态旋转}}]
\begin{enumerate}
\item 对 $A$ 与 $B$ 都施加 $X\mapsto X^{J(h)}$，并模 $X^n+1$ 约减。秘密密钥变为 $S(X^{J(h)})$。
\item 使用 \autoref{subsec:ckks-key-switching} 中的精确 gadget 密钥切换，从 $S(X^{J(h)})$ 切换到 $S(X)$。将新增密钥切换误差计入输出相位。
\end{enumerate}
结果将 $n/2$ 个独立复数槽位向左旋转 $h$ 个位置。
\end{tcolorbox}

\para{半部槽位内的旋转：} 单个 $5^h$ 自同构分别旋转两个轨道，但并不禁止使用掩码与若干自同构组合构造其他置换。CKKS 的后半部由共轭关系确定，而完全分裂的 BFV 打包允许独立选择全部 $n$ 个槽位。


\subsubsection{示例}
\label{subsubsec:ckks-rotation-ex}

本小节给出以下两个示例：
\begin{enumerate}
\item 基于更新后的编码与解码矩阵 $\hathat W$ 和 $\hathat W^*$，将输入向量 $\vec{v}$ 编码为明文多项式 $M(X)$
\item 将编码后的明文 $M(X)$ 更新为 $M(X^{J(h)})$，使输入向量 $\vec{v}$ 的所有元素向左旋转 $h$ 个位置
\end{enumerate}

$ $

仍使用 \autoref{subsubsec:ckks-encoding-ex} 中相同的输入向量 $\vec{v}$：$\vec{v} = (1.1 + 4.3i, 3.5 - 1.4i)$。

回顾一下，$\vec{v}$ 编码得到的明文多项式如下：

$\Delta M(X) = 2355 + 1195X + 1485X^2 + 2933X^3 \in \mathcal{R}_{\langle 4 \rangle}=\mathbb{R}[X]/(X^4+1)$


假设希望将输入向量 $\vec{v}$ 向左旋转 1 个位置：

$\vec{v}^{\langle h=1 \rangle} = (3.5 - 1.4i, 1.1 + 4.3i)$

$ $

因此，将 $\Delta M(X)$ 更新为 $\Delta M(X^{J(1)})$：

$\Delta M(X^{J(1)}) = \Delta M(X^{5}) = 2355 + 1195(X^5) + 1485(X^5)^2 + 2933(X^5)^3$

$= 2355 + 1195X^5 + 1485X^{10} + 2933X^{15}$

$= 2355 + 1195X\cdot(-1) + 1485X^{2}\cdot(-1)\cdot(-1) + 2933X^{3}\cdot(-1)\cdot(-1)\cdot(-1)$

$= 2355 - 1195X + 1485X^2 - 2933X^3$

$ $

旋转后的\textit{正序}厄米输入向量计算如下：

$\dfrac{\hathat W^* \cdot \Delta \vec{m}}{\Delta} = \begin{bmatrix}
1,\omega,\omega^2, \omega^3\\
1,\omega^5, \omega^{10},\omega^{15}\\
1,\overline{\omega}, \overline{\omega^2}, \overline{\omega^3}\\
1,\overline{\omega^5}, \overline{\omega^{10}}, \overline{\omega^{15}}
\end{bmatrix} \cdot \begin{bmatrix}
2355\\- 1195\\1485\\-2933
\end{bmatrix}\cdot \dfrac{1}{1024}$

$= \dfrac{\hathat W^* \cdot \Delta \vec{m}}{\Delta} = \begin{bmatrix}
1,\omega,\omega^2, \omega^3\\
1,\omega^5, \omega^{2},\omega^{7}\\
1,\overline{\omega}, \overline{\omega^2}, \overline{\omega^3}\\
1,\overline{\omega^5}, \overline{\omega^{2}}, \overline{\omega^{7}}
\end{bmatrix} \cdot \begin{bmatrix}
2355\\- 1195\\1485\\-2933
\end{bmatrix}\cdot \dfrac{1}{1024}$

$= \dfrac{\hathat W^* \cdot \Delta \vec{m}}{\Delta} = \begin{bmatrix}
1,\omega,\omega^2, \omega^3\\
1,\omega^5, \omega^{2},\omega^{7}\\
1,\omega^7, \omega^6, \omega^5\\
1,\omega^3, \omega^{6}, \omega
\end{bmatrix} \cdot \begin{bmatrix}
2355\\- 1195\\1485\\-2933
\end{bmatrix}\cdot \dfrac{1}{1024}$

$=
\begin{bmatrix}
1,\dfrac{\sqrt{2}}{2} + \dfrac{i\sqrt{2}}{2},i, -\dfrac{\sqrt{2}}{2} + \dfrac{i\sqrt{2}}{2}\\
1,-\dfrac{\sqrt{2}}{2} - \dfrac{i\sqrt{2}}{2}, i,\dfrac{\sqrt{2}}{2} - \dfrac{i\sqrt{2}}{2}\\
1,\dfrac{\sqrt{2}}{2} - \dfrac{i\sqrt{2}}{2}, -i, -\dfrac{\sqrt{2}}{2} - \dfrac{i\sqrt{2}}{2}\\
1,-\dfrac{\sqrt{2}}{2} + \dfrac{i\sqrt{2}}{2}, -i, \dfrac{\sqrt{2}}{2} + \dfrac{i\sqrt{2}}{2}
\end{bmatrix}$
$\cdot \begin{bmatrix}
2.2998046875\\-1.1669921875\\1.4501953125\\-2.8642578125
\end{bmatrix}$

$ $

$ \approx (3.500 - 1.4003i, \text{ } 1.0997 + 4.3007i,\text{ } 3.500 + 1.4003i, \text{ } 1.0997 - 4.3007i)$

$ $

提取上述厄米向量的前 $\dfrac{n}{2} = 2$ 个元素，即可恢复输入向量：

$(3.500 - 1.4003i, \text{ } 1.0997 + 4.3007i)$

$ \approx (3.5 - 1.4i, \text{ } 1.1 + 4.3i)$
$ = \vec{v}^{\langle h=1 \rangle}$ \textcolor{red}{  $\rhd$ 原始输入向量 $\vec{v}$ 向左旋转 1 个位置}



$ $

实际求值时，对密文分量施加自同构，得到
\[
(A(X^{J(1)}),B(X^{J(1)})).
\]
此密文的秘密密钥为 $S(X^{J(1)})$。在使用原密钥解密前，应先将其密钥切换回 $S(X)$。解码结果近似于旋转后的向量 $\vec v^{\langle h=1\rangle}$，误差包含原编码误差及新增的密钥切换误差。

$ $


\para{源代码：}运行以下文件可执行 CKKS 输入向量槽位同态旋转的示例：\href{https://github.com/wang2000star/fhe-textbook-cn/blob/chinese-edition/source%20code/ckks.py}{\underline{此 Python 脚本}}.





\subsection{对 CKKS 编码的进一步思考}
\label{subsec:ckks-encoding-contemplate}

\para{CKKS 编码为何使用 $\bm{(\mu=2n)}$ 次分圆多项式：} 在 $X^n+1$ 的 $n$ 个不同根上求值，会得到可逆的范德蒙德矩阵。在域上，任意一组互异求值点都给出可逆范德蒙德矩阵；单位根还使求值与模 $X^n+1$ 的乘法相容，提供共轭配对，并支持高效傅里叶变换。二次幂分圆环是一种有用的实现选择，而不是因为只有它们的范德蒙德矩阵可逆。

采用槽位顺序 $\omega^{5^j}$ 时，施加 $X\mapsto X^{5^h}$ 会置换求值矩阵（即解码矩阵）的行。逆变换才是编码矩阵，不能将两者的方向混淆。共轭配对将独立复数坐标数降为 $n/2$，同时保留 $n$ 个实系数坐标。


\subsection{同态共轭}
\label{subsec:ckks-conjugation}

对于实系数多项式 $M$ 和 $\zeta_j=\omega^{J(j)}$，共轭满足
\[
\overline{M(\zeta_j)}=M(\overline{\zeta_j})=M(\zeta_j^{-1}).
\]
因此，正序完整向量由 $(\vec v,\overline{\vec v})$ 变为 $(\overline{\vec v},\vec v)$。应使用与编码一致的选定根 $J(j)=5^j\bmod 2n$，不能将其顺序替换为 $1,3,\ldots,n-1$。

\begin{tcolorbox}[title={\textbf{\tboxlabel{\ref*{subsec:ckks-conjugation}} CKKS 同态共轭}}]
\begin{enumerate}
\item 对密文的两个分量施加环自同构 $X\mapsto X^{-1}=X^{2n-1}$，并对 $X^n+1$ 取模。
\item 从 $S(X^{-1})$ 密钥切换到 $S(X)$，同时计入新增密钥切换误差。
\end{enumerate}
\end{tcolorbox}

\subsection{密文的稀疏打包}
\label{subsec:ckks-sparse-packing}

在 \autoref{subsec:ckks-encoding-decoding} 中，我们学习了 CKKS 编码方案：将具有 $\dfrac{n}{2}$ 个槽位的输入向量（即 $\dfrac{n}{2}$ 槽位输入向量）编码为 $(n-1)$ 次多项式。虽然多项式环的次数 $n$ 作为安全参数在 CKKS 设置阶段便已固定，但某些应用需要的输入向量槽位数可能少于 $\dfrac{n}{2}$。假设在 $\frac{n}{2}$ 个槽位中仅需使用 $\frac{n'}{2}$ 个槽位，其中 $n'\geq2$ 且整除 $n$，则相应的输入向量与编码后的多项式具有以下特殊性质：


\begin{tcolorbox}[title={\textbf{\tboxlabel{\ref*{subsec:ckks-sparse-packing}} CKKS 的稀疏打包多项式与密文}}]

假设一个 $\frac{n}{2}$ 槽位输入向量被编码为 $\mathbb{R}[X] / (X^n + 1)$ 中的多项式，并设 $n'\geq2$ 且整除 $n$。定义多项式 $M(X) \in \mathbb{R}[X] / (X^n + 1)$：仅指数为 $\frac{n}{n'}$ 的倍数的项可以具有非零系数，其余项的系数均为零（即 $M(X) = c_0 + c_{\frac{n}{n'}}X^{\frac{n}{n'}} + c_{\frac{2n}{n'}}X^{\frac{2n}{n'}} + \cdots + c_{n - \frac{n}{n'}}X^{n - \frac{n}{n'}}$）。于是可将 $M(X)$ 写为某个 $M_Y(Y) \in \mathbb{R}[Y] / (Y^{n'} + 1)$，其中 $Y = X^{\frac{n}{n'}}$，因此 $M(X) = M_Y(X^{\frac{n}{n'}})$。

$ $

那么，以下结论成立：

\begin{enumerate}
    \item 每个 $M_Y(Y) \in \mathbb{R}[Y]/(Y^{n'} + 1)$ 都同构地映射（即解码）为某个 $\dfrac{n}{2}$ 槽位输入向量，该向量由连续 $\dfrac{n'}{2}$ 个槽位值重复 $\dfrac{n}{n'}$ 次构成。
    \item 反过来，如果一个 $\dfrac{n}{2}$ 槽位输入向量由前 $\dfrac{n'}{2}$ 个连续槽位值重复 $\dfrac{n}{n'}$ 次构成，那么该向量将编码为某个 $M_Y(Y) \in \mathbb{R}[Y] / (Y^{n'} + 1)$，即 $\mathbb{R}[X] / (X^n + 1)$ 中指数不是 $\dfrac{n}{n'}$ 的倍数的各项系数均为零的多项式。
\end{enumerate}

\end{tcolorbox}

假设 $n'\geq2$ 且整除二次幂 $n$。下面证明正向包含关系。反向结论还需要注意：子环 $\mathbb{R}[Y]/(Y^{n'}+1)$ 与重复槽位向量空间的实维数都为 $n'$。单射的典范嵌入将前者映入后者，因此维数相等保证它满射到该子空间。仅指出整个编码映射是同构，而不论证这些子空间，尚不足以证明反向结论。

\subsubsection{正向证明：稀疏打包密文的解码}

下面证明，对于每个 $M_Y(Y) \in \mathbb{R}[Y] / (Y^{n'} + 1)$（即 $\mathbb{R}[X] / (X^n + 1)$ 中仅指数为 $\dfrac{n}{n'}$ 的倍数的项可能具有非零系数、其余项系数均为零的多项式），其解码结果是某个 $\dfrac{n}{2}$ 槽位输入向量，由前 $\dfrac{n'}{2}$ 个连续槽位值重复 $\dfrac{n}{n'}$ 次构成。

要将 $M(X)$ 解码为输入向量，应从 $X^n + 1$ 的 $n$ 个不同本原 $(\mu=2n)$ 次单位根中选取 $\frac{n}{2}$ 个，对 $M(X)$ 求值。求值向量为：

$\bm{(} \text{ }
M(\omega^{J(0)}), \text{ } M(\omega^{J(1)}), \cdots,  M(\omega^{J(\frac{n}{2}-1)})\bm{)}$

\textcolor{red}{ $\rhd$ 其中 $\omega=e^{\frac{i\pi}{n}}$ 是本原（$\mu=2n$）次单位根的基值与生成元}

$ $

由于 $M(X) = M_Y(X^{\frac{n}{n'}})$，上述求值等价于计算：

$\bm{(} \text{ }
M_Y((\omega^{J(0)})^{\frac{n}{n'}}), \text{ } M_Y((\omega^{J(1)})^{\frac{n}{n'}}), \cdots,  M_Y((\omega^{J(\frac{n}{2}-1)})^{\frac{n}{n'}}) \bm{)}$

$ $

$= \bm{(} \text{ }
M_Y((\omega^{\frac{n}{n'}})^{J(0)}), \text{ } M_Y((\omega^{\frac{n}{n'}})^{J(1)}), \cdots,  M_Y((\omega^{\frac{n}{n'}})^{J(\frac{n}{2}-1)})\bm{)}$

$ $

$= \bm{(} \text{ }
M_Y(\xi^{J(0)}), \text{ } M_Y(\xi^{J(1)}), \cdots,  M_Y(\xi^{J(\frac{n}{2}-1)})\bm{)}$

\textcolor{red}{ $\rhd$ 其中 $\xi = e^{\frac{i\pi}{n'}}$ 是本原 $(\mu=2n')$ 次单位根的基值与生成元}

$ $

注意 $\xi = \omega^{\frac{n}{n'}}$。因此，上述对 $M_Y(Y)$ 的求值，会将 $M_Y(Y)$ 在 $\frac{n'}{2}$ 个不同的本原 $(\mu=2n')$ 次单位根处的求值结果重复输出 $\dfrac{n}{n'}$ 次。



\subsection{模数自举}
\label{subsec:ckks-bootstrapping}

\textbf{- 参考文献：} \href{https://eprint.iacr.org/2018/153.pdf}{Bootstrapping for Approximate Homomorphic Encryption}~\cite{ckks}

在标准的分层求值安排中，重缩放最终会耗尽可用模数链。自举从较小模数 $q_0$ 开始，将其提升到工作模数 $q_L$，在内部消耗若干层后，返回模数为某个 $q_l>q_0$ 的密文。输出不一定处于最大模数 $q_L$。

取 $(A,B)\in R_{q_0}^2$ 的中心化代表元，并假设 $[B+AS]_{q_0}=\Delta M+E$。其未取模相位为
\[
B+AS=\Delta M+E+q_0K,
\]
其中 $K\in R$。这是关于整数提升的恒等式，并不是已经完成取模的解密恒等式。自举近似消除 $q_0K$ 项，保留原有数值误差，并引入新的误差。

\subsubsection{总体思路}
\label{subsubsec:ckks-bootstrapping-high-level}

自举的第一步，是直接将密文 $(A, B)$ 的模数从 $q_0$ 改为 $q_L$。此时，执行带模 $q_L$ 约减的解密，将输出：

$\textsf{RLWE}^{-1}\textbf{(}\textsf{ct} = (A, B)\textbf{)} \bmod q_L$

$ = B + A\cdot S \bmod q_L$

$= \Delta M + E + q_0 K \bmod q_L$

这种解释要求 $\|B+AS\|_\infty<q_L/2$。对于中心化输入代表元，有 $\|B+AS\|_\infty\leq(q_0/2)(1+\|S\|_1)$，因此可以通过选择足够大的 $q_L$ 保证该条件。仅指出秘密系数是三元值，并不能给出与维数无关的界；$\|S\|_1$ 可以达到 $n$。


在 $\Delta M + E + q_0 K \bmod q_L$ 中，$q_0K$ 项不再被模 $q_0$ 约减消去，因此密文解密所得明文多项式的每个 $i$ 次项都会得到被扰动的系数 $\Delta m_i + e_i + q_0\cdot k_i \bmod q_L$，而不是 $\Delta m_i + e_i \bmod q_L$。所以，现在需要消去各系数中的多余项 $q_0\cdot k_i \bmod q_L$，提取出纯粹的明文系数 $\Delta m_i + e_i$。



\begin{figure}[h!]
    \centering
  \includegraphics[width=1.0\linewidth]{figures/modulo-reduction-sine.png}
  \caption{正弦函数 $f(x) = \dfrac{q_0}{2\pi}\cdot \sin \left(\dfrac{2\pi x}{q_0}\right)$ 满足 $f(\Delta m_i + e_i + q_0k_i) \approx \Delta m_i + e_i$（假设 $\Delta m_i + e_i \ll q_0$）\href{https://eprint.iacr.org/2018/153.pdf}{（来源）}}
  \label{fig:modulo-reduction-sine}
\end{figure}

该函数以 $q_0$ 为周期。在 $x=kq_0$ 附近，它近似的是局部中心化余数 $x-kq_0$，而不是全局函数 $x$。它在 $kq_0$ 处的导数为 $1$。

$ $

结合这两个性质，对于输入 $x = \Delta m_i + e_i + q_0k_i$，有

\[
\begin{aligned}
f(\Delta m_i+e_i+q_0k_i)
&=\frac{q_0}{2\pi}\sin\left(\frac{2\pi(\Delta m_i+e_i+q_0k_i)}{q_0}\right)\\
&=\frac{q_0}{2\pi}\sin\left(\frac{2\pi(\Delta m_i+e_i)}{q_0}\right)\\
&\approx\Delta m_i+e_i.
\end{aligned}
\]

$ $

，前提是 $|\Delta m_i+e_i|/q_0$ 对于目标精度足够小。这是额外的自举参数条件，并非普通 CKKS 解密正确性的直接结果。令 $r_i=\Delta m_i+e_i$，正弦本身的近似误差满足
\[
\left|\frac{q_0}{2\pi}\sin(2\pi r_i/q_0)-r_i\right|
\leq\frac{2\pi^2}{3q_0^2}|r_i|^3.
\]
多项式近似和同态求值还会产生额外误差。由于密文乘法按槽位进行，算法先施加线性的系数到槽位变换，在槽位中求值非线性近似，再施加逆线性变换，流程见 \autoref{tab:ckks-bootstrapping-procedure}。

\begin{table}
\begin{tabular}{|lp{0.78\linewidth}|}
\hline
1 & \textsf{ModRaise}：将密文系数提升到 $q_L$，得到相位 $\Delta M+E+q_0K$。\\
2 & \textsf{CoeffToSlot}：同态施加逆嵌入，将 $n$ 个实相位系数放入两个密文。\\
3 & \textsf{EvalExp}：在明确的输入区间及误差预算下，求值周期模约减函数的非线性多项式近似。\\
4 & \textsf{SlotToCoeff}：将近似系数重新打包为一个密文，处于剩余模数 $q_l>q_0$，并包含全部累积误差。\\
\hline
\end{tabular}
\caption{CKKS 自举流程的高层描述}
\label{tab:ckks-bootstrapping-procedure}
\end{table}

\subsubsection{总体思路的数学展开}
\label{subsubsec:ckks-bootstrapping-high-level-correctness}

下述恒等式描述密文所表示的明文值，并不是求值过程中暴露秘密密钥或执行解密的操作。令 $T=\Delta M+E+q_0K$ 为提升后的相位，并仅将 $\mathsf{Enc}(T)$ 用作表示该相位的密文的简写。

\para{\textsf{\textbf{\underline{ModRaise}：}}} 将 $(A,B)$ 的系数提升到工作模数 $q_L$。在上述界成立时，提升后的密文具有原始相位 $T$。它可以直接作为同态求值的输入，无需在明文中计算 $AS$。

\para{\textsf{\textbf{\underline{CoeffToSlot}：}}} 令 $\sigma$ 为典范嵌入，$\pi$ 选出 $n/2$ 个独立槽位。实线性的逆嵌入将槽位 $\pi\sigma(T)$ 映为系数向量 $(T_i)_{i=0}^{n-1}$。其同态实现使用旋转、共轭、加法以及与编码矩阵常量的乘法。将 $n$ 个实系数存入两个密文，每个密文保存 $n/2$ 个系数，并记录这些变换的尺度与误差。

\para{\textsf{\textbf{\underline{EvalExp}：}}} 在所需输入区间上，同态求值近似 $f(x)=q_0\sin(2\pi x/q_0)/(2\pi)$ 的多项式电路 $\widehat f$。这是非线性计算，不是线性变换，也不是与单个编码多项式相乘。对于每个系数，
\[
\widehat f(T_i)=\Delta m_i+e_i+e_{b,i},
\]
其中 $e_{b,i}$ 包含正弦误差、多项式近似误差及全部求值误差。

\para{\textsf{\textbf{\underline{SlotToCoeff}：}}} 对所得两个密文施加逆线性打包操作。在所记录的输出尺度和模数 $q_l$ 下，最终相位表示为 $\Delta M+E+E_b$；$E_b$ 还包含最后一次变换的误差。

\para{\textsf{\textbf{\underline{无需秘密密钥的自举}：}}} 实际电路始终操作提升后的密文，不执行解密再加密；对于带噪声的随机化加密，也不存在 $\mathsf{Enc}(\mathsf{Dec}(\textsf{ct}))=\textsf{ct}$ 这样的恒等式。公开求值密钥支持旋转、共轭和重线性化。输出近似表示原始解码值，具有较大的可用模数，同时带有额外误差。

\subsubsection{细节：\textsf{CoeffToSlot}}
\label{subsubsec:ckks-bootstrapping-coefftoslot-details}

令 $\zeta_j=\omega^{J(j)}$，定义 $n/2\times n$ 求值矩阵 $U$，其中 $U_{j,k}=\zeta_j^k$。若 $\vec t$ 是 $T$ 的实系数向量，则其独立槽位为 $\vec z=U\vec t$。由单位根的正交性，
\[
\vec t=\frac{1}{n}(\overline U^{\,T}\vec z+U^T\overline{\vec z}).
\]
按系数索引的前半部与后半部，将 $U=(U_0\;U_1)$ 分为两个方阵块，则
\[
\vec t_0=\frac{1}{n}(\overline{U_0}^{\,T}\vec z+U_0^T\overline{\vec z}),\qquad
\vec t_1=\frac{1}{n}(\overline{U_1}^{\,T}\vec z+U_1^T\overline{\vec z}).
\]
同态施加这两个实线性映射，可得到槽位分别保存两半系数的密文，并产生编码和求值误差。共轭按 \autoref{subsec:ckks-conjugation} 计算；矩阵与向量乘法采用 \autoref{subsec:bfv-matrix-multiplication} 中的编码对角线和旋转方法。这些运算由密文与明文乘法及加法组成，层级消耗取决于选定的尺度安排。这里的记号不表示在明文中直接将矩阵乘以密文多项式的系数列表。

\subsubsection{细节：\textsf{EvalExp}}
\label{subsubsec:ckks-bootstrapping-evalexp-details}

利用正弦函数 $f(x) = \dfrac{q_0}{2\pi}\cdot \sin\left(\dfrac{2\pi x}{q_0}\right)$，通过计算 $f(\Delta m_i + e_i + q_0k_i) \approx \Delta m_i + e_i$，从 $\Delta m_i + e_i + q_0k_i$ 中近似消去 $q_0k_i$。当相对于 $q_0$ 而言，$\Delta m_i + e_i$ 非常接近 $x=0$（即 $\Delta m_i + e_i \ll q_0$）时，该近似有效。不过，对 $q_0k_i$ 的消去仍然是近似的（即 $\approx \Delta m_i + e_i$），因为在 $x=0$ 附近，$f(x)$ 为 $y \approx x$，而并非精确的 $y = x$。

这里的一个问题是，需要以 \textsf{ct\textsubscript{s1}} 与 \textsf{ct\textsubscript{s2}} 为输入（即 $f(\textsf{ct\textsubscript{s1}})$ 与 $f(\textsf{ct\textsubscript{s2}})$），同态求值 $f(x)$；但 FHE 仅支持 $(+, \cdot)$ 运算，而正弦曲线无法仅用 $(+, \cdot)$ 精确表示。因此，使用泰勒级数（\autoref{sec:taylor-series}）来近似正弦函数 $f(x)$：

$f(x) = f(a) + \dfrac{f'(a)}{1!}(x-a) + \dfrac{f''(a)}{2!}(x-a)^2 + \dfrac{f'''(a)}{3!}(x-a)^3 + \cdots = \sum\limits_{d=0}^{\infty}\dfrac{f^{(d)}(a)}{d!}(x-a)^d$

$ $

如果在 $x = 0$ 附近近似 $f(x)$，则近似多项式为：

$f(x) = \dfrac{q_0}{2\pi}\cdot \sin\left(\dfrac{2\pi}{q_0}\cdot 0\right) + \dfrac{q_0}{2\pi}\cdot\dfrac{2\pi}{q_0}\cdot\dfrac{\cos\left(\dfrac{2\pi}{q_0} \cdot 0\right)}{1!}\cdot x + \dfrac{q_0}{2\pi}\cdot \left(\dfrac{2\pi}{q_0}\right)^{2}\cdot\dfrac{-\sin\left(\dfrac{2\pi}{q_0} \cdot 0\right)}{2!}\cdot x^2 + \cdots$

$ = \dfrac{q_0}{2\pi}\cdot \sum\limits_{j=0}^{\infty}\left(\dfrac{(-1)^j}{(2j + 1)!}\cdot \left(\dfrac{2\pi x}{q_0}\right)^{2j + 1}\right)$

$ \approx \dfrac{q_0}{2\pi}\cdot \sum\limits_{j=0}^{h}\left(\dfrac{(-1)^j}{(2j + 1)!}\cdot \left(\dfrac{2\pi x}{q_0}\right)^{2j + 1}\right) = \hat f(x)$

$ $

其中，$\hat{f}(x)$ 是 $(2h + 1)$ 次多项式。

对于 $A,B$ 的中心化代表元及小秘密 $S$，整数相位满足 $B+AS=\Delta M+E+q_0K$。其系数范数至多为 $(q_0/2)(1+\|S\|_1)$，而对 $K$ 的界还包含 $\|\Delta M+E\|_\infty/q_0$。应选取 $k_{\max}$ 来约束整个归一化输入 $|x|/q_0$，包括残余消息和误差，而不只是 $|k_i|$。稀疏秘密可缩小该区间，但必须按其实际分布进行安全性分析。

归一化近似区间由 $\theta=2\pi x/q_0$ 决定，因此当 $|x|/q_0$ 固定时，单独增大 $q_0$ 不会扩大该区间。对于 $d$ 次多项式，霍纳求值使用 $d$ 次乘法，标准的并行幂构造可在多项式工作量下实现对数乘法深度。计算成本一般并不随次数指数增长。缩小输入后反复平方，是利用指数函数的特殊结构来减少工作量和深度：

\[
\begin{cases}
e^{i\cdot \theta} = \cos \theta + i \cdot \sin \theta\\
(e^{i \cdot \theta})^2 = e^{i\cdot 2\theta}\\
\end{cases}
\]

令 $\theta=2\pi x/q_0$。先在较小区间 $|\theta|/2^r\leq2\pi k_{\max}/2^r$ 上近似 $e^{i\theta/2^r}$，再将结果平方 $r$ 次，以近似 $e^{i\theta}$。所需泰勒次数取决于该归一化区间及目标最终误差。平方会放大初始近似误差并增加算术误差，因此初始容差必须考虑全部 $r$ 次平方。

在 $x = 0$ 附近使用 $d_0$ 次泰勒级数，可将 $e^{\frac{2\pi i x}{2^r q_0}}$ 近似为：

$f_e(x) = e^{\frac{2\pi i x}{2^r q_0}} \approx \sum\limits_{d=0}^{d_0}\dfrac{1}{d!}\left(\dfrac{2\pi i x}{2^rq_0}\right)^d = \hat{f}_e(x)$

$ $

然后，将 $\hat{f}_e$ 连续平方 $r$ 次，得到：

$(\hat{f}_e(x))^{2^r} \approx (f_e(x))^{2^r} = e^{i\frac{2\pi x}{q_0}} = e^{i\theta}$

$ $

接着，依据欧拉公式 $e^{i\cdot \theta} = \cos \theta + i \cdot \sin \theta$，可以推导如下关系：

$ $

$\overline{e^{i\cdot \theta}} = \cos \theta + \overline{i \cdot \sin \theta} $

$ e^{-i\cdot \theta} = \cos \theta - i \cdot \sin \theta$

$ e^{i\cdot \theta} - e^{-i\cdot \theta} = (\cos \theta + i \cdot \sin \theta) - (\cos \theta - i \cdot \sin \theta) = 2i\sin \theta$

$\sin \theta = \dfrac{-i}{2} \cdot (e^{i\cdot \theta} - e^{-i\cdot \theta})$

$\dfrac{q_0}{2\pi}\cdot \sin \theta = \dfrac{q_0}{2\pi}\cdot \dfrac{-i}{2} \cdot (e^{i\cdot \theta} - e^{-i\cdot \theta})$

$ $

代入 $\theta = \dfrac{2\pi x}{q_0}$，最终得到：

$\dfrac{q_0}{2\pi}\cdot \sin \left(\dfrac{2\pi x}{q_0}\right) = \dfrac{q_0}{2\pi}\cdot \dfrac{-i}{2} \cdot (e^{i\cdot \frac{2\pi x}{q_0}} - e^{-i\cdot \frac{2\pi x}{q_0}})$

$ $

利用上述最终关系，\textsf{EvalExp} 步骤按如下方式同态求值 $\dfrac{q_0}{2\pi}\cdot \sin \left(\dfrac{2\pi x}{q_0}\right)$ 的近似值，其中 $x = \Delta m_i + e_i + q_0k_i$：

\begin{enumerate}
\item 同态求值指数函数近似，得到 $\widehat z\approx e^{2\pi ix/q_0}$，并给出误差界。

\item 施加同态共轭（\autoref{subsec:ckks-conjugation}），得到 $\overline{\widehat z}$ 的近似，并计入密钥切换误差。

\item 将两者相减，再乘以编码常数 $-iq_0/(4\pi)$，以近似 $\frac{q_0}{2\pi}\sin(2\pi x/q_0)$。因子 $q_0/(2\pi)$ 不可省略；若仅乘以 $-i/2$，得到的将是未缩放的正弦。
\end{enumerate}

$ $

\textsf{EvalExp} 的结果是两个密文，其输入向量槽位存储了 $M_c$ 自举后的系数，即通过模约减 $q_0$ 将 $\Delta m_i + e_i + q_0k_i$ 转换为 $\Delta m_i + e_i + e_{bi}$。注意，$e_{bi}$ 是由以下三个因素产生的自举误差：（1）\textsf{CoeffToSlot}、\textsf{EvalExp} 和 \textsf{SlotToCoeff} 步骤固有的同态 $(+, \cdot)$ 运算噪声；（2）\textsf{EvalExp} 步骤中用泰勒多项式近似指数函数 $e^{i\theta}$ 产生的误差；（3）\textsf{EvalExp} 步骤的正弦曲线误差，因为曲线在 $x=0$ 附近并非精确等于 $y=x$，而只是 $y \approx x$。

注意，由于 \textsf{CoeffToSlot} 的输出被分为两个密文（\textsf{ct\textsubscript{s1}} 与 \textsf{ct\textsubscript{s2}}），\textsf{EvalExp} 的输出也包含两个密文：\textsf{ct\textsubscript{b1}} 与 \textsf{ct\textsubscript{b2}}。\textsf{ct\textsubscript{b1}} 的输入向量槽位存储 $(\Delta m_i + e_i + e_{bi})_{i=0}^{\frac{n}{2} - 1}$，而 \textsf{ct\textsubscript{b2}} 的输入向量槽位存储 $(\Delta m_i + e_i + e_{bi})_{i=\frac{n}{2}}^{n - 1}$。


\subsubsection{细节：\textsf{SlotToCoeff}}
\label{subsubsec:ckks-bootstrapping-slottocoeff-details}

采用 \autoref{subsubsec:ckks-bootstrapping-coefftoslot-details} 中相同的矩阵块 $U_0,U_1$，令两个 EvalExp 输出向量为 $\vec b_0,\vec b_1$，分别近似保存前半部与后半部系数。计算
\[
\vec z_{\mathrm{out}}=U_0\vec b_0+U_1\vec b_1.
\]
所得单个密文具有 $n/2$ 个复数槽位，表示的多项式具有由这两半拼接而成的 $n$ 个实系数。其隐含共轭槽位为 $\overline{\vec z_{\mathrm{out}}}$，并不是对同一系数多项式的另一次加密。同态变换会引入编码、密钥切换和尺度调整误差。因此，输出原始相位在模 $q_l$ 下为 $\Delta_{\mathrm{out}}M+E_{\mathrm{old,out}}+E_b$，其中全部尺度和误差都应沿具体实现传播。若计算安排恢复输入尺度，可以将其简写为 $\Delta M+E+E_b$。

\subsubsection{通过稀疏打包密文降低自举开销}
\label{subsubsec:ckks-bootstrapping-time-reduction}

假设 $2\leq n'<n$ 是一个真偶因子，令 $d=n/n'$。重复 $n'/2$ 个独立槽位会得到关于 $Y=X^d$ 的消息多项式，它有 $n'$ 个可能非零的系数。若起初只有一个非零槽位块、其余槽位均补零，可用 $\log_2 d$ 轮倍增来复制，每轮移位为 $-(n'/2)2^j$，其中 $0\leq j<\log_2d$。

模数提升后，仅有消息稀疏性还不够：相位 $\Delta M(Y)+E(X)+q_0K(X)$ 中的 $E$ 和 $K$ 一般并不稀疏。因此，原始 CKKS 自举论文第 5.2 节在 ModRaise 之后、CoeffToSlot 之前加入部分求和：
\[
\textsf{ct}\leftarrow\textsf{ct}+\textsf{Rotate}(\textsf{ct},(n'/2)2^j),
\qquad 0\leq j<\log_2d.
\]
这种类似迹映射的操作会消去次数不被 $d$ 整除的系数位置，并将保留系数乘以 $d$，同时产生求值误差。必须在记录的尺度和模约减周期中计入其归一化。若省略这一步，就不能据此丢弃提升后的其他系数。

较小的典范嵌入使用 $Y^{n'}+1$ 的根，具有 $n'$ 个实系数坐标，而不是 $n/n'$ 个。二次幂的真因子满足 $n'\leq n/2$，所以这些系数值可以装入一个密文。相应的较小系数与槽位变换以及模约减求值能够降低成本。实际收益取决于变换算法、旋转和归一化，并非所有参数下都恰好减半。

\subsubsection{总结}
\label{subsubsec:ckks-bootstrapping-summary}

\begin{tcolorbox}[title={\textbf{\tboxlabel{\ref*{subsubsec:ckks-bootstrapping-summary}} CKKS 自举}}]
\begin{enumerate}
\item 输入 $(A,B)\in R_{q_0}^2$，其中心化相位为 $\Delta M+E$，并满足所需输入区间界。
\item 将密文系数提升到 $q_L$。在模 $q_L$ 下满足不回绕界时，原始相位为 $\Delta M+E+q_0K$。RNS 实现需要计算相应的基扩展。
\item 执行 CoeffToSlot：完全打包输入使用两个密文，稀疏输入则采用上述正确归一化的流程。
\item 求值模约减的非线性多项式近似，同时预算近似误差与同态求值误差。
\item 执行 SlotToCoeff。输出处于剩余模数 $q_l>q_0$，并附有所记录的尺度和累积误差。采用等尺度简写时，其相位为 $\Delta M+E+E_b$，而不是精确的 $\Delta M+E$。
\end{enumerate}
\textbf{局限：} 自举不能恢复输入中已经损失的精度。后续数值计算仍须满足消息、误差和模数预算。
\end{tcolorbox}

\para{BFV 与 CKKS 自举的比较：} 精确方案的自举可通过求值精确解密函数消除先前的加密噪声，前提是满足其正确性条件。CKKS 会保留先前的数值近似误差，并增加新误差。无论固定模数的 BFV，还是一组固定 CKKS 参数，都不能脱离相应的噪声与精度条件支持任意无界工作。这里的 CKKS 流程直接使用提升后的密文；不需要单独的加密秘密内积，并不是因为所有中间尺度始终固定。

\subsubsection{降低自举噪声}
\label{subsubsec:ckks-bootstrapping-noise-reduction}

如 \autoref{subsubsec:ckks-bootstrapping-evalexp-details} 所述，自举过程会产生三类噪声：

%\renewcommand{\labelenumi}{\underline{Type-\arabic{enumi} Noise:}}
\begin{itemize}
\item \underline{第 1 类噪声}：\textsf{CoeffToSlot}、\textsf{EvalExp} 和 \textsf{SlotToCoeff} 步骤固有的同态 $(+, \cdot)$ 运算噪声
\item \underline{第 2 类噪声}：\textsf{EvalExp} 步骤中指数函数 $e^{i\theta}$ 的近似误差
\item \underline{第 3 类噪声}：\textsf{EvalExp} 步骤中的正弦曲线误差，即曲线在 $x=0$ 附近并非精确等于 $y=x$，而只是 $y \approx x$
\end{itemize}


$ $

对于选定的输入区间和参数，三类误差都必须有明确界。有限次多项式求值通常仍有非零近似误差，可以通过提高次数和计算精度来减小。改进模约减近似能减少正弦本身的误差，但不能省略对算术误差和近似误差的共同分析。

\begin{figure}[h!]
    \centering
  \includegraphics[width=1.0\linewidth]{figures/arc-sine.png}
  \caption{用于减小近似误差的反正弦曲线 \href{https://www.google.com/search?sca_esv=c744cb070de47b7e&sxsrf=ADLYWIKKaF93fcAFvJC4ehJso1E5qFTaow:1732546684173&q=arcsin(sin(x))&source=lnms&fbs=AEQNm0DmKhoYsBCHazhZSCWuALW8l8eUs1i3TeMYPF4tXSfZ95GzcfXnm5XYTvJV_9Qreh2py964ICpZJthXkELijctC8pFBYULoa3-fvQmwK0VJF0ntzsbN_W2CCJL9N57SWFNwWI58jCKaBJSdPgkprHQVK8H1PYOYWXMHTCCV-rDbC44rR6ANM870jZCZRtTKwFWtMIe2&sa=X&ved=2ahUKEwjZy6qt3_eJAxXbhlYBHXD2IvEQ0pQJegQIDxAB&cshid=1732546705369143&biw=1280&bih=635&dpr=1.5}{（来源）}}
  \label{fig:arc-sine}
\end{figure}

$ $

\para{$\bm{\arcsin(\sin(x))}$ 近似（\href{https://dl.acm.org/doi/10.1007/978-3-030-77870-5_22}{EUROCRYPT 2021}~\cite{cryptoeprint:2020/552}）：} 精确复合函数仅在 $[-\pi/2,\pi/2]$ 等主值区间上等于 $x$，其周期延拓的斜率交替变化。在所需的各个周期整数倍附近，它可以表示目标局部余数。同态实现需在指定区间上近似各组成函数，例如采用极小极大多项式。必须说明其近似误差、区间限制和额外深度；不能从该恒等式推出固定消耗三层的一般结论。

$ $

\para{Meta-BTS（\href{https://dl.acm.org/doi/10.1145/3548606.3560696}{CCS 2022}~\cite{cryptoeprint:2022/1167}）：} Meta-BTS 通过反复纠正前一次自举的误差来提高精度。其纠正思路可概括为：
\begin{enumerate}
\item 计算 $c_1=\mathsf{BTS}(c)$，表示 $m+\eta$。
\item 对齐模数与尺度，再构造 $d=\mathsf{ModDrop}(c_1)-c$，表示 $\eta$。应以自举后的密文减去原密文，才能使随后相减的符号正确。若不同时调整尺度，不能用重缩放替代 ModDrop。
\item 将小误差归一化到基础自举支持的输入区间，对其自举，再撤销归一化，得到表示 $\eta+\eta'$ 的 $\widehat d$。
\item 对齐参数，计算 $c_1-\widehat d$，表示 $m-\eta'$。
\end{enumerate}
归一化、精度及剩余层级是 Meta-BTS 定量保证的必要组成部分。上述概要并不声称，仅重复未经尺度调整的自举就一定能提高精度。

\para{噪声处理的局限：} 在满足算法的归一化与精度假设时，误差纠正可以减少自举新增误差。对固定的参数和实现，剩余算术误差与近似误差会限制可达到的精度。上述概要不能推出与输入无关的指数收敛或无限精度提升，也不能恢复原密文中已经损失的信息。





\clearpage

\section{BGV 方案}
\label{sec:bgv}

与 BFV 类似，BGV 方案面向整数的同态加法和乘法。BGV 保证精确的加密与解密，而 CKKS 采用近似计算。从这一角度看，BGV 与 BFV 相似。不过，两者的主要区别在于：BFV 将明文值放在最高有效位区域，将噪声放在低位（最低有效位）区域；BGV 则相反，将明文放在低位区域，将噪声放在高位区域。具体来说，BFV 用 $\Delta$ 缩放明文多项式，而 BGV 用 $\Delta$ 缩放噪声多项式。因此，两种方案采用略有差异的方式在密文中存储和管理明文与噪声。

BGV 在明文编码、密文与明文加法、密文与密文加法、密文与明文乘法以及输入向量旋转方面，采用与 BFV 几乎相同的策略。另一方面，由于 BGV 缩放的是噪声而不是明文，其加密与解密过程与 BFV 略有不同。此外，BFV 的密文乘法本身不设置次数层级，而 BGV 的密文与密文乘法采用分层机制，如 CKKS 一样将模数切换到较低层级，因此可执行的层数受到限制。BGV 的模数切换和自举也与 BFV 部分不同。



\begin{tcolorbox}[
    title = \textbf{所需背景知识},    % box title
    colback = white,    % light background; tweak to taste
    colframe = black,  % frame colour
    boxrule = 0.8pt,     % line thickness
    left = 1mm, right = 1mm, top = 1mm, bottom = 1mm % inner padding
]

\begin{itemize}
\item \autoref{sec:modulo}: \nameref{sec:modulo}
\item \autoref{sec:group}: \nameref{sec:group}
\item \autoref{sec:field}: \nameref{sec:field}
\item \autoref{sec:order}: \nameref{sec:order}
\item \autoref{sec:polynomial-ring}: \nameref{sec:polynomial-ring}
\item \autoref{sec:decomp}: \nameref{sec:decomp}
\item \autoref{sec:roots}: \nameref{sec:roots}
\item \autoref{sec:cyclotomic}: \nameref{sec:cyclotomic}
\item \autoref{sec:cyclotomic-polynomial-integer-ring}: \nameref{sec:cyclotomic-polynomial-integer-ring}
\item \autoref{sec:matrix}: \nameref{sec:matrix}
\item \autoref{sec:euler}: \nameref{sec:euler}
\item \autoref{sec:modulus-rescaling}: \nameref{sec:modulus-rescaling}
\item \autoref{sec:chinese-remainder}: \nameref{sec:chinese-remainder}
\item \autoref{sec:polynomial-interpolation}: \nameref{sec:polynomial-interpolation}
\item \autoref{sec:ntt}: \nameref{sec:ntt}
\item \autoref{sec:lattice}: \nameref{sec:lattice}
\item \autoref{sec:rlwe}: \nameref{sec:rlwe}
\item \autoref{sec:glwe}: \nameref{sec:glwe}
\item \autoref{sec:glwe-add-cipher}: \nameref{sec:glwe-add-cipher}
\item \autoref{sec:glwe-add-plain}: \nameref{sec:glwe-add-plain}
\item \autoref{sec:glwe-mult-plain}: \nameref{sec:glwe-mult-plain}
\item \autoref{subsec:modulus-switch-rlwe}: \nameref{subsec:modulus-switch-rlwe}
\item \autoref{sec:glwe-key-switching}: \nameref{sec:glwe-key-switching}
\item \autoref{sec:bfv}: \nameref{sec:bfv}
\item \autoref{sec:ckks}: \nameref{sec:ckks}
\end{itemize}
\end{tcolorbox}

\clearpage

\subsection{编码与解码}
\label{subsec:bgv-encoding-decoding}

BGV 的明文编码与 \autoref{subsubsec:bfv-encoding-summary} 中总结~\ref{subsubsec:bfv-encoding-summary} 所述的 BFV 编码几乎相同，唯一区别在于：BGV 不像 BFV 那样将尺度因子 $\Delta = \left\lfloor\dfrac{q}{t}\right\rfloor$ 施加到明文多项式 $M(X)$ 上，而是在每次加密产生新密文时，将自身的尺度因子 $\Delta = t$ 施加到噪声多项式 $E(X)$ 上（详见 \autoref{subsec:bgv-enc-dec}）。

下面给出 BGV 的编码与解码过程。

\begin{tcolorbox}[title={\textbf{\tboxlabel{\ref*{subsec:bgv-encoding-decoding}} BGV 的编码与解码}}]



\textbf{\underline{输入}：} 一个具有 $n$ 个槽位、各元素为模 $t$ 整数的向量 $\vec{v} = (v_0, v_1, \cdots, v_{n-1}) \in \mathbb{Z}_t^n$

$ $

\textbf{\underline{编码}：}

通过变换 $\vec{m} = \dfrac{\hathat W \cdot I_n^R \cdot \vec{v}}{n}$，将 $\vec{v} \in \mathbb{Z}_t^n$ 转换为 $\vec{m} \in \mathbb{Z}_t^n$。

其中 $\hathat W$ 是 $n$ 维向量空间的一组基，构造如下：

{\footnotesize{$\hathat W = \begin{bmatrix}
1 & 1 & \cdots & 1 & 1 & 1 & \cdots & 1\\
(\omega^{J(\frac{n}{2} - 1)}) & (\omega^{J(\frac{n}{2} - 2)}) & \cdots & (\omega^{J(0)}) & (\omega^{J_*(\frac{n}{2} - 1)}) & (\omega^{J_*(\frac{n}{2} - 2)}) & \cdots & (\omega^{J_*(0)})\\
(\omega^{J(\frac{n}{2} - 1)})^2 & (\omega^{J(\frac{n}{2} - 2)})^2 & \cdots & (\omega^{J(0)})^2 & (\omega^{J_*(\frac{n}{2} - 1)})^2 & (\omega^{J_*(\frac{n}{2} - 2)})^2 & \cdots & (\omega^{J_*(0)})^2 \\
\vdots & \vdots & \ddots & \vdots & \vdots & \ddots & \vdots & \vdots \\
(\omega^{J(\frac{n}{2} - 1)})^{n-1} & (\omega^{J(\frac{n}{2} - 2)})^{n-1} & \cdots & (\omega^{J(0)})^{n-1} & (\omega^{J_*(\frac{n}{2} - 1)})^{n-1} & (\omega^{J_*(\frac{n}{2} - 2)})^{n-1} & \vdots  & (\omega^{J_*(0)})^{n-1}
\end{bmatrix}$}}

$ $

其中 $\omega$ 是模 $t$ 的本原 $2n$ 次单位根（因此 $t \equiv 1 \bmod 2n$）。这意味着 $\omega = g^{\frac{t - 1}{2n}} \bmod t$（$g$ 是 $\mathbb{Z}_t^{\times}$ 的生成元，见 \autoref{subsubsec:poly-vector-transformation-modulus}）。最终输出为 $M = \sum\limits_{i=0}^{n-1} m_iX^i \text{ } \in \mathbb{Z}_t[X] / (X^n + 1)$，
在后续的加密和解密过程中，也可以将其视为

$M = \sum\limits_{i=0}^{n-1} m_iX^i \text{ } \in \mathbb{Z}_q[X] / (X^n + 1)$，因为新编码得到的系数 $m_i$ 保证小于任何满足 $q = \{q_0, q_1, \cdots, q_L\}$ 的 $q$。

$ $

\textbf{\underline{解码}：} 对明文多项式 $M = \sum\limits_{i=0}^{n-1} m_iX^i$，
计算 $\vec{v} = \hathat W^* \cdot \vec{m}$，其中

$\hathat{W}^* = \begin{bmatrix}
1 & (\omega^{J(0)}) & (\omega^{J(0)})^2 & \cdots & (\omega^{J(0)})^{n-1}\\
1 & (\omega^{J(1)}) & (\omega^{J(1)})^2 & \cdots & (\omega^{J(1)})^{n-1}\\
1 & (\omega^{J(2)}) & (\omega^{J(2)})^2 & \cdots & (\omega^{J(2)})^{n-1}\\
\vdots & \vdots & \vdots & \ddots & \vdots \\
1 & (\omega^{J(\frac{n}{2}-1)}) & (\omega^{J(\frac{n}{2}-1)})^2 & \cdots & (\omega^{J(\frac{n}{2}-1)})^{n-1}\\
1 & (\omega^{J_*(0)}) & (\omega^{J_*(0)})^2 & \cdots & (\omega^{J_*(0)})^{n-1}\\
1 & (\omega^{J_*(1)}) & (\omega^{J_*(1)})^2 & \cdots & (\omega^{J_*(1)})^{n-1}\\
1 & (\omega^{J_*(2)}) & (\omega^{J_*(2)})^2 & \cdots & (\omega^{J_*(2)})^{n-1}\\
\vdots & \vdots & \vdots & \ddots & \vdots \\
1 & (\omega^{J_*(\frac{n}{2}-1)}) & (\omega^{J_*(\frac{n}{2}-1)})^2 & \cdots & (\omega^{J_*(\frac{n}{2}-1)})^{n-1}\\
\end{bmatrix}$

\end{tcolorbox}


\subsection{加密与解密}
\label{subsec:bgv-enc-dec}

BGV 的加密与解密方案与 BFV（\autoref{subsec:bfv-enc-dec} 中的总结~\ref*{subsec:bfv-enc-dec}）非常相似，但存在一个小区别：BFV 用 $\Delta$ 缩放明文多项式 $M(X)$，BGV 则用 $\Delta$ 缩放噪声多项式 $E(X)$。在 BFV 中，每个编码后的明文多项式 $M(X)$ 都乘以 $ \Delta = \left\lfloor \dfrac{q}{t} \right\rfloor$；这会将每个明文系数移到最高有效位区域，而将噪声留在最低有效位区域。另一方面，BGV 不缩放明文多项式 $M(X)$，而是将每次新生成的噪声 $E(X)$ 乘以 $\Delta = t$，得到 $\Delta E(X)$。这种缩放策略通过乘以 $\Delta = t$ 将噪声（即 $e_i$）移到最高有效位区域，同时将明文值（即 $m_i$），也就是 $M(X)$ 的各个系数，保留在最低有效位区域。

此外，BGV 的密文模数 $q$ 与 CKKS 一样采用分层结构：$q \in \{q_0, q_1, \cdots, q_L\}$，其中各层满足 $q_l = \prod\limits_{i=0}^l w_i$（各个 $w_i$ 是 CRT 模数）。

BGV 的加密与解密过程如下：


\begin{tcolorbox}[title={\textbf{\tboxlabel{\ref*{subsec:bgv-enc-dec}} BGV 加密与解密}}]

\textbf{\underline{初始设置}：}

\begin{itemize}
\item 明文模数 $t = p$（为素数）。

\item 密文模数 $q$ 与 CKKS 一样采用分层结构：$q \in \{q_0, q_1, \cdots, q_L\}$，其中各层满足 $q_l = \prod\limits_{i=0}^l w_i$（各个 $w_i$ 是 CRT 模数），并且满足 $q_l \equiv 1 \bmod t$（将在 \autoref{subsec:bgv-modulus-switch} 解释）。

\item 噪声尺度因子 $\Delta=t$。
\item 私钥 $S \xleftarrow{\$} \mathcal{R}_{\langle n, \textit{tern} \rangle}$。多项式 $S$ 的系数为三元值（即 $\{-1, 0, 1\}$）。

\end{itemize}

\par\noindent\rule{\textwidth}{0.4pt}

\textbf{\underline{加密输入}：} $M \in \mathcal{R}_{\langle n, q \rangle}$、$A \xleftarrow{\$} \mathcal{R}_{\langle n, q \rangle}$、$E \xleftarrow{\chi_\sigma} \mathcal{R}_{\langle n, q \rangle}$


\begin{enumerate}
%\item Scale up $M \rightarrow \Delta M \text { } \in \mathcal{R}_{\langle n, q\rangle}$

\item 计算 $B = -A \cdot S + M + \Delta E \text{ } \in \mathcal{R}_{\langle n,q \rangle}$

\item $\textsf{RLWE}_{S,\sigma}(M + \Delta E) = (A, B) \text{ } \in \mathcal{R}_{\langle n,q \rangle}^2$ 

\end{enumerate}

\par\noindent\rule{\textwidth}{0.4pt}

\textbf{\underline{解密输入}：} $\textsf{ct} = (A, B) \text{ } \in \mathcal{R}_{\langle n,q \rangle}^2$

\begin{enumerate}
\item $\textsf{RLWE}^{-1}_{S,\sigma}(\textsf{ct}) = B + A \cdot S  = M + \Delta E \pmod q$

\item 计算 $M + \Delta E \bmod t$，得到 $M$。\textcolor{red}{ $\rhd$ 即将 $M + \Delta E$ 对 $t$ 取模}

\end{enumerate}

$ $

最终输出为 $M(X) \in \mathbb{Z}_t[X] / (X^n + 1)$。

$ $

\textbf{\underline{正确解密的条件}：}

%\begin{enumerate}
%\item 
每个含缩放噪声与明文的系数 $\Delta e_i + m_i$，都不应超出密文当前乘法层级 $l$ 对应的密文模数 $q_l$ 的范围（即 $\Delta e_i + m_i < q_l$；采用有符号模表示时为 $|\Delta e_i + m_i| < \frac{q_l}{2}$）。

%\end{enumerate}

\end{tcolorbox}

在解密末尾恢复明文时，BFV 将明文和噪声向低位移动，从而通过舍入消除噪声；BGV 则计算 $\text{ mod } p$。由于 $E$ 的每个系数都是 $t$ 的倍数（即 $\Delta$），这一取模操作会消去累计的噪声。最后，低位区域仅保留明文多项式 $m_i$ 的各个系数，不再包含噪声 $e_i$。



\subsection{密文与密文加法}
\label{subsec:bgv-add-cipher}

BGV 的密文与密文加法与 BFV 完全相同（见 \autoref{subsec:bfv-add-cipher} 中的总结~\ref*{subsec:bfv-add-cipher}）。


\begin{tcolorbox}[title={\textbf{\tboxlabel{\ref*{subsec:bgv-add-cipher}} BGV 密文与密文加法}}]
$\textsf{RLWE}_{S, \sigma}(M^{\langle 1 \rangle} + \Delta E^{\langle 1 \rangle} ) + \textsf{RLWE}_{S, \sigma}(M^{\langle 2 \rangle} + \Delta E^{\langle 2 \rangle}) $

$ = ( A^{\langle 1 \rangle}, \text{ } B^{\langle 1 \rangle}) + (A^{\langle 2 \rangle}, \text{ } B^{\langle 2 \rangle}) $

$ = ( A^{\langle 1 \rangle} + A^{\langle 2 \rangle}, \text{ } B^{\langle 1 \rangle} + B^{\langle 2 \rangle} ) $

$= \textsf{RLWE}_{S, \sigma}\bm((M^{\langle 1 \rangle} + M^{\langle 2 \rangle})  + \Delta E^{\langle 1 \rangle} + \Delta E^{\langle 2 \rangle} \bm)$
\end{tcolorbox}

\subsection{密文与明文加法}
\label{subsec:bgv-add-plain}

BGV 的密文与明文加法与 BFV 几乎相同（见 \autoref{subsec:bfv-add-plain} 中的总结~\ref*{subsec:bfv-add-plain}）。区别在于，待加入的明文多项式 $\Lambda(X)$ 不乘以 $\Delta$ 进行放大，而是保持为 $\Lambda(X)$。

\begin{tcolorbox}[title={\textbf{\tboxlabel{\ref*{subsec:bgv-add-plain}} BGV 密文与明文加法}}]
$\textsf{RLWE}_{S, \sigma}(M + \Delta E) + \Lambda $

$=  (A, \text{ } B) + \Lambda$

$=  (A, \text{ } B + \Lambda)$

$= \textsf{RLWE}_{S, \sigma}\bm((M + \Lambda) + \Delta E \bm)$
\end{tcolorbox}


\subsection{密文与明文乘法}
\label{subsec:bgv-mult-plain}




BGV 的密文与明文乘法与 BFV 完全相同（见 \autoref{subsec:bfv-mult-plain} 中的总结~\ref*{subsec:bfv-mult-plain}）。


\begin{tcolorbox}[title={\textbf{\tboxlabel{\ref*{subsec:bgv-mult-plain}} BGV 密文与明文乘法}}]
$\textsf{RLWE}_{S, \sigma}(M + \Delta E) \cdot \Lambda$

$= (A, \text{ } B) \cdot \Lambda$

$= (A \cdot \Lambda, \text{ }  B \cdot \Lambda )$

$= \textsf{RLWE}_{S, \sigma}((M \cdot \Lambda) + \Delta E\cdot \Lambda )$
\end{tcolorbox}

注意，BGV 的密文与明文乘法不消耗乘法层级。

\subsection{\textsf{ModDrop}}
\label{subsec:bgv-moddrop}

BGV 的 \textsf{ModDrop} 与 CKKS 的 \textsf{ModDrop} 类似（见 \autoref{subsec:ckks-moddrop} 中的总结~\ref*{subsec:ckks-moddrop}）。回顾 CKKS 的密文解密关系如下：

$M + \Delta E = A \cdot S + B \bmod q_l$

$M + \Delta E = A \cdot S + B - K\cdot q_l$ \textcolor{red}{ $\rhd$ 其中 $K\cdot q_l$ 表示对 $q_l$ 取模}

BGV 的 \textsf{ModDrop} 通过将密文 $(A, B)$ 更新为新密文 $\bm(A' = A \bmod q_{l-1}$、$B' = B \bmod q_{l-1})$，按 $q_l \rightarrow q_{l-1}$ 降低模数。执行 \textsf{ModDrop} 后，密文模数按 $q_l \rightarrow q_{l-1}$ 降低，但仍保持如下解密关系：

$A' \cdot S + B' - K\cdot q_l $

$= (A \bmod q_{l-1}) \cdot S + (B \bmod q_{l-1}) - K\cdot q_l$

$= (A - K_A\cdot q_{l-1}) \cdot S + (B - K_B\cdot q_{l-1}) - K\cdot q_l$ 

$= A\cdot S + B - (K_A \cdot S + K_B + K\dfrac{q}{q_{l-1}})\cdot q_{l-1}$
\textcolor{red}{ $\rhd$ 其中 $\dfrac{q}{q_{l-1}} = w_l$（即 $q_L$ 的第 $l$ 个素数因子）}

$= A\cdot S + B - K'\cdot q_{l-1}$ \textcolor{red}{ $\rhd$ 其中 $K' = K_A\cdot S + K_B + K\dfrac{q}{q_{l-1}}$、$\dfrac{q}{q_{l-1}} = w_l$}

$= A\cdot S + B \bmod q_{l-1}$

$= M + \Delta E$ \textcolor{red}{ $\rhd$ 因为 $M + \Delta E < q_0 < q_{l-1}$}

$ $

如上所示，$(A', B') \bmod q_{l-1}$ 仍能解密为相同的含缩放误差明文 $M + \Delta E$。不过，由于密文模数与噪声之比减小，噪声预算（即允许的噪声阈值）也会减小。

$ $

BGV 的 \textsf{ModDrop} 可总结如下：



\begin{tcolorbox}[title={\textbf{\tboxlabel{\ref*{subsec:bgv-moddrop}} BGV 的 \textsf{ModDrop}}}]

给定处于第 $l$ 个乘法层级的 BGV 密文 $\textsf{RLWE}_{S, \sigma}(M + \Delta E) = (A, B) \bmod q_l$，\textsf{ModDrop} 操作为：

$(A', B') \bmod q_{l-1} = (A \bmod q_{l-1}, B \bmod q_{l-1})$

$ $

完成后，密文的乘法层级降低 1，噪声尺度因子 $\Delta$ 和明文均不变，但噪声预算（允许的噪声阈值）减小。

\end{tcolorbox}



\subsection{模数切换}
\label{subsec:bgv-modulus-switch}

\noindent \textbf{- 参考资料 1：}
\href{https://eprint.iacr.org/2020/1481}{HElib 的设计与实现：一个同态加密库}~\cite{bgv-modswitch1}


\noindent \textbf{- 参考资料 2：}
\href{https://eprint.iacr.org/2011/277.pdf}{无需自举的全同态加密}~\cite{bgv-modswitch2}


\noindent \textbf{- 参考资料 3：}
\href{https://eprint.iacr.org/2012/099.pdf}{AES 电路的同态求值}~\cite{bgv-modswitch3}

回顾模数切换的要求：将密文模数从 $q$ 改为 $\hat q$ 后，密文仍应解密得到相同明文 $M$。BGV 的模数切换与 RLWE 的模数切换（\autoref{subsec:modulus-switch-rlwe}）类似，但有一个额外要求，因为 BGV 将尺度因子 $\Delta$ 施加于噪声 $E$，而非明文 $M$。BFV 和 CKKS 的解密只需要通过舍入消除低位区域的噪声；而 BGV 的明文位于低位区域，解密必须通过模 $t$（即明文模数）的约减去除高位区域的噪声。更具体地说，BGV 从 $(A, B) \bmod q_{l}$ $\rightarrow$ $(\hat{A}, \hat{B}) \bmod \hat{q}$ 的模数切换应满足解密关系 $((\hat{A} \cdot S + \hat{B}) \bmod \hat{q}) \bmod t = M$。在 BGV 模数切换中，$\hat{q}$ 不必属于密文预设的乘法层级，$\hat{q}$ 只需满足 $\hat{q} < q_l$ 和 $\hat{q} \equiv 1 \bmod t$。BGV 的模数切换过程如下：

$ $

\begin{enumerate}


\item 输入密文为 $\textsf{ct} = (A, B) \bmod q_l$。按如下方式计算新多项式 $A'$ 和 $B'$：

$(A', B') = \left(\left\lceil\dfrac{\hat{q}}{q_l}\cdot A\right\rfloor, \left\lceil\dfrac{\hat{q}}{q_l}\cdot B\right\rfloor\right) \pmod{\hat{q}}$

$ $

并按如下方式计算舍入误差 $\epsilon_A, \epsilon_B$：

$\epsilon_A = \dfrac{\hat{q}}{q_l}\cdot A - A'$

$\epsilon_B = \dfrac{\hat{q}}{q_l}\cdot B - B'$


$ $

将其改写为：

$\hat{q} A = q_{l} A' + q_{l}\epsilon_A = q_{l} A' + \epsilon'_A$ \textcolor{red}{ $\rhd$ 记 $\epsilon'_A = q_{l}\epsilon_A$，其中 $\epsilon'_A \in \mathbb{Z}_{q_l}$}

$\hat{q} B = q_{l} B' + q_{l}\epsilon_B = q_{l} B' + \epsilon'_B$ \textcolor{red}{ $\rhd$ 记 $\epsilon'_B = q_{l}\epsilon_B$，其中 $\epsilon'_B \in \mathbb{Z}_{q_l}$}

$ $

\item 按如下方式计算新多项式 $H_A$ 和 $H_B$：

$H_A = q_l^{-1}\cdot\epsilon'_A \bmod t$

$H_B = q_l^{-1}\cdot\epsilon'_B \bmod t$

$ $

\item 定义最终完成模数切换的密文 $\hat{\textsf{ct}}$ 如下：

$\hat{\textsf{ct}} = (\hat{A}, \hat{B}) = (A' + H_A, \text{ } B' + H_B) \bmod \hat{q}$

$ $

注意，仅计算 $A' + H_A$ 和 $B' + H_B$ 所得的结果可能超出范围 $\mathbb{Z}_{\hat{q}}$，因为 $A', B' \in \mathbb{Z}_{\hat{q}}$ 且 $H_A, H_B \in \mathbb{Z}_t$。因此，需要将 $A' + H_A$ 和 $B' + H_B$ 对 $\hat{q}$ 取模，得到 $\hat{A} \in \mathbb{Z}_{\hat{q}}$ 和 $\hat{B} \in \mathbb{Z}_{\hat{q}}$。

$ $

\item 接下来验证 $\hat{\textsf{ct}}$ 是满足 BGV 解密关系的有效密文。首先，可推导 $\textsf{ct} = (A, B)$、$A' + H_A$ 和 $B' + H_B$ 之间的关系如下：

$\hat{q}\cdot \textsf{ct} \bmod t $

$= (\hat{q}A, \hat{q}B) \bmod t$

$= (q_{l} A' + \epsilon'_A, \text{ } q_{l} B' + \epsilon'_B) \bmod t$ \textcolor{red}{ $\rhd$ 代入步骤 1 的结果：$\hat{q} A = q_{l}A' + \epsilon'_A$，\text{ } $\hat{q} B = q_{l}B' + \epsilon'_B$}

$ = (q_{l} A' + q_lH_A, \text{ } q_{l} B' + q_lH_B) \bmod t$ 
\textcolor{red}{ $\rhd$ 代入步骤 2 的结果：$H_A = q_l^{-1}\cdot\epsilon'_A \bmod t$，\text{ } $H_B = q_l^{-1}\cdot\epsilon'_B \bmod t$}

$ $

$ = q_l\cdot(A' + H_A, \text{ } B' + H_B) \bmod t$

%$ = q_l \cdot \hat{\textsf{ct}} \bmod t$

$ $

因此，$\hat{q}\cdot\textsf{ct} = q_l\cdot(A' + H_A, \text{ } B' + H_B) \bmod t$。在 BGV 中，$q_l \equiv q_2 \equiv \cdots \equiv \hat{q} \equiv \cdots q_L \equiv 1 \bmod t$，于是有：

$A \equiv A' + H_A \bmod t$

$B \equiv B' + H_B \bmod t$




$ $

\item 根据 $\textsf{ct}$ 的解密关系，可以推导 $\hat{\textsf{ct}}$ 的解密关系如下：

$M = (A\cdot S + B \bmod q_l) \bmod t$ \textcolor{red}{ $\rhd$ $\textsf{ct} = (A, B) \bmod q_l$ 的 BGV 解密关系}

$= (A\cdot S + B - K\cdot q_l) \bmod t$ \textcolor{red}{ $\rhd$ 其中 $K\cdot q_l$ 表示模 $q_l$ 约减}

$= ((A' + H_A)\cdot S + B' + H_B - K\cdot q_{l}) \bmod t$ \textcolor{red}{ $\rhd$ 代入步骤 4 的结果：$A \equiv A' + H_A \bmod t$，\text{ } $B \equiv B' + H_B \bmod t$}

$= ((A' + H_A)\cdot S + B' + H_B - K\cdot \hat{q}) \bmod t$ \textcolor{red}{ $\rhd$ 因为 BGV 中 $q_0\equiv q_1\equiv \cdots q_L \equiv 1 \bmod t$，且所选 $\hat{q}$ 满足 $\hat{q} \equiv 1 \bmod t$}

$ $

\item 现在，只要证明 $(A' + H_A)\cdot S + B' + H_B - K\cdot \hat{q} = (A' + H_A)\cdot S + B' + H_B \bmod \hat{q}$（即 $K\cdot\hat{q}$ 将 $(A' + H_A)\cdot S + B' + H_B$ 对 $\hat{q}$ 取模），就足以得到以下结论：

$(\hat{A}\cdot S + \hat{B} \bmod \hat{q}) \bmod t$

$= ((A' + H_A)\cdot S + B' + H_B \bmod \hat{q}) \bmod t$

$= M$ \textcolor{red}{ $\rhd$ 即 $(\hat{A}, \hat{B}) = (A' + H_A, B' + H_B) \bmod \hat{q}$ 是解密得到 $M$ 的有效密文}

$ $
\begin{comment}
\item 以下关系也成立：

$(A' + H_A)\cdot S + B' + H_B \bmod \hat{q}$

$ = |A' + H_A|_{\hat{q}}\cdot S + |B' + H_B|_{\hat{q}} \bmod \hat{q}$ 

\textcolor{red}{ $\rhd$ 其中 $|A' + H_A|_{\hat{q}} = A' + H_A \bmod \hat{q}$，\text{ } $|B' + H_B|_{\hat{q}} = B' + H_B \bmod \hat{q}$}

$ $

$ = \hat{A}\cdot S + \hat{B} \bmod \hat{q}$ \textcolor{red}{ $\rhd$ 代入步骤 3：$\hat{A} = A' + H_A \bmod \hat{q}$，\text{ } $\hat{B} = B' + H_B \bmod \hat{q}$}

$ $

因此，证明 $(A' + H_A)\cdot S + B' + H_B - K\cdot \hat{q} = (A' + H_A)\cdot S + B' + H_B \bmod \hat{q}$ 等价于证明 $(A' + H_A)\cdot S + B' + H_B - K\cdot \hat{q} = \hat{A}\cdot S + \hat{B} \bmod \hat{q}$。

$ $
\end{comment}


\item 下面证明 $(A' + H_A)\cdot S + B' + H_B - K\cdot \hat{q} = (A' + H_A)\cdot S + B' + H_B \bmod \hat{q}$：

$(A' + H_A)\cdot S + B' + H_B - K\cdot \hat{q}$ 


$ = (\dfrac{\hat{q}}{q_{l}}\cdot A - \dfrac{\epsilon'_A}{q_{l}} + H_A)\cdot S + (\dfrac{\hat{q}}{q_{l}}\cdot B - \dfrac{\epsilon'_B}{q_{l}} + H_B) - K\cdot \hat{q}$ 

\textcolor{red}{ $\rhd$ 代入步骤 1 的结果：$A' = \dfrac{\hat{q}}{q_{l}}\cdot A - \dfrac{\epsilon'_A}{q_{l}}, \text{ } B' = \dfrac{\hat{q}}{q_{l}}\cdot B - \dfrac{\epsilon'_B}{q_{l}}$}

$ $

$ $

$ = \left(\dfrac{\hat{q}}{q_{l}}\cdot A\cdot S + \dfrac{\hat{q}}{q_{l}}\cdot B  - K\cdot \hat{q}\right) + H_A\cdot S + H_B - \dfrac{\epsilon'_A}{q_{l}}\cdot S - \dfrac{\epsilon'_B}{q_{l}}$ \textcolor{red}{ $\rhd$ 重新排列各项}

$ = \dfrac{\hat{q}}{q_{l}}\cdot(A\cdot S + B  - K\cdot q_{l}) + H_A\cdot S + H_B - \dfrac{\epsilon'_A}{q_{l}}\cdot S - \dfrac{\epsilon'_B}{q_{l}}$ \textcolor{red}{ $\rhd$ 提取公因子 $\dfrac{\hat{q}}{q_{l}}$}

$ = \dfrac{\hat{q}}{q_{l}}\cdot(A\cdot S + B \bmod q_l) + H_A\cdot S + H_B - \dfrac{\epsilon'_A\cdot S + \epsilon'_B}{q_{l}}$ \textcolor{red}{ $\rhd$ 因为 $A\cdot S + B - K\cdot q_{l} = A\cdot S + B \bmod q_l$}

$ $

要成功解密，上式所得多项式的每个系数都必须位于 $\mathbb{Z}_{\hat{q}}$ 范围内（即 $K\cdot\hat{q}$ 成功将 $(A' + H_A)\cdot S + B' + H_B$ 对 $\hat{q}$ 取模）。第一项 $\dfrac{\hat{q}}{q_{l}}\cdot(A\cdot S + B \bmod q_l)$ 可视为将原密文 \textsf{ct} 中的噪声及明文消息按 $\dfrac{\hat{q}}{q_l}$ 缩小，保证处于 $\mathbb{Z}_{\hat{q}}$ 范围内。第二项 $H_A \cdot S$ 的系数也较小，因为 $H_A \in \mathbb{Z}_t^n$ 且 $S \in \{-1, 0, 1\}^n$。第三项 $H_B$ 的系数同样较小，因为 $H_B \in \mathbb{Z}_t^n$。最后一项 $- \dfrac{\epsilon'_A\cdot S + \epsilon'_B}{q_{l}}$ 的系数也较小，因为 $\dfrac{\epsilon'_A}{q_l}$ 和 $\dfrac{\epsilon'_B}{q_l}$ 均为 $\in \mathbb{Z}_{\frac{q_l}{\hat{q}}}$。

$ $

因此，在上述误差阈值成立的条件下，有 $(A' + H_A)\cdot S + B' + H_B - K\cdot \hat{q} = (A' + H_A)\cdot S + B' + H_B \bmod \hat{q}$。

$ $

\item 最后，将步骤 6 和步骤 7 的结果结合：

$\bm((\hat{A}\cdot S + \hat{B}) \bmod \hat{q} \bm) \bmod t$

$= \bm((A' + H_A)\cdot S + B' + H_B \bmod \hat{q}\bm) \bmod t$ \textcolor{red}{ $\rhd$ 因为 $\hat{A} \equiv A' + H_A \bmod \hat{q}$ 且 $\hat{B} \equiv B' + H_B \bmod \hat{q}$}

$ = ((A' + H_A)\cdot S + B' + H_B - K\cdot \hat{q}) \bmod t$ \textcolor{red}{ $\rhd$ 代入步骤 7}

$ = M$ \textcolor{red}{ $\rhd$ 代入步骤 5}

$ $

所以，解密 $(\hat{A}, \hat{B}) \bmod \hat{q}$ 得到消息 $M$。

\end{enumerate}

$ $

BGV 的模数切换可总结如下：



\begin{tcolorbox}[title={\textbf{\tboxlabel{\ref*{subsec:bgv-modulus-switch}} BGV 的模数切换}}]

设当前密文模数为 $q_l$，新密文模数为 $\hat{q}$，满足 $q_l \equiv \hat{q} \equiv 1 \bmod t$ 和 $\hat{q} < q_l$。因此，$\hat{q}$ 可以属于 BGV 密文乘法层级对应的模数集合 $q_0, q_1, \cdots, q_L$，也可以不属于该集合。

$ $

BGV 的模数切换 $q_l \rightarrow \hat{q}$ 等价于按如下方式将 $(A, B) \bmod q_l$ 更新为 $(\hat A, \hat B) \bmod \hat{q}$：


$(A', B') = \left(\left\lceil\dfrac{\hat{q}}{q_l}\cdot A\right\rfloor, \left\lceil\dfrac{\hat{q}}{q_l}\cdot B\right\rfloor\right) \in \mathcal{R}_{\langle n, \hat{q} \rangle}^2$

$ $

$\epsilon'_A = \hat{q}\cdot A - q_l\cdot A'$ \textcolor{red}{ $\rhd$ 其中 $\epsilon'_A \in \mathbb{Z}_{q_l}$}

$\epsilon'_B = \hat{q}\cdot B - q_l\cdot B'$ \textcolor{red}{ $\rhd$ 其中 $\epsilon'_B \in \mathbb{Z}_{q_l}$}


$ $

$H_A = q_l^{-1}\cdot\epsilon'_A \bmod t$

$H_B = q_l^{-1}\cdot\epsilon'_B \bmod t$

$ $

$\hat{\textsf{ct}} = (\hat{A}, \hat{B}) = (A' + H_A, B' + H_B) \bmod \hat{q}$

$ $

BGV 的模数切换用于密文与密文乘法（详见 \autoref{subsec:bgv-mult-cipher}）。完成模数切换后，噪声尺度因子 $\Delta = t$ 保持不变，私钥 $S$ 也保持不变。噪声大致按 $\dfrac{\hat{q}}{q_l}$ 缩小，但噪声与密文模数之比并不会因此下降。要减小这一比值，需要使用模数自举（详见 \autoref{subsec:bgv-bootstrapping}）。

\end{tcolorbox}



\subsubsection{模数切换与 \textsf{ModDrop} 的区别}
\label{subsubsec:bgv-moddrop-vs-modswitch}

对于 CKKS（\autoref{subsubsec:ckks-mult-cipher-rescale}），模数切换与 \textsf{ModDrop} 的区别在于：前者将明文尺度因子按 $\dfrac{q_l}{q_{l-1}} \approx \dfrac{1}{\Delta}$ 缩小，而后者不影响明文尺度因子。

类似地，在 BGV 中，从 $q_l \rightarrow q_{l-1}$ 开始的模数切换（重缩放）和 \textsf{ModDrop} 都会按 $q_l \rightarrow q_{l-1}$ 降低密文模数。关键区别在于，重缩放还将噪声的缩放量按 $\dfrac{q_l}{q_{l-1}} \approx \dfrac{1}{\Delta}$ 减小，而 \textsf{ModDrop} 保持该缩放量不变。因此，重缩放只用于密文与密文乘法（详见 \autoref{subsec:bgv-mult-cipher}），用于按 $\Delta^2 \rightarrow \Delta$ 缩小中间密文中噪声的缩放量。另一方面，当应用确定某个密文不会再参与密文与密文乘法时，即无须继续降低其模数以支持乘法时，可以用 \textsf{ModDrop} 缩短模运算时间。

CKKS 与 BGV 模数切换的主要区别是：前者将明文的缩放量大致按 $\dfrac{1}{\Delta}$ 缩小，而后者将噪声的缩放量大致按 $\dfrac{1}{\Delta}$ 缩小。

\para{源代码：} 可以运行\href{https://github.com/fhetextbook/fhe-textbook/blob/main/source%20code/bgv.py}{\underline{此 Python 脚本}}，查看 BGV 模数切换的示例。


\subsection{密文与密文乘法}
\label{subsec:bgv-mult-cipher}

\noindent \textbf{- 参考资料：}
\href{https://www.inferati.com/blog/fhe-schemes-bgv}{BGV 加密方案简介}

由于 BGV 与 CKKS 一样使用分层密文模数链，其密文与密文乘法也与 CKKS 相同（见 \autoref{subsec:ckks-mult-cipher} 中的总结~\ref*{subsec:ckks-mult-cipher}），唯一区别是重缩放步骤使用 BGV 的模数切换（\autoref{subsec:bgv-modulus-switch}）。

\begin{tcolorbox}[title={\textbf{\tboxlabel{\ref*{subsec:bgv-mult-cipher}} BGV 密文与密文乘法}}]

假设有以下两个 RLWE 密文：

$\textsf{RLWE}_{S, \sigma}(M^{\langle 1 \rangle} + \Delta E^{\langle 1 \rangle}) = (A^{\langle 1 \rangle}, B^{\langle 1 \rangle})$，\text{ } 其中 $B^{\langle 1 \rangle} = -A^{\langle 1 \rangle} \cdot S +  M^{\langle 1 \rangle} + \Delta E^{\langle 1 \rangle}$

$\textsf{RLWE}_{S, \sigma}(M^{\langle 2 \rangle} + \Delta E^{\langle 2 \rangle}) = (A^{\langle 2 \rangle}, B^{\langle 2 \rangle})$，\text{ } 其中 $B^{\langle 2 \rangle} = -A^{\langle 2 \rangle} \cdot S +  M^{\langle 2 \rangle} + \Delta E^{\langle 2 \rangle}$

$ $

两个密文的乘法过程如下：

$ $

\begin{enumerate}
\item \textbf{\underline{基本乘法}}

计算以下各项：

$ $

$D_0 = B^{\langle 1 \rangle}\cdot B^{\langle 2 \rangle}$

$D_1 = A^{\langle 1 \rangle}\cdot B^{\langle 2 \rangle} + A^{\langle 2 \rangle}\cdot B^{\langle 1 \rangle}$

$D_2 = A^{\langle 1 \rangle} \cdot A^{\langle 2 \rangle}$

$ $

解密关系满足：$M^{\langle 1 \rangle}M^{\langle 2 \rangle} + \Delta\cdot (M^{\langle 1 \rangle}E^{\langle 2 \rangle} + M^{\langle 2 \rangle}E^{\langle 1 \rangle}) + \Delta^2E^{\langle 1 \rangle}E^{\langle 2 \rangle}$

$ = \underbrace{B^{\langle 1 \rangle}\cdot B^{\langle 2 \rangle}}_{D_0}  + \underbrace{(B^{\langle 2 \rangle}\cdot A^{\langle 1 \rangle} + B^{\langle 1 \rangle}\cdot A^{\langle 2 \rangle})}_{D_1} \cdot S + \underbrace{(A^{\langle 1 \rangle} \cdot A^{\langle 2 \rangle})}_{D_2} \cdot \underbrace{S \cdot S}_{S^2}$

$= D_0 + D_1\cdot S + D_2 \cdot S^2$

$ $

\item \textbf{\underline{重线性化}}

$\textsf{RLWE}_{S, \sigma}\bm(M^{\langle 1 \rangle}M^{\langle 2 \rangle} + \Delta\cdot (M^{\langle 1 \rangle}E^{\langle 2 \rangle} + M^{\langle 2 \rangle}E^{\langle 1 \rangle}) + \Delta^2E^{\langle 1 \rangle}E^{\langle 2 \rangle}\bm) $

$= \textsf{RLWE}_{S, \sigma}\bm{(}\text{ }D_0 + D_1\cdot S + D_2\cdot S^2\text{ }\bm{)}$

$  \approx C_\alpha + C_\beta, \text{其中} \text{ } C_\alpha = (D_1, D_0), \text{ }\text{ }\text{ } C_\beta = \bm{\langle}  \text{ } \textsf{Decomp}^{\beta, l}(D_2), \textsf{RLev}_{S, \sigma}^{\beta, l}(S^2)  \text{ } \bm{\rangle}$ 

$ $

\item \textbf{\underline{（可选）重缩放}}

为了抑制按 $\Delta^2$ 缩放的噪声分量，使其回到因子 $\Delta$，按照总结~\ref*{subsec:bgv-modulus-switch}（\autoref{subsec:bgv-modulus-switch}）给出的 BGV 模数切换，将密文 $(A, B)$ 更新为 $(\hat A, \hat B)$，从而按 $q \rightarrow \hat q$ 切换密文模数。

$ $

完成上述从 $(A, B)$ 到 $(\hat A, \hat B)$ 的更新后，噪声尺度因子 $\Delta = t$ 和明文 $M$ 保持不变；我们已在 \autoref{subsec:bgv-modulus-switch} 证明 $((\hat A \cdot S + \hat B) \bmod \hat q) \bmod t = M$。

$ $



\end{enumerate}

\para{交换重线性化与重缩放的顺序：} 重线性化和重缩放的顺序可以交换。先执行重缩放可以减小密文模数，从而加快随后重线性化的执行速度。

\end{tcolorbox}

\para{可选重缩放的细节：} 重缩放前，密文所含内容为 $M^{\langle 1 \rangle}M^{\langle 2 \rangle} + \Delta\cdot (M^{\langle 1 \rangle}E^{\langle 2 \rangle} + M^{\langle 2 \rangle}E^{\langle 1 \rangle}) + \Delta^2E^{\langle 1 \rangle}E^{\langle 2 \rangle} + \epsilon$，其中 $\epsilon$ 是重线性化误差。因此，每次密文与密文乘法后，噪声的缩放量都会平方，变为 $\Delta^2, \Delta^4, \Delta^8, \cdots$。为降低这种指数式噪声增长速度，可以在乘法层级 $l$ 的重线性化结束时，选择将密文按 $w_l = \dfrac{q_l}{q_{l-1}} > \Delta$ 重缩放，从而抑制噪声增长，并使噪声尺度因子维持为 $\Delta$。重缩放后，密文按 $w_l$ 缩小，并加入新噪声 $\epsilon_2$。重缩放前，噪声大致按因子 $\Delta = t$ 增长，因为最大的噪声项是 $\Delta^2 E^{\langle 1 \rangle} E^{\langle 2 \rangle}$；重缩放则通过因子 $w_l$ 降低这一增长量，再引入新的常量噪声 $\epsilon_2$。因此，若 $w_l$ 显著大于 $\Delta = t$，所得噪声会小于 $\Delta E^{\langle 1 \rangle}$ 和 $\Delta E^{\langle 2 \rangle}$。基于这一原因，设计 BGV 模数链时，要求各个 $w_l$ 显著大于 $\Delta = t$，以便在每次密文与密文乘法时有效抑制噪声增长，同时确保按 $t$ 取模能得到明文 $M$（详见 \autoref{subsec:bgv-modulus-switch}）。虽然每次重缩放都会新增常量噪声项 $\epsilon_2$，但它会成为重缩放后密文的一部分，并在之后的重缩放中继续按因子 $w_{l-1}$ 缩小。因此，BGV 在密文与密文乘法后执行重缩放，能够有效抑制噪声增长。

上述降噪设计不适用于 CKKS，因为 CKKS 使用尺度因子 $\Delta$ 缩放消息 $M$，而非噪声 $E$；因此，为在多次密文与密文乘法之间保持相同的明文尺度因子 $\Delta$，CKKS 要求各层满足 $w_l \approx \Delta$。由于这一设计差别，CKKS 在每次密文与密文乘法后都会增加噪声。


\subsection{同态密钥切换}
\label{subsec:bgv-key-switching}

BGV 的同态密钥切换将 RLWE 密文对应的私钥从 $S$ 改为 $S'$。该过程与 BFV 的密钥切换完全相同（见 \autoref{subsec:bfv-key-switching} 中的总结~\ref*{subsec:bfv-key-switching}）。

\begin{tcolorbox}[title={\textbf{\tboxlabel{\ref*{subsec:bgv-key-switching}} BGV 的密钥切换}}]
$\textsf{RLWE}_{S',\sigma}(M + \Delta E) = (0, B) + \bm{\langle} \textsf{Decomp}^{\beta, l}(A), \text{ } \textsf{RLev}_{S', \sigma}^{\beta, l}(S) \bm{\rangle}$
\end{tcolorbox}



\subsection{输入向量槽位的同态旋转}
\label{subsec:bgv-rotation}


BGV 对输入向量槽位的同态旋转与 BFV 完全相同（见 \autoref{subsec:bfv-rotation} 中的总结~\ref*{subsec:bfv-rotation}）。


\begin{tcolorbox}[title={\textbf{\tboxlabel{\ref*{subsec:bgv-rotation}} BGV 输入向量槽位的同态旋转}}]


假设有如下 BGV 密文和密钥切换密钥：

\[
\textsf{RLWE}_{S, \sigma}(M + \Delta E) = (A, B), \quad \textsf{RLev}_{S, \sigma}^{\beta, l}(S(X^{J(h)}))
\]

将原输入向量 $\vec{v}$ 的前半部分元素在该半部分内部循环左移 $h$ 个位置，同时将 $\vec{v}$ 的后半部分元素在该半部分内部循环左移 $h$ 个位置，步骤如下：

\begin{enumerate}
\item 将 $A(X)$、$B(X)$ 更新为 $A(X^{J(h)})$、$B(X^{J(h)})$。
\item 执行以下从 $S(X^{J(h)})$ 到 $S(X)$ 的密钥切换（参见 \autoref{subsec:bfv-key-switching}）：

$\textsf{RLWE}_{S(X),\sigma}\bm{(} M(X^{J(h)}) + \Delta E(X^{J(h)})\bm{)} $

$= \bm{(} 0, B(X^{J(h)}) \bm{)} \text{ } + \text{ } \bm{\langle}  \textsf{Decomp}^{\beta, l}\bm{(}A(X^{J(h)})\bm{)}, \text{ } \textsf{RLev}_{S(X), \sigma}^{\beta, l}\bm{(}S(X^{J(h)})\bm{)} \bm{\rangle}$
\end{enumerate}


\end{tcolorbox}




\subsection{模数自举}
\label{subsec:bgv-bootstrapping}

\noindent \textbf{- 参考资料：}
\href{https://eprint.iacr.org/2022/1363.pdf}{重新审视 BGV 与 BFV 的自举}~\cite{cryptoeprint:2022/1363}

BGV 的自举同时具有 BFV 和 CKKS 自举的一些特点：其目标与 CKKS 相同，而内部技术更接近 BFV。与 CKKS 一样，BGV 自举将已消耗的密文模数按 $q_l \rightarrow q_L$ 重置；严格来说，结果是 $q_l \rightarrow q_{l'}$，满足 $l < l' < L$，因为步骤 2$\sim$6 中的自举运算会消耗乘法层级。这一模数变化既重置乘法层级，也降低噪声与密文模数之比。为实现这一目标，可能会想到直接采用 CKKS 自举中的 \textsf{ModRaise}。但这不能直接套用，因为 CKKS 在升模后通过正弦近似消除 $q_0$ 的倍数形式的溢出，而 BGV 是精确加密方案，不允许对明文值作近似。因此，BGV 采用 BFV 中的数位提取来消除模数溢出。为使用数位提取，BGV 也必须像 BFV 一样，将明文模数改成专门选取的 $p^\varepsilon$。为将明文模数与密文模数分别配置为所需的 $p^\varepsilon$ 和 $q_L$，BGV 与 BFV 一样采用同态解密技术。


BGV 自举的技术细节如下。

$ $

设有 RLWE 密文 $(A, B)  = \textsf{RLWE}_{S, \sigma}(M + \Delta E) \bmod q_l$，其中 $A\cdot S + B = M + \Delta E$、$\Delta = t = p$（为素数），$q_l$ 是当前乘法层级的密文模数。

$ $

\begin{enumerate}

\item \textbf{\underline{模数切换} \boldmath$q_l \rightarrow \hat{q}$：} BFV 自举首先按 $q \rightarrow p^{\varepsilon-1}$ 切换密文模数，其中 $q \gg p^\varepsilon > t = p$。BGV 自举则将密文模数切换为特殊模数 $\hat{q}$，使其满足 $\hat{q} \equiv 1 \bmod p^\varepsilon$ 和 $\hat{q} > p^\varepsilon$（$p^\varepsilon$ 将在下一步说明）。为了能按 $q_l \rightarrow \hat{q}$ 切换到该特殊模数，构成 $\hat{q}$ 的素因子必须与 $q_{0}, \ldots, q_L \bmod t$ 同余，这样才能使用 \autoref{subsec:bgv-modulus-switch} 所述技术完成 $q_l\cdot \hat{q} \rightarrow \hat{q}$ 的模数切换。最终，该步骤将密文 $(A, B) \bmod q_l$ 转换为 $(\hat{A}, \hat{B}) \bmod \hat{q}$，同时保持明文模数（即噪声尺度因子）不变。

$ $

\item \textbf{\underline{密文系数乘以 \boldmath$p^{\varepsilon-1}$}：}
将常量 $p^{\varepsilon-1}$ 乘到密文多项式的每个系数上，得到密文 $p^{\varepsilon-1} \cdot (\hat{A}, \hat{B}) = (A', B') \bmod \hat{q}$，其中 $A' = p^{\varepsilon-1}\hat{A}$ 且 $B' = p^{\varepsilon-1}\hat{B}$。该操作将原解密关系 $\hat{A}\cdot S + \hat{B} = M + p E + K\hat{q}$ 更新为 $A'\cdot S + B' = p^{\varepsilon-1} M + p^\varepsilon E + K'\hat{q}$（其中 $||K'||_{\infty} \leq n + 1$）。注意，明文模数（即噪声尺度因子）已按 $p \rightarrow p^\varepsilon$ 发生变化。选择 $\varepsilon$ 时，BGV 还要求满足 $\hat{q} > p^\varepsilon$ 和 $\hat{q}  \equiv 1 \bmod p^\varepsilon$。

$ $

\item \textbf{\underline{\textsf{ModRaise}}：} 将 $(\hat{A}, \hat{B}) \bmod \hat{q}$ 升模为 $(\hat{A}, \hat{B}) \bmod q_L$，其中 $\hat{q} \ll q_L$。升模后密文的解密关系如下：

$\hat{A}\cdot S + \hat{B} = p^{\varepsilon-1}M + p^\varepsilon E + K'\hat{q} \bmod q_L$

$ $

注意，$K'\hat{q}$ 是 $\hat{q}$ 的倍数形式的溢出项，由于 $K'\hat{q} \ll q_L$，该项不会被模 $q_L$ 约减消去。CKKS 自举中的 \textsf{ModRaise}（\autoref{subsubsec:ckks-bootstrapping-high-level}）也存在同样情况：按 $q_0 \rightarrow q_L$ 重置密文模数的代价是引入 $Kq_0$ 溢出，之后通过 \textsf{EvalExp} 的同态近似正弦求值将其去除（\autoref{subsubsec:ckks-bootstrapping-evalexp-details}）。类似地，BGV 升模后的密文 $(\hat{A}, \hat{B}) \bmod q_L$ 为 $\textsf{RLWE}_{S, \sigma}(p^{\varepsilon-1}M + p^\varepsilon E + K'\hat{q}) \bmod q_L$，即对 $p^{\varepsilon-1}M + p^\varepsilon E + K'\hat{q}$ 的加密。后续将像 BFV 自举一样，使用数位提取同态地消除 $K'\hat{q}$。BGV 采用数位提取而非正弦近似求值，是因为它与 BFV 一样属于精确加密方案，而不像 CKKS 那样采用近似计算。

$ $

\item \textbf{\textsf{\underline{CoeffToSlot}}：} 该步骤与 CKKS 和 BFV 中的同名步骤相同：将多项式 $p^{\varepsilon-1} M + p^\varepsilon E + \hat{q}K'$ 的系数移入新密文的输入向量槽位。记该多项式为 $Z = p^{\varepsilon-1} M + p^\varepsilon E + \hat{q}K'$，并将 $Z$ 的第 $i$ 个系数记为 $z_i$。为完成 \textsf{CoeffToSlot}，同态计算 $Z \cdot n^{-1} \cdot \hathat W \cdot I_n^R$，所得密文各槽位便分别存储多项式 $Z$ 的各个 $z_i$。

$ $

\item \textbf{\underline{数位提取}：} 此时，各个输入向量槽位存储 $p^{\varepsilon-1} M + p^\varepsilon E + \hat{q}K'$ 的相应系数，即 $p^{\varepsilon-1} m_i + p^\varepsilon e_i + \hat{q}k'_i$。回顾最低乘法层级的密文模数 $\hat{q}$ 和同态乘法因子 $p^\varepsilon$ 的选取方式：它们满足 $\hat{q} \equiv 1 \bmod p^\varepsilon$，或对某个正整数 $k^{\langle \hat{q} \rangle }$ 满足 $\hat{q} = k^{\langle \hat{q} \rangle } \cdot p^\varepsilon + 1$。因此，有以下关系：

$p^{\varepsilon-1} m_i + p^\varepsilon e_i + \hat{q}k'_i$

$ = p^{\varepsilon-1} m_i + p^\varepsilon e_i + k'_i\cdot(k^{\langle \hat{q} \rangle } \cdot p^\varepsilon + 1) $ \textcolor{red}{ $\rhd$ 代入 $\hat{q} = k^{\langle \hat{q} \rangle } \cdot p^\varepsilon + 1$}

$ = p^{\varepsilon-1} m_i + k'_i + p^\varepsilon  \cdot (e_i + k'_i\cdot k^{\langle \hat{q} \rangle }) $ \textcolor{red}{ $\rhd$ 重新排列各项}

$ = p^{\varepsilon-1} m_i + k'_i + p^\varepsilon \cdot k_i^{\langle \hat{q} + \varepsilon \rangle } $ \textcolor{red}{ $\rhd$ 其中 $ k_i^{\langle \hat{q} + \varepsilon \rangle } = e_i + k'_i\cdot k^{\langle \hat{q} \rangle } $}

$ \equiv p^{\varepsilon-1} m_i + k'_i \pmod{p^\varepsilon}$

$ $

为了消去上式中的 $k'_i$（其中 $|k'_i| \leq n + 1$），使用与 BFV 自举（\autoref{subsubsec:bfv-bootstrapping-digit-extraction}）相同的数位提取多项式 $G_{\varepsilon, v}$：

$z_i = d_0 + \left(\sum\limits_{j=\varepsilon'}^{\varepsilon-1} d_* p^j\right)$ \textcolor{red}{ $\rhd$ 其中 $d_0 \in \mathbb{Z}_p$，\text{ } $d_*$ 可为任意整数，\text{ } 且 $1 \leq \varepsilon' \leq \varepsilon$}

$F_{\varepsilon'}(z_i) \equiv d_0 \bmod p^{\varepsilon'+1}$

$G_{\varepsilon}(z_i) \equiv (z_i - \underbrace{F_{\varepsilon-1} \circ F_{\varepsilon-2} \circ \cdots \circ F_1}_{\varepsilon - 1 \text{次}} (z_i)) \cdot |p^{-1}|_{q}$

$ $

对 $\{\varepsilon, \varepsilon-1, \varepsilon-2, \cdots, 2\}$ 递归求值数位提取多项式 $G_{\varepsilon}$，共执行 $\varepsilon-1$ 次。这会将 $z_i$ 的最低 $\varepsilon-1$ 个以 $p$ 为基的数位置零并右移，如下所示：

$G_{2} \circ G_{3} \circ \cdots \circ G_{\varepsilon-1} \circ G_{\varepsilon} (z_i)$

$= m_i + k_i''p \bmod q_{l'}$

$ $

其中 $k_i''p^\varepsilon$ 是 $p^\varepsilon$ 的某个倍数，包含原有的 $p^\varepsilon$ 倍数溢出项，以及数位提取过程中新增的 $p^\varepsilon$ 倍数溢出项。由于同态求值多项式 $G_{\varepsilon}$ 需要若干次密文与密文乘法，会消耗乘法层级，因此数位提取将密文模数按 $q_L \rightarrow q_{l'}$ 降低，其中 $l'$ 是小于 $L$ 的整数。该步骤的输出是各输入向量槽位中的 $m_i + k''_ip \bmod q_{l'}$。

$ $

\para{\underline{如何处理除以 $p$}：} 与 BFV（\autoref{subsubsec:bfv-bootstrapping-digit-extraction}）不同，BGV 的尺度因子设置不同，因此无法仅靠重新解释尺度因子完成除法。在 BGV 的每轮数位提取中，都应将槽位值显式地同态乘以 $|p^{-1}|_{q^{\langle i \rangle}}$，其中 $q^{\langle i \rangle}$ 是第 $i$ 轮数位提取时的密文模数。设 $(A^{\langle g_1 \rangle}, B^{\langle g_1 \rangle}) \bmod q^{\langle 1 \rangle}$ 是第一轮数位提取中、执行 $p$ 逆元乘法前的密文，则实际的除以 $p$ 运算等价于执行如下密文与明文乘法（\autoref{subsec:bfv-mult-plain}）：

$(A^{\langle g_1 \rangle}, B^{\langle g_1 \rangle}) \cdot \textsf{Encode}(p^{-1}) \bmod q^{\langle 1 \rangle}$

$= (A^{\langle g_1 \rangle}, B^{\langle g_1 \rangle}) \cdot p^{-1} \bmod q^{\langle 1 \rangle}$ \textcolor{red}{ $\rhd$ 因为 $\textsf{Encode}(p^{-1}) = p^{-1} + 0X + 0X^2 + \cdots + 0X^{n-1}$}

$= (p^{-1}A^{\langle g_1 \rangle}, p^{-1}B^{\langle g_1 \rangle}) \bmod q^{\langle 1 \rangle}$

$ $

因此，将两个 BGV 密文多项式 $A^{\langle g_1 \rangle}$ 和 $B^{\langle g_1 \rangle}$ 的系数乘以 $|p^{-1}|_{q^{\langle 1 \rangle}}$，可得到以下效果：

$(|p^{-1}|_{q^{\langle 1 \rangle}} \cdot A^{\langle g_1 \rangle})\cdot S + (|p^{-1}|_{q^{\langle 1 \rangle}} \cdot B^{\langle g_1 \rangle}) \bmod q^{\langle 1 \rangle}$

$= |p^{-1}|_{q^{\langle 1 \rangle}} \cdot (p^{\varepsilon - 1}M + \left\lfloor{K'}\right\rfloor_p + K''p^{\varepsilon}) \bmod q^{\langle 1 \rangle}$


$= p^{\varepsilon - 2}M + \left\lfloor\dfrac{K'}{p}\right\rfloor + K''p^{\varepsilon-1} \bmod q^{\langle 1 \rangle} $

$ $

用文字表述，将 $(A^{\langle g_1 \rangle}, B^{\langle g_2 \rangle})$ 的所有多项式系数乘以 $p$ 的逆元具有两个效果：（1）将明文和噪声按 $p$ 缩小；（2）将模数残余项 $K'$ 减为 $\left\lfloor\dfrac{K'}{p}\right\rfloor$，即右移一个以 $p$ 为基的数位。

$ $

BGV 的数位提取等价于：从输入 $z_i$ 开始递归求值 $G_{\varepsilon}$，每轮将 $\varepsilon$ 减少 1，共执行 $\varepsilon-1$ 轮，如下所示：

\textbf{输入：} $p^{\varepsilon - 1}M + K' + K''p^{\varepsilon} \bmod \hat{q}$

\textbf{第 1 轮：} $G_{\varepsilon}(z_i) \xRightarrow{\text{效果}}  p^{\varepsilon - 2}M + \left\lfloor\dfrac{K'}{p}\right\rfloor + K''^{\langle 1 \rangle}p^{\varepsilon-1} \bmod q^{\langle 1 \rangle}$

\textbf{第 2 轮：} $G_{\varepsilon-1} \circ G_{\varepsilon}(z_i) \xRightarrow{\text{效果}}  p^{\varepsilon - 3}M + \left\lfloor\dfrac{K'}{p^2}\right\rfloor + K''^{\langle 2 \rangle}p^{\varepsilon-2} \bmod q^{\langle 2 \rangle}$

\textbf{第 3 轮：} $G_{\varepsilon-2} \circ G_{\varepsilon-1} \circ G_{\varepsilon}(z_i)  \xRightarrow{\text{效果}}  p^{\varepsilon - 4}M + \left\lfloor\dfrac{K'}{p^3}\right\rfloor + K''^{\langle 3 \rangle}p^{\varepsilon-3} \bmod q^{\langle 3 \rangle}$

$\vdots$

\textbf{第 $\bm{\varepsilon - 1}$ 轮：} $G_{2}\circ \cdots \circ G_{\varepsilon} (z_i) \xRightarrow{\text{效果}}  M + K''^{\langle \varepsilon - 1 \rangle}p \bmod q^{\langle \varepsilon - 1 \rangle}$

$ $

如上所示，整个数位提取过程对明文槽位中的数值产生两个效果：（1）将消息 $p^{\varepsilon - 1}M$ 缩小为 $M$；（2）将模数残余项 $K'$ 置零，即 $\left\lfloor\dfrac{K'}{p^{\varepsilon - 1}}\right\rfloor = 0$。各轮中 $K''$ 会更新为 $K''^{\langle 1 \rangle}, K''^{\langle 2 \rangle}, \cdots, K''^{\langle \varepsilon -1 \rangle}$，原因在于第 $i$ 轮的函数 $F_{\epsilon'}$ 和 $G_{\epsilon}$ 都在模 $p^{\varepsilon - i}$ 下求值，因此每轮都会生成新的 $p^{\varepsilon - i}$ 倍数形式的溢出残余项。

$ $












$\varepsilon - 1$ 轮乘以 $p$ 逆元的操作还带来一个好处：将明文 $p^{\varepsilon-1}m_i \rightarrow m_i$ 和噪声 $k''p^{\varepsilon} \rightarrow k''p$ 缩小，使噪声尺度因子达到标准 BGV 密文所需的 $p$。

$ $

\item \textbf{\textsf{\underline{SlotToCoeff}}：} 该步骤与 BFV 的同名步骤相同：通过同态乘以 $\hathat W^*$，将输入向量槽位中的 $m_i + k''_ip$ 移回多项式系数位置。完成后，由于步骤 $4\sim6$ 的同态运算消耗了若干乘法层级，密文模数变为小于 $q_L$ 的某个值 $q_{l'}$。所得明文系数为 $m_i + k'''_ip$；由于同态执行 \textsf{SlotToCoeff} 又引入了额外噪声，有 $k''' > k''$。此时，明文模数（即噪声尺度因子）为标准 BGV 密文所使用的 $p$。将相应多项式分别记为 $M$ 和 $K'''$，则输出密文为 $\textsf{RLWE}_{S, \sigma}(M + K'''p) \bmod q_{l'}$。



\end{enumerate}



\subsubsection{讨论}

\para{执行模数切换 $q_l \rightarrow \hat{q}$ 的原因：} BGV 通过 $q_l \rightarrow \hat{q}$ 的模数切换来消除自举中的 $q_l$ 倍数溢出。完成模数切换 $q_l \rightarrow \hat{q}$ 并执行 \textsf{ModRaise} 后，加密明文中会出现 $K'\hat{q}$ 溢出项。利用特殊关系 $\hat{q} \equiv 1 \bmod p^\varepsilon$（其中 $p^\varepsilon$ 是明文模数），从明文模数的角度看，该项可约减为 $K'$。

$ $

\para{$\varepsilon$ 的选取：} $\varepsilon$ 越大，$p^{\varepsilon-1}M$ 与 $K'$ 之间以 $p$ 为基的数位间隔就越大，因此解密失败（即未能将 $K'$ 置零）的可能性越小。不过，较大的 $\varepsilon$ 也意味着数位提取成本更高。

$ $


\para{推广到 $\Delta = p^r$：} 与 BFV 自举（\autoref{subsubsec:bfv-bootstrapping-summary} 中的总结~\ref*{subsubsec:bfv-bootstrapping-summary}）一样，可以将明文模数（即噪声尺度因子）推广为 $p^r$，其中 $p$ 为素数，$r$ 可为任意正整数。







\subsubsection{BFV、BGV 和 CKKS 自举的比较}
\label{subsubsec:bootstrapping-differences}

BFV 和 BGV 的自举都使用数位提取，但目的不同。在 BFV 中，数位提取用于消除低位区域的噪声；在 BGV 中，它用于消除按 $\hat{q}\rightarrow q_L$ 执行 \textsf{ModRaise} 后，在低位区域产生的模数残余值。此外，BGV 每轮数位提取都显式地将密文多项式系数乘以 $|p^{-1}|_{q^{\langle i \rangle}}$。BFV 则省略这种显式乘以 $p$ 逆元的操作，只在概念上重新解释尺度因子。

CKKS 自举不包含数位提取，因为它没有需要从低位区域消去的数值。对于按 $q_0 \rightarrow q_L$ 执行 \textsf{ModRaise} 后产生的模数残余值，CKKS 使用 \textsf{EvalExp} 同态求值近似正弦函数来消除。虽然 CKKS 自举将模数重置为较大的值，但它不会降低噪声与消息之比；这一比值仍会在同态运算中持续增大。

%In the case of BFV's bootstrapping, we need homomorphic decryption (\autoref{subsubsec:bfv-bootstrapping-homomorphic-decryption}) because we need to simultaneously change the ciphertext's plaintext scaling factor from $p^{\varepsilon-1} \rightarrow \left\lfloor\dfrac{q}{p}\right\rfloor$ and the ciphertext modulus from $p^\varepsilon \rightarrow q$.For the similar reason, in the case of BGV's bootstrapping, we used digit extraction to change the noise scaling factor from  $\hat{q} \rightarrow q_L$ ((or technically, $\hat{q} \rightarrow q_{l'}$ where $\hat{q} < q_{l'} < q_L$) while keeping the noise scaling factor (i.e., the plaintext modulus) $\Delta = p^\varepsilon$. On the other hand, in the case of CKKS's bootstrapping, \textsf{ModRaise} instead of homomorphic decryption is sufficient because we only need to change the ciphertext modulus from $q_0 \rightarrow q_L$ (or technically, $q_0 \rightarrow q_{l'}$ where $q_0 < q_{l'} < q_L$) while keeping the same plaintext scaling factor $\Delta \approx \dfrac{q_l}{q_{l-1}}$.  




\clearpage


\section{采用 RNS 的 FHE 方案}
\label{sec:rns}
BFV、BGV 和 CKKS 的参数可能需要超过机器字长的系数模数。剩余数系统（RNS，\autoref{subsec:crt-application}）使用较小且两两互素模数下的剩余类来表示这些系数，使精确模算术可以由机器字长运算实现，而不缩小总模数或改变基础安全参数。所需模数大小取决于方案、电路、精度和安全目标，没有普遍适用的固定比特长度。

CRT 使模加法与模乘法易于实现，但改变基以及实现缩放舍入还需要额外算法。本节针对 BFV、BGV 和 CKKS 介绍这些原语。常见 TFHE 实现采用可由一个机器字表示的环面系数，因此不需要这种特定的大整数分解；这一实现选择，并不是禁止其他 TFHE 变体采用 RNS 技术的定理。

本节中，$[x]_q=|x|_q$ 表示区间 $[-q/2,q/2)$ 中的唯一整数代表元（\autoref{subsec:modulo-centered}）。除非另有说明，RNS 模数均为两两互素的奇整数；使用素数模数很方便，但中国剩余定理本身并不要求它们是素数。对于奇数 $q$，最小与最大代表元分别是 $-(q-1)/2$ 和 $(q-1)/2$。必须区分这些整数代表元的界与同余关系，这对近似和精确基转换都至关重要。

\begin{tcolorbox}[
    title = \textbf{所需背景知识},    % box title
    colback = white,    % light background; tweak to taste
    colframe = black,  % frame colour
    boxrule = 0.8pt,     % line thickness
    left = 1mm, right = 1mm, top = 1mm, bottom = 1mm % inner padding
]

\begin{itemize}
\item \autoref{sec:modulo}: \nameref{sec:modulo}
\item \autoref{sec:group}: \nameref{sec:group}
\item \autoref{sec:field}: \nameref{sec:field}
\item \autoref{sec:polynomial-ring}: \nameref{sec:polynomial-ring}
\item \autoref{sec:decomp}: \nameref{sec:decomp}
\item \autoref{sec:modulus-rescaling}: \nameref{sec:modulus-rescaling}
\item \autoref{sec:chinese-remainder}: \nameref{sec:chinese-remainder}
\item \autoref{sec:polynomial-interpolation}: \nameref{sec:polynomial-interpolation}
\item \autoref{sec:ntt}: \nameref{sec:ntt}
\item \autoref{sec:lattice}: \nameref{sec:lattice}
\item \autoref{sec:rlwe}: \nameref{sec:rlwe}
\item \autoref{sec:glwe}: \nameref{sec:glwe}
\item \autoref{sec:glwe-add-cipher}: \nameref{sec:glwe-add-cipher}
\item \autoref{sec:glwe-add-plain}: \nameref{sec:glwe-add-plain}
\item \autoref{sec:glwe-mult-plain}: \nameref{sec:glwe-mult-plain}
\item \autoref{subsec:modulus-switch-rlwe}: \nameref{subsec:modulus-switch-rlwe}
\item \autoref{sec:glwe-key-switching}: \nameref{sec:glwe-key-switching}
\item \autoref{sec:bfv}: \nameref{sec:bfv}
\item \autoref{sec:ckks}: \nameref{sec:ckks}
\item \autoref{sec:bgv}: \nameref{sec:bgv}
\end{itemize}
\end{tcolorbox}

\clearpage

\subsection{快速基转换：\textsf{FastBConv}}
\label{subsec:rns-fastbconv}

\noindent \textbf{- 参考文献：} \href{https://eprint.iacr.org/2016/510}{A Full RNS Variant of FV-like Somewhat Homomorphic Encryption Schemes}~\cite{rns-bfv}；\href{https://eprint.iacr.org/2022/657}{BASALISC: Programmable Hardware Accelerator for BGV FHE}~\cite{rns-bfv2}。

令 $q=\prod_{i=1}^kq_i$、$b=\prod_{j=1}^lb_j$，所有基模数两两互素。输入剩余类 $x_i=[x]_{q_i}$ 确定中心化整数 $x=[x]_q$。精确基转换应输出 $([x]_{b_j})_{j=1}^l$。快速基转换省略一次大模数约减，而在每个输出坐标使用同一个未经约减的 CRT 和：
\[
y_i=q/q_i,\qquad z_i=y_i^{-1}\bmod q_i,\qquad
T=\sum_{i=1}^k[x_iz_i]_{q_i}y_i.
\]
\begin{tcolorbox}[title={\textbf{\tboxlabel{\ref*{subsec:rns-fastbconv}} 快速基转换：\textsf{FastBConv}}}]
\[
\textsf{FastBConv}(x,q,b)=([T]_{b_1},\ldots,[T]_{b_l}).
\]
存在整数 $u$ 使 $T=x+uq$。对于上述奇模数中心化约定，
\[
|T|<kq/2,\quad |x|<q/2,\quad |u|<(k+1)/2,
\quad |u|\leq\lfloor k/2\rfloor.
\]
因此输出表示 $x+uq\pmod b$，不一定表示 $x\pmod b$。记该界为 $U_k=\lfloor k/2\rfloor$。
\end{tcolorbox}

\begin{myproof}
每个加项的绝对值小于 $q/2$，且 CRT 同余关系给出 $T\equiv x\pmod q$。对 $u=(T-x)/q$ 的界必须同时考虑 $T$ 与 $x$；只约束 $T$，不能将同一区间直接用作 $uq$ 的界。所有输出剩余类都由共同的 $T$ 计算，因此误差是在整个目标基上一致的倍数 $uq$。其他剩余类约定会得到不同的界，包括~\cite{rns-bfv} 采用的无符号约定。
\end{myproof}

这是基转换的近似实现，而不是中国剩余定理失效。后续算法必须消除该 $q$ 的倍数，将其通过除法转化为受控误差，或明确将其计入正确性界。

\subsection{小蒙哥马利约减算法：\textsf{SmallMont}}
\label{subsec:rns-smallmont}

Small Montgomery 约减限制快速基转换引入的 $q$ 的倍数，一般不会使基转换变为精确操作。下面的中心化剩余类变体明确列出所需条件；原始无符号剩余类算法及其参数界见~\cite{rns-bfv} 的算法 2 和引理 4。

令 $x=[x]_q$，选取与 $q$、$b$ 均互素的奇数辅助模数 $b_\alpha$，并计算 $[b_\alpha x]_q$ 的快速 CRT 提升：
\[
T=[b_\alpha x]_q+uq=b_\alpha x+(v+u)q,
\qquad |u|\leq U_k.
\]
输入由同一个 $T$ 在各个 $b_j$ 及 $b_\alpha$ 下的剩余类组成。
\begin{tcolorbox}[title={\textbf{\tboxlabel{\ref*{subsec:rns-smallmont}} Small Montgomery 约减：\textsf{SmallMont}}}]
\begin{enumerate}
\item 使用给定的 $b_\alpha$ 剩余类，计算 $a=[Tq^{-1}]_{b_\alpha}$。
\item 对 $j=1,\ldots,l$，输出 $r_j=[(T_{b_j}-qa)b_\alpha^{-1}]_{b_j}$。
\end{enumerate}
结果表示 $r=(T-qa)/b_\alpha=x+u'q\pmod b$。若 $b_\alpha>U_k+1/2$，则 $u'\in\{-1,0,1\}$。
\end{tcolorbox}

\begin{myproof}
因为 $a\equiv v+u\pmod{b_\alpha}$，商 $u'=(v+u-a)/b_\alpha$ 是整数。中心化代表元给出 $|v|<(b_\alpha+1)/2$ 和 $|a|<b_\alpha/2$，因此
\[
|u'|<1+\frac{U_k+1/2}{b_\alpha}<2.
\]
这就证明了所述的三种可能整数值。分子 $T-qa$ 被 $b_\alpha$ 整除，可以在每个 $b_j$ 下通过乘以 $b_\alpha^{-1}\bmod b_j$ 计算其商的剩余类。
\end{myproof}

计算过程中必须始终保留辅助剩余类。先重构 $[T]_b$，再将该值对 $b_\alpha$ 取模，通常会丢失原来的 $b_\alpha$ 剩余类。此外，$v$ 可能随 $b_\alpha$ 线性增长；声称其增长更慢并不能证明该界。

\subsubsection{使用 \textsf{SmallMont} 改进 \textsf{FastBConv}}
\label{subsubsec:rns-smallmont-fastbconv}

满足 \autoref{subsec:rns-smallmont} 的参数条件时，
\[
\textsf{SmallMont}\bigl(\textsf{FastBConv}([b_\alpha x]_q,q,bb_\alpha),bb_\alpha,b_\alpha,q\bigr)
=[x+u'q]_b,\qquad u'\in\{-1,0,1\}.
\]
预先乘以 $b_\alpha$ 是必要步骤。结果仍是中心化 $x$ 的近似提升，与 \autoref{subsec:rns-fastbconvex} 的 \textsf{FastBConvEx} 不同。


\subsection{基于 RNS 的 ModRaise：\textsf{ModRaise\textsubscript{RNS}}}
\label{subsec:rns-modraise}

\noindent \textbf{- 参考文献：} \href{https://eprint.iacr.org/2018/931}{A Full RNS Variant of Approximate Homomorphic Encryption}~\cite{rns-ckks}，第 3.2 节。

从 $q$ 到 $qb$ 的基扩展保留模各个 $q_i$ 的剩余类，并附加模各个 $b_j$ 的剩余类。使用快速基转换时，它选择一个近似整数提升：
\begin{tcolorbox}[title={\textbf{\tboxlabel{\ref*{subsec:rns-modraise}} \textsf{ModRaise\textsubscript{RNS}}}}]
\[
\textsf{ModRaise}_{\mathrm{RNS}}(x,q,qb)
=(x_1,\ldots,x_k,\textsf{FastBConv}(x,q,b)).
\]
组合后的向量表示 $x+uq\pmod{qb}$，其中 $|u|\leq U_k$。采用 \autoref{subsec:rns-smallmont} 中的 Montgomery 预处理，可以将其限制为满足 $|u'|\leq1$ 的适当提升 $x+u'q$。
\end{tcolorbox}

这是整数代表元的选择，并不是从 $\mathbb Z_q$ 到 $\mathbb Z_{qb}$ 的典范环嵌入。若逐系数作用于密文分量，将新代表元写为 $A+qU_A$ 和 $B+qU_B$，则其相位改变
\[
q(U_B+U_AS).
\]
因此，分量界 $\|U_A\|_\infty,\|U_B\|_\infty\leq U_k$ 给出的相位界是 $qU_k(1+\|S\|_1)$，而不只是 $qU_k$。整数相位 $B+AS$ 中还可能已经含有 $q$ 的倍数。这些项必须由后续乘法或自举算法处理，不能视为独立的新 RLWE 误差。

\subsection{基于 RNS 的 ModDrop：\textsf{ModDrop\textsubscript{RNS}}}
\label{subsec:rns-moddrop}

\textsf{ModDrop\textsubscript{RNS}} 从模 $qb$ 的表示中删除对应于 $b$ 的各因子的剩余类：
\begin{tcolorbox}[title={\textbf{\tboxlabel{\ref*{subsec:rns-moddrop}} \textsf{ModDrop\textsubscript{RNS}}}}]
\[
(x_{q_1},\ldots,x_{q_k},x_{b_1},\ldots,x_{b_l})
\longmapsto(x_{q_1},\ldots,x_{q_k}).
\]
这是精确的模 $q$ 约减，不进行缩放或舍入。
\end{tcolorbox}

对于 CKKS 或 BGV，只有原中心化整数相位的每个系数都已经位于 $[-q/2,q/2)$ 内，才能保持同一相位。此时 CKKS 保持尺度不变，BGV 保持明文模数不变。否则，相位回绕可能改变解码消息。对于 BFV，即使相位不变，也通常不能保持消息，因为标准尺度从 $\lfloor qb/t\rfloor$ 变为 $\lfloor q/t\rfloor$。例如，初始模数为 $105$、$t=7$，无噪声相位 $15$ 编码 $m=1$；降至模数 $35$ 后相位仍为 $15$，但标准 BFV 解码得到 $3$。需要保持消息时，应采用 BFV 专用的缩放舍入操作。

\subsection{基于 RNS 的模数切换：\textsf{ModSwitch\textsubscript{RNS}}}
\label{subsec:rns-modswitch}

令整数 $x$ 以 $q$ 与 $b$ 的组合基给出。对于奇数 $b$，定义 $r=[x]_b$，则 $(x-r)/b=\lfloor x/b\rceil$。下面的操作近似除以 $b$，并删除其基因子：
\begin{tcolorbox}[title={\textbf{\tboxlabel{\ref*{subsec:rns-modswitch}} \textsf{ModSwitch\textsubscript{RNS}}}}]
\begin{enumerate}
\item 使用 \textsf{FastBConv} 将 $r$ 的剩余类从基 $b$ 转换到基 $q$。其共同整数提升是 $r+ub$，其中 $|u|\leq U_l$。
\item 对每个 $i=1,\ldots,k$，输出
\[
y_i=[b^{-1}(x_{q_i}-(r+ub)_{q_i})]_{q_i}.
\]
\end{enumerate}
结果表示
\[
y=\left\lfloor\frac{x}{b}\right\rceil-u\pmod q,\qquad |u|\leq U_l.
\]
\end{tcolorbox}

代入 $x-r=b\lfloor x/b\rceil$ 即可得到该式。此同余关系不要求 $b\ll q$。只删除一个奇模数时，$l=1$ 且 $u=0$，结果就是精确的最近整数除法。删除多个因子时，必须将额外误差 $-u$ 计入密文相位界。作用于 $(A,B)$ 后，总扰动还包含分量误差与 $S$ 的乘积，不能只用单个系数的 $u$ 作为其界。

CKKS 可将其用作重缩放原语，同时跟踪实际除以 $b$ 后的尺度和全部舍入误差。BFV 和 BGV 需要各自的明文保持规则。特别地，普通的缩放舍入并不是 \autoref{subsec:rns-modswitch-bgv} 中的 BGV 模数切换。

\subsubsection{\textsf{ModSwitch\textsubscript{RNS}}、\textsf{ModRaise\textsubscript{RNS}} 与 \textsf{ModDrop\textsubscript{RNS}} 的比较}

\textsf{ModRaise\textsubscript{RNS}} 选择进入更大组合基的提升，\textsf{ModDrop\textsubscript{RNS}} 精确约减到更小的子基，\textsf{ModSwitch\textsubscript{RNS}} 则还会对所表示整数作近似除法。这是不同的操作。它们能否保持密文的明文，取决于具体方案、尺度或明文模数，以及完整的中心化相位界。




\subsection{基于 RNS 的解密}
\label{subsec:rns-dec}

本小节介绍如何高效解密采用 RNS 表示的 BFV、CKKS 与 BGV 密文。


\subsubsection{BFV 解密：$\textsf{Dec}_{\textsf{RNS}}^{\textsf{BFV}}$}
\label{subsubsec:rns-dec-bfv}

采用本书的密文约定 $B=-AS+\Delta M+E\pmod q$。对每个中心化相位系数 $x=[B+AS]_q$，标准 BFV 解码操作是 $[\lfloor tx/q\rceil]_t$；除以 $\Delta=\lfloor q/t\rfloor$ 并不是代数上相同的运算。下面的冗余模数纠正在满足舍入余量时实现缩放舍入，参见~\cite{rns-bfv} 第 3 节。

令 $q=\prod_{i=1}^kq_i$ 为奇数、$\gcd(q,t)=1$，并选择与 $qt$ 互素的奇数 $\gamma$。记
\[
m_0=\left\lfloor\frac{tx}{q}\right\rceil,\quad
\delta=\frac{tx}{q}-m_0,\quad v=[t\gamma x]_q.
\]
$v$ 的快速 CRT 提升为 $v+uq$，其中 $|u|\leq U_k$。计算其模 $\gamma$ 和 $t$ 的剩余类，并分别乘以对应模数下的 $-q^{-1}$，得到 $(y_\gamma,y_t)$。于是
\[
y\equiv-q^{-1}(v+uq)
\equiv\left\lfloor\frac{t\gamma x}{q}\right\rceil-u
=\gamma m_0+\eta\pmod{\gamma t},
\qquad \eta=\lfloor\gamma\delta\rceil-u.
\]
若 $|\eta|<\gamma/2$，其中心化剩余类恰为 $y_\gamma=\eta$，因此
\[
[m_0]_t=[(y_t-y_\gamma)\gamma^{-1}]_t.
\]
一个充分条件是 $\gamma|\delta|+1/2+U_k<\gamma/2$。

\begin{tcolorbox}[title={\textbf{\tboxlabel{\ref*{subsubsec:rns-dec-bfv}} $\textsf{Dec}_{\textsf{RNS}}^{\textsf{BFV}}$}}]
\begin{enumerate}
\item 计算相位剩余类 $x_i=[B+AS]_{q_i}$。
\item 选择满足互素条件和舍入余量界的 $\gamma$，将 $[t\gamma x]_q$ 快速转换到两个模数 $\gamma,t$。
\item 将两个剩余类分别乘以对应模数下的 $-q^{-1}$，得到 $(y_\gamma,y_t)$，其中 $y_\gamma$ 取中心化代表元。
\item 输出 $[(y_t-y_\gamma)\gamma^{-1}]_t$。
\end{enumerate}
\end{tcolorbox}

BFV 正确解密还要求 $[m_0]_t$ 等于预期明文。仅选择任意互素的 $\gamma$ 并不足够，而且纠正阈值是 $\gamma/2$，不是 $\gamma$。任意接近舍入边界的值需要相应更强的精度或纠正策略。这不是一个适用于任意输入与模数的无条件精确舍入原语。\cite{rns-bfv-hps} 中其他缩放舍入构造也规定了各自的精度和误差条件。


\subsubsection{CKKS 与 BGV 解密}
\label{subsubsec:rns-dec-ckks-bgv}

在当前 RNS 基中计算 $[B+AS]_q$。CKKS 按照已跟踪的尺度解码其中心化整数代表元，BGV 则将该代表元对明文模数取模。CRT 重构是一种有效实现。只有整个中心化相位都位于较小模数的一半以内，才能在解码前直接降到单个小模数。当前层级的有效密文不一定满足这个更强的界。必要时，应使用符合方案的重缩放或模数切换并分析误差，不能假设单纯 \textsf{ModDrop} 到最低层始终保持消息。

\subsection{BGV 的 RNS 模数切换：$\textsf{ModSwitch}_{\textsf{RNS}}^{\textsf{BGV}}$}
\label{subsec:rns-modswitch-bgv}

BGV 模数切换在降低密文模数的同时，必须保持明文剩余类。\autoref{subsec:rns-modswitch} 中的普通舍入原语不保证这一条件。也不能通过令 $t=\hat q$ 调用 $\textsf{Dec}_{\textsf{RNS}}^{\textsf{BFV}}$ 来获得无条件的舍入操作：若 $\hat q$ 整除 $q_l$，所需逆元 $q_l^{-1}\bmod\hat q$ 根本不存在。

下面给出一种常用模数链步骤的直接 RNS 实现。令 $q=Qw$，满足 $\gcd(w,Q)=\gcd(w,t)=1$ 和 $w\equiv1\pmod t$。对任一密文分量的每个中心化系数 $a$，定义
\[
d=t[t^{-1}a]_w,\qquad a'=\frac{a-d}{w}\pmod Q.
\]
由于 $d\equiv a\pmod w$ 且 $d\equiv0\pmod t$，这是整数除法，而且在对 $Q$ 约减前满足 $a'\equiv a\pmod t$。在每个保留模数 $q_i\mid Q$ 下，通过下式计算：
\[
a'_i=[w^{-1}(a_i-d)]_{q_i}.
\]
若 $w$ 含有多个 RNS 因子，必须重构中心化剩余类 $[t^{-1}a]_w$，或采用精确且具有已证明界限的转换过程；单因子步骤不需要大整数重构。

令输入相位为 $M+tE$，并使用 \autoref{subsec:bgv-modulus-switch} 中同一个溢出多项式。约减前的输出相位为
\[
Y=\frac{M+tE}{w}-\frac{d_B+d_AS}{w}.
\]
由于 $q=Qw\equiv Q\pmod t$，该相位在模 $t$ 下与 $M$ 同余。一个充分的不回绕条件是
\[
\frac{\|M+tE\|_\infty}{w}+\frac{t}{2}(1+\|S\|_1)<\frac Q2.
\]
在此条件下，输出是相同明文的有效 BGV 密文。更一般的模数链可以跟踪明文纠正因子；这里的同余假设是明确声明的特例，并不是对所有 BGV 的限制。


\subsection{精确快速基转换：\textsf{FastBConvEx}}
\label{subsec:rns-fastbconvex}

与 \textsf{SmallMont} 不同，带冗余模数的精确基转换在满足明确界限时会消除全部 CRT 歧义。这是~\cite{rns-bfv} 第 4.4 节引理 6 使用的 Shenoy--Kumaresan 型纠正。令 $b=\prod_{j=1}^lb_j$，假设已知整数 $x$ 在每个 $b_j$ 以及额外互素奇模数 $b_\alpha$ 下的剩余类。写成
\[
\widehat x=[x]_b=x+\mu b,\qquad
T=\widehat x+ub=x+(\mu+u)b,
\]
其中 $T$ 是两次快速转换共同使用的未经约减的 CRT 和，且 $|u|\leq U_l$。要求 $|\mu+u|<b_\alpha/2$；若已知 $|\mu|\leq\lambda$，可以采用充分条件 $b_\alpha>2(\lambda+U_l)$。

\begin{tcolorbox}[title={\textbf{\tboxlabel{\ref*{subsec:rns-fastbconvex}} 精确快速基转换：\textsf{FastBConvEx}}}]
\begin{enumerate}
\item 保留 $x$ 的基 $b$ 剩余类，并单独保留给定剩余类 $x_\alpha=[x]_{b_\alpha}$。
\item 计算 $\gamma=[(\textsf{FastBConv}(\widehat x,b,b_\alpha)-x_\alpha)b^{-1}]_{b_\alpha}$。
\item 在目标基的各坐标中输出 $[\textsf{FastBConv}(\widehat x,b,q)-\gamma b]_q$。
\end{enumerate}
\end{tcolorbox}

\begin{myproof}
两次转换都使用 $T$。因此在所述界限下，$\gamma=[\mu+u]_{b_\alpha}=\mu+u$，从而 $T-\gamma b=x$，输出恰为 $[x]_q$。这是代数恒等式，不会遗留新的近似误差。若不满足该界，纠正结果可能相差 $bb_\alpha$ 的倍数，而该倍数在模 $q$ 下不一定为零。
\end{myproof}

\para{为何需要中心化（即有符号）剩余类表示：} 必须将纠正项 $\gamma$ 解释为有界有符号整数 $\mu+u$。当 $|\mu+u|<b_\alpha/2$ 时，中心化约定提供了这种解释。无符号实现也可以使用，但必须显式恢复符号；直接把无符号剩余类当作该整数是不正确的。所需条件是可量化的范围界，而不是未经量化的“$x$ 远小于模数”。


\subsection{乘法分解：\textsf{DecompMult\textsubscript{RNS}}}
\label{subsec:rns-decompmult}

RNS 为密钥切换和重线性化提供了一个整数 gadget。令 $q=\prod_iq_i$，定义
\[
g_i=\frac q{q_i}\left[\left(\frac q{q_i}\right)^{-1}\right]_{q_i},\qquad
A_i=[A]_{q_i},\qquad A\equiv\sum_iA_ig_i\pmod q.
\]
这里将每个 $A_i$ 逐系数提升为小整数多项式。将其乘以 $R_q$ 中的密文时，必须在 RNS 基的\emph{每一个}模数下约减同一组小整数。如果只保留第 $i$ 个坐标、其余填零，就会得到不同的乘数和误差界。

\begin{tcolorbox}[title={\textbf{\tboxlabel{\ref*{subsec:rns-decompmult}} \textsf{DecompMult\textsubscript{RNS}}}}]
从 $S$ 切换到 $S'$ 时，生成满足下式的求值密钥 $(a_i,b_i)$：
\[
b_i+a_iS'\equiv g_iS+E_i\pmod q.
\]
于是 $\sum_iA_i(a_i,b_i)$ 在模 $q$ 下具有相位 $AS+\sum_iA_iE_i$。加上 $(0,B)$ 即完成密钥切换。在 BGV 中应以 $tE_i$ 替代 $E_i$，使新增误差被明文模数整除。

重线性化时，将密钥中的 $g_iS$ 替换为 $g_iS^2$，分解 $D_2$，再按照本书的 $(A,B)$ 顺序，将结果加到 $(D_1,D_0)$ 上。所得相位为
\[
D_0+D_1S+D_2S^2+\sum_iD_{2,i}E_i\pmod q.
\]
\end{tcolorbox}

对于负循环多项式，$\|A_iE_i\|_\infty\leq n\|A_i\|_\infty\|E_i\|_\infty$，因此新增误差的一个明确界为 $n\sum_i(q_i/2)\|E_i\|_\infty$。也可以把模逆元放入数位 $[A_i(q/q_i)^{-1}]_{q_i}$，并采用权重 $q/q_i$，如~\cite{rns-bfv} 第 4.5 节所示。两种选择都有效，但不能混用它们的数位与密钥。这些恒等式中的加密符号表示相位关系，而不是重新随机采样的密文分布相等。


\subsection{将 RNS 技术应用于 FHE 运算}
\label{subsec:rns-application}

本小节说明如何使用此前介绍的 RNS 基本操作，处理 BFV、CKKS 与 BGV 中采用 RNS 表示的密文运算。


\subsubsection{多项式加法与乘法}
\label{subsubsec:rns-application-basic}

RNS 分解的是系数模数，因此标量加法与标量乘法分别在每个 $q_i$ 下进行（\autoref{subsec:crt-application}）。多项式加法逐系数进行，但多项式乘法仍需要模 $X^n+1$ 的负循环卷积。只有在模数支持适当 NTT 并完成变换后，乘法才变成逐点乘法。下面的标量恒等式描述了这些多项式运算内部使用的系数算术。

$P^{\langle 1 \rangle} = \sum\limits_{a=0}^{n-1} c^{\langle 1 \rangle}_a \cdot X^a \in \mathcal{R}_{\langle n, q\rangle}$

$P^{\langle 2 \rangle} = \sum\limits_{b=0}^{n-1} c^{\langle 2 \rangle}_b \cdot X^b \in \mathcal{R}_{\langle n, q\rangle}$

$ $

在 FHE 方案的 RNS 变体中，将每个多项式系数表示为如下 RNS 剩余向量：

$c^{\langle 1 \rangle}_a = (c^{\langle 1 \rangle}_{a,1}, c^{\langle 1 \rangle}_{a,2}, \cdots, c^{\langle 1 \rangle}_{a,k}) \in \mathbb{Z}_{q_1} \times \mathbb{Z}_{q_2} \times \cdots \times \mathbb{Z}_{q_k}$ \textcolor{red}{ $\rhd$ 对于 $a \in [0, n - 1]$}

$c^{\langle 2 \rangle}_b = (c^{\langle 2 \rangle}_{b,1}, c^{\langle 2 \rangle}_{b,2}, \cdots, c^{\langle 2 \rangle}_{b,k}) \in \mathbb{Z}_{q_1} \times \mathbb{Z}_{q_2} \times \cdots \times \mathbb{Z}_{q_k}$ \textcolor{red}{ $\rhd$ 对于 $b \in [0, n - 1]$}

$ $

在上述 RNS 设置下，对两个多项式执行加法或乘法时，每次系数相加都通过两个 RNS 剩余向量的逐元素加法计算：

$c^{\langle 1 \rangle}_a + c^{\langle 2 \rangle}_b \equiv \sum\limits_{i=1}^k (c^{\langle 1 \rangle}_{a,i} + c^{\langle 2 \rangle}_{b,i}) y_i z_i \bmod q $ \textcolor{red}{ $\rhd$ 其中 $y_i = \dfrac{q}{q_i}$、$z_i = |y_i^{-1}|_{q_i}$}

$\Longleftrightarrow (c^{\langle 1 \rangle}_{a,1} + c^{\langle 2 \rangle}_{b,1}, \text{ } c^{\langle 1 \rangle}_{a,2} + c^{\langle 2 \rangle}_{b,2}, \cdots, \text{ } c^{\langle 1 \rangle}_{a,k} + c^{\langle 2 \rangle}_{b,k}) \in \mathbb{Z}_{q_1} \times \mathbb{Z}_{q_2} \times \cdots \times \mathbb{Z}_{q_k}$

$ $

类似地，每次系数相乘都通过两个 RNS 剩余向量的逐元素乘法计算：

$c^{\langle 1 \rangle}_a \cdot c^{\langle 2 \rangle}_b \equiv \sum\limits_{i=1}^k (c^{\langle 1 \rangle}_{a,i} \cdot c^{\langle 2 \rangle}_{b,i}) y_i z_i \bmod q $

$\Longleftrightarrow (c^{\langle 1 \rangle}_{a,1} \cdot c^{\langle 2 \rangle}_{b,1}, \text{ } c^{\langle 1 \rangle}_{a,2} \cdot c^{\langle 2 \rangle}_{b,2}, \cdots, \text{ } c^{\langle 1 \rangle}_{a,k} \cdot c^{\langle 2 \rangle}_{b,k}) \in \mathbb{Z}_{q_1} \times \mathbb{Z}_{q_2} \times \cdots \times \mathbb{Z}_{q_k}$

$ $

利用上述同构，可以基于较小的 RNS 剩余（例如每个 30 位），高效地对大多项式系数（例如 1000 位）执行密文与明文相加、密文与密文相加以及密文与明文相乘。


\subsubsection{密钥切换}
\label{subsubsec:rns-application-key-switching}

使用 \autoref{subsec:rns-decompmult} 的求值密钥和提升数位，计算
\[
(A',B')=(0,B)+\textsf{DecompMult}_{\mathrm{RNS}}(A,S_{\langle S',\mathrm{RNS}\rangle}).
\]
若输入相位为 $B+AS$，则输出相位在模 $q$ 下为 $B+AS+E_{\mathrm{ks}}$，其中 $E_{\mathrm{ks}}=\sum_iA_iE_i$。BFV 和 CKKS 保持当前消息尺度，同时增加该误差。BGV 采用误差被 $t$ 整除的求值密钥，使相位仍在模 $t$ 下与同一消息同余。全部运算均可在当前 RNS 基中执行，但相位界仍取决于具体方案。


\subsubsection{输入槽位旋转}
\label{subsubsec:rns-rotation}

对于 $J(h)=5^h\bmod2n$，对两个密文分量施加自同构 $X\mapsto X^{J(h)}$。将指数模 $2n$ 约减，并使用 $X^{n+i}=-X^i$。在系数表示中，这相当于置换整个 RNS 剩余类向量，并在所有模数下对相应向量取负。该步骤本身不引入近似误差。

所得相位是原相位在 $X^{J(h)}$ 处的代入结果，秘密密钥变为 $S(X^{J(h)})$。随后使用 \autoref{subsubsec:rns-application-key-switching} 的 RNS 密钥切换，返回 $S(X)$。此步骤会增加密钥切换误差，BGV 的误差含有因子 $t$。产生的槽位置换取决于根的排列：CKKS 具有 $n/2$ 个独立复数槽位，而完全分裂的 BFV/BGV 打包在 $5$ 的幂作用下具有两个各含 $n/2$ 个标量槽位的轨道。RNS 不改变该置换。

\subsubsection{BFV 的密文与密文相乘}
\label{subsubsec:rns-cipher-mult-bfv}

\cite{rns-bfv} 第 4.2--4.5 节的 BEHZ RNS 乘法依次执行基扩展、张量乘法、近似缩放舍入、精确转换回原基，以及重线性化。辅助基和误差界是算法的一部分。HPS 构造~\cite{rns-bfv-hps} 第 2.2--2.4 节及第 4.5 节使用不同的转换与缩放技术，不能混淆这些变体。

下面给出本节原语的中心化剩余类版本。令 $b=\prod_jb_j$，使用冗余模数 $b_\alpha$ 进行末尾的精确转换，并使用\emph{另一个不同的}临时 Montgomery 模数 $\widetilde m$ 控制初始提升。所有基模数都必须满足所需互素条件。
\begin{enumerate}
\item \textbf{基扩展。} 对每个分量系数 $x$，将 $[\widetilde m x]_q$ 快速转换到组合辅助基 $bb_\alpha\widetilde m$，再以 $\widetilde m$ 为临时模数执行 \textsf{SmallMont}。这样得到模 $bb_\alpha$ 下的提升 $\widehat x=x+uq$，其中 $|u|\leq1$，同时保留原基 $q$ 的剩余类。辅助剩余类 $b_\alpha$ 留作后续精确转换，不能像 Montgomery 模数那样被删除。
\item \textbf{张量乘法。} 在组合基 $qbb_\alpha$ 中计算负循环整数多项式乘积
\[
D_0=\widehat B_1\widehat B_2,\quad
D_1=\widehat A_1\widehat B_2+\widehat A_2\widehat B_1,\quad
D_2=\widehat A_1\widehat A_2.
\]
选择辅助范围，使预期整数 $tD_j$ 能够无回绕表示，例如验证 $\|tD_j\|_\infty<qbb_\alpha/2$。系数卷积、$n$、$t$ 和分量提升的界都会影响这一条件；仅有 $b>\Delta$ 并不足够。
\item \textbf{近似缩放舍入。} 将各个 $D_j$ 乘以 $t$，并执行删除基 $q$ 的 \textsf{ModSwitch\textsubscript{RNS}}。输出在基 $bb_\alpha$ 中表示
\[
Z_j=\left\lfloor\frac{tD_j}{q}\right\rceil-u_j
=\frac{tD_j}{q}+\rho_j,\quad
\|u_j\|_\infty\leq U_k,\quad
\|\rho_j\|_\infty\leq U_k+\frac12.
\]
这是近似缩放舍入，不是精确舍入。
\item \textbf{精确返回基 $q$。} 对每个 $Z_j$，使用保留的 $b_\alpha$ 剩余类执行 \textsf{FastBConvEx}。每个系数都必须满足 \autoref{subsec:rns-fastbconvex} 中的界 $|\mu+u|<b_\alpha/2$。对 $tD_j$ 的辅助范围界，不能代替这一独立的纠正界。
\item \textbf{重线性化。} 按 $(A,B)$ 顺序计算 $(Z_1,Z_0)+\textsf{DecompMult}_{\mathrm{RNS}}(Z_2,S^2_{\langle S,\mathrm{RNS}\rangle})$，此步骤增加求值密钥误差 $E_{\mathrm{ks}}$。
\end{enumerate}

令 $H_i=\widehat B_i+\widehat A_iS$ 表示提升后的整数相位。输出相位满足
\[
F\equiv\frac tqH_1H_2+\rho_0+\rho_1S+\rho_2S^2+E_{\mathrm{ks}}\pmod q.
\]
不能将三个分量的舍入替换为对 $H_1H_2$ 的单次舍入。误差界必须包含依赖秘密密钥的项、输入误差、提升产生的倍数，以及 $t\Delta^2/q$ 与 $\Delta$ 的差异。例如，将中心化输出相位写为 $\Delta M_{\mathrm{out}}+E_{\mathrm{out}}$，一个充分的解码条件是
\[
\|\Delta M_{\mathrm{out}}+E_{\mathrm{out}}\|_\infty<q/2,\qquad
\left\|\frac tqE_{\mathrm{out}}+\left(\frac{t\Delta}{q}-1\right)M_{\mathrm{out}}\right\|_\infty<1/2.
\]
其中 $M_{\mathrm{out}}=[M_1M_2]_t$。这些是绝对误差条件；仅将误差与消息大小比较并不正确，消息系数为零时尤其如此。相应的完整参数分析见 \autoref{subsubsec:bfv-mult-cipher-summary} 和所引论文。


\subsubsection{CKKS 的密文与密文相乘}
\label{subsubsec:rns-cipher-mult-ckks}

直接在当前基中计算张量积，使用适当的 RNS 求值密钥重线性化，并在尺度与噪声预算要求时重缩放（\autoref{subsubsec:ckks-mult-cipher-summary}）。在使用近似解释前，乘积相位必须适合当前模数；仅凭层级编号不能保证这一点。

若输入尺度为 $\Delta_1,\Delta_2$，删除的模数为 $w$，则新尺度是 $\Delta_1\Delta_2/w$。单素数重缩放能精确实现分量的最近整数除法，但仍产生相位误差 $R_B+R_AS$。多因子的近似转换还可能引入其他有界误差。BFV 专用的因子 $t$、用于张量积的大范围辅助基，以及精确返回固定 BFV 模数的最后步骤，不属于这一 CKKS 流程。不过，混合密钥切换仍可能使用自己的辅助基。


\subsubsection{BGV 的密文与密文相乘}
\label{subsubsec:rns-cipher-mult-bgv}

张量乘法与 RNS 重线性化得到模 $t$ 下的同一明文乘积，误差仍被 $t$ 整除（\autoref{subsec:bgv-mult-cipher}）。求值密钥误差也必须是 $t$ 的倍数。如需模数切换，应使用 \autoref{subsec:rns-modswitch-bgv} 的 $\textsf{ModSwitch}_{\mathrm{RNS}}^{\mathrm{BGV}}$，或显式跟踪纠正因子的更一般同余保持变体，不能用普通 CKKS 重缩放代替。因子 $t$ 保持固定；正确性和有效的噪声降低取决于完整的输出相位界。


\subsubsection{BFV 自举}
\label{subsubsec:rns-bfv-bootstrapping}

RNS 是密文系数模数的表示方式，并不禁止明文模数为 $p^\varepsilon$。素数幂也可以作为一个 CRT 模数，只要它与其他基模数互素。无效的只是把重复因子 $p,p,\ldots,p$ 当作不同且互素的 CRT 坐标。

\cite{cryptoeprint:2022/1363} 第 3.2--3.4 节的 BFV 自举构造采用公开预处理
\[
c_i'=\left[\left\lfloor\frac{p^\varepsilon}{q}c_i\right\rceil\right]_{p^\varepsilon},
\]
随后使用自举密钥同态计算 $[c_0'+c_1'S]_{p^\varepsilon}$。这里 $(c_0,c_1)=(B,A)$，该同态计算的明文模数为 $p^\varepsilon$。线性变换与数位提取实现将明文模数降至 $p^r$ 的舍入解密映射，之后施加逆变换。方案层面的讨论见 \autoref{subsubsec:bfv-bootstrapping-summary}。

这种系数到槽位构造要求明文素数与分圆指标 $\mu$ 满足 $\gcd(p,\mu)=1$（所引论文第 2.3 节）。对于 $\mu=2n$ 且 $n$ 为二次幂的情形，应使用奇数 $p$。二进制数位提取规则本身不能补足该环所缺少的明文环 CRT 分解。

公开预处理可以从 RNS 重构密文系数并精确舍入，也可以使用精度足够且经过证明的 RNS 缩放舍入例程。此后，同态电路的密文算术可以为工作密文模数使用另一套 RNS 基。若采用近似基转换，必须先把其误差计入舍入解密界，才能声称正确性。仅有 BEHZ/HPS 乘法例程并不能证明新的自举构造，而关于因子不互素的错误论证也不能成为用未经证明的基扩展与重新解释链替换该预处理的理由。

对于理想化输入相位 $(q/p^r)M+E$，分量舍入贡献 $R_B+R_AS$，舍入解密误差包含 $(p^\varepsilon/q)E+R_B+R_AS$。若实际加密尺度为 $\lfloor q/p^r\rfloor$，在该表达式中必须把尺度差异计入 $E$。恢复 $M$ 要求明确验证该误差小于 $p^{\varepsilon-r}/2$，并且整个同态电路的中间密文都必须有效。

\subsubsection{CKKS 自举}
\label{subsubsec:rns-ckks-bootstrapping}

采用 \autoref{subsubsec:ckks-bootstrapping-summary} 的 CKKS 流程：模数提升、系数到槽位变换、近似模约减，以及槽位到系数变换。稀疏自举还需要 \autoref{subsubsec:ckks-bootstrapping-time-reduction} 中说明的部分求和操作。该基本 CKKS 构造没有单独显式计算秘密密钥内积的同态步骤，也没有普遍适用的末尾重新解释步骤。

若近似 RNS 提升使分量分别增加 $q_0U_A,q_0U_B$，提升后的相位为
\[
\Delta M+E+q_0(K+U_B+U_AS).
\]
逐系数的额外提升界会经过秘密多项式放大。即使 $U_A,U_B$ 的每个系数都属于 $\{-1,0,1\}$，相应相位倍数也未必如此。周期模约减电路的逼近区间必须覆盖全部溢出，误差预算必须包含线性变换、重缩放、密钥切换、多项式逼近，以及保留的输入误差。仅有周期性，并不能使有限的同态多项式成为精确模约减函数。


\subsubsection{BGV 自举}
\label{subsubsec:rns-bgv-bootstrapping}

遵循 \autoref{subsec:bgv-bootstrapping} 的 BGV 构造。首先进行有效的模数切换，使 $\hat q\equiv1\pmod{p^\varepsilon}$，再计算中心化公开系数 $A'=[p^{\varepsilon-r}A]_{\hat q}$、$B'=[p^{\varepsilon-r}B]_{\hat q}$。该特殊模数不必是当前 RNS 基的成员：精确公开重构配合同余保持的模数切换，是一种有效实现途径。

使用明文模数 $p^\varepsilon$、足够大工作密文模数 $Q$ 下的自举密钥，同态计算 $B'+A'S$。RNS 表示的是 $Q$，不要求明文模数等于某个 RNS 因子。随后实现适当的系数与槽位变换，以及采用平衡舍入的数位提取，同时验证 $K'$ 和所有同态中间相位的界。

数位提取期间，每次具有正确性依据的对 $p$ 精确除法，都按 $p^j\rightarrow p^{j-1}$ 改变明文模数；经过 $\varepsilon-r$ 次后，明文模数变为 $p^r$。不能再附加一次乘以 $p^{-(\varepsilon-r)}$ 的操作，否则会对消息重复相除。此外，当 $\hat q$ 不整除 $Q$ 时，基扩展原语 $q\rightarrow qb$ 不能直接描述任意的 $\hat q\rightarrow Q$ 转换。采用这类提升的其他自举构造需要独立的正确性证明，仅有 RNS 表示并不能建立该证明。


\subsubsection{RNS 运算对噪声的影响}
\label{subsubsec:rns-noise}

应跟踪每个原语的实际作用。\textsf{FastBConv} 选择与原值相差源模数有界倍数的提升；\textsf{SmallMont} 限制该倍数，但通常不将其消除。\textsf{ModSwitch\textsubscript{RNS}} 在除法后将基转换偏差变为有界加性误差。\textsf{FastBConvEx} 在范围界成立时是精确的，不需要 \textsf{SmallMont} 来去除误差。\textsf{ModDrop\textsubscript{RNS}} 是精确的剩余类约减，但新相位或尺度条件不成立时，可能使密文失效。

对于分量扰动 $\delta_A,\delta_B$，引起的相位扰动为 $\delta_B+\delta_AS$，其系数界是 $\|\delta_B\|_\infty+\|\delta_A\|_\infty\|S\|_1$。高阶密文还会相应涉及 $S$ 的幂。这些扰动最终被抵消、近似消除，还是保留为解密误差，取决于使用它们的方案和运算，不能用单一的笼统噪声结论涵盖全部 RNS 原语。

\subsubsection{RNS 基本操作的 Python 源代码}
\label{subsubsec:rns-source-code}

我们提供了一个 \href{https://github.com/wang2000star/fhe-textbook-cn/blob/chinese-edition/source%20code/rns.py}{\underline{Python 脚本}}，实现以下 RNS 基本操作示例：\textsf{FastBConv}、
\textsf{ModRaise\textsubscript{RNS}}、
\textsf{ModDrop\textsubscript{RNS}}、
\textsf{ModSwitch\textsubscript{RNS}}、
\textsf{SmallMont} 和
\textsf{FastBConvEx}。







\clearpage

\section{FHE 方案比较与总结}
%\subsection{Scheme Comparison}
%\label{subsec:scheme-comparison}

下面对 TFHE、BFV、CKKS 和 BGV 进行总结与比较：

\begin{comment}
\clearpage



\end{comment}

\captionsetup{justification=raggedright,singlelinecheck=false,labelfont=bf}

\begin{table}[h]
\begin{tabular}{|c||l|}
\hline
&\textbf{困难问题基础}\\\hline\hline
\textbf{TFHE}&LWE\\\hline
\textbf{BFV}&\\
\textbf{CKKS}&RLWE\\
\textbf{BGV}&\\\hline
\end{tabular}
\caption{\textbf{困难问题基础}}
\end{table}

\begin{table}[h]
\begin{tabular}{|c||l|}
\hline
&\textbf{基本数据类型}\\\hline\hline
\textbf{TFHE}&向量\\\hline
\textbf{BFV}&\\
\textbf{CKKS}&多项式\\
\textbf{BGV}&\\\hline
\end{tabular}
\caption{\textbf{基本数据类型}}
\end{table}

\begin{table}[h]
\begin{tabular}{|c||l|}
\hline
&\multicolumn{1}{c|}{\textbf{明文}}\\\hline\hline
\textbf{TFHE}&数值 $m \in \mathbb{Z}_t$ \text{ } \textcolor{red}{ $\rhd$ $t$ 为 2 的幂}\\\hline
\textbf{BFV}&多项式 $M \in \mathbb{Z}_t[X]/X^n+1$ \text{ } \textcolor{red}{ $\rhd$ $t$ 为素数，$n$ 为 2 的幂}\\\hline
\textbf{CKKS}&多项式 $M \in \mathbb{R}[X]/X^n+1$ \text{ }\textcolor{red}{ $\rhd$ $n$ 为 2 的幂}\\\hline
\textbf{BGV}&多项式 $M \in \mathbb{Z}_t[X]/X^n+1$ \text{ } \textcolor{red}{ $\rhd$ $t$ 为素数，$n$ 为 2 的幂}\\\hline
\end{tabular}
\caption{\textbf{明文}}
\end{table}

\begin{table}[h]
\begin{tabular}{|c||l|}
\hline
&\multicolumn{1}{c|}{\textbf{私钥}}\\\hline\hline
\textbf{TFHE}&向量 $\vec{s} \xleftarrow{\$} \mathbb{Z}_2^k$，由 $\{0, 1\}$ 采样 \textcolor{red}{ $\rhd$ $\$$ 为均匀随机分布}\\\hline
\textbf{BFV}&\\
\textbf{CKKS}&多项式 $S \xleftarrow{\$} \mathbb{Z}_3[X]/X^n+1$，其中 $\mathbb{Z}_3 = \{-1, 0, 1\}$\\
\textbf{BGV}&\\\hline
\end{tabular}
\caption{\textbf{私钥}}
\end{table}

\begin{table}[h]
\begin{tabular}{|c||l|}
\hline
&\multicolumn{1}{c|}{\textbf{密文}}\\\hline\hline
\textbf{TFHE}&（向量 $\vec{a}$，数值 $b$）= $(\vec{a} \xleftarrow{\$} \mathbb{Z}_q^k$，$\text{ } b \in \mathbb{Z}_q)$ \text{ } \textcolor{red}{ $\rhd$ $q \gg t$，且 $t$ 整除 $q$}\\\hline
\textbf{BFV}&\\
\textbf{CKKS}&（多项式 $A, B$）= $(A \xleftarrow{\$} \mathbb{Z}_q[X]/X^n+1$，\text{ } $B \in \mathbb{Z}_q[X]/X^n+1)$ \text{ } \textcolor{red}{ $\rhd$ $q \gg t$}\\
\textbf{BGV}&\\\hline
\end{tabular}
\caption{\textbf{密文}}
\end{table}


\begin{table}[h]
\begin{tabular}{|c||l|}
\hline
&\multicolumn{1}{c|}{\textbf{噪声}}\\\hline\hline
\textbf{TFHE}&数值 $e \xleftarrow{\chi} \mathbb{Z}_q$ \text{ } \textcolor{red}{ $\rhd$ $\chi$ 为高斯随机分布}\\\hline
\textbf{BFV}&\\
\textbf{CKKS}&多项式 $E \xleftarrow{\chi} \mathbb{Z}_q[X]/X^n+1$\\
\textbf{BGV}&\\\hline
\end{tabular}
\caption{\textbf{噪声}}
\end{table}


\begin{table}[h]
\begin{tabular}{|c||l|}
\hline
&\multicolumn{1}{c|}{\textbf{尺度因子}}\\\hline\hline
\textbf{TFHE}&用于 $\Delta m$，\text{ } 其中 $\Delta = \dfrac{q}{t}$ \text{ } \textcolor{red}{ $\rhd$ $t$ 整除 $q$}\\\hline
\textbf{BFV}&用于 $\Delta M$，\text{ } 其中 $\Delta = \left\lfloor\dfrac{q}{t}\right\rfloor$ \text{ } \textcolor{red}{ $\rhd$ $t$ 为素数}\\\hline
\textbf{CKKS}&用于 $\Delta M$，\text{ } 其中 $\Delta \cdot ||M||_\infty \ll q_0$\\
& \text{ } \textcolor{red}{ $\rhd$ $q_0$ 为最低乘法层级的密文模数}\\\hline
\textbf{BGV}&用于 $\Delta E$，\text{ } 其中 $\Delta = t$ \text{ } \textcolor{red}{ $\rhd$ $t$ 为素数}\\\hline
\end{tabular}
\caption{\textbf{尺度因子}}
\end{table}


\begin{table}[h]
\begin{tabular}{|c||l|}
\hline
&\multicolumn{1}{c|}{\textbf{加密}}\\\hline\hline
\textbf{TFHE}&$(\vec{a}, b)$ \text{ } 其中 $\vec{a} \xleftarrow{\$} \mathbb{Z}_q^k$、$\text{ } b = \Delta m + e - a\cdot s \bmod q$，\text{ } $e \xleftarrow{\chi} \mathbb{Z}_q$\\
&\text{ } \text{ } \textcolor{red}{ $\rhd$ 每次使用 $e$ 后即丢弃}\\\hline
\textbf{BFV}&$(A, B)$ \text{ } 其中 $A \xleftarrow{\$} \mathbb{Z}_q[X] / (X^n + 1)$、$\text{ } B = \Delta M + E - A\cdot S \bmod q$，\text{ } $E \xleftarrow{\chi} \mathbb{Z}_q[X]/(X^n+1)$\\
\textbf{CKKS}&\text{ } \text{ }\textcolor{red}{ $\rhd$ 每次使用 $E$ 后即丢弃}\\\hline
\textbf{BGV}&$(A, B)$ \text{ } 其中 $A \xleftarrow{\$} \mathbb{Z}_q[X] / (X^n + 1)$、$\text{ } B = M + \Delta E - A\cdot S \bmod q$，\text{ } $E \xleftarrow{\chi} \mathbb{Z}_q[X]/(X^n+1)$\\
&\text{ } \text{ }\textcolor{red}{ $\rhd$ 每次使用 $E$ 后即丢弃}\\\hline
\end{tabular}
\caption{\textbf{加密}}
\end{table}

\begin{table}[h]
\begin{tabular}{|c||l|}
\hline
&\multicolumn{1}{c|}{\textbf{加密关系}}\\\hline\hline
\textbf{TFHE}&$b + a\cdot s = \Delta m + e \pmod q$，\text{ } 其中 $\Delta = \dfrac{q}{t}$ \text{ } \textcolor{red}{ $\rhd$ $t$ 整除 $q$}\\\hline
\textbf{BFV}&$B + A\cdot S = \Delta M + E \pmod q$，\text{ } 其中 $\Delta = \left\lfloor\dfrac{q}{t}\right\rfloor$ \text{ } \textcolor{red}{ $\rhd$ $t$ 为素数}\\\hline
\textbf{CKKS}&$B + A\cdot S = \Delta M + E \pmod q$，\text{ } 其中 $\Delta \cdot ||M||_\infty \ll q_0$\\
& \text{ } \textcolor{red}{ $\rhd$ $q_0$ 为最低乘法层级的密文模数}\\\hline
\textbf{BGV}&$B + A\cdot S = M + \Delta E \pmod q$，\text{ } 其中 $\Delta = t$ \text{ } \textcolor{red}{ $\rhd$ $t$ 为素数}\\\hline
\end{tabular}
\caption{\textbf{加密关系}}
\end{table}


\begin{table}[h]
\begin{tabular}{|c||l|}
\hline
&\multicolumn{1}{c|}{\textbf{解密公式}}\\\hline\hline
\textbf{TFHE}&$m = \left\lceil\dfrac{(b + a\cdot s \bmod q)}{\Delta}\right\rfloor \bmod t$\\
&$\textcolor{white}{m} = \left\lceil\dfrac{(\Delta m + e)}{\Delta}\right\rfloor \bmod t$ \text{ } \textcolor{red}{ $\rhd$ 若 $e < \dfrac{\Delta}{2}$，则 $e$ 被消去}\\\hline
\textbf{BFV}&$M = \left\lceil\dfrac{(B + A\cdot S \bmod q)}{\Delta}\right\rfloor  \bmod t$\\
&$\textcolor{white}{M} = \left\lceil\dfrac{(\Delta M + E)}{\Delta}\right\rfloor  \bmod t$ \text{ } \textcolor{red}{ $\rhd$ 若 $||E||_\infty < \dfrac{\Delta}{2}$，则 $E$ 被消去}\\\hline
\textbf{CKKS}&$M = \left\lceil\dfrac{(B + A\cdot S \bmod q)}{\Delta}\right\rfloor_{\frac{1}{\Delta}}$\\
&$\textcolor{white}{M} = \left\lceil\dfrac{\Delta M + E}{\Delta}\right\rfloor_{\frac{1}{\Delta}}$ \text{ } \textcolor{red}{ $\rhd$ 最终残余噪声为 $\dfrac{E}{\Delta}$（增大 $\Delta$ 可减小它）}\\\hline
\textbf{BGV}&$M = (B + A\cdot S \bmod q) \bmod t$\\
&$\textcolor{white}{M} = (M + \Delta E) \bmod t$ \text{ } \textcolor{red}{ $\rhd$ 若 $\Delta E < q$，则 $E$ 被消去}\\\hline
\end{tabular}
\caption{\textbf{解密公式}}
\end{table}

\begin{table}[h]
\begin{tabular}{|c||l|}
\hline
&\multicolumn{1}{c|}{\textbf{密文模数}}\\\hline\hline
\textbf{TFHE}&单个数值 $q$ \text{ } \textcolor{red}{ $\rhd$ $q \gg t$，且 $t$ 整除 $q$}\\\hline
\textbf{BFV}&单个数值 $q$ \text{ } \textcolor{red}{ $\rhd$ $q \gg t$}\\\hline
\textbf{CKKS}&具有 $L$ 个乘法层级的模数链 $\{q_0, q_1, \cdots, q_L\}$\\
& \text{ } \textcolor{red}{ $\rhd$ 各层满足 $q_i = \prod\limits_{j=0}^{l}w_j$，每个 $w_j$ 为 CRT 模数}\\
&\text{ } \text{ } \text{ } \textcolor{red}{满足 $w_0 \gg \Delta \cdot ||M||_{\infty}, \text{ } w_j \approx \Delta$（对于 $1 \leq j \leq L$）}\\\hline
\textbf{BGV}&具有 $L$ 个乘法层级的模数链 $\{q_0, q_1, \cdots, q_L\}$\\
&\text{ } \textcolor{red}{ $\rhd$ 各层满足 $q_i = \prod\limits_{j=0}^{l}w_j$，每个 $w_j$ 为 CRT 模数}\\
&\text{ } \text{ } \text{ } \textcolor{red}{满足 $w_0 \equiv w_1 \equiv \cdots \equiv w_L \bmod t$}\\\hline
\end{tabular}
\caption{\textbf{密文模数}}
\end{table}



\begin{table}[h]
\begin{tabular}{|c||l|}
\hline
&\multicolumn{1}{c|}{\textbf{密文与密文加法}}\\\hline\hline
\textbf{TFHE}&- 密文 $\textsf{LWE}_{\vec{s}, \sigma}(\Delta m_1) = (\vec{a}_1, b_1) = (a_{1,0}, a_{1, 1}, \cdots, a_{1, k-1}, b_1) \bmod q$\\
&- 密文 $\textsf{LWE}_{\vec{s}, \sigma}(\Delta m_2) = (\vec{a}_2, b_2) = (a_{2,0}, a_{2, 1}, \cdots, a_{2, k-1}, b_2) \bmod q$\\
&$\textsf{LWE}_{\vec{s}, \sigma}(\Delta (m_1 + m_2)) = (\vec{a}_1+\vec{a}_2, b_1+b_2) \bmod q$\\\hline
\textbf{BFV}&- 密文 $\textsf{RLWE}_{S, \sigma}(\Delta M_1) = (A_1, B_1) \bmod q$\\
&- 密文 $\textsf{RLWE}_{S, \sigma}(\Delta M_2) = (A_2, B_2) \bmod q$\\
&$\textsf{RLWE}_{S, \sigma}(\Delta (M_1 + M_2)) =(A_1+A_2,B_1+B_2) \bmod q$\\\hline
\textbf{CKKS}&- 密文 $\textsf{RLWE}_{S, \sigma}(\Delta M_1) = (A_1, B_1) \bmod q_l$\\
&- 密文 $\textsf{RLWE}_{S, \sigma}(\Delta M_2) = (A_2, B_2) \bmod q_l$\\
&$\textsf{RLWE}_{S, \sigma}(\Delta (M_1 + M_2)) =(A_1+A_2,B_1+B_2) \bmod q_l$\\\hline
\textbf{BGV}&- 密文 $\textsf{RLWE}_{S, \sigma}(M_1) = (A_1, B_1) \bmod q_l$\\
&- 密文 $\textsf{RLWE}_{S, \sigma}(M_2) = (A_2, B_2) \bmod q_l$\\
&$\textsf{RLWE}_{S, \sigma}(M_1 + M_2) =(A_1+A_2,B_1+B_2) \bmod q_l$\\\hline
\end{tabular}
\caption{\textbf{密文与密文加法}}
\end{table}

\begin{table}[h]
\begin{tabular}{|c||l|}
\hline
&\multicolumn{1}{c|}{\textbf{密文与明文加法}}\\\hline\hline
\textbf{TFHE}&- 密文 $\textsf{LWE}_{\vec{s}, \sigma}(\Delta m_1) = (\vec{a}_1, b_1) = (a_{1,0}, a_{1, 1}, \cdots, a_{1, k-1}, b_1) \bmod q$\\
&- 明文数值 $c \in \mathbb{Z}_t$\\
&$\textsf{LWE}_{\vec{s}, \sigma}(\Delta (m_1 + c)) = (\vec{a}_1, b_1+\Delta c) \bmod q$\\\hline
\textbf{BFV}&- 密文 $\textsf{RLWE}_{S, \sigma}(\Delta M_1) = (A_1, B_1) \bmod q$\\
&- 明文多项式 $C \in \mathbb{Z}_t[X]/(X^n + 1)$\\
&$\textsf{RLWE}_{S, \sigma}(\Delta (M_1 + C)) =(A_1,B_1+\Delta C) \bmod q$\\\hline
\textbf{CKKS}&- 密文 $\textsf{RLWE}_{S, \sigma}(\Delta M_1) = (A_1, B_1)  \bmod q_l$\\
&- 明文多项式 $C \in \mathbb{R}[X]/(X^n + 1)$\\
&$\textsf{RLWE}_{S, \sigma}(\Delta (M_1 + C)) =(A_1,B_1+\Delta C) \bmod q_l$\\\hline
\textbf{BGV}&- 密文 $\textsf{RLWE}_{S, \sigma}(M_1) = (A_1, B_1) \bmod q_l$\\
&- 明文多项式 $C \in \mathbb{Z}_t[X]/(X^n + 1)$\\
&$\textsf{RLWE}_{S, \sigma}(M_1 + C) =(A_1,B_1+C) \bmod q_l$\\\hline
\end{tabular}
\caption{\textbf{密文与明文加法}}
\end{table}

\begin{table}[h]
\begin{tabular}{|c||l|}
\hline
&\multicolumn{1}{c|}{\textbf{密文与明文乘法}}\\\hline\hline
\textbf{TFHE}&- 密文 $\textsf{LWE}_{\vec{s}, \sigma}(\Delta m_1) = (\vec{a}_1, b_1) = (a_{1,0}, a_{1, 1}, \cdots, a_{1, k-1}, b_1) \bmod q$\\
&- 明文数值 $c \in \mathbb{Z}_t$\\
&$\textsf{LWE}_{\vec{s}, \sigma}(\Delta (m_1 \cdot c)) = (\vec{a}_1 \cdot c, \text{ } b_1 \cdot c) \bmod q$\\\hline
\textbf{BFV}&- 密文 $\textsf{RLWE}_{S, \sigma}(\Delta M_1) = (A_1, B_1) \bmod q$\\
&- 明文多项式 $C \in \mathbb{Z}_t[X]/(X^n + 1)$\\
&$\textsf{RLWE}_{S, \sigma}(\Delta (M_1 \cdot C)) =(A_1 \cdot C, \text{ } B_1 \cdot C)$\\\hline
\textbf{CKKS}&- 密文 $\textsf{RLWE}_{S, \sigma}(\Delta M_1) = (A_1, B_1)  \bmod q_l$\\
&- 明文多项式 $C \in \mathbb{R}[X]/(X^n + 1)$\\
&1. \underline{基本乘法}\\
&\textcolor{white}{1. } $\textsf{ct} = \textsf{RLWE}_{S, \sigma}(\Delta^2 (M_1 \cdot C)) =(A_1 \cdot \Delta C, \text{ } B_1 \cdot \Delta  C)  \bmod q_l$\\
&2. 按 $\dfrac{1}{\Delta}$ \underline{重缩放}：$\left\lceil\dfrac{\textsf{ct}}{\Delta}\right\rfloor = \textsf{RLWE}_{S, \sigma}(\Delta M_1 C) \bmod q_{l-1}$\\\hline
\textbf{BGV}&- 密文 $\textsf{RLWE}_{S, \sigma}(M_1) = (A_1, B_1) \bmod q_l$\\
&- 明文多项式 $C \in \mathbb{Z}_t[X]/(X^n + 1)$\\
&$\textsf{RLWE}_{S, \sigma}(M_1 \cdot C) =(A_1 \cdot C, \text{ } B_1\cdot C)  \bmod q_l$\\\hline
\end{tabular}
\caption{\textbf{密文与明文乘法}}
\end{table}



\begin{table}[h]
\begin{tabular}{|c||l|}
\hline
&\multicolumn{1}{c|}{\textbf{密文与密文乘法}}\\\hline\hline
\textbf{TFHE}&- 密文 $\textsf{LWE}_{\vec{s}, \sigma}(\Delta m_1) = (\vec{a}_1, b_1) = (a_{1,0}, a_{1, 1}, \cdots, a_{1, k-1}, b_1) \bmod q$\\
&- 密文 $\textsf{LWE}_{\vec{s}, \sigma}(\Delta m_2) = (\vec{a}_2, b_2) = (a_{2,0}, a_{2, 1}, \cdots, a_{2, k-1}, b_2)  \bmod q$\\
&1. \underline{可编程自举}:\\ 
&\textcolor{white}{1. } 将 $\textsf{LWE}_{\vec{s}, \sigma}(\Delta m_2)$ 转换为 $\textsf{GSW}_{\vec{s}, \sigma}^{\beta, l}(m_2)$。\\
&2. \underline{同态乘法}:\\
&\textcolor{white}{2. } 计算 $\textsf{LWE}_{\vec{s}, \sigma}(\Delta m_1) \cdot \textsf{GSW}_{\vec{s}, \sigma}^{\beta, l}(m_2)$\\
&\text{ }$ = \sum\limits_{i=0}^{k-1} \langle \textsf{Decomp}^{\beta, l}(a_{1,i}),  \textsf{Lev}_{\vec{s}, \sigma}^{\beta, l}(-s_i\cdot m_2) \rangle + \langle \textsf{Decomp}^{\beta, l}(b_{1}),  \textsf{Lev}_{\vec{s}, \sigma}^{\beta, l}(m_2) \rangle$ \\
&\text{ } $ = \textsf{LWE}_{\vec{s}, \sigma}(\Delta m_1 m_2)$\\\hline
\textbf{BFV}&- 密文 $\textsf{RLWE}_{S, \sigma}(\Delta M_1) = (A_1, B_1)  \bmod q$\\
&- 密文 $\textsf{RLWE}_{S, \sigma}(\Delta M_2) = (A_2, B_2)  \bmod q$\\
&1. \underline{\textsf{ModRaise}}：从 $q$ 升至 $Q = q \cdot \Delta$\\
&\textcolor{white}{1. } - 密文 $\textsf{RLWE}_{S, \sigma}(\Delta M_1) = (A_1, B_1)  \bmod Q$\\
&\textcolor{white}{1. } - 密文 $\textsf{RLWE}_{S, \sigma}(\Delta M_2) = (A_2, B_2)  \bmod Q$\\
&2. \underline{多项式乘法}: \\
& \textcolor{white}{1. } $(A_1A_2, \text{ } A_1B_2 + A_2B_1, \text{ } B_1B_2) \equiv (D_2, D_1, D_0) \pmod Q$\\
&3. \underline{重线性化}: $\textsf{ct}_\alpha = (D
_1, D_0),$\text{ } $\textsf{ct}_\beta = \bm{\langle} \textsf{Decomp}^{\beta, l}(D_2), \text{ } \textsf{RLev}_{S, \sigma}^{\beta, l}( S^2) \bm{\rangle}$\\
&\textcolor{white}{3. 重线性化: } $\textsf{ct}_\alpha + \textsf{ct}_\beta = \textsf{ct}_{\alpha + \beta} = \textsf{RLWE}_{S, \sigma}(\Delta^2 M_1 M_2) \bmod Q$\\
&4. 按 $\dfrac{1}{\Delta}$ \underline{重缩放}：$\left\lceil\dfrac{\textsf{ct}_{\alpha+\beta}}{\Delta}\right\rfloor = \textsf{RLWE}_{S, \sigma}(\Delta M_1 M_2) \bmod q$\\\hline
\textbf{CKKS}&- 密文 $\textsf{RLWE}_{S, \sigma}(\Delta M_1) = (A_1, B_1)  \bmod q_l$\\
&- 密文 $\textsf{RLWE}_{S, \sigma}(\Delta M_2) = (A_2, B_2)  \bmod q_l$\\
&1. \underline{多项式乘法}: \\
& \textcolor{white}{1. } $(A_1A_2, \text{ } A_1B_2 + A_2B_1, \text{ } B_1B_2) \equiv (D_2, D_1, D_0) \pmod{q_l}$\\
&2. \underline{重线性化}: $\textsf{ct}_\alpha = (D
_1, D_0),$\text{ } $\textsf{ct}_\beta = \bm{\langle} \textsf{Decomp}^{\beta, l}(D_2), \text{ } \textsf{RLev}_{S, \sigma}^{\beta, l}( S^2) \bm{\rangle}$\\
&\textcolor{white}{3. 重线性化: } $\textsf{ct}_\alpha + \textsf{ct}_\beta = \textsf{ct}_{\alpha + \beta} = \textsf{RLWE}_{S, \sigma}(\Delta^2 M_1 M_2) \bmod q_l$\\
&3. 按 $\dfrac{1}{\Delta}$ \underline{重缩放}：$\left\lceil\dfrac{\textsf{ct}_{\alpha+\beta}}{\Delta}\right\rfloor = \textsf{RLWE}_{S, \sigma}(\Delta M_1 M_2) \bmod q_{l-1}$\\\hline
\textbf{BGV}&- 密文 $\textsf{RLWE}_{S, \sigma}(M_1) = (A_1, B_1)  \bmod q_l$\\
&- 密文 $\textsf{RLWE}_{S, \sigma}( M_2) = (A_2, B_2)  \bmod q_l$\\
&1. \underline{多项式乘法}: \\
& \textcolor{white}{1. } $(A_1A_2, \text{ } A_1B_2 + A_2B_1, \text{ } B_1B_2) \equiv (D_2, D_1, D_0) \pmod{q_l}$\\
&2. \underline{重线性化}: $\textsf{ct}_\alpha = (D
_1, D_0),$\text{ } $\textsf{ct}_\beta = \bm{\langle} \textsf{Decomp}^{\beta, l}(D_2), \text{ } \textsf{RLev}_{S, \sigma}^{\beta, l}( S^2) \bm{\rangle}$\\
&\textcolor{white}{3. 重线性化: } $\textsf{ct}_\alpha + \textsf{ct}_\beta = \textsf{ct}_{\alpha + \beta} = \textsf{RLWE}_{S, \sigma}( M_1 M_2) \bmod q_l$\\
&3. （可选）按 $\dfrac{1}{\Delta}$ \underline{重缩放}：$\left\lceil\dfrac{\textsf{ct}_{\alpha+\beta}}{\Delta}\right\rfloor_t = \textsf{RLWE}_{S, \sigma}(M_1 M_2) \bmod q_{l-1}$\\
&\textcolor{red}{ $\rhd$ $\lceil\rfloor_t$ 表示舍入到 $t$ 的最近倍数}\\
&\textcolor{red}{ $\rhd$ 重缩放可降低之后的噪声增长速度}\\\hline
\end{tabular}
\caption{\textbf{密文与密文乘法}}
\end{table}

\begin{table}[h]
\begin{tabular}{|c||l|}
\hline
&\multicolumn{1}{c|}{\textbf{最大乘法次数（不执行自举）}}\\\hline\hline
\textbf{TFHE}&使用可编程自举时不受限制（不使用自举则无法实现）\\\hline
\textbf{BFV}&不受限制（原文表述；请参见表后的译者注）\\\hline
\textbf{CKKS}&次数受模数链长度限制\\\hline
\textbf{BGV}&次数受模数链长度限制\\\hline
\end{tabular}
\caption{\textbf{最大乘法次数（不执行自举）}}
\end{table}
\noindent\textbf{译者注：}上表沿用原文的比较口径。BFV 在固定参数且不进行自举时仍受噪声预算约束；“不受限制”不能理解为可以正确执行任意深度的计算。关于 BGV 噪声管理等表述的进一步核对事项，见仓库的 \texttt{TRANSLATION.md}。


\begin{table}[h]
\begin{tabular}{|c||l|}
\hline
&\multicolumn{1}{c|}{\textbf{密钥切换}}\\\hline\hline
\textbf{TFHE}&密钥切换 $\vec{s} \rightarrow \vec{s}_{'}$：\\
&$\textsf{LWE}_{\vec{s}_{'},\sigma}(\Delta m) = b + a\cdot \textsf{LWE}_{\vec{s}_{'}, \sigma}(s)$\\
&\textcolor{white}{$\textsf{LWE}_{\vec{s}_{'},\sigma}(\Delta m) $} $ = b + \bm{\langle} \textsf{Decomp}^{\beta, l}(\vec{a}), \text{ } \textsf{Lev}_{\vec{s}_{'}, \sigma}^{\beta, l}(\vec{s}) \bm{\rangle}$\\\hline
\textbf{BFV}&密钥切换 $S \rightarrow S'$：\\
\textbf{CKKS}&$\textsf{RLWE}_{S',\sigma}(\Delta M) = B + A\cdot \textsf{RLWE}_{S', \sigma}(S)$\\
&\textcolor{white}{$\textsf{RLWE}_{S',\sigma}(\Delta M) $} $ = B + \bm{\langle} \textsf{Decomp}^{\beta, l}(A), \text{ } \textsf{RLev}_{S', \sigma}^{\beta, l}(S) \bm{\rangle}$\\\hline
\textbf{BGV}&密钥切换 $S \rightarrow S'$：\\
&$\textsf{RLWE}_{S',\sigma}(M) = B + A\cdot \textsf{RLWE}_{S', \sigma}(S)$\\
&\textcolor{white}{$\textsf{RLWE}_{S',\sigma}(M) $} $ = B + \bm{\langle} \textsf{Decomp}^{\beta, l}(A), \text{ } \textsf{RLev}_{S', \sigma}^{\beta, l}(S) \bm{\rangle}$\\\hline
\end{tabular}
\caption{\textbf{密钥切换}}
\end{table}


\begin{table}[h]
\begin{tabular}{|c||l|}
\hline
&\multicolumn{1}{c|}{\textbf{模数丢弃 (\textsf{ModDrop})}}\\\hline\hline
\textbf{CKKS}&- 乘法层级为 $l$ 的密文：$\textsf{RLWE}_{S, \sigma}(\Delta M) = (A, B)  \bmod q_l$\\
\textbf{BGV}&- 乘法层级为 $l-1$ 的密文：\\
&\textcolor{white}{ - } $\textsf{RLWE}_{S, \sigma}(\Delta M) = (A', B') = (A \bmod q_{l-1}, B \bmod q_{l-1})$\\\hline
\end{tabular}
\caption{\textbf{模数丢弃 (\textsf{ModDrop})}}
\end{table}



\begin{table}[h]
\begin{tabular}{|c||l|}
\hline
&\multicolumn{1}{c|}{\textbf{明文编码与解码}}\\\hline\hline
\textbf{TFHE}&不需要，因为每个明文是单个数值\\\hline
\textbf{BFV}&为支持批处理，须将输入槽位转换为多项式系数：\\
\textbf{CKKS}&- 将输入槽位 $\vec{v}$ 编码为多项式系数：$\vec{m} = n^{-1}\cdot\vec{v}\cdot I_n^R \cdot \hathat W$\\
\textbf{BGV}&- 将多项式系数 $\vec{m}$ 解码为输入槽位：$\vec{v} = \vec{m} \cdot \hathat W^{*}$\\\hline
\end{tabular}
\caption{\textbf{明文编码与解码}}
\end{table}


\begin{table}[h]
\begin{tabular}{|c||l|}
\hline
&\multicolumn{1}{c|}{\textbf{输入槽位旋转}}\\\hline\hline
\textbf{TFHE}&不适用，因为明文是单个数值（即单个槽位）\\\hline
\textbf{BFV}&给定 $\textsf{ct} = \textsf{RLWE}_{S(X), \sigma}\bm(\Delta M(X)\bm) = \bm(A(X), B(X)\bm)$，\\
&\text{ } 将输入槽位向左旋转 $h$ 个位置：\\
\textbf{CKKS}&1. 将 $\textsf{ct}$ 更新为 $ \textsf{RLWE}_{S(X^{J(h)}), \sigma}\bm(\Delta M(X^{J(h)})\bm) = \bm(A(X^{J(h)})$、$B(X^{J(h)})\bm)$\\
&\text{ } \text{ }（其中 $J(h) = 5^h \bmod 2n$）\\
&2. 密钥切换：将 $\textsf{RLWE}_{S(X^{J(h)}), \sigma}\bm(\Delta M(X^{J(h)})\bm)$ 转为 $\textsf{RLWE}_{S(X), \sigma}\bm(\Delta M(X^{J(h)})\bm)$。\\\hline
\textbf{BGV}&给定 $\textsf{ct} = \textsf{RLWE}_{S(X), \sigma}\bm(M(X)\bm) = \bm(A(X), B(X)\bm)$，\\
&\text{ } 将输入槽位向左旋转 $h$ 个位置：\\
&1. 将 $\textsf{ct}$ 更新为 $\textsf{RLWE}_{S(X^{J(h)}), \sigma}\bm(M(X^{J(h)})\bm) = \bm(A(X^{J(h)})$、$B(X^{J(h)})\bm)$\\
&2. 密钥切换：将 $\textsf{RLWE}_{S(X^{J(h)}), \sigma}\bm( M(X^{J(h)})\bm)$ 转为 $\textsf{RLWE}_{S(X), \sigma}\bm( M(X^{J(h)})\bm)$。\\\hline
\end{tabular}
\caption{\textbf{输入槽位旋转}}
\end{table}


\begin{table}[h]
\begin{tabular}{|c||l|}
\hline
&\multicolumn{1}{c|}{\textbf{自举目标}}\\\hline\hline
\textbf{TFHE}&重置噪声。\\
\textbf{BFV}&\\\hline
\textbf{CKKS}&按 $q_0 \rightarrow q_L$ 重置密文模数（严格地说，重置为满足 $q_0 < q_{l'} < q_L$ 的 $q_{l'}$）。\\
\textbf{BGV}&\\\hline
\end{tabular}
\caption{\textbf{自举目标}}
\end{table}

\begin{table}[h]
\begin{tabular}{|c||l|}
\hline
&\multicolumn{1}{c|}{\textbf{自举细节}}\\\hline\hline
\textbf{TFHE}&1. 按 $q \rightarrow 2n$ \underline{切换模数}，转换 $\textsf{LWE}_{\vec{s}, \sigma}(\Delta m) \rightarrow \textsf{LWE}_{\vec{s}, \sigma}(\hat\Delta m) = ( \vec{\hat a}, \hat b) \bmod 2n $，\\
&\text{ } \text{ } 其中 $\hat\Delta = \dfrac{2n}{t}$。\text{ } \textcolor{red}{ $\rhd$ $t$ 整除 $2n$}\\
&2. \underline{盲旋转}：同态地将 GLWE 加密的查找表\\
&\text{ } \text{ } 多项式 $\textsf{GLWE}_{\vec{S},\sigma}(\Delta V)$ 向左旋转 $\Delta m + e$ 个位置，\\
&\text{ } \text{ } 即按如下方式同态计算 $\textsf{GLWE}_{\vec{S}, \sigma}(\Delta V_k)$：\\
&\text{ } \text{ } $\textsf{GLWE}_{\vec{S}, \sigma}(\Delta V_0) = \textsf{GLWE}_{\vec{S},\sigma}(\Delta V) \cdot X^{-\hat b}$\\
&\text{ } \text{ } $\textsf{GLWE}_{\vec{S}, \sigma}(\Delta V_i) = \textsf{GLWE}_{\vec{S}, \sigma}(\Delta V_{i-1}) \cdot X^{\hat a_i s_{i-1}}$\\
&\text{ } \text{ } \text{ } $ = \textsf{GGSW}_{\vec{S},\sigma}^{\beta, l}(s_{i-1})\cdot \bm(\textsf{GLWE}_{\vec{S}, \sigma}(\Delta V_{i-1}) \cdot X^{\hat{a}_{i-1}} - \textsf{GLWE}_{\vec{S}, \sigma}(\Delta V_{i-1})\bm) + \textsf{GLWE}_{\vec{S}, \sigma}(\Delta V_{i-1})$\\
&3. \underline{系数提取}：旋转后的加密多项式 $V_k$ 的常数项为\\
&\text{ } \text{ } $\Delta m$。从 $\textsf{GLWE}_{\vec{S}, \sigma}(\Delta V_k)$ 中提取该加密常数项，得到 $\textsf{LWE}_{\vec{s}, \sigma}(\Delta m)$。\\\hline
\textbf{BFV}&1. 按 $q \rightarrow p^\varepsilon$ \underline{切换模数}，转换 $\textsf{RLWE}_{S, \sigma}(\Delta M) \rightarrow \textsf{RLWE}_{S, \sigma}(p^{\varepsilon-1} M) \bmod p^\varepsilon$\\
&2. \underline{同态解密}:\\
&\text{ } \text{ } $B + A\cdot \textsf{RLWE}_{S, \sigma}(\Delta' S) = \textsf{RLWE}_{S, \sigma}(\Delta' \cdot(p^{\varepsilon-1}M + E + Kp^\varepsilon)) \bmod q$，其中 $\Delta' = \left\lfloor\dfrac{q}{p^\varepsilon}\right\rfloor$\\
&3. \underline{\textsf{CoeffToSlot}}：将密文乘以 $n^{-1}\cdot I_n^R \cdot \hathat W$，\\
&\text{ } \text{ } 把 $p^{\varepsilon-1} M + E + Kp^\varepsilon$ 的明文系数移到输入槽位。\\
&4. \underline{数位提取}：给定数位提取多项式 $G_{\varepsilon}(x)$，\\
&\text{ } \text{ } 同态计算：\\
&\text{ } \text{ } $G_{2} \circ G_{3}\circ \cdots \circ G_{\varepsilon-1} (p^{\varepsilon-1} M + E + Kp^\varepsilon)$\\
&\text{ } \text{ } 将输出 $M + K^{\langle \varepsilon - 1 \rangle}p$ 存储在明文槽位中。\\
&\text{ } \text{ } 通过重新解释尺度因子处理乘以 $p$ 逆元的操作。\\
&5. \underline{\textsf{SlotToCoeff}}：将密文乘以 $\hathat W^*$，把 $M + K^{\langle \varepsilon - 1 \rangle}p$ 移回\\
&\text{ } \text{ } 多项式系数位置，得到 $\textsf{RLWE}_{S, \sigma}(\Delta'\cdot(M + K^{\langle \varepsilon - 1 \rangle}p)) \bmod q$\\
&\text{ } \text{ } $  = \textsf{RLWE}_{S, \sigma}(\Delta M) \bmod q$\\
\hline
\end{tabular}
\caption{\textbf{自举细节: TFHE, BFV}}
\end{table}

\begin{table}[h]
\begin{tabular}{|c||l|}
\hline
&\multicolumn{1}{c|}{\textbf{自举细节}}\\\hline\hline
\textbf{CKKS}&1. \underline{\textsf{ModRaise}：} 将密文 $(A, B) \bmod q_0$ 视为 $(A, B) \bmod q_L$\\
%&2. \underline{Homomorphic Decryption } (without scaling down by $\Delta$): \\
%&\text{ } \text{ } $B + A \cdot \textsf{RLWE}_{S, \sigma}(S) \bmod q_L = B + \langle \textsf{Decomp}^{\beta, l}(A), \textsf{RLev}_{S, \sigma}^{\beta, l}(S) \rangle$\\
%&$\text{ } \text{ } = \textsf{RLWE}_{S, \sigma}(\Delta M + E + Kq_0) \bmod q_L$\\
&2. \underline{\textsf{CoeffToSlot}}：将 $\Delta M + E + Kq_0$ 的系数移到输入槽位。\\
&3. \underline{\textsf{EvalExp}}：在 $\Delta M + E + Kq_0$ 处同态求值\\
&\text{ } \text{ } 周期为 $q_0$ 的正弦函数的多项式近似，\\
&\text{ } \text{ } 得到在明文槽位中加密 $\Delta M + E$ 的密文。\\
&4. \underline{\textsf{SlotToCoeff}}：将 $\Delta M + E$ 移到多项式系数位置，\\
&\text{ } \text{ } 得到 $\textsf{RLWE}_{S, \sigma}(\Delta M + E)$。\\\hline
\textbf{BGV}&1. 按 $q_{l} \rightarrow \hat{q}$ \underline{切换模数}，转换\\
&\text{ } \text{ } $\textsf{RLWE}_{S, \sigma}(M) = (A, B) \bmod q_l \rightarrow \textsf{RLWE}_{S, \sigma}(M) = (\hat{A}, \hat{B}) \bmod \hat{q}$\\ 
&\text{ } \text{ } 其中 $\hat{q} \equiv 1 \bmod p^\varepsilon$\\
&2. \underline{密文系数乘以 $p^{\varepsilon-1}$}：\\
&\text{ } \text{ } 计算 $(p^{\varepsilon-1}\hat{A}, p^{\varepsilon-1}\hat{B}) = (A', B') \bmod \hat{q}$（其中 $\hat{q} \equiv 1 \bmod p^\varepsilon$）\\
&\text{ } \text{ } 得到密文 $\textsf{RLWE}_{S, \sigma}(p^{\varepsilon-1}M) \bmod \hat{q}$，其噪声为 $p^\varepsilon E$。\\
&3. \underline{\textsf{ModRaise}}: $(A', B') \bmod \hat{q} \rightarrow (A', B') \bmod q_L$\\
&\text{ } \text{ } 这就是密文 $\textsf{RLWE}_{S, \sigma}(p^{\varepsilon-1}M + p^\varepsilon E + K'\hat{q}) \bmod q_L$。\\
&4. \underline{\textsf{CoeffToSlot}}：将密文乘以 $n^{-1}\cdot I_n^R \cdot \hathat W$，\\
&\text{ } \text{ } 把 $p^{\varepsilon-1}M + p^\varepsilon E + K'\hat{q}$ 的明文系数移到输入槽位。\\
&5. \underline{数位提取}：给定数位提取多项式 $G_{\varepsilon}(x)$，\\
&\text{ } \text{ } 同态计算：\\
&\text{ } \text{ } $\bm(\text{ } G_{2} \circ G_{3}\circ \cdots \circ G_{\varepsilon} (p^{\varepsilon-1} M + p^\varepsilon E + K'\hat{q})\text{ } \bm)$\\
&\text{ } \text{ } 将输出 $M + K''p$ 存储在明文槽位中。\\
&\text{ } \text{ } 为处理第 $i$ 轮中乘以 $p$ 逆元的操作，\\
&\text{ } \text{ } 将所有密文多项式系数乘以 $|p^{-1}|_{q^{\langle i \rangle}}$，使明文部分\\
&\text{ } \text{ } 从 $Mp^{\varepsilon - i} + K'''p^{\varepsilon - i + 1} \pmod{q^{\langle i \rangle}}$ 更新为 $M^{\varepsilon - i - 1} + K^{\langle i \rangle}p^{\varepsilon - i} \pmod{q^{\langle i \rangle}}$\\
&\text{ } \text{ } \textcolor{red}{ $\rhd$ 其中 $q^{\langle i \rangle}$ 为该轮对应的密文模数}\\
&6. \underline{\textsf{SlotToCoeff}}：将密文乘以 $\hathat W^*$，把 $M + K^{\langle \varepsilon - 1 \rangle}p$ 移回\\
&\text{ } \text{ } 多项式系数位置，得到 $\textsf{RLWE}_{S, \sigma}(M + K^{\langle \textit{last} \rangle}p)$\\
&\text{ } \text{ } $ = \textsf{RLWE}_{S, \sigma}(M + \Delta K^{\langle \textit{last} \rangle}) \bmod q_{l'}$ \textcolor{red}{ $\rhd$ 其中 $\Delta = p$，最终噪声为 $K^{\langle \textit{last} \rangle}$}\\\hline
\end{tabular}
\caption{\textbf{自举细节: CKKS, BGV}}
\end{table}

\begin{table}[h]
\begin{tabular}{|c||l|}
\hline
&\multicolumn{1}{c|}{\textbf{噪声管理}}\\\hline\hline
\textbf{TFHE}&通过自举重置噪声。\\
\textbf{BFV}&\\\hline
\textbf{CKKS}&- 噪声持续增长，因为自举仅重置\\
&模数链。为减缓噪声增长，应增大明文的\\
&尺度因子 $\Delta$。\\
&- CKKS 的 \textsf{EvalExp} 不能用数位提取多项式去除噪声，\\
&因为其明文是实数，而非模数环中的元素。\\\hline
\textbf{BGV}&BGV 的模数切换具有重置噪声的特殊性质，而其\\
&自举重置模数链，使模数切换能够持续进行。\\\hline
\end{tabular}
\caption{\textbf{噪声管理}}
\end{table}



\clearpage

\section{Python FHE 演示库}
\texttt{source code/} 中的 Python 文件用于教学演示。维护中的命令行示例为 \texttt{tfhe\_example.py}、\texttt{bfv\_example.py}、\texttt{ckks\_example.py} 和 \texttt{bgv\_example.py}，支持的操作见 README。上游源码位于 \url{https://github.com/fhetextbook/fhetextbook.github.io}，经审校的双语分支位于 \url{https://github.com/wang2000star/fhe-textbook-cn}。

应区分已实现的操作、功能模拟和未实现的函数。TFHE 的 \texttt{tfhe\_mock\_pbs} 会读取秘密密钥，解密、查询表格，再重新加密；其“自举密钥”字典直接包含该秘密密钥。因此，它是客户端功能模拟，不能当作公开求值密钥或安全的服务器端可编程自举或电路自举实现。默认密文乘法示例演示的是张量乘积及密钥切换构造，并未对标准 TFHE 自举进行性能测试。BFV、CKKS 和 BGV 的自举函数尚未实现，会明确抛出 \texttt{NotImplementedError}。

示例使用很小的教学维数、有界误差采样和 Python/NumPy 伪随机数来展示代数关系；这些并非经过验证的密码学参数集或生产用途的随机数生成器。示例运行成功仅验证所选输入，不能据此认定安全等级、任意深度正确性或优化 FHE 库的性能。较早的独立编码和 RNS 脚本属于代数实验，其覆盖范围与限制见学术审校报告。


%\input{scratch}
\clearpage

\bibliographystyle{unsrt}
\bibliography{z-bibfile}



%\clearpage
%\input{scratch-real}

\end{document}


%https://www.youtube.com/watch?v=vYKdh5oQ4Zw
